\documentclass[11pt]{report}
 
% --- GEOMETRIA E STRUTTURA DELLA PAGINA ---
\usepackage[a4paper, left=2.8cm, right=2.8cm, top=3cm, bottom=3cm, headheight=15pt, headsep=8mm]{geometry}
\usepackage{fancyhdr}
\usepackage{setspace}
\usepackage{lscape}
\usepackage{titlesec}
\usepackage{etoolbox}
\usepackage[section]{placeins}
 
% --- LINGUA, CODIFICA E TIPOGRAFIA (FONT) ---
\usepackage[italian, english]{babel}
\usepackage[utf8]{inputenc}
\usepackage[T1]{fontenc}
\usepackage{lmodern}
\usepackage{enumitem}
\usepackage{microtype}
 
% --- MATEMATICA E SIMBOLI ---
\usepackage{amsmath, amssymb, amsfonts}
\usepackage{amsthm}
\usepackage{bm}
\usepackage{siunitx}
\sisetup{per-mode=symbol}
 
% --- GRAFICA, DIAGRAMMI E CIRCUITI ---
\usepackage{graphicx}
\usepackage{caption}
\captionsetup[figure]{font=small,labelfont=bf,labelsep=period,justification=raggedright,singlelinecheck=false,position=bottom,skip=10pt}
\usepackage{tikz}
\usepackage{pgfplots}
\usepackage[american]{circuitikz}
 
\usetikzlibrary{shapes.geometric, arrows, arrows.meta, positioning, calc}
\usetikzlibrary{patterns, decorations.pathmorphing, backgrounds, automata}
\pgfplotsset{compat=1.18}
\usepgfplotslibrary{statistics}

\tikzset{
    block/.style = {draw, thick, rectangle, minimum height=1.2cm, minimum width=2cm, align=center},
    sum/.style = {draw, circle, thick, minimum size=8mm, node distance=1cm},
    input/.style = {coordinate},
    output/.style = {coordinate},
    arrow/.style={->, thick, black},
    io/.style = {trapezium, trapezium left angle=70, trapezium right angle=110, minimum width=1cm, minimum height=0.5cm, draw, align=center}
}
 
% --- TABELLE E INDICE ---
\usepackage{tabularx}
\usepackage{booktabs}
\usepackage{tocloft}
\usepackage[hidelinks,bookmarks=true,bookmarksopen=true]{hyperref}
\hypersetup{
  pdftitle={Mathematical Methods for Machine Learning},
  pdfauthor={Tommaso Vescio},
  pdfsubject={A self-contained compendium of mathematical methods for machine learning},
  pdfkeywords={mathematical analysis, linear algebra, optimization, probability, differential equations, machine learning},
  pdfcreator={LaTeX with hyperref}
}
\setcounter{tocdepth}{2}
\setcounter{secnumdepth}{2}
\setlength{\cftbeforechapskip}{1em}
\setlength{\cftbeforesecskip}{0.28em}
\setlength{\cftbeforesubsecskip}{0.10em}
\setlength{\cftchapnumwidth}{2.5em}
\cftsetindents{section}{1.5em}{3.1em}
\cftsetindents{subsection}{3.4em}{4em}
\renewcommand{\cftsecfont}{\bfseries}
\renewcommand{\cftsecpagefont}{\bfseries}
\renewcommand{\cftsubsecfont}{\normalfont}
\renewcommand{\cftsubsecpagefont}{\normalfont}
\renewcommand{\cftchapafterpnum}{\par\nobreak}
\renewcommand{\cftsecafterpnum}{\par\nobreak}
\cftsetpnumwidth{2.5em}
\cftsetrmarg{3.5em}
 
% --- CONFIGURAZIONE STILI TEOREMI ---
\newtheoremstyle{boldplain}
  {14pt plus 3pt minus 2pt}{14pt plus 3pt minus 2pt}{\itshape}{}{\bfseries}{.}{0.7em}
  {\thmname{#1}\thmnumber{ #2}\thmnote{ #3}}
\newtheoremstyle{bolddefinition}
  {14pt plus 3pt minus 2pt}{14pt plus 3pt minus 2pt}{\normalfont}{}{\bfseries}{.}{0.7em}
  {\thmname{#1}\thmnumber{ #2}\thmnote{ #3}}
\newtheoremstyle{boldremark}
  {14pt plus 3pt minus 2pt}{14pt plus 3pt minus 2pt}{\normalfont}{}{\bfseries}{.}{0.7em}
  {\thmname{#1}\thmnumber{ #2}\thmnote{ #3}}
 
\theoremstyle{bolddefinition}
\newtheorem{definition}{Definition}[section]
\newtheorem{example}{Example}[section]
\theoremstyle{boldplain}
\newtheorem{theorem}{Theorem}[section]
\newtheorem{lemma}{Lemma}[section]
\newtheorem{corollary}{Corollary}[section]
\newtheorem{proposition}{Proposition}[section]
\theoremstyle{boldremark}
\newtheorem{remark}{Remark}[section]

 
\setlength{\parindent}{0pt}
\setlength{\parskip}{0.7em plus 0.15em minus 0.1em}
\setstretch{1.0}
\setdisplayskipstretch{1}
\raggedbottom
\clubpenalty=10000
\widowpenalty=10000
\displaywidowpenalty=10000
\setlength{\emergencystretch}{1em}

\appto\normalsize{%
  \setlength{\abovedisplayskip}{16pt plus 3pt minus 2pt}%
  \setlength{\belowdisplayskip}{16pt plus 3pt minus 2pt}%
  \setlength{\abovedisplayshortskip}{10pt plus 2pt minus 1pt}%
  \setlength{\belowdisplayshortskip}{14pt plus 3pt minus 2pt}%
}
\AtBeginDocument{\normalsize}
\setlength{\jot}{6pt}

\titleformat{\chapter}[display]
  {\normalfont\Huge\bfseries\raggedright}
  {\normalfont\large\scshape\chaptertitlename\ \thechapter}
  {24pt}{}
\titlespacing*{\chapter}{0pt}{54pt}{44pt}
\titlespacing*{\section}{0pt}{2.2\baselineskip plus 4pt minus 2pt}{0.85\baselineskip}
\titlespacing*{\subsection}{0pt}{1.6\baselineskip plus 3pt minus 2pt}{0.65\baselineskip}

\setlist[itemize,enumerate,description]{
  topsep=0.8em plus 0.2em minus 0.1em,
  itemsep=0.45em plus 0.1em minus 0.05em,
  parsep=0.2em,
  partopsep=0.2em
}
\BeforeBeginEnvironment{proof}{\par\addvspace{0.6\baselineskip}}
\AfterEndEnvironment{proof}{\par\addvspace{0.6\baselineskip}}

\setlength{\textfloatsep}{22pt plus 4pt minus 2pt}
\setlength{\intextsep}{20pt plus 4pt minus 2pt}
\setlength{\floatsep}{16pt plus 3pt minus 2pt}
\setlength{\abovecaptionskip}{10pt}
\setlength{\belowcaptionskip}{6pt}
\renewcommand{\topfraction}{0.9}
\renewcommand{\bottomfraction}{0.8}
\renewcommand{\textfraction}{0.12}
\renewcommand{\floatpagefraction}{0.75}
\setcounter{topnumber}{3}
\setcounter{bottomnumber}{2}
\setcounter{totalnumber}{4}
 
% --- COMANDI AGGIUNTIVI PER ALGEBRA/ANALISI ---
\newcommand{\bv}{\bm{v}}
\newcommand{\bu}{\bm{u}}
\newcommand{\bx}{\bm{x}}
\newcommand{\by}{\bm{y}}
\newcommand{\bW}{\bm{W}}
\newcommand{\bM}{\bm{M}}
\newcommand{\bX}{\bm{X}}
\newcommand{\R}{\mathbb{R}}
\newcommand{\Z}{\mathbb{Z}}
\newcommand{\C}{\mathbb{C}}
\newcommand{\N}{\mathbb{N}}
\newcommand{\Prob}{\mathbb{P}}
\newcommand{\calH}{\mathcal{H}}
\newcommand{\calS}{\mathcal{S}}
\newcommand{\calF}{\mathcal{F}}
\newcommand{\calB}{\mathcal{B}}
\newcommand{\bepsilon}{\boldsymbol{\epsilon}}
\DeclareMathOperator{\diag}{diag}
\DeclareMathOperator{\tr}{tr}
\DeclareMathOperator{\var}{Var}
\DeclareMathOperator{\cov}{Cov}
\DeclareMathOperator{\E}{\mathbb{E}}
\newcommand{\dd}{\mathrm{d}}
\newcommand{\dalembert}{\mathop{\Box}\nolimits}

% --- PALETTE DEDICATA ALLE FIGURE (coerente con lo stile del libro) ---
\definecolor{figblue}{RGB}{0,90,170}
\definecolor{figred}{RGB}{190,30,45}
\definecolor{figgreen}{RGB}{30,140,80}
\tikzset{
  vecA/.style={draw=figblue, very thick, -{Stealth[length=2.8mm]}},
  vecB/.style={draw=figred, very thick, -{Stealth[length=2.8mm]}},
  vecC/.style={draw=figgreen, very thick, -{Stealth[length=2.8mm]}},
  gridline/.style={draw=black!30, thin},
}

% ==========================================
%            INIZIO DOCUMENTO
% ==========================================
\begin{document}

\pagenumbering{roman} 
\pagestyle{plain}     

% =========================================================================
% --- FRONTESPIZIO COMPATTO (STILE GATSBY UNIT UCL) ---
% =========================================================================
\begin{center}
    {\linespread{1.1}\selectfont
    {\huge Mathematical Methods for Machine Learning\\} 
    }
    \vspace{0.4cm}
    {\large{Tommaso Vescio $\,\cdot\,$ Updated: September 2026 (Original: June 2023)} \par}
\end{center}

\vspace{0.2cm}

\hrule
\vspace{0.15cm}
{\large{\textit{Mathematical Methods for Machine Learning} develops a unified pathway from the foundations of mathematical analysis to advanced methods for modeling and inference. Building on the Analysis course taught by \textit{Professor Monica Conti (Politecnico di Milano)}, this expanded compendium combines rigorous treatments of calculus, linear and functional analysis, differential equations, optimization, Fourier methods, and probability with the operational tools required for biophysical modeling and probabilistic inference in complex systems.} \par}

% --- INDICE ---
\renewcommand{\contentsname}{\LARGE\textbf{Contents}} 
\tableofcontents 
\newpage

% --- CONFIGURAZIONE PAGINE ---
\pagenumbering{arabic} 
\pagestyle{fancy}
\fancyhf{}              
\fancyhead[R]{\small\nouppercase{\rightmark}}
\fancyfoot[C]{\thepage}

% --- INCLUSIONE DEI FILE DELLE LEZIONI ---
\chapter{Real and Complex Numbers: Foundations of Analysis}

\section{Real Number Systems and the Real Line}

\subsection{Natural Numbers, Integers, and Rational Numbers}

The real number system, $\mathbb{R}$, can be approached through a sequence of progressively larger number systems, each containing the preceding one.
The natural numbers form the starting point:
\[
\mathbb{N} = \{0, 1, 2, 3, \dots\}.
\]
Including their negatives gives the integers:
\[
\mathbb{Z} = \{0, 1, -1, 2, -2, \dots\}.
\]
Allowing quotients of integers with nonzero denominators then gives the rational numbers, $\mathbb{Q}$:
\[
\mathbb{Q} = \left\{ \pm \frac{p}{q} : p \in \mathbb{N}, \ q \in \mathbb{N}, \ q \neq 0 \right\}.
\]

\subsection{Irrational, Algebraic, and Transcendental Numbers}

The rational numbers do not exhaust the real line. Real numbers that cannot be expressed as ratios of integers are called irrational numbers.
The following theorem provides a fundamental example.

\begin{theorem}[Irrationality of the Square Root of Two]
The square root of $2$ is irrational.
\end{theorem}

\begin{proof}
Suppose, for a contradiction, that $\sqrt{2}$ is rational. Then we can write $\sqrt{2} = p/q$, for integers $p, q$ with $q \neq 0$. Assume that $p/q$ is in lowest terms; in particular, $p$ and $q$ cannot both be even.
Squaring both sides gives $2 q^2 = p^2$. Thus $p^2$ is even. Since the square of an odd integer is odd, $p$ must also be even.
Writing $p = 2k$ for an integer $k$ and substituting gives
\[
\begin{aligned}
2q^2 &= 4k^2  \\
\Rrightarrow\quad  q^2 &= 2k^2.
\end{aligned}
\]

Hence $q^2$ is even, so $q$ is even as well. Both $p$ and $q$ are therefore even, contradicting the assumption that $p/q$ is in lowest terms.
This proves the theorem.
\end{proof}

The real numbers comprise both rational and irrational numbers. Familiar irrational numbers include $\sqrt{2}$, $\pi$, Euler's number $e$,
and the golden ratio $\varphi$. For two lengths $a$ and $b$ in that ratio, the defining proportion is
\[
\varphi = \frac{a+b}{a} = \frac{a}{b}.
\]

The numbers $\sqrt{2}$ and $\varphi$ also belong to the class of algebraic numbers. A real number $x$ is algebraic if it is a root of a
nonzero polynomial with integer coefficients $a_k \in \mathbb{Z}$, for $k = 0, \dots, n$:
\[
\sum_{k=0}^{n} a_k x^k = 0.
\]

Both $\sqrt{2}$ and $\varphi$ satisfy this definition. By contrast, $\pi$ is transcendental: it is not algebraic.
The historical results on $\pi$ established between 1794 and 1882 illustrate the distinction between irrationality and transcendence.
In particular, irrational algebraic numbers form a proper subclass of the irrational numbers.

\subsection{The Real Line and Interval Notation}

The real line represents $\mathbb{R}$ geometrically: each real number corresponds to a point, and each point corresponds to a real number.
Its absence of gaps reflects the completeness property discussed below. Intervals are particularly important subsets of the real line;
the following notation will be used throughout the text.
\[
\begin{aligned}
\relax
[a,b] &= \{x \in \mathbb{R} : a \leq x \leq b\}, \\
(a,b) &= \{x \in \mathbb{R} : a < x < b\}, \\
[a,b) &= \{x \in \mathbb{R} : a \leq x < b\}, \\
(a, \infty) &= \{x \in \mathbb{R} : x > a\}, \\
(-\infty, b] &= \{x \in \mathbb{R} : x \leq b\}.
\end{aligned}
\]

\begin{remark}
The symbol $\infty$ is not a real number. In interval notation, it indicates that an interval is unbounded in the corresponding direction.
\end{remark}

Figure~\ref{fig:real-intervals} distinguishes included and excluded endpoints and shows how rays represent unbounded intervals.

\begin{figure}[!htbp]
\centering
\begin{tikzpicture}[x=1.3cm,y=1cm,>=Stealth]
  \foreach \rowheight in {0,-1.2,-2.4,-3.6,-4.8} {
    \draw[->,black!35] (0,\rowheight) -- (6,\rowheight);
  }
  \node[left] at (-0.3,0) {$[a,b]$};
  \node[left] at (-0.3,-1.2) {$(a,b)$};
  \node[left] at (-0.3,-2.4) {$[a,b)$};
  \node[left] at (-0.3,-3.6) {$(a,\infty)$};
  \node[left] at (-0.3,-4.8) {$(-\infty,b]$};
  \foreach \rowheight in {0,-1.2,-2.4} {
    \draw[figblue,very thick] (1,\rowheight) -- (4.5,\rowheight);
    \node[below=5pt] at (1,\rowheight) {$a$};
    \node[below=5pt] at (4.5,\rowheight) {$b$};
  }
  \fill[figblue] (1,0) circle (2.5pt) (4.5,0) circle (2.5pt);
  \draw[figblue,thick,fill=white] (1,-1.2) circle (2.5pt) (4.5,-1.2) circle (2.5pt);
  \fill[figblue] (1,-2.4) circle (2.5pt);
  \draw[figblue,thick,fill=white] (4.5,-2.4) circle (2.5pt);
  \draw[figgreen,very thick,->] (1,-3.6) -- (5.8,-3.6);
  \draw[figgreen,thick,fill=white] (1,-3.6) circle (2.5pt);
  \node[below=5pt] at (1,-3.6) {$a$};
  \draw[figgreen,very thick,->] (4.5,-4.8) -- (0.2,-4.8);
  \fill[figgreen] (4.5,-4.8) circle (2.5pt);
  \node[below=5pt] at (4.5,-4.8) {$b$};
\end{tikzpicture}
\caption[Intervals on the real line]{Closed, open, and half-open intervals, followed by two unbounded rays. Filled endpoints belong to the set; hollow endpoints do not. Colored arrows indicate continuation without bound, rather than an endpoint at infinity.}
\label{fig:real-intervals}
\end{figure}

\section{Field Structure, Order, and Completeness}

The real number system can also be characterized axiomatically. Three groups of properties describe its algebraic operations, order, and completeness.

\begin{definition}[The Real Number System]
$\mathbb{R}$ is a complete ordered field.
\end{definition}

\subsection{Axioms of Addition and Multiplication}

We begin with the notion of a field. Addition and multiplication are binary operations on $\mathbb{R}$ satisfying the following laws.
\begin{itemize}
\item \textit{Associative laws}: for all real numbers $a$, $b$, and $c$, both identities hold:

\[
\begin{aligned}
a + (b+c) &= (a+b) + c \\
a(bc) &= (ab)c.
\end{aligned}
\]

\item \textit{Commutative laws}: for all real numbers $a$, $b$, and $c$, both identities hold:

\[
\begin{aligned}
a + b &= b + a \\
ab &= ba.
\end{aligned}
\]

\item \textit{Identity laws}: for every real number $a$, the additive and multiplicative identities satisfy

\[
\begin{aligned}
a + 0 &= a \\
a \cdot 1 &= a.
\end{aligned}
\]

\item \textit{Inverse laws}: for every real number $a$, there is a unique real number $-a$ such that $a + (-a) = 0$, and, if $a \neq 0$, a unique real number $a^{-1}$ such that $a \, a^{-1} = 1$.

\item \textit{Distributive law}: for all real numbers $a$, $b$, and $c$, multiplication distributes over addition:

\[
a(b+c) = ab + ac.
\]

\end{itemize}

\subsection{Order and Compatibility with Arithmetic}

The field $\mathbb{R}$ is ordered: its strict order relation $<$ compares any two distinct real numbers and is compatible with addition and multiplication.
The order satisfies the following properties.
\begin{itemize}
\item \textit{Transitivity}: if $x < y$ and $y < z$ for real numbers $x, y, z$, then $x < z$.
\item \textit{Trichotomy}: for any two real numbers $x, y$, exactly one of the following relations holds:

\[
x < y, \quad x = y, \quad y < x.
\]

\end{itemize}

Compatibility with the field operations means that:
\begin{itemize}
\item \textit{Addition}: if $x < y$, then $x + z < y + z$ for every real number $z$.
\item \textit{Multiplication}: if $x < y$ and $z > 0$, then $xz < yz$.
\end{itemize}

\subsection{Completeness and the Least Upper Bound Property}

Completeness is expressed by the \emph{Axiom of Continuity}. Among its equivalent formulations, we use the least upper bound property.
To state it, we first introduce the notion of an upper bound. Let $A \subset \mathbb{R}$ be nonempty.

\begin{definition}[Upper Bounds]
A real number $b$ is an upper bound of $A$ if $x \leq b$ for every $x \in A$. A set with an upper bound is said to be bounded above.
\end{definition}

The Axiom of Continuity can now be stated precisely.

\begin{theorem}[Axiom of Continuity: Existence of the Least Upper Bound]
Every nonempty subset $A \subset \mathbb{R}$ that admits an upper bound has a least upper bound in $\mathbb{R}$.
\end{theorem}

This property distinguishes the real numbers from the rational numbers: it does not hold in $\mathbb{Q}$. For example, the set $A = \{ x \in \mathbb{Q} : x \geq 0, \ x^2 < 2 \} \subset \mathbb{Q}$ is bounded above in $\mathbb{Q}$, for example by $2$, but has no least upper bound in $\mathbb{Q}$. Such a bound would have to be $\sqrt{2}$,
which is irrational. In $\mathbb{R}$, the least upper bound guaranteed by the Axiom of Continuity has a special name and notation.

\subsection{Suprema, Maxima, and Infima}

\begin{definition}[The Supremum]
Let $A \subset \mathbb{R}$ be nonempty and bounded above. Its supremum, denoted $b = \sup A$, is its least upper bound.
Equivalently, $b = \sup A$ if and only if $x \leq b$ for every $x \in A$ and, for every $\varepsilon > 0$, there exists
$x_\varepsilon \in A$ such that $x_\varepsilon > b - \varepsilon$. The first condition makes $b$ an upper bound; the second ensures that no smaller number is an upper bound.
\end{definition}

Two distinctions are useful.
\begin{enumerate}
\item \textit{Maximum}: a real number $b$ is the maximum of $A$, denoted $b = \max A$, if $x \leq b$ for every $x \in A$ and $b \in A$.
A supremum need not belong to the set. For example, $A = [0,1)$ has supremum $1$ but no maximum, since $1 \notin A$.
Whenever a maximum exists, it equals the supremum: $\max A = \sup A$.
\item \textit{Unbounded sets}: for a set $A \subset \mathbb{R}$ that is not bounded above, the convention is to write $+\infty := \sup A$.
This extends the notation beyond sets with a finite supremum.
\end{enumerate}

Reversing the inequalities gives the equivalent greatest lower bound formulation of the Axiom of Continuity.

\begin{theorem}[Axiom of Continuity: Existence of the Greatest Lower Bound]
Every nonempty subset $A \subset \mathbb{R}$ that admits a lower bound has a greatest lower bound, called the infimum of $A$, in $\mathbb{R}$.
\end{theorem}

In summary, $\mathbb{R}$ is a field because its arithmetic operations satisfy the field axioms. It is ordered because these operations are compatible
with its order. In particular, addition preserves order for every $z$:
\[
x < y \quad \text{implies} \quad x+z < y+z.
\]
Likewise, the product of nonnegative numbers is nonnegative:
\[
x \geq 0 \quad \text{together with} \quad y \geq 0 \quad \text{implies} \quad xy \geq 0.
\]

Finally, $\mathbb{R}$ is complete because it satisfies the Axiom of Continuity.
Together, these properties characterize the real number system.

Figure~\ref{fig:atlas-supremum} contrasts a supremum that is not attained with a maximum that belongs to the set.

\begin{figure}[!htbp]
\centering
\begin{tikzpicture}[>=Stealth]
\draw[->,black!45] (-0.3,0) -- (6.5,0);
\draw[figblue,very thick] (0,0) -- (5,0);
\fill[figblue] (0,0) circle (2pt);
\draw[figblue,thick,fill=white] (5,0) circle (2.5pt);
\node[below] at (0,0) {$0$};
\node[below] at (5,0) {$1$};
\node[left] at (-0.5,0) {$[0,1)$};
\draw[dashed,figgreen] (3.5,-0.2) -- (3.5,0.7);
\node[above,figgreen] at (3.5,0.7) {$1-\varepsilon$};
\fill[figred] (4.25,0) circle (2.5pt);
\draw[figred] (4.25,0.1) -- (4.25,1.3);
\node[above,figred] at (4.25,1.3) {$1-\varepsilon/2$};
\draw[->,black!45] (-0.3,-1.7) -- (6.5,-1.7);
\draw[figblue,very thick] (0,-1.7) -- (5,-1.7);
\fill[figblue] (0,-1.7) circle (2pt) (5,-1.7) circle (2.5pt);
\node[below] at (0,-1.7) {$0$};
\node[below] at (5,-1.7) {$1$};
\node[left] at (-0.5,-1.7) {$[0,1]$};
\end{tikzpicture}
\caption{The sets $[0,1)$ (top) and $[0,1]$ (bottom) both have supremum $1$. Only the second contains that endpoint and therefore has a maximum. For the illustrated choice $0<\varepsilon<1$, the highlighted point $1-\varepsilon/2$ lies within $\varepsilon$ of the supremum in the first set.}
\label{fig:atlas-supremum}
\end{figure}

\section{Decimal Expansions and Geometric Series}

\subsection{Decimal Representations as Infinite Series}

Real numbers can also be represented in decimal notation. Every real number $x$ can be written in the form $x = \pm b_0.b_1 b_2 \ldots$, where $b_0 \in \mathbb{N}$ gives the digits before the decimal point, and the digits $b_i \in \{0, 1, \dots, 9\}$ form the fractional part.
A decimal expansion may terminate, repeat periodically, or continue indefinitely without repeating. The first two examples below are repeating decimals; the third is the nonrepeating expansion of $\pi$, which has been computed to more than thirty-one trillion digits:
\[
\begin{aligned}
0.33333\ldots &= 0.\overline{3}, \\
0.142857142857\ldots &= 0.\overline{142857}, \\
\pi &= 3.1415926535897\ldots.
\end{aligned}
\]
Thus, real numbers admit finite, repeating, or nonrepeating decimal representations. To give this notation a precise meaning, we interpret a decimal expansion as the value of an infinite series, defined as the limit of its partial sums:
\[
\begin{aligned}
x &= b_0.b_1 b_2 \ldots \\
&= b_0 + \frac{b_1}{10} + \frac{b_2}{100} + \cdots + \frac{b_n}{10^n} + \cdots \\
&= \sum_{k=0}^{\infty} \frac{b_k}{10^k} \\
&:= \lim_{n \to \infty} \sum_{k=0}^{n} \frac{b_k}{10^k}.
\end{aligned}
\]

\subsection{Finite Geometric Sums}

The treatment of repeating decimals relies on the formula for the sum of a finite geometric progression.

\begin{theorem}[Geometric Sum Formula]
Given any $q \neq 1$,
\[
\sum_{k=0}^{n-1} q^k = \frac{1 - q^n}{1-q}.
\]

\end{theorem}

\begin{proof}
Multiplying both sides of the identity \[
\sum_{k=0}^{n-1} q^k = q^0 + q^1 + q^2 + \cdots + q^{n-1}
\] by $1-q$, we obtain
\[
\begin{aligned}
(1-q) \sum_{k=0}^{n-1} q^k &= (1-q)(q^0 + q^1 + q^2 + \cdots + q^{n-1}) \\
&= (q^0 + q^1 + \cdots + q^{n-1}) - (q^1 + q^2 + \cdots + q^n).
\end{aligned}
\]
All of the intermediate terms cancel in pairs, leaving only
\[
q^0 - q^n = 1 - q^n,
\]
which proves the formula.
\end{proof}

\subsection{Infinite Geometric Series}

When $|q| < 1$, taking the limit $n \to \infty$ in the finite sum gives the corresponding infinite geometric series.

\begin{theorem}[Geometric Series Formula]
Given any $q \in \mathbb{R}$ with $|q| < 1$,
\[
\sum_{k=0}^{\infty} q^k = \frac{1}{1-q}.
\]

\end{theorem}

\begin{proof}
By definition,
\[
\sum_{k=0}^{\infty} q^k = \lim_{n \to \infty} \sum_{k=0}^{n-1} q^k.
\]
By the Geometric Sum Formula,
\[
\sum_{k=0}^{n-1} q^k = (1 - q^n)/(1-q).
\]
Since $|q| < 1$, the geometric sequence satisfies
\[
\lim_{n \to \infty} q^n = 0,
\]
and the result follows by taking the limit in the finite-sum formula.
\end{proof}

Figure~\ref{fig:atlas-geometric-area} gives an area interpretation of a convergent geometric series.

\begin{figure}[!htbp]
\centering
\begin{tikzpicture}[>=Stealth]
\fill[figblue!28] (0,0) rectangle (2,4);
\fill[figgreen!28] (2,2) rectangle (4,4);
\fill[figred!25] (3,0) rectangle (4,2);
\fill[figblue!50] (2,0) rectangle (3,1);
\draw[thick] (0,0) rectangle (4,4);
\draw (2,0) -- (2,4) (2,2) -- (4,2) (3,0) -- (3,2) (2,1) -- (3,1);
\node at (1,2) {$1/2$};
\node at (3,3) {$1/4$};
\node at (3.5,1) {$1/8$};
\node at (2.5,0.5) {$1/16$};
\node at (2.5,1.5) {$\cdots$};
\draw[<->] (0,-0.4) -- (4,-0.4) node[midway,below] {$1$};
\end{tikzpicture}
\caption{A unit square is successively partitioned into regions of area $1/2$, $1/4$, $1/8$, and $1/16$. The uncolored remainder has area $1/16$ after four steps; continuing the construction makes the remainder tend to zero and illustrates the sum of the geometric series with ratio $1/2$.}
\label{fig:atlas-geometric-area}
\end{figure}

\subsection{Why Repeating Nines Equal One}

The Geometric Series Formula resolves a familiar question about decimal notation: whether $0.\overline{9}$ is equal to $1$.

\begin{theorem}[Repeating Nines: $0.999\ldots=1$]
The repeating decimal $0.\overline{9}=0.999\ldots$ is equal to $1$; that
is,
\[
0.\overline{9}=1.
\]
\end{theorem}

\begin{proof}
By the definition of decimal notation given above, we may write
\[
\begin{aligned}
0.\overline{9} &= 0 + \frac{9}{10} + \frac{9}{100} + \cdots + \frac{9}{10^n} + \cdots \\
&= \frac{9}{10} \sum_{k=0}^{\infty} \frac{1}{10^k}.
\end{aligned}
\]
Applying the Geometric Series Formula with $q = 1/10$, we obtain
\[
\begin{aligned}
\sum_{k=0}^{\infty} (1/10)^k &= 1/(1 - 1/10) \\
&= 10/9,
\end{aligned}
\]
and hence
\[
\begin{aligned}
0.\overline{9} &= \frac{9}{10} \cdot \frac{10}{9} \\
&= 1.
\end{aligned}
\]

\end{proof}

\section{Rational Approximation and Uncountability}

\subsection{Density of the Rational Numbers}

Decimal expansions show that every real number can be approximated arbitrarily well by rational numbers with finite decimal representations.

\begin{theorem}[Density of $\mathbb{Q}$ in $\mathbb{R}$]
Every real number can be approximated to any desired degree of accuracy by rational numbers with finite decimal representations.
\end{theorem}

\begin{proof}
Let $x\in\mathbb{R}$. For each positive integer $n$, define
\[
r_n=\frac{\lfloor 10^n x\rfloor}{10^n}.
\]
The numerator is an integer, so $r_n$ is rational and has a finite
decimal representation. From the defining property of the floor function,
\[
\lfloor 10^n x\rfloor\leq 10^n x<\lfloor 10^n x\rfloor+1.
\]
Dividing by $10^n$ gives
\[
0\leq x-r_n<10^{-n}.
\]
This argument includes $x=0$, negative real numbers, and terminating
decimals; in the terminating case equality may occur on the left. Since
$10^{-n}\to0$, the sequence $r_n$ converges to $x$.
\end{proof}

The same series methods show that the repeating decimal $8.4\overline{35}$ equals
\[
(8435 - 84)/990.
\]

More generally, every repeating decimal represents a rational number, as can be shown by extending the calculation used for $0.\overline{9}$.

\subsection{Countability and Cantor's Diagonal Argument}

Although both $\mathbb{Q}$ and $\mathbb{R}$ are infinite, they differ in size, or \emph{cardinality}.
The rational numbers are countable: there is a bijection, a function that is both one-to-one and onto, from $\mathbb{N}$ to $\mathbb{Q}$.
Consequently, the rational numbers can be listed in an infinite sequence.

The real numbers, by contrast, are uncountable. No one-to-one function from $\mathbb{R}$ to $\mathbb{N}$ exists, so no sequence can list all real numbers.
Even the interval $[0,1]$ is uncountable.

\begin{theorem}[Cantor's Uncountability Theorem]
The set $[0,1]$ is uncountable.
\end{theorem}

\textit{Cantor's diagonal argument}, introduced in 1891, explains the underlying idea.
Consider the set $T$ of infinite binary sequences, whose entries are either zero or one, and suppose that all its elements have been listed.
Construct a new sequence by choosing its $n$-th digit to differ from the $n$-th digit of the $n$-th listed sequence, for every $n$.
The new sequence differs from every entry in the list, contradicting the assumed enumeration.
The full details connecting this construction to the interval are beyond the scope of this discussion.

\section{Complex Numbers and Polynomial Equations}

\subsection{Cartesian Representation and Arithmetic}

Complex numbers extend the real number system to allow solutions of polynomial equations such as $x^2 + 1 = 0$, which has no solution in $\mathbb{R}$.

\begin{definition}[Complex Numbers in Cartesian Form]
A complex number is an expression of the form $a + bi$, where $a$ and $b$ are real numbers, and $i$ is a symbol satisfying $i^2 = -1$. The complex
number $a+bi$ can also be represented by the ordered pair $(a,b)$ and plotted as a point in the \emph{Argand plane}.
In particular, the complex number $i = 0 + 1 \cdot i$ is identified with the point $(0,1)$.
\end{definition}

For $a+bi$, the number $a$ is the real part and $b$ is the imaginary part. For example, $4 - 3i$ has real part $4$ and imaginary part $-3$.
Two complex numbers $a+bi$ and $c+di$ are equal exactly when $a = c$ and $b = d$.
In the Argand plane, the horizontal axis is the real axis and the vertical axis is the imaginary axis. Addition and subtraction are performed separately on the real and imaginary parts:
\[
\begin{aligned}
(a+bi) + (c+di) &= (a+c) + (b+d)i, \\
(a+bi) - (c+di) &= (a-c) + (b-d)i.
\end{aligned}
\]
For instance,
\[
(1-i) + (4+7i) = 5 + 6i.
\]

The product of two complex numbers is defined in such a way that the usual commutative and distributive laws continue to hold, together with the
defining relation $i^2 = -1$; expanding the product $(a+bi)(c+di)$ using distributivity and substituting $i^2 = -1$ yields
\[
(a+bi)(c+di) = (ac - bd) + (ad+bc) i.
\]
As an illustration,
\[
\begin{aligned}
(-1+3i)(2-5i) &= (-1)(2-5i) + 3i(2-5i) \\
&= -2 + 5i + 6i - 15 i^2 \\
&= 13 + 11i.
\end{aligned}
\]
Here the last equality uses $i^2 = -1$.

\subsection{Conjugation, Modulus, and Division}

Division uses the complex conjugate to make the denominator real, much as rationalization removes a square root from a denominator.

\begin{definition}[Complex Conjugation]
For a complex number $z = a+bi$, its complex conjugate is defined to be $\overline{z} = a - bi$.
\end{definition}

To divide by a nonzero complex number, multiply the numerator and denominator by the conjugate of the denominator. For example, to express $(-1+3i)/(2+5i)$ in the form $a+bi$, use the conjugate $2-5i$. The numerator becomes $(-1+3i)(2-5i)$, which was evaluated above. Dividing by the resulting real denominator gives $13/29 + 11i/29$. Geometrically, the complex conjugate $\overline{z}$ is the reflection of $z$ in the real axis. The complex conjugate obeys several elementary
algebraic properties, namely
\[
\begin{aligned}
\overline{z+w} &= \overline{z} + \overline{w} \\
\overline{zw} &= \overline{z}\,\overline{w}.
\end{aligned}
\]
Both identities follow directly from the definitions of the sum and product of complex numbers. The modulus, or absolute value, $|z|$ of a complex
number $z = a+bi$ is defined as its distance from the origin in the Argand plane, that is,
\[
|z| = \sqrt{a^2+b^2}.
\]
It is a direct consequence of the definitions that
\[
\begin{aligned}
z \overline{z} &= (a+bi)(a-bi) \\
&= a^2 + abi - abi - b^2 i^2 \\
&= a^2 + b^2.
\end{aligned}
\]
Consequently, $z\overline{z} = |z|^2$, which explains the division rule. For a nonzero divisor, it can be written as
\[
\begin{aligned}
z/w &= z\overline{w}/(w\overline{w}) \\
&= z\overline{w}/|w|^2.
\end{aligned}
\]

Figure~\ref{fig:complex-conjugation} illustrates why conjugation preserves the modulus while reversing the sign of the imaginary part.

\begin{figure}[!htbp]
\centering
\begin{tikzpicture}[scale=1.2,>=Stealth]
  \draw[->,black!65] (-0.6,0) -- (4.7,0) node[right] {$\mathrm{Re}$};
  \draw[->,black!65] (0,-2.1) -- (0,2.1) node[above] {$\mathrm{Im}$};
  \draw[gridline,dashed] (0,1.6) -- (3,1.6) -- (3,-1.6) -- (0,-1.6);
  \draw[vecA] (0,0) -- (3,1.6);
  \draw[vecB] (0,0) -- (3,-1.6);
  \fill[figblue] (3,1.6) circle (2pt);
  \fill[figred] (3,-1.6) circle (2pt);
  \node[above right,figblue] at (3,1.6) {$z=a+bi$};
  \node[below right,figred] at (3,-1.6) {$\overline{z}=a-bi$};
  \node[below left] at (0,0) {$0$};
  \node[below right] at (3,0) {$a$};
  \node[left] at (0,1.6) {$b$};
  \node[left] at (0,-1.6) {$-b$};
\end{tikzpicture}
\caption[Complex conjugation as reflection]{A complex number and its conjugate are mirror images across the real axis. They have the same real component and opposite imaginary components; their equal distances from the origin express the identity $|\overline{z}|=|z|$.}
\label{fig:complex-conjugation}
\end{figure}

\subsection{Quadratic Equations and Complex Roots}

Since $i^2 = -1$, one may regard $i$ as a square root of $-1$; however, one also has
\[
\begin{aligned}
(-i)^2 &= i^2 \\
&= -1.
\end{aligned}
\]
Consequently, $-i$ is likewise a square root of $-1$. By convention, $i$ is designated the principal square root of $-1$, and one writes $\sqrt{-1} = i$;
more generally, for any positive number $c$, one writes $\sqrt{-c} = \sqrt{c}\, i$.

With this convention, the usual derivation of the quadratic formula for an equation with a nonzero leading coefficient, $ax^2 + bx + c = 0$, remains valid even when the discriminant $b^2 - 4ac$ is negative. The roots are
\[
x = \frac{-b \pm \sqrt{b^2 - 4ac}}{2a}.
\]

\begin{example}
For $x^2+x+1=0$ we have $a=b=c=1$, hence
\[
\Delta=b^2-4ac=1-4=-3.
\]
Using $\sqrt{-3}=i\sqrt{3}$ in the quadratic formula gives
\[
x=\frac{-1\pm\sqrt{-3}}{2}
=\frac{-1\pm i\sqrt{3}}{2}.
\]
Substitution provides a direct check: each root satisfies
$x^2+x+1=0$.
\end{example}

The two roots in this example are complex conjugates. In general, the nonreal roots of a quadratic equation with real coefficients $a$, $b$, $c$
occur in conjugate pairs, while each real root coincides with its own conjugate.

Allowing complex solutions ensures that every quadratic equation has a solution. More generally, every polynomial equation of degree at least one,
\[
a_n x^n + a_{n-1} x^{n-1} + \cdots + a_1 x + a_0 = 0,
\]
has a complex solution. This result, proved by Gauss, is known as the \textit{Fundamental Theorem of Algebra}.

\section{Polar Representation and the Complex Exponential}

\subsection{Polar Coordinates: Modulus and Argument}

A complex number $z = a+bi$ corresponds to the point $(a,b)$ in the Argand plane. This point can also be described by polar coordinates
$(r,\theta)$, with $r \geq 0$, related to the Cartesian coordinates by
\[
\begin{aligned}
a &= r \cos \theta \\
b &= r \sin \theta.
\end{aligned}
\]
This gives the \emph{polar form} of $z$:
\[
\begin{aligned}
z &= a+bi \\
&= (r\cos\theta) + (r\sin\theta) i \\
&= r(\cos\theta + i \sin\theta),
\end{aligned}
\]
The modulus is given by
\[
\begin{aligned}
r &= |z| \\
&= (a^2+b^2)^{1/2}.
\end{aligned}
\]

When the real part is nonzero, the argument also satisfies $\tan \theta = b/a$.
For a nonzero complex number, the angle $\theta$ is called its \emph{argument}, written $\theta = \arg(z)$.
The argument is determined only up to an integer multiple of $2\pi$.

\begin{example}
We express $z = 1+i$ and $w = \sqrt{3} - i$ in polar form.
For $z$, we have
\[
\begin{gathered}
\begin{aligned}
r &= |z| \\
&= (1^2+1^2)^{1/2} \\
&= \sqrt{2}
\end{aligned} \\
\tan \theta = 1.
\end{gathered}
\]
Since $z$ lies in the first quadrant, we take $\theta = \pi/4$, giving
\[
z = \sqrt{2} ( \cos(\pi/4) + i \sin(\pi/4) ).
\]
For $w$, we have
\[
\begin{gathered}
\begin{aligned}
r &= |w| \\
&= (3+1)^{1/2} \\
&= 2
\end{aligned} \\
\tan \theta = -1/\sqrt{3}.
\end{gathered}
\]
Since $w$ lies in the fourth quadrant, we take $\theta = -\pi/6$, giving
\[
w = 2 ( \cos(-\pi/6) + i \sin(-\pi/6) ).
\]

\end{example}

Figure~\ref{fig:ch2-polar-form} shows the geometric meaning of modulus, argument, and Cartesian components.

\begin{figure}[!htbp]
\centering
\begin{tikzpicture}[scale=1.3, >=Stealth]
  \draw[->] (-0.5,0) -- (3.6,0) node[right] {$\mathrm{Re}$};
  \draw[->] (0,-0.5) -- (0,3.1) node[above] {$\mathrm{Im}$};
  \draw[figblue, thick, dashed] (0,0) circle (2.5cm);
  \node[figblue] at (2.9,2.15) {$|z|=r$};
  \coordinate (Z) at (50:2.5);
  \draw[vecB] (0,0) -- (Z) node[right, xshift=2pt] {$z = a+bi = r\,e^{i\theta}$};
  \draw[black!60, dashed] (Z) -- (Z |- 0,0) node[below] {$a = r\cos\theta$};
  \draw[black!60, dashed] (Z) -- (0,0 |- Z) node[left] {$b = r\sin\theta$};
  \draw[figgreen, thick] (0.6,0) arc (0:50:0.6);
  \node[figgreen] at (0.88,0.4) {$\theta$};
  \node[circle,fill=figred,inner sep=1.3pt] at (Z) {};
\end{tikzpicture}
\caption[Polar representation of a complex number]{Polar representation of
  $z = a+bi = r e^{i\theta} \in \C$ in the Argand--Gauss plane: modulus $r=|z|$,
  argument $\theta$, and Cartesian projections $a=r\cos\theta$, $b=r\sin\theta$.}
\label{fig:ch2-polar-form}
\end{figure}

\subsection{Arithmetic, Powers, and Roots in Polar Form}

Polar form reveals the geometry of multiplication and division. Given two complex numbers in polar form,
\[
\begin{aligned}
z_1 &= r_1(\cos\theta_1 + i\sin\theta_1) \\
z_2 &= r_2(\cos\theta_2 + i\sin\theta_2).
\end{aligned}
\]
Expanding their product and applying the addition formulas for sine and cosine yields $z_1 z_2 = r_1 r_2 \big[ \cos(\theta_1+\theta_2) + i \sin(\theta_1+\theta_2) \big]$, so multiplication multiplies the moduli and adds the arguments. Similarly, the subtraction formulas for sine and cosine show that division
divides the moduli and subtracts the arguments. For $z_2 \neq 0$, the quotient is
\[
\frac{z_1}{z_2} = \frac{r_1}{r_2} \big[ \cos(\theta_1-\theta_2) + i \sin(\theta_1-\theta_2) \big].
\]

Taking $z_1 = 1$, so that $r_1 = 1$ and $\theta_1 = 0$, and writing $z_2 = z$, gives the reciprocal formula. Any nonzero complex number in the following polar form has the reciprocal shown alongside it:
\[
\begin{aligned}
z &= r(\cos\theta + i \sin \theta), \\
\frac{1}{z} &= \frac{1}{r} \big( \cos\theta - i \sin\theta \big).
\end{aligned}
\]

\begin{example}
Using the polar representations of $1+i$ and $\sqrt{3}-i$ found above, the multiplication formula gives
\[
\begin{aligned}
(1+i)(\sqrt{3}-i) &= 2\sqrt{2} \left[ \cos(\pi/4 - \pi/6) + i \sin(\pi/4-\pi/6) \right] \\
&= 2\sqrt{2} \left[ \cos(\pi/12) + i \sin(\pi/12) \right].
\end{aligned}
\]
To verify the result in Cartesian form, multiply directly:
\[
\begin{aligned}
(1+i)(\sqrt{3}-i)
&=\sqrt{3}-i+i\sqrt{3}-i^2\\
&=(1+\sqrt{3})+i(\sqrt{3}-1).
\end{aligned}
\]
The two forms agree because
\[
2\sqrt{2}\cos(\pi/12)=1+\sqrt{3},\qquad
2\sqrt{2}\sin(\pi/12)=\sqrt{3}-1.
\]

\end{example}

Repeated multiplication in polar form gives a rule for powers. Consider a complex number in polar form:
\[
z = r(\cos\theta + i\sin\theta).
\]
Its square and cube are given by the following identities:
\[
\begin{aligned}
z^2 &= r^2(\cos 2\theta + i \sin 2\theta) \\
z^3 &= z^2 z \\
&= r^3(\cos 3\theta + i \sin 3\theta).
\end{aligned}
\]
The general result is named after the French mathematician Abraham de Moivre.

\begin{theorem}[De Moivre's Theorem]
Let $n$ be a positive integer. For a complex number in polar form, its $n$-th power is given by
\[
\begin{aligned}
z &= r(\cos\theta + i\sin\theta), \\
z^n &= \big[ r (\cos\theta + i \sin \theta) \big]^n \\
&= r^n \big( \cos n\theta + i \sin n\theta \big).
\end{aligned}
\]

\end{theorem}

Thus, taking the $n$-th power raises the modulus to the $n$-th power and multiplies the argument by $n$.

\begin{example}
We compute
\[
\left( \frac{1}{2} + \frac{1}{2}i \right)^{10}.
\]
The calculation uses the following factorization and the polar form of $1+i$ found above:
\[
\begin{aligned}
\frac{1}{2} + \frac{1}{2}i &= \frac{1}{2}(1+i), \\
1+i &= \sqrt{2}\big(\cos(\pi/4) + i\sin(\pi/4)\big).
\end{aligned}
\]
De Moivre's Theorem gives the tenth power from this modulus and argument.
Indeed,
\[
\frac12+\frac12i
=\frac{1}{\sqrt2}\left(\cos\frac{\pi}{4}+i\sin\frac{\pi}{4}\right),
\]
and therefore
\[
\begin{aligned}
\left(\frac12+\frac12i\right)^{10}
&=\left(\frac{1}{\sqrt2}\right)^{10}
\left(\cos\frac{10\pi}{4}+i\sin\frac{10\pi}{4}\right)\\
&=\frac1{32}\left(\cos\frac{5\pi}{2}+i\sin\frac{5\pi}{2}\right)\\
&=\frac{i}{32}.
\end{aligned}
\]
\end{example}

De Moivre's Theorem also gives the $n$-th roots of a complex number. An $n$-th root of $z$ is a complex number $w$ satisfying $w^n = z$.
For nonzero $z$, write both numbers in polar form:
\[
\begin{aligned}
w &= s(\cos\varphi + i\sin\varphi) \\
z &= r(\cos\theta + i\sin\theta).
\end{aligned}
\]
Applying De Moivre's Theorem to $w^n$, the equation $w^n = z$ becomes
\[
s^n(\cos n\varphi + i \sin n\varphi) = r(\cos\theta + i \sin\theta).
\]

Comparing moduli and arguments on both sides, and recalling that sine and cosine have period $2\pi$, gives the following relations for some integer $k$:
\[
\begin{aligned}
s &= r^{1/n} \\
n\varphi &= \theta + 2k\pi.
\end{aligned}
\]

This gives $\varphi = (\theta + 2k\pi)/n$. Each index $k = 0, 1, 2, \dots, n-1$ yields a distinct value of $w$, giving the following description of the $n$-th roots of a complex number.

\begin{theorem}[Roots of a Complex Number]
Let $n$ be a positive integer and let $z$ be a nonzero complex number with the polar representation below. For each index $k = 0, 1, \dots, n-1$, the second formula gives one of its exactly $n$ distinct $n$-th roots:
\[
\begin{aligned}
z &= r(\cos\theta + i\sin\theta), \\
w_k &= r^{1/n} \left[ \cos\left( \frac{\theta + 2k\pi}{n} \right) + i \sin\left( \frac{\theta+2k\pi}{n} \right) \right].
\end{aligned}
\]

\end{theorem}

All these roots have modulus $r^{1/n}$ and therefore lie on a circle of radius $r^{1/n}$ in the complex plane.
Their successive arguments differ by $2\pi/n$, so the roots are equally spaced around the circle.

\begin{example}
We find the sixth roots of $z = -8$. In trigonometric form,
\[
z = 8(\cos\pi + i\sin\pi).
\]
Applying the formula with $n = 6$, we obtain the six roots by taking $k = 0, 1, 2, 3, 4, 5$ in
\[
8^{1/6} \left[ \cos\left( \frac{\pi + 2k\pi}{6} \right) + i \sin\left( \frac{\pi+2k\pi}{6} \right) \right].
\]
All six resulting points lie on the circle of radius $8^{1/6} = \sqrt{2}$ in the complex plane, equally spaced around it.
Their arguments and Cartesian forms are
\[
\begin{aligned}
w_0&=\sqrt2\,e^{i\pi/6}
=\frac{\sqrt6}{2}+i\frac{\sqrt2}{2},\\
w_1&=\sqrt2\,e^{i\pi/2}=i\sqrt2,\\
w_2&=\sqrt2\,e^{i5\pi/6}
=-\frac{\sqrt6}{2}+i\frac{\sqrt2}{2},\\
w_3&=\sqrt2\,e^{i7\pi/6}
=-\frac{\sqrt6}{2}-i\frac{\sqrt2}{2},\\
w_4&=\sqrt2\,e^{i3\pi/2}=-i\sqrt2,\\
w_5&=\sqrt2\,e^{i11\pi/6}
=\frac{\sqrt6}{2}-i\frac{\sqrt2}{2}.
\end{aligned}
\]
For every $k$, raising $w_k$ to the sixth power multiplies its argument
by $6$ and raises its modulus $\sqrt{2}$ to the sixth power, yielding
$(\sqrt{2})^6=8$ and hence $w_k^6=-8$.
\end{example}

Figure~\ref{fig:sixth-roots-minus-eight} displays these six roots and the equal angular spacing between successive ones.

\begin{figure}[!htbp]
\centering
\begin{tikzpicture}[scale=2.1,>=Stealth]
  \draw[->,black!55] (-1.85,0) -- (1.95,0) node[right] {$\mathrm{Re}$};
  \draw[->,black!55] (0,-1.85) -- (0,1.95) node[above] {$\mathrm{Im}$};
  \draw[figblue,thick] (0,0) circle ({sqrt(2)});
  \draw[figblue!45,dashed] (30:{sqrt(2)}) -- (90:{sqrt(2)}) -- (150:{sqrt(2)})
    -- (210:{sqrt(2)}) -- (270:{sqrt(2)}) -- (330:{sqrt(2)}) -- cycle;
  \foreach \angle/\rootindex in {30/0,90/1,150/2,210/3,270/4,330/5} {
    \draw[gridline,dashed] (0,0) -- (\angle:{sqrt(2)});
    \fill[figblue] (\angle:{sqrt(2)}) circle (1.3pt);
  }
  \node[above right] at (30:{sqrt(2)}) {$w_0$};
  \node[above right] at (90:{sqrt(2)}) {$w_1$};
  \node[above left] at (150:{sqrt(2)}) {$w_2$};
  \node[below left] at (210:{sqrt(2)}) {$w_3$};
  \node[below right] at (270:{sqrt(2)}) {$w_4$};
  \node[below right] at (330:{sqrt(2)}) {$w_5$};
  \draw[figgreen,thick,->] (30:0.55) arc[start angle=30,end angle=90,radius=0.55];
  \node[figgreen] at (60:0.85) {$\pi/3$};
  \node[below left] at (0,0) {$0$};
\end{tikzpicture}
\caption[The six sixth roots of minus eight]{The solutions of $w^6=-8$ form the vertices of a regular hexagon on the circle of radius $\sqrt{2}$. Starting at argument $\pi/6$, the labels follow increasing argument in steps of $\pi/3$. Dashed segments emphasize the symmetry and do not represent additional roots.}
\label{fig:sixth-roots-minus-eight}
\end{figure}

\subsection{The Complex Exponential and Euler's Formula}


We now define $e^z$ for a complex number $z = x+iy$. Power series, studied in detail in later chapters, also allow complex terms.
Motivated by the Maclaurin series for $e^x$, we define the complex exponential by
\[
\begin{aligned}
e^z &:= \sum_{n=0}^{\infty} \frac{z^n}{n!} \\
&= 1 + z + \frac{z^2}{2!} + \frac{z^3}{3!} + \cdots.
\end{aligned}
\]
This has the same fundamental algebraic properties as the real exponential. In particular, $e^{z+w} = e^z e^w$, for all complex numbers $z$ and $w$.


Substituting $z = iy$ for real $y$ into the defining series and using $i^2 = -1$, $i^3 = -i$, $i^4 = 1$, together with the periodic recurrence of higher powers of $i$, one obtains
\[
\begin{aligned}
e^{iy} &= 1 + iy + \frac{(iy)^2}{2!} + \frac{(iy)^3}{3!} + \frac{(iy)^4}{4!} + \cdots \\
&= \left(1 - \frac{y^2}{2!} + \frac{y^4}{4!} - \cdots\right) + i \left( y - \frac{y^3}{3!} + \frac{y^5}{5!} - \cdots \right).
\end{aligned}
\]

The two series in parentheses are the Maclaurin series for $\cos y$ and $\sin y$, respectively. This gives Euler's formula.

\begin{theorem}[Euler's Formula]
For every real number $y$,
\[
e^{iy} = \cos y + i \sin y.
\]

\end{theorem}

Combining Euler's formula with the multiplicative property gives
\[
\begin{aligned}
e^{x+iy} &= e^x e^{iy} \\
&= e^x (\cos y + i \sin y).
\end{aligned}
\]
This expresses the exponential of $x+iy$ using the real exponential and trigonometric functions.

\begin{example}
We evaluate two particular complex exponentials using Euler's formula. First,
\[
\begin{aligned}
e^{i\pi} &= \cos\pi + i \sin\pi \\
&= -1 + i \cdot 0 \\
&= -1.
\end{aligned}
\]
This is an identity linking five fundamental mathematical constants. Second, the general formula for $e^{x+iy}$ gives
\[
\begin{aligned}
e^{1+i\pi/2} &= e^1 (\cos(\pi/2) + i \sin(\pi/2)) \\
&= e(0 + i \cdot 1) \\
&= ie.
\end{aligned}
\]

\end{example}

\enlargethispage{3\baselineskip}
Euler's formula also gives a short proof of De Moivre's Theorem. Writing $z = re^{i\theta}$, we obtain
\[
\begin{aligned}
\big[ r(\cos\theta + i\sin\theta) \big]^n &= \big( r e^{i\theta} \big)^n \\
&= r^n e^{in\theta} \\
&= r^n (\cos n\theta + i \sin n\theta),
\end{aligned}
\]
so De Moivre's Theorem follows directly from the laws of exponents for the complex exponential.

\newpage
\chapter{Convergence Theory and Numerical Series}

\section{Real Sequences and Their Basic Properties}

\subsection{Sequences as Functions and Their Graphs}

A sequence is a list of real numbers arranged in a definite order, which we denote by
\[
a_1, a_2, a_3, \dots, a_n, \dots
\]

The real number $a_n$ occupying the $n$-th position in this ordered list is called the $n$-th term of the sequence. More formally, a sequence is a
rule that assigns to every positive integer $n \in \mathbb{N}$ a unique real number $a_n \in \mathbb{R}$; equivalently, a sequence is simply a function
\[
\mathbb{N} \to \mathbb{R}, \qquad n \mapsto a_n
\]

Throughout this text a sequence will be denoted either by $\{a_n\}_{n \in \mathbb{N}}$, by $\{a_n\}$, or, when no ambiguity arises, simply by $a_n$.
Several elementary examples illustrate this notion.
\begin{align*}
a_n =& 1/n \quad \rightarrow \quad 1, 1/2, 1/3, \dots, 1/n, \dots, \\
a_n =& (-1)^n \quad \rightarrow \quad -1, 1, -1, 1, \dots, \\
a_n =& n^2 \quad \rightarrow \quad 1, 4, 9, \dots, \\
a_n =& n/(n+1) \quad \rightarrow \quad 1/2, 2/3, 3/4, 4/5, \dots.
\end{align*}

A sequence can be visualized by plotting its graph. Since a sequence is, by definition, a function whose domain is the set of positive integers,
its graph consists of an isolated collection of points with coordinates $(n, a_n)$, rather than a continuous curve.


\subsection{Boundedness, Monotonicity, and Eventual Properties}

Before studying the limiting behavior of a sequence, it is useful to introduce two qualitative properties that a sequence may possess, namely
boundedness and monotonicity. A sequence $\{a_n\}$ is said to be bounded above if there exists a real number $M$ such that $a_n \leq M$ for
every $n \in \mathbb{N}$, and it is said to be bounded below if there exists a real number $m$ such that $m \leq a_n$ for every $n \in \mathbb{N}$.
When a sequence is simultaneously bounded above and bounded below, it is simply called a bounded sequence. A closely related but distinct notion is
that of monotonicity.
\begin{itemize}
\item \textit{Increasing sequences}: if $a_n \leq a_{n+1}$ for every integer $n$, the sequence $\{a_n\}$ of real numbers is called increasing.
\item \textit{Decreasing sequences}: it is called decreasing if $a_n \geq a_{n+1}$ for every integer $n$.
\end{itemize}

A sequence is called monotonic if it is either increasing or decreasing. It is also convenient to introduce
terminology that will recur throughout the discussion of limits: a sequence $\{a_n\}$ is said to
satisfy a property \emph{eventually} if that property holds from some index onward, that is, for every sufficiently large $n$.

\section{Limits and Monotone Convergence}

\subsection{The Limit of a Sequence}

The central goal of the theory of sequences is to study the behavior of the terms $a_n$ as the index $n$ becomes
arbitrarily large; this behavior is encoded by the limit
\[
\lim_{n \to \infty} a_n
\]

Let $L \in \mathbb{R}$ be a given real number. The sequence $\{a_n\}$ converges to $L$ if its terms can be made as close to $L$ as desired by taking $n$ sufficiently large. We use either of the following equivalent notations:
\[
\begin{gathered}
\lim_{n \to \infty} a_n = L, \\
a_n \to L.
\end{gathered}
\]
The following definition makes this intuitive idea precise.

\begin{definition}
The statement $a_n \to L$ means that, for every $\varepsilon > 0$, there exists $N \in \mathbb{N}$ for which the following implication holds:
\[
n > N \implies |a_n - L| < \varepsilon.
\]

\end{definition}

An analogous definition applies when the terms grow without bound. The sequence $\{a_n\}$ diverges to positive infinity if its terms can be made as large as desired by taking $n$ sufficiently large. The two equivalent notations are
\[
\begin{gathered}
\lim_{n \to \infty} a_n = +\infty, \\
a_n \to +\infty.
\end{gathered}
\]

\begin{definition}
The statement $a_n \to +\infty$ means that, for every $M > 0$, there exists $N \in \mathbb{N}$ for which the following implication holds:
\[
n > N \implies a_n > M.
\]

\end{definition}

These notions allow us to classify sequences according to their limiting behavior. A sequence
$\{a_n\}$ is said to be convergent if, for some finite real number $L$,
\[
\lim_{n \to \infty} a_n = L.
\]

It is said to diverge to infinity if one of the following alternatives holds:
\[
\begin{gathered}
\lim_{n \to \infty} a_n = +\infty, \\
\lim_{n \to \infty} a_n = -\infty.
\end{gathered}
\]

If none of these alternatives holds, the sequence has no limit, even in the extended real line. Two representative examples are $a_n = (-1)^n$ and $a_n = \cos n$,
both of which oscillate indefinitely without settling toward any particular value. A particularly instructive
family of sequences is provided by the geometric sequences, which are sequences of the form $a_n = q^n, \qquad n \in \mathbb{N}$ for a fixed real number $q$, generating the terms $1, q, q^2, q^3, q^4, \dots$. The limiting behavior of a
geometric sequence depends entirely on the value of the common ratio $q$, according to the following
classification.
\[
\lim_{n \to \infty} q^n =
\begin{cases}
0 & \text{if } 0 \leq q < 1, \\
1 & \text{if } q = 1, \\
+\infty & \text{if } q > 1, \\
\text{no limit} & \text{if } q \leq -1
\end{cases}
\]

This result can be established directly from the $\varepsilon$--$N$ definition of limit, and its verification is
a valuable exercise for the reader.

Figure~\ref{fig:atlas-oscillating-sequence} shows why two subsequences with different limits prevent convergence.

\begin{figure}[!htbp]
\centering
\begin{tikzpicture}[>=Stealth]
\begin{axis}[width=11cm,height=5.8cm,axis lines=middle,ticklabel style={font=\small},label style={font=\small},legend style={font=\scriptsize,draw=black!20,fill=white},xlabel={$n$},ylabel={$a_n$},xmin=0,xmax=13,ymin=-1.5,ymax=1.5,xtick={2,4,6,8,10,12},ytick={-1,1}]
\addplot[black!30,dashed,domain=0:12.5] {1};
\addplot[black!30,dashed,domain=0:12.5] {-1};
\addplot[only marks,figblue,mark=*,mark size=2pt] coordinates {(2,1)(4,1)(6,1)(8,1)(10,1)(12,1)};
\addplot[only marks,figred,mark=square*,mark size=2pt] coordinates {(1,-1)(3,-1)(5,-1)(7,-1)(9,-1)(11,-1)};
\end{axis}
\end{tikzpicture}
\caption{The sequence $a_n=(-1)^n$ alternates between two horizontal levels. The even subsequence is constantly $1$ and the odd subsequence is constantly $-1$, so the full sequence cannot approach a single limit.}
\label{fig:atlas-oscillating-sequence}
\end{figure}

\subsection{The Monotone Convergence Theorem and Euler's Number}

One of the most important results concerning monotonic sequences is the following theorem, which guarantees
convergence purely on the basis of boundedness and monotonicity, without requiring any explicit knowledge of the
limiting value.

\begin{theorem}[Monotone Sequence Theorem]
Every bounded monotonic sequence is convergent.
\end{theorem}

\begin{remark}
If the sequence under consideration is increasing, it will be shown below that it converges precisely to its
least upper bound, whose existence is guaranteed by the Completeness Axiom of the real numbers.
\end{remark}

\begin{proof}
Suppose that $\{a_n\}$ is an increasing sequence; an entirely analogous argument, based on the greatest lower
bound, applies when $\{a_n\}$ is decreasing. Since $\{a_n\}$ is bounded, the set $\{a_n : n \in \mathbb{N}\}$
admits an upper bound, and therefore, by the Completeness Axiom, it admits a least upper bound, which we denote
by $L$. Fix $\varepsilon > 0$; we shall show that the terms of the sequence eventually lie in the interval
$(L - \varepsilon, L + \varepsilon)$. Since $L - \varepsilon$ is not an upper bound of the set, there exists an
index $N$ such that $a_N > L - \varepsilon$. Because the sequence is increasing, we have $a_n \geq a_N$ for every
$n > N$, and hence $a_n > L - \varepsilon$ for every $n > N$. On the other hand, since $L$ is an upper bound of
the set, we know that $a_n \leq L$ for every integer $n$. For every $n > N$, combining these two facts gives
\[
L - \varepsilon < a_n < L + \varepsilon,
\]
which is precisely the statement that
\[
\lim_{n \to \infty} a_n = L.
\]

\end{proof}

The conclusions of this theorem may be summarized as follows. If $\{a_n\}$ is increasing and bounded, then
\[
\lim_{n \to \infty} a_n = \sup_n \{a_n\}.
\]
If instead $\{a_n\}$ is decreasing and bounded, then
\[
\lim_{n \to \infty} a_n = \inf_n \{a_n\}.
\]
An analogous argument establishes a companion result for unbounded monotonic sequences.

\begin{theorem}[Monotone Divergence Theorem]
Every unbounded monotonic sequence is divergent.
\end{theorem}

Figure~\ref{fig:ch3-monotone-convergence} illustrates convergence to the supremum for a bounded increasing sequence.

\begin{figure}[!htbp]
\centering
\begin{tikzpicture}
  \begin{axis}[
    width=13cm, height=7cm,
    xlabel={$n$}, ylabel={$a_n$},
    axis lines=middle,
    xmin=0, xmax=16.3, ymin=0, ymax=3.6,
    ytick={3}, yticklabels={$L=\sup\nolimits_n a_n$},
    xtick=\empty,
    ]
    \addplot[figred, thick, dashed, domain=0:14] {3};
    \addplot[figblue, mark=*, mark size=2pt, only marks] coordinates {
      (1,0.6) (2,1.3) (3,1.85) (4,2.25) (5,2.5) (6,2.68) (7,2.8)
      (8,2.88) (9,2.93) (10,2.96) (11,2.98) (12,2.99) (13,3.0)
    };
    \draw[figgreen, dashed] (axis cs:0,3.15) -- (axis cs:14,3.15)
      node[pos=1,anchor=west,figgreen,font=\scriptsize] {$L+\varepsilon$};
    \draw[figgreen, dashed] (axis cs:0,2.85) -- (axis cs:14,2.85)
      node[pos=1,anchor=west,figgreen,font=\scriptsize] {$L-\varepsilon$};
    \draw[black!50,dotted] (axis cs:6.4,0) -- (axis cs:6.4,3.6)
      node[pos=1,above,font=\scriptsize,black!70]{$N$};
  \end{axis}
\end{tikzpicture}
\caption[Convergence of a bounded monotone sequence]{Increasing and bounded
  sequence $(a_n)$: the band $(L-\varepsilon,L+\varepsilon)$ contains all terms from a certain
  index $N$ onwards.}
\label{fig:ch3-monotone-convergence}
\end{figure}

A remarkable application of the Monotone Sequence Theorem concerns the sequence
\[
\left( 1 + \frac{1}{n} \right)^n, \qquad n \in \mathbb{N}
\]

It can be shown that this sequence is monotone increasing and bounded above; consequently, by the Monotone
Convergence Theorem, it admits a finite limit. This limit is a fundamental mathematical constant known as Euler's
number, defined by
\[
\begin{aligned}
e &:= \lim_{n \to \infty} \left( 1 + \frac{1}{n} \right)^n \\
&\approx 2.7182818284\ldots
\end{aligned}
\]

The origin of this constant is closely tied to a classical problem in finance. Jacob Bernoulli discovered it in
1683 while studying a question concerning compound interest. Consider an account that starts with one dollar and
pays one hundred percent interest per year; if the interest is credited only once, at the end of the year, the
value of the account will be two dollars. The natural question is what happens if the interest is computed and
credited more frequently throughout the year. If the interest is credited twice per year, the rate applicable to
each six-month period is fifty percent, so the initial dollar is multiplied by $1.5$ twice, yielding $1 \times (1.5)^2 = 2.25$ dollars at the end of the year. Compounding quarterly yields $1 \times (1.25)^4 = 2.4414\ldots$ dollars, while compounding monthly yields $1 \times (1 + 1/12)^{12} = 2.613035\ldots$ dollars. More generally,
if there are $n$ compounding intervals within the year, the interest rate applicable to each interval is
$100\%/n$, and the value of the account at the end of the year is precisely $1 \times (1 + 1/n)^n$ dollars,
which is exactly the term of the sequence considered above.

\subsection{Uniqueness and Boundedness of Limits}

Two natural questions arise once the notion of limit has been introduced: how many limits may a given sequence
possess, and what qualitative properties must a convergent sequence necessarily satisfy. The first question is
answered by the following uniqueness theorem.

\begin{theorem}[Uniqueness of the Limit]
Whenever a sequence has a finite or infinite limit, that limit is unique.
\end{theorem}

\begin{proof}
Suppose, for the sake of contradiction, that $a_n \to L_1$ and $a_n \to L_2$ with $L_1, L_2 \in \mathbb{R}$ and
$L_1 \neq L_2$. Set
\[
d := |L_1 - L_2| > 0
\]

By the definition of limit, there exist integers $N_1 \in \mathbb{N}$ and $N_2 \in \mathbb{N}$ such that the first bound holds for every $n > N_1$ and the second for every $n > N_2$:
\[
\begin{aligned}
|a_n - L_1| &< \frac{d}{2}, \\
|a_n - L_2| &< \frac{d}{2}.
\end{aligned}
\]
Consequently, for every $n > \max\{N_1, N_2\}$ we obtain, using the triangle inequality,
\[
|L_1 - L_2| = |L_1 - a_n + a_n - L_2| \leq |a_n - L_1| + |a_n - L_2| < \frac{d}{2} + \frac{d}{2} = d
\]

We have thus shown that $|L_1 - L_2| < d$, which contradicts the definition of $d$ and establishes the theorem.
\end{proof}

The second question concerns the relationship between convergence and boundedness.

\begin{theorem}
If a sequence is convergent, then it is bounded.
\end{theorem}

\begin{proof}
Suppose $a_n \to L$. Applying the definition of limit with $\varepsilon = 1$, we find an integer
$N \in \mathbb{N}$ such that $|a_n - L| < 1$ for every $n > N$. In particular, for such $n$ we obtain
\[
|a_n| = |a_n - L + L| \leq |a_n - L| + |L| < |L| + 1
\]

Combining this estimate with the finitely many terms $a_1, \dots, a_N$ gives a bound valid for every $n \in \mathbb{N}$:
\[
|a_n| \leq \max\{|a_1|, |a_2|, \dots, |a_N|, |L| + 1\},
\]
which shows that the sequence is bounded.
\end{proof}

\begin{remark}
The converse of this theorem is false: the sequence $\{(-1)^n\}$ is bounded, yet it is not convergent.
\end{remark}

\section{Asymptotic Comparison and Limit Laws}

\subsection{Comparing Infinities and Infinitesimals}

We call a sequence $\{a_n\}$ infinite if
\[
\lim_{n \to \infty} a_n = +\infty.
\]

A natural question that arises
once several infinite sequences are available is which one grows faster; for instance, one may ask which is the
stronger infinity between $n$ and $n^2$, or, equivalently, which of the two diverges to $+\infty$ more rapidly.
In order to compare two infinite sequences $a_n \to +\infty$ and $b_n \to +\infty$, it is customary to examine
the limit of their quotient,
\[
\lim_{n \to \infty} \frac{a_n}{b_n}
\]

If this limit equals $+\infty$, we say that $a_n$ is an infinite of higher order than $b_n$, and we write
$b_n \prec a_n$. If instead the limit equals $0$, we say that $a_n$ is an infinite of lower order than $b_n$,
and we write $a_n \prec b_n$. A third possibility, of particular importance, occurs when the quotient tends to
one.

\begin{definition}
If \[
\lim_{n \to \infty} a_n / b_n = 1,
\] we say that $a_n$ is asymptotically equivalent to $b_n$, and we write
$a_n \sim b_n$. In this case $a_n$ and $b_n$ are said to be infinities of the same order.
\end{definition}

This comparison mechanism applies in particular to powers of $n$. To compare $n^a$ with $n^b$ for exponents
$a > b$, one computes the limit of the quotient $n^b / n^a \to 0$, from which it follows that $n^a$ is an
infinity of higher order than $n^b$; symbolically, $n^b \prec n^a$ whenever $b < a$. More generally, a sequence
$a_n$ is said to be an infinity of order $\beta$ if \[
\lim_{n \to \infty} \frac{a_n}{n^\beta} = k
\] for suitable constants $k > 0$ and $\beta > 0$; the expression $k n^\beta$ is then called the leading part of
$a_n$. For example, the sequence $a_n = -2n^3 + n^2$ satisfies $a_n \sim -2n^3$, since the cubic term dominates
the quadratic one as $n \to \infty$. A more refined classification is obtained by comparing infinities of
essentially different kinds, such as logarithms, powers, exponentials, and factorials. The following collection
of special limits, which we state without proof, provides a hierarchy of increasingly rapid growth. For every
choice of real constants $\alpha, \beta > 0$ and $b > 1$, one has
\[
\begin{gathered}
\lim_{n \to \infty} \frac{(\ln n)^\beta}{n^\alpha} = 0, \qquad
\lim_{n \to \infty} \frac{n^\alpha}{b^n} = 0, \\
\lim_{n \to +\infty} \frac{b^n}{n!} = 0, \qquad
\lim_{n \to +\infty} \frac{n!}{n^n} = 0
\end{gathered}
\]

These four special limits together yield the following hierarchy of infinities, ordered from the slowest to the
fastest growing:
\[
(\ln n)^\beta \prec n^\alpha \prec b^n \prec n! \prec n^n
\]

An entirely dual theory applies to sequences that tend to zero. We call a sequence $\{\varepsilon_n\}$
infinitesimal if \[
\lim_{n \to \infty} \varepsilon_n = 0;
\] a natural question, in analogy with the case of
infinities, is which of two infinitesimals, for instance $1/n$ and $1/n^2$, tends to zero more rapidly. A useful
family of infinitesimal sequences is generated by composing an elementary function with an infinitesimal argument,
as in the sequences $\ln(1 + 1/n)$, $\sin(1/n^2)$, and $e^{1/n} - 1$. The fact that these sequences are indeed
infinitesimal follows from the following general rule.

\begin{theorem}[Composition Rule]
Let $f$ be continuous at $a$ relative to its domain. If $a_n\to a$ and
$a_n\in\operatorname{Dom}(f)$ eventually, then $f(a_n)\to f(a)$.
\end{theorem}

In particular, this rule guarantees that $\sin a_n \to \sin a$ and
$e^{a_n}\to e^a$. For logarithms it requires $a>0$ and $a_n>0$
eventually. For arbitrary real $\alpha$, the representation
$x^\alpha=e^{\alpha\ln x}$ gives $a_n^\alpha\to a^\alpha$ when $a>0$
and $a_n>0$ eventually; integer and some rational exponents admit larger
real domains. Just as for infinities, two
infinitesimal sequences $a_n \to 0$ and $b_n \to 0$ are compared by examining the limit of their quotient
\[
\lim_{n \to \infty} a_n / b_n.
\]

If this limit equals zero, we say that $a_n$ is an infinitesimal of higher order
than $b_n$; if the limit equals infinity, we say that $a_n$ is an infinitesimal of lower order than $b_n$.
As before, a special case deserves its own definition.

\begin{definition}
If \[
\lim_{n \to \infty} a_n / b_n = 1,
\] we say that $a_n$ is asymptotically equivalent to $b_n$, and we write
$a_n \sim b_n$. In this case $a_n$ and $b_n$ are said to be infinitesimals of the same order.
\end{definition}

A number of special limits, again stated without proof and valid for an
arbitrary infinitesimal sequence $\{\varepsilon_n\}$, establish that
certain elementary infinitesimals are asymptotically equivalent to
$\varepsilon_n$ itself. For every $\alpha>0$, the complete family is
\[
\begin{aligned}
\lim_{n \to \infty} \frac{\ln(1 + \varepsilon_n)}{\varepsilon_n} = 1
&\ \longleftrightarrow\ \ln(1 + \varepsilon_n) \sim \varepsilon_n, \\
\lim_{n \to \infty} \frac{e^{\varepsilon_n} - 1}{\varepsilon_n} = 1
&\ \longleftrightarrow\ e^{\varepsilon_n} - 1 \sim \varepsilon_n, \\
\lim_{n \to \infty}
\frac{(1 + \varepsilon_n)^\alpha - 1}{\varepsilon_n} = \alpha
&\ \longleftrightarrow\
(1 + \varepsilon_n)^\alpha - 1 \sim \alpha\varepsilon_n, \\
\lim_{n \to \infty} \frac{\sin \varepsilon_n}{\varepsilon_n} = 1
&\ \longleftrightarrow\ \sin \varepsilon_n \sim \varepsilon_n, \\
\lim_{n \to \infty} \frac{1 - \cos \varepsilon_n}{(\varepsilon_n)^2} = \frac{1}{2}
&\ \longleftrightarrow\
1 - \cos \varepsilon_n \sim \frac{1}{2}\varepsilon_n^2.
\end{aligned}
\]

\subsection{Comparison Theorems for Limits of Sequences}

The comparison of sequences with respect to order of magnitude, discussed above, relies on a more fundamental
set of results concerning the compatibility of limits with inequalities. The first such result concerns the sign
of the terms of a convergent sequence.

\begin{theorem}[Sign Preserving Property]
Let $a_n \to L$. If $L > 0$, then $a_n > 0$ for every $n$ sufficiently large. Conversely, if $a_n > 0$ for every
$n$ sufficiently large, then $L \geq 0$.
\end{theorem}

A second fundamental result allows one to pass inequalities between sequences to the limit.

\begin{theorem}[Comparison Theorem I]
If $a_n \geq b_n$ for every $n$ and both sequences possess a limit, then
\[
\lim_{n \to \infty} a_n \geq \lim_{n \to \infty} b_n
\]
In particular,
\[
\begin{gathered}
\lim_{n \to \infty} b_n = +\infty\quad\text{implies}\quad \lim_{n \to \infty} a_n = +\infty, \\
\lim_{n \to \infty} a_n = -\infty\quad\text{implies}\quad \lim_{n \to \infty} b_n = -\infty.
\end{gathered}
\]

\end{theorem}

A third comparison result, often called the Squeeze Theorem, is particularly useful when the limit of a sequence
cannot be computed directly but can be bounded between two sequences with a common, known limit.

\begin{theorem}[Comparison Theorem II, Squeeze Theorem]
Suppose the following inequalities hold for every $n$ sufficiently large:
\[
c_n \leq b_n \leq a_n.
\]
If also $a_n \to L$ and $c_n \to L$, then $b_n \to L$.
\end{theorem}

An immediate and frequently used consequence of the Squeeze Theorem is the following.

\begin{corollary}
If $a_n \to 0$ and $b_n$ is bounded, then $a_n \cdot b_n \to 0$.
\end{corollary}

The proofs of the Sign Preserving Property and of the above corollary are left to the reader as instructive
exercises.

\subsection{Algebraic Limit Laws}

Having established the qualitative comparison theorems, we now turn to the quantitative problem of computing
limits explicitly. The following theorem collects the elementary algebraic operations under which convergence
is preserved.

\begin{theorem}[Algebraic Limit Laws]
Let $\{a_n\}$ and $\{b_n\}$ be two convergent sequences, with $a_n \to a$ and $b_n \to b$. Then
\[
a_n + b_n \to a + b, \qquad a_n \cdot b_n \to a \cdot b, \qquad \frac{a_n}{b_n} \to \frac{a}{b} \ \text{ if } b \neq 0
\]

\end{theorem}

We present the proof of the product law, as it illustrates a technique of general usefulness.

\begin{proof}[Proof of the product law]
We must show that $|a_n b_n - ab| \to 0$. Adding and subtracting the term $a b_n$, we obtain
\[
|a_n b_n - ab| = |a_n b_n - a b_n + a b_n - ab| \leq |b_n| \, |a_n - a| + |a| \, |b_n - b|
\]

Since $\{b_n\}$ is convergent, it is bounded, so there exists $M \geq 1$ such that $|b_n| \leq M$ for every
$n \in \mathbb{N}$. Hence
\[
|a_n b_n - ab| \leq M |a_n - a| + |a| \, |b_n - b|
\]

Fix $\varepsilon > 0$. By the convergence $a_n \to a$, there exists $N_1$ such that
$|a_n - a| < \varepsilon / (2M)$ for every $n > N_1$. Likewise, by the convergence $b_n \to b$, there exists
$N_2$ such that
\[
|b_n-b|<\frac{\varepsilon}{2(|a|+1)}
\]
 for every $n>N_2$. This choice is valid also when $a=0$. Combining these two estimates, we obtain
$|a_n b_n - ab| < \varepsilon$ for every $n > \max\{N_1, N_2\}$, which proves the claim.
\end{proof}

\subsection{The Algebra of Infinity and Indeterminate Forms}

The algebraic limit laws stated above apply, strictly speaking, only to finite limits, and one may naturally ask
how similar operations behave when one or both of the sequences diverge to infinity. The answer must be handled
with considerable care, since some conclusions remain valid while others are entirely unjustified. The following
list records several valid rules. They are shorthand for statements about
limits, not arithmetic identities involving a number called infinity. For
a finite real number $a$, one has
\[
a + \infty = +\infty, \qquad a - \infty = -\infty, \qquad +\infty + \infty = +\infty, \qquad -\infty - \infty = -\infty
\]
For products, the sign must be retained:
\[
\begin{aligned}
a>0 &: \quad a(+\infty)=+\infty,\qquad a(-\infty)=-\infty,\\
a<0 &: \quad a(+\infty)=-\infty,\qquad a(-\infty)=+\infty,\\
a\in\mathbb{R} &: \quad \frac{a}{+\infty}
=\frac{a}{-\infty}=0.
\end{aligned}
\]

The shorthand notation
\[
a + \infty = +\infty
\]
 should be understood in the following operational sense: if
$a_n \to a$ and $b_n \to +\infty$, then $a_n + b_n \to +\infty$, and analogous interpretations apply to the
remaining rules. By contrast, certain combinations of infinite or infinitesimal limits do not determine the
resulting limit uniquely, and are accordingly called indeterminate forms. These are
\[
+\infty - \infty, \qquad 0 \cdot \infty, \qquad \frac{0}{0}, \qquad \frac{\infty}{\infty}
\]

In each of these situations, essentially any outcome is possible depending on the specific sequences involved,
as illustrated by the three sequences
\[
b_n = \frac{n^3}{n^2 + 1}, \qquad c_n = \frac{1 + 4n^2}{5 + 3n + n^2}, \qquad d_n = \frac{4 + n^2}{3n - n^5}.
\]
All of these expressions present the indeterminate form $\infty / \infty$ yet exhibit different limiting behaviors.
A useful heuristic for resolving such indeterminate forms consists of identifying, within each expression, the
dominant term, that is, the infinity of highest order, and factoring it out; the reader is encouraged to practice
this technique on limits such as
\[
\lim_{n \to \infty} \frac{e^{1/n} - 1}{\sin(2/n)}, \qquad
\lim_{n \to \infty} \frac{\ln(1 + 1/n)}{n \sin(1/n^2)}, \qquad
\lim_{n \to \infty} \frac{1 - \cos(1/n)}{1/n^2}
\]
which involve infinitesimal sequences and can be resolved using the special limits collected in the previous
subsection. A systematic and powerful method for handling such indeterminate quotients and products is furnished
by the following substitution principle, which allows one to replace each factor by any sequence to which it is
asymptotically equivalent.

An expression such as $a/0$ is not assigned an infinite value. Its
one-sided behavior depends on the signs of both numerator and denominator;
for example, $1/x\to+\infty$ as $x\to0^+$ but
$1/x\to-\infty$ as $x\to0^-$. Thus a two-sided limit need not exist.

\begin{theorem}[Substitution Principle]
Assume that all quotients below are eventually defined. If $a_n \sim A_n$, $b_n \sim B_n$, and $c_n \sim C_n$, then
\[
\frac{a_n \cdot b_n}{c_n} \sim \frac{A_n \cdot B_n}{C_n}.
\]
Consequently, whenever either of the following limits exists in the extended real sense, the other exists and has the same value:
\[
\lim_{n \to \infty} \frac{a_n \cdot b_n}{c_n},
\qquad
\lim_{n \to \infty} \frac{A_n \cdot B_n}{C_n}.
\]

\end{theorem}

The reader is invited to revisit the exercises considered so far and to solve them anew by exploiting this
substitution principle. Besides the indeterminate forms already discussed, further indeterminate forms arise
when studying limits of the type \[
\lim_{n \to \infty} a_n^{b_n},
\] namely the forms $1^{\infty}$, $0^0$, and
$\infty^0$. To handle these cases, one takes the logarithm, writing
\[
\begin{aligned}
a_n^{b_n} &= e^{\ln a_n^{b_n}} \\
&= e^{b_n \ln a_n}.
\end{aligned}
\]
The exponent $b_n \ln a_n$ itself reduces to one of the previously discussed indeterminate
forms: if $a_n \to 1$ and $b_n \to \infty$, then $b_n \ln a_n$ has the form $\infty \cdot 0$; if $a_n \to 0$ and
$b_n \to 0$, then $b_n \ln a_n$ has the form $0 \cdot (-\infty)$; and if $a_n \to \infty$ and $b_n \to 0$, then
$b_n \ln a_n$ has the form $0 \cdot \infty$. Once it has been shown that $b_n \ln a_n \to L$, possibly with
$L = \pm \infty$, one concludes that
\[
\begin{aligned}
\lim_{n \to \infty} a_n^{b_n} &= \lim_{n \to \infty} e^{b_n \ln a_n} \\
&= e^L
\end{aligned}
\]

A particularly important special limit of the form $1^\infty$ generalizes the definition of Euler's number
encountered earlier: if $b_n \to \infty$, then
\[
\lim_{n \to \infty} \left( 1 + \frac{1}{b_n} \right)^{b_n} = e
\]

Summarizing the discussion of this section, the principal tools available for resolving indeterminate forms in
the computation of limits of sequences are algebraic manipulation and simplification, the hierarchy of infinities,
the collection of special limits, and the substitution principle based on asymptotic equivalence.

We close this section by proving one of the special limits invoked above, namely $\sin \varepsilon_n / \varepsilon_n \to 1$
for an infinitesimal sequence $\varepsilon_n \to 0$.

\begin{proof}
For every sufficiently large $n$ with $\varepsilon_n\neq0$, put
$\theta_n=|\varepsilon_n|\in(0,\pi/2)$. The following inequalities hold:
\[
\begin{aligned}
\sin \theta_n &\leq \theta_n \\
&\leq \tan \theta_n
\end{aligned}
\]
Dividing throughout by $\sin \theta_n$, which is positive, we obtain
\[
\begin{aligned}
1 &\leq \frac{\theta_n}{\sin \theta_n} \\
&\leq \frac{\tan \theta_n}{\sin \theta_n} \\
&= \frac{1}{\cos \theta_n}
\end{aligned}
\]
Taking reciprocals reverses the inequalities and yields
\[
\begin{aligned}
\cos \theta_n &\leq \frac{\sin \theta_n}{\theta_n} \\
&\leq 1
\end{aligned}
\]

Since sine is odd,
$\sin\varepsilon_n/\varepsilon_n=\sin\theta_n/\theta_n$. Moreover,
$\cos\theta_n\to1$, so the Squeeze Theorem gives the claim.
\end{proof}

\section{Infinite Series and Partial Sums}

\subsection{From Infinite Decimals to Infinite Series}

The theory of series arises naturally from a question that is implicit in ordinary decimal notation: what
precisely do we mean when we express a number as an infinite decimal expansion? For instance, what does it mean
to write
\[
\pi = 3.14159265358979323846264338327950288 \cdots?
\]

The convention underlying decimal notation is that any real number can be represented as an infinite sum, so
that
\[
\pi = 3 + \frac{1}{10} + \frac{4}{10^2} + \frac{1}{10^3} + \frac{5}{10^4} + \cdots,
\] where the ellipsis indicates that the sum continues indefinitely, and where the more terms one adds, the closer
the resulting partial sum approaches the true value of $\pi$. This idea of summing infinitely many terms is
formalized by the notion of an infinite series. If one attempts to add the terms of an infinite sequence
$\{a_n\}$, one obtains an expression of the form
\[
a_0 + a_1 + a_2 + a_3 + \cdots + a_n + \cdots
\]
which is called an infinite series, or simply a series, and is denoted for brevity by the symbol
\[
\sum_{n=0}^{\infty} a_n.
\]

It is important to emphasize that this expression is, at this stage, merely a formal
symbol and not yet a rigorous definition of what it means for the series to have a value; the precise definition requires the auxiliary notion of partial sum.

\begin{definition}[Partial Sum]
Given a sequence $\{a_n\}_{n \in \mathbb{N}}$, its $k$-th partial sum is defined by
\[
s_k = a_0 + a_1 + a_2 + \cdots + a_k.
\]
Thus $s_k$ is the sum of the terms $a_i$ up to index $k$. The resulting sequence $\{s_k\}_{k \in \mathbb{N}}$ is called the sequence of partial sums associated with $\{a_n\}$.
\end{definition}

With this notion in hand, the series of general term $a_n$ can be rigorously defined as the limit of its sequence of partial sums.

\begin{definition}[Series]
The series of term $a_n$, denoted by
\[
a_0 + a_1 + a_2 + \cdots + a_n + \cdots = \sum_{n=0}^{\infty} a_n,
\]
is defined as the limit of the sequence of its partial sums,
\[
\begin{aligned}
\sum_{n=0}^{\infty} a_n &:= \lim_{k \to \infty} s_k \\
&= \lim_{k \to \infty} \sum_{n=0}^{k} a_n.
\end{aligned}
\]

\end{definition}

\subsection{Geometric and Basel Series}

Just as for sequences, the limiting behavior of the sequence of partial sums determines the convergence behavior of a series.

\begin{definition}[Convergence and Divergence of a Series]
If the limit \[
\lim_{k \to \infty} s_k
\] exists and is finite, we say that the series \[
\sum_{n=0}^{\infty} a_n
\] is convergent. If \[
\lim_{k \to \infty} s_k = +\infty
\] or $-\infty$, we say that the series diverges to the corresponding infinity. If the limit does not exist, the series is also called divergent.
\end{definition}

The paradigmatic example illustrating this classification is the geometric series.

\begin{definition}[Geometric Series]
Given a real number $q$, the series \[
\sum_{n=0}^{\infty} q^n
\] is called the geometric series of common ratio $q$.
\end{definition}

The convergence behavior of a geometric series depends entirely on the value of $q$:
\[
\begin{cases}
\text{divergent} & \text{if } q \geq 1, \\
\text{convergent} & \text{if } |q| < 1, \\
\text{divergent} & \text{if } q \leq -1
\end{cases}.
\]
Moreover, whenever $|q| < 1$, the sum of the series admits the explicit closed-form expression
\[
\sum_{n=0}^{\infty} q^n = \frac{1}{1 - q}.
\]

This result follows from the classical formula for the partial sums of a geometric progression: for every $q \neq 1$, one has
\[
\sum_{n=0}^{k} q^n = \frac{1 - q^{k+1}}{1 - q}.
\]

The value of the series is then obtained by letting $k \to \infty$ in this explicit expression for the partial sums, using the results on the limit of $q^n$ established in the previous section.

Having established this concrete example, the general goal of the theory of series becomes clear: given an arbitrary sequence $\{a_n\}$, one wishes to determine the convergence behavior of the associated series
\[
\sum_{n=0}^{\infty} a_n,
\]
that is, to decide whether it is convergent or divergent, and, if it converges, to compute or at least estimate its sum.


A celebrated example illustrating the subtlety of this problem is the Basel series, defined as the sum of the reciprocals of the squares of all positive integers,
\[
\sum_{n=1}^{\infty} \frac{1}{n^2} = 1 + \frac{1}{4} + \frac{1}{9} + \cdots + \frac{1}{n^2} + \cdots.
\]

It will be shown below that this series is convergent; however, unlike the geometric series, no explicit closed-form formula is available for its partial sums, which raises the question of how one might determine the exact value of the sum. Computing the exact value of a convergent series can indeed be a very difficult problem in general. It is a celebrated fact, established by Euler in 1735 in his solution to the so-called Basel Problem, that
\[
\sum_{n=1}^{\infty} \frac{1}{n^2} = \frac{\pi^2}{6}.
\]
More generally, Euler obtained an explicit formula for the sum \[
\sum_{n=1}^{\infty} 1/n^\alpha =: \zeta(\alpha)
\] whenever the exponent $\alpha$ is a positive even integer. In contrast, for odd integers $\alpha>1$, no general closed-form formula in terms of familiar elementary constants is known. In particular, $\zeta(3)$, called Apéry's constant, has been known to be irrational since 1978.

\subsection{Necessary Conditions, Finite Changes, and Algebraic Operations}

Our principal objective in the remainder of this section is to develop systematic techniques for determining the convergence behavior of a given series without necessarily computing its exact sum. A first useful observation is that altering a finite number of terms of a series does not affect its convergence or divergence, since \[
\sum_{n=0}^{\infty} a_n = a_0 + a_1 + \cdots + a_N + \sum_{n=N+1}^{\infty} a_n
\] for any fixed integer $N$; naturally, the numerical values of the two series appearing on either side of this identity generally differ when both are convergent, even though their convergence behavior is the same. Because of this observation, when one is interested only in convergence and not in the precise value of the sum, it is customary to denote the series simply by \[
\sum a_n
\] or
\[
\sum_n a_n,
\]
omitting the starting index. A first necessary condition for convergence, applicable to series with arbitrary terms, is provided by the following theorem.

\begin{theorem}[Necessary Condition for Convergence]
If the series of term $a_n$ is convergent, then
\[
\lim_{n \to \infty} a_n = 0.
\]

\end{theorem}

\begin{proof}
By assumption, there exists $s \in \mathbb{R}$ such that
\[
\lim_{n \to \infty} s_n = s.
\]
Here the $n$-th partial sum is defined by
\[
s_n = a_1 + \cdots + a_{n-1} + a_n.
\]
The partial sum with index $n-1$ is
\[
s_{n-1} = a_1 + \cdots + a_{n-1}.
\]
Shifting the index does not change the limit, so
\[
\lim_{n \to \infty} s_{n-1} = s.
\]
Subtracting the two sequences, we find $s_n - s_{n-1} = a_n$, and hence
\[
\begin{aligned}
\lim_{n \to \infty} a_n &= \lim_{n \to \infty} s_n - \lim_{n \to \infty} s_{n-1} \\
&= s - s \\
&= 0,
\end{aligned}
\]
which proves the theorem.
\end{proof}

\begin{remark}
The condition \[
\lim_{n \to \infty} a_n = 0
\] is necessary but by no means sufficient for the convergence of the series \[
\sum a_n.
\] The harmonic series, discussed below, provides the standard counterexample: its general term $a_n = 1/n$ tends to zero, yet the series \[
\sum 1/n
\] is divergent.
\end{remark}


Convergent series inherit certain algebraic properties from the corresponding limit laws for sequences, as recorded in the following theorem.

\begin{theorem}
If \[
\sum a_n
\] and \[
\sum b_n
\] are convergent series, then so are \[
\sum c \cdot a_n,
\] for any constant $c \in \mathbb{R}$, and \[
\sum (a_n + b_n).
\] Moreover,
\[
\sum_{n=0}^{\infty} c \cdot a_n = c \sum_{n=0}^{\infty} a_n, \qquad \sum_{n=0}^{\infty} (a_n + b_n) = \sum_{n=0}^{\infty} a_n + \sum_{n=0}^{\infty} b_n.
\]

\end{theorem}

These properties follow directly from the corresponding algebraic limit laws for sequences, applied to the associated sequences of partial sums. As a worked application, consider
\[
\sum_{n=1}^{\infty} ( (-1)^n / 3^n + 1/2^n )
\]
Both components are geometric series. Since their ratios are $-1/3$ and
$1/2$, respectively, both converge absolutely, and
\[
\begin{aligned}
\sum_{n=1}^{\infty}\frac{(-1)^n}{3^n}
&=\sum_{n=1}^{\infty}\left(-\frac13\right)^n
=\frac{-1/3}{1-(-1/3)}=-\frac14,\\
\sum_{n=1}^{\infty}\frac1{2^n}
&=\frac{1/2}{1-1/2}=1.
\end{aligned}
\]
Consequently,
\[
\sum_{n=1}^{\infty}\left(\frac{(-1)^n}{3^n}+\frac1{2^n}\right)
=-\frac14+1=\frac34.
\]
Moreover, if \[
\sum a_n
\] converges and \[
\sum b_n
\] diverges, then
\[
\sum(a_n+b_n)
\] must diverge: otherwise subtracting the convergent series
\[
\sum a_n
\] would imply that \[
\sum b_n
\] converges, a contradiction.

\section{Positive Series and Convergence Tests}

\subsection{Harmonic and Generalized Harmonic Series}

We now return to the two celebrated series introduced earlier, providing convergence and divergence proofs that do not require an explicit formula for the partial sums.

\begin{theorem}
The Basel series
\[
\sum_{n=1}^{\infty} 1/n^2 = 1 + 1/4 + 1/9 + \cdots + 1/n^2 + \cdots
\]
is convergent.
\end{theorem}

\begin{proof}
Although the partial sums of this series cannot be computed explicitly, they can be compared with those of a suitable telescopic series. For every $k \geq 2$ we have
\[
s_k = \sum_{n=1}^{k} \frac{1}{n^2} < 1 + \sum_{n=2}^{k} \frac{1}{n(n-1)} = 1 + \sum_{n=2}^{k} \left( \frac{1}{n-1} - \frac{1}{n} \right).
\]
The sum on the right-hand side telescopes, so that
\[
1 + \left(1 - \frac{1}{2}\right) + \left(\frac{1}{2} - \frac{1}{3}\right) + \left(\frac{1}{3} - \frac{1}{4}\right) + \cdots + \left(\frac{1}{k-1} - \frac{1}{k}\right) = 1 + 1 - \frac{1}{k}.
\]

Together with $s_1=1<2$, this shows that $s_k<2$ for every $k\geq1$, so that the sequence $\{s_k\}$ is bounded. Since $\{s_k\}$ is also increasing, being the sequence of partial sums of a series with positive terms, the Monotone Convergence Theorem guarantees that it possesses a finite limit; by definition of series, this proves that the Basel series is convergent.
\end{proof}

\begin{theorem}
The harmonic series
\[
\sum_{n=1}^{\infty} 1/n = 1 + 1/2 + 1/3 + \cdots + 1/n + \cdots
\]
is divergent.
\end{theorem}

\begin{proof}
As with the Basel series, the partial sums of the harmonic series cannot be computed explicitly, but they can be compared with those of a computable divergent series. Using the elementary inequality $x \geq \ln(1 + x)$, valid for every $x>-1$, we estimate the $k$-th partial sum from below as
\[
\begin{aligned}
s_k &= \sum_{n=1}^{k} \frac{1}{n} \\
&\geq \sum_{n=1}^{k} \ln \left(1 + \frac{1}{n}\right) \\
&= \sum_{n=1}^{k} \big( \ln(n+1) - \ln n \big).
\end{aligned}
\]
The right-hand side is again a telescopic sum,
\[
\begin{aligned}
\sum_{n=1}^{k} \big( \ln(n+1) - \ln n \big) &= (\ln 2 - \ln 1) + (\ln 3 - \ln 2) + \cdots + \big(\ln(k+1) - \ln k\big) \\
&= \ln(k+1).
\end{aligned}
\]
Hence $s_k \geq \ln(k+1)$ for every $k$. Since
\[
\lim_{k \to \infty} \ln(k+1) = \infty,
\]
it follows that
\[
\lim_{k \to \infty} s_k = \infty,
\]
which proves that the harmonic series is divergent.
\end{proof}

The harmonic series is the special case $\alpha = 1$ of a broader family known as the generalized harmonic series, defined for a real parameter $\alpha$ as
\[
\sum_{n=1}^{\infty} 1/n^\alpha.
\]

This series is convergent for $\alpha > 1$ and divergent for $\alpha \leq 1$; the cases $\alpha > 2$ and $\alpha < 1$ follow directly from the Comparison Test discussed in the next subsection, while the boundary case $\alpha = 1$ is precisely the harmonic series treated above.

The divergence of the harmonic series, despite the fact that its general term $a_n = 1/n$ tends to zero, provides the promised counterexample confirming that the necessary condition for convergence established earlier is not sufficient: although \[
\lim_{n \to \infty} 1/n = 0,
\] the series \[
\sum 1/n
\] is nonetheless divergent.

\subsection{Comparison and Asymptotic Comparison Tests}

A series is called a positive series if its general term satisfies $a_n > 0$ for every $n$. Positive series enjoy a particularly simple qualitative property, recorded in the following theorem.

\begin{theorem}
A positive series either converges to a finite value or diverges to $+\infty$.
\end{theorem}

\begin{proof}
The sequence $s_n$ of partial sums of a positive series is monotone increasing, since $s_{n+1} = s_n + a_{n+1}$ and $a_{n+1} > 0$ implies $s_{n+1} > s_n$ for every integer $n$. By the Monotone Convergence Theorem, $s_n$ has a finite limit if it is bounded above and tends to $+\infty$ otherwise, which proves the theorem.
\end{proof}

A first practical criterion, useful for quickly identifying divergence, is the following.

\begin{theorem}[Divergence Test for Positive Series]
If $a_n > 0$ eventually and \[
\lim_{n \to \infty} a_n \neq 0,
\] then the series \[
\sum a_n
\] diverges.
\end{theorem}

As an application, one may consider the series
\[
\sum_{n=1}^{\infty} (4n+5)/(3n+1),
\]
whose general term does not tend to zero and which is therefore divergent by the Divergence Test.

A more refined tool for determining the convergence behavior of a positive series by comparison with another, better-understood series is the Comparison Test.

\begin{theorem}[Comparison Test]
Given two positive sequences $\{a_n\}$ and $\{b_n\}$ such that $a_n \leq b_n$ for every $n$ sufficiently large, the following implications hold: if \[
\sum a_n
\] is divergent, then \[
\sum b_n
\] is divergent; and if \[
\sum b_n
\] is convergent, then \[
\sum a_n
\] is convergent.
\end{theorem}

\begin{proof}
By assumption, $a_n \leq b_n$ for every $n$, and consequently \[
\sum_{n=1}^{k} a_n \leq \sum_{n=1}^{k} b_n
\] for every $k$. Passing to the limit as $k \to \infty$ and invoking the Comparison Theorem for limits of sequences, we obtain
\[
\lim_{k \to \infty} \sum_{n=1}^{k} a_n \leq \lim_{k \to \infty} \sum_{n=1}^{k} b_n.
\]
If \[
\sum a_n
\] is known to be divergent, then \[
\lim_{k \to \infty} \sum_{n=1}^{k} a_n = +\infty,
\] and the inequality above forces \[
\lim_{k \to \infty} \sum_{n=1}^{k} b_n = +\infty
\] as well, so that \[
\sum b_n
\] is divergent too. The second assertion of the theorem is proved analogously.
\end{proof}

As an application of the Comparison Test, one may consider the series
\[
\sum_{n=1}^{\infty} 1/(n^4 + n - 1),
\]
whose convergence follows by comparison with the convergent generalized harmonic series of exponent four.

A variant of the Comparison Test, often more convenient in practice, replaces the pointwise inequality between the two sequences with an asymptotic equivalence.

\begin{theorem}[Asymptotic Comparison Test]
Given two positive sequences $\{a_n\}$ and $\{b_n\}$ such that $a_n \sim b_n$ as $n \to \infty$, the series \[
\sum a_n
\] and \[
\sum b_n
\] have the same convergence behavior: either both converge or both diverge.
\end{theorem}

\begin{proof}
The asymptotic equivalence gives the limit
\[
\lim_{n \to \infty} a_n / b_n = 1
\]
Consequently, there exists an integer $N$ such that, for every $n>N$,
\[
\frac{1}{2} < \frac{a_n}{b_n} < 2.
\]
Equivalently, $(1/2)b_n<a_n<2b_n$. We now invoke the ordinary Comparison Test twice. If \[
\sum b_n
\] converges, then so does \[
\sum 2 b_n,
\] and hence \[
\sum a_n
\] converges by the Comparison Test, since $a_n < 2 b_n$. Conversely, if \[
\sum b_n
\] diverges, then so does \[
\sum (1/2) b_n,
\] and hence \[
\sum a_n
\] diverges by the Comparison Test, since $a_n > (1/2) b_n$. This establishes the equivalence of their convergence behavior.
\end{proof}

As an application, one may consider the series \[
\sum_{n=1}^{\infty} (4n+5)/(n^4 + 3n + 1),
\] whose general term is asymptotically equivalent to $4/n^3$, from which convergence follows immediately by the Asymptotic Comparison Test together with the convergence of the generalized harmonic series of exponent three. It is worth noting that, beyond the comparison-based criteria discussed here, additional tests are available for positive series, most notably the Ratio Test and the Root Test, which will not be developed further in the present text.

\section{Absolute and Alternating Convergence}

\subsection{Series with Terms of Arbitrary Sign}

The theory developed so far applies specifically to positive series, for which the sequence of partial sums is automatically monotone. When the general term $a_n$ is allowed to have arbitrary sign, this monotonicity is lost. The archetypal example is the series \[
\sum_{n=0}^{\infty} (-1)^n,
\] whose partial sums satisfy $s_{2n} = 1$ and $s_{2n+1} = 0$ for every integer $n \geq 0$; since the sequence $\{s_n\}$ oscillates between these two values without converging, the series diverges. In general, when $a_n$ has arbitrary sign, the sequence of partial sums need not be monotone and may fail to converge.

To extend the comparison-based techniques developed for positive series to this more general setting, it is natural to associate with an arbitrary series \[
\sum a_n
\] the auxiliary positive series \[
\sum |a_n|,
\] obtained by taking absolute values of the terms. This leads to a stronger notion of convergence.

\begin{definition}[Absolute Convergence]
We say that the series \[
\sum a_n
\] is absolutely convergent if the series of absolute values \[
\sum |a_n|
\] is convergent.
\end{definition}

\begin{theorem}
If a series is absolutely convergent, then it is convergent.
\end{theorem}

\begin{proof}
Observe first that
\[
\begin{aligned}
0 &\leq a_n + |a_n| \\
&\leq 2 |a_n|.
\end{aligned}
\]
This holds for every $n$. If \[
\sum a_n
\] is absolutely convergent, then \[
\sum |a_n|
\] is convergent, and hence \[
\sum 2 |a_n|
\] is convergent as well. By the Comparison Test, applied to the positive series \[
\sum (a_n + |a_n|),
\] we conclude that this series is convergent. Finally, writing \[
\sum_n a_n = \sum_n (a_n + |a_n|) - \sum_n |a_n|
\] as the difference of two convergent series, we conclude that \[
\sum a_n
\] is itself convergent, as required.
\end{proof}

This theorem provides a powerful sufficient criterion for the convergence of series with arbitrary sign, applicable for instance to
\[
\sum_n (-1)^n n^2 / n^n
\]
and to
\[
\sum_n \sin n / n^2,
\]
both of which can be shown to be absolutely convergent.

\subsection{Alternating Series and the Leibniz Test}

A particularly important class of series with terms of arbitrary sign is provided by the alternating series. The alternating harmonic series,
\[
\sum_{n=1}^{\infty} (-1)^{n-1} \frac{1}{n} = 1 - \frac{1}{2} + \frac{1}{3} - \frac{1}{4} + \cdots,
\]
furnishes a natural example: it is not absolutely convergent, since the associated series of absolute values is precisely the divergent harmonic series, yet the question of whether the alternating series itself converges remains open at this point. In general, an alternating series is defined as follows.

\begin{definition}
A series of the form \[
\sum_{n=1}^{\infty} (-1)^{n-1} b_n = b_1 - b_2 + b_3 - \cdots,
\] where every term $b_n$ is positive, is called an alternating series.
\end{definition}

The convergence of alternating series is governed by a simple and elegant criterion, known as the Alternating Series Test or Leibniz's Test.

\begin{theorem}[Leibniz's Test]
Assume that the alternating series \[
\sum_{n=1}^{\infty} (-1)^{n-1} b_n = b_1 - b_2 + b_3 - b_4 + \cdots
\] satisfies the following two conditions: first, \[
\lim_{n \to \infty} b_n = 0;
\] second, the sequence $\{b_n\}$ decreases monotonically. Then the alternating series converges.
\end{theorem}

As a direct application of this criterion, the alternating harmonic series \[
\sum_{n=1}^{\infty} (-1)^{n-1} / n = 1 - 1/2 + 1/3 - 1/4 + \cdots
\] is convergent, since its general term $b_n = 1/n$ decreases monotonically to zero. The reader is encouraged to apply Leibniz's Test to establish the convergence of the series
\[
\begin{gathered}
\sum_{n=1}^{\infty} (-1)^{n-1} \, 2n/(n+4) \\
\sum_{n=1}^{\infty} (-1)^n \, n^2/(n^3+1),
\end{gathered}
\]
and to consider whether the series
\[
\sum_{n=1}^{\infty} (-1)^n (2 + (-1)^n)/n
\]
can be analyzed by means of Leibniz's Test, and, more fundamentally, whether this series is convergent at all, since one of the two hypotheses of the test may fail to hold.

\subsection{Estimating the Sum of an Alternating Series}

A partial sum of any convergent series can, in principle, be used as an approximation to the value of the entire sum; however, such an approximation is of little practical use unless one can also estimate its accuracy. For alternating series satisfying the hypotheses of Leibniz's Test, such an estimate is available in a remarkably simple form.

\begin{theorem}[Estimation Theorem]
Under the hypotheses of Leibniz's Test, let $S$ denote the sum of the series
\[
\sum_{n=1}^{\infty} (-1)^{n-1} b_n.
\]
Then the $k$-th partial sum \[
S_k = \sum_{n=1}^{k} (-1)^{n-1} b_n
\] approximates $S$ with an error $E_k$ bounded by the magnitude of the first omitted term, that is,
\[
\begin{aligned}
E_k &= |S_k - S| \\
&\leq b_{k+1}.
\end{aligned}
\]

\end{theorem}

As an application of this estimation theorem, consider \[
\sum_{n=0}^{\infty} (-1)^n / n!,
\] which equals $e^{-1}$. To guarantee an error below $5\times10^{-4}$,
which is sufficient for correct rounding to three decimal places, choose
$k$ so that $b_{k+1} = 1/(k+1)!$ is smaller than $5\times10^{-4}$. The first such convenient choice is
$k=6$, because
\[
\frac1{7!}=\frac1{5040}\approx1.98\times10^{-4}<5\times10^{-4}.
\]
The corresponding partial sum is
\[
\begin{aligned}
S_6
&=1-1+\frac1{2!}-\frac1{3!}+\frac1{4!}-\frac1{5!}+\frac1{6!}\\
&=\frac{53}{144}\\
&\approx0.3680556.
\end{aligned}
\]
Since $|e^{-1}-S_6|\leq1/7!$, the required value is therefore
\[
e^{-1}=0.368\quad\text{to three decimal places}.
\]

Figure~\ref{fig:atlas-alternating-bounds} illustrates the nested upper and lower estimates supplied by alternating partial sums.

\begin{figure}[!htbp]
\centering
\begin{tikzpicture}[>=Stealth]
\begin{axis}[width=11cm,height=5.8cm,axis lines=middle,ticklabel style={font=\small},label style={font=\small},legend style={font=\scriptsize,draw=black!20,fill=white},xlabel={$n$},ylabel={$S_n$},xmin=0,xmax=15,ymin=0.45,ymax=1.06,xtick={2,4,6,8,10,12,14}]
\addplot[black!65,dashed,domain=0:15] {ln(2)};
\addplot[figred,thick,mark=*] coordinates {(1,1.00000000) (3,0.83333333) (5,0.78333333) (7,0.75952381) (9,0.74563492) (11,0.73654401) (13,0.73013376)};
\addplot[figblue,thick,mark=square*] coordinates {(2,0.50000000) (4,0.58333333) (6,0.61666667) (8,0.63452381) (10,0.64563492) (12,0.65321068) (14,0.65870518)};
\node[above left] at (axis cs:14,0.693147) {$S$};
\end{axis}
\end{tikzpicture}
\caption{Partial sums of the alternating harmonic series bracket its sum
$S=\log 2$. This is a visual case study of the general alternating-series
error mechanism, not the preceding numerical example for
\[
e^{-1}=\sum_{n=0}^{\infty}(-1)^n/n!.
\] Odd partial sums decrease from
above and even partial sums increase from below.}
\label{fig:atlas-alternating-bounds}
\end{figure}

\newpage
\chapter{Topology of the Real Line, Limits, and Continuity}

\section{Real Functions and Their Properties}

\subsection{Definitions, Graphs, and Qualitative Properties}

A real-valued function $f : D \to \mathbb{R}$, with $D \subseteq \mathbb{R}$, is a rule that assigns to each $x \in D$ exactly one value $f(x) \in \mathbb{R}$. The input variable is called the independent variable, and the domain $D = \operatorname{Dom}(f)$ is the set of all admissible inputs. The output is called the dependent variable, and the range of $f$ is the set of all values actually attained by $f$, that is,
\[
\begin{aligned}
\operatorname{Im}(f) &= f(D) \\
&= \{ f(x) : x \in \operatorname{Dom}(f) \} \subset \mathbb{R}.
\end{aligned}
\]

The most common method for visualizing a function $f : D \to \mathbb{R}$ is its graph, defined as the set of ordered input--output pairs $\{ (x, f(x)) : x \in D \}$. Since the graph of a real-valued function is a curve in the $xy$-plane, it is natural to ask which curves in the plane actually arise as graphs of functions. This question is answered by the following geometric criterion.

\begin{theorem}[Vertical Line Test]
A curve in the $xy$-plane is the graph of a function of $x$ if and only if no vertical line intersects the curve more than once.
\end{theorem}


Several qualitative features recur throughout the study of elementary functions, the first of which concerns symmetry. If a function $f$ satisfies $f(-x) = f(x)$ for every number $x$ in its domain, then $f$ is called an even function; if instead $f(-x) = -f(x)$ for every such $x$, then $f$ is called an odd function. This algebraic property has a direct geometric significance: the graph of an even function is symmetric with respect to the $y$-axis, so that once the portion of the graph corresponding to $x > 0$ has been plotted, the entire graph is obtained simply by reflecting this portion about the $y$-axis.

A second important qualitative feature is periodicity. A function $f$ is called periodic if there exists a positive constant $T$ such that $f(x + T) = f(x)$ for every value of $x$ in the domain of $f$; the smallest value of $T$ for which this identity holds is called the period of $f$, and the reciprocal $1/T$, called the frequency, measures the number of occurrences of the repeating pattern per unit of the independent variable. Nature abounds with examples of cyclic or recurring behavior that are naturally modeled by periodic functions, ranging from breathing to the beating of the heart.

A third feature, of a rather different nature, concerns the monotone behavior of a function on a given set. A function $f$ is called strictly increasing on a set $I$ if, for all $x_1, x_2 \in I$ such that $x_1 < x_2$, one has $f(x_1) < f(x_2)$; it is called strictly decreasing on $I$ if instead $f(x_1) > f(x_2)$ whenever $x_1 < x_2$ in $I$.

\subsection{Injectivity and Inverse Functions}

Closely related to monotonicity is the notion of injectivity. A function $f$ is called one-to-one, or injective, if it never takes the same value twice, that is, if $f(x_1) \neq f(x_2)$ whenever $x_1 \neq x_2$. As in the case of the Vertical Line Test, this property admits a simple geometric characterization.

\begin{theorem}[Horizontal Line Test]
A function is one-to-one if and only if no horizontal line intersects its graph more than once.
\end{theorem}

Injective functions play a distinguished role because they alone admit an inverse. Let $f : A \to B$ be a one-to-one function with domain $A$ and range $B$. Then for every $y \in B$ there exists a unique $x \in A$ such that $f(x) = y$, that is,
\begin{equation*}
\forall y \in B \ \exists! \, x \in A : f(x) = y. \tag{1}
\end{equation*}

This means that if an element $y$ of $B$ is used as input, one obtains one and only one output $x$; in this way it becomes possible to construct a new function, called the inverse function of $f$. A function satisfying condition (1) is said to be invertible on $A$, and one has the general equivalence that $f$ is invertible on $A$ if and only if $f$ is injective on $A$. The reader is encouraged to explore further the connections between injectivity and monotonicity, since a strictly monotone function is automatically injective, although the converse implication does not hold in general.

\begin{definition}[Inverse Function]
Let $f$ be an injective function with domain $A$ and range $B$, so that $f : A \to B$. The inverse function of $f$, denoted by $f^{-1}$, is the function from $B$ to $A$,
\[
f^{-1} : B \to A,
\]
defined by the rule that $f^{-1}(y)$ is the unique solution $x$ of the equation $f(x) = y$; equivalently,
\[
f^{-1}(y) = x \iff f(x) = y.
\]

\end{definition}

The graph of the inverse function $f^{-1}$ is obtained from the graph of $f$ by reflecting it about the line $y=x$. For example, starting from
\[
y=x^3+2,
\]
we subtract $2$ and take the real cube root:
\[
y-2=x^3 \quad\Longrightarrow\quad x=\sqrt[3]{y-2}.
\]
Interchanging $x$ and $y$ gives
\[
f^{-1}(x)=\sqrt[3]{x-2}.
\]
The two compositions verify the result:
\[
f(f^{-1}(x))=(\sqrt[3]{x-2})^3+2=x,
\qquad
f^{-1}(f(x))=\sqrt[3]{x^3}=x.
\]

Figure~\ref{fig:atlas-inverse-reflection} relates inverse functions to reflection across the diagonal.

\begin{figure}[!htbp]
\centering
\begin{tikzpicture}[>=Stealth]
\begin{axis}[width=11cm,height=7cm,axis lines=middle,ticklabel style={font=\small},label style={font=\small},legend style={font=\scriptsize,draw=black!20,fill=white},xlabel={$x$},ylabel={$y$},xmin=0,xmax=4.7,ymin=0,ymax=4.7,axis equal image]
\addplot[black!40,dashed,domain=0:4.5] {x};
\addplot[figblue,very thick,domain=0:2.12,samples=100] {x^2};
\addplot[figred,very thick,domain=0:4.5,samples=150] {sqrt(x)};
\addplot[only marks,figblue,mark=*] coordinates {(2,4)};
\addplot[only marks,figred,mark=*] coordinates {(4,2)};
\draw[black!40,dotted] (axis cs:2,4) -- (axis cs:4,2);
\node[left,figblue] at (axis cs:2,4) {$(2,4)$};
\node[below,figred] at (axis cs:4,2) {$(4,2)$};
\end{axis}
\end{tikzpicture}
\caption{The graph of $f(x)=x^2$, restricted to nonnegative inputs, and the graph of its inverse $f^{-1}(x)=\sqrt{x}$ are reflections across $y=x$. The marked points $(2,4)$ and $(4,2)$ exchange their coordinates.}
\label{fig:atlas-inverse-reflection}
\end{figure}

\section{A Catalog of Elementary Functions}

\subsection{Power, Root, and Real-Exponent Functions}

We now turn to a systematic review of the elementary functions most frequently encountered in analysis, beginning with power functions. A function of the form $f(x) = x^n$, where $n \in \mathbb{N}$, is called a power function. The general shape of the graph of $f$ depends critically on whether the exponent $n$ is even or odd: if $n$ is even, then $f$ is an even function and its graph resembles the parabola $y = x^2$, whereas if $n$ is odd, then $f$ is an odd function and its graph resembles that of $y = x^3$. A further useful property, valuable for sketching graphs, is that the greater the power $n$, the flatter the graph near the origin and the steeper the graph for $x$ beyond $1$; in other words, close to the origin the functions with lower powers dominate, while far from the origin the functions with higher powers dominate.

The inverse of the power function $y = x^n$, for $n \in \mathbb{N}$, is called a root function and is denoted by
\[
f(x) = \sqrt[n]{x} = x^{1/n}, \qquad \text{so that} \qquad y = \sqrt[n]{x} \iff y^n = x.
\]

When $n$ is odd, the function $y = x^n$ is one-to-one on all of $\mathbb{R}$, and consequently $\operatorname{Dom}(f) = \mathbb{R}$. When $n$ is even, however, the function $y = x^n$ fails to be injective on $\mathbb{R}$, but it becomes injective once the domain is restricted to $[0, \infty)$; since the corresponding range is likewise $[0, \infty)$, one has $\operatorname{Dom}(f) = [0, \infty)$ in this case.


Having defined $x^n$ for natural exponents $n$ and its inverse for a fixed natural number $n$, it is natural to ask how the expression $b^\alpha$ can be defined for an arbitrary real exponent $\alpha$, given a fixed base $b \in \mathbb{R}$ with $b > 0$. The construction proceeds in successive stages of generality. If $\alpha = -n$ for some $n \in \mathbb{N}$, one sets $b^{-n} := 1/b^n$. If $\alpha = 1/n$, one sets
\[
b^{1/n} := \sqrt[n]{b}.
\]
If $\alpha \in \mathbb{Q}$, that is, if $\alpha = m/n$ with $m \in \mathbb{Z}$ and $n > 0$, one sets
\[
\begin{aligned}
b^{\alpha} &:= (b^m)^{1/n} \\
&= (b^m)^{1/n}.
\end{aligned}
\]

The remaining and most delicate case concerns an irrational exponent $\alpha$, which we illustrate by constructing the value
\[
3^{\,\sqrt{2}}.
\]
The construction proceeds in four steps. First, one writes the decimal expansion
\[
\sqrt{2} = 1.414213562373095048801688724209\ldots.
\]
Second, one constructs the sequence of finite decimal truncations approximating $\sqrt{2}$, namely
\[
1, 1.4, 1.41, 1.414, 1.4142, \ldots,
\]
and defines the corresponding sequence of powers
\[
3^1, 3^{1.4}, 3^{1.41}, 3^{1.414}, 3^{1.4142}, \ldots;
\]
since each exponent appearing in this sequence is rational, every term is already well defined by the earlier stages of the construction, as in the example
\[
\begin{aligned}
3^{1.4} &= 3^{14/10} \\
&= (3^7)^{1/5}.
\end{aligned}
\]
Third, one considers the set
\[
A = \{3^1, 3^{1.4}, 3^{1.41}, 3^{1.414}, 3^{1.4142}, \ldots\}
\]
and observes that $A$ is bounded above, an explicit upper bound being $3^2$. Fourth, by the Axiom of Completeness, the supremum of $A$ exists, and one defines
\[
3^{\,\sqrt{2}} := \sup A.
\]

The same construction applies verbatim to define quantities such as $2^\pi$, $\pi^e$, and $e^\pi$; more generally, for any irrational number $x$ and any base $b > 1$, one sets
\[
b^{\alpha} := \sup \{ b^q : q \in \mathbb{Q}, \ q < \alpha \}.
\]

Once powers with arbitrary real exponents have been defined in this way, they obey the familiar laws of exponents. For any $a, b > 0$ and any $r, s \in \mathbb{R}$, one has
\[
a^0 = 1, \qquad 1^r = 1, \qquad a^{r+s} = a^r a^s, \qquad (ab)^r = a^r b^r, \qquad (a^r)^s = a^{rs}.
\]

Building on this construction, for any real number $\alpha > 0$ one may define the real power function $f(x) = x^\alpha$ with domain $[0, \infty)$. The general shape of the graph when $\alpha > 1$ resembles that of the parabola $y = x^2$ restricted to $x \geq 0$, whereas the graph takes the form of a root function when $\alpha \in (0,1)$. When the exponent has the form $\alpha = m/n$ with $n$ odd, the function can be further extended by symmetry to the whole real line $\mathbb{R}$.

\subsection{Exponential and Logarithmic Functions}

A function of a fundamentally different kind is obtained by placing the variable in the exponent rather than in the base. Given any real number $b > 0$ with $b \neq 1$, the exponential function with base $b$ is defined by $f(x) = b^x$. Its qualitative behavior falls into two distinct families depending on the value of the base: when $0 < b < 1$, the function is decreasing and decays to $0$, whereas when $b > 1$, the function is increasing and grows rapidly to infinity. All such graphs pass through the common point $(0,1)$, since $b^0 = 1$ for every $b \neq 0$; moreover, as the base $b$ increases, the corresponding exponential function grows more rapidly for $x > 0$.

Since, for $b > 0$ and $b \neq 1$, the exponential function $f(x) = b^x$ is injective, it admits an inverse function defined on its range $(0, \infty)$. This inverse is called the logarithmic function with base $b$, denoted by
\[
\log_b x : (0, \infty) \to \mathbb{R}, \qquad y = \log_b x \iff b^y = x \quad (x > 0).
\]

As an immediate consequence of this definition of inverse function, one has the two identities $\log_b(b^x) = x$ for every $x \in \mathbb{R}$, and $b^{\log_b x} = x$ for every $x > 0$. The laws of powers give the following logarithmic identities. In the first four, take $x, y > 0$ and any $\lambda \in \mathbb{R}$; the change-of-base identity also requires $a > 0$ and $a \neq 1$:
\[
\begin{aligned}
\log_b 1 &= 0, \\
\log_b(xy) &= \log_b x + \log_b y, \\
\log_b \frac{x}{y} &= \log_b x - \log_b y, \\
\log_b(x^\lambda) &= \lambda \log_b x, \\
\log_b x &= \frac{\log_a x}{\log_a b}.
\end{aligned}
\]

Among all possible choices of base $b$, the most convenient and mathematically natural choice is Euler's number $e$; the resulting logarithm is called the natural logarithm and is given the special notation $\log_e x = \ln x$. By definition, one thus has $\ln(e^x) = x$ for every $x \in \mathbb{R}$, and $e^{\ln x} = x$ for every $x > 0$. It is also customary to denote the logarithm with basis ten simply by $\log x$, so that $\log_{10} x = \log x$.

Although $\log_b x$, for $b > 1$, is an increasing function, it is worth emphasizing that it grows extremely slowly once $x > 1$; in fact, $\log_b x$ grows more slowly than any positive power $x^\alpha$, however small the exponent $\alpha$ may be. This extremely slow growth can be illustrated by a striking numerical example. Suppose one attempts to compute the value $\ln n$ recursively, starting from $\ln 2$ and proceeding through the chain
\[
\ln 2 \to \ln 3 \to \cdots \to \ln n,
\]
so that reaching $\ln n$ requires $n$ elementary steps. The fastest computers currently available can perform approximately $10^{18}$ such computations per second, and one may ask how long such a computer would take to make the natural logarithm reach the value $100$. Since $\ln x$ reaches the value $100$ precisely when
\[
\begin{aligned}
x &= e^{100} \\
&\approx 2.688\times10^{43},
\end{aligned}
\]
the computer would require approximately
\[
\frac{2.688\times10^{43}}{10^{18}}=2.688\times10^{25}
\]
seconds, or about $8.52\times10^{17}$ years. This duration should be compared with the age of the universe, approximately $14$ billion years; converting to seconds using
\[
14 \times 10^9 \times 365.25 \times 24 \times 3600 \approx 4.42 \times 10^{17}
\]
seconds, the hypothetical computation would take about $6.1\times10^7$ times the age of the universe, vividly illustrating the extraordinary slowness of logarithmic growth.

\subsection{Trigonometric Functions and Their Inverses}

The final family of elementary functions to be reviewed consists of the trigonometric functions, namely $y = \sin x$ and $y = \cos x$, together with the tangent function
\[
\begin{aligned}
y &= \tan x \\
&= \sin x / \cos x,
\end{aligned}
\]
defined for $x \neq \pi/2 + k\pi$ with $k \in \mathbb{Z}$. It is important to recall that in this context the independent variable $x$ represents an angle measured in radians, abbreviated as rad; since $\pi \, \text{rad} = 180^\circ$, one has the approximate conversion $1 \, \text{rad} \approx 57.3^\circ$.

Although the tangent function fails to be injective on its natural domain, it becomes injective once the domain is restricted to the interval $(-\pi/2, \pi/2)$, on which its range is the whole real line $\mathbb{R}$. The inverse of the resulting restriction
\[
\tan : \left( -\frac{\pi}{2}, \frac{\pi}{2} \right) \to \mathbb{R}
\]
is called the arctangent function, denoted by $\arctan$ or $\tan^{-1}$, and is defined as the function
\[
\arctan : \mathbb{R} \to \left( -\frac{\pi}{2}, \frac{\pi}{2} \right), \qquad y = \arctan x \iff \tan y = x.
\]

A similar construction applies to the sine function. Although $\sin x$ is not injective on its natural domain, it becomes injective once restricted to the interval $[-\pi/2, \pi/2]$, on which it defines a one-to-one function onto $[-1, 1]$. The inverse of the restriction
\[
\sin : \left[ -\frac{\pi}{2}, \frac{\pi}{2} \right] \to [-1, 1]
\]
is called the arcsine function, denoted by $\arcsin$ or $\sin^{-1}$, and is defined as the function
\[
\arcsin : [-1, 1] \to \left[ -\frac{\pi}{2}, \frac{\pi}{2} \right], \qquad y = \arcsin x \iff \sin y = x.
\]

An entirely analogous construction, based on restricting the domain of the cosine function to $[0, \pi]$, defines the arccosine function $\arccos x$ as the inverse of $\cos x$ on that restricted domain.

\section{Topology and the Definition of a Limit}

\subsection{Function Limits: Two Motivating Examples}

The theory of limits of functions is best introduced by contrasting two seemingly similar examples. Consider first the function $(\sin x)/x$, which is not defined at $x = 0$, since the denominator vanishes there; nevertheless, numerical experimentation with a calculator strongly suggests that, intuitively,
\[
\lim_{x \to 0} (\sin x)/x = 1.
\]

Consider next the function $\sin(\pi/x)$, which is likewise undefined at $x = 0$; numerical experimentation in this second case suggests instead that the limit
\[
\lim_{x \to 0} \sin(\pi/x)
\]
does not exist at all, since the function oscillates increasingly rapidly as $x$ approaches zero. These two examples illustrate the need for a precise mathematical framework capable of distinguishing between such fundamentally different behaviors, and motivate the systematic development of the theory of limits undertaken in this section.

This development pursues two complementary goals. The first goal concerns limits of a function at infinity: given a function $f : (r, \infty) \to \mathbb{R}$ for some real number $r$, one wishes to define the limit of $f(x)$ as $x$ tends to $+\infty$, that is, as $x$ becomes larger and larger, allowing for the three possible outcomes
\[
\begin{gathered}
\lim_{x \to \infty} f(x) = +\infty, \\
\lim_{x \to \infty} f(x) = L \in \mathbb{R}, \\
\lim_{x \to \infty} f(x) = -\infty.
\end{gathered}
\]

An entirely analogous notion is required when $f : (-\infty, -r) \to \mathbb{R}$, in which case one wishes to define
\[
\lim_{x \to -\infty} f(x);
\]
and when the limiting value $L$ is finite, it is furthermore of interest to distinguish between the limit taken from above and the limit taken from below. The second goal concerns limits of a function at a finite point: given a function
\[
f : (a-r, a) \cup (a, a+r) \to \mathbb{R},
\]
defined in a neighborhood of a point $a$ except possibly at $a$ itself, one wishes to define the limit of $f(x)$ as $x$ approaches $a$, again allowing for the three possible outcomes $+\infty$, a finite value $L$, or $-\infty$; in this setting it is likewise of interest to distinguish the one-sided limits
\[
\begin{gathered}
\lim_{x \to a^+} f(x) \\
\lim_{x \to a^-} f(x).
\end{gathered}
\]

\subsection{Neighbourhoods and Accumulation Points}

In order to formulate a single definition of limit capable of covering all of the cases described above, whether the point approached and the limiting value are finite or infinite, it is convenient to introduce a small amount of topological terminology. We call the extended real line the set
\[
\overline{\mathbb{R}} = \mathbb{R} \cup \{+\infty\} \cup \{-\infty\},
\]
obtained by adjoining to the ordinary real line two ideal points at infinity. A neighborhood of a point $a \in \mathbb{R}$ is an interval of the form
\[
U(a) = (a-r, a+r),
\]
for some $r > 0$. A neighborhood of $+\infty$ is a semi-line of the form $U(+\infty) = (r, +\infty)$, for some $r > 0$, and a neighborhood of $-\infty$ is a semi-line of the form $U(-\infty) = (-\infty, -r)$, for some $r > 0$. Given a point $x_0 \in \overline{\mathbb{R}}$, it is convenient to introduce the notation $\dot{U}(x_0)$ for the corresponding neighborhood with the point $x_0$ itself removed, namely
\[
\dot{U}(x_0) =
\begin{cases}
(a-r, a) \cup (a, a+r) & \text{if } x_0 = a \in \mathbb{R}, \\
(r, +\infty) & \text{if } x_0 = +\infty, \\
(-\infty, -r) & \text{if } x_0 = -\infty.
\end{cases}
\]

Finally, given a finite point $a \in \mathbb{R}$, we call right-neighborhood of $a$ an interval of the form $U^+(a) = (a, a+r)$, for some $r > 0$, and we call left-neighborhood of $a$ an interval of the form $U^-(a) = (a-r, a)$, for some $r > 0$.

\subsection{The Universal Definition of Limit}

Equipped with this topological vocabulary, we can now formulate a single, universal definition of limit that simultaneously covers finite and infinite points, as well as finite and infinite limiting values. The ingredients of the definition are a limiting value $L \in \overline{\mathbb{R}}$, a point $x_0 \in \overline{\mathbb{R}}$ that is approached, and a function $f$ defined in a punctured neighborhood of $x_0$, that is, a function $f : \dot{U}(x_0) \to \mathbb{R}$.

\begin{definition}[Universal Definition of Limit]
We say that the limit of $f(x)$ as $x$ tends to $x_0$ is $L$, and we write
\[
\lim_{x \to x_0} f(x) = L,
\]
if $f(x)$ can be made arbitrarily close to $L$ by taking $x$ sufficiently close to $x_0$, namely if for every neighborhood $U(L)$ of $L$ there exists a neighborhood $U(x_0)$ of $x_0$ such that
\[
x \in \dot{U}(x_0) \implies f(x) \in U(L).
\]

\end{definition}

This single definition, once the appropriate neighborhoods are unpacked according to whether $x_0$ and $L$ are finite or infinite, reproduces every one of the classical $\varepsilon$--$\delta$ and $\varepsilon$--$N$ formulations of limit encountered in elementary analysis, as will be verified explicitly in the following subsections.

Two further refinements of the universal definition are frequently useful. When the limiting value is finite, $L \in \mathbb{R}$, one may define the one-sided limits from above and from below, denoted
\[
\begin{gathered}
\lim_{x \to x_0} f(x) = L^+ \\
\lim_{x \to x_0} f(x) = L^-,
\end{gathered}
\]
by replacing the condition $f(x) \in U(L)$ in the universal definition with $f(x) \in U^+(L)$ and $f(x) \in U^-(L)$, respectively. Analogously, when the point approached is finite, $x_0 \in \mathbb{R}$, one may define the limits from the right and from the left, denoted
\[
\begin{gathered}
\lim_{x \to x_0^+} f(x) = L \\
\lim_{x \to x_0^-} f(x) = L,
\end{gathered}
\]
by replacing the condition $x \in U(x_0)$ with $x \in U^+(x_0)$ and $x \in U^-(x_0)$, respectively.

\section{Limits at Finite and Infinite Points}

\subsection{Finite and Infinite Limits at Infinity}


We now specialize the universal definition to the case of a function $f : (r, \infty) \to \mathbb{R}$, with the aim of describing the behavior of $f(x)$ as $x$ becomes larger and larger, that is, as $x$ tends to $+\infty$. As already anticipated, three qualitatively different outcomes are possible for the limit
\[
\lim_{x \to \infty} f(x),
\]
namely $+\infty$, a finite value $L \in \mathbb{R}$, or $-\infty$; we now make each of these precise.

\begin{definition}[Finite Limit at Infinity]
Let $L \in \mathbb{R}$. We say that the limit of $f(x)$ as $x$ tends to $+\infty$ is $L$, and we write
\[
\lim_{x \to +\infty} f(x) = L,
\]
if the values of $f(x)$ can be made arbitrarily close to $L$ by taking $x$ sufficiently large, that is, if for every $\varepsilon > 0$ there exists $N > 0$ such that
\[
x > N \implies |f(x) - L| < \varepsilon.
\]

\end{definition}

Whenever a function admits a finite limit at $+\infty$, its graph exhibits a characteristic geometric feature, made precise by the following definition.

\begin{definition}[Horizontal Asymptote]
The line $y = L$ is called a horizontal asymptote of the graph of $y = f(x)$ if
\[
\lim_{x \to +\infty} f(x) = L.
\]

\end{definition}

The complementary case, in which the values of $f(x)$ grow without bound, is described by the following definition.

\begin{definition}[Infinite Limit at Infinity]
We say that the limit of $f(x)$ as $x$ tends to $+\infty$ is $+\infty$, and we write
\[
\lim_{x \to +\infty} f(x) = +\infty,
\]
if the values of $f(x)$ can be made arbitrarily large by taking $x$ sufficiently large, that is, if for every $M > 0$ there exists $N > 0$ such that $x > N \implies f(x) > M$.
\end{definition}

Among functions that diverge to $+\infty$, it is useful to distinguish those that grow faster than a linear function from those that grow more slowly. If
\[
\lim_{x \to +\infty} f(x)/x = +\infty,
\]
we say that $f$ is superlinear as $x \to +\infty$; the function $f(x) = x^2$ provides a standard example. If instead
\[
\lim_{x \to +\infty} f(x)/x = 0,
\]
we say that $f$ is sublinear as $x \to +\infty$; the function $f(x) = \sqrt{x}$ provides a standard example of this second behavior.

Functions that are neither superlinear nor sublinear, but instead grow at essentially a linear rate, may possess a slant asymptote, generalizing the notion of horizontal asymptote to the case of nonzero slope.

\begin{definition}[Slant Asymptote]
We say that the line $y = mx + q$, with $m \neq 0$, is a slant asymptote as $x \to +\infty$ for the function $f$ if
\[
\lim_{x \to +\infty} [f(x) - (mx+q)] = 0.
\]

\end{definition}

\begin{remark}
There is no slant asymptote when $f$ is either superlinear or sublinear as $x \to +\infty$. Indeed, a necessary condition for the existence of a slant asymptote is that $f(x) \sim mx$ as $x \to +\infty$; however, this asymptotic equivalence alone is not sufficient to guarantee the existence of a slant asymptote, as the definition additionally requires the finite limit of the difference $f(x) - mx$ to exist.
\end{remark}

The search for a slant asymptote can be organized into a systematic two-step procedure. The first step consists of computing
\[
\lim_{x \to +\infty} f(x)/x;
\]
if this limit fails to exist, or equals $0$, or is infinite, then $f$ has no slant asymptote. If instead
\[
\lim_{x \to +\infty} f(x)/x = m
\]
for some nonzero real number $m$, one proceeds to the second step, which consists of computing
\[
\lim_{x \to +\infty} [f(x) - mx];
\]
if this second limit fails to exist or is infinite, then again $f$ has no slant asymptote, whereas if
\[
\lim_{x \to +\infty} [f(x) - mx] = q \in \mathbb{R},
\]
then the line $y = mx + q$ is indeed the slant asymptote of $f$ as $x \to +\infty$. Consider now
\[
f(x)=-x+e^{-x}.
\]
Since $e^{-x}\to0$, one has $f(x)\to-\infty$. The asymptote coefficients are
\[
\begin{aligned}
m&=\lim_{x\to+\infty}\frac{-x+e^{-x}}{x}
=-1+\lim_{x\to+\infty}\frac{e^{-x}}{x}=-1,\\
q&=\lim_{x\to+\infty}\bigl(f(x)-mx\bigr)
=\lim_{x\to+\infty}e^{-x}=0.
\end{aligned}
\]
Thus the slant asymptote is $y=-x$. The precise definition of
\[
\lim_{x \to +\infty} f(x) = -\infty
\]
is: for every $M<0$ there exists $R>0$ such that $x>R$ implies
$f(x)<M$.

An entirely analogous set of definitions applies to the behavior of a function $f : (-\infty, r) \to \mathbb{R}$ as $x$ tends to $-\infty$, that is, as $x$ becomes larger and larger in absolute value while remaining negative, again allowing for the three possible outcomes
\[
\lim_{x \to -\infty} f(x) \in \{+\infty, L, -\infty\}.
\]

\subsection{Finite and Infinite Limits at a Finite Point}


We now turn to the second main class of limits, namely those taken as the argument $x$ approaches a finite point $a \in \mathbb{R}$. Given such a point $a$, we consider a function $f$ defined in a neighborhood of $a$, except possibly at $a$ itself. For some $r > 0$, its domain can be written as
\[
f : (a-r, a) \cup (a, a+r) \to \mathbb{R}.
\]
We investigate the two-sided limit
\[
\lim_{x \to a} f(x)
\]
together with the one-sided limits
\[
\begin{gathered}
\lim_{x \to a^+} f(x) \\
\lim_{x \to a^-} f(x).
\end{gathered}
\]

\begin{definition}[Infinite Limit at a Finite Point]
We write
\[
\lim_{x \to a} f(x) = +\infty
\]
if the values of $f(x)$ can be made arbitrarily large by taking $x$ sufficiently close to $a$, that is, if for every $M > 0$ there exists $\delta > 0$ such that
\[
0 < |x - a| < \delta \implies f(x) > M.
\]

\end{definition}

An analogous definition, with the inequality reversed, applies to
\[
\lim_{x \to a} f(x) = -\infty.
\]
In either case, the corresponding geometric feature of the graph is described as follows.

\begin{definition}
The line $x = a$ is called a vertical asymptote of the curve $y = f(x)$.
\end{definition}

\begin{remark}
More precisely, the line $x = a$ is called a vertical asymptote of the curve $y = f(x)$ if at least one of the four possible one-sided infinite limits holds, namely
\[
\begin{gathered}
\lim_{x \to a^+} f(x) = +\infty, \\
\lim_{x \to a^-} f(x) = +\infty, \\
\lim_{x \to a^+} f(x) = -\infty, \\
\lim_{x \to a^-} f(x) = -\infty.
\end{gathered}
\]

\end{remark}

We turn now to the case of a finite limiting value at a finite point, which recovers the classical $\varepsilon$--$\delta$ definition familiar from elementary analysis.

\begin{definition}[Finite Limit at a Finite Point]
Let $L \in \mathbb{R}$. We write
\[
\lim_{x \to a} f(x) = L
\]
if the values of $f(x)$ can be made arbitrarily close to $L$ by taking $x$ sufficiently close to $a$, but distinct from $a$, that is, if for every $\varepsilon > 0$ there exists $\delta > 0$ such that
\[
0 < |x - a| < \delta \implies |f(x) - L| < \varepsilon.
\]

\end{definition}

The corresponding one-sided versions of this definition are obtained by restricting the range of $x$ to one side of $a$. We write
\[
\lim_{x \to a^+} f(x) = L
\]
if for every $\varepsilon > 0$ there exists $\delta > 0$ such that $a < x < a + \delta$ implies $|f(x) - L| < \varepsilon$, and we write
\[
\lim_{x \to a^-} f(x) = L
\]
if for every $\varepsilon > 0$ there exists $\delta > 0$ such that $a - \delta < x < a$ implies $|f(x) - L| < \varepsilon$. These one-sided notions are related to the two-sided limit by the following characterization.

\begin{remark}

\[
\lim_{x \to a} f(x) = L
\]
if and only if
\[
\begin{aligned}
\lim_{x \to a^+} f(x) &= \lim_{x \to a^-} f(x) \\
&= L.
\end{aligned}
\]

\end{remark}

\subsection{Continuity at a Point}

The notion of limit at a finite point allows us to state, at this stage in a preliminary form, the definition of continuity, which will be studied in depth subsequently.

\begin{definition}[Continuity]
Let $f : D \to \mathbb{R}$ be such that $D$ contains an open interval of the form $(a-r, a+r)$ for some $r > 0$. We say that $f$ is continuous at $x = a$ if
\[
\lim_{x \to a} f(x) = f(a).
\]
We say that $f$ is right-continuous at $a$ if
\[
\lim_{x \to a^+} f(x) = f(a),
\]
and left-continuous at $a$ if
\[
\lim_{x \to a^-} f(x) = f(a).
\]

\end{definition}

Figure~\ref{fig:ch4-eps-delta} illustrates how the horizontal and vertical neighborhoods enter the definition of a limit.

\begin{figure}[!htbp]
\centering
\begin{tikzpicture}
  \begin{axis}[
    width=13cm,height=8cm,
    axis lines=middle,
    xlabel={$x$}, ylabel={$f(x)$},
    xmin=-0.5,xmax=4.3, ymin=-0.5,ymax=4.3,
    xtick={2},xticklabels={$a$},
    ytick={2.3},yticklabels={$f(a)=L$},
    ]
    \addplot[figblue,thick,domain=0.2:4,samples=100]
      {0.15*(x-1)^3 - 0.3*(x-1)^2 + 0.9*(x-1) + 2.3};
    \fill[figgreen,opacity=0.15] (axis cs:1.4,-0.5) rectangle (axis cs:2.6,4.3);
    \draw[figgreen,dashed] (axis cs:1.4,-0.5) -- (axis cs:1.4,4.3)
        node[pos=0,below,font=\scriptsize]{$a-\delta$};
    \draw[figgreen,dashed] (axis cs:2.6,-0.5) -- (axis cs:2.6,4.3)
        node[pos=0,below,font=\scriptsize]{$a+\delta$};
    \fill[figred,opacity=0.15] (axis cs:-0.5,1.75) rectangle (axis cs:4.3,2.85);
    \draw[figred,dashed] (axis cs:-0.5,1.75) -- (axis cs:4.3,1.75)
        node[pos=0.97,above,font=\scriptsize]{$L-\varepsilon$};
    \draw[figred,dashed] (axis cs:-0.5,2.85) -- (axis cs:4.3,2.85)
        node[pos=0.97,below,font=\scriptsize]{$L+\varepsilon$};
    \draw[black!50,dotted] (axis cs:2,-0.5) -- (axis cs:2,2.3) -- (axis cs:-0.5,2.3);
    \node[circle,fill=black,inner sep=1.2pt] at (axis cs:2,2.3) {};
  \end{axis}
\end{tikzpicture}
\caption[$\varepsilon$--$\delta$ definition of limit/continuity]{The vertical band
  $(a-\delta,a+\delta)$ (green) is mapped by $f$ entirely inside the horizontal band
  $(L-\varepsilon,L+\varepsilon)$ (red).}
\label{fig:ch4-eps-delta}
\end{figure}

\subsection{General Theorems for Limits of Functions}

Throughout this section, $x_0$ denotes a point of the extended real line $\overline{\mathbb{R}}$, and $f$ denotes a function defined in a neighborhood of $x_0$, except possibly at $x_0$ itself.

One of the most powerful tools for the study of limits of functions is a theorem that reduces the computation of such limits to the computation of limits of sequences, a theory already developed in detail. This reduction rests on the following equivalence.

\begin{theorem}[Limit by Sequences]
Let $L \in \overline{\mathbb{R}}$ and $x_0 \in \overline{\mathbb{R}}$. The following two statements are equivalent: first,
\[
\lim_{x \to x_0} f(x) = L;
\]
second, for every sequence $\{a_n\}$ such that $a_n \to x_0$, one has
\[
\lim_{n \to \infty} f(a_n) = L.
\]

\end{theorem}

A first and immediate consequence of this theorem concerns the nonexistence of limits: in order to show that the limit
\[
\lim_{x \to x_0} f(x)
\]
does not exist, it suffices to exhibit two sequences $a_n \to x_0$ and $b_n \to x_0$ such that
\[
\lim_{n \to \infty} f(a_n) \neq \lim_{n \to \infty} f(b_n).
\]
This technique can be applied, for instance, to prove that
\[
\lim_{x \to 0} \sin(\pi/x)
\]
does not exist, by exhibiting two sequences approaching $0$ along which the sine function takes different limiting values, in line with the motivating example discussed at the outset.

The principal significance of the Theorem on Limits by Sequences is that it shows the computation of limits of real functions to be, in essence, the same problem as the computation of limits of sequences. As a consequence, every result already established for limits of sequences carries over, essentially without modification, to limits of functions; we now list the corresponding theorems for functions, each obtained as a direct transcription of its counterpart for sequences.

We now prove the equivalence stated above between limits of functions and limits of sequences. We restrict attention, for definiteness, to the finite case $x_0 = a \in \mathbb{R}$ and $L \in \mathbb{R}$; the remaining cases, involving infinite points or infinite limiting values, are established by entirely analogous arguments and are left as exercises.

\begin{proof}[Proof of the Sequential Characterization]
We first prove that the second statement implies the first. We argue by contradiction, assuming that $f(x) \not\to L$ as $x \to a$. This means that there exists $\varepsilon > 0$ such that, for every $\delta > 0$, there exists a point $x = x_\delta$ satisfying $0 < |x - a| < \delta$ and $|f(x) - L| \geq \varepsilon$. Choosing $\delta = 1/n$ for each $n \in \mathbb{N}$, we obtain in this way a sequence $\{x_n\} \subset D$ such that
\[
\begin{gathered}
|x_n - a| < \frac{1}{n} \\
|f(x_n) - L| \geq \varepsilon.
\end{gathered}
\]

This construction produces a sequence satisfying $x_n \to a$ but $f(x_n) \not\to L$, in direct contradiction with the second statement, thereby completing the proof of this implication.

We now prove the converse implication, namely that the first statement implies the second. Assume that
\[
\lim_{x \to a} f(x) = L,
\]
and let $\{x_n\} \subset D$ be an arbitrary sequence such that $x_n \to a$; we must show that $f(x_n) \to L$. Fix $\varepsilon > 0$. By the first statement, there exists $\delta > 0$ such that, for every $x \in D$,
\[
|x - a| < \delta \implies |f(x) - L| < \varepsilon.
\]

Since $x_n \to a$, there exists $N \in \mathbb{N}$ such that $|x_n - a| < \delta$ for every $n > N$, and consequently $|f(x_n) - L| < \varepsilon$ for every $n > N$. Since $\varepsilon > 0$ was arbitrary, this shows that $f(x_n) \to L$, as claimed, and completes the proof of the theorem.
\end{proof}

\begin{theorem}[Uniqueness]
If $f$ has a limit at $x_0$, then the value of this limit is unique.
\end{theorem}

\begin{theorem}[Sign Preserving Property]
Suppose that $L > 0$ or $L = +\infty$, and that
\[
\lim_{x \to x_0} f(x) = L.
\]
Then there exists a neighborhood $U(x_0)$ such that $f(x) > 0$ for every $x \in \dot{U}(x_0)$.
\end{theorem}

\begin{theorem}[Comparison Theorem I]
If $f(x) \geq g(x)$ throughout $\dot{U}(x_0)$, then
\[
\lim_{x \to x_0} f(x) \geq \lim_{x \to x_0} g(x).
\]
In particular,
\[
\begin{gathered}
\lim_{x \to x_0} g(x) = +\infty\quad\text{implies}\quad \lim_{x \to x_0} f(x) = +\infty, \\
\lim_{x \to x_0} f(x) = -\infty\quad\text{implies}\quad \lim_{x \to x_0} g(x) = -\infty.
\end{gathered}
\]

\end{theorem}

\begin{theorem}[Comparison Theorem II, Squeeze Theorem]
Suppose that the following inequalities hold throughout $\dot{U}(x_0)$:
\[
\begin{aligned}
f(x) &\leq h(x) \\
&\leq g(x).
\end{aligned}
\]
Assume also that the bounding functions have the same limit:
\[
\begin{aligned}
\lim_{x \to x_0} f(x) &= \lim_{x \to x_0} g(x) \\
&= L.
\end{aligned}
\]
Then the middle function has that limit as well:
\[
\lim_{x \to x_0} h(x) = L.
\]

\end{theorem}

As an application of the Squeeze Theorem, one may compute the limit
\[
\lim_{x \to 0} x^2 \sin(1/x):
\]
since the sine function is bounded between $-1$ and $1$, the product $x^2 \sin(1/x)$ is squeezed between $-x^2$ and $x^2$, both of which tend to $0$ as $x \to 0$, so that the limit equals $0$.

A further essential tool concerns the interchange of limits with the application of an elementary function, generalizing the Composition Rule already established for sequences.

\begin{theorem}[Composition Rule]
Let $g$ be defined on a set $D$ and continuous at $L\in D$. Suppose that
\[
\lim_{x \to x_0} f(x) = L \in \mathbb{R}.
\]
If $f(x)\in D$ for every $x$ sufficiently close to $x_0$, then
The limit of the composition is given by
\[
\begin{aligned}
\lim_{x \to x_0} g(f(x)) &= g\left( \lim_{x \to x_0} f(x) \right) \\
&= g(L).
\end{aligned}
\]

\end{theorem}

In particular, as $x \to x_0$, one has $\sin f(x) \to \sin L$ and
$e^{f(x)}\to e^L$ without additional domain restrictions. For logarithms
one must have $L>0$ and $f(x)>0$ eventually, in which case
$\ln f(x)\to\ln L$. For an arbitrary real exponent $\alpha$, the formula
$[f(x)]^\alpha\to L^\alpha$ follows from
$y^\alpha=e^{\alpha\ln y}$ when $L>0$ and $f(x)>0$ eventually. Rational or
integer exponents may admit larger real domains, but those cases must be
treated according to the natural domain of the particular power function.
The extended-limit conventions $e^{+\infty}=+\infty$,
$e^{-\infty}=0$, and $\ln(+\infty)=+\infty$ are also useful; the notation
$\ln0=-\infty$ refers only to the one-sided limit as the argument tends to
$0$ through positive values.

The algebraic operations of sum, product, and quotient likewise behave well under limits, exactly as in the theory of sequences, subject to the usual restrictions concerning indeterminate forms.

\begin{theorem}[Algebraic Limit Laws]
Let $f$ and $g$ be two functions having a limit at $x_0$, say
\[
\begin{gathered}
\lim_{x \to x_0} f(x) = a \\
\lim_{x \to x_0} g(x) = b,
\end{gathered}
\]
with $a, b \in \mathbb{R}$. Then, as $x \to x_0$,
\[
f(x) \pm g(x) \to a \pm b, \qquad f(x) \cdot g(x) \to a \cdot b, \qquad \frac{f(x)}{g(x)} \to \frac{a}{b},
\]
where the quotient conclusion requires $b\neq0$ (and hence $g(x)\neq0$
eventually).
\end{theorem}

\section{Asymptotic Methods and Special Limits}

\subsection{Infinities, Infinitesimals, and Order of Growth}

Let $f$ and $g$ be two functions such that, as $x \to x_0$, either both $f(x) \to +\infty$ and $g(x) \to +\infty$, or both $f(x) \to 0$ and $g(x) \to 0$. As in the theory of sequences, it is useful to compare the relative rates at which $f$ and $g$ approach their common limiting behavior.

\begin{definition}
If
\[
\lim_{x \to x_0} f(x)/g(x) = 1,
\]
we say that $f$ is asymptotically equivalent to $g$ as $x \to x_0$, and we write $f \sim g$ as $x \to x_0$; in this case $f$ and $g$ are said to be infinities, respectively infinitesimals, of the same order. If the limit in question equals $+\infty$, we say that $f$ is an infinite of higher order than $g$ as $x \to x_0$, and we write $g \prec f$. If instead the limit equals $0$, we say that $f$ is an infinitesimal of higher order than $g$ as $x \to x_0$.
\end{definition}

As a first illustration of this comparison, one may examine how the family of test functions $x^\alpha$ compares with itself for different values of $\alpha$, thereby recovering the hierarchy of powers familiar from the theory of infinities of sequences.

A more refined classification assigns to each infinity, respectively infinitesimal, a definite order, together with an explicit leading term.

\begin{definition}[Order of Infinity]
Let $|f(x)|\to\infty$ as $x\to+\infty$. We say that $f$ is an infinity of order $\beta$ as $x \to +\infty$ if
\[
\lim_{x \to +\infty} \frac{f(x)}{k x^\beta} = 1, \qquad \text{equivalently} \qquad f(x) \sim k x^\beta,
\]
for suitable numbers $k\neq0$ and $\beta>0$. The expression $kx^\beta$
is then called the leading part of $f$ as $x\to+\infty$. The sign of
$k$ records whether $f$ tends to $+\infty$ or $-\infty$.
\end{definition}

For $f(x)=x^2-2x^3$, the leading part as $x\to+\infty$ is
$-2x^3$, because
\[
\frac{x^2-2x^3}{-2x^3}=1-\frac{1}{2x}\longrightarrow1.
\]

\begin{definition}[Order of Infinitesimal]
Let $a \in \mathbb{R}$, and let $f$ be infinitesimal as $x \to a$, that is,
\[
\lim_{x \to a} f(x) = 0.
\]
We say that $f$ is infinitesimal of order $\gamma$ as $x \to a$ if
\[
\lim_{x \to a} \frac{f(x)}{h (x-a)^\gamma} = 1, \qquad \text{equivalently} \qquad f(x) \sim h (x-a)^\gamma,
\]
for suitable numbers $h \neq 0$ and $\gamma > 0$. The expression $h(x-a)^\gamma$ is then called the leading part of $f$ as $x \to a$.
\end{definition}

For the same function, the leading part as $x\to0$ is $x^2$, since
\[
\frac{x^2-2x^3}{x^2}=1-2x\longrightarrow1.
\]
Thus the dominant term is $x^2$ near the origin and $-2x^3$ at
positive infinity.

\subsection{Indeterminate Forms, Substitution, and Euler's Number}

The comparison of orders of growth is greatly facilitated by a collection of special limits, which we now record, organized according to the type of indeterminate form to which they apply.

We first consider indeterminate forms of type $\infty/\infty$. For every value of the real constants $\alpha, \beta > 0$ and $b > 1$, one has
\[
\lim_{x \to +\infty} \frac{(\ln x)^\beta}{x^\alpha} = 0, \qquad \lim_{x \to +\infty} \frac{x^\alpha}{b^x} = 0,
\]
which together yield the hierarchy of infinities as $x \to +\infty$,
\[
(\ln x)^\beta \prec x^\alpha \prec b^x.
\]

It is furthermore useful to record that $x^a \prec x^b$ whenever $a < b$, extending the comparison of powers already familiar from the theory of sequences. We next consider indeterminate forms of type $0/0$, all valid as $x \to 0$. These are
\[
\frac{\ln(1+x)}{x} \to 1 \ \longleftrightarrow\ \ln(1+x) \sim x, \qquad \frac{e^x - 1}{x} \to 1 \ \longleftrightarrow\ e^x - 1 \sim x,
\]
and, for every $\alpha \in \mathbb{R}$,
\[
\frac{(1+x)^\alpha - 1}{x} \to \alpha \ \longleftrightarrow\ (1+x)^\alpha - 1 \sim \alpha x,
\]
together with
\[
\frac{\sin x}{x} \to 1 \ \longleftrightarrow\ \sin x \sim x, \qquad \frac{1 - \cos x}{x^2} \to \frac{1}{2} \ \longleftrightarrow\ 1 - \cos x \sim \frac{1}{2} x^2.
\]

Finally, we record a useful family of limits of the indeterminate form $0 \cdot \infty$, valid as $x \to 0^+$: for every $\alpha > 0$ one has
\[
\lim_{x \to 0^+} x^\alpha \ln x = 0,
\]
and more generally, for every pair of constants $\alpha, \beta > 0$,
\[
\lim_{x \to 0^+} x^\alpha |\ln x|^\beta = 0.
\]

These special limits, together with the hierarchy of infinities, provide the essential building blocks for resolving a wide range of indeterminate forms, as may be illustrated by computing the limits
\[
\lim_{x \to +\infty} \frac{e^x - x^2 + 1}{\sqrt{x} + x^{100}}, \qquad \lim_{x \to 0} \frac{x^3 + 7x^2 + 4x}{x^5 - x}, \qquad \lim_{x \to 0} \frac{\sin x^2}{\sin^2 x}, \qquad \lim_{x \to 0} \frac{\ln(1+3x)}{x^2 + 2x}.
\]


A particularly efficient method for evaluating limits involving products and quotients of asymptotically equivalent functions is furnished by the Substitution Principle, entirely analogous to the corresponding result already established for sequences.

\begin{theorem}[Substitution Principle]
If, as $x \to x_0$, one has $f(x) \sim F(x)$, $g(x) \sim G(x)$, and $h(x) \sim H(x)$, then
\[
\lim_{x \to x_0} \frac{f(x) \cdot g(x)}{h(x)} = \lim_{x \to x_0} \frac{F(x) \cdot G(x)}{H(x)},
\]
provided that the limit on the right-hand side exists.
\end{theorem}

A further class of indeterminate forms arises when studying limits of the type
\[
\lim_{x \to x_0} [f(x)]^{g(x)}.
\]

As in the corresponding treatment for sequences, it is convenient to take the logarithm, writing
\[
\begin{aligned}
f(x)^{g(x)} &= e^{\ln [f(x)]^{g(x)}} \\
&= e^{g(x) \ln f(x)}.
\end{aligned}
\]
If it can be shown that $g(x) \ln f(x) \to L$, possibly with $L = \pm \infty$, then one concludes that
\[
[f(x)]^{g(x)} = e^{g(x) \ln f(x)} \to e^L.
\]

The indeterminate forms to which this technique applies are precisely $1^\infty$, $0^0$, and $\infty^0$. For example, let
\[
L=\lim_{x \to 0^+} ((1+x)^{1/2})^{1/\sin x}.
\]
Taking logarithms gives
\[
\begin{aligned}
\ln L
&=\lim_{x\to0^+}\frac{\ln(1+x)}{2\sin x}\\
&=\frac12\left(\lim_{x\to0^+}\frac{\ln(1+x)}{x}\right)
\left(\lim_{x\to0^+}\frac{x}{\sin x}\right)\\
&=\frac12.
\end{aligned}
\]
Exponentiating yields $L=e^{1/2}=\sqrt e$. A particularly important special case of the indeterminate form $1^\infty$ generalizes the definition of Euler's number already encountered in the theory of sequences.

\begin{theorem}[Special Limit of Type $1^\infty$]
For every $k \in \mathbb{R}$,
\[
\lim_{x \to \infty} \left( 1 + \frac{k}{x} \right)^x = e^k.
\]

\end{theorem}

As a related pair of limits, one also has
\[
\lim_{x \to -\infty} (1 + 1/x)^x = e,
\]
in perfect analogy with the case $x \to +\infty$, together with the equivalent formulation obtained by the substitution $x = 1/t$, namely
\[
\lim_{x \to 0} (1+x)^{1/x} = e.
\]

The mathematical constant $e$ is fundamentally defined as the limit of a specific sequence, and this definition serves as the foundation for a broader family of convergence results involving exponential functions and indeterminate forms. Recall that
\[
e := \lim_{n \to \infty} \left( 1 + \frac{1}{n} \right)^n.
\]

This definition, originally stated for the sequence of natural numbers $n$, extends naturally to any sequence diverging to positive infinity, as the following theorem shows.

\begin{theorem}
Let $\{b_n\}$ be a sequence of real numbers such that $b_n \to +\infty$ as $n \to \infty$. Then
\[
\lim_{n \to \infty} \left( 1 + \frac{1}{b_n} \right)^{b_n} = e.
\]

\end{theorem}

\begin{proof}
Since $b_n \to +\infty$, for every $n \in \mathbb{N}$ one may identify an integer $k = k(n)$, depending on $n$, such that $k < b_n \leq k+1$. From this inequality one derives corresponding bounds on the reciprocal term, namely
\[
\frac{1}{k+1} \leq \frac{1}{b_n} < \frac{1}{k}.
\]

Observe that, as $n \to \infty$, one has $k \to \infty$ as well, since $b_n$ diverges. Using the monotonicity of the function $x \mapsto (1+x)^y$ with respect to positive bases and exponents, one obtains the two-sided bound
\[
\begin{aligned}
\left( 1 + \frac{1}{k+1} \right)^k &\leq \left( 1 + \frac{1}{b_n} \right)^{b_n} \\
&\leq \left( 1 + \frac{1}{k} \right)^{k+1}.
\end{aligned}
\]

We analyze the limits of the lower and upper bounds separately. For the lower bound, we write
\[
\left( 1 + \frac{1}{k+1} \right)^k = \left( 1 + \frac{1}{k+1} \right)^{k+1} \cdot \left( 1 + \frac{1}{k+1} \right)^{-1}.
\]

As $n \to \infty$, and hence $k \to \infty$, the first factor converges to $e$ by the standard definition, while the second factor converges to $1$; consequently,
\[
\begin{aligned}
\lim_{n \to \infty} (1+(1/(k+1)))^k &= e \cdot 1 \\
&= e.
\end{aligned}
\]
Analogously, for the upper bound, we write
\[
\left( 1 + \frac{1}{k} \right)^{k+1} = \left( 1 + \frac{1}{k} \right)^k \cdot \left( 1 + \frac{1}{k} \right),
\]
and, as $k \to \infty$, the first factor converges to $e$ while the second converges to $1$, so that
\[
\begin{aligned}
\lim_{n \to \infty} (1+(1/k))^{k+1} &= e \cdot 1 \\
&= e.
\end{aligned}
\]

Since both the lower and upper bounds converge to the common value $e$, the Squeeze Theorem guarantees that the sequence $(1+1/b_n)^{b_n}$ likewise converges to $e$, which proves the theorem.
\end{proof}

More generally, for any real number $x \in \mathbb{R}$, the exponential function $e^x$ can itself be defined by means of the analogous limit
\[
e^x := \lim_{n \to \infty} \left( 1 + \frac{x}{n} \right)^n.
\]

\subsection{Fundamental Notable Limits and Operational Rules}

\begin{remark}[The logarithmic limit]

Let $\{\varepsilon_n\}$ be a sequence of real numbers such that $\varepsilon_n \to 0$ as $n \to \infty$. We now present five fundamental limits, involving trigonometric, logarithmic, and exponential functions evaluated at $\varepsilon_n$, all of which present the indeterminate form $0/0$ and are of crucial importance for evaluating derivatives and analyzing asymptotic behavior throughout analysis. These limits establish that the natural logarithm of $1+\varepsilon_n$ is asymptotically equivalent to $\varepsilon_n$ itself; that $e^{\varepsilon_n}-1$ is likewise asymptotically equivalent to $\varepsilon_n$; that, for any $\alpha > 0$, the expression $(1+\varepsilon_n)^\alpha - 1$ is asymptotically equivalent to $\alpha \varepsilon_n$; that $\sin(\varepsilon_n)$ is asymptotically equivalent to $\varepsilon_n$; and that $1-\cos(\varepsilon_n)$ is asymptotically equivalent to $\varepsilon_n^2/2$. We now provide a rigorous proof of each of these five statements in turn, building each proof on the results already established.


We aim to prove that
\[
\lim_{n \to \infty} \frac{\ln(1+\varepsilon_n)}{\varepsilon_n} = 1.
\]

First suppose that $\varepsilon_n>0$. Setting
$b_n:=1/\varepsilon_n$, the generalized limit established in the preceding
subsection gives
\[
\lim_{n \to \infty} (1+\varepsilon_n)^{1/\varepsilon_n} = e.
\]

Taking the natural logarithm of both sides, and invoking the continuity of the logarithm, we obtain
\[
\begin{aligned}
\lim_{n \to \infty} \ln\!\left( (1+\varepsilon_n)^{1/\varepsilon_n} \right)
&= \ln\!\left( \lim_{n \to \infty} (1+\varepsilon_n)^{1/\varepsilon_n} \right) \\
&= \ln(e) = 1.
\end{aligned}
\]

Using the logarithmic identity $\ln(a^b) = b \ln(a)$, the expression inside the limit on the left-hand side may be rewritten as
\[
\begin{aligned}
\ln\!\left( (1+\varepsilon_n)^{1/\varepsilon_n} \right) &= \frac{1}{\varepsilon_n} \ln(1+\varepsilon_n) \\
&= \frac{\ln(1+\varepsilon_n)}{\varepsilon_n},
\end{aligned}
\]
from which the desired conclusion follows immediately.

If instead $\varepsilon_n<0$, define
\[
\delta_n=-\frac{\varepsilon_n}{1+\varepsilon_n}>0.
\]
For all sufficiently large $n$, $\delta_n$ is defined, tends to zero, and
$1+\varepsilon_n=(1+\delta_n)^{-1}$. Therefore
\[
\frac{\ln(1+\varepsilon_n)}{\varepsilon_n}
=(1+\delta_n)\frac{\ln(1+\delta_n)}{\delta_n}\longrightarrow1.
\]
For a sequence whose signs vary, apply the two arguments to its positive
and negative subsequences; terms with $\varepsilon_n=0$ may be omitted
when evaluating the punctured quotient. Every subsequence has the same
limit, so the original quotient tends to $1$.

\end{remark}

\begin{remark}[The exponential limit]

We aim to prove that
\[
\lim_{n \to \infty} \frac{e^{\varepsilon_n} - 1}{\varepsilon_n} = 1.
\]

Setting $\delta_n := e^{\varepsilon_n} - 1$, we observe that, since $\varepsilon_n \to 0$, one has $e^{\varepsilon_n} \to 1$, and consequently $\delta_n \to 0$. Taking the natural logarithm of the relation $e^{\varepsilon_n} = 1 + \delta_n$ yields $\varepsilon_n = \ln(1+\delta_n)$, and substituting into the original quotient gives
\[
\frac{e^{\varepsilon_n} - 1}{\varepsilon_n} = \frac{\delta_n}{\ln(1+\delta_n)}.
\]

This expression is precisely the reciprocal of the quotient analyzed in the proof of the Logarithmic Limit; since that limit was shown to equal $1$, we conclude that
\[
\begin{aligned}
\lim_{n \to \infty} \frac{e^{\varepsilon_n} - 1}{\varepsilon_n}
&= \lim_{n \to \infty} \frac{\delta_n}{\ln(1+\delta_n)} \\
&= \frac{1}{\displaystyle \lim_{n \to \infty} \frac{\ln(1+\delta_n)}{\delta_n}}
= \frac{1}{1} = 1.
\end{aligned}
\]

\end{remark}

\begin{remark}[The power-function limit]

We aim to prove that, for every $\alpha > 0$,
\[
\lim_{n \to \infty} \frac{(1+\varepsilon_n)^\alpha - 1}{\varepsilon_n} = \alpha.
\]
Rewriting the power in terms of the exponential and logarithmic functions,
\[
(1+\varepsilon_n)^\alpha = e^{\alpha \ln(1+\varepsilon_n)},
\]
and setting $u_n := \alpha \ln(1+\varepsilon_n)$, we observe that $u_n \to 0$ as $\varepsilon_n \to 0$, since $\ln(1+\varepsilon_n) \to 0$. Multiplying and dividing the quotient of interest by $u_n$, for $n$ large enough that $u_n \neq 0$, we obtain
\[
\begin{aligned}
\frac{(1+\varepsilon_n)^\alpha - 1}{\varepsilon_n} &= \frac{e^{u_n}-1}{\varepsilon_n} \\
&= \frac{e^{u_n}-1}{u_n} \cdot \frac{u_n}{\varepsilon_n},
\end{aligned}
\]
and, substituting back the definition of $u_n$, the second factor becomes
\[
\begin{aligned}
\frac{u_n}{\varepsilon_n} &= \frac{\alpha \ln(1+\varepsilon_n)}{\varepsilon_n} \\
&= \alpha \cdot \frac{\ln(1+\varepsilon_n)}{\varepsilon_n}.
\end{aligned}
\]

As $n \to \infty$, the first factor $(e^{u_n}-1)/u_n$ converges to $1$ by the Exponential Limit, since $u_n \to 0$, while the second factor converges to $\alpha$ by the Logarithmic Limit. Multiplying these two limits together, we conclude that
\[
\begin{aligned}
\lim_{n \to \infty} \frac{(1+\varepsilon_n)^\alpha - 1}{\varepsilon_n} &= 1 \cdot \alpha \\
&= \alpha.
\end{aligned}
\]

\end{remark}

\begin{remark}[The sine limit]

We aim to prove that
\[
\lim_{n \to \infty} \frac{\sin(\varepsilon_n)}{\varepsilon_n} = 1.
\]

For every sufficiently large $n$ with $\varepsilon_n\neq0$, set
$\theta_n=|\varepsilon_n|\in(0,\pi/2)$. A classical geometric argument,
comparing the areas of a triangle, a circular sector, and a larger triangle
in the unit circle, establishes
$\sin\theta<\theta<\tan\theta$ for $\theta\in(0,\pi/2)$. Hence
\[
1 < \frac{\theta_n}{\sin(\theta_n)} < \frac{1}{\cos(\theta_n)}.
\]
Taking reciprocals, which reverses the direction of both inequalities, gives
\[
\cos(\theta_n) < \frac{\sin(\theta_n)}{\theta_n} < 1.
\]

Because sine is odd,
\[
\frac{\sin(\varepsilon_n)}{\varepsilon_n}
=\frac{\sin(\theta_n)}{\theta_n}.
\]
As $n\to\infty$, $\theta_n\to0$ and $\cos\theta_n\to1$; the Squeeze
Theorem therefore yields
\[
\lim_{n \to \infty} \sin(\varepsilon_n)/\varepsilon_n = 1.
\]

\end{remark}

\begin{remark}[The cosine limit]

We aim to prove that
\[
\lim_{n \to \infty} \frac{1-\cos(\varepsilon_n)}{\varepsilon_n^2} = \frac{1}{2}.
\]

Using the Pythagorean identity $\sin^2\theta + \cos^2\theta = 1$, which yields $1-\cos^2\theta = \sin^2\theta$, we rewrite the numerator $1-\cos(\varepsilon_n)$ by multiplying and dividing by the conjugate expression $1+\cos(\varepsilon_n)$,
\[
\begin{aligned}
1 - \cos(\varepsilon_n)
&= \frac{\big(1-\cos(\varepsilon_n)\big)\big(1+\cos(\varepsilon_n)\big)}{1+\cos(\varepsilon_n)} \\
&= \frac{1-\cos^2(\varepsilon_n)}{1+\cos(\varepsilon_n)}
= \frac{\sin^2(\varepsilon_n)}{1+\cos(\varepsilon_n)}.
\end{aligned}
\]
Substituting this expression into the original quotient, we obtain
\[
\begin{aligned}
\frac{1-\cos(\varepsilon_n)}{\varepsilon_n^2} &= \frac{\sin^2(\varepsilon_n)}{\varepsilon_n^2 \big(1+\cos(\varepsilon_n)\big)} \\
&= \left( \frac{\sin(\varepsilon_n)}{\varepsilon_n} \right)^2 \cdot \frac{1}{1+\cos(\varepsilon_n)}.
\end{aligned}
\]

By the Sine Limit established above, the first factor converges to $1^2 = 1$; and, since $\varepsilon_n \to 0$ implies $\cos(\varepsilon_n) \to 1$, the second factor converges to $1/(1+1) = 1/2$. Multiplying these two limits together yields the final result,
\[
\begin{aligned}
\lim_{n \to \infty} \frac{1-\cos(\varepsilon_n)}{\varepsilon_n^2} &= 1 \cdot \frac{1}{2} \\
&= \frac{1}{2}.
\end{aligned}
\]

\end{remark}

\section{Continuity and Continuous Functions on Intervals}

\subsection{Continuity and Types of Discontinuity}

We now turn from the classification of elementary functions to the fundamental analytic property of continuity. Let $f : D \to \mathbb{R}$ be a function such that the domain $D$ contains an open interval of the form $(a-r, a+r)$ for some $r > 0$; in other words, $D$ contains a full neighborhood of the point $a$.

\begin{definition}
We say that $f$ is continuous at $x = a$ if
\[
\lim_{x \to a} f(x) = f(a).
\]

\end{definition}

Unwinding the definition of limit, this condition can be restated in an equivalent, purely quantitative form: $f$ is continuous at $x = a$ if and only if for every $\varepsilon > 0$ there exists $\delta > 0$ such that
\[
|x - a| < \delta \implies |f(x) - f(a)| < \varepsilon.
\]

This equivalence relies on the notion of one-sided limits, and it is useful to recall the precise relationship between the two-sided limit and its one-sided counterparts. Recalling that
\[
\lim_{x \to a} f(x) = L
\]
means that for every $\varepsilon > 0$ there exists $\delta > 0$ such that $0 < |x-a| < \delta$ implies $|f(x) - L| < \varepsilon$, one has the following characterization.

\begin{remark}

\[
\lim_{x \to a} f(x) = L
\]
if and only if
\[
\begin{aligned}
\lim_{x \to a^+} f(x) &= \lim_{x \to a^-} f(x) \\
&= L.
\end{aligned}
\]

\end{remark}

Combining this remark with the definition of continuity yields a practical four-part test for continuity at a point, which is often more convenient to apply than the $\varepsilon$--$\delta$ definition directly. A function $f$ is continuous at $x = a$ if and only if the right-hand limit
\[
\lim_{x \to a^+} f(x)
\]
is finite, the left-hand limit
\[
\lim_{x \to a^-} f(x)
\]
is finite, the two one-sided limits coincide, say
\[
\begin{aligned}
\lim_{x \to a^+} f(x) &= \lim_{x \to a^-} f(x) \\
&= : L,
\end{aligned}
\]
and this common value $L$ coincides exactly with $f(a)$.


When one or more of the conditions in the continuity test fail, the function is said to be discontinuous at the point in question, and it is useful to classify the resulting types of discontinuity according to which condition fails. A removable discontinuity occurs when the one-sided limits are finite and equal to one another, but the common value differs from $f(a)$, the value of the function at the point itself. A jump discontinuity occurs when the one-sided limits are both finite but fail to coincide. An infinite discontinuity occurs when the one-sided limits exist, in the extended sense, but are infinite. Finally, a discontinuity is said to be of the second type whenever at least one of the one-sided limits fails to exist even in the extended real sense.

Figure~\ref{fig:atlas-discontinuities} compares a removable discontinuity with a jump discontinuity.

\begin{figure}[!htbp]
\centering
\begin{tikzpicture}[>=Stealth]
\begin{axis}[width=5.8cm,height=5.2cm,axis lines=middle,ticklabel style={font=\small},label style={font=\small},legend style={font=\scriptsize,draw=black!20,fill=white},xlabel={$x$},ylabel={$f(x)$},xmin=-1.2,xmax=1.2,ymin=-0.3,ymax=2.5,xtick={-1,0,1},ytick={1,2}]
\addplot[figblue,very thick,domain=-1:1] {x+1};
\addplot[only marks,mark=*,mark options={fill=white,draw=figblue},mark size=3pt] coordinates {(0,1)};
\addplot[only marks,figred,mark=*] coordinates {(0,2)};
\end{axis}
\begin{axis}[width=5.8cm,height=5.2cm,axis lines=middle,ticklabel style={font=\small},label style={font=\small},legend style={font=\scriptsize,draw=black!20,fill=white},at={(6.2cm,0)},xlabel={$x$},ylabel={$f(x)$},xmin=-1.2,xmax=1.2,ymin=-0.3,ymax=2.5,xtick={-1,0,1},ytick={1,2}]
\addplot[figblue,very thick] coordinates {(-1,0)(0,0)};
\addplot[figred,very thick] coordinates {(0,1)(1,1)};
\addplot[only marks,mark=*,mark options={fill=white,draw=figblue},mark size=3pt] coordinates {(0,0)};
\addplot[only marks,figred,mark=*] coordinates {(0,1)};
\end{axis}
\end{tikzpicture}
\caption{Left: the graph agrees with $x+1$ away from zero, but its assigned value at zero is $2$; the hole at height $1$ can be repaired by redefining that value. Right: a unit jump at zero has unequal one-sided limits, so changing the value at the jump cannot restore continuity.}
\label{fig:atlas-discontinuities}
\end{figure}

\subsection{Continuity of Elementary Functions}

Having classified the possible failures of continuity, we turn to establishing that continuity is, in fact, the typical behavior exhibited by elementary functions. This is a consequence of a small number of general theorems governing how continuity is preserved under algebraic operations, composition, and inversion.

\begin{theorem}[Algebra of Continuous Functions]
Let $f$ and $g$ be continuous at the point $a$. Then the functions $f + g$, $f \cdot g$, and, provided $g(a) \neq 0$, the quotient $f/g$ are all continuous at $a$.
\end{theorem}

\begin{theorem}[Continuity of Compositions]
If $f$ is continuous at $a$ and $g$ is continuous at $f(a)$, then the composition $g \circ f$, defined by
\[
(g \circ f)(x) = g(f(x)),
\]
is continuous at $a$.
\end{theorem}

\begin{theorem}[Continuity of the Inverse]
Let $I$ be an interval and let $f : I \to J$ be continuous and injective on $I$, with range $J$. Then the inverse function $f^{-1} : J \to I$ is continuous on $J$.
\end{theorem}

Starting from the elementary facts that the constant function, the identity function $f(x) = x$, the exponential function $f(x) = e^x$, and the sine function $f(x) = \sin x$ are all continuous, the three theorems above allow one to deduce that essentially every elementary function encountered in analysis is continuous on its natural domain. This far-reaching consequence can be organized according to the following chain of deductions.

\begin{itemize}
\item \textit{Power functions}: the continuity of $x^n$ for $n \in \mathbb{N}$ follows from repeated application of the product rule for continuous functions.
\item \textit{Rational functions}: the continuity of polynomial functions follows from the sum and product rules, while the continuity of rational functions additionally requires the quotient rule.
\item \textit{Root and fractional power functions}: the continuity of the root functions $x^{1/n}$ follows from the Inverse Theorem applied to the power function $x^n$, and the continuity of $x^{m/n}$ follows by composition of a power function with a root function.
\item \textit{Logarithmic and hyperbolic functions}: the continuity of the exponential function implies, via the Inverse Theorem, the continuity of the logarithmic functions, and consequently also the continuity of the hyperbolic functions, which are defined in terms of exponentials.
\item \textit{Real power functions}: for an arbitrary real exponent
$\alpha$, the representation $x^\alpha=e^{\alpha\ln x}$ proves continuity
on $(0,\infty)$. Integer powers extend continuously to all of
$\mathbb{R}$, and rational powers may have larger real domains depending
on the parity of the denominator; these extensions must be considered
separately.
\item \textit{Trigonometric functions}: the continuity of $\sin x$ implies, by composition, the continuity of $\cos x$; the continuity of $\tan x$ then follows from the quotient rule, and the continuity of the inverse trigonometric functions follows from the Inverse Theorem.
\end{itemize}

A further useful reformulation of continuity, particularly convenient when working with compositions of continuous functions, concerns the interchange of limits and function evaluation. If $f$ and $g$ are two continuous functions, then the limit may be taken inside the function, in the sense that
\[
\begin{aligned}
\lim_{x \to a} f(g(x)) &= f(g(a)) \\
&= f\left( \lim_{x \to a} g(x) \right).
\end{aligned}
\]

Recalling the characterization of limits by means of sequences, this property admits an equivalent sequential formulation: $f$ is continuous at $x = a$ if and only if, for every sequence $\{x_n\} \subset D$ such that $x_n \to a$, one has
\[
\lim_{n \to \infty} f(x_n) = f\left( \lim_{n \to \infty} x_n \right).
\]

\subsection{Extreme and Intermediate Value Theorems on Compact Intervals}


We now turn to a collection of results that reveal the particularly rich structure enjoyed by continuous functions when their domain is a closed and bounded interval. The first such result guarantees both the boundedness of the function and the attainment of its extreme values.

\begin{theorem}[Weierstrass' Theorem]
Let $f$ be continuous on a closed bounded interval $[a,b]$. Then $f$ is bounded on $[a,b]$, and $f$ attains its extreme values; that is, there exist two points $c, d \in [a,b]$ for which the following bounds hold for every $x \in [a,b]$:
\[
f(d) \leq f(x) \leq f(c).
\]

\end{theorem}

The hypotheses of this theorem, namely continuity together with a closed and bounded domain, are all essential, as illustrated by several simple examples of functions that fail to attain extreme values once one of these hypotheses is dropped. The function $f(x) = x$ on $[0, \infty)$ is not even bounded, since the domain is unbounded. The function $f(x) = x/(1+x)$ on $[0, \infty)$ is bounded but never attains its least upper bound, which equals $1$ but is only approached asymptotically as $x \to \infty$. The function $f(x) = 1/x$ on $(0,1]$ is not bounded, since it diverges as $x \to 0^+$, reflecting the fact that the domain is not closed. Finally, the function $f(x) = 1 - x$ on $(0,1]$ is bounded but never attains its least upper bound, which equals $1$ but is only approached as $x \to 0^+$, again because the domain fails to be closed. These examples show precisely why Weierstrass' Theorem requires both continuity and a closed bounded interval: only under these combined hypotheses is the attainment of extreme values guaranteed.


A second fundamental property of continuous functions on closed bounded intervals concerns the existence of roots whenever the function changes sign.

\begin{theorem}[Bolzano's Theorem]
Let $f$ be a continuous function on a closed bounded interval $[a,b]$. Then
\[
f(a) \cdot f(b) < 0 \implies \exists c \in (a,b) : f(c) = 0.
\]

\end{theorem}

\begin{remark}
Bolzano's Theorem depends crucially on, and is in fact equivalent to, the completeness of the real number system.
\end{remark}

This result admits a natural generalization, known as the Intermediate Value Theorem, which asserts that a continuous function on a closed bounded interval attains every value lying between its minimum and its maximum.

\begin{theorem}[Intermediate Value Theorem]
Let $f$ be a continuous function on a closed bounded interval $[a,b]$. Then $f$ takes on every value between its maximum value $M$ and its minimum value $m$; that is,
\[
\forall k \in [m, M] \ \exists c \in [a,b] : f(c) = k.
\]

\end{theorem}

\begin{proof}
Since $f$ is continuous on $[a,b]$, Weierstrass' Theorem guarantees that $f$ attains its extreme values $m$ and $M$ at certain points $x_m, x_M \in [a,b]$. For every $x \in [a,b]$, this gives
\[
m = f(x_m) \leq f(x) \leq f(x_M) = M.
\]

Given any $k \in (m, M)$, consider the auxiliary function $g(x) = f(x) - k$ on the interval with endpoints $x_m$ and $x_M$. This function is continuous, being the difference of a continuous function and a constant, and satisfies $g(x_m) < 0$ and $g(x_M) > 0$. Applying Bolzano's Theorem to $g$, we obtain a point $c$ such that $g(c) = 0$, that is, $f(c) = k$, which completes the proof.
\end{proof}

We conclude this section by presenting a constructive proof of Bolzano's Theorem, based on the classical bisection method, which iteratively narrows down an interval containing a root of $f$.

\begin{proof}[Proof of Bolzano's Theorem by Bisection]
Set $a_0 = a$, $b_0 = b$, and $L = b_0 - a_0$, and assume without loss of generality that $f(a_0) < 0$ and $f(b_0) > 0$. The starting configuration therefore consists of an interval $[a_0, b_0]$ satisfying $f(a_0) < 0$, $f(b_0) > 0$, and $b_0 - a_0 = L$. Consider the midpoint $(a_0 + b_0)/2$ of this interval, and distinguish three possibilities. If
\[
f( (a_0+b_0)/2 ) = 0,
\]
then $c = (a_0+b_0)/2$ is a root of $f$, and the proof is complete. If instead $f( (a_0+b_0)/2 ) < 0$, set $a_1 = (a_0+b_0)/2$ and $b_1 = b_0$. Finally, if $f( (a_0+b_0)/2 ) > 0$, set $a_1 = a_0$ and $b_1 = (a_0+b_0)/2$. In either of the latter two cases, one obtains a new interval $[a_1, b_1] \subset [a_0, b_0]$ satisfying
\[
f(a_1) < 0, \qquad f(b_1) > 0, \qquad b_1 - a_1 = \frac{L}{2}.
\]

This bisection procedure can be iterated. At a generic step $k$, it
produces an interval $[a_k,b_k]\subset[a_{k-1},b_{k-1}]$ satisfying
\[
f(a_k)<0,\qquad f(b_k)>0,\qquad b_k-a_k=\frac{L}{2^k}.
\]
Consequently, after $n$ steps one obtains $[a_n,b_n]$ such that
\[
f(a_n) < 0, \qquad f(b_n) > 0, \qquad b_n - a_n = \frac{L}{2^n},
\]
and, by construction, the resulting intervals are nested,
\[
[a_n, b_n] \subset [a_{n-1}, b_{n-1}] \subset \cdots \subset [a_1, b_1] \subset [a_0, b_0].
\]

We now analyze the sequences $\{a_n\}$ and $\{b_n\}$ separately. The sequence $\{a_n\}$ is monotone increasing and bounded above, since $a_n \leq b$ for every $n$; by the Monotone Convergence Theorem for real sequences, it therefore converges to a finite limit $\alpha \in \mathbb{R}$, and moreover $\alpha \leq b$. Similarly, the sequence $\{b_n\}$ is monotone decreasing and bounded below, since $b_n \geq a$ for every $n$; again by the Monotone Convergence Theorem, it converges to a finite limit $\beta \in \mathbb{R}$, with $\beta \geq a$. Consider now the sequence $b_n - a_n$. On one hand, by the algebra of limits,
\[
\lim_{n \to \infty} (b_n - a_n) = \beta - \alpha.
\]
On the other hand, by construction,
\[
\begin{aligned}
\lim_{n \to \infty} (b_n - a_n) &= \lim_{n \to \infty} L / 2^n \\
&= 0.
\end{aligned}
\]

Since the limit of a sequence is unique, these two expressions must coincide, so that $\beta - \alpha = 0$, that is, $\alpha = \beta$. Set
\[
\begin{aligned}
c &= \alpha \\
&= \beta,
\end{aligned}
\]
and note that $c \in [a,b]$, since $\alpha \leq b$ and $\beta \geq a$. Since $f$ is continuous, the sequential characterization of continuity gives
\[
\lim_{n \to \infty} a_n = c \implies \lim_{n \to \infty} f(a_n) = f(c).
\]
Because $f(a_n) < 0$ for every $n$, the Sign Preserving Property implies $f(c) \leq 0$. Analogously,
\[
\lim_{n \to \infty} b_n = c \implies \lim_{n \to \infty} f(b_n) = f(c),
\]
and since $f(b_n) > 0$ for every $n$, the same property implies $f(c) \geq 0$. Combining the two inequalities, we obtain
\[
0 \leq f(c) \leq 0 \implies f(c) = 0,
\]
which shows that $c$ is a root of $f$ in $[a,b]$ and completes the proof.
\end{proof}

\newpage
\chapter{Differential Calculus and Taylor Approximations}

\section{Derivatives and Differentiability}

\subsection{The Derivative: Definition, Geometry, and Notation}

Let $f : D \to \mathbb{R}$ be a function defined in a neighborhood of a point $a \in \mathbb{R}$. The central object of differential calculus is the derivative of $f$ at $a$, which measures the instantaneous rate of change of $f$ near that point.

\begin{definition}
The derivative of a function $f$ at a point $a$, denoted by $f'(a)$, is defined as
\[
\begin{aligned}
f'(a) &= \lim_{h \to 0} \frac{f(a+h) - f(a)}{h} \\
&= \lim_{x \to a} \frac{f(x) - f(a)}{x - a},
\end{aligned}
\]
provided that this limit exists and is finite.
\end{definition}

For any given $h \in \mathbb{R}$, it is convenient to introduce the notation
\[
R_a(h) = \frac{f(a+h) - f(a)}{h}
\]
for the ratio appearing in the definition above, so that
\[
f'(a) = \lim_{h \to 0} R_a(h),
\]
provided the limit exists and is finite. The derivative admits several complementary interpretations, the first of which is geometric. Given a curve $C$ with equation $y = f(x)$, one seeks the tangent line to $C$ at the point $P = (a, f(a))$. To construct this tangent line, one considers a nearby point $Q = (x, f(x))$ with $x \neq a$, and computes the slope of the secant line through $P$ and $Q$,
\[
m_Q = \frac{f(x) - f(a)}{x - a}.
\]

Letting $Q$ approach $P$ along the curve, that is, letting $x$ approach $a$, one defines the tangent line $t$ to be the limiting position of the secant line, whenever such a limiting position exists.

\begin{definition}
The tangent line to the curve $C$ at the point $P$ is the line through $P$ with slope
\[
m = \lim_{x \to a} \frac{f(x) - f(a)}{x - a},
\]
provided that this limit exists and is finite.
\end{definition}

This limit is precisely the derivative $f'(a)$, so that the derivative may be identified with the slope of the tangent line to the graph of $f$ at the point $(a, f(a))$.

A second interpretation of the derivative arises from physics. Suppose an object moves along a straight line according to an equation of motion $s = f(t)$, where $s$ denotes the displacement of the object from the origin at time $t$. In order to determine the instantaneous velocity of the object at a given time $t = a$, one first considers the time interval from $t = a$ to $t = a + h$, over which the change in position is $f(a+h) - f(a)$; the average velocity over this interval is
\[
v_{\text{mean}} = \frac{\text{displacement}}{\text{time}} = \frac{f(a+h) - f(a)}{h}.
\]

Computing the average velocity over shorter and shorter time intervals, one defines the instantaneous velocity at time $t = a$ as the limit of these average velocities,
\[
v(a) = \lim_{h \to 0} \frac{f(a+h) - f(a)}{h},
\]
which again is recognized as the derivative $f'(a)$. A third and more general interpretation concerns rates of change. Suppose $y$ is a quantity depending on another quantity $x$, so that $y$ is a function of $x$ and one writes $y = f(x)$. If $x$ changes from $x_1$ to $x_2$, the change in $x$, also called the increment of $x$, is $\Delta x = x_2 - x_1$, and the corresponding change in $y$ is
\[
\Delta y = f(x_2) - f(x_1).
\]
The difference quotient, or incremental ratio,
\[
\frac{\Delta y}{\Delta x} = \frac{f(x_2) - f(x_1)}{x_2 - x_1}
\]
is called the average rate of change of $y$ with respect to $x$ over the interval $[x_1, x_2]$. Considering the average rate of change over smaller and smaller intervals, by letting $\Delta x$ approach zero, one defines the instantaneous rate of change of $y$ with respect to $x$ at $x_1$ as
\[
\lim_{\Delta x \to 0} \frac{\Delta y}{\Delta x} = \lim_{\Delta x \to 0} \frac{f(x_1 + \Delta x) - f(x_1)}{\Delta x},
\]
which is once again recognized as the derivative $f'(x_1)$. This general notion of instantaneous rate of change has broad applicability outside physics; for instance, biomedical scientists have studied the physiological changes in the body resulting from alcohol consumption by means of the law
\[
C(t) = 0.0225 \, t \, e^{-0.0467 t}
\]
milligrams per milliliter, and one may ask how quickly the blood alcohol concentration is increasing after ten minutes, a question that is naturally answered by evaluating the derivative $C'(10)$.

Several notations for the derivative coexist in the literature, namely $f'(x)$, $y'$, $df/dx$, $dy/dx$, $(d/dx) f(x)$, and $Df(x)$. Among these, the symbol $dy/dx$, known as Leibniz notation and dating back to 1672, deserves a specific comment: it should not be regarded as an actual ratio of two quantities $dy$ and $dx$, but simply as a synonym for $f'(x)$. Nonetheless, this notation is extremely useful and suggestive, since it directly invokes the definition of the derivative as a limit of a ratio of increments,
\[
\frac{dy}{dx} = \lim_{\Delta x \to 0} \frac{\Delta y}{\Delta x}.
\]

Figure~\ref{fig:atlas-secant-tangent} shows how secant slopes approach the derivative at a fixed point.

\begin{figure}[!htbp]
\centering
\begin{tikzpicture}[>=Stealth]
\begin{axis}[width=11cm,height=5.8cm,axis lines=middle,ticklabel style={font=\small},label style={font=\small},legend style={font=\scriptsize,draw=black!20,fill=white},xlabel={$x$},ylabel={$y$},xmin=0,xmax=2.4,ymin=-0.3,ymax=5,xtick={1,1.5,2},ytick={1,2,3,4}]
\addplot[figblue,very thick,domain=0:2.2,samples=100] {x^2};
\addplot[figred,dashed,thick,domain=0.6:2.15] {3*x-2};
\addplot[figred!55,dashdotted,thick,domain=0.6:2.15] {2.5*x-1.5};
\addplot[figgreen,thick,domain=0.4:2.1] {2*x-1};
\addplot[only marks,mark=*,figblue] coordinates {(1,1)(1.5,2.25)(2,4)};
\node[below right] at (axis cs:1,1) {$(1,1)$};
\end{axis}
\end{tikzpicture}
\caption{For $f(x)=x^2$ at $x=1$, the secants through the second points $x=2$ and $x=1.5$ have slopes $3$ and $2.5$. They approach the tangent slope $2$ as the second point tends to the first. The tangent is drawn in green.}
\label{fig:atlas-secant-tangent}
\end{figure}

\subsection{Differentiability and Continuity}

\begin{definition}
A function $f$ is differentiable at $a$ if $f'(a)$ exists. We say that $f$ is differentiable on an interval $I$ if it is differentiable at every point of $I$.
\end{definition}

Differentiability is a strictly stronger requirement than continuity, as recorded in the following theorem.

\begin{theorem}
If $f$ is differentiable at $a$, then $f$ is continuous at $a$.
\end{theorem}

\begin{remark}
The converse of this theorem is not true: the function $f(x) = |x|$ is continuous but not differentiable at $x = 0$.
\end{remark}

\begin{proof}
We wish to prove that
\[
\lim_{x \to a} f(x) = f(a),
\]
which is equivalent to showing that
\[
\lim_{x \to a} [f(x) - f(a)] = 0.
\]
Writing
\[
f(x) - f(a) = \frac{f(x) - f(a)}{x - a} \cdot (x-a)
\]
and passing to the limit as $x \to a$, using the differentiability of $f$ at $a$, we obtain
\[
\begin{aligned}
\lim_{x \to a} [f(x) - f(a)] &= \lim_{x \to a} \frac{f(x) - f(a)}{x-a} \cdot \lim_{x \to a} (x-a) \\
&= f'(a) \cdot 0 \\
&= 0,
\end{aligned}
\]
which proves the claim.
\end{proof}

There are essentially four distinct ways in which a function can fail to be differentiable at a point $x = a$. The most familiar is the presence of an angle in the graph, as exemplified by $f(x) = |x|$ at the origin: here $f$ is continuous at $a$, and the right-derivative and left-derivative, respectively defined by
\[
f'_+(a) := \lim_{h \to 0^+} \frac{f(a+h) - f(a)}{h}, \qquad f'_-(a) := \lim_{h \to 0^-} \frac{f(a+h) - f(a)}{h},
\]
both exist and are finite, but differ from one another. A second way is the presence of a discontinuity of $f$ itself at $a$, which automatically precludes differentiability. A third way is the presence of a vertical tangent or a cusp, in which case $f$ is continuous at $a$ but the right-derivative and left-derivative are both infinite.

\subsection{Derivative Functions and Higher Derivatives}

Given any number $x$ for which the limit
\[
f'(x) = \lim_{h \to 0} [f(x+h) - f(x)]/h
\]
exists and is finite, one may assign to $x$ the corresponding number $f'(x)$, obtaining a rule $x \mapsto f'(x)$. In this way, $f'$ may itself be regarded as a new function, called the derivative of $f$. The domain of $f'$ is the set
\[
\{x \in \operatorname{Dom}(f) : f'(x) \text{ is finite}\},
\]
and this set may be strictly smaller than the domain of $f$ itself. The graph of $f'$ can be derived from the graph of $f$ by exploiting the geometric fact that $f'(x)$ represents the slope of the tangent line to the graph of $f$ at the point $(x, f(x))$.


If $y = f(x)$ is a differentiable function, then its derivative $y' = f'(x)$ is itself a function of $x$, and it may happen that $f'$ possesses a derivative of its own, denoted by $(f')' = f''$. This new function is called the second derivative of $f$, since it is obtained by differentiating the derivative of $f$. Using Leibniz notation, the second derivative of $f$ is written as
\[
\begin{aligned}
y'' &= f''(x) \\
&= \frac{d}{dx}\left( \frac{dy}{dx} \right) \\
&= \frac{d^2 y}{dx^2}.
\end{aligned}
\]

A particularly natural physical interpretation of the second derivative is as a rate of change of a rate of change, the most familiar instance of which is acceleration. If $s = s(t)$ is the position function of an object moving along a straight line, its first derivative represents the velocity of the object as a function of time,
\[
\begin{aligned}
v(t) &= s'(t) \\
&= ds/dt;
\end{aligned}
\]
the instantaneous rate of change of velocity with respect to time is then called the acceleration of the object,
\[
\begin{aligned}
a(t) &= v'(t) \\
&= \frac{d^2 s}{dt^2}.
\end{aligned}
\]

The process of differentiation can be iterated further. The third derivative is defined as the derivative of the second derivative, $(f'')' = f'''$, and is written in Leibniz notation as
\[
\begin{aligned}
y''' &= f'''(x) \\
&= \frac{d}{dx}\left( \frac{d^2 y}{dx^2} \right) \\
&= \frac{d^3 y}{dx^3}.
\end{aligned}
\]

More generally, the $n$-th derivative of $f$ is denoted by $f^{(n)}$ and is obtained from $f$ by differentiating $n$ times in succession,
\[
\begin{aligned}
y^{(n)} &= f^{(n)}(x) \\
&= \frac{d^n y}{dx^n}.
\end{aligned}
\]

As will be seen in later developments of the theory, second and higher derivatives play a fundamental role in representing functions as infinite sums of power functions.

\section{Differentiation Rules}

\subsection{General Rules for Combining Derivatives}

When new functions are formed from previously known functions by means of addition, subtraction, multiplication, division, or composition, their derivatives can be computed directly in terms of the derivatives of the original functions. The theory of differentiation is governed by a small number of such rules, namely the rule for the sum, difference, and multiplication by a constant, known collectively as linearity, together with the product rule, the quotient rule, the chain rule for composition, and the rule for the derivative of the inverse function.

\begin{theorem}[Linearity]
Let $f$ and $g$ be differentiable at the point $x$, and let $c \in \mathbb{R}$. Then the functions $f + g$ and $c \cdot f$ are differentiable at $x$, and
\[
(f+g)'(x) = f'(x) + g'(x), \qquad (c \cdot f)'(x) = c \cdot f'(x).
\]

\end{theorem}

These rules follow immediately by applying the algebraic limit laws to the corresponding incremental ratios.

\begin{theorem}[Product Rule]
Let $f$ and $g$ be differentiable at the point $x$. Then the product function $f \cdot g$ is differentiable at $x$, and
\[
(f \cdot g)'(x) = f'(x) \cdot g(x) + f(x) \cdot g'(x).
\]

\end{theorem}

\begin{theorem}[Quotient Rule]
Let $f$ and $g$ be differentiable at the point $x$. If $g(x) \neq 0$, then the quotient $f/g$ is differentiable at $x$, and
\[
\left( \frac{f}{g} \right)'(x) = \frac{f'(x) \cdot g(x) - f(x) \cdot g'(x)}{[g(x)]^2}.
\]

\end{theorem}

\begin{theorem}[Chain Rule]
Let $g$ be differentiable at the point $x$, and let $f$ be differentiable at the point $y = g(x)$. Then the composition $f \circ g$ is differentiable at $x$, and
\[
(f \circ g)'(x) = f'(g(x)) \cdot g'(x).
\]

\end{theorem}

\begin{theorem}[Derivative of the Inverse]
Let $f : (a,b) \to \mathbb{R}$ be continuous and injective. If $f$ is differentiable at the point $c$ and $f'(c) \neq 0$, then the inverse function $f^{-1}$ is differentiable at the point $d = f(c)$, and
\[
(f^{-1})'(d) = \frac{1}{f'(c)}.
\]

\end{theorem}

We conclude this subsection with the proof of the product rule, since it illustrates the standard technique of adding and subtracting an auxiliary term in order to isolate the incremental ratios of the individual factors.

\begin{proof}[Proof of the Product Rule]
We must compute the limit of the incremental ratio of $f \cdot g$,
\[
R_h = \frac{f(x+h) g(x+h) - f(x) g(x)}{h}.
\]
Adding and subtracting the term $f(x+h) g(x)$, we obtain
\[
f(x+h) g(x+h) - f(x) g(x) = f(x+h) [g(x+h) - g(x)] + g(x) [f(x+h) - f(x)],
\]
so that
\[
R_h = f(x+h) \, \frac{g(x+h) - g(x)}{h} + g(x) \, \frac{f(x+h) - f(x)}{h}.
\]

In light of the differentiability of $f$ and $g$ at $x$, we may now pass to the limit as $h \to 0$. Observing that $f(x+h) \to f(x)$ by the continuity of $f$ at $x$, which is itself a consequence of the differentiability of $f$ at $x$, and that
\[
\frac{g(x+h) - g(x)}{h} \to g'(x), \qquad \frac{f(x+h) - f(x)}{h} \to f'(x),
\]
we conclude that
\[
\lim_{h \to 0} R_h = f(x) g'(x) + g(x) f'(x),
\]
which is precisely the statement of the product rule.
\end{proof}

\subsection{Derivatives of the Elementary Functions}

Combining the general differentiation rules with the definition of derivative, one can compute explicitly the derivative of each of the elementary functions cataloged earlier. Here the parameter $\alpha$ denotes an arbitrary real number, and every derivative formula stated below is understood to hold at every point $x$ where the right-hand side is well defined. We begin with the exponential function. Its derivative is given by
\[
\frac{d}{dx} e^x = e^x.
\]
Indeed, writing the incremental ratio and factoring out $e^x$, we have
\[
\lim_{h \to 0} \frac{e^{x+h} - e^x}{h} = e^x \lim_{h \to 0} \frac{e^h - 1}{h},
\]
and the result follows from the special limit
\[
\lim_{h \to 0} (e^h - 1)/h = 1
\]
established in the theory of infinitesimal sequences. The derivative of the natural logarithm is given by
\[
\frac{d}{dx} \ln x = \frac{1}{x}.
\]
Indeed, we compute
\[
\begin{aligned}
\lim_{h \to 0} \frac{\ln(x+h) - \ln x}{h} &= \lim_{h \to 0} \frac{\ln\!\left( \frac{x+h}{x} \right)}{h} \\
&= \frac{1}{x} \lim_{h \to 0} \frac{\ln\!\left(1 + \frac{h}{x}\right)}{\frac{h}{x}},
\end{aligned}
\]
and the conclusion again follows from the corresponding special limit for the logarithm. The derivatives of the sine and cosine functions are given by
\[
\frac{d}{dx} \sin x = \cos x, \qquad \frac{d}{dx} \cos x = -\sin x.
\]
To establish the first identity, we must compute
\[
\lim_{h \to 0} \frac{\sin(x+h) - \sin x}{h}.
\]
Using the addition formula
\[
\sin(x+h) = \sin x \cos h + \cos x \sin h,
\]
we obtain
\[
\lim_{h \to 0} \frac{\sin x \cos h + \cos x \sin h - \sin x}{h} = \sin x \cdot \lim_{h \to 0} \frac{\cos h - 1}{h} + \cos x \cdot \lim_{h \to 0} \frac{\sin h}{h}.
\]
Exploiting the special limits
\[
\begin{gathered}
\lim_{h \to 0} (\sin h)/h = 1 \\
\lim_{h \to 0} \frac{\cos h - 1}{h} = \lim_{h \to 0} (-h) \, \frac{1 - \cos h}{h^2} = 0,
\end{gathered}
\]
we conclude that the derivative of $\sin x$ equals $\cos x$; the derivative of $\cos x$ is obtained by a completely analogous computation. The derivative of a real power function is given by
\[
\frac{d}{dx} x^\alpha = \alpha \, x^{\alpha - 1}, \qquad \alpha \in \mathbb{R}.
\]
Indeed, one rewrites the increment $(x+h)^\alpha - x^\alpha$ as $x^\alpha [ ( 1 + h/x )^\alpha - 1 ]$, so that
\[
\frac{(x+h)^\alpha - x^\alpha}{h} = x^{\alpha - 1} \, \frac{\left(1 + \frac{h}{x}\right)^\alpha - 1}{\frac{h}{x}},
\]
and hence
\[
\begin{aligned}
\lim_{h \to 0} \frac{(x+h)^\alpha - x^\alpha}{h} &= x^{\alpha - 1} \lim_{h \to 0} \frac{\left(1 + \frac{h}{x}\right)^\alpha - 1}{\frac{h}{x}} \\
&= x^{\alpha - 1} \cdot \alpha,
\end{aligned}
\]
where the last limit is again computed by means of the corresponding special limit for infinitesimal sequences.

Using the product rule, one easily obtains the derivative of $\tan x$ from the derivatives of $\sin x$ and $\cos x$; using instead the rule for the derivative of the inverse function, one proves the derivatives of the inverse trigonometric functions,
\[
D \arcsin x = \frac{1}{\sqrt{1 - x^2}}, \qquad D \arccos x = -\frac{1}{\sqrt{1 - x^2}}, \qquad D \arctan x = \frac{1}{1 + x^2}.
\]

\begin{proof}
Let $f(x) = \sin x$. By definition, $y = \arcsin x$ if and only if $x = \sin y$. Applying the rule for the derivative of the inverse function, we obtain
\[
\begin{aligned}
D f^{-1}(x) &= \frac{1}{f'(y)} \\
&= \frac{1}{\cos y} \\
&= \frac{1}{\sqrt{1 - \sin^2 y}} \\
&= \frac{1}{\sqrt{1 - x^2}},
\end{aligned}
\]
where the positive square root is selected since $\cos y \geq 0$ on the range of $\arcsin$. The formula for $\arccos x$ is obtained analogously. For $f(x) = \tan x$, one instead has $y = \arctan x$ if and only if $x = \tan y$, and hence
\[
\begin{aligned}
D f^{-1}(x) &= \frac{1}{f'(y)} \\
&= \frac{1}{\frac{1}{\cos^2 y}} \\
&= \cos^2 y \\
&= \frac{\cos^2 y}{\cos^2 y + \sin^2 y} \\
&= \frac{1}{1 + \tan^2 y} \\
&= \frac{1}{1 + x^2},
\end{aligned}
\]
which completes the proof.
\end{proof}

\section{The Main Theorems of Differential Calculus}

Having established the basic computational rules for derivatives, we now turn to three fundamental theoretical results, namely Fermat's Theorem, Rolle's Theorem, and Lagrange's Theorem, together with several applications to optimization, the behavior of graphs, and the computation of limits.

\subsection{Fermat's Theorem, Critical Points, and Optimization}

We begin by making precise the notions of local extremum that underlie the theory. Let $f$ be defined in a neighborhood of a point $c \in \mathbb{R}$.

\begin{definition}
The point $c$ is a local minimum point for $f$ if there exists a neighborhood $U(c)$ such that $f(x) \geq f(c)$ for every $x \in U(c)$, and $c$ is a local maximum point for $f$ if there exists a neighborhood $U(c)$ such that $f(x) \leq f(c)$ for every $x \in U(c)$. In either case, the corresponding value $f(c)$ is called a local minimum, respectively local maximum, value for $f$.
\end{definition}

The following theorem provides a necessary condition, expressed in terms of the derivative, for the existence of a local extremum.

\begin{theorem}[Fermat's Theorem]
If $f$ has a local maximum or a local minimum at $c$, and $f$ is differentiable at $c$, then $f'(c) = 0$.
\end{theorem}

\begin{definition}
A point $c$ in the domain of $f$ such that $f'(c) = 0$ is called a critical point of $f$.
\end{definition}

Fermat's Theorem may thus be rephrased as follows: if $f$ has a local maximum or minimum at $c$, then either $f$ is not differentiable at $c$, or $c$ is a critical point of $f$.

\begin{proof}
Suppose, for the sake of definiteness, that $f$ has a local maximum at $c$; the case of a local minimum is treated analogously. Then $f(c) \geq f(x)$ for every $x$ sufficiently close to $c$, which implies that for $h$ sufficiently close to $0$,
\begin{equation*}
f(c) \geq f(c+h) \implies f(c+h) - f(c) \leq 0. \tag{$*$}
\end{equation*}
If $h > 0$ is sufficiently small, dividing this inequality by $h$ preserves its direction, giving
\[
\frac{f(c+h) - f(c)}{h} \leq 0.
\]

Taking the right-hand limit as $h \to 0^+$ and invoking the Sign Preserving Property for limits, we obtain
\[
\lim_{h \to 0^+} \frac{f(c+h) - f(c)}{h} \leq 0 \implies f'_+(c) \leq 0.
\]

Analogously, if $h < 0$ is sufficiently small, dividing inequality $(*)$ by $h$ reverses its direction, giving
\[
\frac{f(c+h) - f(c)}{h} \geq 0,
\]
so that taking the left-hand limit yields
\[
\lim_{h \to 0^-} \frac{f(c+h) - f(c)}{h} \geq 0 \implies f'_-(c) \geq 0.
\]
Since $f$ is differentiable at $c$, we have
\[
\begin{aligned}
f'_+(c) &= f'_-(c) \\
&= f'(c),
\end{aligned}
\]
and the only value compatible with the two inequalities $f'(c) \leq 0$ and $f'(c) \geq 0$ is $f'(c) = 0$, as claimed.
\end{proof}


Fermat's Theorem provides the theoretical foundation for a systematic method of finding the absolute extreme values of a function on a closed bounded interval. Let $f$ be a continuous function on a closed bounded interval $[a,b]$; the problem is to find the absolute maximum and minimum values of $f$ on $[a,b]$, that is, two points $c, d \in [a,b]$ such that
\[
\begin{aligned}
f(c) &\leq f(x) \\
&\leq f(d)
\end{aligned}
\]
for every $x \in [a,b]$. Weierstrass' Theorem guarantees that such points exist, but it does not, by itself, indicate how to locate them; this is precisely the role played by Fermat's Theorem.

The resulting procedure, known as the Closed Interval Method, consists of three steps. The first step is to find the critical points of $f$ inside the open interval $(a,b)$, including any points where $f$ fails to be differentiable, and to evaluate $f$ at each of these points. The second step is to evaluate $f$ at the two endpoints of the interval, namely at $a$ and at $b$. The third step is to compare all the values obtained in the first two steps: the largest of these values is the absolute maximum value of $f$ on $[a,b]$, and the smallest is the absolute minimum value.

As an illustration of this method, one may consider the function $g(x) = x^4 - 4x^3$ on the interval $[0,5]$, and ask what its maximum and minimum values are on this interval, as well as what happens on the open interval $(0,5)$ and on the whole real line $\mathbb{R}$. A further classical application concerns an open box constructed from a sheet of metal measuring twenty centimetres by twelve centimetres, from each corner of which a square of side length $x$ is removed before the sides are folded upward; the problem is to determine the value of $x$ that maximizes the resulting volume of the box.

\subsection{Rolle's Theorem and Lagrange's Theorem}

A second fundamental theorem, known as the Mean Value Theorem, generalizes Fermat's Theorem by relating the derivative of a function at an interior point to the average rate of change of the function over an entire interval.

\begin{theorem}[Lagrange's Theorem]
Let $f$ be a function satisfying the following two hypotheses: $f$ is continuous on the closed interval $[a,b]$, and $f$ is differentiable on the open interval $(a,b)$. Then there exists a number $c \in (a,b)$ such that
\[
f'(c) = \frac{f(b) - f(a)}{b - a}.
\]

\end{theorem}

A particularly important special case of Lagrange's Theorem, in which the values of $f$ at the two endpoints coincide, is traditionally stated as a separate result.

\begin{theorem}[Rolle's Theorem]
Let $f$ be a function satisfying the following three hypotheses: $f$ is continuous on the closed interval $[a,b]$, $f$ is differentiable on the open interval $(a,b)$, and $f(a) = f(b)$. Then there exists a number $c \in (a,b)$ such that $f'(c) = 0$.
\end{theorem}

\begin{proof}
The proof distinguishes three cases according to the behavior of $f$ on $(a,b)$. In the first case, $f$ is constant on $[a,b]$; then $f' \equiv 0$ throughout, so the number $c$ may be taken to be any point of $(a,b)$. In the second case, $f$ is not constant and there exists some $x \in (a,b)$ such that
\[
f(x) > f(a) = f(b);
\]
by Weierstrass' Theorem, $f$ attains a maximum value somewhere on $[a,b]$, and since this maximum value exceeds $f(a) = f(b)$, it must be attained at a point $c$ inside the open interval $(a,b)$, so that $c$ is a local maximum of $f$. Since $f$ is differentiable at $c$, Fermat's Theorem yields $f'(c) = 0$. In the third case, $f$ is not constant and there exists some $x \in (a,b)$ such that
\[
f(x) < f(a) = f(b);
\]
by Weierstrass' Theorem, $f$ attains a minimum value somewhere on $[a,b]$, and since this minimum value is smaller than $f(a) = f(b)$, it must be attained at a point $c$ inside $(a,b)$. Again, Fermat's Theorem yields $f'(c) = 0$, which completes the proof in all three cases.
\end{proof}

\begin{proof}[Proof of Lagrange's Theorem]
We define the auxiliary function
\[
F(x) = f(x) - \frac{f(b) - f(a)}{b - a} (x - a).
\]

Observe that $F$ is continuous on $[a,b]$ and differentiable on $(a,b)$, being a linear combination of $f$ and of an affine function of $x$. Moreover, a direct computation shows that
\[
\begin{aligned}
F(a) &= f(a) \\
&= F(b),
\end{aligned}
\]
so that $F$ satisfies the hypotheses of Rolle's Theorem on $[a,b]$. Applying Rolle's Theorem to $F$, we find a point $c \in (a,b)$ such that $F'(c) = 0$. For every $x \in (a,b)$, the derivative of $F$ is
\[
F'(x) = f'(x) - \frac{f(b) - f(a)}{b-a}.
\]
Setting $x = c$ yields
\[
0 = F'(c) = f'(c) - \frac{f(b) - f(a)}{b-a} \implies f'(c) = \frac{f(b) - f(a)}{b-a},
\]
which establishes the theorem.
\end{proof}

Figure~\ref{fig:atlas-mean-value} illustrates the parallel chord and tangent in the Mean Value Theorem.

\begin{figure}[!htbp]
\centering
\begin{tikzpicture}[>=Stealth]
\begin{axis}[width=11cm,height=5.8cm,axis lines=middle,ticklabel style={font=\small},label style={font=\small},legend style={font=\scriptsize,draw=black!20,fill=white},xlabel={$x$},ylabel={$f(x)$},xmin=-0.15,xmax=2.3,ymin=-0.4,ymax=4.7,xtick={0,2},ytick={1,2,3,4}]
\addplot[figblue,very thick,domain=0:2.15,samples=100] {x^2};
\addplot[figred,thick] coordinates {(0,0)(2,4)};
\addplot[figgreen,thick,domain=0.5:2] {2*x-1};
\addplot[only marks,mark=*,figblue] coordinates {(0,0)(1,1)(2,4)};
\draw[dotted] (axis cs:1,0) -- (axis cs:1,1);
\node[below] at (axis cs:1,0) {$c$};
\end{axis}
\end{tikzpicture}
\caption{For $f(x)=x^2$ on $[0,2]$, the endpoint chord has slope $2$. The tangent at $c=1$ has the same slope, as required by the Mean Value Theorem. The red chord records average change; the green tangent records instantaneous change.}
\label{fig:atlas-mean-value}
\end{figure}

\section{L'H\^opital's Rule and Its Applications}

\subsection{The Rule and Its Proof}

Having established the main theorems of differential calculus, we now turn to a natural question: how can derivatives be used to compute limits, particularly limits presenting an indeterminate form? The following rule, based ultimately on Lagrange's Theorem and a generalization thereof due to Cauchy, provides a powerful tool for this purpose.

\begin{theorem}[L'H\^opital's Rule]
Let $c \in \mathbb{R}$, and let $f, g$ be differentiable in a punctured neighborhood $\dot{U}(c)$ of $c$. Assume that the following three conditions hold: first,
\[
\begin{aligned}
\lim_{x \to c} f(x) &= \lim_{x \to c} g(x) \\
&= 0;
\end{aligned}
\]
second, $g'(x) \neq 0$ throughout $\dot{U}(c)$; and third,
\[
\lim_{x \to c} f'(x)/g'(x) = \ell
\]
for some $\ell \in \mathbb{R}$. Then
\[
\lim_{x \to c} \frac{f(x)}{g(x)} = \ell.
\]

\end{theorem}

\begin{remark}
The rule remains valid for indeterminate forms of the type $\infty / \infty$, for one-sided limits, and for limits taken as $x \to \pm \infty$.
\end{remark}

The proof of L'H\^opital's Rule rests on a generalization of Lagrange's Theorem due to Cauchy, which we now state.

\begin{theorem}[Cauchy's Theorem]
Let $f, g$ be continuous on $[a,b]$ and differentiable on $(a,b)$, with $g' \neq 0$ on $(a,b)$. Then there exists $c \in (a,b)$ such that
\[
\frac{f(b) - f(a)}{g(b) - g(a)} = \frac{f'(c)}{g'(c)}.
\]

\end{theorem}

\begin{proof}[Proof of L'H\^opital's Rule]
We present a short proof under the additional assumption that $f$ and $g$ are continuous at $c$, so that, in particular,
\[
\begin{aligned}
f(c) &= g(c) \\
&= 0.
\end{aligned}
\]

Fix any $x > c$ sufficiently close to $c$. Since $g' \neq 0$ on $(c,x)$, Cauchy's Theorem may be applied to the interval $[c,x]$, yielding a point $\xi_x \in (c,x)$ such that
\[
\begin{aligned}
\frac{f(x)}{g(x)} &= \frac{f(x) - f(c)}{g(x) - g(c)} \\
&= \frac{f'(\xi_x)}{g'(\xi_x)}.
\end{aligned}
\]

As $x \to c$, one has $\xi_x \to c$ as well, since $\xi_x$ lies between $c$ and $x$; the third hypothesis of the theorem then guarantees that the right-hand side tends to $\ell$, and the conclusion follows.
\end{proof}

\subsection{Repeated Application and Indeterminate Forms}

As an illustration of the rule, consider the limit
\[
\lim_{x \to 0} (\sin x - x)/x^3.
\]

Taking the quotient of the derivatives of numerator and denominator, one obtains $(\cos x - 1)/(3x^2)$; applying L'H\^opital's Rule once more to this new quotient yields $-\sin x / (6x)$, and one further application yields $-\cos x / 6$. In summary, in light of L'H\^opital's Rule applied three times in succession,
\[
\begin{aligned}
\lim_{x \to 0} \frac{\sin x - x}{x^3}
&= \lim_{x \to 0} \frac{\cos x - 1}{3x^2} \\
&= \lim_{x \to 0} \frac{-\sin x}{6x}
= \lim_{x \to 0} \frac{-\cos x}{6}
= -\frac{1}{6}.
\end{aligned}
\]

It is worth noting that the computation may in fact be stopped one step earlier, at the quotient $-\sin x / (6x)$, since the value $-1/6$ then follows directly from the special limit
\[
\lim_{x \to 0} (\sin x)/x = 1
\]
established previously, without a further application of the rule. L'H\^opital's Rule is particularly useful for establishing the hierarchy of infinities as $x \to +\infty$, namely the relation $(\ln x)^\beta \prec x^\alpha \prec e^x$ for $\alpha > 0$; consider, for instance, the comparison $\ln x \prec x^\alpha$. Indeed, for $\alpha > 0$, applying L'H\^opital's Rule gives
\[
\lim_{x \to +\infty} \frac{\ln x}{x^\alpha} \xrightarrow{\text{l'H\^opital}} \lim_{x \to +\infty} \frac{\frac{1}{x}}{\alpha x^{\alpha - 1}} = \lim_{x \to +\infty} \frac{1}{\alpha x^\alpha} = 0.
\]

The rule is likewise useful for resolving indeterminate forms of type $0 \cdot \infty$; for instance, the special limit
\[
\lim_{x \to 0^+} x^\alpha (\ln x)^\beta = 0,
\]
in the particular case $\beta = 1$, can be established by rewriting $x^\alpha \ln x$ as the quotient $\ln x / x^{-\alpha}$, which presents the indeterminate form $-\infty / \infty$, and applying L'H\^opital's Rule:
\[
\begin{aligned}
\lim_{x \to 0^+} x^\alpha \ln x
&= \lim_{x \to 0^+} \frac{\ln x}{x^{-\alpha}} \\
&\xrightarrow{\text{l'H\^opital}} \lim_{x \to 0^+} \frac{\frac{1}{x}}{-\alpha x^{-\alpha - 1}} \\
&= \lim_{x \to 0^+} \left( -\frac{x^\alpha}{\alpha} \right) = 0.
\end{aligned}
\]

Despite its evident usefulness, L'H\^opital's Rule must be applied with a degree of caution, and it is important to be aware of its limitations. In particular, the rule cannot be used to prove the special limits from which it itself indirectly draws, since some of those limits are logically prior to the derivatives of the elementary functions involved. Moreover, the rule cannot be applied to limits such as
\[
\begin{gathered}
\lim_{x \to 1^+} x/(x-1) \\
\lim_{x \to +\infty} (x + \sin x)/x,
\end{gathered}
\]
since in the first case the hypothesis of a $0/0$ or $\infty/\infty$ indeterminate form is not satisfied, while in the second case the derivative of the denominator does not lead to a resolvable limit in the required sense. Finally, there exist situations in which the rule is simply inconclusive, such as the limit
\[
\lim_{x \to \infty} (x^2+1)^{1/2}/x,
\]
where repeated differentiation of numerator and denominator fails to terminate in a determinate expression, so that alternative methods must be used instead.

\section{Monotonicity, Concavity, and Graph Analysis}

\subsection{Monotonicity}

The sign of the first derivative of a function is directly related to its monotone behavior, as expressed by the following test, which is itself a consequence of Lagrange's Theorem.

\begin{theorem}[Increasing/Decreasing Test]
Let $f$ be differentiable on an interval $I$. If $f'(x) = 0$ for every $x \in I$, then $f$ is constant on $I$. If $f' > 0$ on $I$, then $f$ is increasing on $I$. If $f' < 0$ on $I$, then $f$ is decreasing on $I$.
\end{theorem}

\begin{proof}
Let $x_1 < x_2$ be any two numbers in $I$. Applying Lagrange's Theorem to the interval $[x_1, x_2]$, we find a number $c$ between $x_1$ and $x_2$ such that
\[
f(x_2) - f(x_1) = f'(c)(x_2 - x_1).
\]

If $f' > 0$ on $I$, this identity immediately shows that $f(x_2) - f(x_1) > 0$, proving that $f$ is increasing on $I$; the remaining two cases are established by an entirely analogous argument.
\end{proof}

In view of this test, it follows that a local maximum or minimum can occur at any point $c$ where the first derivative changes sign, and this observation is made precise in the following criterion.

\begin{theorem}[First Derivative Test]
Assume that $f$ is continuous at a point $c$ and differentiable in a punctured neighborhood $\dot{U}(c)$. If $f'$ changes sign from positive to negative at $c$, then $f$ has a local maximum at $c$. If $f'$ changes sign from negative to positive at $c$, then $f$ has a local minimum at $c$. If $f'$ does not change sign at $c$, then $f$ has neither a local maximum nor a local minimum at $c$.
\end{theorem}

\begin{remark}
It should not be assumed that all local minima look essentially alike upon close inspection of the graph: there exist local minima at which the function fails to be decreasing on the left and increasing on the right of the point in question, so that the First Derivative Test, while sufficient, does not capture every possible local extremum.
\end{remark}

A further useful consequence of Lagrange's Theorem provides a sufficient condition for differentiability at a point, based on the limiting behavior of the derivative nearby.

\begin{theorem}
Assume $f$ is continuous at a point $c$ and differentiable in a punctured neighborhood $\dot{U}(c)$. If
\[
\lim_{x \to c} f'(x) = L \in \mathbb{R},
\]
then $f$ is differentiable at $c$ and $f'(c) = L$.
\end{theorem}

\begin{proof}
By Lagrange's Theorem applied to the interval $[c, c+h]$, for $h$ sufficiently small, there exists a point $x_h \in (c, c+h)$ such that
\[
\frac{f(c+h) - f(c)}{h} = f'(x_h).
\]
Since $x_h \to c$ as $h \to 0$, it follows that
\[
\begin{aligned}
f'_+(c) &= \lim_{h \to 0^+} \frac{f(c+h) - f(c)}{h} \\
&= \lim_{h \to 0^+} f'(x_h) \\
&= L,
\end{aligned}
\]
and an analogous argument for $h \to 0^-$ shows that $f'_-(c) = L$ as well, so that $f'(c) = L$, as claimed.
\end{proof}

\begin{remark}
The value $L = \infty$ is also admissible in this theorem: in that case one cannot conclude that $f$ is differentiable at $c$, but the value $f'(c) = \infty$ is naturally interpreted as indicating the presence of a vertical tangent to the graph at the point $(c, f(c))$, as illustrated by the functions
\[
\begin{gathered}
f(x) = x^{1/3} \\
f(x) = \sqrt[3]{x^2(6-x)}.
\end{gathered}
\]

It is also worth noting that, as a consequence of the Darboux Property, the derivative of a function, regarded as a function in its own right, can never exhibit a jump discontinuity or a removable discontinuity.
\end{remark}

\subsection{Concavity and Inflection Points}

Beyond monotonicity, the shape of a graph is further characterized by its concavity, which describes whether the graph lies above or below its tangent lines.

\begin{definition}
Given a function $f$ that is differentiable on an interval $I$, we say that $f$ is convex, or concave upward, on $I$ if, for every $c \in I$, the graph of $f$ lies above the tangent line to the graph at the point $(c, f(c))$. Equivalently, convexity is characterized by the following inequality, required for every $x, c \in I$:
\[
f \text{ convex on } I \iff f(x) \geq f(c) + f'(c)(x-c).
\]

We say that $f$ is concave, or concave downward, on $I$ whenever the opposite inequality holds throughout $I$.
\end{definition}

\begin{definition}
A point $P = (c, f(c))$ on a curve $y = f(x)$ is called an inflection point if $f$ is continuous at $c$ and the curve changes from concave upward to concave downward, or from concave downward to concave upward, at $c$.
\end{definition}

It is worth noting that at an inflection point where $f$ is differentiable, the curve crosses its own tangent line, in contrast with the behavior at a generic point of the graph. The relationship between concavity and the monotonicity of the derivative is made precise by the following theorem.

\begin{theorem}
Let $f$ be differentiable on an interval $I$. Then the following two statements are equivalent: first, $f$ is convex on $I$; second, $f'$ is increasing on $I$.
\end{theorem}

\begin{proof}
We first show that the second statement implies the first. Take $c \in I$ and $x > c$. By Lagrange's Theorem applied to the interval $[c,x]$, there exists $y \in (c,x)$ such that
\[
\frac{f(x) - f(c)}{x - c} = f'(y).
\]
Since $f'$ is increasing, we have $f'(y) \geq f'(c)$, and hence
\[
\begin{aligned}
f(x) - f(c) &= f'(y)(x-c) \\
&\geq f'(c)(x-c),
\end{aligned}
\]
so that
\[
f(x) \geq f'(c)(x-c) + f(c),
\]
which is precisely the definition of convexity at $c$, from the right; the case $x < c$ is treated analogously, establishing convexity on all of $I$.

We now show that the first statement implies the second. Let $x < y$ be two points of $I$. By the definition of convexity, the value $f(y)$ is greater than the value of the tangent line to the graph at $(x, f(x))$ evaluated at $y$, that is,
\[
f(y) \geq f(x) + f'(x)(y-x), \qquad \text{hence} \qquad f'(x) \leq \frac{f(y) - f(x)}{y - x}.
\]

Similarly, the value $f(x)$ is greater than the value of the tangent line to the graph at $(y, f(y))$ evaluated at $x$, that is,
\[
f(x) \geq f(y) + f'(y)(x-y), \qquad \text{hence} \qquad f'(y) \geq \frac{f(y) - f(x)}{y - x}.
\]

Combining the two inequalities, we immediately deduce that $f'(x) \leq f'(y)$, which shows that $f'$ is increasing on $I$ and completes the proof of the equivalence.
\end{proof}

As a direct consequence of this equivalence, one obtains a practical criterion for concavity expressed in terms of the second derivative.

\begin{theorem}[Concavity Test]
Let $f$ be twice differentiable on an interval $I$. If $f'' > 0$ on $I$, then $f$ is convex on $I$. If $f'' < 0$ on $I$, then $f$ is concave, or concave downward, on $I$.
\end{theorem}

In view of this test, it follows that a point of inflection can occur at any point where the second derivative changes sign; more precisely, inflection points are to be sought among those points $c$ at which either $f''(c) = 0$ or $f$ fails to be twice differentiable. These principles, together with the Increasing/Decreasing Test, the First Derivative Test, and the Closed Interval Method discussed earlier, provide a complete set of tools for sketching the graph of a given function, as may be illustrated by sketching the graphs of
\[
\begin{gathered}
g(x) = x^4 - 4x^3 \\
f(x) = \sqrt[3]{x^2(6-x)}.
\end{gathered}
\]

\section{Taylor Approximation and Asymptotic Expansions}

\subsection{Local Polynomial Approximation}

A fundamental geometric fact about differentiable functions is that, upon zooming in toward a point on their graph, the graph increasingly resembles its tangent line at that point. This observation suggests that the tangent line provides a natural first-order approximation to a differentiable function near a given point, and it is this idea that we now develop rigorously. Let $f$ be differentiable at a point $a$, and consider the tangent line approximation
\[
L(x) = f(a) + f'(a)(x-a).
\]

This linear function is, in a precise sense, the best linear approximation to $y = f(x)$ near $x = a$: among all first-degree polynomials of the form $P(x) = c_0 + c_1 x$, the function $L(x)$ is the one that fits the graph of $f$ most closely near the point $(a, f(a))$. This distinguished status is characterized algebraically by the fact that $L(x)$ is the unique first-degree polynomial satisfying the two conditions $P(a) = f(a)$ and $P'(a) = f'(a)$, that is, the unique linear polynomial that agrees with $f$ both in value and in first derivative at the point $a$.

The precise sense in which $L(x)$ approximates $f(x)$ is clarified by the following consequence of the very definition of differentiability.

\begin{theorem}
Let $f$ be differentiable at $a$. Then
\[
\lim_{x \to a} \frac{f(x) - L(x)}{x - a} = 0.
\]

\end{theorem}

This result may be interpreted as saying that the error
\[
E(x) = f(x) - L(x)
\]
approaches zero faster than $(x-a)$ itself as $x \to a$; adopting the standard Landau notation, we write $E(x) = o(x-a)$ as $x \to a$. With this notation, the conclusion of the theorem may be restated as
\[
f(x) = f(a) + f'(a)(x-a) + o(x-a), \qquad x \to a,
\]
or, equivalently, setting $x = a + h$,
\[
f(a+h) = f(a) + f'(a) h + o(h), \qquad h \to 0.
\]

Having established that the tangent line provides a first-order approximation to $f$ near $a$, it is natural to ask whether a better approximation can be obtained by allowing a higher-degree polynomial. Specifically, one may attempt to approximate the curve $y = f(x)$ near $(a, f(a))$ not by a straight line but by a parabola, that is, by a second-degree, or quadratic, approximation of the form
\[
P(x) = c_0 + c_1 x + c_2 x^2.
\]

A concrete numerical experiment with the function $y = \cos x$ centered at $a = 0$ illustrates vividly how such a quadratic approximation improves upon the linear one.

Inspired by this example, we search for a polynomial $P(x)$ satisfying the three conditions $P(a) = f(a)$, $P'(a) = f'(a)$, and
\[
P''(a) = f''(a).
\]

It turns out that there is a unique quadratic function satisfying these three conditions simultaneously, namely
\[
T_2(x) = f(a) + f'(a)(x-a) + \frac{1}{2} f''(a) (x-a)^2.
\]

This function is called the second-degree Taylor polynomial of $f$ centered at $a$, and it is customarily denoted by $T_2(x)$. A graphical comparison, for instance between the function $f(x) = \cos x$ and its second-degree Taylor polynomial centered at $a = 0$, makes it apparent that the quadratic approximation is considerably more accurate than the linear one over a wider range of values of $x$; this observation is not merely qualitative but reflects a precise general fact, stated in the following theorem.

\begin{theorem}[Taylor--Peano, Degree Two]
If $f$ is twice differentiable at $a$, then
\[
\lim_{x \to a} \frac{f(x) - T_2(x)}{(x-a)^2} = 0.
\]

\end{theorem}

As in the linear case, this result may be interpreted in terms of the error term: the error
\[
E_2(x) = f(x) - T_2(x)
\]
approaches zero faster than $(x-a)^2$ as $x \to a$, that is,
\[
E_2(x) = o((x-a)^2)
\]
as $x \to a$. The conclusion of the theorem may accordingly be rewritten as
\[
f(x) = f(a) + f'(a)(x-a) + \frac{1}{2} f''(a)(x-a)^2 + o\big((x-a)^2\big), \qquad x \to a,
\]
or, setting $x = a+h$,
\[
f(a+h) = f(a) + f'(a) h + \frac{1}{2} f''(a) h^2 + o(h^2), \qquad h \to 0.
\]

\begin{proof}
We prove the result for $a = 0$; the general case follows by an entirely analogous argument after a translation of the variable. Setting
\[
g(x) = f(x) - T_2(x),
\]
the thesis becomes
\[
\lim_{x \to 0} g(x)/x^2 = 0.
\]
Observe that, by the very construction of $T_2$, we have
\[
\begin{aligned}
g(0) &= g'(0) \\
&= g''(0) \\
&= 0.
\end{aligned}
\]

As a first step, we apply Lagrange's Theorem to $g$ on the interval between $0$ and $x$: there exists $\tau \in (0,1)$ such that
\begin{equation*}
\frac{g(x) - g(0)}{x - 0} = g'(\tau x) \implies g(x) = g'(\tau x) \, x. \tag{1}
\end{equation*}
As a second step, we exploit the linear approximation of the function $g'$ at the point $0$, namely
\[
g'(y) = g'(0) + g''(0) y + o(y), \qquad y \to 0.
\]
Setting $y = \tau x$ and recalling that
\[
\begin{aligned}
g'(0) &= g''(0) \\
&= 0,
\end{aligned}
\]
we find
\begin{equation*}
g'(\tau x) = o(\tau x). \tag{2}
\end{equation*}
Combining $(1)$ and $(2)$, we obtain $g(x) = o(\tau x) \, x$, and hence
\[
\lim_{x \to 0} \frac{g(x)}{x^2} = \lim_{x \to 0} \frac{o(\tau x) \, x}{x^2} = \lim_{x \to 0} \frac{o(\tau x)}{\tau x} \cdot \tau = 0,
\]
which completes the proof.
\end{proof}

Figure~\ref{fig:ch5-taylor} compares a function with its linear and quadratic approximations near the expansion point.

\begin{figure}[!htbp]
\centering
\begin{tikzpicture}
  \begin{axis}[
    width=12.5cm,height=7.2cm,
    axis lines=middle,
    xlabel={$x$},ylabel={$y$},
    xmin=-2.6,xmax=2.6,ymin=-2.35,ymax=1.45,
    xtick={-2,0,2},xticklabels={$-2$,$a=0$,$2$},
    ytick={-2,-1,1},
    ticklabel style={font=\small},
    label style={font=\small},
    legend style={
      at={(0.5,1.03)},anchor=south,
      legend columns=3,
      font=\scriptsize,draw=black!30,fill=white,
      legend cell align=left,column sep=0.8em
    }
    ]
    \addplot[figblue,very thick,domain=-2.5:2.5,samples=180] {cos(deg(x))};
    \addlegendentry{$f(x)=\cos x$}
    \addplot[figred,thick,dashed,domain=-2.5:2.5] {1};
    \addlegendentry{$T_1(x)=1$}
    \addplot[figgreen,thick,dash dot,domain=-2.5:2.5,samples=180] {1-x^2/2};
    \addlegendentry{$T_2(x)=1-x^2/2$}
    \node[circle,fill=black,inner sep=1.25pt] at (axis cs:0,1) {};
  \end{axis}
\end{tikzpicture}
\caption[Linear and quadratic Taylor approximation]{Comparison between $f(x)=\cos x$, its
  first-degree Taylor polynomial $T_1(x)=1$, and its second-degree Taylor polynomial
  $T_2(x)=1-x^2/2$, all centered at $a=0$. The quadratic polynomial follows the curvature of $f$ and therefore
  remains accurate on a wider neighborhood of the expansion point.}
\label{fig:ch5-taylor}
\end{figure}

\subsection{Taylor's Theorem and Remainder Formulas}

The construction of the second-degree Taylor polynomial generalizes naturally to any degree $n$, provided that the function under consideration is differentiable a sufficient number of times.

\begin{definition}
Let $f$ be differentiable $n$ times at $a$. The polynomial
\[
T_n(x) := f(a) + f'(a)(x-a) + \frac{f''(a)}{2}(x-a)^2 + \cdots + \frac{f^{(n)}(a)}{n!}(x-a)^n
\]
is called the $n$-th degree Taylor polynomial of $f$ centered at $a$. In the special case $a = 0$, it is also called the $n$-th degree Maclaurin polynomial of $f$.
\end{definition}

By construction, the polynomial $T_n$ satisfies
\[
T_n^{(k)}(a) = f^{(k)}(a)
\]
for every $k = 1, \dots, n$, and it is in fact the unique polynomial of degree $n$ possessing this property. As the degree $n$ increases, the corresponding Taylor polynomial provides an increasingly accurate approximation to $f$ near $a$, a phenomenon that becomes visually apparent when several Taylor polynomials of increasing degree are plotted together against the original function.

This improvement in accuracy with increasing degree is made precise by the following generalization of the Taylor--Peano theorem stated above for the case $n=2$.

\begin{theorem}[Taylor--Peano Formula]
Let $f$ be differentiable $n$ times at $a$. Then
\[
\lim_{x \to a} \frac{f(x) - T_n(x)}{(x-a)^n} = 0.
\]

\end{theorem}

As before, this result may be phrased in terms of the error term: the error
\[
E_n(x) = f(x) - T_n(x)
\]
approaches zero faster than $(x-a)^n$ as $x \to a$, that is,
\[
E_n(x) = o((x-a)^n)
\]
as $x \to a$. The conclusion of the theorem is thus written as
\[
f(x) = f(a) + f'(a)(x-a) + \cdots + \frac{1}{n!} f^{(n)}(a)(x-a)^n + o\big((x-a)^n\big), \qquad x \to a,
\]
or, setting $x = a + h$,
\[
f(a+h) = f(a) + f'(a) h + \cdots + \frac{1}{n!} f^{(n)}(a) h^n + o(h^n), \qquad h \to 0.
\]

An instructive feature of the Taylor--Peano formula is that the remainder term $o(x^n)$ appearing at one order can itself be expanded further, revealing additional terms of the expansion at the next order; in this sense, one may zoom into the approximation as much as desired. For instance, as $x \to 0$ one has the first-order expansion $\sin x = x + o(x)$; zooming further into the remainder reveals that
\[
o(x) = -(1/6) x^3 + o(x^3),
\]
so that
\[
\sin x = x - \frac{1}{6} x^3 + o(x^3).
\]
Zooming in once again, the next term shows that
\[
o(x^3) = (1/5!) x^5 + o(x^5),
\]
yielding the still more refined expansion
\[
\sin x = x - \frac{1}{6} x^3 + \frac{1}{120} x^5 + o(x^5).
\]

This technique of successively refined Maclaurin expansions provides a powerful and systematic method for computing limits presenting an indeterminate form. As an illustration, consider the limit
\[
\lim_{x \to 0} (\sin x - x)/x^3.
\]
The Maclaurin expansion of $\sin x$ to third order gives
\[
\sin x = x - (1/6) x^3 + o(x^3),
\]
and hence
\[
\begin{aligned}
\lim_{x \to 0} \frac{\sin x - x}{x^3} &= \lim_{x \to 0} \frac{-\frac{1}{6} x^3 + o(x^3)}{x^3} \\
&= -\frac{1}{6} + \lim_{x \to 0} \frac{o(x^3)}{x^3}.
\end{aligned}
\]

Recalling the very definition of the symbol "small $o$", the second limit on the right-hand side vanishes by definition, and we conclude that
\[
\lim_{x \to 0} \frac{\sin x - x}{x^3} = -\frac{1}{6}.
\]

The Taylor--Peano formula, discussed above, describes the local behavior of a function near a point $a$, in the sense that it controls the error only asymptotically as $x \to a$, without providing any explicit bound valid for $x$ at a fixed, finite distance from $a$. A complementary result, known as the Taylor--Lagrange formula, instead provides an explicit expression for the remainder term, valid throughout an entire interval, and is therefore of a global rather than purely local character.

\begin{theorem}[Taylor--Lagrange Formula]
Let $f$ be differentiable $n+1$ times on an interval $[a,x]$. Then there exists a point $z \in (a,x)$ such that
\[
\begin{aligned}
f(x) &= T_n(x) + E_n(x), \\
E_n(x) &= \frac{f^{(n+1)}(z)}{(n+1)!} (x-a)^{n+1}.
\end{aligned}
\]

The expression for $E_n(x)$ given above is called the Lagrange form of the remainder term, in contrast with the Peano form
\[
E_n(x) = o((x-a)^n)
\]
furnished by the earlier theorem.
\end{theorem}

The global character of the Taylor--Lagrange formula becomes particularly valuable whenever an upper bound is available for the derivatives of $f$ throughout the interval in question, since in that case the Lagrange remainder can be estimated explicitly, yielding quantitative error bounds rather than merely qualitative asymptotic information. This feature is exploited decisively in the following application.

\subsection{Taylor Expansions of Elementary Functions}

It follows from the general theory developed above that, for every $n \in \mathbb{N}$,
\[
e^x = \sum_{k=0}^{n} \frac{x^k}{k!} + o(x^n), \qquad x \to 0,
\]
which shows that the Maclaurin polynomials of the exponential function reconstruct $e^x$ near $a = 0$ with arbitrary precision, in the local sense described by the Taylor--Peano formula. A far more striking fact, however, is that the Maclaurin polynomials of $e^x$ reconstruct the function not merely near $a=0$, but throughout its entire domain, in the sense made precise by the following theorem.

\begin{theorem}
For every $x \in \mathbb{R}$,
\[
e^x = \sum_{n=0}^{\infty} \frac{x^n}{n!}.
\]

\end{theorem}

Setting $x = 1$ in this identity yields the celebrated relation between Euler's number and the factorial series,
\[
e = \sum_{n=0}^{\infty} \frac{1}{n!}.
\]

\begin{proof}
Let $x > 0$ be fixed; the case $x \leq 0$ is handled by an entirely analogous argument. By the Taylor--Lagrange formula applied to $f(x) = e^x$ centered at $a = 0$, we have
\begin{equation*}
e^x - T_n(x) = \frac{f^{(n+1)}(z)}{(n+1)!} x^{n+1} \tag{$*$}
\end{equation*}
for some $z$ with $0 < z < x$, where $f^{(k)}(z) = e^z$ for every $k \geq 0$, and
\[
T_n(x) = \sum_{k=0}^{n} x^k / k!.
\]
Since $z < x$, we have $e^z < e^x$, and from $(*)$ we obtain the bound
\[
|e^x - T_n(x)| \leq e^x \, \frac{x^{n+1}}{(n+1)!}.
\]
Recalling that
\[
\lim_{n \to \infty} q^{n+1}/(n+1)! = 0
\]
for every fixed real number $q$, a fact established in the theory of sequences, we conclude that
\[
\lim_{n \to \infty} [e^x - T_n(x)] = 0,
\]
Equivalently, the limiting representation is
\[
\begin{aligned}
e^x &= \lim_{n \to \infty} T_n(x) \\
&= \sum_{n=0}^{\infty} \frac{x^n}{n!},
\end{aligned}
\]
where the last equality follows directly from the definition of series as the limit of its partial sums.
\end{proof}

The phenomenon just illustrated for the exponential function is an instance of a more general theory, that of Taylor series, or Maclaurin series in the case $a = 0$: for a sufficiently regular function $f$, the sequence of Maclaurin polynomials can reconstruct the entire function over its domain, not merely a local approximation near a single point. A useful intuitive image for this phenomenon is that the Taylor series of a function discovers, so to speak, the mathematical DNA encoded within the function at a single point, and allows the entire function to be rebuilt from this single piece of local information. This intuition is well illustrated by the sequence of Taylor polynomials of increasing degree for the function $y = \sin x$ centered at $a = 0$, which approximate the sine function over an ever wider range of values of $x$ as the degree increases.

The systematic use of Maclaurin expansions for the computation of limits, illustrated earlier for the limit
\[
\lim_{x \to 0} (\sin x - x)/x^3,
\]
rests on a small collection of algebraic manipulation rules governing the symbol $o(x^n)$. We now record the precise definition of this symbol, together with its basic algebraic properties.

\begin{definition}
Given $n \in \mathbb{N}$, the symbol $o(x^n)$, read "small $o$ of $x^n$", denotes a generic, unspecified quantity depending on $x$, vanishing as $x \to 0$, with the property that
\[
\lim_{x \to 0} \frac{o(x^n)}{x^n} = 0.
\]

\end{definition}

As $x \to 0$, the symbol $o(x^n)$ obeys the following algebraic rules, which together constitute what is customarily called the algebra of small $o$. First, $x^n = o(x^m)$ for any integer $m < n$, reflecting the fact that a higher power of $x$ is itself negligible compared to a lower power as $x \to 0$. Second, $c \cdot o(x^n) = o(x^n)$ for any constant $c$, so that multiplying a small-$o$ quantity by a fixed constant does not change its order. Third,
\[
o(x^n) \pm o(x^n) = o(x^n),
\]
so that the sum or difference of two quantities of the same order remains of that order. Fourth,
\[
x^n \cdot o(x^m) = o(x^{n+m}),
\]
so that multiplying a small-$o$ quantity by an explicit power increases its order accordingly. Fifth,
\[
o(x^n) \cdot o(x^m) = o(x^{n+m}),
\]
extending the previous rule to the product of two small-$o$ quantities. Sixth,
\[
(o(x^n))^m = o(x^{nm}),
\]
governing the behavior of small-$o$ quantities under exponentiation. Seventh,
\[
o(x^n + o(x^n)) = o(x^n),
\]
which shows that the symbol $o(x^n)$ is stable under composition with itself.

\subsection{Applications to Indeterminate Forms}

The algebra of small $o$ underlies a general method for computing limits of an indeterminate form $0/0$, of the type
\[
\lim_{x \to 0} f(x)/g(x).
\]

If $f$ and $g$ possess derivatives of every order at $x = 0$, the Maclaurin theorem allows one to write $f(x) = a x^n + o(x^n)$ for some integer $n$ and some nonzero constant $a$, meaning that $f$ is an infinitesimal of order $n$ with leading part $a x^n$, and analogously $g(x) = b x^m + o(x^m)$ for suitable $m$ and nonzero $b$. Substituting these expansions into the quotient, one obtains
\[
\begin{aligned}
\lim_{x \to 0} \frac{f(x)}{g(x)} &= \lim_{x \to 0} \frac{a x^n + o(x^n)}{b x^m + o(x^m)} \\
&= \lim_{x \to 0} \frac{x^n \left( a + \dfrac{o(x^n)}{x^n} \right)}{x^m \left( b + \dfrac{o(x^m)}{x^m} \right)},
\end{aligned}
\]
and it then suffices to apply the definition of small $o$, according to which the two fractions $o(x^n)/x^n$ and $o(x^m)/x^m$ both tend to zero as $x \to 0$, in order to resolve the limit explicitly in terms of $n$, $m$, $a$, and $b$.

\begin{remark}[Exercises on Taylor expansions and limits]
The techniques developed in this section admit a broad range of applications, several of which are proposed here as exercises for the reader. One may first show, for $x \to 0$ and any two integers $m < n$, that $x^n = o(x^m)$, and subsequently deduce the multiplicative rule
\[
x^n \cdot o(x^m) = o(x^{n+m})
\]
directly from the definition of small $o$. A second exercise consists of finding the Taylor polynomial of degree six for $f(x) = \sin x$ centered at $a = \pi/3$, thereby practicing the construction of Taylor polynomials at a center other than the origin. A third group of exercises, relying on the algebra of small $o$ together with the standard table of Maclaurin expansions, asks the reader to find the Maclaurin polynomial of degree four for $f(x) = \ln(\cos x)$; to determine the order of infinitesimal and the leading part, as $x \to 0$, of the function
\[
f(x) = (\sinh x)^2 - \sin(x^2);
\]
and to compute the limit
\[
\lim_{x \to 0} \frac{\tan x - \sin x}{x^2 (x+1) + \ln(1 - x^2)}
\]
by means of suitable Maclaurin expansions of numerator and denominator. A further set of exercises concerns the qualitative behavior of graphs near points where several derivatives vanish. One may first sketch a local graph of $f(x) = \ln(\cos x)$ close to $x = 0$, using the Maclaurin expansion obtained above. A second, more theoretical exercise concerns the classification of critical points by means of higher-order derivatives: suppose that
\[
\begin{gathered}
f'(a) = \cdots = f^{(n-1)}(a) = 0 \\
\begin{aligned}
f^{(n)}(a) &= k \\
&\neq 0
\end{aligned}
\end{gathered}
\]
for some integer $n \geq 2$; one is asked to prove that if $n$ is odd, then $a$ is an inflection point with horizontal tangent, while if $n$ is even, then $a$ is a local minimum when $k > 0$ and a local maximum when $k < 0$. A useful hint for this exercise, in the illustrative case $n = 2$, is that it suffices to know the values $f(a)$, $f'(a)$, and $f''(a)$: indeed, if $f''(a) > 0$, then as $x \to a$ one may write
\[
\begin{aligned}
f(x) &= f(a) + f'(a)(x-a) + \frac{f''(a)}{2}(x-a)^2 + o\big((x-a)^2\big) \\
&= f(a) + f'(a)(x-a)
+ (x-a)^2 \left[ \frac{f''(a)}{2} + \frac{o\big((x-a)^2\big)}{(x-a)^2} \right],
\end{aligned}
\]
and since the bracketed term is positive for $x$ sufficiently close to $a$, it follows that
\[
f(x) \geq f(a) + f'(a)(x-a)
\]
near $a$, that is, the graph of $f$ stays above its tangent line, confirming that $a$ is a local minimum when $n=2$ and $k>0$. The general case follows a similar pattern: for $h$ small, one has
\[
f(a+h) - f(a) = h^n \left[ \frac{k}{n!} + \frac{o(h^n)}{h^n} \right],
\]
and the sign of the bracketed expression, which approaches $k/n!$ as $h \to 0$, together with the sign of $h^n$, which depends on the parity of $n$, together determine the nature of the critical point.

A final group of exercises draws together the theory of Taylor series with the theory of series developed earlier. Using the table of Maclaurin expansions, one may find the sum of the series
\[
\sum_{n=0}^{\infty} (-1)^n x^{4n}/n!,
\]
for an arbitrary real number $x$, and of the numerical series
\[
\sum_{n=1}^{\infty} (-1)^{n-1} 3^n/(n \cdot 5^n),
\]
by recognizing each as a Maclaurin series evaluated at a suitable point. A second exercise asks the reader to use Maclaurin polynomials together with the Taylor--Lagrange Theorem to compute the numerical value of $e^{\pi/3}$ with a precision of $10^{-2}$; the standard approach is to approximate $e^{\pi/3}$ by the partial sum
\[
T_n(x) = \sum_{k=0}^{n} x^k/k!
\]
evaluated at $x = \pi/3$, choosing $n$ large enough that the Lagrange remainder satisfies $|E_n| < 10^{-2}$, using the bound
\[
0 < E_n(\pi/3) \leq \max_{z \in [0, \pi/3]} \frac{f^{(n+1)}(z)}{(n+1)!} \left( \frac{\pi}{3} \right)^{n+1} \leq e^{\pi/3} \, \frac{1}{(n+1)!} \left( \frac{\pi}{3} \right)^{n+1},
\]
and selecting the smallest integer $n$ for which the right-hand side falls below $10^{-2}$. A final, more theoretical exercise asks the reader to prove that Euler's number $e$ is irrational, a classical result that can be established using the series representation of $e$ together with a careful estimate of the remainder term of its Taylor--Lagrange expansion.
\end{remark}

\newpage
\chapter{Integral Calculus}

\section{Area, Distance, and Integral Approximation}

\subsection{Approximating Areas under Nonnegative Curves}

Let $f : [a,b] \to \mathbb{R}$ be a nonnegative bounded function. Our goal is to compute, or more precisely to define in a rigorous way, the area of the region
\[
S = \{ (x,y) \in \mathbb{R}^2 : x \in [a,b], \ 0 \leq y \leq f(x) \}
\]
lying between the graph of $f$ and the horizontal axis. Before attempting a general definition, it is instructive to consider the concrete example $f(x) = x^2$ on the interval $[0,1]$, and to ask how the area beneath this parabola might be approximated.

The natural idea is to divide $S$ into $n$ vertical strips $S_i$ and to approximate each strip by a rectangle. Two immediate choices for the height of the approximating rectangle present themselves: one may take the height equal to the value of $f$ at the right endpoint of the strip, or equal to the value of $f$ at the left endpoint of the strip. These two choices produce two sequences of approximations, which we may denote by $R_n$ and $L_n$ respectively. As $n$ increases, both $L_n$ and $R_n$ appear, from geometric inspection, to become increasingly accurate approximations to the true area of $S$; this suggests that the area $A$ should be defined as the common limit of the sums of the areas of the approximating rectangles, as the number of strips tends to infinity.

More generally, each strip $S_i$ may be approximated by a rectangle having the same base as the strip, but whose height is given by $f(x_i^*)$, where $x_i^*$ denotes an arbitrary point chosen within the $i$-th subinterval $[x_{i-1}, x_i]$; this leads to the approximation
\[
A \approx \sum_{i=1}^{n} f(x_i^*) \, \Delta x.
\]

In the particular example under consideration, since $f$ is increasing, one has the sandwich inequality
\[
L_n \leq \sum_{i=1}^{n} f(x_i^*) \, \Delta x \leq R_n
\]
for every choice of sample points $x_i^*$, which shows that
\[
A = \lim_{n \to +\infty} \sum_{i=1}^{n} f(x_i^*) \, \Delta x
\]
regardless of the particular choice of the points $x_i^*$ within their respective subintervals.


We now apply this idea to a general region $S$ of the form introduced above. One begins by subdividing the interval $[a,b]$ into $n$ strips $S_i$ of equal width $\Delta x = (b-a)/n$, and one then approximates each strip $S_i$ by a rectangle having the same base as the strip and height equal to $f(x_i^*)$, where $x_i^*$ is an arbitrary point chosen within the $i$-th subinterval $[x_{i-1}, x_i]$; the points
\[
x_1^*, x_2^*, \ldots, x_n^*
\]
obtained in this way are called the sample points. It can be proved that the limit
\[
\lim_{n \to \infty} \sum_{i=1}^{n} f(x_i^*) \, \Delta x
\]
always exists, and it can furthermore be shown that this limit takes the same value regardless of the particular choice of sample points. On the basis of this result, one defines the area $A$ of the region $S$ to be the limit of the sums of the areas of the approximating rectangles, that is,
\[
A = \lim_{n \to \infty} \sum_{i=1}^{n} f(x_i^*) \, \Delta x.
\]

\subsection{The Distance Problem}

A closely related problem, arising from kinematics rather than geometry, is the distance problem: given the velocity of an object at every instant of a certain time period, one wishes to determine the total distance travelled by the object during that period. If the velocity remains constant throughout the interval, the distance problem is solved immediately by the elementary formula
\[
\text{distance} = \text{velocity} \times \text{time}.
\]
The situation becomes considerably more delicate when the velocity
\[
\begin{aligned}
v &= f(t) \\
&\geq 0,
\end{aligned}
\]
so that the object always moves in the positive direction, varies with time over the interval
\[
\begin{aligned}
a &\leq t \\
&\leq b.
\end{aligned}
\]

To address this more general situation, one measures the velocity on a collection of very small subintervals of $[a,b]$, of the form $[t_{i-1}, t_i]$; within each such subinterval, an approximate value for the velocity is provided by $f(t_i^*)$, for an arbitrarily chosen point
\[
t_i^* \in [t_{i-1}, t_i],
\]
so that the total distance travelled is approximated by the sum
\[
\sum_{i=1}^{n} f(t_i^*) \, \Delta t.
\]

As the velocity is measured more and more frequently, these estimates become increasingly accurate, which suggests that the exact distance $d$ travelled should be defined as the limit of this expression,
\[
d = \lim_{n \to \infty} \sum_{i=1}^{n} f(t_i^*) \, \Delta t.
\]

This limit is of precisely the same nature as the one encountered in the area problem, which motivates the introduction of a single, unified concept capable of encompassing both situations, namely the definite integral.

\section{The Definite Integral}

\subsection{Riemann Sums and the Definition of the Integral}


In order to formulate the definition of the definite integral with full rigor, several preliminary notions are required. A partition $P$ of an interval $[a,b]$ is a finite sequence of numbers of the form
\[
a = x_0 < x_1 < x_2 < \cdots < x_n = b;
\]
the interval $[x_{i-1}, x_i]$ is called the $i$-th subinterval of the partition. The mesh, or norm, of a partition is defined to be the length of its longest subinterval, that is, writing $\Delta x_i := x_i - x_{i-1}$, one sets
\[
m_P := \max\{ \Delta x_i : 1 \leq i \leq n \}.
\]

A tagged partition of an interval $[a,b]$ is a partition together with a finite sequence of sample points
\[
x_1^*, x_2^*, \ldots, x_n^*,
\]
subject to the requirement that $x_i^* \in [x_{i-1}, x_i]$ for every index $i$. Given a tagged partition, the Riemann sum of $f$ with respect to this partition, with nodes $x_0 < x_1 < \cdots < x_n$ and sample points $x_1^*, x_2^*, \ldots, x_n^*$, is the quantity
\[
\sum_{i=1}^{n} f(x_i^*) \, \Delta x_i.
\]

\begin{definition}[Integrable Function]
Let $f$ be a bounded function defined on $[a,b]$. We say that $f$ is integrable on $[a,b]$ if the limit
\[
\lim_{m_P \to 0} \sum_{i=1}^{n} f(x_i^*) \, \Delta x_i
\]
exists, is finite, and gives the same value for every tagged partition $P$ of $[a,b]$, regardless of the choice of subdivision points and sample points, provided only that the mesh $m_P$ tends to zero.
\end{definition}

\begin{definition}[Definite Integral]
Let $f$ be an integrable function on $[a,b]$. The definite integral of $f$ from $a$ to $b$ is the limit of the Riemann sums of $f$ over any possible tagged partition of $[a,b]$, as the partitions become arbitrarily fine, that is,
\[
\int_a^b f(x) \, dx := \lim_{m_P \to 0} \sum_{i=1}^{n} f(x_i^*) \, \Delta x_i.
\]

\end{definition}

It is important to state the precise meaning of the limit appearing in this definition, since the index set over which the limit is taken, namely the collection of all tagged partitions, is considerably more general than the index set $n \in \mathbb{N}$ used in the theory of sequences. Precisely, the definition asserts that for every $\varepsilon > 0$ there exists $\delta > 0$ such that, for every tagged partition $P$ with $m_P < \delta$, one has
\[
\left| \int_a^b f(x) \, dx - \sum_{i=1}^{n} f(x_i^*) \, \Delta x_i \right| < \varepsilon.
\]

Several remarks concerning notation and convention are in order. The definite integral is a real number, and in particular it does not depend on the variable $x$ used to describe the integrand; this variable is therefore called a dummy or bound variable. Although the definition of the integral implicitly assumes that $a < b$, the characterization as a limit of Riemann sums continues to make sense even when $b \leq a$; in that case, one adopts the conventions
\[
\int_b^a f(x) \, dx := -\int_a^b f(x) \, dx, \qquad \int_a^a f(x) \, dx = 0.
\]
When $f \geq 0$ on $[a,b]$, the definite integral satisfies
\[
\int_a^b f(x) \, dx \geq 0
\]
and can be interpreted geometrically as the area under the curve $y = f(x)$ between $a$ and $b$. When $f$ takes on both positive and negative values on $[a,b]$, its definite integral can instead be interpreted as a net area, that is, as a difference of areas,
\[
\int_a^b f(x) \, dx = A_1 - A_2,
\]
where $A_1$ denotes the area of the region lying above the horizontal axis and below the graph of $f$, and $A_2$ denotes the area of the region lying below the horizontal axis and above the graph of $f$. Finally, returning to the distance problem discussed earlier, if $v = f(t)$ is the velocity of an object, then the total distance travelled during a time interval
\[
\begin{aligned}
a &\leq t \\
&\leq b,
\end{aligned}
\]
regardless of whether the object reverses direction, is given by
\[
d = \int_a^b |f(t)| \, dt.
\]

Figure~\ref{fig:ch6-riemann-sums} compares left-endpoint, right-endpoint, and midpoint approximations to the same area.

\begin{figure}[!htbp]
\centering
\begin{tikzpicture}
  \begin{axis}[
    width=6.3cm,height=6cm,
    axis lines=middle, xlabel={$x$}, ylabel={},
    xmin=-0.3,xmax=4.3, ymin=-0.3,ymax=4.0,
    xtick={0,0.57,1.14,1.71,2.29,2.86,3.43,4}, xticklabels=\empty, ytick=\empty,
    ]
    \foreach \i in {0,...,6}{
      \pgfmathsetmacro{\xa}{\i*0.5714}
      \pgfmathsetmacro{\xb}{(\i+1)*0.5714}
      \pgfmathsetmacro{\h}{0.5+2.2*sin(deg(0.55*\xa))+0.3*\xa}
      \addplot[fill=figblue,opacity=0.3,draw=figblue] coordinates
        {(\xa,0) (\xa,\h) (\xb,\h) (\xb,0)} -- cycle;
    }
    \addplot[black,thick,domain=0:4,samples=150]{0.5+2.2*sin(deg(0.55*x))+0.3*x};
  \end{axis}
  \begin{axis}[
    at={(6.9cm,0)},
    width=6.3cm,height=6cm,
    axis lines=middle, xlabel={$x$}, ylabel={},
    xmin=-0.3,xmax=4.3, ymin=-0.3,ymax=4.0,
    xtick={0,0.57,1.14,1.71,2.29,2.86,3.43,4}, xticklabels=\empty, ytick=\empty,
    ]
    \foreach \i in {0,...,6}{
      \pgfmathsetmacro{\xa}{\i*0.5714}
      \pgfmathsetmacro{\xb}{(\i+1)*0.5714}
      \pgfmathsetmacro{\h}{0.5+2.2*sin(deg(0.55*\xb))+0.3*\xb}
      \addplot[fill=figred,opacity=0.3,draw=figred] coordinates
        {(\xa,0) (\xa,\h) (\xb,\h) (\xb,0)} -- cycle;
    }
    \addplot[black,thick,domain=0:4,samples=150]{0.5+2.2*sin(deg(0.55*x))+0.3*x};
  \end{axis}
  \begin{axis}[
    at={(0,-5.4cm)},
    width=6.3cm,height=6cm,
    axis lines=middle, xlabel={$x$}, ylabel={},
    xmin=-0.3,xmax=4.3, ymin=-0.3,ymax=4.0,
    xtick={0,0.57,1.14,1.71,2.29,2.86,3.43,4},
    xticklabels={$x_0$,$x_1$,$x_2$,$x_3$,$x_4$,$x_5$,$x_6$,$x_7$},
    xticklabel style={font=\tiny}, ytick=\empty,
    ]
    \foreach \i in {0,...,6}{
      \pgfmathsetmacro{\xa}{\i*0.5714}
      \pgfmathsetmacro{\xb}{(\i+1)*0.5714}
      \pgfmathsetmacro{\xm}{(\xa+\xb)/2}
      \pgfmathsetmacro{\h}{0.5+2.2*sin(deg(0.55*\xm))+0.3*\xm}
      \addplot[fill=figgreen,opacity=0.3,draw=figgreen] coordinates
        {(\xa,0) (\xa,\h) (\xb,\h) (\xb,0)} -- cycle;
    }
    \addplot[black,thick,domain=0:4,samples=150]{0.5+2.2*sin(deg(0.55*x))+0.3*x};
  \end{axis}
  \begin{axis}[
    at={(6.9cm,-5.4cm)},
    width=6.3cm,height=6cm,
    axis lines=middle, xlabel={$x$}, ylabel={},
    xmin=-0.3,xmax=4.3, ymin=-0.3,ymax=4.0, xtick=\empty, ytick=\empty,
    ]
    \addplot[black,thick,domain=0:4,samples=150]{0.5+2.2*sin(deg(0.55*x))+0.3*x};
    \addplot[fill=black,opacity=0.12,draw=none,domain=0:4,samples=150] {0.5+2.2*sin(deg(0.55*x))+0.3*x} \closedcycle;
  \end{axis}
\end{tikzpicture}
\caption[Riemann sums: left, right, midpoint]{Top left: left-endpoint sums (blue). Top right: right-endpoint sums (red). Bottom left: midpoint sums (green). All three use the same partition $x_0<\dots<x_7$. As the partition is refined, the sums converge to the area under the curve, shaded in the bottom-right panel.}
\label{fig:ch6-riemann-sums}
\end{figure}

\subsection{Which Functions Are Integrable?}

Not every bounded function is integrable in the sense just defined; the function that takes the value $0$ at every rational number and the value $1$ at every irrational number, restricted to $[0,1]$, provides a standard example of a bounded function that fails to be integrable, since the corresponding Riemann sums can be made arbitrarily close to either $0$ or $1$ depending on the choice of sample points, so that no common limiting value exists. Fortunately, the following theorem shows that the vast majority of functions encountered in practice are indeed integrable.

\begin{theorem}
If $f$ is continuous on $[a,b]$, or if $f$ is monotone on $[a,b]$, or if $f$ has only a finite number of jump discontinuities on $[a,b]$, then $f$ is integrable on $[a,b]$; that is, the definite integral
\[
\int_a^b f(x) \, dx
\]
exists.
\end{theorem}

Integrability is furthermore preserved under composition with a continuous function, as recorded in the following result.

\begin{theorem}[Composition]
If $f$ is integrable on $[a,b]$ and $\varphi$ is continuous on an interval containing $f([a,b])$, then the composition $\varphi \circ f$ is integrable on $[a,b]$.
\end{theorem}

In particular, this theorem shows that both $|f|$ and $f^2$ are integrable on $[a,b]$ whenever $f$ itself is integrable, since the absolute value function and the squaring function are continuous.

\subsection{Properties of the Definite Integral}

The definite integral obeys a number of algebraic properties, closely paralleling those established earlier for series, which greatly facilitate its computation. Let $f$ and $g$ be two integrable functions on $[a,b]$, and let $\lambda \in \mathbb{R}$ be an arbitrary constant. Then the functions $f + g$ and $\lambda \cdot f$ are also integrable on $[a,b]$, and one has the following linearity properties.
\begin{itemize}
\item \textit{Constant functions}: for the constant function equal to $\lambda$ throughout $[a,b]$, one has

\begin{align*}
\int_a^b \lambda \, dx = \lambda (b-a).
\end{align*}

\item \textit{Additivity of the integrand}: for the sum of two integrable functions,

\begin{align*}
\int_a^b [f(x) + g(x)] \, dx = \int_a^b f(x) \, dx + \int_a^b g(x) \, dx.
\end{align*}

\item \textit{Homogeneity}: for the product of an integrable function with a scalar,

\begin{align*}
\int_a^b \lambda f(x) \, dx = \lambda \int_a^b f(x) \, dx.
\end{align*}

\end{itemize}

\begin{remark}
The product $f \cdot g$ of two integrable functions on $[a,b]$ is itself an integrable function. This follows as a consequence of the linearity properties just stated, together with the algebraic identity
\[
f \cdot g = \frac{1}{2}(f+g)^2 - \frac{1}{2}(f^2 + g^2),
\]
which expresses the product as a linear combination of squares of integrable functions, each of which is integrable by the Composition Theorem.
\end{remark}

A further fundamental property describes how integrals over adjacent intervals combine.

\begin{theorem}[Additivity over Adjacent Intervals]
Let $c$ be a number such that $a < c < b$. Then $f$ is integrable on $[a,c]$ and on $[c,b]$, and
\[
\int_a^c f(x) \, dx + \int_c^b f(x) \, dx = \int_a^b f(x) \, dx.
\]

\end{theorem}

When the integrand possesses a symmetry property, the corresponding integral over a symmetric interval can often be simplified considerably, as recorded in the following theorem.

\begin{theorem}[Integrals of Symmetric Functions]
Let $f$ be integrable on $[-a,a]$. If $f$ is even, then
\[
\int_{-a}^{a} f(x) \, dx = 2 \int_0^a f(x) \, dx.
\]
If $f$ is odd, then
\[
\int_{-a}^{a} f(x) \, dx = 0.
\]

\end{theorem}

Finally, a collection of comparison properties allows one to bound the value of a definite integral without computing it explicitly.

\begin{itemize}
\item \textit{Nonnegativity}: if $f(x) \geq 0$ for every $x \in [a,b]$, then

\[
\int_a^b f(x) \, dx \geq 0.
\]

\item \textit{Monotonicity}: if $f(x) \geq g(x)$ for every $x \in [a,b]$, then

\[
\int_a^b f(x) \, dx \geq \int_a^b g(x) \, dx.
\]

\item \textit{Triangle inequality}: the integral of $f$ is bounded by the integral of its absolute value,

\begin{align*}
\left| \int_a^b f(x) \, dx \right| \leq \int_a^b |f(x)| \, dx.
\end{align*}

\item \textit{Boundedness}: if

\[
k \leq f(x) \leq K
\]
throughout $[a,b]$, then
\begin{align*}
k(b-a) \leq \int_a^b f(x) \, dx \leq K(b-a).
\end{align*}

\end{itemize}

Figure~\ref{fig:atlas-signed-area} distinguishes signed integration from total geometric area.

\begin{figure}[!htbp]
\centering
\begin{tikzpicture}[>=Stealth]
\begin{axis}[width=11cm,height=5.8cm,axis lines=middle,ticklabel style={font=\small},label style={font=\small},legend style={font=\scriptsize,draw=black!20,fill=white},xlabel={$x$},ylabel={$\sin x$},xmin=0,xmax=6.6,ymin=-1.3,ymax=1.3,xtick={0,3.14159,6.28318},xticklabels={$0$,$\pi$,$2\pi$},ytick={-1,1}]
\addplot[draw=none,fill=figblue!25,domain=0:pi,samples=80] {sin(deg(x))} \closedcycle;
\addplot[draw=none,fill=figred!25,domain=pi:2*pi,samples=80] {sin(deg(x))} \closedcycle;
\addplot[figblue,very thick,domain=0:2*pi,samples=160] {sin(deg(x))};
\node at (axis cs:1.57,0.45) {$+$};
\node at (axis cs:4.71,-0.45) {$-$};
\end{axis}
\end{tikzpicture}
\caption{The graph of $\sin x$ has equal positive and negative areas on $[0,2\pi]$. Its signed integral is zero, whereas the integral of its absolute value is $4$. Blue and red shading indicate the opposite signs of the contributions.}
\label{fig:atlas-signed-area}
\end{figure}

\subsection{The Integral Mean Value Theorem}

A particularly useful consequence of the comparison properties, combined with the Intermediate Value Theorem for continuous functions, is the following result, which asserts that a continuous function attains its own average value at some point of the interval.

\begin{theorem}[Integral Mean Value Theorem]
Let $f$ be continuous on $[a,b]$. Then there exists $c \in [a,b]$ such that
\[
f(c) = \frac{1}{b-a} \int_a^b f(x) \, dx.
\]
This value is called the integral mean value of $f$ on $[a,b]$.
\end{theorem}

\begin{proof}
Let $m$ and $M$ denote respectively the minimum and maximum values of $f$ on $[a,b]$, guaranteed to exist by Weierstrass' Theorem, so that
\[
\begin{aligned}
m &\leq f(x) \\
&\leq M
\end{aligned}
\]
for every $x \in [a,b]$. By the boundedness property of the definite integral, it follows that
\[
m(b-a) \leq \int_a^b f(x) \, dx \leq M(b-a),
\]
and hence, setting
\[
y := \frac{1}{b-a} \int_a^b f(x) \, dx,
\]
we obtain
\[
\begin{aligned}
m &\leq y \\
&\leq M.
\end{aligned}
\]

Applying the Intermediate Value Theorem for continuous functions on $[a,b]$, we deduce the existence of a point $c \in [a,b]$ such that $f(c) = y$, which is precisely the desired conclusion.
\end{proof}

\section{The Fundamental Theorem of Calculus}

\subsection{The Integral Function and Its Derivative}

We now come to the central result of the entire theory, which reveals a deep and unexpected connection between the operations of differentiation and integration. Let $f$ be integrable on $[a,b]$, and fix any point $x \in [a,b]$. The quantity
\[
\int_a^x f(t) \, dt
\]
is a well-defined real number depending on the choice of $x$; letting $x$ vary throughout $[a,b]$ defines a new function. For every $x \in [a,b]$, set
\[
g(x) = \int_a^x f(t) \, dt.
\]

\begin{definition}[Integral Function]
The map $x \mapsto g(x)$ defined above is called the integral function, or area function, of $f$.
\end{definition}

The first part of the Fundamental Theorem of Calculus asserts that, whenever $f$ is continuous, its integral function is not merely continuous but is in fact differentiable, and that its derivative recovers the original function $f$.

\begin{theorem}[Fundamental Theorem of Calculus, Part I]
If $f$ is continuous on $[a,b]$, then the area function
\[
g(x) = \int_a^x f(t) \, dt
\]
is continuous and differentiable on $[a,b]$, and its derivative is given by $g'(x) = f(x)$ for every $x \in [a,b]$.
\end{theorem}

\begin{proof}
Fix $x$ and $h$ such that $x+h \in [a,b]$; we assume without loss of generality that $h > 0$. By the additivity property of the definite integral, we may write
\[
\begin{aligned}
g(x+h) - g(x)
&= \int_a^{x+h} f(t) \, dt - \int_a^x f(t) \, dt \\
&= \left[ \int_a^x f(t) \, dt + \int_x^{x+h} f(t) \, dt \right]
- \int_a^x f(t) \, dt \\
&= \int_x^{x+h} f(t) \, dt.
\end{aligned}
\]

Since $f$ is continuous, we may apply the Integral Mean Value Theorem on the interval $[x, x+h]$, obtaining a point $c_h \in [x, x+h]$ such that
\[
\int_x^{x+h} f(t) \, dt = f(c_h) \cdot h.
\]

As $h \to 0$, the point $c_h$, being squeezed between $x$ and $x+h$, satisfies $c_h \to x$; since $f$ is continuous, it follows that $f(c_h) \to f(x)$. Therefore,
\[
\lim_{h \to 0} \frac{g(x+h) - g(x)}{h} = \lim_{h \to 0} \frac{1}{h} \int_x^{x+h} f(t) \, dt = \lim_{h \to 0} f(c_h) = f(x).
\]
This shows that $g$ is differentiable at $x$ and that $g'(x) = f(x)$, as claimed.
\end{proof}

Using Leibniz notation for derivatives, the conclusion of the first part of the Fundamental Theorem of Calculus may be restated compactly as follows: if $f$ is continuous on $[a,b]$, then for every $x \in (a,b)$,
\[
\frac{d}{dx} \int_a^x f(t) \, dt = f(x).
\]

This identity is of great practical importance, since it means that the entire apparatus of differential calculus can be brought to bear on the study of the area function $g$, even when $g$ itself cannot be expressed in terms of elementary functions.

A useful extension of this result concerns integral functions whose limits of integration are themselves differentiable functions of $x$, rather than $x$ itself; given a continuous function $f$ and two differentiable functions $k$ and $h$, one may seek a formula for the derivative of
\[
\int_{h(x)}^{k(x)} f(t) \, dt,
\]
which can be obtained by combining the Fundamental Theorem of Calculus with the Chain Rule and the additivity property of the integral. It is also worth noting that continuity of $f$, while sufficient, is not strictly necessary for the continuity of the area function: if $f$ is merely assumed to be integrable, but not necessarily continuous, on $[a,b]$, one can still prove that the corresponding area function $g$ remains continuous on $[a,b]$, although in general it need no longer be differentiable.

As applications of these ideas, consider the sine integral
\[
\operatorname{Si}(x) = \int_0^x (\sin t)/t \, dt,
\]
where the integrand is continuously extended at $t=0$ by the value $1$.
The Fundamental Theorem gives
\[
\operatorname{Si}'(x)=\frac{\sin x}{x}\quad(x\neq0),
\qquad \operatorname{Si}'(0)=1.
\]
For the function
\[
G(x)=\int_0^{x^4} \sin(t^2) \, dt,
\]
the upper limit is $u(x)=x^4$. Combining the Fundamental Theorem with
the Chain Rule yields
\[
G'(x)=\sin\bigl((x^4)^2\bigr)\,4x^3
=4x^3\sin(x^8).
\]
Finally, let
\[
F(x) = \int_0^x f(t) \, dt
\]
for a continuous function $f$. If $f$ is even, then the substitution
$t=-u$ gives $F(-x)=-F(x)$, so $F$ is odd. If $f$ is odd, the same
substitution gives $F(-x)=F(x)$, so $F$ is even.

\subsection{Antiderivatives}

The Fundamental Theorem of Calculus naturally leads to the notion of an antiderivative, which plays a central role in the practical evaluation of definite integrals.

\begin{definition}
We say that $F$ is an antiderivative of $f$ on an interval $I$ if $F$ is differentiable on $I$ and $F'(x) = f(x)$ for every $x \in I$.
\end{definition}

\begin{remark}
If $F$ is an antiderivative of $f$ on $I$, then for any constant $c \in \mathbb{R}$ one has
\[
\begin{aligned}
\frac{d}{dx} [F(x) + c] &= \frac{d}{dx} [F(x)] + \frac{d}{dx} c \\
&= f(x) + 0 \\
&= f(x),
\end{aligned}
\]
so that $F + c$ is also an antiderivative of $f$ on $I$. Thus, whenever $f$ admits an antiderivative, it in fact admits infinitely many, which shows that the antiderivative of a given function is never unique.
\end{remark}

The following theorem shows that, on an interval, this non-uniqueness is fully accounted for by the addition of an arbitrary constant, and no other source of ambiguity is present.

\begin{theorem}
Given any two antiderivatives $F$ and $G$ of the same function $f$ on an interval $I$, the difference $F - G$ is constant on $I$.
\end{theorem}

\begin{proof}
For every $x \in I$ we have
\[
\begin{aligned}
\frac{d}{dx} [F(x) - G(x)] &= \frac{d}{dx}[F(x)] - \frac{d}{dx}[G(x)] \\
&= f(x) - f(x) \\
&= 0.
\end{aligned}
\]

We have thus shown that the function $F - G$ has vanishing derivative throughout the interval $I$, and hence, by the Increasing/Decreasing Test applied in both directions, $F - G$ must be constant on $I$; that is, there exists a constant $C \in \mathbb{R}$ such that $G(x) - F(x) = C$ for every $x \in I$.
\end{proof}

This theorem may be rephrased as follows: if $F$ is an antiderivative of $f$ on the interval $I$, then every other antiderivative of $f$ on $I$ is of the form $G(x) = F(x) + C$ for some arbitrary constant $C$.

\begin{remark}
The conclusion of this theorem is false, in general, on a domain that is not an interval. For instance, the most general antiderivative of $f(x) = -1/x^2$, defined for $x \neq 0$, is given by
\[
F(x) =
\begin{cases}
\dfrac{1}{x} + c_1 & x < 0, \\[1mm]
\dfrac{1}{x} + c_2 & x > 0,
\end{cases}
\]
for arbitrary and, crucially, independent constants $c_1, c_2 \in \mathbb{R}$, since the domain $\mathbb{R} \setminus \{0\}$ consists of two disjoint intervals rather than a single connected interval.
\end{remark}

\subsection{Existence of Antiderivatives}

It is natural to ask whether every function admits an antiderivative, and the Fundamental Theorem of Calculus provides a complete answer for the class of continuous functions: every continuous function $f$ has an antiderivative, and one such antiderivative $F$ is given by the definite integral of $f$ with variable upper limit,
\[
F(x) = \int_a^x f(t) \, dt,
\]
precisely as established in the first part of the Fundamental Theorem of Calculus. It should be emphasized, however, that although such an antiderivative always exists for a continuous function, it cannot always be expressed in terms of elementary functions; well-known examples of continuous functions whose antiderivatives escape elementary expression include $e^{-x^2}$, $\sin(x^2)$, $(\sin x)/x$, and $1/\ln x$.

It is also possible for non-continuous functions to possess antiderivatives, as illustrated by the function
\[
F(x) =
\begin{cases}
x^2 \sin(1/x) & x \neq 0, \\
0 & x = 0.
\end{cases}
\]
Computing its derivative directly from the definition, one finds
\[
F'(x) =
\begin{cases}
2x \sin(1/x) - \cos(1/x) & x \neq 0, \\
0 & x = 0.
\end{cases}
\]

The function $f(x) := F'(x)$ obtained in this way is not continuous at $x = 0$, since $\cos(1/x)$ oscillates without limit as $x \to 0$; nonetheless, $F$ is by construction an antiderivative of $f$ on every interval containing $x = 0$, showing that continuity of $f$ is sufficient, but not necessary, for the existence of an antiderivative.

On the other hand, not every function admits an antiderivative: a function with a jump discontinuity, such as
\[
f(x) =
\begin{cases}
-1 & x < 0, \\
1 & x \geq 0,
\end{cases}
\]
cannot be the derivative of any function on an interval containing the origin. This impossibility is a direct consequence of the following remarkable theorem, which shows that derivatives, even of merely differentiable functions, always satisfy an intermediate value property analogous to that of continuous functions, even when the derivative itself fails to be continuous.

\begin{theorem}[Darboux's Theorem]
Let $F : [a,b] \to \mathbb{R}$ be a differentiable function, with derivative function $F'(x) = f(x)$ for every $x \in [a,b]$. Then $f$ has the intermediate value property, that is, if $\lambda$ is any real number lying between $f(a)$ and $f(b)$, then there exists a point $c$ between $a$ and $b$ such that $f(c) = \lambda$.
\end{theorem}

\begin{proof}
Assume that $f(a) < f(b)$, and take any $\lambda$ such that $f(a) < \lambda < f(b)$; the case $f(a) > f(b)$ is treated analogously. Define the auxiliary function $H(x) = F(x) - \lambda x$ on $[a,b]$, whose derivative is $H'(x) = f(x) - \lambda$. Since $H$ is continuous on $[a,b]$, being differentiable, it attains its absolute minimum somewhere on $[a,b]$ by Weierstrass' Theorem; we claim that this minimum is attained at an interior point $c$ of $(a,b)$, rather than at one of the two endpoints. Indeed,
\[
H'(a) = f(a) - \lambda < 0
\]
shows that $H$ is decreasing near $a$, so the minimum cannot occur at $a$, while
\[
H'(b) = f(b) - \lambda > 0
\]
shows that $H$ is increasing near $b$, so the minimum cannot occur at $b$ either. Since the minimum is attained at an interior point $c \in (a,b)$, Fermat's Theorem yields $H'(c) = 0$. But $H'(c) = f(c) - \lambda$, so $f(c) - \lambda = 0$, that is, $f(c) = \lambda$, which proves the theorem.
\end{proof}

\section{Evaluating Definite Integrals}

\subsection{The Fundamental Theorem of Calculus, Part II}

Having established the existence of antiderivatives for continuous functions, we now present the second part of the Fundamental Theorem of Calculus, which furnishes an extraordinarily efficient method for evaluating definite integrals, provided that an antiderivative of the integrand can be found explicitly.

\begin{theorem}[Fundamental Theorem of Calculus, Part II]
If $f$ is integrable on $[a,b]$ and has an antiderivative $F$ on $[a,b]$, then
\[
\int_a^b f(x) \, dx = F(b) - F(a).
\]
In particular, this conclusion holds for every continuous function $f$.
\end{theorem}

It is worth pausing to appreciate how remarkable this result is: the definite integral
\[
\int_a^b f(x) \, dx,
\]
defined by means of a rather complicated limiting procedure involving the values of $f$ at every point of $[a,b]$, can nonetheless be computed simply by evaluating a suitable function $F$ at only the two endpoints $a$ and $b$.

\begin{proof}[Proof for continuous $f$]
Define the area function of $f$ by
\[
g(x) = \int_a^x f(t) \, dt
\]

By the first part of the Fundamental Theorem of Calculus, we know that $g' = f$, that is, $g$ is itself an antiderivative of $f$. By the definition of $g$, we have
\[
g(b) - g(a) = \int_a^b f(t) \, dt - \int_a^a f(t) \, dt = \int_a^b f(t) \, dt.
\]

If $F$ is any other antiderivative of $f$ on $[a,b]$, then $g$ and $F$ differ by a constant, so there exists $C \in \mathbb{R}$ such that $F(x) = g(x) + C$ for every $x \in [a,b]$. Consequently,
\[
F(b) - F(a) = [g(b) + C] - [g(a) + C] = g(b) - g(a) = \int_a^b f(t) \, dt,
\]
which completes the proof.
\end{proof}

\subsection{Differentiation and Integration as Inverse Processes}

The two parts of the Fundamental Theorem of Calculus, taken together, reveal that differentiation and integration are, in a precise sense, inverse operations of one another. If $f$ is continuous on $[a,b]$, the first part asserts that
\[
\frac{d}{dx} \int_a^x f(t) \, dt = f(x),
\]
while the second part asserts that, for any function $F$ that is differentiable with continuous derivative on $[a,b]$,
\[
\int_a^b F'(x) \, dx = F(b) - F(a).
\]

The first identity states that if a function $f$ is integrated and the result is subsequently differentiated, one recovers the original function $f$; the second identity states that if a function $F$ is first differentiated and the result is subsequently integrated, one recovers the original function $F$, in the form $F(b) - F(a)$. Together, these two facts justify the assertion that differentiation and integration are inverse processes, in the sense that each undoes what the other does.

This close relationship motivates the introduction of a companion notion to the definite integral, applicable to an entire interval rather than to a specific pair of endpoints.

\begin{definition}[Indefinite Integral]
Given a function $f : I \to \mathbb{R}$, where $I$ is an interval, we call the indefinite integral of $f$ on $I$ the set of all antiderivatives of $f$ on $I$. This set is denoted by the symbol
\[
\int f(x) \, dx.
\]

\end{definition}

\subsection{Change of Variable in the Definite Integral}

A further practical tool for the evaluation of definite integrals is provided by the change of variable formula, which allows a given integral to be transformed into a possibly simpler one through a suitable substitution.

\begin{theorem}[Change of Variable]
Let $f$ be continuous on $[a,b]$, and let $\varphi : [\alpha, \beta] \to [a,b]$ be an invertible function from $[\alpha,\beta]$ onto $[a,b]$. Assume further that $\varphi$ is differentiable, with $\varphi'$ continuous on $[\alpha,\beta]$. Then
\[
\int_a^b f(x) \, dx = \int_\alpha^\beta f(\varphi(t)) \cdot \varphi'(t) \, dt,
\]
where $\varphi(\alpha) = a$ and $\varphi(\beta) = b$.
\end{theorem}

\begin{proof}
Let $F$ be any antiderivative of $f$ on $[a,b]$. By the Chain Rule, the composition $F \circ \varphi$ is an antiderivative of $(f \circ \varphi) \, \varphi'$ on $[\alpha,\beta]$. Applying the second part of the Fundamental Theorem of Calculus to the right-hand side, we obtain
\[
\begin{aligned}
\int_\alpha^\beta f(\varphi(t)) \, \varphi'(t) \, dt
&= (F \circ \varphi)(\beta) - (F \circ \varphi)(\alpha) \\
&= F(\varphi(\beta)) - F(\varphi(\alpha)) \\
&= F(b) - F(a).
\end{aligned}
\]

Applying the second part of the Fundamental Theorem of Calculus once more, this time to the left-hand side, we likewise obtain
\[
\int_a^b f(x) \, dx = F(b) - F(a).
\]
Since both integrals equal $F(b) - F(a)$, the two integrals coincide, which proves the theorem.
\end{proof}

\section{Antiderivatives of Elementary Functions}

\subsection{Domains, Decomposition, and Chain-Rule Integrals}

The primitives, or antiderivatives, of elementary functions are derived directly from the differentiation rules established for those functions. It is essential, however, to exercise mathematical rigor concerning the domains on which these primitives are valid. A primitive $F$ of a function $f$ on an interval $(a,b)$ must, by definition, be differentiable, and consequently continuous, at every point of that interval; standard integration formulas are therefore valid only on intervals on which the corresponding primitive is indeed well defined and differentiable.

A particularly instructive example of this subtlety is furnished by the function $f(x) = 1/x$. Although it is common to state that the primitive of $1/x$ is $\ln|x|$, this compact notation requires careful interpretation. The function $\ln|x|$ is neither defined nor differentiable at $x=0$, and consequently it cannot serve as a primitive for $1/x$ on any interval containing the origin. The function $1/x$ is, in fact, integrable only on intervals that do not include zero, that is, on intervals of the form $(a,b)$ with either $0 \leq a < b$ or $a < b \leq 0$. On an interval where $x > 0$, the correct primitive is $\ln(x)$, while on an interval where $x < 0$, the correct primitive is $\ln(-x)$; the compact notation $\ln|x|$ serves merely as a convenient shorthand for describing both cases simultaneously, but the underlying domain restrictions must always be respected when applying this formula.


The linearity of the operation of differentiation translates directly into a corresponding linearity for the operation of integration. Recall that the derivative of a linear combination of functions equals the corresponding linear combination of their derivatives, that is,
\[
\frac{d}{dx} \big[ \alpha F(x) + \beta G(x) \big] = \alpha F'(x) + \beta G'(x).
\]

Conversely, if $f(x)$ and $g(x)$ are functions admitting primitives $F(x)$ and $G(x)$ respectively, then the linear combination $\alpha f(x) + \beta g(x)$ admits the primitive $\alpha F(x) + \beta G(x)$. This observation yields the integration formula by decomposition,
\[
\int \big[ \alpha f(x) + \beta g(x) \big] \, dx = \alpha \int f(x) \, dx + \beta \int g(x) \, dx.
\]


The Chain Rule for differentiation provides a powerful tool for identifying primitives of composite functions. Let $F$ be a primitive of $f$ on an interval, and let $g$ be a differentiable function; the derivative of the composite function $F(g(x))$ is then given by
\[
\begin{aligned}
\frac{d}{dx} F(g(x)) &= F'(g(x)) \cdot g'(x) \\
&= f(g(x)) \cdot g'(x).
\end{aligned}
\]

Integrating both sides of this identity with respect to $x$ yields the substitution formula in its differential form,
\[
\int f(g(x)) \, g'(x) \, dx = F(g(x)) + C.
\]

This formula allows for the direct computation of what are commonly called almost immediate integrals, namely integrals whose integrand matches the form $f(g(x)) g'(x)$; whenever the inner function $g(x)$ and its derivative $g'(x)$ can be recognized within the integrand, the corresponding primitive can be written down immediately, without any further algebraic manipulation.

\begin{example}
Consider the integral
\[
\int x^3 (x^4+1)^5 \, dx.
\]

Setting $g(x) = x^4+1$, one has $g'(x) = 4x^3$, and the factor $x^3$ appearing in the integrand is proportional to $g'(x)$; taking $f(t) = t^5$, whose primitive is $F(t) = t^6/6$, and adjusting for the constant of proportionality, one obtains
\[
\int x^3 (x^4+1)^5 \, dx = \frac{1}{4} \int 4x^3 (x^4+1)^5 \, dx = \frac{1}{4} F(x^4+1) + C = \frac{1}{24} (x^4+1)^6 + C.
\]

\end{example}

By choosing specific elementary functions for $f(y)$ in the general formula above, one obtains a small table of frequently occurring almost immediate integrals. If $f(y) = y^\alpha$ with $\alpha \neq -1$, then
\[
F(y) = y^{\alpha+1}/(\alpha+1),
\]
yielding the power rule for almost immediate integrals,
\[
\int [g(x)]^\alpha \, g'(x) \, dx = \frac{[g(x)]^{\alpha+1}}{\alpha+1} + C.
\]
If instead $f(y) = 1/y$, then $F(y) = \ln|y|$, yielding the corresponding logarithmic rule,
\[
\int \frac{g'(x)}{g(x)} \, dx = \ln|g(x)| + C.
\]

\subsection{Integration by Substitution}

The method of integration by substitution generalizes the technique of almost immediate integrals to situations in which the integrand does not directly match the form $f(g(x)) g'(x)$, and is used whenever an integral proves difficult to evaluate in its original variable $x$. The method proceeds by introducing a new variable $t$ through a change of variables $x = g(t)$, in the hope that the resulting integral becomes simpler to evaluate. Formally, if $x = g(t)$, then $dx = g'(t) \, dt$, and the integral transforms according to
\[
\int f(x) \, dx = \int f(g(t)) \, g'(t) \, dt.
\]

Once the integral on the right-hand side has been evaluated with respect to $t$, one substitutes back $t = g^{-1}(x)$ in order to express the final result in terms of the original variable $x$.

\begin{example}
To compute
\[
\int \sqrt{2x+1} \, dx,
\]
we set $u = 2x+1$, so that $du = 2 \, dx$, that is, $dx = (1/2) \, du$. The integral becomes
\[
\int \sqrt{u} \cdot \frac{1}{2} \, du = \frac{1}{2} \int u^{1/2} \, du = \frac{1}{2} \cdot \frac{2}{3} u^{3/2} + C = \frac{1}{3} (2x+1)^{3/2} + C.
\]

\end{example}

\subsection{Integration by Parts}

Integration by parts is derived from the Leibniz rule for the differentiation of a product of two functions. Recall that
\[
\frac{d}{dx} \big[ f(x) g(x) \big] = f'(x) g(x) + f(x) g'(x).
\]
Integrating both sides of this identity with respect to $x$ yields
\[
f(x) g(x) = \int f'(x) g(x) \, dx + \int f(x) g'(x) \, dx,
\]
and rearranging the terms produces the formula for integration by parts,
\[
\int f(x) g'(x) \, dx = f(x) g(x) - \int f'(x) g(x) \, dx,
\]
often written more compactly, using differential notation, as
\[
\int u \, dv = uv - \int v \, du.
\]

\begin{example}
Consider the integral
\[
\int x \sin(x) \, dx.
\]

Setting $f(x) = x$ and $g'(x) = \sin(x)$, one has $f'(x) = 1$, and one may choose $g(x) = -\cos(x)$. Applying the formula for integration by parts,
\[
\int x \sin(x) \, dx = x(-\cos(x)) - \int 1 \cdot (-\cos(x)) \, dx = -x\cos(x) + \int \cos(x) \, dx,
\]
and evaluating the remaining elementary integral yields the final result, $-x\cos(x) + \sin(x) + C$.
\end{example}

This technique proves particularly useful for integrals involving products of polynomials with trigonometric functions, such as $\sin(mx)$ or $\cos(mx)$, or with exponential functions of the form $e^{mx}$; in some cases, the formula for integration by parts must be applied repeatedly, in succession, before the integral is finally resolved.

\subsection{Rational Functions and Partial Fractions}

Rational functions are quotients of polynomials, and the strategy required to integrate them depends critically on the relative degrees of the numerator and denominator, as well as on the algebraic structure of the denominator.


For integrals of the form
\[
\int (ax+b)/(cx+d) \, dx,
\]
one must first perform polynomial division whenever the degree of the numerator is greater than or equal to the degree of the denominator. In the present case, since both numerator and denominator have degree one, the fraction may be rewritten directly as
\[
\frac{ax+b}{cx+d} = \frac{a}{c} + \frac{b - \tfrac{ad}{c}}{cx+d}.
\]

Applying the linearity of the integral, together with the substitution $u = cx+d$, for which $du = c \, dx$, one obtains the explicit antiderivative
\[
\int \frac{ax+b}{cx+d} \, dx = \frac{a}{c} x + \frac{bc-ad}{c^2} \ln|cx+d| + C.
\]

For integrals of the form
\[
\int P(x)/(ax^2+bx+c) \, dx,
\]
where $P(x)$ is a polynomial of degree less than two, the appropriate technique depends on the sign of the discriminant $\Delta = b^2 - 4ac$ of the denominator; we now examine the two resulting cases in turn.

If the discriminant is positive, $\Delta > 0$, the denominator factors into two distinct real roots, and the appropriate technique is partial fraction decomposition. As an illustration, consider the integral
\[
\int 1/(x^2-x-2) \, dx.
\]
The roots of $x^2-x-2=0$ are $x=2$ and $x=-1$, so that the integrand decomposes as
\[
\frac{1}{(x-2)(x+1)} = \frac{A}{x-2} + \frac{B}{x+1}
\]
for suitable constants $A$ and $B$; clearing denominators gives
\[
1 = A(x+1) + B(x-2),
\]
and setting $x=2$ yields $A = 1/3$, while setting $x=-1$ yields $B = -1/3$. Integrating the resulting decomposition term by term, one finds
\[
\int \left( \frac{1/3}{x-2} - \frac{1/3}{x+1} \right) dx = \frac{1}{3} \ln|x-2| - \frac{1}{3} \ln|x+1| + C.
\]

If instead the discriminant is negative, $\Delta < 0$, the denominator cannot be factored into real linear terms, and the appropriate technique is to complete the square within the denominator. For a general integral of the form
\[
\int 1/(ax^2+bx+c) \, dx,
\]
one first manipulates the numerator, if it is of first degree, by splitting the integral into a part resembling $f'(x)/f(x)$ together with a constant remainder; for a purely constant numerator, one relies on the fundamental formula
\[
\int \frac{1}{x^2+a^2} \, dx = \frac{1}{a} \arctan\!\left( \frac{x}{a} \right) + C.
\]
As an illustration, consider the integral
\[
\int 1/(4x^2+2x+1) \, dx.
\]
Completing the square in the denominator, one finds
\[
\begin{aligned}
4x^2+2x+1 &= 4\left(x^2 + \frac{1}{2}x\right)+1 \\
&= 4\left[\left(x+\frac{1}{4}\right)^2 - \frac{1}{16}\right]+1 \\
&= 4\left(x+\frac{1}{4}\right)^2 + \frac{3}{4}.
\end{aligned}
\]
Setting $u = x+1/4$, so that $du = dx$, the integral becomes
\[
\int \frac{1}{4u^2 + 3/4} \, du = \int \frac{1}{\tfrac{3}{4}\left( \tfrac{16}{3} u^2 + 1 \right)} \, du,
\]
and the further substitution $v=4u/\sqrt{3}$ gives
\[
\begin{aligned}
\int \frac{\dd u}{4u^2+3/4}
&=\frac{1}{\sqrt{3}}\int\frac{\dd v}{1+v^2}\\
&=\frac{1}{\sqrt{3}}\arctan v+C.
\end{aligned}
\]
Substituting back first $v=4u/\sqrt{3}$ and then $u=x+1/4$
produces the explicit result
\[
\boxed{\int\frac{\dd x}{4x^2+2x+1}
=\frac{1}{\sqrt{3}}
\arctan\!\left(\frac{4x+1}{\sqrt{3}}\right)+C.}
\]


The two cases examined above extend to a general strategy applicable to an arbitrary rational function. If the degree of the numerator is greater than or equal to the degree of the denominator, one performs polynomial long division first, obtaining a polynomial part, which is integrated immediately term by term, together with a proper rational function, that is, one in which the degree of the numerator is strictly smaller than the degree of the denominator. The resulting proper rational function is then decomposed into partial fractions according to the factorization of the denominator into linear and irreducible quadratic factors, according to the general decomposition
\[
\frac{P(x)}{Q(x)} = \sum_i \frac{A_i}{(x-r_i)^{k_i}} + \sum_j \frac{B_j x + C_j}{(x^2+p_j x+q_j)^{m_j}}.
\]

Each term of this decomposition is then integrated separately, using the logarithmic rule established above for the linear factors, and the arctangent and substitution techniques established above for the irreducible quadratic factors.

\section{Improper Integrals}

\subsection{Improper Integrals of Type I}

In defining the definite integral
\[
\int_a^b f(x) \, dx,
\]
we have dealt exclusively with a bounded function $f$ on a bounded interval $[a,b]$. We now extend this concept to the case in which the interval of integration is infinite, of the form $[a, \infty)$; an integral of this kind is called an improper integral, and is denoted by
\[
\int_a^{\infty} f(x) \, dx.
\]

To motivate the appropriate definition, consider the region $S$ that lies under the graph of $y = f(x)$, with $f \geq 0$, above the horizontal axis, and to the right of $x = a$, that is,
\[
S = \{ (x,y) \in \mathbb{R}^2 : x \geq a, \ 0 \leq y \leq f(x) \}.
\]

One might initially suspect that, since $S$ extends infinitely far to the right, its area must necessarily be infinite; a concrete example shows that this intuition is, in fact, mistaken. Consider the infinite region lying under the curve $y = 1/x^2$, above the horizontal axis, and to the right of the line $x = 1$. The area of the portion of this region lying to the left of the line $x = b$ is given by
\[
A(b) = \int_1^b \frac{1}{x^2} \, dx = \left[ -\frac{1}{x} \right]_{x=1}^{x=b} = 1 - \frac{1}{b}.
\]
The limiting area is determined by
\[
\lim_{b \to \infty} A(b) = 1
\]

Thus, the area of the shaded region approaches the finite value $1$ as $b \to \infty$, and it is therefore natural to say that the area of the entire infinite region $S$ equals $1$, despite the unbounded extent of the region itself.

Guided by this example, we now formulate the general definition of an improper integral for a bounded function $f$, not necessarily positive, defined on $[a, \infty)$.

\begin{definition}[Improper Integral of Type I]
Suppose that the following integral exists for every $t > a$:
\[
\int_a^t f(x) \, dx.
\]
If the limit on the right exists, we define the improper integral by
\[
\int_a^{\infty} f(x) \, dx := \lim_{t \to \infty} \int_a^t f(x) \, dx.
\]

If the limit is finite, we say that the improper integral is convergent; if the limit is infinite, we say that it is divergent.
\end{definition}

The improper integral introduced in this definition retains a geometric interpretation as an area whenever $f$ is a positive function. Thus, if $f \geq 0$ and the integral
\[
\int_a^{\infty} f(x) \, dx
\]
is convergent, one defines the area of the region
\[
S = \{ (x,y) \in \mathbb{R}^2 : x \geq a, \ 0 \leq y \leq f(x) \}
\]
to be
\[
A(S) = \int_a^{\infty} f(x) \, dx.
\]
It is instructive to contrast the convergent example above with the closely related integral
\[
\int_1^{\infty} (1/x) \, dx,
\]
for which
\[
\int_1^R\frac{\dd x}{x}=[\ln x]_1^R=\ln R\longrightarrow+\infty.
\]
Thus it diverges, and the corresponding area is infinite.

\begin{remark}
Both $1/x^2$ and $1/x$ tend to $0$ as $x \to \infty$, yet the corresponding improper integrals exhibit opposite behavior: the former converges, while the latter diverges. This shows that the values of $1/x$ simply do not decrease fast enough, as $x \to \infty$, for the corresponding improper integral to attain a finite value.
\end{remark}

More generally, for $p\neq1$,
\[
\int_1^{\infty} x^{-p} \, dx
\]
satisfies
\[
\begin{aligned}
\int_1^R x^{-p}\,\dd x
&=\left[\frac{x^{1-p}}{1-p}\right]_1^R\\
&=\frac{R^{1-p}-1}{1-p}.
\end{aligned}
\]
If $p>1$, then $R^{1-p}\to0$ and the integral converges to
$1/(p-1)$. If $p<1$, it diverges; the case $p=1$ is the logarithmic
case above. Finally,
\[
\int_0^{\infty} e^{-x} \, dx
\]
satisfies
\[
\int_0^R e^{-x}\,\dd x=[-e^{-x}]_0^R=1-e^{-R}\longrightarrow1,
\]
so it converges to $1$.

Figure~\ref{fig:atlas-improper-tail} visualizes the remainder left when a convergent improper integral is truncated.

\begin{figure}[!htbp]
\centering
\begin{tikzpicture}[>=Stealth]
\begin{axis}[width=11cm,height=5.8cm,axis lines=middle,ticklabel style={font=\small},label style={font=\small},legend style={font=\scriptsize,draw=black!20,fill=white},xlabel={$x$},ylabel={$1/x^2$},xmin=0,xmax=8.5,ymin=0,ymax=1.15,xtick={1,3,5,7},ytick={0.5,1}]
\addplot[draw=none,fill=figred!25,domain=3:8,samples=100] {1/x^2} \closedcycle;
\addplot[figblue,very thick,domain=1:8,samples=130] {1/x^2};
\draw[dashed,figred] (axis cs:3,0) -- (axis cs:3,0.32);
\node[above,figred] at (axis cs:3,0.32) {$R=3$};
\draw[->,figred] (axis cs:6,0.09) -- (axis cs:8,0.09);
\end{axis}
\end{tikzpicture}
\caption{For the integrable function $1/x^2$ on $[1,\infty)$, truncation at $R=3$ leaves a tail with integral $1/3$. Only the portion of that tail up to $x=8$ is shaded; the arrow indicates its continuation to infinity.}
\label{fig:atlas-improper-tail}
\end{figure}

\subsection{Convergence Tests and the Integral Test}

Just as for series, it is frequently more convenient to establish the convergence or divergence of an improper integral by comparison with another, better-understood integral, rather than by direct computation. The following two comparison results serve precisely this purpose.

\begin{theorem}[Comparison Test]
Let $f, g \geq 0$ be two bounded functions defined on $[a, \infty)$, with $f(x) \geq g(x)$ for every $x$ sufficiently large.
\[
\begin{gathered}
\text{If }\int_a^{\infty} f(x) \, dx
\text{ is convergent, then }\int_a^{\infty} g(x) \, dx
\text{ is convergent.} \\
\text{If }\int_a^{\infty} g(x) \, dx
\text{ is divergent, then }\int_a^{\infty} f(x) \, dx
\text{ is divergent.}
\end{gathered}
\]

\end{theorem}

As an application of this theorem, one may show that the integral
\[
\int_0^{\infty} e^{-x^2} \, dx
\]
is convergent, by comparing the integrand with a suitable exponential function for large values of $x$; it can in fact be shown, using tools from multivariable calculus, that the exact value of this integral is
\[
\sqrt{\pi}/2,
\]
although this precise value lies beyond the scope of the comparison method itself. A refinement of the Comparison Test, analogous to the Asymptotic Comparison Test established earlier for series, replaces the pointwise inequality between the two functions with an asymptotic equivalence as $x \to \infty$.

\begin{theorem}[Asymptotic Comparison Test]
Let $f, g \geq 0$ be two bounded functions defined on $[a, \infty)$. If $f \sim g$ as $x \to \infty$, then the integral
\[
\begin{gathered}
\int_a^{\infty} f(x) \, dx
\text{ is convergent if and only if }
\int_a^{\infty} g(x) \, dx
\text{ is convergent, and} \\
\int_a^{\infty} f(x) \, dx
\text{ is divergent if and only if }
\int_a^{\infty} g(x) \, dx
\text{ is divergent.}
\end{gathered}
\]

\end{theorem}

As an application, one may show that the integral
\[
\int_2^{\infty} \ln\!\left(1 + 1/x^2\right) dx
\]
is convergent, by observing that the integrand is asymptotically equivalent to $1/x^2$ as $x \to \infty$, and invoking the convergence of
\[
\int_1^{\infty} x^{-2} \, dx
\]
established earlier.

The theory of improper integrals is intimately connected with the theory of series, since both arise from a common geometric idea, namely the approximation of an area by a sum of rectangles. This connection can be exploited in two directions.

On one hand, one may use an improper integral to bound the sum of a series from above. Consider rectangles of unit width erected above the intervals $[n, n+1]$, with height equal to $1/n^2$; the sum of the areas of these rectangles is precisely
\[
\sum_{n=1}^{\infty} 1/n^2.
\]

If the first rectangle is excluded, the total area of the remaining rectangles is smaller than the area under the curve $y = 1/x^2$ for $x \geq 1$, which is precisely the value of the integral
\[
\int_1^{\infty} 1/x^2 \, dx.
\]
This comparison yields the explicit numerical bound
\[
\sum_{n=1}^{\infty} \frac{1}{n^2} \leq 1 + \int_1^{\infty} \frac{1}{x^2} \, dx = 2.
\]

On the other hand, an analogous comparison can be used to establish the divergence of a series. In order to estimate
\[
\sum_{n=1}^{\infty} 1/\sqrt{n},
\]
one instead uses rectangles whose tops lie above the curve
\[
y = 1/\sqrt{x},
\]
choosing the height of each rectangle equal to the value of the function at the left endpoint of the corresponding interval; this comparison yields
\[
\sum_{n=1}^{\infty} \frac{1}{\sqrt{n}} \geq \int_1^{\infty} \frac{1}{\sqrt{x}} \, dx = \infty,
\]
so that the series diverges. These two examples are particular instances of a general principle, known as the Integral Test, which relates the convergence of a series directly to the convergence of an associated improper integral.

\begin{theorem}[Integral Test for Series]
Let $f$ be a continuous, positive, decreasing function on $[b, \infty)$, and set $a_n = f(n)$. Then the series
\[
\sum_n a_n
\]
is convergent if and only if the improper integral
\[
\int_b^{\infty} f(x) \, dx
\]
is convergent.
\end{theorem}

\begin{remark}
It is not necessary that $f$ be decreasing throughout its entire domain; what matters is that $f$ be ultimately decreasing, that is, decreasing for $x$ larger than some fixed number $N$. It should also be emphasized that one should not infer from the Integral Test that the sum of the series is equal to the numerical value of the corresponding integral; the theorem provides information only about the character of the series, and not about the precise value of its sum.
\end{remark}

As applications of the Integral Test, one may study the convergence or divergence of the series
\[
\begin{gathered}
\sum_{n=2}^{\infty} 1/[n (\ln n)^q] \\
\sum_{n=1}^{\infty} 1/n^p,
\end{gathered}
\]
for positive parameters $p, q > 0$, by comparing each series with the corresponding improper integral and applying the criteria for convergence of
\[
\int^{\infty} x^{-p} \, dx
\]
established earlier.

The definition of improper integral given above extends naturally to a left semi-infinite interval, and to the entire real line. If
\[
\int_t^b f(x) \, dx
\]
exists for every number $t \leq b$, then
\[
\int_{-\infty}^{b} f(x) \, dx := \lim_{t \to -\infty} \int_t^b f(x) \, dx,
\]
provided that this limit exists. If both
\[
\begin{gathered}
\int_{-\infty}^{c} f(x) \, dx \\
\int_c^{\infty} f(x) \, dx
\end{gathered}
\]
are convergent for some real number $c$, one then defines
\[
\int_{-\infty}^{\infty} f(x) \, dx := \int_{-\infty}^{c} f(x) \, dx + \int_c^{\infty} f(x) \, dx,
\]
and it can be shown that the resulting value does not depend on the particular choice of the splitting point $c$.

\subsection{Improper Integrals of Type II}

A second, geometrically distinct type of improper integral arises not from an unbounded interval of integration, but from an unbounded integrand. Suppose that $f$ is a positive, continuous function defined on a finite interval $[a,b)$, but possesses a vertical asymptote at $b$; let $S$ denote the unbounded region lying under the graph of $f$ and above the horizontal axis, between $a$ and $b$. The area of the portion of $S$ lying between $a$ and $t$, for $t < b$, is given by
\[
A(t) = \int_a^t f(x) \, dx.
\]

If it happens that $A(t)$ approaches a definite finite number $A$ as $t \to b^-$, one says that the area of the region $S$ equals $A$, and one writes
\[
\int_a^b f(x) \, dx := \lim_{t \to b^-} \int_a^t f(x) \, dx.
\]

This geometric idea can be used to define an improper integral of Type II even when $f$ is not necessarily a positive function.

\begin{definition}[Improper Integral of Type II]
If $f$ is continuous on $[a,b)$ and has a vertical asymptote at $x = b$, then
\[
\int_a^b f(x) \, dx := \lim_{t \to b^-} \int_a^t f(x) \, dx,
\]
provided that this limit exists as a finite number. The improper integral is called convergent if the corresponding limit exists and is finite, and divergent if the limit is infinite.
\end{definition}

An entirely analogous definition applies when the vertical asymptote occurs at the left endpoint rather than at the right: if $f$ is continuous on $(a,b]$ and has a vertical asymptote at $x = a$, one sets
\[
\int_a^b f(x) \, dx := \lim_{t \to a^+} \int_t^b f(x) \, dx,
\]
whenever this limit exists as a finite number. Finally, if $f$ is continuous on $[a,b] \setminus \{c\}$ and has a vertical asymptote at an interior point $c$, one defines the improper integral over the entire interval $[a,b]$ by splitting at the singularity,
\[
\int_a^b f(x) \, dx := \int_a^c f(x) \, dx + \int_c^b f(x) \, dx,
\]
provided that both improper integrals on the right-hand side are convergent.

\begin{remark}
From this point onward, whenever the symbol
\[
\int_a^b f(x) \, dx
\]
is encountered, one must first examine the behavior of $f$ on $[a,b]$ in order to determine whether the expression represents an ordinary definite integral or an improper integral, since the two require different definitions and different tests for existence.
\end{remark}

We now evaluate the standard endpoint examples directly. First,
\[
\begin{aligned}
\int_0^1\frac{\dd x}{x}
&=\lim_{t\to0^+}\int_t^1\frac{\dd x}{x}\\
&=\lim_{t\to0^+}[\ln x]_t^1\\
&=\lim_{t\to0^+}(-\ln t)=+\infty,
\end{aligned}
\]
so the integral diverges. By contrast,
\[
\begin{aligned}
\int_0^1\ln x\,\dd x
&=\lim_{t\to0^+}[x\ln x-x]_t^1\\
&=-1-\lim_{t\to0^+}(t\ln t-t)\\
&=-1,
\end{aligned}
\]
because $t\ln t\to0$ as $t\to0^+$. More generally, for $p\neq1$,
\[
\int_0^1 x^{-p} \, dx
\]
satisfies
\[
\begin{aligned}
\int_0^1x^{-p}\,\dd x
&=\lim_{t\to0^+}\left[\frac{x^{1-p}}{1-p}\right]_t^1\\
&=\lim_{t\to0^+}\frac{1-t^{1-p}}{1-p}.
\end{aligned}
\]
If $p<1$, then $t^{1-p}\to0$ and the value is $1/(1-p)$; if
$p>1$, then $t^{1-p}\to+\infty$ and the integral diverges. The case
$p=1$ is the logarithmic divergence computed above. Thus the integral
converges exactly for $p<1$ and diverges for $p\geq1$. This behavior is
opposite to that of
\[
\int_1^{\infty} x^{-p} \, dx
\]
which converges exactly for $p>1$.

\subsection{Improper Integrals and Probability}


A particularly important application of improper integrals arises in the theory of probability, where continuous random variables are described by means of density functions defined on the entire real line. Random variables of this kind occur throughout the sciences: examples include the time one waits on hold before an agent of a company answers a telephone call, the test scores obtained by an individual on an aptitude test, the heights and weights of individuals drawn from a homogeneous population, and the lifetime of a randomly chosen battery of a given type. In the case of a battery, one might ask, for instance, what the probability is that the battery lasts between $100$ and $200$ hours; if $X$ denotes the lifetime of the battery, this probability is denoted by $P(100 \leq X \leq 200)$.

Every continuous random variable $X$ possesses a probability density function $f$, meaning that the probability that $X$ lies between two values $a$ and $b$ is found by integrating $f$ from $a$ to $b$,
\[
P(a \leq X \leq b) = \int_a^b f(x) \, dx.
\]

Since probabilities are measured on a scale from $0$ to $1$, a density function must satisfy $f \geq 0$ together with the normalization condition
\[
\int_{-\infty}^{\infty} f(x) \, dx = 1,
\]
expressed precisely as an improper integral over the whole real line. If
the first absolute moment is finite,
\[
\int_{-\infty}^{\infty}|x|f(x)\,\dd x<\infty,
\]
the distribution has a finite mean, defined by
\[
\mu = \int_{-\infty}^{\infty} x f(x) \,\dd x.
\]
Not every probability density satisfies this condition; the Cauchy
distribution is a standard example for which no finite mean exists.


Many important random phenomena encountered in practice, such as test scores on aptitude tests, the heights and weights of individuals from a homogeneous population, or the annual rainfall recorded at a given location, are well modeled by a normal distribution. This means that the probability density function of the corresponding random variable $X$ belongs to the family of functions
\[
f(x) = \frac{1}{\sigma \sqrt{2\pi}} \, e^{-\frac{(x-\mu)^2}{2\sigma^2}}.
\]

The mean of this distribution equals $\mu$, a fact that can be verified directly from the definition of the mean given above; the positive constant $\sigma$ is called the standard deviation, and it measures how spread out the values of $X$ tend to be around the mean $\mu$. The fact that
\[
\int_{-\infty}^{\infty} f(x) \, dx = 1
\]
for this particular density function can be verified rigorously by methods belonging to multivariable calculus, and lies beyond the scope of the present discussion.

As an application of the normal distribution, one may consider intelligence quotient scores, which are known to be distributed normally with mean $100$ and standard deviation $15$; one may then ask what percentage of the population has an intelligence quotient score between $85$ and $115$, and what percentage has a score above $140$, both questions being answered by evaluating suitable definite or improper integrals of the corresponding density function. As a further application, one may consider a company whose average customer waiting time before a call is answered by a representative is five minutes, and ask for the probability that a call is answered during the first minute, as well as the probability that a customer waits more than five minutes to be answered, both of which again reduce to the evaluation of integrals of an appropriate probability density function.

\newpage
\chapter{Applied Linear Algebra}

\section{Matrices, Products, and Linear Systems}

\subsection{Matrices and Basic Operations}


A matrix is a rectangular array of $m \cdot n$ numbers organized into $m$ horizontal rows and $n$ vertical columns,
\[
A =
\begin{pmatrix}
a_{11} & a_{12} & \cdots & a_{1n} \\
a_{21} & a_{22} & \cdots & a_{2n} \\
\vdots & \vdots & \ddots & \vdots \\
a_{m1} & a_{m2} & \cdots & a_{mn}
\end{pmatrix}.
\]

Each element, or entry, of the matrix is denoted by a variable carrying two subscripts: the symbol $a_{ij}$ represents the entry located at the $i$-th row and $j$-th column, where the row index $i$ ranges from $1$ to $m$ and the column index $j$ ranges from $1$ to $n$. Each row of the matrix may itself be regarded as a vector with $n$ components,
\[
\mathbf{a}_i = [a_{i1} \ a_{i2} \ \cdots \ a_{in}]
\]
for $i = 1, \dots, m$, and each column may likewise be regarded as a vector with $m$ components,
\[
\mathbf{a}^j = (a_{1j}, a_{2j}, \dots, a_{mj})
\]
for $j = 1, \dots, n$. A matrix $A$ as above is denoted compactly by $A = [a_{ij}]$. The set of all matrices with real entries of dimension $m \times n$ is denoted by $M_{m,n}$, or simply by $M_n$ in the special case $m = n$, in which case the matrices are called square matrices. The matrix all of whose entries equal zero is called the zero matrix, and is denoted by $O \in M_{m,n}$. Finally, given $n$ vectors $\mathbf{b}^1, \dots, \mathbf{b}^n \in \mathbb{R}^m$, the notation
\[
[\mathbf{b}^1 | \cdots | \mathbf{b}^n] \in M_{m,n}
\]
denotes the matrix obtained by placing these vectors into the columns of the matrix, in the given order.


Certain matrix configurations recur frequently enough throughout the theory to deserve dedicated terminology. A matrix is said to be in row echelon form when the leading coefficient of every nonzero row, that is, the first nonzero entry counted from the left, also called the pivot of that row, lies strictly to the right of the leading coefficient of the row directly above it, and all rows consisting entirely of zeros, if any, are collected at the bottom of the matrix. These two conditions together imply that every entry lying in a column below a leading coefficient must itself be zero, so that the lower left-hand corner of a matrix in row echelon form is filled with zeros.

Among matrices with an equal number of rows and columns, several further special types are worth distinguishing. A square matrix is a matrix belonging to $M_n$ for some $n$, that is, one having the same number of rows and columns. A diagonal matrix is a square matrix all of whose off-diagonal entries vanish; the particular diagonal matrix having every diagonal entry equal to $1$ is called the identity matrix, and is denoted by $I$. An upper triangular matrix is a square matrix all of whose entries below the main diagonal vanish, while a lower triangular matrix is a square matrix all of whose entries above the main diagonal vanish.


The set of matrices of a fixed dimension can be equipped with two natural algebraic operations, entirely analogous to the corresponding operations on vectors of $\mathbb{R}^n$.

\begin{definition}[Matrix Addition]
Given two matrices of the same dimension, $A = [a_{ij}]$ and $B = [b_{ij}]$, both belonging to $M_{m,n}$, their sum $A+B$ is the $m \times n$ matrix whose $(i,j)$ entry is the sum of the corresponding entries of $A$ and $B$, that is, $A + B = [a_{ij} + b_{ij}]$, for every $i = 1, \dots, m$ and $j = 1, \dots, n$.
\end{definition}

\begin{definition}[Scalar Multiplication]
Given a matrix $A = [a_{ij}] \in M_{m,n}$ and a real number $t \in \mathbb{R}$, the scalar product $tA$ is the $m \times n$ matrix obtained by multiplying every entry of $A$ by $t$, that is, $tA = [t \, a_{ij}]$.
\end{definition}

These two operations satisfy a collection of algebraic properties entirely analogous to those already familiar from ordinary arithmetic on real numbers. Given any three matrices $A, B, C \in M_{m,n}$ and any two real numbers $s, t$, one has that addition is associative,
\[
A + (B+C) = (A+B) + C,
\]
and commutative, $A + B = B + A$; that the zero matrix $O \in M_{m,n}$ satisfies $A + O = A$; that the matrix $-A = [-a_{ij}]$ is the additive inverse of $A$, satisfying $A + (-A) = O$; and that scalar multiplication satisfies $s(tA) = (st)A$, $1 \cdot A = A$, $t(A+B) = tA + tB$, and $(s+t)A = sA + tA$. Taken together, these properties show that $M_{m,n}$, equipped with the two operations of matrix addition and scalar multiplication, forms a vector space of dimension $m \cdot n$.

\subsection{Matrix Products and Linear Systems}

Before defining the product of two general matrices, it is instructive to consider the simpler case in which one of the two factors is a vector. Given a row vector $\mathbf{a} = [a_1 \ a_2 \ \cdots \ a_n]$ and a column vector
\[
\mathbf{x} = (x_1, x_2, \dots, x_n)
\]
having the same number $n$ of components, the product $\mathbf{a}\mathbf{x}$ is defined as the ordinary dot product in $\mathbb{R}^n$,
\[
\mathbf{a}\mathbf{x} = [a_1 \ a_2 \ \cdots \ a_n]
\begin{pmatrix} x_1 \\ x_2 \\ \vdots \\ x_n \end{pmatrix}
= a_1 x_1 + a_2 x_2 + \cdots + a_n x_n.
\]
This notion extends naturally to the product of a general matrix with a column vector.

\begin{definition}[Matrix-Vector Product]
Given a matrix $A = [a_{ij}] \in M_{m,n}$ and a column vector
\[
\mathbf{x} = (x_1, x_2, \dots, x_n) \in \mathbb{R}^n,
\]
the product $A\mathbf{x}$ is the vector with $m$ components given by
\[
A\mathbf{x} =
\begin{pmatrix} \mathbf{a}_1 \mathbf{x} \\ \mathbf{a}_2 \mathbf{x} \\ \vdots \\ \mathbf{a}_m \mathbf{x} \end{pmatrix}
=
\begin{pmatrix}
a_{11} x_1 + a_{12} x_2 + \cdots + a_{1n} x_n \\
a_{21} x_1 + a_{22} x_2 + \cdots + a_{2n} x_n \\
\vdots \\
a_{m1} x_1 + a_{m2} x_2 + \cdots + a_{mn} x_n
\end{pmatrix} \in \mathbb{R}^m,
\]
Here the $i$-th row of $A$ is
\[
\mathbf{a}_i = [a_{i1} \ a_{i2} \ \cdots \ a_{in}]
\]

\end{definition}

\begin{example}
Consider the following matrix and vector:
\[
\begin{gathered}
A = \begin{pmatrix} 0 & 1 & 1 \\ 2 & 2 & 1 \end{pmatrix} \in M_{2,3} \\
\mathbf{x} = (2,1,3) \in \mathbb{R}^3
\end{gathered}
\]
Their product is
\[
A\mathbf{x} = \begin{pmatrix} 0\cdot 2 + 1 \cdot 1 + 1 \cdot 3 \\ 2 \cdot 2 + 2 \cdot 1 + 1 \cdot 3 \end{pmatrix} = \begin{pmatrix} 4 \\ 9 \end{pmatrix} \in \mathbb{R}^2.
\]

\end{example}

The matrix-vector product enjoys an important algebraic property, namely linearity with respect to the vector argument, which lies at the very heart of the entire theory of linear algebra.

\begin{theorem}[Linearity of the Matrix-Vector Product]
Given $A \in M_{m,n}$, the matrix-vector product $A\mathbf{x}$ is linear in $\mathbf{x}$; that is, $A(\mathbf{x}+\mathbf{y}) = A\mathbf{x} + A\mathbf{y}$ for every $\mathbf{x}, \mathbf{y} \in \mathbb{R}^n$, and $A(t\mathbf{x}) = t(A\mathbf{x})$ for every $t \in \mathbb{R}$ and $\mathbf{x} \in \mathbb{R}^n$.
\end{theorem}

Given a matrix $A \in M_{m,n}$, one may associate with it the function $L : \mathbb{R}^n \to \mathbb{R}^m$ defined by $L(\mathbf{x}) = A\mathbf{x}$. The linearity property just stated shows precisely that $L$ is a linear map, in the sense that
\[
\begin{gathered}
L(\mathbf{x} + \mathbf{y}) = L(\mathbf{x}) + L(\mathbf{y}) \\
L(t\mathbf{x}) = t L(\mathbf{x})
\end{gathered}
\]
for every $\mathbf{x}, \mathbf{y} \in \mathbb{R}^n$ and $t \in \mathbb{R}$; this correspondence between matrices and linear maps, easily verified by direct computation, will be developed systematically in a later section.

A particularly useful way of viewing the matrix-vector product emerges from rewriting the explicit expression for $A\mathbf{x}$ given above. Since
\[
A\mathbf{x} =
x_1 \begin{pmatrix} a_{11} \\ a_{21} \\ \vdots \\ a_{m1} \end{pmatrix}
+ x_2 \begin{pmatrix} a_{12} \\ a_{22} \\ \vdots \\ a_{m2} \end{pmatrix}
+ \cdots +
x_n \begin{pmatrix} a_{1n} \\ a_{2n} \\ \vdots \\ a_{mn} \end{pmatrix},
\]
one recognizes that the product $A\mathbf{x}$ is precisely the linear combination of the columns $\mathbf{a}^j$ of $A$, with coefficients given by the components of $\mathbf{x}$, that is,
\[
A\mathbf{x} = x_1 \mathbf{a}^1 + x_2 \mathbf{a}^2 + \cdots + x_n \mathbf{a}^n.
\]

This alternative viewpoint proves indispensable when relating the matrix-vector product to the theory of linear systems, developed in the following subsection, and to the theory of linear independence, developed subsequently.

The general form of a linear system of $m$ equations in $n$ unknowns $x_1, x_2, \dots, x_n$ is
\[
\begin{cases}
a_{11} x_1 + a_{12} x_2 + \cdots + a_{1n} x_n = b_1, \\
a_{21} x_1 + a_{22} x_2 + \cdots + a_{2n} x_n = b_2, \\
\quad \vdots \\
a_{m1} x_1 + a_{m2} x_2 + \cdots + a_{mn} x_n = b_m,
\end{cases}
\]
where the coefficients $a_{ij}$ and the constants $b_i$ are assigned numbers. In light of the definition of the matrix-vector product, this system can be written compactly as the single vector equation $A\mathbf{x} = \mathbf{b}$ in $\mathbb{R}^m$, having set
\[
A = \begin{pmatrix} a_{11} & \cdots & a_{1n} \\ \vdots & \ddots & \vdots \\ a_{m1} & \cdots & a_{mn} \end{pmatrix}, \qquad
\mathbf{b} = \begin{pmatrix} b_1 \\ \vdots \\ b_m \end{pmatrix}, \qquad
\mathbf{x} = \begin{pmatrix} x_1 \\ \vdots \\ x_n \end{pmatrix}.
\]

The matrix $A$ is called the coefficient matrix of the system, and the matrix obtained by appending the column $\mathbf{b}$ to $A$, denoted $[A | \mathbf{b}]$, is called the augmented matrix of the system.

A solution of the linear system is an assignment of values to the variables $x_1, \dots, x_n$ satisfying every equation simultaneously; the set of all such solutions is called the solution set,
\[
S = \{ \mathbf{x} = (x_1, \dots, x_n) \in \mathbb{R}^n : A\mathbf{x} = \mathbf{b} \}.
\]

A linear system may exhibit one of exactly three possible behaviors: it may have infinitely many solutions, it may have a single unique solution, or it may have no solution at all. A system possessing no solution is said to be inconsistent, while a system possessing at least one solution is said to be consistent, or solvable.

\begin{example}[Application: Linear Independence]
Given $n$ vectors $\mathbf{v}^1, \dots, \mathbf{v}^n \in \mathbb{R}^m$, one may efficiently compute their linear combinations by first assembling the matrix
\[
A = [\mathbf{v}^1 | \cdots | \mathbf{v}^n] \in M_{m,n}.
\]
The coefficients can be collected into a vector
\[
\mathbf{x} = (x_1, \dots, x_n) \in \mathbb{R}^n.
\]

For every such coefficient vector, the linear combination of the vectors $\mathbf{v}^i$ with coefficients $x_i$ is given precisely by $A\mathbf{x}$. Consequently, the vectors $\mathbf{v}^1, \dots, \mathbf{v}^n$ are linearly independent if and only if the homogeneous system $A\mathbf{x} = \vec{0}_m$ has the unique solution $\mathbf{x} = \vec{0}_n$, reducing the study of linear independence to the study of a homogeneous linear system.
\end{example}

The product of two matrices is defined only when the number of columns of the first matrix equals the number of rows of the second, and the resulting product generalizes the matrix-vector product introduced above.

\begin{definition}[Matrix Multiplication]
Given $A = [a_{ij}] \in M_{m,n}$ and $B = [b_{ij}] \in M_{n,r}$, the product of $A$ times $B$, denoted $AB$, is the $m \times r$ matrix obtained by taking, for each entry, the dot product of the corresponding row of $A$ with the corresponding column of $B$; that is, $AB \in M_{m,r}$, with $[AB]_{ij} = \mathbf{a}_i \cdot \mathbf{b}^j$, so that the $(i,j)$ entry of $AB$ is the scalar product of the $i$-th row of $A$ with the $j$-th column of $B$.
\end{definition}

\begin{example}
Consider the following pair of matrices:
\[
\begin{gathered}
A = \begin{pmatrix} 0 & 1 & 1 \\ 2 & 2 & 1 \end{pmatrix} \in M_{2,3} \\
B = \begin{pmatrix} 1 & 2 \\ 2 & 0 \\ 3 & 1 \end{pmatrix} \in M_{3,2}.
\end{gathered}
\]
The product $AB$ is the $2 \times 2$ matrix
\[
AB = \begin{pmatrix} 4 & 3 \\ 9 & 5 \end{pmatrix},
\]
where the entries are computed as
\[
\begin{aligned}
c_{11} &= 0\cdot 2 + 1\cdot 1 + 1 \cdot 3 \\
&= 4,
\end{aligned}
\]

\[
\begin{aligned}
c_{12} &= 0 \cdot 0 + 1 \cdot 2 + 1 \cdot 1 \\
&= 3,
\end{aligned}
\]

\[
\begin{gathered}
\begin{aligned}
c_{21} &= 2 \cdot 2 + 2 \cdot 1 + 1 \cdot 3 \\
&= 9
\end{aligned} \\
\begin{aligned}
c_{22} &= 2 \cdot 0 + 2 \cdot 2 + 1 \cdot 1 \\
&= 5.
\end{aligned}
\end{gathered}
\]

\end{example}

Matrix multiplication satisfies several properties analogous to those of ordinary multiplication of numbers, namely associativity, $A(BC) = (AB)C$ for every $A \in M_{m,r}$, $B \in M_{r,n}$, $C \in M_{n,p}$; distributivity over addition, $A(B+C) = AB + AC$ for every $A \in M_{m,r}$ and $B, C \in M_{r,n}$, and $(A+B)C = AC + BC$ for every $A, B \in M_{m,r}$ and $C \in M_{r,n}$; and compatibility with scalar multiplication,
\[
\begin{aligned}
A(tB) &= (tA)B \\
&= t(AB),
\end{aligned}
\]
for every $t \in \mathbb{R}$, $A \in M_{m,r}$, $B \in M_{r,n}$.

\begin{remark}
In marked contrast with the multiplication of ordinary numbers, which is independent of the order of the factors, matrix multiplication is not commutative: even for square matrices, it may happen that $AB \neq BA$. For example, with
\[
\begin{gathered}
\begin{pmatrix} 0 & 1 \\ 0 & 0 \end{pmatrix} \\
\begin{pmatrix} 0 & 0 \\ 1 & 0 \end{pmatrix},
\end{gathered}
\]
one computes
\[
\begin{pmatrix} 0 & 1 \\ 0 & 0 \end{pmatrix}\begin{pmatrix} 0 & 0 \\ 1 & 0 \end{pmatrix} = \begin{pmatrix} 1 & 0 \\ 0 & 0 \end{pmatrix},
\]
while
\[
\begin{pmatrix} 0 & 0 \\ 1 & 0 \end{pmatrix}\begin{pmatrix} 0 & 1 \\ 0 & 0 \end{pmatrix} = \begin{pmatrix} 0 & 0 \\ 0 & 1 \end{pmatrix},
\]
and these two products clearly differ. Moreover, the product of two matrices can vanish identically even when neither factor is the zero matrix: taking
\[
A = B = \begin{pmatrix} 0 & 1 \\ 0 & 0 \end{pmatrix},
\]
one finds $AB = O$, despite $A \neq O$.
\end{remark}


A further operation on matrices, of frequent use throughout the theory, consists of interchanging rows and columns.

\begin{definition}[Transpose]
The transpose of an $m \times n$ matrix $A$ is the $n \times m$ matrix $A^{\top}$ obtained by turning rows into columns and vice versa, that is, $[A^{\top}]_{ij} = a_{ji}$, for every $i = 1, \dots, n$ and $j = 1, \dots, m$.
\end{definition}
\subsection{Solving Linear Systems by Gaussian Elimination}

When the coefficient matrix of a homogeneous linear system is already in row echelon form, the corresponding system can be solved directly by a simple procedure known as back-substitution. Consider, for instance, the homogeneous system $U\mathbf{x} = \vec{0}$, where
\[
U = \begin{pmatrix} 2 & 1 & 1 \\ 0 & \tfrac{3}{2} & \tfrac{1}{2} \\ 0 & 0 & 0 \end{pmatrix},
\]
corresponding explicitly to the system
\[
\begin{cases}
2x + y + z = 0, \\
\tfrac{3}{2} y + \tfrac{1}{2} z = 0, \\
0 = 0.
\end{cases}
\]

This system is readily solved by back-substitution: one first sets $z = t$, in correspondence with the column of $U$ containing no pivot; one then obtains $y$ from the second equation; and finally one obtains $x$ from the first equation. With $t \in \mathbb{R}$ arbitrary, the resulting solution set is
\[
\begin{pmatrix} x \\ y \\ z \end{pmatrix} = t \begin{pmatrix} -\tfrac{1}{3} \\ -\tfrac{1}{3} \\ 1 \end{pmatrix}.
\]

For a general coefficient matrix, not necessarily already in row echelon form, one first reduces the matrix to row echelon form by means of a systematic procedure known as the Gauss elimination method, before applying back-substitution as above.

The Gauss elimination method consists of a sequence of elementary row operations applied to a matrix, with the goal of transforming it into row echelon form, that is, of filling the lower left-hand corner of the matrix with zeros as far as possible. Three types of elementary row operations are permitted: swapping two rows; multiplying a row by a nonzero number; and adding a multiple of one row to another row. Using these three operations, any matrix can always be transformed into a matrix in row echelon form.

Once, in addition, every leading coefficient of the resulting matrix equals $1$, and every column containing a leading coefficient has zeros in all its remaining entries, the matrix is said to be in reduced row echelon form. This final reduced form is unique, in the sense that it does not depend on the particular sequence of elementary row operations used to reach it.

\begin{example}
Applying the Gauss elimination method to the matrix
\[
\begin{pmatrix} 1 & 1 & -1 & 1 \\ 1 & 3 & 1 & 9 \\ 3 & 11 & 5 & 35 \end{pmatrix},
\]
one first subtracts the first row from the second, and three times the first row from the third, obtaining
\[
\begin{pmatrix} 1 & 3 & 1 & 9 \\ 0 & -2 & -2 & -8 \\ 0 & 0 & 0 & 0 \end{pmatrix}
\]
after also swapping the first two rows for convenience; multiplying the second row by $-(1/2)$ and adding it to the third yields the row echelon form
\[
\begin{pmatrix} 1 & 3 & 1 & 9 \\ 0 & 1 & 1 & 4 \\ 0 & 0 & 0 & 0 \end{pmatrix}.
\]

\end{example}

The Gauss elimination method is the principal tool for finding the solution set of a linear system with an arbitrary coefficient matrix $A$. Its validity rests on the fact that, whenever an elementary operation is performed on the equations of a system, the solution set of the system remains unchanged; and, crucially, an elementary operation on the equations of a system corresponds precisely to an elementary row operation on the corresponding coefficient matrix. Consequently, if $A$ is the coefficient matrix of the homogeneous system $A\mathbf{x} = \vec{0}$, and $U$ denotes the matrix obtained by applying the Gauss algorithm to $A$, then
\[
A\mathbf{x} = \vec{0} \iff U\mathbf{x} = \vec{0},
\]
so that the original system and the reduced system share exactly the same solution set.

\begin{example}
Consider the homogeneous system
\[
\begin{cases}
2x + y + z = 0, \\
x + 2y + z = 0, \\
x + y + \tfrac{2}{3} z = 0,
\end{cases}
\]
with coefficient matrix
\[
A = \begin{pmatrix} 2 & 1 & 1 \\ 1 & 2 & 1 \\ 1 & 1 & \tfrac{2}{3} \end{pmatrix}.
\]
Applying the Gauss elimination method to $A$ yields the reduced matrix
\[
U = \begin{pmatrix} 2 & 1 & 1 \\ 0 & \tfrac{3}{2} & \tfrac{1}{2} \\ 0 & 0 & 0 \end{pmatrix},
\]
and solving $U\mathbf{x} = \vec{0}$ by back-substitution exactly as in the earlier example yields the following solution set, with $t \in \mathbb{R}$ arbitrary:
\[
\begin{pmatrix} x \\ y \\ z \end{pmatrix} = t \begin{pmatrix} -\tfrac{1}{3} \\ -\tfrac{1}{3} \\ 1 \end{pmatrix}.
\]

\end{example}

The Gauss elimination method extends immediately to nonhomogeneous systems of the general form $A\mathbf{x} = \mathbf{b}$. Letting $[A | \mathbf{b}]$ denote the augmented matrix of the system, and $[U | \mathbf{b}']$ the matrix obtained by applying the Gauss algorithm to $[A | \mathbf{b}]$, one has the equivalence
\[
A\mathbf{x} = \mathbf{b} \iff U\mathbf{x} = \mathbf{b}',
\]
so that the two systems again share the same solution set.

\begin{example}
Consider the nonhomogeneous system
\[
\begin{cases}
2x + y + z = 6, \\
x + 2y + z = 6, \\
x + y + \tfrac{2}{3} z = 4.
\end{cases}
\]
The coefficient matrix and the augmented matrix are
\[
\begin{gathered}
A = \begin{pmatrix} 2 & 1 & 1 \\ 1 & 2 & 1 \\ 1 & 1 & \tfrac{2}{3} \end{pmatrix} \\
[A|\mathbf{b}] = \begin{pmatrix} 2 & 1 & 1 & 6 \\ 1 & 2 & 1 & 6 \\ 1 & 1 & \tfrac{2}{3} & 4 \end{pmatrix}.
\end{gathered}
\]
Applying the Gauss algorithm to $[A|\mathbf{b}]$ yields the reduced augmented matrix
\[
[U|\mathbf{b}'] = \begin{pmatrix} 2 & 1 & 1 & 6 \\ 0 & \tfrac{3}{2} & \tfrac{1}{2} & 3 \\ 0 & 0 & 0 & 0 \end{pmatrix};
\]
solving $U\mathbf{x} = \mathbf{b}'$ by back-substitution, setting $z = t$ for the column with no pivot, and then obtaining $y$ and $x$ in succession, yields the following solution set, with $t \in \mathbb{R}$ arbitrary:
\[
\begin{pmatrix} x \\ y \\ z \end{pmatrix} = \begin{pmatrix} 2 \\ 2 \\ 0 \end{pmatrix} + t \begin{pmatrix} -\tfrac{1}{3} \\ -\tfrac{1}{3} \\ 1 \end{pmatrix}.
\]

As a related exercise, applying the Gauss elimination method to the similar system with third equation $x + y + 5z = 4$ in place of $x+y+(2/3)z=4$ shows that this modified system has no solutions at all.
\end{example}

The general procedure just illustrated applies to an arbitrary matrix $A \in M_{m,n}$ and vector $\mathbf{b} \in \mathbb{R}^m$: letting $[A|\mathbf{b}]$ denote the augmented matrix and $[U|\mathbf{b}']$ the corresponding matrix obtained from the Gauss algorithm, the two systems $A\mathbf{x} = \mathbf{b}$ and $U\mathbf{x} = \mathbf{b}'$ have the same solution set, that is, they are equivalent. The original system is consistent if and only if the number of pivots of $U$ equals the number of pivots of $[U|\mathbf{b}'];$ denoting this common number of pivots by $p$, the solution set depends on $n - p$ free parameters, corresponding precisely to the number of columns of $U$ that contain no pivot.

\section{Determinants, Inverses, and Rank}

\subsection{Determinants: Definition and Recursive Computation}

Given a square matrix, the determinant is a scalar value that encodes certain fundamental properties of the linear transformation described by the matrix. In the simplest case, the determinant admits an explicit closed-form expression.

\begin{definition}[Determinant of a $2\times 2$ Matrix]
Given the matrix
\[
A = \begin{pmatrix} a & b \\ c & d \end{pmatrix} \in M_2,
\]
the determinant of $A$ is the number $\operatorname{Det}(A) = ad - bc$.
\end{definition}

For matrices of larger size, the determinant is computed recursively, by means of a procedure known as Laplace expansion, which reduces the computation of an $n \times n$ determinant to the computation of several $(n-1) \times (n-1)$ determinants. Given $A \in M_n$, denote by $A_{ij} \in M_{n-1}$ the submatrix obtained by deleting the $i$-th row and the $j$-th column of $A$. The real number
\[
C_{ij} = (-1)^{i+j} \operatorname{Det}(A_{ij})
\]
is called the cofactor of $a_{ij}$. It can be proved that the sum of the products of the elements of any row, or any column, of $A$ with their respective cofactors always yields the same value, regardless of which row or column is chosen; this common value is, by definition, the determinant of $A$.

\begin{theorem}[Laplace Expansion for the Determinant]
Fix any row $i \in \{1, \dots, n\}$. Then
\[
\begin{aligned}
\operatorname{Det}(A) &= a_{i1} C_{i1} + a_{i2} C_{i2} + \cdots + a_{in} C_{in} \\
&= \sum_{j=1}^{n} (-1)^{i+j} a_{ij} \operatorname{Det}(A_{ij}).
\end{aligned}
\]
Analogously, fixing any column $j$,
\[
\begin{aligned}
\operatorname{Det}(A) &= a_{1j} C_{1j} + a_{2j} C_{2j} + \cdots + a_{nj} C_{nj} \\
&= \sum_{i=1}^{n} (-1)^{i+j} a_{ij} \operatorname{Det}(A_{ij}).
\end{aligned}
\]

\end{theorem}

Several elementary properties of the determinant follow readily from this recursive definition. If $A$ has a row, or a column, consisting entirely of zeros, then $\operatorname{Det}(A) = 0$. The determinant is invariant under transposition,
\[
\operatorname{Det}(A^{\top}) = \operatorname{Det}(A).
\]

If $A$ is a triangular matrix, whether upper or lower, then its determinant equals the product of its diagonal entries,
\[
\operatorname{Det}(A) = a_{11} a_{22} \cdots a_{nn}.
\]

Finally, the determinant of a product of two square matrices of the same order factors as the product of the individual determinants.

\begin{theorem}[Binet's Theorem]
If $A$ and $B$ are square matrices of the same order, then
\[
\operatorname{Det}(AB) = \operatorname{Det}(A) \cdot \operatorname{Det}(B).
\]

\end{theorem}

It can be shown that the determinant is, in fact, the unique function $M_n \to \mathbb{R}$ satisfying three fundamental properties, and this axiomatic characterization is sometimes taken as the very definition of the determinant. First, $\operatorname{Det}(I) = 1$. Second, the determinant is alternating: whenever two rows of a matrix are identical, its determinant equals zero. Third, the determinant is multilinear, that is, it is a linear function of each individual row separately; for instance, linearity in the first row means that
\[
\operatorname{Det}\begin{pmatrix} \mathbf{a}_1 + t \mathbf{b}_1 \\ \mathbf{a}_2 \\ \vdots \\ \mathbf{a}_n \end{pmatrix} = \operatorname{Det}\begin{pmatrix} \mathbf{a}_1 \\ \mathbf{a}_2 \\ \vdots \\ \mathbf{a}_n \end{pmatrix} + t \cdot \operatorname{Det}\begin{pmatrix} \mathbf{b}_1 \\ \mathbf{a}_2 \\ \vdots \\ \mathbf{a}_n \end{pmatrix}
\]
for every $t \in \mathbb{R}$.

On the basis of these three fundamental properties, one readily deduces how the determinant changes when elementary row operations are performed. If two rows of a matrix are interchanged, the determinant changes sign; and if any linear combination of the other rows is added to a given row, the determinant remains unchanged. These two facts allow one to compute the determinant of an arbitrary matrix $A$ by relating it to the determinant of a triangular matrix $U$ obtained from $A$ by Gaussian elimination, using only these two operations: if an even number of row interchanges is used, or none at all, then
\[
\operatorname{Det}(A) = \operatorname{Det}(U);
\]
whereas if an odd number of row interchanges is used, then
\[
\operatorname{Det}(A) = -\operatorname{Det}(U).
\]

This relationship yields a particularly efficient criterion for determining when a matrix has nonzero determinant. Since $\operatorname{Det}(A) \neq 0$ if and only if $\operatorname{Det}(U) \neq 0$, and since
\[
\operatorname{Det}(U) = u_{11} u_{22} \cdots u_{nn}
\]
for a triangular matrix, it follows that $\operatorname{Det}(A) \neq 0$ if and only if every diagonal entry $u_{ii}$ of $U$ is nonzero, that is, if and only if $U$ has exactly $n$ pivots. Consequently, $\operatorname{Det}(A) \neq 0$ if and only if $\operatorname{Rank}(A) = n$.

\begin{definition}
A square matrix $A$ such that $\operatorname{Det}(A) \neq 0$ is called a nonsingular matrix.
\end{definition}

The nonsingularity of a matrix, as characterized by a nonvanishing determinant, has several important consequences for the associated linear system and linear transformation. A square linear system $A\mathbf{x} = \mathbf{b}$ has a unique solution if and only if $\operatorname{Det}(A) \neq 0$; in particular, if $\mathbf{b} = \vec{0}$, the only solution is the null vector. The columns, and equally the rows, of a square matrix are linearly independent vectors if and only if $\operatorname{Det}(A) \neq 0$. Finally, if $A$ is the representative matrix of a linear transformation $L : \mathbb{R}^n \to \mathbb{R}^n$, then $L$ is injective if and only if $L$ is surjective, and both properties hold if and only if $\operatorname{Det}(A) \neq 0$.

Figure~\ref{fig:atlas-determinant-area} interprets the determinant as an oriented area scale factor.

\begin{figure}[!htbp]
\centering
\begin{tikzpicture}[>=Stealth]
\begin{scope}[scale=1.25]
\fill[figblue!15] (0,0) rectangle (1,1);
\draw[black!35] (0,0) rectangle (1,1);
\draw[vecA] (0,0) -- (1,0) node[below] {$e_1$};
\draw[vecB] (0,0) -- (0,1) node[left] {$e_2$};
\node at (0.5,0.5) {$1$};
\end{scope}
\draw[->,thick] (2,0.8) -- (3,0.8) node[midway,above] {$A$};
\begin{scope}[shift={(4,0)},scale=1.25]
\fill[figgreen!15] (0,0) -- (2,0.5) -- (3,2.5) -- (1,2) -- cycle;
\draw[black!35] (2,0.5) -- (3,2.5) -- (1,2);
\draw[vecA] (0,0) -- (2,0.5) node[below right] {$Ae_1$};
\draw[vecB] (0,0) -- (1,2) node[above left] {$Ae_2$};
\node at (1.5,1.25) {$7/2$};
\end{scope}
\end{tikzpicture}
\caption{The unit square on the left is mapped to the parallelogram spanned by $(2,1/2)$ and $(1,2)$ on the right. Its area is $7/2$, the absolute determinant of the matrix with those columns. The positive determinant preserves the orientation of the ordered basis.}
\label{fig:atlas-determinant-area}
\end{figure}

\subsection{The Inverse of a Matrix}

\begin{definition}[Invertible Matrix]
Let $A \in M_n$. We say that $A$ is invertible if there exists a matrix $B \in M_n$ such that
\[
\begin{aligned}
AB &= BA \\
&= I.
\end{aligned}
\]
When such an inverse exists, it is unique, and is denoted by $A^{-1}$.
\end{definition}

\begin{theorem}
If $A$ is invertible, then $\operatorname{Det}(A) \neq 0$.
\end{theorem}

\begin{proof}
Since $A A^{-1} = I$, Binet's Theorem yields
\[
\begin{aligned}
\operatorname{Det}(A) \operatorname{Det}(A^{-1}) &= \operatorname{Det}(AA^{-1}) \\
&= \operatorname{Det}(I) \\
&= 1,
\end{aligned}
\]
which shows in particular that $\operatorname{Det}(A) \neq 0$, and moreover that
\[
\operatorname{Det}(A^{-1}) = 1/\operatorname{Det}(A).
\]

\end{proof}

The converse implication also holds, so that nonsingularity is not merely necessary but also sufficient for invertibility.

\begin{theorem}
If $\operatorname{Det}(A) \neq 0$, then $A$ is invertible.
\end{theorem}

\begin{proof}
Since $\operatorname{Det}(A) \neq 0$, for each $i = 1, \dots, n$ the system $A\mathbf{x} = \mathbf{e}_i$, where $\mathbf{e}_i$ denotes the $i$-th canonical basis vector, has a unique solution, which we call $\mathbf{b}^i$. Defining $B = [\mathbf{b}^1 | \cdots | \mathbf{b}^n]$, one observes that
\[
\begin{aligned}
AB &= [A\mathbf{b}^1 | \cdots | A\mathbf{b}^n] \\
&= [\mathbf{e}_1 | \cdots | \mathbf{e}_n] \\
&= I,
\end{aligned}
\]
establishing that $A$ is invertible with inverse $B$.
\end{proof}

In practice, the inverse of a nonsingular matrix can be computed by applying the Gauss elimination method to the augmented matrix $[A | I]$ until a reduced matrix of the form $[I | B]$ is obtained; the resulting matrix $B$ is then the inverse of $A$. An explicit formula for the inverse, expressed directly in terms of cofactors, is also available.

\begin{theorem}[Laplace's Formula for the Inverse of a Matrix]
The inverse of a nonsingular matrix $A$ is given by
\[
A^{-1} = \frac{1}{\operatorname{Det}(A)} C^{\top},
\]
where $C^{\top}$ denotes the transpose of the matrix of cofactors $C = [C_{ij}]$, with
\[
C_{ij} = (-1)^{i+j} \operatorname{Det}(A_{ij}).
\]

\end{theorem}

\begin{example}
We find the inverse, if it exists, of
\[
A = \begin{pmatrix} 3 & 4 & 2 \\ 1 & 2 & 0 \\ 1 & 0 & 1 \end{pmatrix}.
\]
Expanding along the second column, we compute $\operatorname{Det}(A_{12}) = 1$ and $\operatorname{Det}(A_{22}) = 1$, and hence
\[
\operatorname{Det}(A) = -4 \cdot 1 + 2 \cdot 1 \cdot \ldots,
\]
which, after completing the expansion, equals $2 \neq 0$; thus $A$ is invertible. Computing the remaining cofactors, namely $\operatorname{Det}(A_{11}) = 4$, $\operatorname{Det}(A_{13}) = -4$, $\operatorname{Det}(A_{21}) = 2$, $\operatorname{Det}(A_{23}) = -2$, $\operatorname{Det}(A_{31}) = 4$, $\operatorname{Det}(A_{32}) = 2$, $\operatorname{Det}(A_{33}) = -2$, one obtains the inverse
\[
A^{-1} = \frac{1}{2} \begin{pmatrix} 4 & -2 & 4 \\ -1 & 1 & -2 \\ -4 & 2 & -2 \end{pmatrix}.
\]

\end{example}

\subsection{The Rank of a Matrix}

We now consider matrices $A \in M_{m,n}$ of arbitrary, not necessarily square, dimension. Recall that the rank of $A$, already introduced in connection with the column space of $A$, is the dimension of $\operatorname{Col}(A)$, that is, the maximum number of linearly independent columns of $A$. The determinant provides an alternative, and in many contexts more computationally direct, characterization of this quantity.

\begin{definition}
Given $A \in M_{m,n}$, we call a minor of $A$ the determinant of any square submatrix of $A$, and we call the order of the minor the dimension of the corresponding square submatrix.
\end{definition}

\begin{proposition}
$\operatorname{Rank}(A) = r$ if and only if there exists a nonzero minor of order $r$, and every minor of order greater than $r$ vanishes. Equivalently, the rank of $A$ is the largest order among all nonzero minors of $A$.
\end{proposition}

Directly verifying that every minor of a given order beyond a certain threshold vanishes can be computationally expensive, since the number of such minors grows rapidly with the dimensions of the matrix. The following theorem shows that a much smaller number of minors need actually be checked.

\begin{theorem}[Kronecker's Theorem]
$\operatorname{Rank}(A) = r$ if and only if there is a nonzero minor of order $r$, corresponding to a certain square submatrix $B \in M_r$, and the minor of every square submatrix of order $r+1$ that contains $B$ as a submatrix vanishes.
\end{theorem}

\begin{example}
Consider the matrix
\[
A = \begin{pmatrix} 0 & 1 & 3 & 2 \\ 1 & 0 & -5 & -1 \\ 1 & 2 & 1 & 3 \end{pmatrix}.
\]
Since $A$ has only three rows, $\operatorname{Rank}(A) \leq 3$. The square submatrix
\[
B = \begin{pmatrix} 0 & 1 \\ 1 & 0 \end{pmatrix}
\]
has nonzero determinant, so $\operatorname{Rank}(A) \geq 2$. There are exactly two square submatrices of order $3$ containing $B$ as a submatrix, and direct computation shows that both of their determinants vanish; by Kronecker's Theorem, we conclude that $\operatorname{Rank}(A) = 2$.
\end{example}

\subsection{Applications: Change of Basis and Similar Matrices}

The determinant, and the associated notion of invertibility, play a decisive role in two important constructions of linear algebra, namely the change of basis and the similarity of matrices.

Given any basis $\mathcal{B} = \{\mathbf{b}^1, \dots, \mathbf{b}^n\}$ of $\mathbb{R}^n$, and a vector $\mathbf{v} = (x_1, \dots, x_n)$ expressed in the canonical basis, one may ask how to find the coordinates $\hat{\mathbf{x}} = (\hat{x}_1, \dots, \hat{x}_n)$ of $\mathbf{v}$ with respect to $\mathcal{B}$. Setting
\[
M = [\mathbf{b}^1 | \cdots | \mathbf{b}^n] \in M_n,
\]
called the change of basis matrix, one has $\hat{\mathbf{x}} = M^{-1} \mathbf{x}$. Indeed, writing
\[
\begin{aligned}
\mathbf{v} &= \sum_{i=1}^n x_i \mathbf{e}_i \\
&= I \mathbf{x}
\end{aligned}
\]
using the canonical basis, and
\[
\begin{aligned}
\mathbf{v} &= \sum_{i=1}^n \hat{x}_i \mathbf{b}^i \\
&= M \hat{\mathbf{x}}
\end{aligned}
\]
using the basis $\mathcal{B}$, one obtains $\mathbf{x} = M\hat{\mathbf{x}}$; since $\operatorname{Det}(M) \neq 0$, because the columns of $M$ form a basis and are therefore linearly independent, one may invert this relation to find $\hat{\mathbf{x}} = M^{-1}\mathbf{x}$.

A closely related construction concerns the representation of a linear map in different bases. Let $A \in M_n$ be the representative matrix of a linear map $L : \mathbb{R}^n \to \mathbb{R}^n$ in the canonical basis, and take any basis $\mathcal{B} = \{\mathbf{b}^1, \dots, \mathbf{b}^n\}$ of $\mathbb{R}^n$. Setting again $M = [\mathbf{b}^1 | \cdots | \mathbf{b}^n]$, the representative matrix $B$ of $L$ in the new basis $\mathcal{B}$ is given by $B = M^{-1} A M$. Two matrices $A$ and $B$ related in this way, that is, satisfying $B = M^{-1}AM$ for some invertible matrix $M$, are said to be similar; two $n$-by-$n$ matrices are similar if and only if they represent the same linear transformation of $\mathbb{R}^n$ with respect to two different bases.

\section{Vector Spaces and Fundamental Subspaces}

\subsection{Vector Operations and Vector Spaces}

The foundational object underlying the entire theory developed in this document is the vector. Two fundamental operations govern the algebra of vectors, and it is useful to describe them both geometrically and in coordinates.

Geometrically, if two vectors $\mathbf{u}$ and $\mathbf{v}$ are positioned so that the initial point of $\mathbf{v}$ coincides with the terminal point of $\mathbf{u}$, then their sum $\mathbf{u} + \mathbf{v}$ is defined as the vector running from the initial point of $\mathbf{u}$ to the terminal point of $\mathbf{v}$. The scalar multiple $t\mathbf{v}$, for a real number $t$, is defined as the vector whose length equals the length of $\mathbf{v}$ multiplied by $|t|$; provided $t \neq 0$, the vector $t\mathbf{v}$ points in the same direction as $\mathbf{v}$, with the same orientation if $t$ is positive and the opposite orientation if $t$ is negative.

In coordinates, given two vectors $\mathbf{x} = (x_1, \dots, x_m)$ and $\mathbf{y} = (y_1, \dots, y_m)$ in $\mathbb{R}^m$, their sum, denoted $\mathbf{x}+\mathbf{y}$, is the vector of $\mathbb{R}^m$ given by
\[
\mathbf{x} + \mathbf{y} = (x_1+y_1, \dots, x_m+y_m).
\]
Given a number $t \in \mathbb{R}$ and a vector
\[
\mathbf{x} = (x_1, \dots, x_m) \in \mathbb{R}^m,
\]
the scalar multiple $t\mathbf{x}$ is the vector
\[
t\mathbf{x} = (t x_1, \dots, t x_m) \in \mathbb{R}^m.
\]

Consider the following set of $m$ vectors of $\mathbb{R}^m$,
\[
\mathbf{e}_1 = \begin{pmatrix} 1 \\ 0 \\ \vdots \\ 0 \end{pmatrix}, \qquad \mathbf{e}_2 = \begin{pmatrix} 0 \\ 1 \\ \vdots \\ 0 \end{pmatrix}, \qquad \ldots, \qquad \mathbf{e}_m = \begin{pmatrix} 0 \\ 0 \\ \vdots \\ 1 \end{pmatrix}.
\]

Given any vector $\mathbf{v} = (v_1, \dots, v_m)$, one can write $\mathbf{v} = v_1 \mathbf{e}_1 + \cdots + v_m \mathbf{e}_m$; in this sense, every vector of $\mathbb{R}^m$ can be reached from $\mathbf{e}_1, \dots, \mathbf{e}_m$ using only the two fundamental operations of addition and scalar multiplication.

Several pieces of terminology, central to the entire subject of linear algebra, arise directly from this observation. When a vector is obtained from a set of given vectors using only the two fundamental operations, one says that it is a linear combination of the given vectors. The set of all possible vectors reachable in this way from a given collection of vectors is called the span of those vectors. The vectors $\mathbf{e}_1, \dots, \mathbf{e}_m$ are linearly independent, meaning that none of them belongs to the span of the others; and one says that they form a basis of $\mathbb{R}^m$, capturing the fact that they constitute a set of linearly independent vectors that spans the entire space.

The properties observed above for $\mathbb{R}^m$ are, in fact, common to a much wider class of mathematical objects, axiomatized by the notion of a vector space.

\begin{definition}[Vector Space]
A vector space is a set $V$, whose elements are called vectors, on which two fundamental operations are defined: an internal operation called sum, $V \times V \to V$, $(\mathbf{u}, \mathbf{v}) \mapsto \mathbf{u}+\mathbf{v}$, and an external operation called scalar multiplication, $\mathbb{R} \times V \to V$, $(t, \mathbf{u}) \mapsto t\mathbf{u}$. These two operations must satisfy the following properties, for every $\mathbf{u}, \mathbf{v}, \mathbf{z} \in V$ and every $s, t \in \mathbb{R}$: associativity of addition,
\[
\mathbf{u} + (\mathbf{v}+\mathbf{z}) = (\mathbf{u}+\mathbf{v}) + \mathbf{z};
\]
commutativity of addition, $\mathbf{u}+\mathbf{v} = \mathbf{v}+\mathbf{u}$; existence of a zero element $\vec{0} \in V$ such that $\mathbf{u}+\vec{0} = \mathbf{u}$; existence, for each $\mathbf{u} \in V$, of an element $\mathbf{w} \in V$ such that $\mathbf{u}+\mathbf{w} = \vec{0}$; the associativity relation $s(t\mathbf{u}) = (st)\mathbf{u}$; the identity relation $1\mathbf{u} = \mathbf{u}$; and the two distributive laws $t(\mathbf{u}+\mathbf{v}) = t\mathbf{u}+t\mathbf{v}$ and $(s+t)\mathbf{u} = s\mathbf{u}+t\mathbf{u}$.
\end{definition}

\begin{definition}[Subspace]
A subset $W \subset V$ is called a subspace of $V$ if it is closed under the operations of sum and scalar multiplication, that is, if $\mathbf{u}, \mathbf{v} \in W$ implies $\mathbf{u}+\mathbf{v} \in W$, and $\mathbf{u} \in W$, $t \in \mathbb{R}$ implies $t\mathbf{u} \in W$.
\end{definition}

\begin{remark}
Every subspace necessarily contains the zero vector: if $W$ is a subspace of $V$, then $\vec{0}_V \in W$. Moreover, a subspace $W \subset V$ inherits the vector space structure of $V$: a subspace is itself a vector space, with the same operations restricted to $W$.
\end{remark}

\begin{definition}[Linear Combination]
Given $k$ vectors $\mathbf{v}^1, \mathbf{v}^2, \dots, \mathbf{v}^k \in V$ and $k$ given numbers $a_1, a_2, \dots, a_k \in \mathbb{R}$, the vector
\[
\mathbf{w} = a_1 \mathbf{v}^1 + a_2 \mathbf{v}^2 + \cdots + a_k \mathbf{v}^k \in V
\]
is called the linear combination of $\mathbf{v}^1, \mathbf{v}^2, \dots, \mathbf{v}^k$ with coefficients $a_1, a_2, \dots, a_k$.
\end{definition}

\begin{definition}[Span]
Given $k$ vectors $\mathbf{v}^1, \dots, \mathbf{v}^k \in V$, we denote by $\operatorname{span}(\mathbf{v}^1, \dots, \mathbf{v}^k)$ the set of all possible linear combinations of these vectors, that is,
\[
\operatorname{span}(\mathbf{v}^1, \dots, \mathbf{v}^k) = \left\{ \sum_{i=1}^{k} x_i \mathbf{v}^i : x_1, \dots, x_k \in \mathbb{R} \right\}.
\]

\end{definition}

One readily observes that $\operatorname{span}(\mathbf{v}^1, \dots, \mathbf{v}^k)$ is always a subspace of $V$, since it is closed under both fundamental operations; this subspace is accordingly called the subspace generated by the vectors $\mathbf{v}^1, \dots, \mathbf{v}^k$.

\subsection{Linear Independence, Basis, and Dimension}

\begin{definition}[Linear Dependence]
The vectors $\mathbf{v}^1, \mathbf{v}^2, \dots, \mathbf{v}^k \in V$ are called linearly dependent if there exist scalars $x_1, \dots, x_k$, not all zero, such that
\[
x_1 \mathbf{v}^1 + x_2 \mathbf{v}^2 + \cdots + x_k \mathbf{v}^k = \vec{0}.
\]

\end{definition}

If not every scalar in this relation vanishes, then at least one, say $x_1$, is nonzero, and the defining equation can be rearranged as
\[
\mathbf{v}^1 = -(x_2/x_1) \mathbf{v}^2 - \cdots - (x_k/x_1) \mathbf{v}^k,
\]
exhibiting $\mathbf{v}^1$ as a linear combination of the remaining vectors. In this sense, $\mathbf{v}^1$ is redundant, and one has the equality
\[
\operatorname{span}(\mathbf{v}^1, \mathbf{v}^2, \dots, \mathbf{v}^k) = \operatorname{span}(\mathbf{v}^2, \dots, \mathbf{v}^k),
\]
showing that a linearly dependent set can always be reduced without changing its span.

\begin{definition}[Linear Independence]
The vectors $\mathbf{v}^1, \mathbf{v}^2, \dots, \mathbf{v}^k \in V$ are called linearly independent if the relation
\[
x_1 \mathbf{v}^1 + x_2 \mathbf{v}^2 + \cdots + x_k \mathbf{v}^k = \vec{0}
\]
implies $x_i = 0$ for every $i = 1, \dots, k$.
\end{definition}

A set of linearly independent vectors is, in this sense, minimal for the generation of its span, since no vector in the set is redundant. In the particular case of two vectors, linear dependence admits a simple geometric characterization: two vectors $\mathbf{v}^1$ and $\mathbf{v}^2$ in $V$ are linearly dependent if and only if they are proportional, that is, if and only if $\mathbf{v}^1 = \lambda \mathbf{v}^2$ for some nonzero scalar $\lambda$.

\begin{example}
We investigate whether the vectors $\mathbf{v}^1 = (1,0,0)$, $\mathbf{v}^2 = (1, 1/2, 1)$, and $\mathbf{v}^3 = (5,1,2)$ are linearly independent in $\mathbb{R}^3$. Imposing the condition
\[
x_1 \mathbf{v}^1 + x_2 \mathbf{v}^2 + x_3 \mathbf{v}^3 = \vec{0},
\]
and writing the resulting equation in components, one obtains the linear system
\[
\begin{cases}
x_1 + x_2 + 5 x_3 = 0, \\
\tfrac{1}{2} x_2 + x_3 = 0, \\
x_2 + 2 x_3 = 0.
\end{cases}
\]

This system reduces to $x_1 = -3x_3$ and $x_2 = -2x_3$, with $x_3$ free to take any value in $\mathbb{R}$; consequently, the three vectors are linearly dependent. In particular, fixing $x_3 = 1$ yields the explicit relation $\mathbf{v}^3 = 3\mathbf{v}^1 + 2\mathbf{v}^2$.
\end{example}

The example just presented illustrates the general strategy for investigating the linear independence of a family of vectors in $\mathbb{R}^m$. Given $k$ vectors $\mathbf{a}^1, \dots, \mathbf{a}^k \in \mathbb{R}^m$, with
\[
\mathbf{a}^j = (a_{1j}, a_{2j}, \dots, a_{mj})
\]
for each $j = 1, \dots, k$, one investigates their linear dependence or independence by asking whether the equation
\[
x_1 \mathbf{a}^1 + x_2 \mathbf{a}^2 + \cdots + x_k \mathbf{a}^k = \vec{0}
\]
in $\mathbb{R}^m$ has the unique solution $x_i = 0$ for all $i$. Writing this equation out in components yields precisely a homogeneous linear system of $m$ equations in the $k$ unknowns $x_1, \dots, x_k$; consequently, proving the linear independence of a family of $k$ vectors in $\mathbb{R}^m$ amounts to showing that this associated homogeneous linear system has only the trivial solution, connecting the study of linear independence directly to the theory of linear systems developed earlier.

\begin{definition}[Basis]
Given vectors $\mathbf{b}^1, \mathbf{b}^2, \dots, \mathbf{b}^d \in V$, the set $\{\mathbf{b}^1, \dots, \mathbf{b}^d\}$ is called a basis for $V$ if the vectors $\mathbf{b}^1, \dots, \mathbf{b}^d$ are linearly independent, and if $\{\mathbf{b}^1, \dots, \mathbf{b}^d\}$ is a set of generators of $V$, that is,
\[
\operatorname{span}(\mathbf{b}^1, \dots, \mathbf{b}^d) = V.
\]

\end{definition}

By the very definition of span, every vector $\mathbf{v}$ in the space can be written as a linear combination of the basis vectors; the following theorem shows that the coefficients in this expansion are, in fact, uniquely determined.

\begin{theorem}
Let $\mathcal{B} = \{\mathbf{b}^1, \mathbf{b}^2, \dots, \mathbf{b}^d\}$ be a basis for $V$. Then, for every $\mathbf{v} \in V$, there is a unique set of $d$ scalars $x_1, x_2, \dots, x_d \in \mathbb{R}$ such that
\[
\mathbf{v} = \sum_{i=1}^{d} x_i \mathbf{b}^i.
\]
The scalars $x_1, \dots, x_d$ are called the components, or coordinates, of $\mathbf{v}$ in the basis $\mathcal{B}$.
\end{theorem}

\begin{proof}[Proof of uniqueness]
Suppose, for the sake of contradiction, that there exist a vector $\mathbf{v} \in V$ and two distinct sets of coefficients, $x_1, \dots, x_d$ and $\tilde{x}_1, \dots, \tilde{x}_d$, such that
\[
\begin{aligned}
\mathbf{v} &= \sum_{i=1}^{d} x_i \mathbf{b}^i \\
&= \sum_{i=1}^{d} \tilde{x}_i \mathbf{b}^i.
\end{aligned}
\]
Subtracting the two expressions yields
\[
\sum_{i=1}^{d} (x_i - \tilde{x}_i) \mathbf{b}^i = \vec{0}.
\]

Since the vectors $\mathbf{b}^1, \dots, \mathbf{b}^d$ are linearly independent, this forces $x_i - \tilde{x}_i = 0$ for every $i$, that is, $x_i = \tilde{x}_i$ for every $i$, contradicting the assumed distinctness of the two sets of coefficients.
\end{proof}

The number of vectors in a basis of a given vector space turns out to be an intrinsic property of the space itself, independent of the particular basis chosen, as recorded in the following fundamental theorem.

\begin{theorem}[Dimension Theorem]
If $V$ has a basis consisting of $d \in \mathbb{N}$ vectors, then every basis of $V$ consists of exactly $d$ vectors.
\end{theorem}

In light of this theorem, one says that $V$ has dimension $d$, and writes $\operatorname{Dim}(V) = d$. It follows moreover that, if $V$ has dimension $d$, then any set of $d$ linearly independent vectors of $V$ is automatically a basis of $V$; no separate verification that the set spans $V$ is required. As a basic example, the vector space $\mathbb{R}^m$ has dimension $m$; the three vectors $\mathbf{v}^1$, $\mathbf{v}^2$, $\mathbf{v}^3$ considered in the earlier example, having been shown to be linearly dependent, therefore fail to constitute a basis of $\mathbb{R}^3$.

Several useful facts relate the dimension of a subspace to the dimension of the ambient space. Let $W \subset V$ be a subspace, with $\operatorname{Dim}(V) = m$. Then
\[
\operatorname{Dim}(W) \leq \operatorname{Dim}(V);
\]
if $W$ is a proper subspace of $V$, then $\operatorname{Dim}(W) < m$ strictly; and, conversely, if $\operatorname{Dim}(W) = m$, then necessarily $W = V$. The null subspace $W = \{\vec{0}_V\}$ has dimension $0$, and is in fact the unique subspace of dimension zero. Finally, if
\[
W = \operatorname{span}(\mathbf{v}^1, \dots, \mathbf{v}^k),
\]
then $\operatorname{Dim}(W) = r$, where $r$ denotes the largest number of linearly independent vectors among the generators $\mathbf{v}^1, \dots, \mathbf{v}^k$; these $r$ vectors then form a basis of $W$. It follows that $r \leq m$ and $r \leq k$.

\subsection{Column Space, Kernel, and the Nullity--Rank Theorem}

\begin{definition}
Let $A = [a_{ij}] \in M_{m,n}$ be a given matrix. The column space of $A$ is the subspace of $\mathbb{R}^m$ generated by the columns of $A$, that is,
\[
\operatorname{Col}(A) = \operatorname{span}(\mathbf{a}^1, \dots, \mathbf{a}^n) \subset \mathbb{R}^m,
\]
where $\mathbf{a}^j$ denotes the $j$-th column of $A$. Equivalently,
\[
\operatorname{Col}(A) = \{ A\mathbf{x} : \mathbf{x} \in \mathbb{R}^n \},
\]
in light of the interpretation of the matrix-vector product as a linear combination of columns.
\end{definition}

\begin{definition}
The dimension of $\operatorname{Col}(A)$ is called the rank of $A$, denoted $\operatorname{Rank}(A)$.
\end{definition}

\begin{remark}
The computation of the rank is immediate whenever the matrix is already in row echelon form: the columns of the matrix that contain a pivot are linearly independent, while every remaining column belongs to the span of the pivot columns preceding it. Consequently,
\[
\operatorname{Dim}(\operatorname{Col}(U)) = \operatorname{Rank}(U)
\]
equals the number of pivots of $U$, for a matrix $U$ in row echelon form; this same number of pivots also indicates the maximum number of linearly independent rows of $U$.
\end{remark}

\begin{definition}
Let $A = [a_{ij}] \in M_{m,n}$ be a given matrix. The kernel of $A$ is the subset of $\mathbb{R}^n$ given by
\[
\operatorname{Ker}(A) = \{ \mathbf{x} \in \mathbb{R}^n : A\mathbf{x} = \vec{0}_m \} \subset \mathbb{R}^n.
\]

\end{definition}

\begin{theorem}
$\operatorname{Ker}(A)$ is a vector space.
\end{theorem}

The proof of this theorem consists in verifying directly that $\operatorname{Ker}(A)$ is closed under vector addition and scalar multiplication, a straightforward consequence of the linearity of the matrix-vector product established earlier.

Figure~\ref{fig:atlas-kernel-image} separates the directions preserved by a projection from those annihilated by it.

\begin{figure}[!htbp]
\centering
\begin{tikzpicture}[scale=1.5,>=Stealth]
\draw[->,figblue,thick] (-1,0) -- (4,0) node[right] {$\operatorname{Im}P$};
\draw[->,figred,thick] (0,-0.6) -- (0,2.8) node[above] {$\ker P$};
\draw[->,black!70,thick] (0,0) -- (2.7,2.1);
\draw[->,figgreen,thick,dashed] (2.7,2.1) -- (2.7,0.05);
\fill (2.7,2.1) circle (2pt);
\fill[figblue] (2.7,0) circle (2pt);
\node[above right] at (2.7,2.1) {$(x,y)$};
\node[below] at (2.7,0) {$(x,0)$};
\node[below left] at (0,0) {$0$};
\end{tikzpicture}
\caption{For the projection $P(x,y)=(x,0)$, the horizontal axis is the image and the vertical axis is the kernel. A point and its projection differ by a vertical vector in the kernel. Both subspaces have dimension one, in agreement with rank plus nullity equaling two.}
\label{fig:atlas-kernel-image}
\end{figure}

We now come to one of the central results of linear algebra, relating the dimension of the kernel of a matrix to the dimension of its column space.

\begin{theorem}[Nullity-Rank Theorem]
Given any matrix $A$ with $n$ columns,
\[
\operatorname{Dim}(\operatorname{Ker}(A)) + \operatorname{Dim}(\operatorname{Col}(A)) = n.
\]

\end{theorem}

\begin{proof}
Assume that
\[
\operatorname{Dim}(\operatorname{Ker}(A)) = k,
\]
and let $\{\mathbf{u}^1, \dots, \mathbf{u}^k\}$ be a basis of $\operatorname{Ker}(A)$, so that
\[
\operatorname{Ker}(A) = \operatorname{span}(\mathbf{u}^1, \dots, \mathbf{u}^k).
\]

By an important theorem of linear algebra, one may select $n-k$ further vectors $\mathbf{w}^{k+1}, \dots, \mathbf{w}^n \in \mathbb{R}^n$, linearly independent both among themselves and from the vectors $\mathbf{u}^1, \dots, \mathbf{u}^k$, such that $\{\mathbf{u}^1, \dots, \mathbf{u}^k, \mathbf{w}^{k+1}, \dots, \mathbf{w}^n\}$ is a basis of $\mathbb{R}^n$; note that
\[
\operatorname{span}(\mathbf{w}^{k+1}, \dots, \mathbf{w}^n) \cap \operatorname{Ker}(A) = \{\vec{0}_n\}.
\]

The strategy of the proof is to show that the $n-k$ vectors $A\mathbf{w}^{k+1}, \dots, A\mathbf{w}^n \in \mathbb{R}^m$ form a basis of $\operatorname{Col}(A)$, from which
\[
\operatorname{Dim}(\operatorname{Col}(A)) = n-k
\]
follows immediately, establishing the theorem. We first show that $\{A\mathbf{w}^{k+1}, \dots, A\mathbf{w}^n\}$ generates $\operatorname{Col}(A)$. Given any $\mathbf{x} \in \mathbb{R}^n$, there exist numbers $\alpha_1, \dots, \alpha_n \in \mathbb{R}$ such that
\[
\mathbf{x} = \sum_{j=1}^{k} \alpha_j \mathbf{u}^j + \sum_{j=k+1}^{n} \alpha_j \mathbf{w}^j.
\]
By linearity of the matrix-vector product,
\[
\begin{aligned}
A\mathbf{x} &= \sum_{j=1}^{k} \alpha_j A\mathbf{u}^j + \sum_{j=k+1}^{n} \alpha_j A\mathbf{w}^j \\
&= \sum_{j=k+1}^{n} \alpha_j A\mathbf{w}^j,
\end{aligned}
\]
since $A\mathbf{u}^j = \vec{0}_m$ for each $j = 1, \dots, k$, as $\mathbf{u}^j \in \operatorname{Ker}(A)$. This shows that $A\mathbf{x} \in \operatorname{span}(A\mathbf{w}^{k+1}, \dots, A\mathbf{w}^n)$; since
\[
\operatorname{Col}(A) = \{A\mathbf{x} : \mathbf{x} \in \mathbb{R}^n\},
\]
this establishes the claim. We next show that the vectors $A\mathbf{w}^{k+1}, \dots, A\mathbf{w}^n$ are linearly independent. Take $\beta_1, \dots, \beta_{n-k} \in \mathbb{R}$ such that
\[
\sum_{j=1}^{n-k} \beta_j A\mathbf{w}^{k+j} = \vec{0}_m.
\]
By linearity, this equals
\[
A( \sum_{j=1}^{n-k} \beta_j \mathbf{w}^{k+j} ),
\]
so that the vector
\[
\mathbf{u} := \sum_{j=1}^{n-k} \beta_j \mathbf{w}^{k+j} \in \operatorname{Ker}(A).
\]

But this same vector $\mathbf{u}$ also lies in $\operatorname{span}(\mathbf{w}^{k+1}, \dots, \mathbf{w}^n)$, and we recorded above that this span intersects $\operatorname{Ker}(A)$ only in the zero vector; hence $\mathbf{u} = \vec{0}_n$. Since the vectors $\mathbf{w}^{k+1}, \dots, \mathbf{w}^n$ are themselves linearly independent, the relation
\[
\sum_{j=1}^{n-k} \beta_j \mathbf{w}^{k+j} = \vec{0}_n
\]
forces $\beta_j = 0$ for every $j = 1, \dots, n-k$, establishing the linear independence of $A\mathbf{w}^{k+1},\allowbreak \dots, A\mathbf{w}^n$ and completing the proof.
\end{proof}


Let $A \in M_{m,n}$. A homogeneous system consists in finding $\mathbf{x} \in \mathbb{R}^n$ such that $A\mathbf{x} = \vec{0}_m$. Such a system is always solvable, since the null vector $\mathbf{x} = \vec{0}_n$ is trivially a solution; moreover, the entire solution set coincides precisely with $\operatorname{Ker}(A)$, and hence, by the Nullity-Rank Theorem, its dimension equals $n - \operatorname{Rank}(A)$, where $n$ denotes the number of unknowns.

\begin{theorem}
Let $A \in M_{m,n}$. The solution set of the homogeneous system $A\mathbf{x} = \vec{0}_m$ is a vector space of dimension $n - \operatorname{Rank}(A)$. Setting $r = \operatorname{Rank}(A)$, the following hold: if $r = n$, then $\mathbf{x} = \vec{0}_n$ is the unique solution of the system; if $r < n$, the system has $n-r$ linearly independent solutions $\mathbf{x}^1, \dots, \mathbf{x}^{n-r}$, and every solution is given by
\[
S = \{ \mathbf{x} = t_1 \mathbf{x}^1 + \cdots + t_{n-r} \mathbf{x}^{n-r} : t_i \in \mathbb{R} \}.
\]

The system therefore has infinitely many solutions, depending on $n-r$ free parameters, whenever $r < n$.
\end{theorem}

\begin{remark}
Recall that, when studying homogeneous linear systems by the Gauss elimination method, the solution set was found to depend on $n-p$ parameters, where $p$ denotes the number of pivots of the reduced matrix. Comparing this earlier result with the theorem just stated, in which the solutions depend on $n - \operatorname{Rank}(A)$ parameters, one concludes that $p = \operatorname{Rank}(A)$; that is, the rank of a matrix equals the number of pivots of its row-echelon reduction. This identity has a further important consequence: the number of linearly independent columns of a matrix necessarily coincides with the number of linearly independent rows.
\end{remark}

Let $A \in M_{m,n}$ and $\mathbf{b} \in \mathbb{R}^m$ with $\mathbf{b} \neq \vec{0}_m$. A nonhomogeneous system consists in finding $\mathbf{x} \in \mathbb{R}^n$ such that $A\mathbf{x} = \mathbf{b}$. Recalling that the matrix-vector product equals
\[
A\mathbf{x} = x_1 \mathbf{a}^1 + x_2 \mathbf{a}^2 + \cdots + x_n \mathbf{a}^n,
\]
solving $A\mathbf{x} = \mathbf{b}$ is equivalent to finding coefficients $x_1, \dots, x_n$, if any exist, such that $\mathbf{b}$ can be written as a linear combination of the columns $\mathbf{a}^1, \mathbf{a}^2, \dots, \mathbf{a}^n$ of $A$. This reformulation immediately raises the question of whether such a system is always solvable.

The answer is that the system is solvable, that is, consistent, if and only if $\mathbf{b} \in \operatorname{Col}(A)$, that is, if and only if
\[
\operatorname{Col}([A|\mathbf{b}]) = \operatorname{Col}(A).
\]
Since $\operatorname{Col}(A) \subset \operatorname{Col}([A|\mathbf{b}])$ always holds, the two column spaces coincide if and only if
\[
\operatorname{Dim}(\operatorname{Col}([A|\mathbf{b}])) = \operatorname{Dim}(\operatorname{Col}(A)),
\]
yielding the following solvability condition.

\begin{theorem}[Solvability Condition]
A nonhomogeneous system of linear equations $A\mathbf{x} = \mathbf{b}$ has a solution if and only if the rank of its coefficient matrix $A$ equals the rank of its augmented matrix $[A|\mathbf{b}]$, that is,
\[
\operatorname{Rank}([A|\mathbf{b}]) = \operatorname{Rank}(A).
\]

\end{theorem}

Unlike the homogeneous case, the solution set of a nonhomogeneous system does not, in general, form a vector space, since it need not contain the zero vector; nonetheless, it is closely related to the solution set of the associated homogeneous system.

\begin{theorem}[Structure of the Solution Set]
Given any specific solution $\mathbf{v}_p$ of $A\mathbf{x} = \mathbf{b}$, the entire solution set of the nonhomogeneous system is
\[
S = \{ \mathbf{y} = \mathbf{v} + \mathbf{v}_p : \mathbf{v} \in \mathbb{R}^n, \ A\mathbf{v} = \vec{0}_m \} = \operatorname{Ker}(A) + \mathbf{v}_p.
\]

\end{theorem}

\begin{proof}
It suffices to observe that $\mathbf{y} \in \mathbb{R}^n$ satisfies $A\mathbf{y} = \mathbf{b}$ if and only if
\[
\begin{aligned}
A(\mathbf{y} - \mathbf{v}_p) &= A\mathbf{y} - A\mathbf{v}_p \\
&= \mathbf{b} - \mathbf{b} \\
&= \vec{0}_m.
\end{aligned}
\]

Hence $\mathbf{y} - \mathbf{v}_p \in \operatorname{Ker}(A)$, so that $\mathbf{y}$ equals the sum of the particular solution $\mathbf{v}_p$ and an element of $\operatorname{Ker}(A)$, and conversely every such sum is a solution.
\end{proof}

Combining the solvability condition with the structure theorem for the solution set yields the following comprehensive result, which fully characterizes the possible behaviors of a general linear system.

\begin{theorem}[Rouch\'e--Capelli Theorem]
Let $\mathbf{b} \neq \vec{0}_m$. The nonhomogeneous system $A\mathbf{x} = \mathbf{b}$ has a solution if and only if
\[
\operatorname{Rank}(A) = \operatorname{Rank}([A|\mathbf{b}]).
\]

In this case, setting $r = \operatorname{Rank}(A)$: if $r = n$, the solution is unique; if $r < n$, the system has infinitely many solutions, depending on $n-r$ free parameters, and the solution set is
\[
S = \{ \mathbf{y} = \mathbf{v}_p + t_1 \mathbf{x}^1 + \cdots + t_{n-r} \mathbf{x}^{n-r} : t_i \in \mathbb{R} \},
\]
where $\mathbf{x}^1, \dots, \mathbf{x}^{n-r}$ are $n-r$ linearly independent solutions of the associated homogeneous system $A\mathbf{x} = \vec{0}_m$, and $\mathbf{v}_p$ is any particular solution of the nonhomogeneous system.
\end{theorem}

\section{Linear Maps}

\subsection{Linear Maps and Their Representative Matrix}

Given $A \in M_{m,n}$, one may associate with it the map $L : \mathbb{R}^n \to \mathbb{R}^m$, $L(\mathbf{x}) = A\mathbf{x}$, and we have already observed that this map is linear, satisfying
\[
\begin{gathered}
L(\mathbf{x}+\mathbf{y}) = L(\mathbf{x}) + L(\mathbf{y}) \\
L(t\mathbf{x}) = tL(\mathbf{x})
\end{gathered}
\]
for every $\mathbf{x}, \mathbf{y} \in \mathbb{R}^n$ and $t \in \mathbb{R}$, or, equivalently,
\[
L(s\mathbf{x}+t\mathbf{y}) = sL(\mathbf{x}) + tL(\mathbf{y}).
\]

It is worth noting that the columns of $A$ are precisely the images of the canonical basis vectors, $L(\mathbf{e}_j) = \mathbf{a}^j$ for every $j = 1, \dots, n$; consequently, if $\mathbf{x} = (x_1, \dots, x_n)$, then
\[
\begin{aligned}
L(\mathbf{x}) &= A\mathbf{x} \\
&= \sum_{j=1}^n x_j L(\mathbf{e}_j).
\end{aligned}
\]

Conversely, every linear map between finite-dimensional Euclidean spaces can be represented in this fashion by a suitable matrix.

\begin{theorem}[Representation Theorem for Linear Maps]
Let $L : \mathbb{R}^n \to \mathbb{R}^m$ be any linear map, that is, a map satisfying
\[
L(s\mathbf{x}+t\mathbf{y}) = sL(\mathbf{x}) + tL(\mathbf{y})
\]
for every $\mathbf{x}, \mathbf{y} \in \mathbb{R}^n$ and $s, t \in \mathbb{R}$. Setting $\mathbf{a}^j = L(\mathbf{e}_j)$ for $j = 1, \dots, n$, and
\[
A = [\mathbf{a}^1 | \mathbf{a}^2 | \cdots | \mathbf{a}^n] \in M_{m,n},
\]
called the representative matrix of $L$ in the canonical basis, one has $L(\mathbf{x}) = A\mathbf{x}$ for every $\mathbf{x} \in \mathbb{R}^n$.
\end{theorem}

\begin{proof}
Take any vector
\[
\mathbf{x} = (x_1, \dots, x_n) = \sum_{j=1}^n x_j \mathbf{e}_j \in \mathbb{R}^n.
\]
By the linearity of $L$,
\[
\begin{aligned}
L(\mathbf{x}) &= L( \sum_{j=1}^n x_j \mathbf{e}_j ) \\
&= \sum_{j=1}^n x_j L(\mathbf{e}_j) \\
&= \sum_{j=1}^n x_j \mathbf{a}^j \\
&= A\mathbf{x},
\end{aligned}
\]
where the last equality follows from the interpretation of the matrix-vector product as a linear combination of columns.
\end{proof}

This representation theorem shows that the action of any linear map between Euclidean spaces is represented, once bases have been fixed, by a matrix-vector product.

\begin{example}
Consider the linear transformation of $\mathbb{R}^2$ describing a rotation of $180$ degrees counterclockwise combined with a stretching by a factor of $2$. To find its representative matrix, one determines where the canonical basis vectors land under this transformation: $\mathbf{e}_1$ maps to $(-2,0)$, and $\mathbf{e}_2$ maps to $(0,-2)$. Any vector $\mathbf{x} = x\mathbf{e}_1 + y\mathbf{e}_2$ therefore lands on
\[
x(-2,0) + y(0,-2) = \begin{pmatrix} -2 & 0 \\ 0 & -2 \end{pmatrix} \begin{pmatrix} x \\ y \end{pmatrix},
\]
identifying the representative matrix of the transformation as
\[
\begin{pmatrix} -2 & 0 \\ 0 & -2 \end{pmatrix}.
\]

\end{example}

Figure~\ref{fig:ch8-linear-map} shows how a linear map transforms a coordinate grid and its basis vectors.

\begin{figure}[!htbp]
\centering
\begin{tikzpicture}[scale=1.05,>=Stealth]
  \begin{scope}
    \fill[figblue!8] (0,0) rectangle (1,1);
    \foreach \gridindex in {0,...,2} {
      \draw[gridline] (\gridindex,0) -- (\gridindex,2);
      \draw[gridline] (0,\gridindex) -- (2,\gridindex);
    }
    \draw[->] (-0.4,0) -- (2.6,0) node[right,font=\small] {$x_1$};
    \draw[->] (0,-0.4) -- (0,3.4) node[above,font=\small] {$x_2$};
    \draw[vecB] (0,0) -- (1,0) node[below=4pt,font=\small] {$e_1$};
    \draw[vecC] (0,0) -- (0,1) node[left=4pt,font=\small] {$e_2$};
    \node[below left,font=\small] at (0,0) {$0$};
  \end{scope}
  \draw[->,very thick] (2.9,1.5) -- (4.8,1.5)
    node[midway,above=5pt,font=\small] {$T(\bx)=A\bx$};
  \begin{scope}[xshift=5.5cm]
    \fill[figblue!8] (0,0) -- (1.3,0.4) -- (1.9,1.5) -- (0.6,1.1) -- cycle;
    \foreach \gridindex in {0,...,2} {
      \draw[gridline] ({1.3*\gridindex},{0.4*\gridindex})
        -- ({1.3*\gridindex+1.2},{0.4*\gridindex+2.2});
      \draw[gridline] ({0.6*\gridindex},{1.1*\gridindex})
        -- ({2.6+0.6*\gridindex},{0.8+1.1*\gridindex});
    }
    \draw[->] (-0.4,0) -- (4.2,0) node[right,font=\small] {$y_1$};
    \draw[->] (0,-0.4) -- (0,3.4) node[above,font=\small] {$y_2$};
    \draw[vecB] (0,0) -- (1.3,0.4)
      node[right=4pt,font=\small,fill=white,inner sep=1pt] {$A e_1$};
    \draw[vecC] (0,0) -- (0.6,1.1)
      node[above=4pt,font=\small,fill=white,inner sep=1pt] {$A e_2$};
    \node[below left,font=\small] at (0,0) {$0$};
  \end{scope}
\end{tikzpicture}
\caption[Linear transformation of a grid]{A coordinate grid before and after an
  illustrative linear map $T(\bx)=A\bx$. The shaded unit square on the left maps
  to the shaded parallelogram spanned by the columns $Ae_1,Ae_2$ on the right.
  Matching colors identify each basis vector and its image; both coordinate
  systems have their axes intersecting at the origin.}
\label{fig:ch8-linear-map}
\end{figure}

\subsection{Image, Kernel, and Injectivity/Surjectivity}

A central goal in the study of linear maps is to determine conditions guaranteeing their injectivity or surjectivity. To this end, given a linear map $L : \mathbb{R}^n \to \mathbb{R}^m$ with representative matrix $A$, we define the range, or image, of $L$ as the collection of all images $L(\mathbf{x})$ for $\mathbf{x} \in \mathbb{R}^n$, that is, $\operatorname{Im}(L) = \operatorname{Col}(A)$; and we define the kernel of $L$ as the subspace of vectors whose image is the zero vector,
\[
\operatorname{Ker}(L) = \operatorname{Ker}(A).
\]

\begin{theorem}
The linear map $L : \mathbb{R}^n \to \mathbb{R}^m$ is surjective if and only if $\operatorname{Rank}(A) = m$.
\end{theorem}

This theorem follows immediately from the fact that $L$ is surjective precisely when $\operatorname{Im}(L) = \operatorname{Col}(A)$ equals the entire codomain $\mathbb{R}^m$, which occurs exactly when
\[
\operatorname{Dim}(\operatorname{Col}(A)) = m,
\]
that is, when $\operatorname{Rank}(A) = m$.

\begin{theorem}[Injectivity]
The linear map $L : \mathbb{R}^n \to \mathbb{R}^m$ is one-to-one if and only if $\operatorname{Ker}(L) = \{\vec{0}_n\}$.
\end{theorem}

\begin{proof}
Assume first that $L$ is injective. Since $L(\vec{0}_n) = \vec{0}_m$ by linearity, and $L$ is injective, no other vector can share this image; hence, for every $\mathbf{x} \neq \vec{0}_n$, one has $L(\mathbf{x}) \neq \vec{0}_m$, so that no nonzero vector belongs to $\operatorname{Ker}(L)$, that is,
\[
\operatorname{Ker}(L) = \{\vec{0}_n\}.
\]

Conversely, assume $\operatorname{Ker}(L) = \{\vec{0}_n\}$, and suppose, for the sake of contradiction, that $L(\mathbf{x}_1) = L(\mathbf{x}_2)$ for some $\mathbf{x}_1 \neq \mathbf{x}_2$ in $\mathbb{R}^n$. By linearity,
\[
\begin{aligned}
L(\mathbf{x}_1 - \mathbf{x}_2) &= L(\mathbf{x}_1) - L(\mathbf{x}_2) \\
&= \vec{0}_m,
\end{aligned}
\]
so that $\mathbf{x}_1 - \mathbf{x}_2 \in \operatorname{Ker}(L)$. But the only element of $\operatorname{Ker}(L)$ is $\vec{0}_n$, so $\mathbf{x}_1 - \mathbf{x}_2 = \vec{0}_n$, contradicting the assumption $\mathbf{x}_1 \neq \mathbf{x}_2$; this contradiction establishes injectivity.
\end{proof}

The linear map and its representative matrix have the same kernel:
\[
\operatorname{Ker}(L) = \operatorname{Ker}(A),
\]
The Nullity-Rank Theorem yields
\[
\operatorname{Dim}(\operatorname{Ker}(L)) = n - \operatorname{Rank}(A),
\]
from which one obtains the following corollary.

\begin{corollary}
The linear map $L : \mathbb{R}^n \to \mathbb{R}^m$ is one-to-one if and only if $\operatorname{Rank}(A) = n$.
\end{corollary}


Given a square matrix $A \in M_n$, the linear function $T : \mathbb{R}^n \to \mathbb{R}^n$, $T(\mathbf{x}) = A\mathbf{x}$, represents a geometric transformation of $\mathbb{R}^n$ into itself, called an endomorphism. Combining the injectivity and surjectivity criteria established above with the earlier characterization of nonsingularity in terms of the determinant, one obtains a particularly clean criterion for such maps: the linear function $T : \mathbb{R}^n \to \mathbb{R}^n$ is injective if and only if it is surjective, and both properties hold precisely when the rank of its representative matrix equals $n$, that is, $\operatorname{Rank}(A) = n$; equivalently, as established earlier, this occurs if and only if $\operatorname{Det}(A) \neq 0$.

\section{Spectral and Singular-Value Methods}

\subsection{Eigenvalues, Eigenvectors, and the Characteristic Equation}

Given $A \in M_n$, consider the associated linear transformation $T : \mathbb{R}^n \to \mathbb{R}^n$, $T(\mathbf{x}) = A\mathbf{x}$. Among the many vectors that this transformation acts upon, some special nonzero vectors change, under the action of $T$, by at most a scalar factor, that is, they are simply stretched, or possibly reversed, by the transformation, in the sense that $T(\mathbf{v}) = \lambda \mathbf{v}$ for some scalar $\lambda$. Such a special vector is called an eigenvector, and the corresponding scalar $\lambda$, the factor by which it is stretched, is called the eigenvalue associated with that eigenvector.

\begin{definition}
Let $A \in M_n$ be a given square matrix. A vector $\mathbf{u} \in \mathbb{R}^n$, $\mathbf{u} \neq \vec{0}$, is an eigenvector of $A$ if $A\mathbf{u} = \lambda \mathbf{u}$ for some $\lambda \in \mathbb{R}$. The scalar $\lambda$ is called the eigenvalue corresponding to that eigenvector.
\end{definition}

To determine systematically the pairs $\mathbf{u} \neq \vec{0}$ and $\lambda \in \mathbb{R}$ satisfying this relation, one observes that the defining equation $A\mathbf{u} = \lambda \mathbf{u}$ can be equivalently rewritten as $(A - \lambda I)\mathbf{u} = \vec{0}$, which is a homogeneous linear system in the unknown $\mathbf{u}$, with coefficient matrix $A - \lambda I$. Finding eigenvectors of $A$, for a fixed value of $\lambda$, therefore amounts to finding nontrivial solutions of this homogeneous system, and, by the general theory of linear systems developed earlier, such nontrivial solutions exist if and only if the coefficient matrix $A - \lambda I$ is singular, that is, if and only if its determinant vanishes. This observation yields the following fundamental characterization of eigenvalues.

\begin{theorem}
A scalar $\lambda \in \mathbb{R}$ is an eigenvalue of $A$ if and only if $\lambda$ satisfies the equation $\operatorname{Det}(A - \lambda I) = 0$.
\end{theorem}

Several pieces of terminology accompany this characterization. The function
\[
P(t) = \operatorname{Det}(A - tI)
\]
is a polynomial of degree $n$ in the variable $t$, called the characteristic polynomial of $A$; the equation $P(t) = 0$ is called the characteristic equation of $A$. By the Fundamental Theorem of Algebra, the characteristic polynomial of an $n$-by-$n$ matrix has exactly $n$ complex roots, counted with multiplicity and not necessarily distinct. Given any fixed real eigenvalue $\lambda$ of $A$, one may factor the characteristic polynomial as
\[
P(t) = (t-\lambda)^k Q(t),
\]
for some polynomial $Q$ satisfying $Q(\lambda) \neq 0$ and some positive integer $k$; this integer $k$ is called the algebraic multiplicity of $\lambda$, and is denoted $a_\lambda$.

Associated with each real eigenvalue $\lambda$ of $A$ is a distinguished subspace, obtained as the solution set of the corresponding homogeneous system: the solution set of $(A-\lambda I)\mathbf{x} = \vec{0}$ is the vector space $V_\lambda = \operatorname{Ker}(A - \lambda I)$, called the eigenspace corresponding to $\lambda$, and it contains precisely all the eigenvectors of $A$ associated with $\lambda$, together with the zero vector. The dimension of this eigenspace is called the geometric multiplicity of $\lambda$, denoted $g_\lambda$; by the Nullity-Rank Theorem,
\[
\begin{aligned}
g_\lambda &= \operatorname{Dim}(V_\lambda) \\
&= n - \operatorname{Rank}(A - \lambda I).
\end{aligned}
\]
Figure~\ref{fig:ch8-eigenvectors} highlights the directions preserved by a linear transformation.

\begin{figure}[!htbp]
\centering
\begin{tikzpicture}[scale=1.65,>=Stealth]
  \draw[->] (-2.3,0) -- (2.3,0) node[right,font=\scriptsize]{$x_1$};
  \draw[->] (0,-1.6) -- (0,1.9) node[above,font=\scriptsize]{$x_2$};
  \draw[figblue!35, thick] (-2,-1) -- (2,1);
  \draw[figred!35, thick] (-1.2,1.6) -- (1.2,-1.6);
  \draw[vecA] (0,0) -- (1,0.5)
    node[below right,fill=white,inner sep=1pt,font=\scriptsize]{$\bv_1$};
  \draw[vecB] (0,0) -- (0.6,-0.8)
    node[below right,fill=white,inner sep=1pt,font=\scriptsize]{$\bv_2$};
  \draw[vecA,dashed] (0,0) -- (1.8,0.9)
    node[above right,fill=white,inner sep=1pt,font=\scriptsize]{$A\bv_1=\lambda_1 \bv_1$};
  \draw[vecB,dashed] (0,0) -- (0.3,-0.4);
  \node[figred,fill=white,inner sep=1pt,font=\scriptsize,anchor=north west]
    at (0.42,-0.44) {$A\bv_2=\lambda_2 \bv_2$};
  \draw[black!55,very thick,->] (0,0) -- (-1,1)
    node[left,fill=white,inner sep=1pt,font=\scriptsize]{$\bu$};
  \draw[black!55,very thick,dotted,->] (0,0) -- (-0.4,1.55)
    node[above right,fill=white,inner sep=1pt,font=\scriptsize]{$A\bu$};
  \node[anchor=south east,fill=white,inner sep=1pt,font=\scriptsize]
    at (2.2,-1.5) {$\lambda_1 \approx 1.8,\ \lambda_2\approx 0.5$};
\end{tikzpicture}
\caption[Eigenvectors: invariant directions]{Invariant lines (eigenspaces) in light colors:
  $\bv_1$ is stretched ($\lambda_1>1$), $\bv_2$ is contracted ($0<\lambda_2<1$), both remaining
  on their respective lines; the generic vector $\bu$ (gray) is also rotated.}
\label{fig:ch8-eigenvectors}
\end{figure}

The characterization of eigenvalues in terms of the characteristic equation yields a systematic three-step procedure for determining the eigenvalues and eigenvectors of a given matrix. The first step consists of finding the values of $t \in \mathbb{R}$ satisfying $\operatorname{Det}(A - tI) = 0$; this is a polynomial equation of degree $n$ in $t$, whose real solutions, if any, are precisely the real eigenvalues of $A$, and from whose factorization one determines the algebraic multiplicity $a_\lambda$ of each eigenvalue $\lambda$ found. The second step consists of writing, for each eigenvalue $\lambda$ found in the first step, the corresponding homogeneous linear system with coefficient matrix $A - \lambda I$. The third step consists of solving this homogeneous system $(A - \lambda I)\mathbf{x} = \vec{0}$; every nonzero solution $\mathbf{x}$ of this system is an eigenvector of $A$ corresponding to the eigenvalue $\lambda$, and the dimension of the solution space $\operatorname{Ker}(A - \lambda I)$ determines the geometric multiplicity $g_\lambda$.

\begin{example}
We determine the eigenvalues and eigenvectors of
\[
A = \begin{pmatrix} 2 & 0 \\ 2 & 1 \end{pmatrix}.
\]
We first use the characteristic equation to determine the eigenvalues:
\[
\operatorname{Det}(A - \lambda I) = \begin{vmatrix} 2-\lambda & 0 \\ 2 & 1-\lambda \end{vmatrix} = (1-\lambda)(2-\lambda),
\]
so that the eigenvalues are $\lambda_1 = 1$ and $\lambda_2 = 2$. To determine the eigenvectors corresponding to $\lambda_i$, we solve $A\mathbf{x} = \lambda_i \mathbf{x}$ with $\mathbf{x} = (x,y)$, that is, the system $2x - \lambda_i x = 0$, $2x + y - \lambda_i y = 0$. For $\lambda_1 = 1$, this system yields
\[
\mathbf{x} = (0,t) = t(0,1),
\]
for $t \neq 0$; for $\lambda_2 = 2$, it yields
\[
\mathbf{x} = (t, 2t) = t(1,2),
\]
for $t \neq 0$.
\end{example}

Several general remarks concerning the multiplicities $g_\lambda$ and $a_\lambda$ are worth recording. The geometric multiplicity $g_\lambda$ equals the maximum number of linearly independent eigenvectors associated with the eigenvalue $\lambda$, and satisfies $g_\lambda \geq 1$, since every eigenvalue has at least one associated eigenvector by definition. A further, more difficult result, whose proof we omit, shows that $g_\lambda \leq a_\lambda$ always; this inequality may be strict, and whenever $g_\lambda = a_\lambda$ one says that $\lambda$ is a regular eigenvalue. In particular, if $\lambda$ is simple, that is, if $a_\lambda = 1$, then necessarily $\lambda$ is regular, since
\[
\begin{aligned}
1 &\leq g_\lambda \\
&\leq a_\lambda \\
&= 1
\end{aligned}
\]
forces
\[
\begin{aligned}
g_\lambda &= a_\lambda \\
&= 1.
\end{aligned}
\]

\begin{proposition}
Two eigenvectors corresponding to two different eigenvalues are linearly independent.
\end{proposition}

\begin{proof}
Suppose $A\mathbf{v}^1 = \lambda_1 \mathbf{v}^1$ and $A\mathbf{v}^2 = \lambda_2 \mathbf{v}^2$ for some $\lambda_1 \neq \lambda_2$. To show that $\mathbf{v}^1, \mathbf{v}^2$ are linearly independent, consider a linear combination
\[
c_1 \mathbf{v}^1 + c_2 \mathbf{v}^2 = \vec{0}.
\]
Applying $A$ to both sides, and using linearity together with the eigenvalue equations, yields
\[
\vec{0} = c_1 \lambda_1 \mathbf{v}^1 + c_2 \lambda_2 \mathbf{v}^2.
\]
Multiplying the original relation
\[
c_1 \mathbf{v}^1 + c_2 \mathbf{v}^2 = \vec{0}
\]
by $\lambda_2$ and subtracting the two resulting equations, we obtain
\[
\vec{0} = c_1(\lambda_2 - \lambda_1) \mathbf{v}^1.
\]

Since an eigenvector is, by definition, nonzero, and $\lambda_2 \neq \lambda_1$, this forces $c_1 = 0$; substituting back into the original relation then yields $c_2 = 0$ as well, establishing linear independence.
\end{proof}

\subsection{Diagonalization and Spectral Coordinates}

We now investigate under which conditions a matrix $A \in M_n$ possesses the maximum possible number, namely $n$, of linearly independent eigenvectors. Since the total number of linearly independent eigenvectors of $A$ equals
\[
\sum_\lambda g_\lambda,
\]
summed over the distinct eigenvalues $\lambda$ of $A$, and since this sum cannot exceed $n$, the dimension of the ambient space, one seeks conditions under which the sum
\[
\sum_\lambda g_\lambda
\]
attains exactly the value $n$.

\begin{proposition}
Let $A \in M_n$ be such that $A$ has exactly $n$ real eigenvalues, not necessarily distinct, that is,
\[
\sum_\lambda a_\lambda = n,
\]
and each eigenvalue $\lambda$ is regular, that is, $g_\lambda = a_\lambda$. Then there exists a basis of $\mathbb{R}^n$ made up of eigenvectors of $A$.
\end{proposition}

An important special class of matrices automatically satisfies these hypotheses, as recorded in the following theorem.

\begin{theorem}[Spectral Theorem]
Let $A \in M_n$ be a symmetric matrix. Then $A$ has $n$ real eigenvalues, not necessarily distinct, and $A$ has $n$ linearly independent, mutually orthogonal eigenvectors. Consequently, for every symmetric matrix $A \in M_n$, there exists an orthonormal basis of $\mathbb{R}^n$ made up of eigenvectors of $A$.
\end{theorem}


When a matrix $A \in M_n$ possesses $n$ linearly independent eigenvectors, it can be decomposed in a particularly illuminating way, known as its eigendecomposition. Let $\mathbf{u}^1, \dots, \mathbf{u}^n$ be $n$ linearly independent eigenvectors of $A$, with corresponding eigenvalues $\lambda_1, \dots, \lambda_n$. Construct the nonsingular matrix $P = [\mathbf{u}^1 | \cdots | \mathbf{u}^n]$, called the change of basis matrix, and the diagonal matrix
\[
D = \operatorname{Diag}(\lambda_1, \dots, \lambda_n) = \begin{pmatrix} \lambda_1 & 0 & \cdots & 0 \\ 0 & \lambda_2 & \cdots & 0 \\ \vdots & \vdots & \ddots & \vdots \\ 0 & 0 & \cdots & \lambda_n \end{pmatrix}.
\]
The product $AP$ is computed as
\[
\begin{aligned}
AP &= A[\mathbf{u}^1 | \cdots | \mathbf{u}^n] \\
&= [A\mathbf{u}^1 | \cdots | A\mathbf{u}^n] \\
&= [\lambda_1 \mathbf{u}^1 | \cdots | \lambda_n \mathbf{u}^n],
\end{aligned}
\]
using the defining property of eigenvectors applied to each column separately. On the other hand, the product $PD$ is computed as
\[
PD = [\mathbf{u}^1 | \cdots | \mathbf{u}^n] \begin{pmatrix} \lambda_1 & 0 & \cdots & 0 \\ 0 & \lambda_2 & \cdots & 0 \\ \vdots & \vdots & \ddots & \vdots \\ 0 & 0 & \cdots & \lambda_n \end{pmatrix} = [\lambda_1 \mathbf{u}^1 | \cdots | \lambda_n \mathbf{u}^n],
\]
by the definition of matrix multiplication applied to a diagonal matrix. Since these two computations agree, we conclude that $AP = PD$, and hence, since $P$ is nonsingular,
\[
A = PDP^{-1}.
\]

Summarizing, if the eigenvectors of $A$ form a basis of $\mathbb{R}^n$, then $A$ can be decomposed into the product of the matrix $P$, whose columns are the eigenvectors of $A$, the diagonal matrix $D$, whose diagonal entries are the corresponding eigenvalues of $A$, called the diagonal form of $A$, and the inverse $P^{-1}$ of the eigenvector matrix; this decomposition is called the eigendecomposition of the matrix.

\begin{definition}[Similar Matrices]
Given two matrices $A, B \in M_n$, we say that they are similar if there exists an invertible matrix $M \in M_n$ such that $A = MBM^{-1}$, or equivalently $M^{-1}AM = B$.
\end{definition}

\begin{definition}
A matrix $A \in M_n$ is diagonalizable, over $\mathbb{R}$, if it is similar to a diagonal matrix $D \in M_n$. In this case, $D$ is called the diagonal form of $A$.
\end{definition}

The eigendecomposition constructed above shows immediately that if $\mathbb{R}^n$ has a basis of eigenvectors of $A$, then $A$ is diagonalizable, since $A = PDP^{-1}$ exhibits the required similarity with a diagonal matrix. The converse also holds, essentially by reversing the argument: if $A$ is diagonalizable, so that $A = PDP^{-1}$ for some invertible matrix $P$ and diagonal matrix $D$, then the columns of $P$ are $n$ linearly independent eigenvectors of $A$, with the corresponding diagonal entries of $D$ giving the associated eigenvalues. This yields the following characterization.

\begin{theorem}[Diagonalization Test I]
$A \in M_n$ is diagonalizable over $\mathbb{R}$ if and only if $\mathbb{R}^n$ has a basis of eigenvectors of $A$.
\end{theorem}

Combining this criterion with the earlier discussion of algebraic and geometric multiplicities yields an equivalent, and often more directly verifiable, criterion.

\begin{theorem}[Diagonalization Test II]
$A \in M_n$ is diagonalizable over $\mathbb{R}$ if and only if $A$ has exactly $n$ real eigenvalues, not necessarily distinct, and each eigenvalue is regular.
\end{theorem}

Two particular situations guarantee diagonalizability automatically, and are worth recording as sufficient conditions of frequent practical use. If $A \in M_n$ has $n$ distinct eigenvalues, that is, if its characteristic polynomial has $n$ distinct roots, then $A$ is diagonalizable, since every simple eigenvalue is automatically regular, as observed earlier. If $A \in M_n$ is symmetric, then $A$ is diagonalizable, as guaranteed directly by the Spectral Theorem.

\begin{example}
The matrix
\[
C = \begin{pmatrix} 0 & 1 \\ 0 & 0 \end{pmatrix}
\]
is not diagonalizable: it has a single eigenvalue $\lambda = 0$, of algebraic multiplicity $2$, but the corresponding eigenspace has dimension only $1$, so that this eigenvalue is not regular.
\end{example}

\begin{example}
The matrix
\[
B = \begin{pmatrix} 0 & 1 \\ -1 & 0 \end{pmatrix}
\]
is not diagonalizable over $\mathbb{R}$, since it possesses no real eigenvalues at all: its characteristic polynomial has only complex roots.
\end{example}

\begin{example}
The matrix
\[
A = \begin{pmatrix} 2 & 0 \\ 2 & 1 \end{pmatrix}
\]
considered earlier is diagonalizable, since it has two distinct eigenvalues, $\lambda = 1$ and $\lambda = 2$.
\end{example}

\begin{example}
We investigate whether
\[
A = \begin{pmatrix} -1 & 0 & 1 \\ 1 & -1 & 0 \\ 0 & 1 & 1 \end{pmatrix}
\]
is diagonalizable. Expanding the characteristic polynomial along a convenient row or column yields
\[
\begin{aligned}
P(\lambda) &= \operatorname{Det}(A - \lambda I) \\
&= (1-\lambda)[\lambda(\lambda-1) - 1] + (\lambda - 1) \\
&= (1-\lambda)(\lambda+1)(\lambda-2),
\end{aligned}
\]
so that $A$ has three distinct real eigenvalues, $\lambda_1 = -1$, $\lambda_2 = 1$, $\lambda_3 = 2$; since these are distinct, $A$ is diagonalizable. One may further note that, since $A$ is not itself symmetric, its eigenvectors need not be automatically orthogonal, though a basis of $\mathbb{R}^3$ consisting of its eigenvectors nonetheless exists by the criterion of distinct eigenvalues.
\end{example}

\begin{example}
We investigate whether
\[
A = \begin{pmatrix} 2 & 1 & 0 \\ 0 & 1 & -1 \\ 0 & 2 & 4 \end{pmatrix}
\]
is diagonalizable. A direct computation of the eigenvalues and their multiplicities yields $\lambda_1 = 2$, with algebraic multiplicity $a_{\lambda_1} = 2$ but geometric multiplicity $g_{\lambda_1} = 1$, and $\lambda_2 = 3$, with algebraic and geometric multiplicity both equal to $1$. Since $\lambda_1$ fails to be regular, being
\[
g_{\lambda_1} = 1 < 2 = a_{\lambda_1},
\]
we conclude that $A$ is not diagonalizable. Explicitly, one finds
\[
\begin{gathered}
V_{\lambda_1} = \operatorname{span}((1,0,0)) \\
V_{\lambda_2} = \operatorname{span}((1,1,-2)).
\end{gathered}
\]

\end{example}

\begin{example}
We investigate whether
\[
A = \begin{pmatrix} -3 & 5 & -1 \\ -1 & 3 & -1 \\ -3 & 3 & 1 \end{pmatrix}
\]
is diagonalizable. A direct computation gives the following two eigenvalues and their algebraic and geometric multiplicities:
\[
\begin{aligned}
\lambda_1 &= 1, & a_{\lambda_1} &= g_{\lambda_1} = 1, \\
\lambda_2 &= 2, & a_{\lambda_2} &= g_{\lambda_2} = 2.
\end{aligned}
\]
The sum of the algebraic multiplicities is
\[
\begin{aligned}
a_{\lambda_1} + a_{\lambda_2} &= 1 + 2 \\
&= 3.
\end{aligned}
\]

Both eigenvalues are regular, so we conclude that $A$ is diagonalizable. The corresponding diagonal form $D$ and change of basis matrix $P$ can be shown to be
\[
D = \begin{pmatrix} 1 & 0 & 0 \\ 0 & 2 & 0 \\ 0 & 0 & 2 \end{pmatrix}, \qquad P = \begin{pmatrix} 1 & 1 & 0 \\ 1 & 1 & -1 \\ 1 & 0 & 3 \end{pmatrix}.
\]

\end{example}

As further exercises illustrating these criteria, one may show that the matrix
\[
\begin{pmatrix} 1 & 2 \\ 0 & 1 \end{pmatrix}
\]
is not diagonalizable, since its unique eigenvalue $\lambda=1$ has algebraic multiplicity $2$ but geometric multiplicity only $1$; and one may determine the eigenvalues and eigenvectors of the matrix
\[
A = \begin{pmatrix} 2 & 1 & 0 \\ 0 & 1 & -1 \\ 0 & 2 & 4 \end{pmatrix}
\]
considered above, confirming explicitly that $\lambda_1 = 2$ has algebraic multiplicity $2$ but geometric multiplicity $1$, with eigenspace
\[
V_{\lambda_1} = \operatorname{span}((1,0,0)),
\]
and that $\lambda_2 = 3$ has algebraic and geometric multiplicity both equal to $1$, with eigenspace
\[
V_{\lambda_2} = \operatorname{span}((1,1,-2)),
\]
thereby confirming once more the non-diagonalizability of this matrix.

\subsection{Singular Values, Least Squares, and the Pseudoinverse}


The eigendecomposition studied earlier applies only to square matrices possessing a full set of linearly independent eigenvectors, and even then it treats the domain and codomain of the associated linear map as the same space. The singular value decomposition removes both restrictions: it applies to every real matrix, of any shape, and it reveals how the matrix distorts lengths and directions by relating an orthonormal basis of the domain to an orthonormal basis of the codomain.

Given $A \in M_{m,n}$, the matrix $A^{\top}A \in M_n$ is symmetric, and consequently, by the Spectral Theorem, it possesses $n$ real, nonnegative eigenvalues
\[
\lambda_1 \geq \lambda_2 \geq \cdots \geq \lambda_n \geq 0,
\]
together with an orthonormal basis of eigenvectors $\mathbf{v}_1, \dots, \mathbf{v}_n$ of $\mathbb{R}^n$. Nonnegativity follows from the identity
\[
\begin{aligned}
\mathbf{v}_i^{\top} A^{\top} A \mathbf{v}_i &= \|A\mathbf{v}_i\|^2 \\
&= \lambda_i \|\mathbf{v}_i\|^2 \\
&\geq 0.
\end{aligned}
\]
This observation motivates the following definition.

\begin{definition}[Singular Values]
Given $A \in M_{m,n}$, the singular values of $A$ are the numbers $\sigma_i = \lambda_i^{1/2}$, for $i = 1, \dots, n$, where $\lambda_1 \geq \cdots \geq \lambda_n \geq 0$ are the eigenvalues of $A^{\top}A$.
\end{definition}

The singular values, together with the associated eigenvectors of $A^{\top}A$, assemble into a factorization of $A$ that is valid without any assumption of squareness or diagonalizability.

\begin{theorem}[Singular Value Decomposition]
Given $A \in M_{m,n}$ with $\operatorname{Rank}(A) = r$, there exist an orthogonal matrix $U \in M_m$, an orthogonal matrix $V \in M_n$, and a matrix $\Sigma \in M_{m,n}$ whose only nonzero entries are $\Sigma_{ii} = \sigma_i$ for $i = 1, \dots, r$, arranged so that $\sigma_1 \geq \cdots \geq \sigma_r > 0$, such that $A = U \Sigma V^{\top}$.
\end{theorem}

\begin{proof}
Let $\mathbf{v}_1, \dots, \mathbf{v}_n$ be an orthonormal basis of eigenvectors of $A^{\top}A$, ordered so that the corresponding eigenvalues satisfy $\lambda_1 \geq \cdots \geq \lambda_n \geq 0$, and set
\[
\sigma_i = \lambda_i^{1/2}.
\]
The rank identity gives
\[
\begin{aligned}
\operatorname{Rank}(A^{\top}A) &= \operatorname{Rank}(A) \\
&= r,
\end{aligned}
\]

Therefore, exactly $r$ of the $\lambda_i$ are nonzero, say $\lambda_1, \dots, \lambda_r$. For $i = 1, \dots, r$, define $\mathbf{u}_i = \sigma_i^{-1} A \mathbf{v}_i$; these vectors are orthonormal, since
\[
\begin{aligned}
\mathbf{u}_i \cdot \mathbf{u}_j &= (\sigma_i \sigma_j)^{-1} \mathbf{v}_i^{\top} A^{\top} A \mathbf{v}_j \\
&= (\sigma_i \sigma_j)^{-1} \lambda_j \mathbf{v}_i \cdot \mathbf{v}_j,
\end{aligned}
\]
which equals $1$ if $i=j$ and $0$ otherwise by orthonormality of the $\mathbf{v}_i$. Complete $\mathbf{u}_1, \dots, \mathbf{u}_r$ to an orthonormal basis $\mathbf{u}_1, \dots, \mathbf{u}_m$ of $\mathbb{R}^m$, and let $U = [\mathbf{u}_1 | \cdots | \mathbf{u}_m]$ and $V = [\mathbf{v}_1 | \cdots | \mathbf{v}_n]$, both orthogonal by construction, since their columns form orthonormal bases. For $i \leq r$, one has $A\mathbf{v}_i = \sigma_i \mathbf{u}_i$ by definition of $\mathbf{u}_i$; for $i > r$, one has $\lambda_i = 0$, so
\[
\begin{aligned}
\|A\mathbf{v}_i\|^2 &= \mathbf{v}_i^{\top}A^{\top}A\mathbf{v}_i \\
&= \lambda_i \|\mathbf{v}_i\|^2 \\
&= 0,
\end{aligned}
\]
forcing $A\mathbf{v}_i = \vec{0}$. Consequently $AV = U\Sigma$, where $\Sigma$ has entries $\sigma_i$ on its first $r$ diagonal positions and zero elsewhere, and since $V$ is orthogonal, hence invertible with $V^{-1} = V^{\top}$, one obtains
\[
\begin{aligned}
A &= AVV^{\top} \\
&= U\Sigma V^{\top}.
\end{aligned}
\]

\end{proof}

The rank of $A$ therefore equals the number of nonzero singular values, and the decomposition exposes the four fundamental subspaces of $A$ directly: the columns $\mathbf{u}_1, \dots, \mathbf{u}_r$ of $U$ form a basis of $\operatorname{Col}(A)$, the columns $\mathbf{u}_{r+1}, \dots, \mathbf{u}_m$ form a basis of $\operatorname{Ker}(A^{\top})$, the columns $\mathbf{v}_1, \dots, \mathbf{v}_r$ form a basis of $\operatorname{Col}(A^{\top})$, and the columns $\mathbf{v}_{r+1}, \dots, \mathbf{v}_n$ form a basis of $\operatorname{Ker}(A)$. This last fact deserves justification: since $A\mathbf{v}_i = \vec{0}$ for $i > r$, the span of $\mathbf{v}_{r+1}, \dots, \mathbf{v}_n$ is certainly contained in $\operatorname{Ker}(A)$; and since these $n-r$ vectors are linearly independent, while the Nullity--Rank Theorem gives
\[
\begin{aligned}
\operatorname{Dim}(\operatorname{Ker}(A)) &= n - \operatorname{Rank}(A) \\
&= n-r,
\end{aligned}
\]
the inclusion is in fact an equality. The uniqueness properties of the decomposition are worth recording explicitly. The singular values $\sigma_1 \geq \cdots \geq \sigma_r$ are always uniquely determined by $A$, being the square roots of the nonzero eigenvalues of $A^{\top}A$, themselves an intrinsic invariant of $A$; the vectors $\mathbf{u}_i$ and $\mathbf{v}_i$, however, are unique only up to simultaneous rotation within eigenspaces of repeated singular values, since any orthonormal basis of such an eigenspace serves equally well in the construction above.

\begin{example}
Consider the matrix
\[
A = \begin{pmatrix} 2 & 0 \\ 0 & -3 \end{pmatrix}.
\]
The following product is already diagonal:
\[
A^{\top}A = \begin{pmatrix} 4 & 0 \\ 0 & 9 \end{pmatrix}
\]

Its eigenvalues are $\lambda_1 = 9$, $\lambda_2 = 4$, with eigenvectors $\mathbf{v}_1 = (0,1)$, $\mathbf{v}_2 = (1,0)$, giving singular values $\sigma_1 = 3$, $\sigma_2 = 2$. Computing
\[
\begin{gathered}
\mathbf{u}_1 = \sigma_1^{-1}A\mathbf{v}_1 = (0,-1) \\
\mathbf{u}_2 = \sigma_2^{-1}A\mathbf{v}_2 = (1,0),
\end{gathered}
\]
one obtains the decomposition $A = U\Sigma V^{\top}$ with
\[
U = \begin{pmatrix} 0 & 1 \\ -1 & 0 \end{pmatrix},
\qquad
\Sigma = \begin{pmatrix} 3 & 0 \\ 0 & 2 \end{pmatrix},
\qquad
V = \begin{pmatrix} 0 & 1 \\ 1 & 0 \end{pmatrix}.
\]
As a verification,
\[
\begin{aligned}
U\Sigma V^{\top}
&= \begin{pmatrix} 0 & 1 \\ -1 & 0 \end{pmatrix}\begin{pmatrix} 3 & 0 \\ 0 & 2 \end{pmatrix}\begin{pmatrix} 0 & 1 \\ 1 & 0 \end{pmatrix} \\
&= \begin{pmatrix} 0 & 2 \\ -3 & 0 \end{pmatrix}\begin{pmatrix} 0 & 1 \\ 1 & 0 \end{pmatrix} \\
&= \begin{pmatrix} 2 & 0 \\ 0 & -3 \end{pmatrix} = A,
\end{aligned}
\]
as required.
\end{example}

Beyond its structural role, the singular value decomposition provides the optimal low-rank approximation of a matrix in a precise sense, a fact that underlies its widespread use in data compression and dimensionality reduction. To state the result, one measures the size of a matrix using the operator norm,
\[
\|B\| = \max_{\|\mathbf{x}\|=1} \|B\mathbf{x}\|,
\]
and one first records an elementary but essential fact.

\begin{theorem}
If $A \in M_{m,n}$ has singular values
\[
\sigma_1 \geq \cdots \geq \sigma_r > 0,
\]
then $\|A\| = \sigma_1$.
\end{theorem}

\begin{proof}
Write $A = U\Sigma V^{\top}$ as above. Since $U$ and $V$ are orthogonal, they preserve Euclidean norms, so
\[
\begin{aligned}
\|A\mathbf{x}\| &= \|U\Sigma V^{\top}\mathbf{x}\| \\
&= \|\Sigma V^{\top}\mathbf{x}\|
\end{aligned}
\]
for every $\mathbf{x}$, and as $\mathbf{x}$ ranges over the unit sphere, so does $\mathbf{y} = V^{\top}\mathbf{x}$, by orthogonality of $V$. It therefore suffices to compute
\[
\max_{\|\mathbf{y}\|=1} \|\Sigma \mathbf{y}\|.
\]
Multiplication by $\Sigma$ gives
\[
\Sigma \mathbf{y} = (\sigma_1 y_1, \dots, \sigma_r y_r, 0, \dots, 0),
\]
The squared norm therefore satisfies
\[
\begin{aligned}
\|\Sigma\mathbf{y}\|^2 &= \sum_{i=1}^{r} \sigma_i^2 y_i^2 \\
&\leq \sigma_1^2 \sum_{i=1}^{r} y_i^2 \\
&\leq \sigma_1^2,
\end{aligned}
\]
with equality when $\mathbf{y} = \mathbf{e}_1$, so the maximum equals $\sigma_1$.
\end{proof}

\begin{theorem}[Eckart--Young Theorem]
Let $A = U\Sigma V^{\top}$ be the singular value decomposition of $A \in M_{m,n}$, with singular values $\sigma_1 \geq \cdots \geq \sigma_r > 0$, and for $k < r$, let
\[
A_k = \sum_{i=1}^{k} \sigma_i \mathbf{u}_i \mathbf{v}_i^{\top}.
\]

Then $\operatorname{Rank}(A_k) \leq k$, and $A_k$ minimizes $\|A - B\|$ among all $B \in M_{m,n}$ of rank at most $k$; moreover $\|A - A_k\| = \sigma_{k+1}$.
\end{theorem}

\begin{proof}
Since $A_k$ is a sum of $k$ rank-one matrices $\mathbf{u}_i\mathbf{v}_i^{\top}$, $\operatorname{Rank}(A_k) \leq k$. Writing
\[
A - A_k = \sum_{i=k+1}^{r} \sigma_i \mathbf{u}_i \mathbf{v}_i^{\top},
\]
this is itself the singular value decomposition of $A-A_k$, restricted to indices $k+1, \dots, r$, so by the previous theorem $\|A - A_k\| = \sigma_{k+1}$.

It remains to show that no matrix of rank at most $k$ does better. Let $B \in M_{m,n}$ with $\operatorname{Rank}(B) \leq k$, so that
\[
\operatorname{Dim}(\operatorname{Ker}(B)) \geq n-k
\]
by the Nullity--Rank Theorem. Consider the subspace
\[
W = \operatorname{span}(\mathbf{v}_1, \dots, \mathbf{v}_{k+1}),
\]
of dimension $k+1$; since
\[
\operatorname{Dim}(\operatorname{Ker}(B)) + \operatorname{Dim}(W) \geq (n-k) + (k+1) = n+1 > n,
\]
the two subspaces must intersect nontrivially, so there exists a unit vector $\mathbf{w} \in W \cap \operatorname{Ker}(B)$. Expand this vector in the basis of $W$:
\[
\mathbf{w} = \sum_{i=1}^{k+1} c_i \mathbf{v}_i.
\]
The unit-norm constraint is
\[
\sum_i c_i^2 = 1.
\]
Using orthonormality of the $\mathbf{v}_i$ and $\mathbf{u}_i$, we obtain
\[
\begin{aligned}
\|(A-B)\mathbf{w}\|^2
&= \|A\mathbf{w}\|^2
= \left\| \sum_{i=1}^{k+1} c_i \sigma_i \mathbf{u}_i \right\|^2 \\
&= \sum_{i=1}^{k+1} c_i^2 \sigma_i^2
\geq \sigma_{k+1}^2 \sum_{i=1}^{k+1} c_i^2
= \sigma_{k+1}^2,
\end{aligned}
\]
where the inequality uses $\sigma_i \geq \sigma_{k+1}$ for $i \leq k+1$. Since $\mathbf{w}$ is a unit vector, this shows
\[
\begin{aligned}
\|A-B\| &\geq \sigma_{k+1} \\
&= \|A-A_k\|,
\end{aligned}
\]
establishing the optimality of $A_k$.
\end{proof}

Figure~\ref{fig:ch10-svd} shows the geometric action of the singular value decomposition on the unit circle.

\begin{figure}[!htbp]
\centering
\begin{tikzpicture}[scale=1.5,>=Stealth]
  \begin{scope}
    \draw[->] (-1.6,0) -- (1.6,0);
    \draw[->] (0,-1.6) -- (0,1.6);
    \draw[figblue,thick] (0,0) circle (1);
    \draw[vecB] (0,0) -- (30:1) node[above right,font=\scriptsize]{$\bv_1$};
    \draw[vecC] (0,0) -- (120:1) node[above left,font=\scriptsize]{$\bv_2$};
  \end{scope}
  \draw[->,very thick] (2.0,0.7) -- (3.1,0.7) node[midway,above,font=\scriptsize]{$A=U\Sigma V^\top$};
  \begin{scope}[xshift=5.6cm]
    \draw[->] (-1.9,0) -- (1.9,0);
    \draw[->] (0,-1.3) -- (0,1.3);
    \draw[figblue,thick,rotate=15] (0,0) ellipse (1.6 and 0.7);
    \draw[vecB] (0,0) -- (15:1.6) node[right,font=\scriptsize]{$\sigma_1 \bu_1$};
    \draw[vecC] (0,0) -- (105:0.7) node[above,font=\scriptsize]{$\sigma_2 \bu_2$};
  \end{scope}
\end{tikzpicture}
\caption[Singular Value Decomposition (SVD)]{Left: the unit circle and the orthonormal domain basis $\{\bv_1,\bv_2\}$. Right: the image ellipse, with orthogonal semi-axis vectors $\sigma_1\bu_1$ and $\sigma_2\bu_2$. The map $A=U\Sigma V^\top$ rotates or reflects, scales along perpendicular directions, and rotates or reflects again; the semi-axis lengths are $\sigma_1,\sigma_2$.}
\label{fig:ch10-svd}
\end{figure}

The inverse of a matrix, as defined earlier, exists only for nonsingular square matrices, yet many of the systems arising in applications are either rectangular or singular, and one still wishes to associate to them a canonical generalized inverse. The singular value decomposition supplies exactly such a generalization.

\begin{definition}[Pseudoinverse]
Given $A \in M_{m,n}$ with singular value decomposition $A = U\Sigma V^{\top}$, the Moore--Penrose pseudoinverse of $A$ is the matrix
\[
A^{+} = V \Sigma^{+} U^{\top} \in M_{n,m},
\]
where $\Sigma^{+} \in M_{n,m}$ is obtained from $\Sigma^{\top}$ by replacing every nonzero diagonal entry $\sigma_i$ with $\sigma_i^{-1}$.
\end{definition}

When $A$ is square and nonsingular, every singular value is nonzero and $\Sigma$ is itself square and invertible, with $\Sigma^{+} = \Sigma^{-1}$, so that
\[
\begin{aligned}
A^{+} &= V\Sigma^{-1}U^{\top} \\
&= (U\Sigma V^{\top})^{-1} \\
&= A^{-1};
\end{aligned}
\]
the pseudoinverse indeed extends the ordinary inverse. Its defining property, independent of the particular decomposition used to construct it, is captured by the following characterization, which we now prove directly from the definition above.

\begin{theorem}[Penrose Conditions]
The pseudoinverse $A^{+}$ satisfies the four conditions
\[
AA^{+}A = A, \qquad A^{+}AA^{+} = A^{+}, \qquad (AA^{+})^{\top} = AA^{+}, \qquad (A^{+}A)^{\top} = A^{+}A,
\]
and it is the unique matrix in $M_{n,m}$ satisfying all four simultaneously.
\end{theorem}

\begin{proof}
Write $A = U\Sigma V^{\top}$ and $A^{+} = V\Sigma^{+}U^{\top}$. A direct computation, using $U^{\top}U = I_m$ and $V^{\top}V = I_n$, gives
\[
\begin{aligned}
AA^{+} &= U\Sigma V^{\top}V\Sigma^{+}U^{\top} \\
&= U(\Sigma \Sigma^{+})U^{\top};
\end{aligned}
\]
since $\Sigma\Sigma^{+} \in M_m$ is diagonal with entries $1$ in the first $r$ positions and $0$ elsewhere, it is symmetric, and hence $AA^{+} = U(\Sigma\Sigma^{+})U^{\top}$ is symmetric as a product of a symmetric matrix conjugated by an orthogonal one, establishing $(AA^{+})^{\top} = AA^{+}$. The identical argument, with the roles of $U$ and $V$ exchanged, gives $A^{+}A = V(\Sigma^{+}\Sigma)V^{\top}$ symmetric, establishing $(A^{+}A)^{\top} = A^{+}A$. For the first two conditions, one checks directly that $\Sigma\Sigma^{+}\Sigma = \Sigma$ and $\Sigma^{+}\Sigma\Sigma^{+} = \Sigma^{+}$, since both sides act identically on the first $r$ standard basis vectors, where $\Sigma$ and $\Sigma^{+}$ are mutually inverse scalars, and both sides vanish on the remaining basis vectors, where $\Sigma$ has zero rows or columns; substituting these identities into
\[
\begin{aligned}
AA^{+}A &= U\Sigma V^{\top}V\Sigma^{+}U^{\top}U\Sigma V^{\top} \\
&= U(\Sigma\Sigma^{+}\Sigma)V^{\top} \\
&= U\Sigma V^{\top} \\
&= A,
\end{aligned}
\]
and analogously for $A^{+}AA^{+} = A^{+}$. For uniqueness, suppose $B \in M_{n,m}$ also satisfies all four conditions. One shows $B = A^{+}$ by a sequence of substitutions exploiting the four identities for both $A^{+}$ and $B$; we omit the lengthy but entirely mechanical verification, which proceeds by repeatedly inserting
\[
\begin{aligned}
AA^{+}A &= A \\
&= ABA
\end{aligned}
\]
and the symmetry conditions into the expression for $B$ until it is transformed into $A^{+}$.
\end{proof}

The practical importance of the pseudoinverse stems from its role in solving linear systems that admit no exact solution, or that admit infinitely many.

\begin{theorem}[Least Squares via the Pseudoinverse]
Given $A \in M_{m,n}$ and $\mathbf{b} \in \mathbb{R}^m$, the vector $\mathbf{x}^{*} = A^{+}\mathbf{b}$ minimizes $\|A\mathbf{x} - \mathbf{b}\|$ over all $\mathbf{x} \in \mathbb{R}^n$, and among all such minimizers it is the one of smallest Euclidean norm.
\end{theorem}

\begin{proof}
Write $A = U\Sigma V^{\top}$, and set $\mathbf{y} = V^{\top}\mathbf{x}$ and $\mathbf{c} = U^{\top}\mathbf{b}$; since $U$ and $V$ are orthogonal,
\[
\begin{aligned}
\|A\mathbf{x}-\mathbf{b}\| &= \|U\Sigma V^{\top}\mathbf{x} - \mathbf{b}\| \\
&= \|\Sigma \mathbf{y} - \mathbf{c}\|.
\end{aligned}
\]
Moreover, $\mathbf{x}$ ranges over $\mathbb{R}^n$ exactly as $\mathbf{y}$ does. Multiplication by $\Sigma$ gives
\[
\Sigma \mathbf{y} = (\sigma_1 y_1, \dots, \sigma_r y_r, 0, \dots, 0).
\]
The squared residual is therefore
\[
\|\Sigma\mathbf{y} - \mathbf{c}\|^2 = \sum_{i=1}^r (\sigma_i y_i - c_i)^2 + \sum_{i=r+1}^{m} c_i^2.
\]

Over all choices of $y_1, \dots, y_n$, it is minimized precisely by taking $y_i = c_i/\sigma_i$ for $i \leq r$, which eliminates the first sum entirely, while $y_{r+1}, \dots, y_n$ remain unconstrained by the minimization; the second sum is independent of $\mathbf{y}$ and represents the unavoidable residual. Among all minimizers, the one of smallest norm is obtained by setting the free components $y_{r+1}, \dots, y_n$ to zero, since $\|\mathbf{y}\|^2 = \|\mathbf{x}\|^2$ by orthogonality of $V$, and this choice corresponds precisely to
\[
\begin{aligned}
\mathbf{y} &= \Sigma^{+}\mathbf{c} \\
&= \Sigma^{+}U^{\top}\mathbf{b}.
\end{aligned}
\]
Thus, the minimizer has the equivalent representations
\[
\begin{aligned}
\mathbf{x}^{*} &= V\mathbf{y} \\
&= V\Sigma^{+}U^{\top}\mathbf{b} \\
&= A^{+}\mathbf{b}.
\end{aligned}
\]

\end{proof}

\begin{example}
Consider the matrix
\[
A = \begin{pmatrix} 1 & 0 \\ 0 & 0 \end{pmatrix}.
\]
Its singular values are $\sigma_1 = 1$, $\sigma_2 = 0$, with
\[
\begin{aligned}
U &= V \\
&= I,
\end{aligned}
\]
so
\[
\Sigma^{+} = \begin{pmatrix} 1 & 0 \\ 0 & 0 \end{pmatrix}
\]
and hence $A^{+} = A$. For $\mathbf{b} = (1,1)$, the vector
\[
\mathbf{x}^{*} = A^{+}\mathbf{b} = (1,0)
\]
is the minimum-norm least squares solution of $A\mathbf{x} = \mathbf{b}$: indeed $A\mathbf{x} = (x_1, 0)$ for every $\mathbf{x} = (x_1,x_2)$, so
\[
\|A\mathbf{x}-\mathbf{b}\|^2 = (x_1-1)^2 + 1
\]
is minimized precisely when $x_1=1$, regardless of $x_2$, and among these infinitely many minimizers the one of least norm is $(1,0)$, obtained by additionally setting $x_2=0$.
\end{example}

Figure~\ref{fig:atlas-least-squares-projection} shows least squares as orthogonal projection onto the column space.

\begin{figure}[!htbp]
\centering
\begin{tikzpicture}[scale=1.45,>=Stealth]
\draw[black!45] (-0.6,-0.3) -- (4.1,2.05) node[right] {$\operatorname{Col}A$};
\draw[vecA] (0,0) -- (1,3) node[above] {$b$};
\draw[vecB] (0,0) -- (2,1) node[below right] {$A\hat{x}$};
\draw[vecC] (2,1) -- (1,3) node[midway,right] {$r$};
\draw[black!60] (2.25,1.125) -- (2.125,1.375) -- (1.875,1.25);
\node[below] at (0,0) {$0$};
\end{tikzpicture}
\caption{The vector $b$ is decomposed into its projection $A\hat{x}$ onto the column space and a perpendicular residual $r=b-A\hat{x}$. Orthogonality gives the normal equations $A^{\top}r=0$. Here the column space is represented by a line in the plane.}
\label{fig:atlas-least-squares-projection}
\end{figure}

\newpage
\chapter{Applied Linear Algebra}

\section{Matrices, Products, and Linear Systems}

\subsection{Matrices and Basic Operations}


A matrix is a rectangular array of $m \cdot n$ numbers organized into $m$ horizontal rows and $n$ vertical columns,
\[
A =
\begin{pmatrix}
a_{11} & a_{12} & \cdots & a_{1n} \\
a_{21} & a_{22} & \cdots & a_{2n} \\
\vdots & \vdots & \ddots & \vdots \\
a_{m1} & a_{m2} & \cdots & a_{mn}
\end{pmatrix}.
\]

Each element, or entry, of the matrix is denoted by a variable carrying two subscripts: the symbol $a_{ij}$ represents the entry located at the $i$-th row and $j$-th column, where the row index $i$ ranges from $1$ to $m$ and the column index $j$ ranges from $1$ to $n$. Each row of the matrix may itself be regarded as a vector with $n$ components, $\mathbf{a}_i = [a_{i1} \ a_{i2} \ \cdots \ a_{in}]$ for $i = 1, \dots, m$, and each column may likewise be regarded as a vector with $m$ components, $\mathbf{a}^j = (a_{1j}, a_{2j}, \dots, a_{mj})$ for $j = 1, \dots, n$. A matrix $A$ as above is denoted compactly by $A = [a_{ij}]$. The set of all matrices with real entries of dimension $m \times n$ is denoted by $M_{m,n}$, or simply by $M_n$ in the special case $m = n$, in which case the matrices are called square matrices. The matrix all of whose entries equal zero is called the zero matrix, and is denoted by $O \in M_{m,n}$. Finally, given $n$ vectors $\mathbf{b}^1, \dots, \mathbf{b}^n \in \mathbb{R}^m$, the notation
\[
[\mathbf{b}^1 | \cdots | \mathbf{b}^n] \in M_{m,n}
\]
denotes the matrix obtained by placing these vectors into the columns of the matrix, in the given order.


Certain matrix configurations recur frequently enough throughout the theory to deserve dedicated terminology. A matrix is said to be in row echelon form when the leading coefficient of every nonzero row, that is, the first nonzero entry counted from the left, also called the pivot of that row, lies strictly to the right of the leading coefficient of the row directly above it, and all rows consisting entirely of zeros, if any, are collected at the bottom of the matrix. These two conditions together imply that every entry lying in a column below a leading coefficient must itself be zero, so that the lower left-hand corner of a matrix in row echelon form is filled with zeros.

Among matrices with an equal number of rows and columns, several further special types are worth distinguishing. A square matrix is a matrix belonging to $M_n$ for some $n$, that is, one having the same number of rows and columns. A diagonal matrix is a square matrix all of whose off-diagonal entries vanish; the particular diagonal matrix having every diagonal entry equal to $1$ is called the identity matrix, and is denoted by $I$. An upper triangular matrix is a square matrix all of whose entries below the main diagonal vanish, while a lower triangular matrix is a square matrix all of whose entries above the main diagonal vanish.


The set of matrices of a fixed dimension can be equipped with two natural algebraic operations, entirely analogous to the corresponding operations on vectors of $\mathbb{R}^n$.

\begin{definition}[Matrix Addition]
Given two matrices of the same dimension, $A = [a_{ij}]$ and $B = [b_{ij}]$, both belonging to $M_{m,n}$, their sum $A+B$ is the $m \times n$ matrix whose $(i,j)$ entry is the sum of the corresponding entries of $A$ and $B$, that is, $A + B = [a_{ij} + b_{ij}]$, for every $i = 1, \dots, m$ and $j = 1, \dots, n$.
\end{definition}

\begin{definition}[Scalar Multiplication]
Given a matrix $A = [a_{ij}] \in M_{m,n}$ and a real number $t \in \mathbb{R}$, the scalar product $tA$ is the $m \times n$ matrix obtained by multiplying every entry of $A$ by $t$, that is, $tA = [t \, a_{ij}]$.
\end{definition}

These two operations satisfy a collection of algebraic properties entirely analogous to those already familiar from ordinary arithmetic on real numbers. Given any three matrices $A, B, C \in M_{m,n}$ and any two real numbers $s, t$, one has that addition is associative, $A + (B+C) = (A+B) + C$, and commutative, $A + B = B + A$; that the zero matrix $O \in M_{m,n}$ satisfies $A + O = A$; that the matrix $-A = [-a_{ij}]$ is the additive inverse of $A$, satisfying $A + (-A) = O$; and that scalar multiplication satisfies $s(tA) = (st)A$, $1 \cdot A = A$, $t(A+B) = tA + tB$, and $(s+t)A = sA + tA$. Taken together, these properties show that $M_{m,n}$, equipped with the two operations of matrix addition and scalar multiplication, forms a vector space of dimension $m \cdot n$.

\subsection{Matrix Products and Linear Systems}

Before defining the product of two general matrices, it is instructive to consider the simpler case in which one of the two factors is a vector. Given a row vector $\mathbf{a} = [a_1 \ a_2 \ \cdots \ a_n]$ and a column vector $\mathbf{x} = (x_1, x_2, \dots, x_n)$ having the same number $n$ of components, the product $\mathbf{a}\mathbf{x}$ is defined as the ordinary dot product in $\mathbb{R}^n$,
\[
\mathbf{a}\mathbf{x} = [a_1 \ a_2 \ \cdots \ a_n]
\begin{pmatrix} x_1 \\ x_2 \\ \vdots \\ x_n \end{pmatrix}
= a_1 x_1 + a_2 x_2 + \cdots + a_n x_n.
\]
This notion extends naturally to the product of a general matrix with a column vector.

\begin{definition}[Matrix-Vector Product]
Given a matrix $A = [a_{ij}] \in M_{m,n}$ and a column vector $\mathbf{x} = (x_1, x_2, \dots, x_n) \in \mathbb{R}^n$, the product $A\mathbf{x}$ is the vector with $m$ components given by
\[
A\mathbf{x} =
\begin{pmatrix} \mathbf{a}_1 \mathbf{x} \\ \mathbf{a}_2 \mathbf{x} \\ \vdots \\ \mathbf{a}_m \mathbf{x} \end{pmatrix}
=
\begin{pmatrix}
a_{11} x_1 + a_{12} x_2 + \cdots + a_{1n} x_n \\
a_{21} x_1 + a_{22} x_2 + \cdots + a_{2n} x_n \\
\vdots \\
a_{m1} x_1 + a_{m2} x_2 + \cdots + a_{mn} x_n
\end{pmatrix} \in \mathbb{R}^m,
\]
Here the $i$-th row of $A$ is
\[
\mathbf{a}_i = [a_{i1} \ a_{i2} \ \cdots \ a_{in}]
\]

\end{definition}

\begin{example}
Consider the following matrix and vector:
\[
\begin{gathered}
A = \begin{pmatrix} 0 & 1 & 1 \\ 2 & 2 & 1 \end{pmatrix} \in M_{2,3} \\
\mathbf{x} = (2,1,3) \in \mathbb{R}^3
\end{gathered}
\]
Their product is
\[
A\mathbf{x} = \begin{pmatrix} 0\cdot 2 + 1 \cdot 1 + 1 \cdot 3 \\ 2 \cdot 2 + 2 \cdot 1 + 1 \cdot 3 \end{pmatrix} = \begin{pmatrix} 4 \\ 9 \end{pmatrix} \in \mathbb{R}^2.
\]

\end{example}

The matrix-vector product enjoys an important algebraic property, namely linearity with respect to the vector argument, which lies at the very heart of the entire theory of linear algebra.

\begin{theorem}[Linearity of the Matrix-Vector Product]
Given $A \in M_{m,n}$, the matrix-vector product $A\mathbf{x}$ is linear in $\mathbf{x}$; that is, $A(\mathbf{x}+\mathbf{y}) = A\mathbf{x} + A\mathbf{y}$ for every $\mathbf{x}, \mathbf{y} \in \mathbb{R}^n$, and $A(t\mathbf{x}) = t(A\mathbf{x})$ for every $t \in \mathbb{R}$ and $\mathbf{x} \in \mathbb{R}^n$.
\end{theorem}

Given a matrix $A \in M_{m,n}$, one may associate with it the function $L : \mathbb{R}^n \to \mathbb{R}^m$ defined by $L(\mathbf{x}) = A\mathbf{x}$. The linearity property just stated shows precisely that $L$ is a linear map, in the sense that
\[
\begin{gathered}
L(\mathbf{x} + \mathbf{y}) = L(\mathbf{x}) + L(\mathbf{y}) \\
L(t\mathbf{x}) = t L(\mathbf{x}).
\end{gathered}
\]
This holds for every $\mathbf{x}, \mathbf{y} \in \mathbb{R}^n$ and $t \in \mathbb{R}$; this correspondence between matrices and linear maps, easily verified by direct computation, will be developed systematically in a later section.

A particularly useful way of viewing the matrix-vector product emerges from rewriting the explicit expression for $A\mathbf{x}$ given above. Since
\[
A\mathbf{x} =
x_1 \begin{pmatrix} a_{11} \\ a_{21} \\ \vdots \\ a_{m1} \end{pmatrix}
+ x_2 \begin{pmatrix} a_{12} \\ a_{22} \\ \vdots \\ a_{m2} \end{pmatrix}
+ \cdots +
x_n \begin{pmatrix} a_{1n} \\ a_{2n} \\ \vdots \\ a_{mn} \end{pmatrix}.
\]
One recognizes that the product $A\mathbf{x}$ is precisely the linear combination of the columns $\mathbf{a}^j$ of $A$, with coefficients given by the components of $\mathbf{x}$, that is,
\[
A\mathbf{x} = x_1 \mathbf{a}^1 + x_2 \mathbf{a}^2 + \cdots + x_n \mathbf{a}^n.
\]

This alternative viewpoint proves indispensable when relating the matrix-vector product to the theory of linear systems, developed in the following subsection, and to the theory of linear independence, developed subsequently.

The general form of a linear system of $m$ equations in $n$ unknowns $x_1, x_2, \dots, x_n$ is
\[
\begin{cases}
a_{11} x_1 + a_{12} x_2 + \cdots + a_{1n} x_n = b_1, \\
a_{21} x_1 + a_{22} x_2 + \cdots + a_{2n} x_n = b_2, \\
\quad \vdots \\
a_{m1} x_1 + a_{m2} x_2 + \cdots + a_{mn} x_n = b_m
\end{cases}.
\]
Here the coefficients $a_{ij}$ and the constants $b_i$ are assigned numbers. In light of the definition of the matrix-vector product, this system can be written compactly as the single vector equation $A\mathbf{x} = \mathbf{b}$ in $\mathbb{R}^m$, having set
\[
A = \begin{pmatrix} a_{11} & \cdots & a_{1n} \\ \vdots & \ddots & \vdots \\ a_{m1} & \cdots & a_{mn} \end{pmatrix}, \qquad
\mathbf{b} = \begin{pmatrix} b_1 \\ \vdots \\ b_m \end{pmatrix}, \qquad
\mathbf{x} = \begin{pmatrix} x_1 \\ \vdots \\ x_n \end{pmatrix}.
\]

The matrix $A$ is called the coefficient matrix of the system, and the matrix obtained by appending the column $\mathbf{b}$ to $A$, denoted $[A | \mathbf{b}]$, is called the augmented matrix of the system.

A solution of the linear system is an assignment of values to the variables $x_1, \dots, x_n$ satisfying every equation simultaneously; the set of all such solutions is called the solution set,
\[
S = \{ \mathbf{x} = (x_1, \dots, x_n) \in \mathbb{R}^n : A\mathbf{x} = \mathbf{b} \}.
\]

A linear system may exhibit one of exactly three possible behaviors: it may have infinitely many solutions, it may have a single unique solution, or it may have no solution at all. A system possessing no solution is said to be inconsistent, while a system possessing at least one solution is said to be consistent, or solvable.

\begin{example}[Application: Linear Independence]
Given $n$ vectors $\mathbf{v}^1, \dots, \mathbf{v}^n \in \mathbb{R}^m$, one may efficiently compute their linear combinations by first assembling the matrix
\[
A = [\mathbf{v}^1 | \cdots | \mathbf{v}^n] \in M_{m,n}.
\]
The coefficients can be collected into a vector
\[
\mathbf{x} = (x_1, \dots, x_n) \in \mathbb{R}^n.
\]

For every such coefficient vector, the linear combination of the vectors $\mathbf{v}^i$ with coefficients $x_i$ is given precisely by $A\mathbf{x}$. Consequently, the vectors $\mathbf{v}^1, \dots, \mathbf{v}^n$ are linearly independent if and only if the homogeneous system $A\mathbf{x} = \vec{0}_m$ has the unique solution $\mathbf{x} = \vec{0}_n$, reducing the study of linear independence to the study of a homogeneous linear system.
\end{example}

The product of two matrices is defined only when the number of columns of the first matrix equals the number of rows of the second, and the resulting product generalizes the matrix-vector product introduced above.

\begin{definition}[Matrix Multiplication]
Given $A = [a_{ij}] \in M_{m,n}$ and $B = [b_{ij}] \in M_{n,r}$, the product of $A$ times $B$, denoted $AB$, is the $m \times r$ matrix obtained by taking, for each entry, the dot product of the corresponding row of $A$ with the corresponding column of $B$; that is, $AB \in M_{m,r}$, with $[AB]_{ij} = \mathbf{a}_i \cdot \mathbf{b}^j$, so that the $(i,j)$ entry of $AB$ is the scalar product of the $i$-th row of $A$ with the $j$-th column of $B$.
\end{definition}

\begin{example}
Consider the following pair of matrices:
\[
\begin{gathered}
A = \begin{pmatrix} 0 & 1 & 1 \\ 2 & 2 & 1 \end{pmatrix} \in M_{2,3} \\
B = \begin{pmatrix} 1 & 2 \\ 2 & 0 \\ 3 & 1 \end{pmatrix} \in M_{3,2}.
\end{gathered}
\]
The product $AB$ is the $2\times2$ matrix whose entries are row--column
dot products:
\[
\begin{aligned}
c_{11}&=0\cdot1+1\cdot2+1\cdot3=5,\\
c_{12}&=0\cdot2+1\cdot0+1\cdot1=1,\\
c_{21}&=2\cdot1+2\cdot2+1\cdot3=9,\\
c_{22}&=2\cdot2+2\cdot0+1\cdot1=5.
\end{aligned}
\]
Therefore
\[
AB=\begin{pmatrix}5&1\\9&5\end{pmatrix}.
\]

\end{example}

Matrix multiplication satisfies several properties analogous to those of ordinary multiplication of numbers, namely associativity, $A(BC) = (AB)C$ for every $A \in M_{m,r}$, $B \in M_{r,n}$, $C \in M_{n,p}$; distributivity over addition, $A(B+C) = AB + AC$ for every $A \in M_{m,r}$ and $B, C \in M_{r,n}$, and $(A+B)C = AC + BC$ for every $A, B \in M_{m,r}$ and $C \in M_{r,n}$; and compatibility with scalar multiplication,
\[
\begin{aligned}
A(tB) &= (tA)B \\
&= t(AB).
\end{aligned}
\]
This holds for every $t \in \mathbb{R}$, $A \in M_{m,r}$, $B \in M_{r,n}$.

\begin{remark}
In marked contrast with the multiplication of ordinary numbers, which is independent of the order of the factors, matrix multiplication is not commutative: even for square matrices, it may happen that $AB \neq BA$. For example, with
\[
\begin{gathered}
\begin{pmatrix} 0 & 1 \\ 0 & 0 \end{pmatrix} \quad \text{and} \quad 
\begin{pmatrix} 0 & 0 \\ 1 & 0 \end{pmatrix}.
\end{gathered}
\]
One computes
\[
\begin{pmatrix} 0 & 1 \\ 0 & 0 \end{pmatrix}\begin{pmatrix} 0 & 0 \\ 1 & 0 \end{pmatrix} = \begin{pmatrix} 1 & 0 \\ 0 & 0 \end{pmatrix}.
\]
Meanwhile,
\[
\begin{pmatrix} 0 & 0 \\ 1 & 0 \end{pmatrix}\begin{pmatrix} 0 & 1 \\ 0 & 0 \end{pmatrix} = \begin{pmatrix} 0 & 0 \\ 0 & 1 \end{pmatrix}.
\]
These two products clearly differ. Moreover, the product of two matrices can vanish identically even when neither factor is the zero matrix: taking
\[
A = B = \begin{pmatrix} 0 & 1 \\ 0 & 0 \end{pmatrix}.
\]
One finds $AB = O$, despite $A \neq O$.
\end{remark}


A further operation on matrices, of frequent use throughout the theory, consists of interchanging rows and columns.

\begin{definition}[Transpose]
The transpose of an $m \times n$ matrix $A$ is the $n \times m$ matrix $A^{\top}$ obtained by turning rows into columns and vice versa, that is, $[A^{\top}]_{ij} = a_{ji}$, for every $i = 1, \dots, n$ and $j = 1, \dots, m$.
\end{definition}
\subsection{Solving Linear Systems by Gaussian Elimination}

When the coefficient matrix of a homogeneous linear system is already in row echelon form, the corresponding system can be solved directly by a simple procedure known as back-substitution. Consider, for instance, the homogeneous system $U\mathbf{x} = \vec{0}$, where
\[
U = \begin{pmatrix} 2 & 1 & 1 \\ 0 & \tfrac{3}{2} & \tfrac{1}{2} \\ 0 & 0 & 0 \end{pmatrix},
\]
corresponding explicitly to the system
\[
\begin{cases}
2x + y + z = 0, \\
\tfrac{3}{2} y + \tfrac{1}{2} z = 0, \\
0 = 0
\end{cases}.
\]

This system is readily solved by back-substitution: one first sets $z = t$, in correspondence with the column of $U$ containing no pivot; one then obtains $y$ from the second equation; and finally one obtains $x$ from the first equation. With $t \in \mathbb{R}$ arbitrary, the resulting solution set is
\[
\begin{pmatrix} x \\ y \\ z \end{pmatrix} = t \begin{pmatrix} -\tfrac{1}{3} \\ -\tfrac{1}{3} \\ 1 \end{pmatrix}.
\]

For a general coefficient matrix, not necessarily already in row echelon form, one first reduces the matrix to row echelon form by means of a systematic procedure known as the Gauss elimination method, before applying back-substitution as above.

The Gauss elimination method consists of a sequence of elementary row operations applied to a matrix, with the goal of transforming it into row echelon form, that is, of filling the lower left-hand corner of the matrix with zeros as far as possible. Three types of elementary row operations are permitted: swapping two rows; multiplying a row by a nonzero number; and adding a multiple of one row to another row. Using these three operations, any matrix can always be transformed into a matrix in row echelon form.

Once, in addition, every leading coefficient of the resulting matrix equals $1$, and every column containing a leading coefficient has zeros in all its remaining entries, the matrix is said to be in reduced row echelon form. This final reduced form is unique, in the sense that it does not depend on the particular sequence of elementary row operations used to reach it.

\begin{example}
Applying the Gauss elimination method to the matrix
\[
\begin{pmatrix} 1 & 1 & -1 & 1 \\ 1 & 3 & 1 & 9 \\ 3 & 11 & 5 & 35 \end{pmatrix}.
\]
One first applies $R_2\leftarrow R_2-R_1$ and
$R_3\leftarrow R_3-3R_1$, obtaining
\[
\begin{pmatrix} 1 & 1 & -1 & 1 \\ 0 & 2 & 2 & 8 \\ 0 & 8 & 8 & 32 \end{pmatrix}.
\]
Next, $R_3\leftarrow R_3-4R_2$ and $R_2\leftarrow R_2/2$ give
\[
\begin{pmatrix} 1 & 1 & -1 & 1 \\ 0 & 1 & 1 & 4 \\ 0 & 0 & 0 & 0 \end{pmatrix}.
\]
Finally, $R_1\leftarrow R_1-R_2$ gives the reduced row echelon form
\[
\begin{pmatrix} 1 & 0 & -2 & -3 \\ 0 & 1 & 1 & 4 \\ 0 & 0 & 0 & 0 \end{pmatrix}.
\]

\end{example}

The Gauss elimination method is the principal tool for finding the solution set of a linear system with an arbitrary coefficient matrix $A$. Its validity rests on the fact that, whenever an elementary operation is performed on the equations of a system, the solution set of the system remains unchanged; and, crucially, an elementary operation on the equations of a system corresponds precisely to an elementary row operation on the corresponding coefficient matrix. Consequently, if $A$ is the coefficient matrix of the homogeneous system $A\mathbf{x} = \vec{0}$, and $U$ denotes the matrix obtained by applying the Gauss algorithm to $A$, then
\[
A\mathbf{x} = \vec{0} \iff U\mathbf{x} = \vec{0},
\]
so that the original system and the reduced system share exactly the same solution set.

\begin{example}
Consider the homogeneous system
\[
\begin{cases}
2x + y + z = 0, \\
x + 2y + z = 0, \\
x + y + \tfrac{2}{3} z = 0
\end{cases}.
\]
The coefficient matrix is
\[
A = \begin{pmatrix} 2 & 1 & 1 \\ 1 & 2 & 1 \\ 1 & 1 & \tfrac{2}{3} \end{pmatrix}.
\]
Applying the Gauss elimination method to $A$ yields the reduced matrix
\[
U = \begin{pmatrix} 2 & 1 & 1 \\ 0 & \tfrac{3}{2} & \tfrac{1}{2} \\ 0 & 0 & 0 \end{pmatrix}.
\]
This reduction can be obtained, for instance, through
\[
R_2\leftarrow R_2-\frac12R_1,\qquad
R_3\leftarrow R_3-\frac12R_1,\qquad
R_3\leftarrow R_3-\frac13R_2.
\]
Moreover, solving $U\mathbf{x} = \vec{0}$ by back-substitution exactly as in the earlier example yields the following solution set, with $t \in \mathbb{R}$ arbitrary:
\[
\begin{pmatrix} x \\ y \\ z \end{pmatrix} = t \begin{pmatrix} -\tfrac{1}{3} \\ -\tfrac{1}{3} \\ 1 \end{pmatrix}.
\]

\end{example}

The Gauss elimination method extends immediately to nonhomogeneous systems of the general form $A\mathbf{x} = \mathbf{b}$. Letting $[A | \mathbf{b}]$ denote the augmented matrix of the system, and $[U | \mathbf{b}']$ the matrix obtained by applying the Gauss algorithm to $[A | \mathbf{b}]$, one has the equivalence
\[
A\mathbf{x} = \mathbf{b} \iff U\mathbf{x} = \mathbf{b}',
\]
so that the two systems again share the same solution set.

\begin{example}
Consider the nonhomogeneous system
\[
\begin{cases}
2x + y + z = 6, \\
x + 2y + z = 6, \\
x + y + \tfrac{2}{3} z = 4
\end{cases}.
\]
The coefficient matrix and the augmented matrix are
\[
\begin{gathered}
A = \begin{pmatrix} 2 & 1 & 1 \\ 1 & 2 & 1 \\ 1 & 1 & \tfrac{2}{3} \end{pmatrix} \\
[A|\mathbf{b}] = \begin{pmatrix} 2 & 1 & 1 & 6 \\ 1 & 2 & 1 & 6 \\ 1 & 1 & \tfrac{2}{3} & 4 \end{pmatrix}.
\end{gathered}
\]
Applying the Gauss algorithm to $[A|\mathbf{b}]$ yields the reduced augmented matrix
\[
[U|\mathbf{b}'] = \begin{pmatrix} 2 & 1 & 1 & 6 \\ 0 & \tfrac{3}{2} & \tfrac{1}{2} & 3 \\ 0 & 0 & 0 & 0 \end{pmatrix}.
\]
Solving $U\mathbf{x} = \mathbf{b}'$ by back-substitution, setting $z = t$ for the column with no pivot, and then obtaining $y$ and $x$ in succession, yields the following solution set, with $t \in \mathbb{R}$ arbitrary:
\[
\begin{pmatrix} x \\ y \\ z \end{pmatrix} = \begin{pmatrix} 2 \\ 2 \\ 0 \end{pmatrix} + t \begin{pmatrix} -\tfrac{1}{3} \\ -\tfrac{1}{3} \\ 1 \end{pmatrix}.
\]

For comparison, replace the third equation by $x+y+5z=4$. Subtracting
the first two equations gives $x-y=0$, hence $x=y$. The first equation
then gives $3x+z=6$, whereas the modified third equation gives
$2x+5z=4$. Substituting $z=6-3x$ yields
\[
2x+5(6-3x)=4
\quad\Longrightarrow\quad
-13x=-26.
\]
Therefore
\[
x=y=2,\qquad z=0.
\]
Thus the modified system has the unique solution $(2,2,0)$.
\end{example}

The general procedure just illustrated applies to an arbitrary matrix $A \in M_{m,n}$ and vector $\mathbf{b} \in \mathbb{R}^m$: letting $[A|\mathbf{b}]$ denote the augmented matrix and $[U|\mathbf{b}']$ the corresponding matrix obtained from the Gauss algorithm, the two systems $A\mathbf{x} = \mathbf{b}$ and $U\mathbf{x} = \mathbf{b}'$ have the same solution set, that is, they are equivalent. The original system is consistent if and only if the number of pivots of $U$ equals the number of pivots of $[U|\mathbf{b}'];$ denoting this common number of pivots by $p$, the solution set depends on $n - p$ free parameters, corresponding precisely to the number of columns of $U$ that contain no pivot.

\section{Determinants, Inverses, and Rank}

\subsection{Determinants: Definition and Recursive Computation}

Given a square matrix, the determinant is a scalar value that encodes certain fundamental properties of the linear transformation described by the matrix. In the simplest case, the determinant admits an explicit closed-form expression.

\begin{definition}[Determinant of a $2\times 2$ Matrix]
Given the matrix
\[
A = \begin{pmatrix} a & b \\ c & d \end{pmatrix} \in M_2.
\]
The determinant of $A$ is the number $\operatorname{Det}(A) = ad - bc$.
\end{definition}

For matrices of larger size, the determinant is computed recursively, by means of a procedure known as Laplace expansion, which reduces the computation of an $n \times n$ determinant to the computation of several $(n-1) \times (n-1)$ determinants. Given $A \in M_n$, denote by $A_{ij} \in M_{n-1}$ the submatrix obtained by deleting the $i$-th row and the $j$-th column of $A$. The real number $C_{ij} = (-1)^{i+j} \operatorname{Det}(A_{ij})$ is called the cofactor of $a_{ij}$. It can be proved that the sum of the products of the elements of any row, or any column, of $A$ with their respective cofactors always yields the same value, regardless of which row or column is chosen; this common value is, by definition, the determinant of $A$.

\begin{theorem}[Laplace Expansion for the Determinant]
Fix any row $i \in \{1, \dots, n\}$. Then
\[
\begin{aligned}
\operatorname{Det}(A) &= a_{i1} C_{i1} + a_{i2} C_{i2} + \cdots + a_{in} C_{in} \\
&= \sum_{j=1}^{n} (-1)^{i+j} a_{ij} \operatorname{Det}(A_{ij}).
\end{aligned}
\]
Analogously, fixing any column $j$,
\[
\begin{aligned}
\operatorname{Det}(A) &= a_{1j} C_{1j} + a_{2j} C_{2j} + \cdots + a_{nj} C_{nj} \\
&= \sum_{i=1}^{n} (-1)^{i+j} a_{ij} \operatorname{Det}(A_{ij}).
\end{aligned}
\]

\end{theorem}

Several elementary properties of the determinant follow readily from this recursive definition. If $A$ has a row, or a column, consisting entirely of zeros, then $\operatorname{Det}(A) = 0$. The determinant is invariant under transposition,
\[
\operatorname{Det}(A^{\top}) = \operatorname{Det}(A).
\]

If $A$ is a triangular matrix, whether upper or lower, then its determinant equals the product of its diagonal entries,
\[
\operatorname{Det}(A) = a_{11} a_{22} \cdots a_{nn}.
\]

Finally, the determinant of a product of two square matrices of the same order factors as the product of the individual determinants.

\begin{theorem}[Binet's Theorem]
If $A$ and $B$ are square matrices of the same order, then
\[
\operatorname{Det}(AB) = \operatorname{Det}(A) \cdot \operatorname{Det}(B).
\]

\end{theorem}

It can be shown that the determinant is, in fact, the unique function $M_n \to \mathbb{R}$ satisfying three fundamental properties, and this axiomatic characterization is sometimes taken as the very definition of the determinant. First, $\operatorname{Det}(I) = 1$. Second, the determinant is alternating: whenever two rows of a matrix are identical, its determinant equals zero. Third, the determinant is multilinear, that is, it is a linear function of each individual row separately; for instance, linearity in the first row means that
\[
\operatorname{Det}\begin{pmatrix} \mathbf{a}_1 + t \mathbf{b}_1 \\ \mathbf{a}_2 \\ \vdots \\ \mathbf{a}_n \end{pmatrix} = \operatorname{Det}\begin{pmatrix} \mathbf{a}_1 \\ \mathbf{a}_2 \\ \vdots \\ \mathbf{a}_n \end{pmatrix} + t \cdot \operatorname{Det}\begin{pmatrix} \mathbf{b}_1 \\ \mathbf{a}_2 \\ \vdots \\ \mathbf{a}_n
\end{pmatrix}.
\]
This holds for every $t \in \mathbb{R}$.

On the basis of these three fundamental properties, one readily deduces how the determinant changes when elementary row operations are performed. If two rows of a matrix are interchanged, the determinant changes sign; and if any linear combination of the other rows is added to a given row, the determinant remains unchanged. These two facts allow one to compute the determinant of an arbitrary matrix $A$ by relating it to the determinant of a triangular matrix $U$ obtained from $A$ by Gaussian elimination, using only these two operations: if an even number of row interchanges is used, or none at all, then
\[
\operatorname{Det}(A) = \operatorname{Det}(U);
\]
whereas if an odd number of row interchanges is used, then
\[
\operatorname{Det}(A) = -\operatorname{Det}(U).
\]

This relationship yields a particularly efficient criterion for determining when a matrix has nonzero determinant. Since $\operatorname{Det}(A) \neq 0$ if and only if $\operatorname{Det}(U) \neq 0$, and since $\operatorname{Det}(U) = u_{11} u_{22} \cdots u_{nn}$ for a triangular matrix, it follows that $\operatorname{Det}(A) \neq 0$ if and only if every diagonal entry $u_{ii}$ of $U$ is nonzero, that is, if and only if $U$ has exactly $n$ pivots. Consequently, $\operatorname{Det}(A) \neq 0$ if and only if $\operatorname{Rank}(A) = n$.

\begin{definition}
A square matrix $A$ such that $\operatorname{Det}(A) \neq 0$ is called a nonsingular matrix.
\end{definition}

The nonsingularity of a matrix, as characterized by a nonvanishing determinant, has several important consequences for the associated linear system and linear transformation. A square linear system $A\mathbf{x} = \mathbf{b}$ has a unique solution if and only if $\operatorname{Det}(A) \neq 0$; in particular, if $\mathbf{b} = \vec{0}$, the only solution is the null vector. The columns, and equally the rows, of a square matrix are linearly independent vectors if and only if $\operatorname{Det}(A) \neq 0$. Finally, if $A$ is the representative matrix of a linear transformation $L : \mathbb{R}^n \to \mathbb{R}^n$, then $L$ is injective if and only if $L$ is surjective, and both properties hold if and only if $\operatorname{Det}(A) \neq 0$.

Figure~\ref{fig:atlas-determinant-area} interprets the determinant as an oriented area scale factor.

\begin{figure}[!htbp]
\centering
\begin{tikzpicture}[>=Stealth]
\begin{scope}[scale=1.25]
\fill[figblue!15] (0,0) rectangle (1,1);
\draw[black!35] (0,0) rectangle (1,1);
\draw[vecA] (0,0) -- (1,0) node[below] {$e_1$};
\draw[vecB] (0,0) -- (0,1) node[left] {$e_2$};
\node at (0.5,0.5) {$1$};
\end{scope}
\draw[->,thick] (2,0.8) -- (3,0.8) node[midway,above] {$A$};
\begin{scope}[shift={(4,0)},scale=1.25]
\fill[figgreen!15] (0,0) -- (2,0.5) -- (3,2.5) -- (1,2) -- cycle;
\draw[black!35] (2,0.5) -- (3,2.5) -- (1,2);
\draw[vecA] (0,0) -- (2,0.5) node[below right] {$Ae_1$};
\draw[vecB] (0,0) -- (1,2) node[above left] {$Ae_2$};
\node at (1.5,1.25) {$7/2$};
\end{scope}
\end{tikzpicture}
\caption{The unit square on the left is mapped to the parallelogram spanned by $(2,1/2)$ and $(1,2)$ on the right. Its area is $7/2$, the absolute determinant of the matrix with those columns. The positive determinant preserves the orientation of the ordered basis.}
\label{fig:atlas-determinant-area}
\end{figure}

\subsection{The Inverse of a Matrix}

\begin{definition}[Invertible Matrix]
Let $A \in M_n$. We say that $A$ is invertible if there exists a matrix $B \in M_n$ such that
\[
\begin{aligned}
AB &= BA \\
&= I.
\end{aligned}
\]
When such an inverse exists, it is unique, and is denoted by $A^{-1}$.
\end{definition}

\begin{theorem}
If $A$ is invertible, then $\operatorname{Det}(A) \neq 0$.
\end{theorem}

\begin{proof}
Since $A A^{-1} = I$, Binet's Theorem yields
\[
\begin{aligned}
\operatorname{Det}(A) \operatorname{Det}(A^{-1}) &= \operatorname{Det}(AA^{-1}) \\
&= \operatorname{Det}(I) \\
&= 1.
\end{aligned}
\]
This shows in particular that $\operatorname{Det}(A) \neq 0$, and moreover that
\[
\operatorname{Det}(A^{-1}) = 1/\operatorname{Det}(A).
\]

\end{proof}

The converse implication also holds, so that nonsingularity is not merely necessary but also sufficient for invertibility.

\begin{theorem}
If $\operatorname{Det}(A) \neq 0$, then $A$ is invertible.
\end{theorem}

\begin{proof}
Since $\operatorname{Det}(A) \neq 0$, for each $i = 1, \dots, n$ the system $A\mathbf{x} = \mathbf{e}_i$, where $\mathbf{e}_i$ denotes the $i$-th canonical basis vector, has a unique solution, which we call $\mathbf{b}^i$. Defining $B = [\mathbf{b}^1 | \cdots | \mathbf{b}^n]$, one observes that
\[
\begin{aligned}
AB &= [A\mathbf{b}^1 | \cdots | A\mathbf{b}^n] \\
&= [\mathbf{e}_1 | \cdots | \mathbf{e}_n] \\
&= I.
\end{aligned}
\]
Thus $A$ is invertible with inverse $B$.
\end{proof}

In practice, the inverse of a nonsingular matrix can be computed by applying the Gauss elimination method to the augmented matrix $[A | I]$ until a reduced matrix of the form $[I | B]$ is obtained; the resulting matrix $B$ is then the inverse of $A$. An explicit formula for the inverse, expressed directly in terms of cofactors, is also available.

\begin{theorem}[Laplace's Formula for the Inverse of a Matrix]
The inverse of a nonsingular matrix $A$ is given by $A^{-1} = \frac{1}{\operatorname{Det}(A)} C^{\top}$, where $C^{\top}$ denotes the transpose of the matrix of cofactors $C = [C_{ij}]$, with
\[
C_{ij} = (-1)^{i+j} \operatorname{Det}(A_{ij}).
\]

\end{theorem}

\begin{example}
We find the inverse, if it exists, of
\[
A = \begin{pmatrix} 3 & 4 & 2 \\ 1 & 2 & 0 \\ 1 & 0 & 1 \end{pmatrix}.
\]
Expanding along the first row gives
\[
\begin{aligned}
\operatorname{Det}(A)
&=3\begin{vmatrix}2&0\\0&1\end{vmatrix}
-4\begin{vmatrix}1&0\\1&1\end{vmatrix}
+2\begin{vmatrix}1&2\\1&0\end{vmatrix}\\
&=3(2)-4(1)+2(-2)\\
&=-2\neq0.
\end{aligned}
\]
Thus $A$ is invertible. The cofactor matrix is
\[
C=\begin{pmatrix}
2&-1&-2\\
-4&1&4\\
-4&2&2
\end{pmatrix}.
\]
Consequently, Laplace's formula yields
\[
\begin{aligned}
A^{-1}
&=\frac{1}{-2}C^{\top}\\
&=\begin{pmatrix}
-1&2&2\\[1mm]
\tfrac12&-\tfrac12&-1\\[1mm]
1&-2&-1
\end{pmatrix}.
\end{aligned}
\]
A direct multiplication confirms that $AA^{-1}=I$.

\end{example}

\subsection{The Rank of a Matrix}

We now consider matrices $A \in M_{m,n}$ of arbitrary, not necessarily square, dimension. Recall that the rank of $A$, already introduced in connection with the column space of $A$, is the dimension of $\operatorname{Col}(A)$, that is, the maximum number of linearly independent columns of $A$. The determinant provides an alternative, and in many contexts more computationally direct, characterization of this quantity.

\begin{definition}
Given $A \in M_{m,n}$, we call a minor of $A$ the determinant of any square submatrix of $A$, and we call the order of the minor the dimension of the corresponding square submatrix.
\end{definition}

\begin{proposition}
$\operatorname{Rank}(A) = r$ if and only if there exists a nonzero minor of order $r$, and every minor of order greater than $r$ vanishes. Equivalently, the rank of $A$ is the largest order among all nonzero minors of $A$.
\end{proposition}

Directly verifying that every minor of a given order beyond a certain threshold vanishes can be computationally expensive, since the number of such minors grows rapidly with the dimensions of the matrix. The following theorem shows that a much smaller number of minors need actually be checked.

\begin{theorem}[Kronecker's Theorem]
$\operatorname{Rank}(A) = r$ if and only if there is a nonzero minor of order $r$, corresponding to a certain square submatrix $B \in M_r$, and the minor of every square submatrix of order $r+1$ that contains $B$ as a submatrix vanishes.
\end{theorem}

\begin{example}
Consider the matrix
\[
A = \begin{pmatrix} 0 & 1 & 3 & 2 \\ 1 & 0 & -5 & -1 \\ 1 & 2 & 1 & 3 \end{pmatrix}.
\]
Since $A$ has only three rows, $\operatorname{Rank}(A) \leq 3$. The square submatrix
\[
B = \begin{pmatrix} 0 & 1 \\ 1 & 0
\end{pmatrix}.
\]
This submatrix has nonzero determinant, so $\operatorname{Rank}(A) \geq 2$. There are exactly two square submatrices of order $3$ containing $B$ as a submatrix, and direct computation shows that both of their determinants vanish; by Kronecker's Theorem, we conclude that $\operatorname{Rank}(A) = 2$.
Indeed, using columns $(1,2,3)$ and $(1,2,4)$, respectively,
\[
\begin{aligned}
\begin{vmatrix}0&1&3\\1&0&-5\\1&2&1\end{vmatrix}
&=-\bigl(1+5\bigr)+3(2)=0,\\
\begin{vmatrix}0&1&2\\1&0&-1\\1&2&3\end{vmatrix}
&=-\bigl(3+1\bigr)+2(2)=0.
\end{aligned}
\]
\end{example}

\subsection{Applications: Change of Basis and Similar Matrices}

The determinant, and the associated notion of invertibility, play a decisive role in two important constructions of linear algebra, namely the change of basis and the similarity of matrices.

Given any basis $\mathcal{B} = \{\mathbf{b}^1, \dots, \mathbf{b}^n\}$ of $\mathbb{R}^n$, and a vector $\mathbf{v} = (x_1, \dots, x_n)$ expressed in the canonical basis, one may ask how to find the coordinates $\hat{\mathbf{x}} = (\hat{x}_1, \dots, \hat{x}_n)$ of $\mathbf{v}$ with respect to $\mathcal{B}$. Setting $M = [\mathbf{b}^1 | \cdots | \mathbf{b}^n] \in M_n$, called the change of basis matrix, one has $\hat{\mathbf{x}} = M^{-1} \mathbf{x}$. Indeed, writing
\[
\begin{aligned}
\mathbf{v} &= \sum_{i=1}^n x_i \mathbf{e}_i \\
&= I \mathbf{x}.
\end{aligned}
\]
Using the canonical basis, and
\[
\begin{aligned}
\mathbf{v} &= \sum_{i=1}^n \hat{x}_i \mathbf{b}^i \\
&= M \hat{\mathbf{x}}.
\end{aligned}
\]
Using the basis $\mathcal{B}$, one obtains $\mathbf{x} = M\hat{\mathbf{x}}$; since $\operatorname{Det}(M) \neq 0$, because the columns of $M$ form a basis and are therefore linearly independent, one may invert this relation to find $\hat{\mathbf{x}} = M^{-1}\mathbf{x}$.

A closely related construction concerns the representation of a linear map in different bases. Let $A \in M_n$ be the representative matrix of a linear map $L : \mathbb{R}^n \to \mathbb{R}^n$ in the canonical basis, and take any basis $\mathcal{B} = \{\mathbf{b}^1, \dots, \mathbf{b}^n\}$ of $\mathbb{R}^n$. Setting again $M = [\mathbf{b}^1 | \cdots | \mathbf{b}^n]$, the representative matrix $B$ of $L$ in the new basis $\mathcal{B}$ is given by $B = M^{-1} A M$. Two matrices $A$ and $B$ related in this way, that is, satisfying $B = M^{-1}AM$ for some invertible matrix $M$, are said to be similar; two $n$-by-$n$ matrices are similar if and only if they represent the same linear transformation of $\mathbb{R}^n$ with respect to two different bases.

\section{Vector Spaces and Fundamental Subspaces}

\subsection{Vector Operations and Vector Spaces}

The foundational object underlying the entire theory developed in this document is the vector. Two fundamental operations govern the algebra of vectors, and it is useful to describe them both geometrically and in coordinates.

Geometrically, if two vectors $\mathbf{u}$ and $\mathbf{v}$ are positioned so that the initial point of $\mathbf{v}$ coincides with the terminal point of $\mathbf{u}$, then their sum $\mathbf{u} + \mathbf{v}$ is defined as the vector running from the initial point of $\mathbf{u}$ to the terminal point of $\mathbf{v}$. The scalar multiple $t\mathbf{v}$, for a real number $t$, is defined as the vector whose length equals the length of $\mathbf{v}$ multiplied by $|t|$; provided $t \neq 0$, the vector $t\mathbf{v}$ points in the same direction as $\mathbf{v}$, with the same orientation if $t$ is positive and the opposite orientation if $t$ is negative.

In coordinates, given two vectors $\mathbf{x} = (x_1, \dots, x_m)$ and $\mathbf{y} = (y_1, \dots, y_m)$ in $\mathbb{R}^m$, their sum, denoted $\mathbf{x}+\mathbf{y}$, is the vector of $\mathbb{R}^m$ given by
\[
\mathbf{x} + \mathbf{y} = (x_1+y_1, \dots, x_m+y_m).
\]
Given a number $t \in \mathbb{R}$ and a vector $\mathbf{x} = (x_1, \dots, x_m) \in \mathbb{R}^m$, the scalar multiple $t\mathbf{x}$ is the vector
\[
t\mathbf{x} = (t x_1, \dots, t x_m) \in \mathbb{R}^m.
\]

Consider the following set of $m$ vectors of $\mathbb{R}^m$,
\[
\mathbf{e}_1 = \begin{pmatrix} 1 \\ 0 \\ \vdots \\ 0 \end{pmatrix}, \qquad \mathbf{e}_2 = \begin{pmatrix} 0 \\ 1 \\ \vdots \\ 0 \end{pmatrix}, \qquad \ldots, \qquad \mathbf{e}_m = \begin{pmatrix} 0 \\ 0 \\ \vdots \\ 1 \end{pmatrix}.
\]

Given any vector $\mathbf{v} = (v_1, \dots, v_m)$, one can write $\mathbf{v} = v_1 \mathbf{e}_1 + \cdots + v_m \mathbf{e}_m$; in this sense, every vector of $\mathbb{R}^m$ can be reached from $\mathbf{e}_1, \dots, \mathbf{e}_m$ using only the two fundamental operations of addition and scalar multiplication.

Several pieces of terminology, central to the entire subject of linear algebra, arise directly from this observation. When a vector is obtained from a set of given vectors using only the two fundamental operations, one says that it is a linear combination of the given vectors. The set of all possible vectors reachable in this way from a given collection of vectors is called the span of those vectors. The vectors $\mathbf{e}_1, \dots, \mathbf{e}_m$ are linearly independent, meaning that none of them belongs to the span of the others; and one says that they form a basis of $\mathbb{R}^m$, capturing the fact that they constitute a set of linearly independent vectors that spans the entire space.

The properties observed above for $\mathbb{R}^m$ are, in fact, common to a much wider class of mathematical objects, axiomatized by the notion of a vector space.

\begin{definition}[Vector Space]
A vector space is a set $V$, whose elements are called vectors, on which two fundamental operations are defined: an internal operation called sum, $V \times V \to V$, $(\mathbf{u}, \mathbf{v}) \mapsto \mathbf{u}+\mathbf{v}$, and an external operation called scalar multiplication, $\mathbb{R} \times V \to V$, $(t, \mathbf{u}) \mapsto t\mathbf{u}$. These two operations must satisfy the following properties, for every $\mathbf{u}, \mathbf{v}, \mathbf{z} \in V$ and every $s, t \in \mathbb{R}$: associativity of addition, $\mathbf{u} + (\mathbf{v}+\mathbf{z}) = (\mathbf{u}+\mathbf{v}) + \mathbf{z}$; commutativity of addition, $\mathbf{u}+\mathbf{v} = \mathbf{v}+\mathbf{u}$; existence of a zero element $\vec{0} \in V$ such that $\mathbf{u}+\vec{0} = \mathbf{u}$; existence, for each $\mathbf{u} \in V$, of an element $\mathbf{w} \in V$ such that $\mathbf{u}+\mathbf{w} = \vec{0}$; the associativity relation $s(t\mathbf{u}) = (st)\mathbf{u}$; the identity relation $1\mathbf{u} = \mathbf{u}$; and the two distributive laws $t(\mathbf{u}+\mathbf{v}) = t\mathbf{u}+t\mathbf{v}$ and $(s+t)\mathbf{u} = s\mathbf{u}+t\mathbf{u}$.
\end{definition}

\begin{definition}[Subspace]
A subset $W \subset V$ is called a subspace of $V$ if it is closed under the operations of sum and scalar multiplication, that is, if $\mathbf{u}, \mathbf{v} \in W$ implies $\mathbf{u}+\mathbf{v} \in W$, and $\mathbf{u} \in W$, $t \in \mathbb{R}$ implies $t\mathbf{u} \in W$.
\end{definition}

\begin{remark}
Every subspace necessarily contains the zero vector: if $W$ is a subspace of $V$, then $\vec{0}_V \in W$. Moreover, a subspace $W \subset V$ inherits the vector space structure of $V$: a subspace is itself a vector space, with the same operations restricted to $W$.
\end{remark}

\begin{definition}[Linear Combination]
Given $k$ vectors $\mathbf{v}^1, \mathbf{v}^2, \dots, \mathbf{v}^k \in V$ and $k$ given numbers $a_1, a_2, \dots, a_k \in \mathbb{R}$, the vector $\mathbf{w} = a_1 \mathbf{v}^1 + a_2 \mathbf{v}^2 + \cdots + a_k \mathbf{v}^k \in V$ is called the linear combination of $\mathbf{v}^1, \mathbf{v}^2, \dots, \mathbf{v}^k$ with coefficients $a_1, a_2, \dots, a_k$.
\end{definition}

\begin{definition}[Span]
Given $k$ vectors $\mathbf{v}^1, \dots, \mathbf{v}^k \in V$, we denote by $\operatorname{span}(\mathbf{v}^1, \dots, \mathbf{v}^k)$ the set of all possible linear combinations of these vectors, that is,
\[
\operatorname{span}(\mathbf{v}^1, \dots, \mathbf{v}^k) = \left\{ \sum_{i=1}^{k} x_i \mathbf{v}^i : x_1, \dots, x_k \in \mathbb{R} \right\}.
\]

\end{definition}

One readily observes that $\operatorname{span}(\mathbf{v}^1, \dots, \mathbf{v}^k)$ is always a subspace of $V$, since it is closed under both fundamental operations; this subspace is accordingly called the subspace generated by the vectors $\mathbf{v}^1, \dots, \mathbf{v}^k$.

\subsection{Linear Independence, Basis, and Dimension}

\begin{definition}[Linear Dependence]
The vectors $\mathbf{v}^1, \mathbf{v}^2, \dots, \mathbf{v}^k \in V$ are called linearly dependent if there exist scalars $x_1, \dots, x_k$, not all zero, such that
\[
x_1 \mathbf{v}^1 + x_2 \mathbf{v}^2 + \cdots + x_k \mathbf{v}^k = \vec{0}.
\]

\end{definition}

If not every scalar in this relation vanishes, then at least one, say $x_1$, is nonzero, and the defining equation can be rearranged as $\mathbf{v}^1 = -(x_2/x_1) \mathbf{v}^2 - \cdots - (x_k/x_1) \mathbf{v}^k$, exhibiting $\mathbf{v}^1$ as a linear combination of the remaining vectors. In this sense, $\mathbf{v}^1$ is redundant, and one has the equality
\[
\operatorname{span}(\mathbf{v}^1, \mathbf{v}^2, \dots, \mathbf{v}^k) = \operatorname{span}(\mathbf{v}^2, \dots, \mathbf{v}^k),
\]
showing that a linearly dependent set can always be reduced without changing its span.

\begin{definition}[Linear Independence]
The vectors $\mathbf{v}^1, \mathbf{v}^2, \dots, \mathbf{v}^k \in V$ are called linearly independent if the relation $x_1 \mathbf{v}^1 + x_2 \mathbf{v}^2 + \cdots + x_k \mathbf{v}^k = \vec{0}$ implies $x_i = 0$ for every $i = 1, \dots, k$.
\end{definition}

A set of linearly independent vectors is, in this sense, minimal for the generation of its span, since no vector in the set is redundant. In the particular case of two vectors, linear dependence admits a simple characterization: $\mathbf{v}^1$ and $\mathbf{v}^2$ are linearly dependent if and only if one is a scalar multiple of the other, where the scalar is allowed to be zero. Equivalently, either at least one vector is zero, or both are nonzero and proportional.

\begin{example}
We investigate whether the vectors $\mathbf{v}^1 = (1,0,0)$, $\mathbf{v}^2 = (1, 1/2, 1)$, and $\mathbf{v}^3 = (5,1,2)$ are linearly independent in $\mathbb{R}^3$. Imposing the condition
\[
x_1 \mathbf{v}^1 + x_2 \mathbf{v}^2 + x_3 \mathbf{v}^3 = \vec{0},
\]
and writing the resulting equation in components, one obtains the linear system
\[
\begin{cases}
x_1 + x_2 + 5 x_3 = 0, \\
\tfrac{1}{2} x_2 + x_3 = 0, \\
x_2 + 2 x_3 = 0
\end{cases}.
\]

This system reduces to $x_1 = -3x_3$ and $x_2 = -2x_3$, with $x_3$ free to take any value in $\mathbb{R}$; consequently, the three vectors are linearly dependent. In particular, fixing $x_3 = 1$ yields the explicit relation $\mathbf{v}^3 = 3\mathbf{v}^1 + 2\mathbf{v}^2$.
\end{example}

The example just presented illustrates the general strategy for investigating the linear independence of a family of vectors in $\mathbb{R}^m$. Given $k$ vectors $\mathbf{a}^1, \dots, \mathbf{a}^k \in \mathbb{R}^m$, with $\mathbf{a}^j = (a_{1j}, a_{2j}, \dots, a_{mj})$ for each $j = 1, \dots, k$, one investigates their linear dependence or independence by asking whether the equation $x_1 \mathbf{a}^1 + x_2 \mathbf{a}^2 + \cdots + x_k \mathbf{a}^k = \vec{0}$ in $\mathbb{R}^m$ has the unique solution $x_i = 0$ for all $i$. Writing this equation out in components yields precisely a homogeneous linear system of $m$ equations in the $k$ unknowns $x_1, \dots, x_k$; consequently, proving the linear independence of a family of $k$ vectors in $\mathbb{R}^m$ amounts to showing that this associated homogeneous linear system has only the trivial solution, connecting the study of linear independence directly to the theory of linear systems developed earlier.

\begin{definition}[Basis]
Given vectors $\mathbf{b}^1, \mathbf{b}^2, \dots, \mathbf{b}^d \in V$, the set $\{\mathbf{b}^1, \dots, \mathbf{b}^d\}$ is called a basis for $V$ if the vectors $\mathbf{b}^1, \dots, \mathbf{b}^d$ are linearly independent, and if $\{\mathbf{b}^1, \dots, \mathbf{b}^d\}$ is a set of generators of $V$, that is,
\[
\operatorname{span}(\mathbf{b}^1, \dots, \mathbf{b}^d) = V.
\]

\end{definition}

By the very definition of span, every vector $\mathbf{v}$ in the space can be written as a linear combination of the basis vectors; the following theorem shows that the coefficients in this expansion are, in fact, uniquely determined.

\begin{theorem}
Let $\mathcal{B} = \{\mathbf{b}^1, \mathbf{b}^2, \dots, \mathbf{b}^d\}$ be a basis for $V$. Then, for every $\mathbf{v} \in V$, there is a unique set of $d$ scalars $x_1, x_2, \dots, x_d \in \mathbb{R}$ such that
\[
\mathbf{v} = \sum_{i=1}^{d} x_i \mathbf{b}^i.
\]
The scalars $x_1, \dots, x_d$ are called the components, or coordinates, of $\mathbf{v}$ in the basis $\mathcal{B}$.
\end{theorem}

\begin{proof}[Proof of uniqueness]
Suppose, for the sake of contradiction, that there exist a vector $\mathbf{v} \in V$ and two distinct sets of coefficients, $x_1, \dots, x_d$ and $\tilde{x}_1, \dots, \tilde{x}_d$, such that
\[
\begin{aligned}
\mathbf{v} &= \sum_{i=1}^{d} x_i \mathbf{b}^i \\
&= \sum_{i=1}^{d} \tilde{x}_i \mathbf{b}^i.
\end{aligned}
\]
Subtracting the two expressions yields
\[
\sum_{i=1}^{d} (x_i - \tilde{x}_i) \mathbf{b}^i = \vec{0}.
\]

Since the vectors $\mathbf{b}^1, \dots, \mathbf{b}^d$ are linearly independent, this forces $x_i - \tilde{x}_i = 0$ for every $i$, that is, $x_i = \tilde{x}_i$ for every $i$, contradicting the assumed distinctness of the two sets of coefficients.
\end{proof}

The number of vectors in a basis of a given vector space turns out to be an intrinsic property of the space itself, independent of the particular basis chosen, as recorded in the following fundamental theorem.

\begin{theorem}[Dimension Theorem]
If $V$ has a basis consisting of $d \in \mathbb{N}$ vectors, then every basis of $V$ consists of exactly $d$ vectors.
\end{theorem}

In light of this theorem, one says that $V$ has dimension $d$, and writes $\operatorname{Dim}(V) = d$. It follows moreover that, if $V$ has dimension $d$, then any set of $d$ linearly independent vectors of $V$ is automatically a basis of $V$; no separate verification that the set spans $V$ is required. As a basic example, the vector space $\mathbb{R}^m$ has dimension $m$; the three vectors $\mathbf{v}^1$, $\mathbf{v}^2$, $\mathbf{v}^3$ considered in the earlier example, having been shown to be linearly dependent, therefore fail to constitute a basis of $\mathbb{R}^3$.

Several useful facts relate the dimension of a subspace to the dimension of the ambient space. Let $W \subset V$ be a subspace, with $\operatorname{Dim}(V) = m$. Then $\operatorname{Dim}(W) \leq \operatorname{Dim}(V)$; if $W$ is a proper subspace of $V$, then $\operatorname{Dim}(W) < m$ strictly; and, conversely, if $\operatorname{Dim}(W) = m$, then necessarily $W = V$. The null subspace $W = \{\vec{0}_V\}$ has dimension $0$, and is in fact the unique subspace of dimension zero. Finally, if $W = \operatorname{span}(\mathbf{v}^1, \dots, \mathbf{v}^k)$, then $\operatorname{Dim}(W) = r$, where $r$ denotes the largest number of linearly independent vectors among the generators $\mathbf{v}^1, \dots, \mathbf{v}^k$; these $r$ vectors then form a basis of $W$. It follows that $r \leq m$ and $r \leq k$.

\subsection{Column Space, Kernel, and the Nullity--Rank Theorem}

\begin{definition}
Let $A = [a_{ij}] \in M_{m,n}$ be a given matrix. The column space of $A$ is the subspace of $\mathbb{R}^m$ generated by the columns of $A$, that is, $\operatorname{Col}(A) = \operatorname{span}(\mathbf{a}^1, \dots, \mathbf{a}^n) \subset \mathbb{R}^m$, where $\mathbf{a}^j$ denotes the $j$-th column of $A$. Equivalently,
\[
\operatorname{Col}(A) = \{ A\mathbf{x} : \mathbf{x} \in \mathbb{R}^n \},
\]
in light of the interpretation of the matrix-vector product as a linear combination of columns.
\end{definition}

\begin{definition}
The dimension of $\operatorname{Col}(A)$ is called the rank of $A$, denoted $\operatorname{Rank}(A)$.
\end{definition}

\begin{remark}
The computation of the rank is immediate whenever the matrix is already in row echelon form: the columns of the matrix that contain a pivot are linearly independent, while every remaining column belongs to the span of the pivot columns preceding it. Consequently, $\operatorname{Dim}(\operatorname{Col}(U)) = \operatorname{Rank}(U)$ equals the number of pivots of $U$, for a matrix $U$ in row echelon form; this same number of pivots also indicates the maximum number of linearly independent rows of $U$.
\end{remark}

\begin{definition}
Let $A = [a_{ij}] \in M_{m,n}$ be a given matrix. The kernel of $A$ is the subset of $\mathbb{R}^n$ given by
\[
\operatorname{Ker}(A) = \{ \mathbf{x} \in \mathbb{R}^n : A\mathbf{x} = \vec{0}_m \} \subset \mathbb{R}^n.
\]

\end{definition}

\begin{theorem}
$\operatorname{Ker}(A)$ is a vector space.
\end{theorem}

The proof of this theorem consists in verifying directly that $\operatorname{Ker}(A)$ is closed under vector addition and scalar multiplication, a straightforward consequence of the linearity of the matrix-vector product established earlier.

Figure~\ref{fig:atlas-kernel-image} separates the directions preserved by a projection from those annihilated by it.

\begin{figure}[!htbp]
\centering
\begin{tikzpicture}[scale=1.5,>=Stealth]
\draw[->,figblue,thick] (-1,0) -- (4,0) node[right] {$\operatorname{Im}P$};
\draw[->,figred,thick] (0,-0.6) -- (0,2.8) node[above] {$\ker P$};
\draw[->,black!70,thick] (0,0) -- (2.7,2.1);
\draw[->,figgreen,thick,dashed] (2.7,2.1) -- (2.7,0.05);
\fill (2.7,2.1) circle (2pt);
\fill[figblue] (2.7,0) circle (2pt);
\node[above right] at (2.7,2.1) {$(x,y)$};
\node[below] at (2.7,0) {$(x,0)$};
\node[below left] at (0,0) {$0$};
\end{tikzpicture}
\caption{For the projection $P(x,y)=(x,0)$, the horizontal axis is the image and the vertical axis is the kernel. A point and its projection differ by a vertical vector in the kernel. Both subspaces have dimension one, in agreement with rank plus nullity equaling two.}
\label{fig:atlas-kernel-image}
\end{figure}

We now come to one of the central results of linear algebra, relating the dimension of the kernel of a matrix to the dimension of its column space.

\begin{theorem}[Nullity-Rank Theorem]
Given any matrix $A$ with $n$ columns,
\[
\operatorname{Dim}(\operatorname{Ker}(A)) + \operatorname{Dim}(\operatorname{Col}(A)) = n.
\]

\end{theorem}

\begin{proof}
Assume that $\operatorname{Dim}(\operatorname{Ker}(A)) = k$, and let $\{\mathbf{u}^1, \dots, \mathbf{u}^k\}$ be a basis of $\operatorname{Ker}(A)$, so that
\[
\operatorname{Ker}(A) = \operatorname{span}(\mathbf{u}^1, \dots, \mathbf{u}^k).
\]

By an important theorem of linear algebra, one may select $n-k$ further vectors $\mathbf{w}^{k+1}, \dots, \mathbf{w}^n \in \mathbb{R}^n$, linearly independent both among themselves and from the vectors $\mathbf{u}^1, \dots, \mathbf{u}^k$, such that $\{\mathbf{u}^1, \dots, \mathbf{u}^k, \mathbf{w}^{k+1}, \dots, \mathbf{w}^n\}$ is a basis of $\mathbb{R}^n$; note that
\[
\operatorname{span}(\mathbf{w}^{k+1}, \dots, \mathbf{w}^n) \cap \operatorname{Ker}(A) = \{\vec{0}_n\}.
\]

The strategy of the proof is to show that the $n-k$ vectors $A\mathbf{w}^{k+1}, \dots, A\mathbf{w}^n \in \mathbb{R}^m$ form a basis of $\operatorname{Col}(A)$, from which $\operatorname{Dim}(\operatorname{Col}(A)) = n-k$ follows immediately, establishing the theorem. We first show that $\{A\mathbf{w}^{k+1}, \dots, A\mathbf{w}^n\}$ generates $\operatorname{Col}(A)$. Given any $\mathbf{x} \in \mathbb{R}^n$, there exist numbers $\alpha_1, \dots, \alpha_n \in \mathbb{R}$ such that
\[
\mathbf{x} = \sum_{j=1}^{k} \alpha_j \mathbf{u}^j + \sum_{j=k+1}^{n} \alpha_j \mathbf{w}^j.
\]
By linearity of the matrix-vector product,
\[
\begin{aligned}
A\mathbf{x} &= \sum_{j=1}^{k} \alpha_j A\mathbf{u}^j + \sum_{j=k+1}^{n} \alpha_j A\mathbf{w}^j \\
&= \sum_{j=k+1}^{n} \alpha_j A\mathbf{w}^j.
\end{aligned}
\]
Since $A\mathbf{u}^j = \vec{0}_m$ for each $j = 1, \dots, k$, as $\mathbf{u}^j \in \operatorname{Ker}(A)$. This shows that $A\mathbf{x} \in \operatorname{span}(A\mathbf{w}^{k+1}, \dots, A\mathbf{w}^n)$; since $\operatorname{Col}(A) = \{A\mathbf{x} : \mathbf{x} \in \mathbb{R}^n\}$, this establishes the claim. We next show that the vectors $A\mathbf{w}^{k+1}, \dots, A\mathbf{w}^n$ are linearly independent. Take $\beta_1, \dots, \beta_{n-k} \in \mathbb{R}$ such that
\[
\sum_{j=1}^{n-k} \beta_j A\mathbf{w}^{k+j} = \vec{0}_m.
\]
By linearity, this equals
\[
A( \sum_{j=1}^{n-k} \beta_j \mathbf{w}^{k+j} ),
\]
so that the vector
\[
\mathbf{u} := \sum_{j=1}^{n-k} \beta_j \mathbf{w}^{k+j} \in \operatorname{Ker}(A).
\]

But this same vector $\mathbf{u}$ also lies in $\operatorname{span}(\mathbf{w}^{k+1}, \dots, \mathbf{w}^n)$, and we recorded above that this span intersects $\operatorname{Ker}(A)$ only in the zero vector; hence $\mathbf{u} = \vec{0}_n$. Since the vectors $\mathbf{w}^{k+1}, \dots, \mathbf{w}^n$ are themselves linearly independent, the relation \[
\sum_{j=1}^{n-k} \beta_j \mathbf{w}^{k+j} = \vec{0}_n
\] forces $\beta_j = 0$ for every $j = 1, \dots, n-k$, establishing the linear independence of $A\mathbf{w}^{k+1},\allowbreak \dots, A\mathbf{w}^n$ and completing the proof.
\end{proof}


Let $A \in M_{m,n}$. A homogeneous system consists in finding $\mathbf{x} \in \mathbb{R}^n$ such that $A\mathbf{x} = \vec{0}_m$. Such a system is always solvable, since the null vector $\mathbf{x} = \vec{0}_n$ is trivially a solution; moreover, the entire solution set coincides precisely with $\operatorname{Ker}(A)$, and hence, by the Nullity-Rank Theorem, its dimension equals $n - \operatorname{Rank}(A)$, where $n$ denotes the number of unknowns.

\begin{theorem}
Let $A \in M_{m,n}$. The solution set of the homogeneous system $A\mathbf{x} = \vec{0}_m$ is a vector space of dimension $n - \operatorname{Rank}(A)$. Setting $r = \operatorname{Rank}(A)$, the following hold: if $r = n$, then $\mathbf{x} = \vec{0}_n$ is the unique solution of the system; if $r < n$, the system has $n-r$ linearly independent solutions $\mathbf{x}^1, \dots, \mathbf{x}^{n-r}$, and every solution is given by
\[
S = \{ \mathbf{x} = t_1 \mathbf{x}^1 + \cdots + t_{n-r} \mathbf{x}^{n-r} : t_i \in \mathbb{R} \}.
\]

The system therefore has infinitely many solutions, depending on $n-r$ free parameters, whenever $r < n$.
\end{theorem}

\begin{remark}
Recall that, when studying homogeneous linear systems by the Gauss elimination method, the solution set was found to depend on $n-p$ parameters, where $p$ denotes the number of pivots of the reduced matrix. Comparing this earlier result with the theorem just stated, in which the solutions depend on $n - \operatorname{Rank}(A)$ parameters, one concludes that $p = \operatorname{Rank}(A)$; that is, the rank of a matrix equals the number of pivots of its row-echelon reduction. This identity has a further important consequence: the number of linearly independent columns of a matrix necessarily coincides with the number of linearly independent rows.
\end{remark}

Let $A \in M_{m,n}$ and $\mathbf{b} \in \mathbb{R}^m$ with $\mathbf{b} \neq \vec{0}_m$. A nonhomogeneous system consists in finding $\mathbf{x} \in \mathbb{R}^n$ such that $A\mathbf{x} = \mathbf{b}$. Recalling that the matrix-vector product equals $A\mathbf{x} = x_1 \mathbf{a}^1 + x_2 \mathbf{a}^2 + \cdots + x_n \mathbf{a}^n$, solving $A\mathbf{x} = \mathbf{b}$ is equivalent to finding coefficients $x_1, \dots, x_n$, if any exist, such that $\mathbf{b}$ can be written as a linear combination of the columns $\mathbf{a}^1, \mathbf{a}^2, \dots, \mathbf{a}^n$ of $A$. This reformulation immediately raises the question of whether such a system is always solvable.

The answer is that the system is solvable, that is, consistent, if and only if $\mathbf{b} \in \operatorname{Col}(A)$, that is, if and only if
\[
\operatorname{Col}([A|\mathbf{b}]) = \operatorname{Col}(A).
\]
Since $\operatorname{Col}(A) \subset \operatorname{Col}([A|\mathbf{b}])$ always holds, the two column spaces coincide if and only if
\[
\operatorname{Dim}(\operatorname{Col}([A|\mathbf{b}])) = \operatorname{Dim}(\operatorname{Col}(A)),
\]
yielding the following solvability condition.

\begin{theorem}[Solvability Condition]
A nonhomogeneous system of linear equations $A\mathbf{x} = \mathbf{b}$ has a solution if and only if the rank of its coefficient matrix $A$ equals the rank of its augmented matrix $[A|\mathbf{b}]$, that is,
\[
\operatorname{Rank}([A|\mathbf{b}]) = \operatorname{Rank}(A).
\]

\end{theorem}

Unlike the homogeneous case, the solution set of a nonhomogeneous system does not, in general, form a vector space, since it need not contain the zero vector; nonetheless, it is closely related to the solution set of the associated homogeneous system.

\begin{theorem}[Structure of the Solution Set]
Given any specific solution $\mathbf{v}_p$ of $A\mathbf{x} = \mathbf{b}$, the entire solution set of the nonhomogeneous system is
\[
S = \{ \mathbf{y} = \mathbf{v} + \mathbf{v}_p : \mathbf{v} \in \mathbb{R}^n, \ A\mathbf{v} = \vec{0}_m \} = \operatorname{Ker}(A) + \mathbf{v}_p.
\]

\end{theorem}

\begin{proof}
It suffices to observe that $\mathbf{y} \in \mathbb{R}^n$ satisfies $A\mathbf{y} = \mathbf{b}$ if and only if
\[
\begin{aligned}
A(\mathbf{y} - \mathbf{v}_p) &= A\mathbf{y} - A\mathbf{v}_p \\
&= \mathbf{b} - \mathbf{b} \\
&= \vec{0}_m.
\end{aligned}
\]

Hence $\mathbf{y} - \mathbf{v}_p \in \operatorname{Ker}(A)$, so that $\mathbf{y}$ equals the sum of the particular solution $\mathbf{v}_p$ and an element of $\operatorname{Ker}(A)$, and conversely every such sum is a solution.
\end{proof}

Combining the solvability condition with the structure theorem for the solution set yields the following comprehensive result, which fully characterizes the possible behaviors of a general linear system.

\begin{theorem}[Rouch\'e--Capelli Theorem]
Let $\mathbf{b} \neq \vec{0}_m$. The nonhomogeneous system $A\mathbf{x} = \mathbf{b}$ has a solution if and only if
\[
\operatorname{Rank}(A) = \operatorname{Rank}([A|\mathbf{b}]).
\]

In this case, setting $r = \operatorname{Rank}(A)$: if $r = n$, the solution is unique; if $r < n$, the system has infinitely many solutions, depending on $n-r$ free parameters, and the solution set is
\[
S = \{ \mathbf{y} = \mathbf{v}_p + t_1 \mathbf{x}^1 + \cdots + t_{n-r} \mathbf{x}^{n-r} : t_i \in \mathbb{R} \}.
\]
Here $\mathbf{x}^1, \dots, \mathbf{x}^{n-r}$ are $n-r$ linearly independent solutions of the associated homogeneous system $A\mathbf{x} = \vec{0}_m$, and $\mathbf{v}_p$ is any particular solution of the nonhomogeneous system.
\end{theorem}

\section{Linear Maps}

\subsection{Linear Maps and Their Representative Matrix}

Given $A \in M_{m,n}$, one may associate with it the map $L : \mathbb{R}^n \to \mathbb{R}^m$, $L(\mathbf{x}) = A\mathbf{x}$, and we have already observed that this map is linear, satisfying
\[
\begin{gathered}
L(\mathbf{x}+\mathbf{y}) = L(\mathbf{x}) + L(\mathbf{y}) \\
L(t\mathbf{x}) = tL(\mathbf{x}).
\end{gathered}
\]
This holds for every $\mathbf{x}, \mathbf{y} \in \mathbb{R}^n$ and $t \in \mathbb{R}$, or, equivalently,
\[
L(s\mathbf{x}+t\mathbf{y}) = sL(\mathbf{x}) + tL(\mathbf{y}).
\]

It is worth noting that the columns of $A$ are precisely the images of the canonical basis vectors, $L(\mathbf{e}_j) = \mathbf{a}^j$ for every $j = 1, \dots, n$; consequently, if $\mathbf{x} = (x_1, \dots, x_n)$, then
\[
\begin{aligned}
L(\mathbf{x}) &= A\mathbf{x} \\
&= \sum_{j=1}^n x_j L(\mathbf{e}_j).
\end{aligned}
\]

Conversely, every linear map between finite-dimensional Euclidean spaces can be represented in this fashion by a suitable matrix.

\begin{theorem}[Representation Theorem for Linear Maps]
Let $L : \mathbb{R}^n \to \mathbb{R}^m$ be any linear map, that is, a map satisfying $L(s\mathbf{x}+t\mathbf{y}) = sL(\mathbf{x}) + tL(\mathbf{y})$ for every $\mathbf{x}, \mathbf{y} \in \mathbb{R}^n$ and $s, t \in \mathbb{R}$. Setting $\mathbf{a}^j = L(\mathbf{e}_j)$ for $j = 1, \dots, n$, and $A = [\mathbf{a}^1 | \mathbf{a}^2 | \cdots | \mathbf{a}^n] \in M_{m,n}$, called the representative matrix of $L$ in the canonical basis, one has $L(\mathbf{x}) = A\mathbf{x}$ for every $\mathbf{x} \in \mathbb{R}^n$.
\end{theorem}

\begin{proof}
Take any vector
\[
\mathbf{x} = (x_1, \dots, x_n) = \sum_{j=1}^n x_j \mathbf{e}_j \in \mathbb{R}^n.
\]
By the linearity of $L$,
\[
\begin{aligned}
L(\mathbf{x}) &= L( \sum_{j=1}^n x_j \mathbf{e}_j ) \\
&= \sum_{j=1}^n x_j L(\mathbf{e}_j) \\
&= \sum_{j=1}^n x_j \mathbf{a}^j \\
&= A\mathbf{x},
\end{aligned}
\]
where the last equality follows from the interpretation of the matrix-vector product as a linear combination of columns.
\end{proof}

This representation theorem shows that the action of any linear map between Euclidean spaces is represented, once bases have been fixed, by a matrix-vector product.

\begin{example}
Consider the linear transformation of $\mathbb{R}^2$ describing a rotation of $180$ degrees counterclockwise combined with a stretching by a factor of $2$. To find its representative matrix, one determines where the canonical basis vectors land under this transformation: $\mathbf{e}_1$ maps to $(-2,0)$, and $\mathbf{e}_2$ maps to $(0,-2)$. Any vector $\mathbf{x} = x\mathbf{e}_1 + y\mathbf{e}_2$ therefore lands on
\[
x(-2,0) + y(0,-2) = \begin{pmatrix} -2 & 0 \\ 0 & -2 \end{pmatrix} \begin{pmatrix} x \\ y \end{pmatrix},
\]
identifying the representative matrix of the transformation as
\[
\begin{pmatrix} -2 & 0 \\ 0 & -2 \end{pmatrix}.
\]

\end{example}

Figure~\ref{fig:ch8-linear-map} shows how a linear map transforms a coordinate grid and its basis vectors.

\begin{figure}[!htbp]
\centering
\begin{tikzpicture}[scale=1.05,>=Stealth]
  \begin{scope}
    \fill[figblue!8] (0,0) rectangle (1,1);
    \foreach \gridindex in {0,...,2} {
      \draw[gridline] (\gridindex,0) -- (\gridindex,2);
      \draw[gridline] (0,\gridindex) -- (2,\gridindex);
    }
    \draw[->] (-0.4,0) -- (2.6,0) node[right,font=\small] {$x_1$};
    \draw[->] (0,-0.4) -- (0,3.4) node[above,font=\small] {$x_2$};
    \draw[vecB] (0,0) -- (1,0) node[below=4pt,font=\small] {$e_1$};
    \draw[vecC] (0,0) -- (0,1) node[left=4pt,font=\small] {$e_2$};
    \node[below left,font=\small] at (0,0) {$0$};
  \end{scope}
  \draw[->,very thick] (2.9,1.5) -- (4.8,1.5)
    node[midway,above=5pt,font=\small] {$T(\bx)=A\bx$};
  \begin{scope}[xshift=5.5cm]
    \fill[figblue!8] (0,0) -- (1.3,0.4) -- (1.9,1.5) -- (0.6,1.1) -- cycle;
    \foreach \gridindex in {0,...,2} {
      \draw[gridline] ({1.3*\gridindex},{0.4*\gridindex})
        -- ({1.3*\gridindex+1.2},{0.4*\gridindex+2.2});
      \draw[gridline] ({0.6*\gridindex},{1.1*\gridindex})
        -- ({2.6+0.6*\gridindex},{0.8+1.1*\gridindex});
    }
    \draw[->] (-0.4,0) -- (4.2,0) node[right,font=\small] {$y_1$};
    \draw[->] (0,-0.4) -- (0,3.4) node[above,font=\small] {$y_2$};
    \draw[vecB] (0,0) -- (1.3,0.4)
      node[right=4pt,font=\small,fill=white,inner sep=1pt] {$A e_1$};
    \draw[vecC] (0,0) -- (0.6,1.1)
      node[above=4pt,font=\small,fill=white,inner sep=1pt] {$A e_2$};
    \node[below left,font=\small] at (0,0) {$0$};
  \end{scope}
\end{tikzpicture}
\caption[Linear transformation of a grid]{A coordinate grid before and after an
  illustrative linear map $T(\bx)=A\bx$. The shaded unit square on the left maps
  to the shaded parallelogram spanned by the columns $Ae_1,Ae_2$ on the right.
  Matching colors identify each basis vector and its image; both coordinate
  systems have their axes intersecting at the origin.}
\label{fig:ch8-linear-map}
\end{figure}

\subsection{Image, Kernel, and Injectivity/Surjectivity}

A central goal in the study of linear maps is to determine conditions guaranteeing their injectivity or surjectivity. To this end, given a linear map $L : \mathbb{R}^n \to \mathbb{R}^m$ with representative matrix $A$, we define the range, or image, of $L$ as the collection of all images $L(\mathbf{x})$ for $\mathbf{x} \in \mathbb{R}^n$, that is, $\operatorname{Im}(L) = \operatorname{Col}(A)$; and we define the kernel of $L$ as the subspace of vectors whose image is the zero vector,
\[
\operatorname{Ker}(L) = \operatorname{Ker}(A).
\]

\begin{theorem}
The linear map $L : \mathbb{R}^n \to \mathbb{R}^m$ is surjective if and only if $\operatorname{Rank}(A) = m$.
\end{theorem}

This theorem follows immediately from the fact that $L$ is surjective precisely when $\operatorname{Im}(L) = \operatorname{Col}(A)$ equals the entire codomain $\mathbb{R}^m$, which occurs exactly when $\operatorname{Dim}(\operatorname{Col}(A)) = m$, that is, when $\operatorname{Rank}(A) = m$.

\begin{theorem}[Injectivity]
The linear map $L : \mathbb{R}^n \to \mathbb{R}^m$ is one-to-one if and only if $\operatorname{Ker}(L) = \{\vec{0}_n\}$.
\end{theorem}

\begin{proof}
Assume first that $L$ is injective. Since $L(\vec{0}_n) = \vec{0}_m$ by linearity, and $L$ is injective, no other vector can share this image; hence, for every $\mathbf{x} \neq \vec{0}_n$, one has $L(\mathbf{x}) \neq \vec{0}_m$, so that no nonzero vector belongs to $\operatorname{Ker}(L)$, that is,
\[
\operatorname{Ker}(L) = \{\vec{0}_n\}.
\]

Conversely, assume $\operatorname{Ker}(L) = \{\vec{0}_n\}$, and suppose, for the sake of contradiction, that $L(\mathbf{x}_1) = L(\mathbf{x}_2)$ for some $\mathbf{x}_1 \neq \mathbf{x}_2$ in $\mathbb{R}^n$. By linearity,
\[
\begin{aligned}
L(\mathbf{x}_1 - \mathbf{x}_2) &= L(\mathbf{x}_1) - L(\mathbf{x}_2) \\
&= \vec{0}_m.
\end{aligned}
\]
Consequently, $\mathbf{x}_1 - \mathbf{x}_2 \in \operatorname{Ker}(L)$. But the only element of $\operatorname{Ker}(L)$ is $\vec{0}_n$, so $\mathbf{x}_1 - \mathbf{x}_2 = \vec{0}_n$, contradicting the assumption $\mathbf{x}_1 \neq \mathbf{x}_2$; this contradiction establishes injectivity.
\end{proof}

The linear map and its representative matrix have the same kernel:
\[
\operatorname{Ker}(L) = \operatorname{Ker}(A),
\]
The Nullity-Rank Theorem yields
\[
\operatorname{Dim}(\operatorname{Ker}(L)) = n - \operatorname{Rank}(A),
\]
from which one obtains the following corollary.

\begin{corollary}
The linear map $L : \mathbb{R}^n \to \mathbb{R}^m$ is one-to-one if and only if $\operatorname{Rank}(A) = n$.
\end{corollary}


Given a square matrix $A \in M_n$, the linear function $T : \mathbb{R}^n \to \mathbb{R}^n$, $T(\mathbf{x}) = A\mathbf{x}$, represents a geometric transformation of $\mathbb{R}^n$ into itself, called an endomorphism. Combining the injectivity and surjectivity criteria established above with the earlier characterization of nonsingularity in terms of the determinant, one obtains a particularly clean criterion for such maps: the linear function $T : \mathbb{R}^n \to \mathbb{R}^n$ is injective if and only if it is surjective, and both properties hold precisely when the rank of its representative matrix equals $n$, that is, $\operatorname{Rank}(A) = n$; equivalently, as established earlier, this occurs if and only if $\operatorname{Det}(A) \neq 0$.

\section{Spectral and Singular-Value Methods}

\subsection{Eigenvalues, Eigenvectors, and the Characteristic Equation}

Given $A \in M_n$, consider the associated linear transformation $T : \mathbb{R}^n \to \mathbb{R}^n$, $T(\mathbf{x}) = A\mathbf{x}$. Among the many vectors that this transformation acts upon, some special nonzero vectors change, under the action of $T$, by at most a scalar factor, that is, they are simply stretched, or possibly reversed, by the transformation, in the sense that $T(\mathbf{v}) = \lambda \mathbf{v}$ for some scalar $\lambda$. Such a special vector is called an eigenvector, and the corresponding scalar $\lambda$, the factor by which it is stretched, is called the eigenvalue associated with that eigenvector.

\begin{definition}
Let $A \in M_n$ be a given square matrix. A vector $\mathbf{u} \in \mathbb{R}^n$, $\mathbf{u} \neq \vec{0}$, is an eigenvector of $A$ if $A\mathbf{u} = \lambda \mathbf{u}$ for some $\lambda \in \mathbb{R}$. The scalar $\lambda$ is called the eigenvalue corresponding to that eigenvector.
\end{definition}

To determine systematically the pairs $\mathbf{u} \neq \vec{0}$ and $\lambda \in \mathbb{R}$ satisfying this relation, one observes that the defining equation $A\mathbf{u} = \lambda \mathbf{u}$ can be equivalently rewritten as $(A - \lambda I)\mathbf{u} = \vec{0}$, which is a homogeneous linear system in the unknown $\mathbf{u}$, with coefficient matrix $A - \lambda I$. Finding eigenvectors of $A$, for a fixed value of $\lambda$, therefore amounts to finding nontrivial solutions of this homogeneous system, and, by the general theory of linear systems developed earlier, such nontrivial solutions exist if and only if the coefficient matrix $A - \lambda I$ is singular, that is, if and only if its determinant vanishes. This observation yields the following fundamental characterization of eigenvalues.

\begin{theorem}
A scalar $\lambda \in \mathbb{R}$ is an eigenvalue of $A$ if and only if $\lambda$ satisfies the equation $\operatorname{Det}(A - \lambda I) = 0$.
\end{theorem}

Several pieces of terminology accompany this characterization. The function $P(t) = \operatorname{Det}(A - tI)$ is a polynomial of degree $n$ in the variable $t$, called the characteristic polynomial of $A$; the equation $P(t) = 0$ is called the characteristic equation of $A$. By the Fundamental Theorem of Algebra, the characteristic polynomial of an $n$-by-$n$ matrix has exactly $n$ complex roots, counted with multiplicity and not necessarily distinct. Given any fixed real eigenvalue $\lambda$ of $A$, one may factor the characteristic polynomial as $P(t) = (t-\lambda)^k Q(t)$, for some polynomial $Q$ satisfying $Q(\lambda) \neq 0$ and some positive integer $k$; this integer $k$ is called the algebraic multiplicity of $\lambda$, and is denoted $a_\lambda$.

Associated with each real eigenvalue $\lambda$ of $A$ is a distinguished subspace, obtained as the solution set of the corresponding homogeneous system: the solution set of $(A-\lambda I)\mathbf{x} = \vec{0}$ is the vector space $V_\lambda = \operatorname{Ker}(A - \lambda I)$, called the eigenspace corresponding to $\lambda$, and it contains precisely all the eigenvectors of $A$ associated with $\lambda$, together with the zero vector. The dimension of this eigenspace is called the geometric multiplicity of $\lambda$, denoted $g_\lambda$; by the Nullity-Rank Theorem,
\[
\begin{aligned}
g_\lambda &= \operatorname{Dim}(V_\lambda) \\
&= n - \operatorname{Rank}(A - \lambda I).
\end{aligned}
\]
Figure~\ref{fig:ch8-eigenvectors} highlights the directions preserved by a linear transformation.

\begin{figure}[!htbp]
\centering
\begin{tikzpicture}[scale=1.65,>=Stealth]
  \draw[->] (-2.3,0) -- (2.3,0) node[right,font=\scriptsize]{$x_1$};
  \draw[->] (0,-1.6) -- (0,1.9) node[above,font=\scriptsize]{$x_2$};
  \draw[figblue!35, thick] (-2,-1) -- (2,1);
  \draw[figred!35, thick] (-1.2,1.6) -- (1.2,-1.6);
  \draw[vecA] (0,0) -- (1,0.5)
    node[below right,fill=white,inner sep=1pt,font=\scriptsize]{$\bv_1$};
  \draw[vecB] (0,0) -- (0.6,-0.8)
    node[below right,fill=white,inner sep=1pt,font=\scriptsize]{$\bv_2$};
  \draw[vecA,dashed] (0,0) -- (1.8,0.9)
    node[above right,fill=white,inner sep=1pt,font=\scriptsize]{$A\bv_1=\lambda_1 \bv_1$};
  \draw[vecB,dashed] (0,0) -- (0.3,-0.4);
  \node[figred,fill=white,inner sep=1pt,font=\scriptsize,anchor=north west]
    at (0.42,-0.44) {$A\bv_2=\lambda_2 \bv_2$};
  \draw[black!55,very thick,->] (0,0) -- (-1,1)
    node[left,fill=white,inner sep=1pt,font=\scriptsize]{$\bu$};
  \draw[black!55,very thick,dotted,->] (0,0) -- (-0.4,1.55)
    node[above right,fill=white,inner sep=1pt,font=\scriptsize]{$A\bu$};
  \node[anchor=south east,fill=white,inner sep=1pt,font=\scriptsize]
    at (2.2,-1.5) {$\lambda_1 \approx 1.8,\ \lambda_2\approx 0.5$};
\end{tikzpicture}
\caption[Eigenvectors: invariant directions]{Invariant lines (eigenspaces) in light colors:
  $\bv_1$ is stretched ($\lambda_1>1$), $\bv_2$ is contracted ($0<\lambda_2<1$), both remaining
  on their respective lines; the generic vector $\bu$ (gray) is also rotated.}
\label{fig:ch8-eigenvectors}
\end{figure}

The characterization of eigenvalues in terms of the characteristic equation yields a systematic three-step procedure for determining the eigenvalues and eigenvectors of a given matrix. The first step consists of finding the values of $t \in \mathbb{R}$ satisfying $\operatorname{Det}(A - tI) = 0$; this is a polynomial equation of degree $n$ in $t$, whose real solutions, if any, are precisely the real eigenvalues of $A$, and from whose factorization one determines the algebraic multiplicity $a_\lambda$ of each eigenvalue $\lambda$ found. The second step consists of writing, for each eigenvalue $\lambda$ found in the first step, the corresponding homogeneous linear system with coefficient matrix $A - \lambda I$. The third step consists of solving this homogeneous system $(A - \lambda I)\mathbf{x} = \vec{0}$; every nonzero solution $\mathbf{x}$ of this system is an eigenvector of $A$ corresponding to the eigenvalue $\lambda$, and the dimension of the solution space $\operatorname{Ker}(A - \lambda I)$ determines the geometric multiplicity $g_\lambda$.

\begin{example}
We determine the eigenvalues and eigenvectors of
\[
A = \begin{pmatrix} 2 & 0 \\ 2 & 1 \end{pmatrix}.
\]
We first use the characteristic equation to determine the eigenvalues:
\[
\operatorname{Det}(A - \lambda I) = \begin{vmatrix} 2-\lambda & 0 \\ 2 & 1-\lambda \end{vmatrix} = (1-\lambda)(2-\lambda).
\]
Consequently, the eigenvalues are $\lambda_1 = 1$ and $\lambda_2 = 2$. To determine the eigenvectors corresponding to $\lambda_i$, we solve $A\mathbf{x} = \lambda_i \mathbf{x}$ with $\mathbf{x} = (x,y)$, that is, the system $2x - \lambda_i x = 0$, $2x + y - \lambda_i y = 0$. For $\lambda_1 = 1$, this system yields $\mathbf{x} = (0,t) = t(0,1)$, for $t \neq 0$; for $\lambda_2 = 2$, it yields $\mathbf{x} = (t, 2t) = t(1,2)$, for $t \neq 0$.
\end{example}

Several general remarks concerning the multiplicities $g_\lambda$ and $a_\lambda$ are worth recording. The geometric multiplicity $g_\lambda$ equals the maximum number of linearly independent eigenvectors associated with the eigenvalue $\lambda$, and satisfies $g_\lambda \geq 1$, since every eigenvalue has at least one associated eigenvector by definition. A further, more difficult result, whose proof we omit, shows that $g_\lambda \leq a_\lambda$ always; this inequality may be strict, and whenever $g_\lambda = a_\lambda$ one says that $\lambda$ is a regular eigenvalue. In particular, if $\lambda$ is simple, that is, if $a_\lambda = 1$, then necessarily $\lambda$ is regular, since
\[
\begin{aligned}
1 &\leq g_\lambda \\
&\leq a_\lambda \\
&= 1
\end{aligned}
\]
forces
\[
\begin{aligned}
g_\lambda &= a_\lambda \\
&= 1.
\end{aligned}
\]

\begin{proposition}
Two eigenvectors corresponding to two different eigenvalues are linearly independent.
\end{proposition}

\begin{proof}
Suppose $A\mathbf{v}^1 = \lambda_1 \mathbf{v}^1$ and $A\mathbf{v}^2 = \lambda_2 \mathbf{v}^2$ for some $\lambda_1 \neq \lambda_2$. To show that $\mathbf{v}^1, \mathbf{v}^2$ are linearly independent, consider a linear combination
\[
c_1 \mathbf{v}^1 + c_2 \mathbf{v}^2 = \vec{0}.
\]
Applying $A$ to both sides, and using linearity together with the eigenvalue equations, yields
\[
\vec{0} = c_1 \lambda_1 \mathbf{v}^1 + c_2 \lambda_2 \mathbf{v}^2.
\]
Multiplying the original relation $c_1 \mathbf{v}^1 + c_2 \mathbf{v}^2 = \vec{0}$ by $\lambda_2$ and subtracting the two resulting equations, we obtain
\[
\vec{0} = c_1(\lambda_2 - \lambda_1) \mathbf{v}^1.
\]

Since an eigenvector is, by definition, nonzero, and $\lambda_2 \neq \lambda_1$, this forces $c_1 = 0$; substituting back into the original relation then yields $c_2 = 0$ as well, establishing linear independence.
\end{proof}

\subsection{Diagonalization and Spectral Coordinates}

We now investigate under which conditions a matrix $A \in M_n$ possesses the maximum possible number, namely $n$, of linearly independent eigenvectors. Since the total number of linearly independent eigenvectors of $A$ equals \[
\sum_\lambda g_\lambda,
\] summed over the distinct eigenvalues $\lambda$ of $A$, and since this sum cannot exceed $n$, the dimension of the ambient space, one seeks conditions under which the sum \[
\sum_\lambda g_\lambda
\] attains exactly the value $n$.

\begin{proposition}
Let $A \in M_n$ be such that $A$ has exactly $n$ real eigenvalues, not necessarily distinct, that is, \[
\sum_\lambda a_\lambda = n,
\] and each eigenvalue $\lambda$ is regular, that is, $g_\lambda = a_\lambda$. Then there exists a basis of $\mathbb{R}^n$ made up of eigenvectors of $A$.
\end{proposition}

An important special class of matrices automatically satisfies these hypotheses, as recorded in the following theorem.

\begin{theorem}[Spectral Theorem]
Let $A \in M_n$ be a symmetric matrix. Then $A$ has $n$ real eigenvalues, not necessarily distinct, and $A$ has $n$ linearly independent, mutually orthogonal eigenvectors. Consequently, for every symmetric matrix $A \in M_n$, there exists an orthonormal basis of $\mathbb{R}^n$ made up of eigenvectors of $A$.
\end{theorem}


When a matrix $A \in M_n$ possesses $n$ linearly independent eigenvectors, it can be decomposed in a particularly illuminating way, known as its eigendecomposition. Let $\mathbf{u}^1, \dots, \mathbf{u}^n$ be $n$ linearly independent eigenvectors of $A$, with corresponding eigenvalues $\lambda_1, \dots, \lambda_n$. Construct the nonsingular matrix $P = [\mathbf{u}^1 | \cdots | \mathbf{u}^n]$, called the change of basis matrix, and the diagonal matrix
\[
D = \operatorname{Diag}(\lambda_1, \dots, \lambda_n) = \begin{pmatrix} \lambda_1 & 0 & \cdots & 0 \\ 0 & \lambda_2 & \cdots & 0 \\ \vdots & \vdots & \ddots & \vdots \\ 0 & 0 & \cdots & \lambda_n \end{pmatrix}.
\]
The product $AP$ is computed as
\[
\begin{aligned}
AP &= A[\mathbf{u}^1 | \cdots | \mathbf{u}^n] \\
&= [A\mathbf{u}^1 | \cdots | A\mathbf{u}^n] \\
&= [\lambda_1 \mathbf{u}^1 | \cdots | \lambda_n \mathbf{u}^n].
\end{aligned}
\]
Using the defining property of eigenvectors applied to each column separately. On the other hand, the product $PD$ is computed as
\[
PD = [\mathbf{u}^1 | \cdots | \mathbf{u}^n] \begin{pmatrix} \lambda_1 & 0 & \cdots & 0 \\ 0 & \lambda_2 & \cdots & 0 \\ \vdots & \vdots & \ddots & \vdots \\ 0 & 0 & \cdots & \lambda_n \end{pmatrix} = [\lambda_1 \mathbf{u}^1 | \cdots | \lambda_n \mathbf{u}^n].
\]
By the definition of matrix multiplication applied to a diagonal matrix. Since these two computations agree, we conclude that $AP = PD$, and hence, since $P$ is nonsingular,
\[
A = PDP^{-1}.
\]

Summarizing, if the eigenvectors of $A$ form a basis of $\mathbb{R}^n$, then $A$ can be decomposed into the product of the matrix $P$, whose columns are the eigenvectors of $A$, the diagonal matrix $D$, whose diagonal entries are the corresponding eigenvalues of $A$, called the diagonal form of $A$, and the inverse $P^{-1}$ of the eigenvector matrix; this decomposition is called the eigendecomposition of the matrix.

\begin{definition}[Similar Matrices]
Given two matrices $A, B \in M_n$, we say that they are similar if there exists an invertible matrix $M \in M_n$ such that $A = MBM^{-1}$, or equivalently $M^{-1}AM = B$.
\end{definition}

\begin{definition}
A matrix $A \in M_n$ is diagonalizable, over $\mathbb{R}$, if it is similar to a diagonal matrix $D \in M_n$. In this case, $D$ is called the diagonal form of $A$.
\end{definition}

The eigendecomposition constructed above shows immediately that if $\mathbb{R}^n$ has a basis of eigenvectors of $A$, then $A$ is diagonalizable, since $A = PDP^{-1}$ exhibits the required similarity with a diagonal matrix. The converse also holds, essentially by reversing the argument: if $A$ is diagonalizable, so that $A = PDP^{-1}$ for some invertible matrix $P$ and diagonal matrix $D$, then the columns of $P$ are $n$ linearly independent eigenvectors of $A$, with the corresponding diagonal entries of $D$ giving the associated eigenvalues. This yields the following characterization.

\begin{theorem}[Diagonalization Test I]
$A \in M_n$ is diagonalizable over $\mathbb{R}$ if and only if $\mathbb{R}^n$ has a basis of eigenvectors of $A$.
\end{theorem}

Combining this criterion with the earlier discussion of algebraic and geometric multiplicities yields an equivalent, and often more directly verifiable, criterion.

\begin{theorem}[Diagonalization Test II]
$A \in M_n$ is diagonalizable over $\mathbb{R}$ if and only if $A$ has exactly $n$ real eigenvalues, not necessarily distinct, and each eigenvalue is regular.
\end{theorem}

Two particular situations guarantee diagonalizability automatically, and are worth recording as sufficient conditions of frequent practical use. If $A \in M_n$ has $n$ distinct eigenvalues, that is, if its characteristic polynomial has $n$ distinct roots, then $A$ is diagonalizable, since every simple eigenvalue is automatically regular, as observed earlier. If $A \in M_n$ is symmetric, then $A$ is diagonalizable, as guaranteed directly by the Spectral Theorem.

\begin{example}
The matrix
\[
C = \begin{pmatrix} 0 & 1 \\ 0 & 0
\end{pmatrix}.
\]
This is not diagonalizable: it has a single eigenvalue $\lambda = 0$, of algebraic multiplicity $2$, but the corresponding eigenspace has dimension only $1$, so that this eigenvalue is not regular.
\end{example}

\begin{example}
The matrix
\[
B = \begin{pmatrix} 0 & 1 \\ -1 & 0
\end{pmatrix}.
\]
This is not diagonalizable over $\mathbb{R}$, since it possesses no real eigenvalues at all: its characteristic polynomial has only complex roots.
\end{example}

\begin{example}
The matrix
\[
A = \begin{pmatrix} 2 & 0 \\ 2 & 1
\end{pmatrix}.
\]
This matrix, considered earlier, is diagonalizable because it has two distinct eigenvalues, $\lambda = 1$ and $\lambda = 2$.
\end{example}

\begin{example}
We investigate whether
\[
A = \begin{pmatrix} -1 & 0 & 1 \\ 1 & -1 & 0 \\ 0 & 1 & 1
\end{pmatrix}.
\]
This is diagonalizable. Expanding the characteristic polynomial along a convenient row or column yields
\[
\begin{aligned}
p(\lambda)&=\operatorname{Det}(\lambda I-A)\\
&=\begin{vmatrix}
\lambda+1&0&-1\\
-1&\lambda+1&0\\
0&-1&\lambda-1
\end{vmatrix}\\
&=(\lambda+1)^2(\lambda-1)-1\\
&=\lambda^3+\lambda^2-\lambda-2.
\end{aligned}
\]
This polynomial has no rational roots. Numerically, its roots are
\[
\begin{aligned}
\lambda_1&\approx1.205569,\\
\lambda_{2,3}&\approx-1.102785\pm0.665457i.
\end{aligned}
\]
Thus $A$ has only one real eigenvalue. For a root $\lambda$ of
$p$, solving $(A-\lambda I)\mathbf{x}=0$ gives, for example,
\[
\mathbf{x}=t\bigl(1+\lambda,\,1,\,(1+\lambda)^2\bigr).
\]
Consequently $A$ is not diagonalizable over $\mathbb{R}$, because it
does not possess three real eigenvalues. It is diagonalizable over
$\mathbb{C}$ because its three complex eigenvalues are distinct.
\end{example}

\begin{example}
We investigate whether
\[
A = \begin{pmatrix} 2 & 1 & 0 \\ 0 & 1 & -1 \\ 0 & 2 & 4
\end{pmatrix}.
\]
This is diagonalizable. A direct computation of the eigenvalues and their multiplicities yields $\lambda_1 = 2$, with algebraic multiplicity $a_{\lambda_1} = 2$ but geometric multiplicity $g_{\lambda_1} = 1$, and $\lambda_2 = 3$, with algebraic and geometric multiplicity both equal to $1$. Since $\lambda_1$ fails to be regular, being $g_{\lambda_1} = 1 < 2 = a_{\lambda_1}$, we conclude that $A$ is not diagonalizable. Explicitly, one finds
\[
\begin{gathered}
V_{\lambda_1} = \operatorname{span}((1,0,0)) \\
V_{\lambda_2} = \operatorname{span}((1,1,-2)).
\end{gathered}
\]
The second eigenspace can be checked directly:
\[
\begin{pmatrix}2&1&0\\0&1&-1\\0&2&4\end{pmatrix}
\begin{pmatrix}1\\1\\-2\end{pmatrix}
=\begin{pmatrix}3\\3\\-6\end{pmatrix}
=3\begin{pmatrix}1\\1\\-2\end{pmatrix}.
\]

\end{example}

\begin{example}
We investigate whether
\[
A = \begin{pmatrix} -3 & 5 & -1 \\ -1 & 3 & -1 \\ -3 & 3 & 1
\end{pmatrix}.
\]
This is diagonalizable. Its characteristic polynomial is
\[
\begin{aligned}
p(\lambda)&=\operatorname{Det}(\lambda I-A)\\
&=(\lambda+2)(\lambda-1)(\lambda-2).
\end{aligned}
\]
Hence the eigenvalues are $-2$, $1$, and $2$, each with algebraic
multiplicity one. Solving the three systems gives
\[
V_{-2}=\operatorname{span}\{(3,1,2)\},\qquad
V_1=\operatorname{span}\{(1,1,1)\},\qquad
V_2=\operatorname{span}\{(1,1,0)\}.
\]
The eigenvalues are distinct, so these eigenvectors are linearly
independent and $A$ is diagonalizable. Taking them as the columns of
$P$ gives
\[
D=\begin{pmatrix}-2&0&0\\0&1&0\\0&0&2\end{pmatrix},
\qquad
P=\begin{pmatrix}3&1&1\\1&1&1\\2&1&0\end{pmatrix}.
\]
Since $\det P=-2\neq0$ and $AP=PD$, one has $A=PDP^{-1}$.

\end{example}

For the matrix
\[
J=\begin{pmatrix} 1 & 2 \\ 0 & 1 \end{pmatrix}.
\]
The characteristic polynomial is
\[
\det(\lambda I-J)=(\lambda-1)^2.
\]
For $\lambda=1$,
\[
(J-I)\begin{pmatrix}x\\y\end{pmatrix}
=\begin{pmatrix}0&2\\0&0\end{pmatrix}
\begin{pmatrix}x\\y\end{pmatrix}=\vec0.
\]
This equation forces $y=0$, while $x$ is free. Hence
$V_1=\operatorname{span}\{(1,0)\}$ has dimension one, smaller than the
algebraic multiplicity two, and $J$ is not diagonalizable.

For the matrix
\[
A = \begin{pmatrix} 2 & 1 & 0 \\ 0 & 1 & -1 \\ 0 & 2 & 4
\end{pmatrix}.
\]
One computes
\[
\det(\lambda I-A)
=(\lambda-2)\begin{vmatrix}\lambda-1&1\\-2&\lambda-4\end{vmatrix}
=(\lambda-2)^2(\lambda-3).
\]
For $\lambda=2$, solving $(A-2I)\mathbf{x}=0$ gives $y=z=0$, so its
eigenspace is
\[
V_{\lambda_1} = \operatorname{span}((1,0,0))
\]
It has dimension one although the algebraic multiplicity is two. For
$\lambda=3$, the equations give $x=y$ and $z=-2y$, hence
\[
V_{\lambda_2} = \operatorname{span}((1,1,-2))
\]
This eigenspace also has dimension one. The total number of linearly independent
eigenvectors is therefore only two, confirming that $A$ is not
diagonalizable.

\subsection{Singular Values, Least Squares, and the Pseudoinverse}


The eigendecomposition studied earlier applies only to square matrices possessing a full set of linearly independent eigenvectors, and even then it treats the domain and codomain of the associated linear map as the same space. The singular value decomposition removes both restrictions: it applies to every real matrix, of any shape, and it reveals how the matrix distorts lengths and directions by relating an orthonormal basis of the domain to an orthonormal basis of the codomain.

Given $A \in M_{m,n}$, the matrix $A^{\top}A \in M_n$ is symmetric, and consequently, by the Spectral Theorem, it possesses $n$ real, nonnegative eigenvalues $\lambda_1 \geq \lambda_2 \geq \cdots \geq \lambda_n \geq 0$, together with an orthonormal basis of eigenvectors $\mathbf{v}_1, \dots, \mathbf{v}_n$ of $\mathbb{R}^n$. Nonnegativity follows from the identity
\[
\begin{aligned}
\mathbf{v}_i^{\top} A^{\top} A \mathbf{v}_i &= \|A\mathbf{v}_i\|^2 \\
&= \lambda_i \|\mathbf{v}_i\|^2 \\
&\geq 0.
\end{aligned}
\]
This observation motivates the following definition.

\begin{definition}[Singular Values]
Given $A \in M_{m,n}$, the singular values of $A$ are the numbers $\sigma_i = \lambda_i^{1/2}$, for $i = 1, \dots, n$, where $\lambda_1 \geq \cdots \geq \lambda_n \geq 0$ are the eigenvalues of $A^{\top}A$.
\end{definition}

The singular values, together with the associated eigenvectors of $A^{\top}A$, assemble into a factorization of $A$ that is valid without any assumption of squareness or diagonalizability.

\begin{theorem}[Singular Value Decomposition]
Given $A \in M_{m,n}$ with $\operatorname{Rank}(A) = r$, there exist an orthogonal matrix $U \in M_m$, an orthogonal matrix $V \in M_n$, and a matrix $\Sigma \in M_{m,n}$ whose only nonzero entries are $\Sigma_{ii} = \sigma_i$ for $i = 1, \dots, r$, arranged so that $\sigma_1 \geq \cdots \geq \sigma_r > 0$, such that $A = U \Sigma V^{\top}$.
\end{theorem}

\begin{proof}
Let $\mathbf{v}_1, \dots, \mathbf{v}_n$ be an orthonormal basis of eigenvectors of $A^{\top}A$, ordered so that the corresponding eigenvalues satisfy $\lambda_1 \geq \cdots \geq \lambda_n \geq 0$, and set
\[
\sigma_i = \lambda_i^{1/2}.
\]
The rank identity gives
\[
\begin{aligned}
\operatorname{Rank}(A^{\top}A) &= \operatorname{Rank}(A) \\
&= r,
\end{aligned}
\]

Therefore, exactly $r$ of the $\lambda_i$ are nonzero, say $\lambda_1, \dots, \lambda_r$. For $i = 1, \dots, r$, define $\mathbf{u}_i = \sigma_i^{-1} A \mathbf{v}_i$; these vectors are orthonormal, since
\[
\begin{aligned}
\mathbf{u}_i \cdot \mathbf{u}_j &= (\sigma_i \sigma_j)^{-1} \mathbf{v}_i^{\top} A^{\top} A \mathbf{v}_j \\
&= (\sigma_i \sigma_j)^{-1} \lambda_j \mathbf{v}_i \cdot \mathbf{v}_j.
\end{aligned}
\]
This equals $1$ if $i=j$ and $0$ otherwise by orthonormality of the $\mathbf{v}_i$. Complete $\mathbf{u}_1, \dots, \mathbf{u}_r$ to an orthonormal basis $\mathbf{u}_1, \dots, \mathbf{u}_m$ of $\mathbb{R}^m$, and let $U = [\mathbf{u}_1 | \cdots | \mathbf{u}_m]$ and $V = [\mathbf{v}_1 | \cdots | \mathbf{v}_n]$, both orthogonal by construction, since their columns form orthonormal bases. For $i \leq r$, one has $A\mathbf{v}_i = \sigma_i \mathbf{u}_i$ by definition of $\mathbf{u}_i$; for $i > r$, one has $\lambda_i = 0$, so
\[
\begin{aligned}
\|A\mathbf{v}_i\|^2 &= \mathbf{v}_i^{\top}A^{\top}A\mathbf{v}_i \\
&= \lambda_i \|\mathbf{v}_i\|^2 \\
&= 0.
\end{aligned}
\]
Therefore $A\mathbf{v}_i = \vec{0}$. Consequently $AV = U\Sigma$, where $\Sigma$ has entries $\sigma_i$ on its first $r$ diagonal positions and zero elsewhere. Since $V$ is orthogonal, it is invertible with $V^{-1} = V^{\top}$, and one obtains
\[
\begin{aligned}
A &= AVV^{\top} \\
&= U\Sigma V^{\top}.
\end{aligned}
\]

\end{proof}

The rank of $A$ therefore equals the number of nonzero singular values, and the decomposition exposes the four fundamental subspaces of $A$ directly: the columns $\mathbf{u}_1, \dots, \mathbf{u}_r$ of $U$ form a basis of $\operatorname{Col}(A)$, the columns $\mathbf{u}_{r+1}, \dots, \mathbf{u}_m$ form a basis of $\operatorname{Ker}(A^{\top})$, the columns $\mathbf{v}_1, \dots, \mathbf{v}_r$ form a basis of $\operatorname{Col}(A^{\top})$, and the columns $\mathbf{v}_{r+1}, \dots, \mathbf{v}_n$ form a basis of $\operatorname{Ker}(A)$. This last fact deserves justification: since $A\mathbf{v}_i = \vec{0}$ for $i > r$, the span of $\mathbf{v}_{r+1}, \dots, \mathbf{v}_n$ is certainly contained in $\operatorname{Ker}(A)$; and since these $n-r$ vectors are linearly independent, while the Nullity--Rank Theorem gives
\[
\begin{aligned}
\operatorname{Dim}(\operatorname{Ker}(A)) &= n - \operatorname{Rank}(A) \\
&= n-r.
\end{aligned}
\]
The inclusion is in fact an equality. The uniqueness properties of the decomposition are worth recording explicitly. The singular values $\sigma_1 \geq \cdots \geq \sigma_r$ are always uniquely determined by $A$, being the square roots of the nonzero eigenvalues of $A^{\top}A$, themselves an intrinsic invariant of $A$; the vectors $\mathbf{u}_i$ and $\mathbf{v}_i$, however, are unique only up to simultaneous rotation within eigenspaces of repeated singular values, since any orthonormal basis of such an eigenspace serves equally well in the construction above.

\begin{example}
Consider the matrix
\[
A = \begin{pmatrix} 2 & 0 \\ 0 & -3 \end{pmatrix}.
\]
The following product is already diagonal:
\[
A^{\top}A = \begin{pmatrix} 4 & 0 \\ 0 & 9 \end{pmatrix}
\]

Its eigenvalues are $\lambda_1 = 9$, $\lambda_2 = 4$, with eigenvectors $\mathbf{v}_1 = (0,1)$, $\mathbf{v}_2 = (1,0)$, giving singular values $\sigma_1 = 3$, $\sigma_2 = 2$. Computing
\[
\begin{gathered}
\mathbf{u}_1 = \sigma_1^{-1}A\mathbf{v}_1 = (0,-1) \\
\mathbf{u}_2 = \sigma_2^{-1}A\mathbf{v}_2 = (1,0).
\end{gathered}
\]
One obtains the decomposition $A = U\Sigma V^{\top}$ with
\[
U = \begin{pmatrix} 0 & 1 \\ -1 & 0 \end{pmatrix},
\qquad
\Sigma = \begin{pmatrix} 3 & 0 \\ 0 & 2 \end{pmatrix},
\qquad
V = \begin{pmatrix} 0 & 1 \\ 1 & 0 \end{pmatrix}.
\]
As a verification,
\[
\begin{aligned}
U\Sigma V^{\top}
&= \begin{pmatrix} 0 & 1 \\ -1 & 0 \end{pmatrix}\begin{pmatrix} 3 & 0 \\ 0 & 2 \end{pmatrix}\begin{pmatrix} 0 & 1 \\ 1 & 0 \end{pmatrix} \\
&= \begin{pmatrix} 0 & 2 \\ -3 & 0 \end{pmatrix}\begin{pmatrix} 0 & 1 \\ 1 & 0 \end{pmatrix} \\
&= \begin{pmatrix} 2 & 0 \\ 0 & -3 \end{pmatrix} = A,
\end{aligned}
\]
as required.
\end{example}

Beyond its structural role, the singular value decomposition provides the optimal low-rank approximation of a matrix in a precise sense, a fact that underlies its widespread use in data compression and dimensionality reduction. To state the result, one measures the size of a matrix using the operator norm,
\[
\|B\| = \max_{\|\mathbf{x}\|=1} \|B\mathbf{x}\|,
\]
and one first records an elementary but essential fact.

\begin{theorem}
If $A \in M_{m,n}$ has singular values $\sigma_1 \geq \cdots \geq \sigma_r > 0$, then $\|A\| = \sigma_1$.
\end{theorem}

\begin{proof}
Write $A = U\Sigma V^{\top}$ as above. Since $U$ and $V$ are orthogonal, they preserve Euclidean norms, so
\[
\begin{aligned}
\|A\mathbf{x}\| &= \|U\Sigma V^{\top}\mathbf{x}\| \\
&= \|\Sigma V^{\top}\mathbf{x}\|.
\end{aligned}
\]
This holds for every $\mathbf{x}$, and as $\mathbf{x}$ ranges over the unit sphere, so does $\mathbf{y} = V^{\top}\mathbf{x}$, by orthogonality of $V$. It therefore suffices to compute
\[
\max_{\|\mathbf{y}\|=1} \|\Sigma \mathbf{y}\|.
\]
Multiplication by $\Sigma$ gives
\[
\Sigma \mathbf{y} = (\sigma_1 y_1, \dots, \sigma_r y_r, 0, \dots, 0),
\]
The squared norm therefore satisfies
\[
\begin{aligned}
\|\Sigma\mathbf{y}\|^2 &= \sum_{i=1}^{r} \sigma_i^2 y_i^2 \\
&\leq \sigma_1^2 \sum_{i=1}^{r} y_i^2 \\
&\leq \sigma_1^2.
\end{aligned}
\]
Here equality when $\mathbf{y} = \mathbf{e}_1$, so the maximum equals $\sigma_1$.
\end{proof}

\begin{theorem}[Eckart--Young Theorem]
Let $A = U\Sigma V^{\top}$ be the singular value decomposition of $A \in M_{m,n}$, with singular values $\sigma_1 \geq \cdots \geq \sigma_r > 0$, and for $k < r$, let
\[
A_k = \sum_{i=1}^{k} \sigma_i \mathbf{u}_i \mathbf{v}_i^{\top}.
\]

Then $\operatorname{Rank}(A_k) \leq k$, and $A_k$ minimizes $\|A - B\|$ among all $B \in M_{m,n}$ of rank at most $k$; moreover $\|A - A_k\| = \sigma_{k+1}$.
\end{theorem}

\begin{proof}
Since $A_k$ is a sum of $k$ rank-one matrices $\mathbf{u}_i\mathbf{v}_i^{\top}$, $\operatorname{Rank}(A_k) \leq k$. Writing \[
A - A_k = \sum_{i=k+1}^{r} \sigma_i \mathbf{u}_i \mathbf{v}_i^{\top},
\] this is itself the singular value decomposition of $A-A_k$, restricted to indices $k+1, \dots, r$, so by the previous theorem $\|A - A_k\| = \sigma_{k+1}$.

It remains to show that no matrix of rank at most $k$ does better. Let $B \in M_{m,n}$ with $\operatorname{Rank}(B) \leq k$, so that $\operatorname{Dim}(\operatorname{Ker}(B)) \geq n-k$ by the Nullity--Rank Theorem. Consider the subspace $W = \operatorname{span}(\mathbf{v}_1, \dots, \mathbf{v}_{k+1})$, of dimension $k+1$; since $\operatorname{Dim}(\operatorname{Ker}(B)) + \operatorname{Dim}(W) \geq (n-k) + (k+1) = n+1 > n$, the two subspaces must intersect nontrivially, so there exists a unit vector $\mathbf{w} \in W \cap \operatorname{Ker}(B)$. Expand this vector in the basis of $W$:
\[
\mathbf{w} = \sum_{i=1}^{k+1} c_i \mathbf{v}_i.
\]
The unit-norm constraint is
\[
\sum_i c_i^2 = 1.
\]
Using orthonormality of the $\mathbf{v}_i$ and $\mathbf{u}_i$, we obtain
\[
\begin{aligned}
\|(A-B)\mathbf{w}\|^2
&= \|A\mathbf{w}\|^2
= \left\| \sum_{i=1}^{k+1} c_i \sigma_i \mathbf{u}_i \right\|^2 \\
&= \sum_{i=1}^{k+1} c_i^2 \sigma_i^2
\geq \sigma_{k+1}^2 \sum_{i=1}^{k+1} c_i^2
= \sigma_{k+1}^2.
\end{aligned}
\]
Here the inequality uses $\sigma_i \geq \sigma_{k+1}$ for $i \leq k+1$. Since $\mathbf{w}$ is a unit vector, this shows
\[
\begin{aligned}
\|A-B\| &\geq \sigma_{k+1} \\
&= \|A-A_k\|.
\end{aligned}
\]
This establishes the optimality of $A_k$.
\end{proof}

Figure~\ref{fig:ch10-svd} shows the geometric action of the singular value decomposition on the unit circle.

\begin{figure}[!htbp]
\centering
\begin{tikzpicture}[scale=1.5,>=Stealth]
  \begin{scope}
    \draw[->] (-1.6,0) -- (1.6,0);
    \draw[->] (0,-1.6) -- (0,1.6);
    \draw[figblue,thick] (0,0) circle (1);
    \draw[vecB] (0,0) -- (30:1) node[above right,font=\scriptsize]{$\bv_1$};
    \draw[vecC] (0,0) -- (120:1) node[above left,font=\scriptsize]{$\bv_2$};
  \end{scope}
  \draw[->,very thick] (2.0,0.7) -- (3.1,0.7) node[midway,above,font=\scriptsize]{$A=U\Sigma V^\top$};
  \begin{scope}[xshift=5.6cm]
    \draw[->] (-1.9,0) -- (1.9,0);
    \draw[->] (0,-1.3) -- (0,1.3);
    \draw[figblue,thick,rotate=15] (0,0) ellipse (1.6 and 0.7);
    \draw[vecB] (0,0) -- (15:1.6) node[right,font=\scriptsize]{$\sigma_1 \bu_1$};
    \draw[vecC] (0,0) -- (105:0.7) node[above,font=\scriptsize]{$\sigma_2 \bu_2$};
  \end{scope}
\end{tikzpicture}
\caption[Singular Value Decomposition (SVD)]{Left: the unit circle and the orthonormal domain basis $\{\bv_1,\bv_2\}$. Right: the image ellipse, with orthogonal semi-axis vectors $\sigma_1\bu_1$ and $\sigma_2\bu_2$. The map $A=U\Sigma V^\top$ rotates or reflects, scales along perpendicular directions, and rotates or reflects again; the semi-axis lengths are $\sigma_1,\sigma_2$.}
\label{fig:ch10-svd}
\end{figure}

The inverse of a matrix, as defined earlier, exists only for nonsingular square matrices, yet many of the systems arising in applications are either rectangular or singular, and one still wishes to associate to them a canonical generalized inverse. The singular value decomposition supplies exactly such a generalization.

\begin{definition}[Pseudoinverse]
Given $A \in M_{m,n}$ with singular value decomposition $A = U\Sigma V^{\top}$, the Moore--Penrose pseudoinverse of $A$ is the matrix $A^{+} = V \Sigma^{+} U^{\top} \in M_{n,m}$, where $\Sigma^{+} \in M_{n,m}$ is obtained from $\Sigma^{\top}$ by replacing every nonzero diagonal entry $\sigma_i$ with $\sigma_i^{-1}$.
\end{definition}

When $A$ is square and nonsingular, every singular value is nonzero and $\Sigma$ is itself square and invertible, with $\Sigma^{+} = \Sigma^{-1}$, so that
\[
\begin{aligned}
A^{+} &= V\Sigma^{-1}U^{\top} \\
&= (U\Sigma V^{\top})^{-1} \\
&= A^{-1};
\end{aligned}
\]
the pseudoinverse indeed extends the ordinary inverse. Its defining property, independent of the particular decomposition used to construct it, is captured by the following characterization, which we now prove directly from the definition above.

\begin{theorem}[Penrose Conditions]
The pseudoinverse $A^{+}$ satisfies the four conditions
\[
AA^{+}A = A, \qquad A^{+}AA^{+} = A^{+}, \qquad (AA^{+})^{\top} = AA^{+}, \qquad (A^{+}A)^{\top} = A^{+}A.
\]
Moreover, it is the unique matrix in $M_{n,m}$ satisfying all four simultaneously.
\end{theorem}

\begin{proof}
Write $A = U\Sigma V^{\top}$ and $A^{+} = V\Sigma^{+}U^{\top}$. A direct computation, using $U^{\top}U = I_m$ and $V^{\top}V = I_n$, gives
\[
\begin{aligned}
AA^{+} &= U\Sigma V^{\top}V\Sigma^{+}U^{\top} \\
&= U(\Sigma \Sigma^{+})U^{\top}.
\end{aligned}
\]
Since $\Sigma\Sigma^{+} \in M_m$ is diagonal with entries $1$ in the first $r$ positions and $0$ elsewhere, it is symmetric, and hence $AA^{+} = U(\Sigma\Sigma^{+})U^{\top}$ is symmetric as a product of a symmetric matrix conjugated by an orthogonal one, establishing $(AA^{+})^{\top} = AA^{+}$. The identical argument, with the roles of $U$ and $V$ exchanged, gives $A^{+}A = V(\Sigma^{+}\Sigma)V^{\top}$ symmetric, establishing $(A^{+}A)^{\top} = A^{+}A$. For the first two conditions, one checks directly that $\Sigma\Sigma^{+}\Sigma = \Sigma$ and $\Sigma^{+}\Sigma\Sigma^{+} = \Sigma^{+}$, since both sides act identically on the first $r$ standard basis vectors, where $\Sigma$ and $\Sigma^{+}$ are mutually inverse scalars, and both sides vanish on the remaining basis vectors, where $\Sigma$ has zero rows or columns; substituting these identities into
\[
\begin{aligned}
AA^{+}A &= U\Sigma V^{\top}V\Sigma^{+}U^{\top}U\Sigma V^{\top} \\
&= U(\Sigma\Sigma^{+}\Sigma)V^{\top} \\
&= U\Sigma V^{\top} \\
&= A.
\end{aligned}
\]
Moreover, analogously for $A^{+}AA^{+} = A^{+}$. For uniqueness, suppose $B \in M_{n,m}$ also satisfies all four conditions. One shows $B = A^{+}$ by a sequence of substitutions exploiting the four identities for both $A^{+}$ and $B$; we omit the lengthy but entirely mechanical verification, which proceeds by repeatedly inserting
\[
\begin{aligned}
AA^{+}A &= A \\
&= ABA.
\end{aligned}
\]
The symmetry conditions into the expression for $B$ until it is transformed into $A^{+}$.
\end{proof}

The practical importance of the pseudoinverse stems from its role in solving linear systems that admit no exact solution, or that admit infinitely many.

\begin{theorem}[Least Squares via the Pseudoinverse]
Given $A \in M_{m,n}$ and $\mathbf{b} \in \mathbb{R}^m$, the vector $\mathbf{x}^{*} = A^{+}\mathbf{b}$ minimizes $\|A\mathbf{x} - \mathbf{b}\|$ over all $\mathbf{x} \in \mathbb{R}^n$, and among all such minimizers it is the one of smallest Euclidean norm.
\end{theorem}

\begin{proof}
Write $A = U\Sigma V^{\top}$, and set $\mathbf{y} = V^{\top}\mathbf{x}$ and $\mathbf{c} = U^{\top}\mathbf{b}$; since $U$ and $V$ are orthogonal,
\[
\begin{aligned}
\|A\mathbf{x}-\mathbf{b}\| &= \|U\Sigma V^{\top}\mathbf{x} - \mathbf{b}\| \\
&= \|\Sigma \mathbf{y} - \mathbf{c}\|.
\end{aligned}
\]
Moreover, $\mathbf{x}$ ranges over $\mathbb{R}^n$ exactly as $\mathbf{y}$ does. Multiplication by $\Sigma$ gives
\[
\Sigma \mathbf{y} = (\sigma_1 y_1, \dots, \sigma_r y_r, 0, \dots, 0).
\]
The squared residual is therefore
\[
\|\Sigma\mathbf{y} - \mathbf{c}\|^2 = \sum_{i=1}^r (\sigma_i y_i - c_i)^2 + \sum_{i=r+1}^{m} c_i^2.
\]

Over all choices of $y_1, \dots, y_n$, it is minimized precisely by taking $y_i = c_i/\sigma_i$ for $i \leq r$, which eliminates the first sum entirely, while $y_{r+1}, \dots, y_n$ remain unconstrained by the minimization; the second sum is independent of $\mathbf{y}$ and represents the unavoidable residual. Among all minimizers, the one of smallest norm is obtained by setting the free components $y_{r+1}, \dots, y_n$ to zero, since $\|\mathbf{y}\|^2 = \|\mathbf{x}\|^2$ by orthogonality of $V$, and this choice corresponds precisely to
\[
\mathbf{y} = \Sigma^{+}\mathbf{c} = \Sigma^{+}U^{\top}\mathbf{b}.
\]
Thus, the minimizer has the equivalent representations
\[
\mathbf{x}^{*} = V\mathbf{y}
= V\Sigma^{+}U^{\top}\mathbf{b}
= A^{+}\mathbf{b}.
\]

\end{proof}

Figure~\ref{fig:atlas-least-squares-projection} shows least squares as orthogonal projection onto the column space.

\begin{figure}[!htbp]
\centering
\begin{tikzpicture}[scale=1.45,>=Stealth]
\draw[black!45] (-0.6,-0.3) -- (4.1,2.05) node[right] {$\operatorname{Col}A$};
\draw[vecA] (0,0) -- (1,3) node[above] {$b$};
\draw[vecB] (0,0) -- (2,1) node[below right] {$A\hat{x}$};
\draw[vecC] (2,1) -- (1,3) node[midway,right] {$r$};
\draw[black!60] (2.25,1.125) -- (2.125,1.375) -- (1.875,1.25);
\node[below] at (0,0) {$0$};
\end{tikzpicture}
\caption{The vector $b$ is decomposed into its projection $A\hat{x}$ onto the column space and a perpendicular residual $r=b-A\hat{x}$. Orthogonality gives the normal equations $A^{\top}r=0$. Here the column space is represented by a line in the plane.}
\label{fig:atlas-least-squares-projection}
\end{figure}

\begin{example}
Consider the matrix
\[
A = \begin{pmatrix} 1 & 0 \\ 0 & 0 \end{pmatrix}.
\]
Its singular values are $1$ and $0$, and one may take
$U=V=I$ and $\Sigma=A$. Therefore $\Sigma^+=A$ and hence $A^+=A$.
For $\mathbf{b}=(1,1)$, $\mathbf{x}^{*} = A^{+}\mathbf{b} = (1,0)$ is the minimum-norm least-squares solution. Indeed, for
$\mathbf{x}=(x_1,x_2)$ one has
$\|A\mathbf{x}-\mathbf{b}\|^2=(x_1-1)^2+1$, which is minimized when
$x_1=1$. Among these minimizers, Euclidean norm is smallest for $x_2=0$,
giving $\mathbf{x}^*=(1,0)$.
\end{example}

\newpage
\chapter[Multivariable Calculus, Vector Operators, and Tensors]{Multivariable Differential Calculus, Vector Operators and Tensor Theory}

\section{Multivariable Functions and Partial Derivatives}

\subsection{Definition, Graphs, and Level Curves}

A function $f$ of two variables is a rule that assigns to each ordered pair of real numbers $(x,y)$ belonging to a set $D \subset \mathbb{R}^2$ a unique real number, denoted by $f(x,y)$; symbolically,
\[
f : D \subset \mathbb{R}^2 \to \mathbb{R}, \qquad (x,y) \mapsto f(x,y).
\]

The domain $D$ is a subset of the $xy$-plane, and consists of the set of all admissible inputs of the function. Several elementary examples illustrate this notion, such as
\[
\begin{gathered}
z = (1-x-y)^{1/2}/(x-y), \\[0.5em]
z = \ln(y^2-x), \\[0.5em]
z = 4 - x^2 - y^2,
\end{gathered}
\]
each defined on a suitable domain determined by the requirement that the corresponding algebraic expression be well defined.

One method for visualizing the behavior of a function of two variables is to consider its graph. If $f$ is a function of two variables with domain $D$, the graph of $f$ is defined as the set of all points $(x,y,z) \in \mathbb{R}^3$ such that $z = f(x,y)$ and $(x,y) \in D$; this graph typically takes the form of a surface in three-dimensional space.

A second, complementary method for visualizing such functions, borrowed from the practice of mapmakers, relies on two-dimensional maps annotated with level curves.

\begin{definition}[Level Curve]
The level curves of a function $f$ of two variables are the curves in the $xy$-plane with equations $f(x,y) = k$, where $k$ is a constant belonging to the range of $f$.
\end{definition}

As illustrations of this notion, one may sketch the level-curve map of the function
\[
z = 6 - 3x - 2y,
\]
whose level curves form a family of parallel straight lines, or of the function
\[
g(x,y) = 9 - x^2 - y^2,
\]
whose level curves form a family of concentric circles centered at the origin.

\subsection{Limits of Functions of Two Variables}

Let $f$ be a function of two variables whose domain $D$ includes a neighborhood $U(a,b)$ of a given point $(a,b)$, that is, a disc centered at $(a,b)$.

The notion of limit developed earlier for functions of a single real variable extends naturally to this two-dimensional setting, with the absolute distance $|x-a|$ replaced by the Euclidean distance between the points $(x,y)$ and $(a,b)$.

\begin{definition}
We write \[
\lim_{(x,y) \to (a,b)} f(x,y) = L,
\] and say that the limit of $f(x,y)$ as $(x,y)$ approaches $(a,b)$ is $L$, if for every $\varepsilon > 0$ there exists $\delta > 0$ such that
\[
0 < \sqrt{(x-a)^2 + (y-b)^2} < \delta \implies |f(x,y) - L| < \varepsilon.
\]

Equivalently, the distance between $f(x,y)$ and $L$ can be made arbitrarily small by making the distance from $(x,y)$ to $(a,b)$ sufficiently small, though not equal to zero.
\end{definition}

\begin{remark}
This definition refers exclusively to the Euclidean distance between $(x,y)$ and $(a,b)$, and makes no reference whatsoever to the particular direction along which $(x,y)$ approaches $(a,b)$.

Consequently, if the limit exists, then $f(x,y)$ must approach the very same value $L$ regardless of the specific path along which $(x,y)$ approaches $(a,b)$; this observation provides a standard technique for proving the nonexistence of a two-variable limit, namely by exhibiting two distinct paths of approach along which $f$ tends to two different values. As an illustration of this phenomenon, one may investigate whether the limit \[
\lim_{(x,y) \to (0,0)} f(x,y)
\] exists for the function $f(x,y) = xy/(x^2+y^2)$, and discover, by comparing the limiting values obtained along the two coordinate axes and along the line $y=x$, that no common limiting value exists.
\end{remark}

Figure~\ref{fig:atlas-path-dependent-limit} compares two paths approaching the same point with incompatible limiting values.

\begin{figure}[!htbp]
\centering
\begin{tikzpicture}[>=Stealth]
\begin{axis}[
  width=12.5cm,height=6.2cm,
  axis lines=left,
  ticklabel style={font=\small},label style={font=\small},
  xlabel={$t$},ylabel={$f$},
  xmin=-1.2,xmax=1.2,ymin=-0.15,ymax=0.72,
  xtick={-1,0,1},ytick={0,0.5},
  legend style={at={(0.5,1.02)},anchor=south,legend columns=2,
    font=\scriptsize,draw=black!20,fill=white}
]
\addplot[figblue,very thick] coordinates {(-1,0)(1,0)};
\addlegendentry{$f(t,0)$}
\addplot[figred,very thick] coordinates {(-1,0.5)(1,0.5)};
\addlegendentry{$f(t,t)$}
\draw[black!35,densely dotted] (axis cs:0,-0.15) -- (axis cs:0,0.62);
\addplot[only marks,mark=*,mark options={fill=white,draw=figblue},mark size=3pt,forget plot] coordinates {(0,0)};
\addplot[only marks,mark=*,mark options={fill=white,draw=figred},mark size=3pt,forget plot] coordinates {(0,0.5)};
\node[figblue,below right,font=\scriptsize] at (axis cs:0.04,0) {$0$};
\node[figred,above right,font=\scriptsize] at (axis cs:0.04,0.5) {$1/2$};
\end{axis}
\end{tikzpicture}
\caption{For $f(x,y)=xy/(x^2+y^2)$ away from the origin, the path $(t,0)$ gives value zero and the path $(t,t)$ gives value $1/2$ whenever $t\ne0$. The holes mark the excluded parameter value. These distinct limits rule out a two-variable limit at the origin.}
\label{fig:atlas-path-dependent-limit}
\end{figure}

\subsection{Continuity of Functions of Two Variables}

\begin{definition}
A function $f$ of two variables is called continuous at $(a,b)$ if
\[
\lim_{(x,y) \to (a,b)} f(x,y) = f(a,b).
\]
We say that $f$ is continuous on $D$ if $f$ is continuous at every point of $D$.
\end{definition}

Using the algebraic properties of limits, entirely analogous to those established for functions of a single variable, one readily verifies that sums, differences, products, and quotients of continuous functions of two variables remain continuous on their respective domains. Consequently, any polynomial function of two variables, such as $f(x,y) = x^4 + 5x^3y^2 - 6xy^4 + 7y^2$, is continuous throughout $\mathbb{R}^2$. Likewise, since a rational function is a ratio of two polynomials, it is continuous wherever its denominator does not vanish; the function $f(x,y) = x^2y/(x^2+y^2)$ provides a standard example of such a rational function, continuous throughout $\mathbb{R}^2 \setminus \{(0,0)\}$.

\subsection{Partial Derivatives: Definition and Computation}


Let $f(x,y)$ be a function of two variables, defined on a neighborhood $U(a,b)$ of a point $(a,b)$.

Suppose that we allow only the variable $x$ to vary, while keeping $y$ fixed at the value $b$; this amounts to considering the function of a single variable $g(x) = f(x,b)$.

If $g$ is differentiable at $x=a$, we say that $f$ has a partial derivative with respect to $x$ at $(a,b)$, denoted $f_x(a,b)$ or $\partial f/\partial x \, (a,b)$, defined by $f_x(a,b) = g'(a)$; equivalently, unwinding the definition of ordinary derivative,
\[
\begin{aligned}
f_x(a,b) &= \lim_{h \to 0} \frac{g(a+h) - g(a)}{h} \\[0.5em]
&= \lim_{h \to 0} \frac{f(a+h,b) - f(a,b)}{h}.
\end{aligned}
\]

An entirely analogous construction, obtained by keeping $x$ fixed at the value $a$ and letting $y$ vary, defines the partial derivative of $f$ with respect to $y$ at $(a,b)$, denoted $f_y(a,b)$ or $\partial f/\partial y \, (a,b)$; setting $G(y) = f(a,y)$, one has $f_y(a,b) = G'(b)$, that is,
\[
f_y(a,b) = \lim_{k \to 0} \frac{f(a,b+k) - f(a,b)}{k}.
\]

If the point $(a,b)$ is allowed to vary throughout the domain, the partial derivatives $f_x$ and $f_y$ themselves become functions of the two variables $x$ and $y$,
\[
\begin{aligned}
f_x(x,y) &= \lim_{h \to 0} \frac{f(x+h,y) - f(x,y)}{h}, \\[0.5em]
f_y(x,y) &= \lim_{k \to 0} \frac{f(x,y+k) - f(x,y)}{k}.
\end{aligned}
\]

In practice, partial derivatives are computed according to the following simple rule: to find $f_x$, one regards $y$ as a constant and differentiates the resulting expression with respect to $x$ using the ordinary rules of single-variable differentiation; symmetrically, to find $f_y$, one regards $x$ as a constant and differentiates with respect to $y$. For $f(x,y) = x^3 + x^2 y^3 - 2y^2$. holding $y$ constant gives $f_x(x,y)=3x^2+2xy^3$, whereas holding $x$ constant gives
\[
f_y(x,y)=3x^2y^2-4y.
\]
Differentiating once more also gives
\[
f_{xx}=6x+2y^3,\quad
f_{xy}=f_{yx}=6xy^2,\quad
f_{yy}=6x^2y-4,
\]
which explicitly illustrates Clairaut's equality of the mixed derivatives.

If $f$ is a function of two variables, its partial derivatives $f_x$ and $f_y$ are themselves functions of two variables, and one may accordingly consider their own partial derivatives, namely $(f_x)_x$, $(f_x)_y$, $(f_y)_x$, and $(f_y)_y$, collectively called the second-order partial derivatives of $f$. Writing $z = f(x,y)$, one adopts the standard notation $f_{xy} = \partial^2 f / \partial y\, \partial x$ to indicate that differentiation is performed first with respect to $x$, and subsequently with respect to $y$; the notation $f_{yx}$ indicates the reverse order of differentiation. Although these two mixed partial derivatives are computed in opposite orders, they in fact coincide under mild regularity assumptions, as the following classical result shows.

\begin{theorem}[Clairaut's Theorem]
Suppose $f$ is defined on a neighborhood $U(a,b)$. If the functions $f_{xy}$ and $f_{yx}$ are both continuous on $U(a,b)$, then
\[
f_{xy}(a,b) = f_{yx}(a,b).
\]

\end{theorem}

Partial derivatives occur pervasively in the formulation of partial differential equations expressing fundamental physical laws, and it is instructive to record a few celebrated examples. The Laplace equation, $\Delta u = 0$, where $\Delta = \partial^2/\partial x^2 + \partial^2/\partial y^2$ denotes the two-dimensional Laplacian operator already introduced in the study of vector calculus, has solutions called harmonic functions, which play a central role in problems of heat conduction, fluid flow, and electric potential. If $u(t,x)$ represents the displacement of a vibrating violin string at time $t$ and at a distance $x$ from one end of the string, then $u$ satisfies the wave equation, $\partial^2 u/\partial t^2 = a^2 \, \partial^2 u/\partial x^2$, for a suitable constant $a$ related to the physical properties of the string. Finally, the celebrated Navier--Stokes equation, governing the velocity field $\mathbf{u}$ of an incompressible fluid, takes the form
\[
\partial \mathbf{u}/\partial t + (\mathbf{u} \cdot \nabla)\mathbf{u} - \nu \Delta \mathbf{u} = -\nabla w + \mathbf{g}.
\]
Here $\nu$, $w$, and $\mathbf{g}$ denote, respectively, the viscosity, pressure, and external force density of the fluid.

\section{Scalar and Vector Fields: Differential Operators}
\label{sec:vector-operators}

Vector calculus is a fundamental tool in mathematical analysis, differential geometry, and mathematical physics. It extends the concepts of differential calculus, already developed for functions of a single real variable, to the setting of scalar and vector fields defined on regions of three-dimensional space.

\subsection{Scalar and Vector Fields and First-Order Operators}

Let $\Omega \subseteq \mathbb{R}^3$ be a given open set. Two distinct types of fields naturally arise on such a domain, and it is useful to introduce a unified notation before distinguishing between them. A scalar field is a function $f : \Omega \to \mathbb{R}$ that assigns a single real number to every point $\mathbf{x} = (x,y,z) \in \Omega$; a vector field, by contrast, is a function $\mathbf{F} : \Omega \to \mathbb{R}^3$ that assigns an entire vector, rather than a single number, to every point of $\Omega$, and is accordingly written in components as
\[
\mathbf{F}(x,y,z) = P(x,y,z) \, \mathbf{i} + Q(x,y,z) \, \mathbf{j} + R(x,y,z) \, \mathbf{k} = \begin{pmatrix} P(x,y,z) \\ Q(x,y,z) \\ R(x,y,z) \end{pmatrix}.
\]

In order to express directional derivatives and the various differential operations acting on scalar and vector fields in a compact and systematic notation, it is convenient to introduce the symbolic Nabla operator,
\[
\nabla = \mathbf{i} \frac{\partial}{\partial x} + \mathbf{j} \frac{\partial}{\partial y} + \mathbf{k} \frac{\partial}{\partial z} = \begin{pmatrix} \dfrac{\partial}{\partial x} \\[2mm] \dfrac{\partial}{\partial y} \\[2mm] \dfrac{\partial}{\partial z} \end{pmatrix}.
\]

Although $\nabla$ is not itself a vector in the ordinary sense, but rather a formal symbol representing a vector of partial differentiation operators, it can be manipulated according to many of the familiar rules of vector algebra, and this formal analogy underlies the definitions of the gradient, divergence, and curl introduced in the following subsection.


Three fundamental first-order differential operators act on scalar and vector fields, each capturing a distinct geometric aspect of how the field varies from point to point; we now introduce each of them in turn, together with its defining formula in terms of the Nabla operator.

The gradient of a differentiable scalar field $f$ is a vector field that points, at every point of $\Omega$, in the direction of the maximum rate of increase of $f$, and whose magnitude equals that maximum rate of increase. It is defined by applying the Nabla operator to $f$, yielding
\[
\begin{aligned}
\operatorname{grad} f &= \nabla f \\[0.5em]
&= \frac{\partial f}{\partial x} \, \mathbf{i} + \frac{\partial f}{\partial y} \, \mathbf{j} + \frac{\partial f}{\partial z} \, \mathbf{k}.
\end{aligned}
\]

The divergence, by contrast, acts on a vector field $\mathbf{F}$ rather than on a scalar field, and produces a scalar rather than a vector; it measures the magnitude of the source or sink of the vector field at a given point, that is, the extent to which the field appears to originate from or converge toward that point. It is defined, formally, as the dot product of $\nabla$ with $\mathbf{F}$,
\[
\begin{aligned}
\operatorname{div} \mathbf{F} &= \nabla \cdot \mathbf{F} \\[0.5em]
&= \frac{\partial P}{\partial x} + \frac{\partial Q}{\partial y} + \frac{\partial R}{\partial z}.
\end{aligned}
\]

The curl, finally, again acts on a vector field, but unlike the divergence it returns another vector field rather than a scalar, and it measures the tendency of the field to rotate about a given point. It is defined, formally, as the cross product of $\nabla$ with $\mathbf{F}$, evaluated symbolically as the determinant
\[
\operatorname{curl} \mathbf{F} = \nabla \times \mathbf{F} = \det \begin{pmatrix} \mathbf{i} & \mathbf{j} & \mathbf{k} \\[1mm] \dfrac{\partial}{\partial x} & \dfrac{\partial}{\partial y} & \dfrac{\partial}{\partial z} \\[2mm] P & Q & R \end{pmatrix}.
\]

\subsection{Second-Order Differential Operators}

Beyond the three first-order operators just introduced, it is natural to consider differential operators obtained by applying two differentiations in succession; the most important such operator, arising throughout mathematical physics, is the Laplacian.

The Laplacian is the second-order differential operator obtained by taking
the divergence of the gradient of a scalar field. Some literature writes
it as $\nabla^2 f$; throughout this volume we use $\Delta f$ for the
Laplacian and reserve $\nabla^2 f$ for the Hessian matrix. Thus
\[
\begin{aligned}
\Delta f &= \nabla \cdot (\nabla f) \\[0.5em]
&= \frac{\partial^2 f}{\partial x^2} + \frac{\partial^2 f}{\partial y^2} + \frac{\partial^2 f}{\partial z^2}.
\end{aligned}
\]

Scalar fields satisfying $\Delta f = 0$ throughout their domain are called harmonic functions, and they play a central role in potential theory and in the study of steady-state physical phenomena. The Laplacian operator can furthermore be extended to act on a vector field $\mathbf{F}$, by applying it separately, or component-wise, to each of the three scalar components of the field,
\[
\Delta \mathbf{F} = (\Delta P) \, \mathbf{i} + (\Delta Q) \, \mathbf{j} + (\Delta R) \, \mathbf{k}.
\]

When the theory of vector calculus is extended from ordinary three-dimensional space to the four-dimensional Minkowski spacetime with coordinates $(t,x,y,z)$, a setting frequently employed in the study of wave propagation and in relativity, the Laplacian itself is generalized into a further second-order operator, known as the D'Alembertian, or wave operator, and denoted by $\dalembert$. This operator incorporates, in addition to the spatial Laplacian, a second derivative with respect to time, suitably balanced by the speed of light $c$, or by the value $c=1$ if natural units are adopted; explicitly,
\[
\begin{aligned}
\dalembert f &= \frac{1}{c^2} \frac{\partial^2 f}{\partial t^2} - \Delta f \\
&= \frac{1}{c^2} \frac{\partial^2 f}{\partial t^2}
- \left( \frac{\partial^2 f}{\partial x^2} + \frac{\partial^2 f}{\partial y^2} + \frac{\partial^2 f}{\partial z^2} \right).
\end{aligned}
\]

\begin{remark}
Depending on the metric signature adopted in a given treatment of theoretical physics, either the mostly-minus or the mostly-plus convention, one may instead encounter the opposite sign convention for the D'Alembertian, namely
\[
\dalembert f = \Delta f - \frac{1}{c^2} \frac{\partial^2 f}{\partial t^2}.
\]

Regardless of which sign convention is adopted, the equation $\dalembert f = 0$ represents the homogeneous wave equation, which describes how electromagnetic or acoustic waves propagate through space as time evolves.
\end{remark}

\section{Quadratic Forms, Least Squares, and Multivariable Extrema}

\subsection{Gradients and Hessians of Quadratic Forms}

Quadratic forms occur pervasively in optimization and in statistics, most notably in the least squares problem, and it is therefore essential to compute their gradients and Hessians explicitly.

\begin{definition}[Quadratic Form]
Given a matrix $A \in M_n$, the quadratic form associated with $A$ is the function $Q : \mathbb{R}^n \to \mathbb{R}$ defined by $Q(\mathbf{x}) = \mathbf{x}^{\top} A \mathbf{x}$.

\end{definition}

\begin{theorem}[Gradient and Hessian of a Quadratic Form]
Given $A \in M_n$, the gradient and Hessian of $Q(\mathbf{x}) = \mathbf{x}^{\top}A\mathbf{x}$ are
\[
\begin{aligned}
\nabla Q(\mathbf{x}) &= (A + A^{\top})\mathbf{x}, \\[0.5em]
\nabla^2 Q(\mathbf{x}) &= A + A^{\top}.
\end{aligned}
\]
In particular, if $A$ is symmetric,
\[
\begin{gathered}
\nabla Q(\mathbf{x}) = 2A\mathbf{x} \\[0.5em]
\nabla^2 Q(\mathbf{x}) = 2A.
\end{gathered}
\]

\end{theorem}

\begin{proof}
Writing \[
Q(\mathbf{x}) = \sum_{i,j} a_{ij} x_i x_j
\] and differentiating with respect to $x_k$, the terms with $i=k$, namely \[
\sum_j a_{kj}x_k x_j,
\] contribute \[
\sum_j a_{kj} x_j
\] upon differentiation, while the terms with $j=k$, namely \[
\sum_i a_{ik}x_ix_k,
\] contribute \[
\sum_i a_{ik} x_i;
\] the single term with
\[
\begin{aligned}
i &= j \\
&= k.
\end{aligned}
\]
This is counted correctly by either contribution alone and not twice, since it is a single term $a_{kk}x_k^2$ whose derivative $2a_{kk}x_k$ equals the sum of the two separate contributions $a_{kk}x_k$ from each. One obtains \[
\partial Q / \partial x_k = \sum_j a_{kj}x_j + \sum_i a_{ik}x_i,
\] which is the $k$-th component of $A\mathbf{x} + A^{\top}\mathbf{x}$; since this holds for every $k$,
\[
\nabla Q(\mathbf{x}) = (A+A^{\top})\mathbf{x}.
\]

Differentiating this expression once more with respect to $x_l$, since $A+A^{\top}$ is a constant matrix, yields directly
\[
\partial^2 Q/\partial x_l \partial x_k = (A+A^{\top})_{kl},
\]
Hence the Hessian is
\[
\nabla^2 Q(\mathbf{x}) = A+A^{\top}.
\]

\end{proof}

Together with the linearity of differentiation, this result yields the gradient and Hessian of the general quadratic function
\[
\begin{gathered}
f(\mathbf{x}) = \mathbf{x}^{\top}A\mathbf{x} + \mathbf{b}^{\top}\mathbf{x} + c, \\[0.5em]
\text{namely}\quad \nabla f(\mathbf{x}) = (A+A^{\top})\mathbf{x} + \mathbf{b} \\[0.5em]
\nabla^2 f(\mathbf{x}) = A + A^{\top}.
\end{gathered}
\]
Since $\nabla(\mathbf{b}^{\top}\mathbf{x}) = \mathbf{b}$, directly from the definition of the gradient applied to the linear function \[
\mathbf{b}^{\top}\mathbf{x} = \sum_i b_i x_i,
\] and the constant $c$ contributes nothing to either derivative.

\subsection{Least Squares and the Normal Equations}

\begin{example}[Derivation of the Normal Equations]
Consider the least squares problem of minimizing $f(\mathbf{x}) = \|A\mathbf{x} - \mathbf{b}\|^2$ over $\mathbf{x} \in \mathbb{R}^n$, for $A \in M_{m,n}$ and $\mathbf{b} \in \mathbb{R}^m$. Expanding,
\[
\begin{aligned}
f(\mathbf{x}) &= (A\mathbf{x}-\mathbf{b})^{\top}(A\mathbf{x}-\mathbf{b}) \\[0.5em]
&= \mathbf{x}^{\top}A^{\top}A\mathbf{x} - 2\mathbf{b}^{\top}A\mathbf{x} + \mathbf{b}^{\top}\mathbf{b}.
\end{aligned}
\]
A quadratic function with symmetric matrix $A^{\top}A$ and linear coefficient $-2A^{\top}\mathbf{b}$. Its gradient is $\nabla f(\mathbf{x}) = 2A^{\top}A\mathbf{x} - 2A^{\top}\mathbf{b}$, and its Hessian is the constant matrix $2A^{\top}A$, which is positive semidefinite since
\[
\begin{aligned}
\mathbf{v}^{\top}(A^{\top}A)\mathbf{v} &= \|A\mathbf{v}\|^2 \\
&\geq 0.
\end{aligned}
\]
This holds for every $\mathbf{v}$, so that every critical point of $f$ is automatically a global minimum. Setting $\nabla f(\mathbf{x}) = \vec{0}$ yields the normal equations $A^{\top}A\mathbf{x} = A^{\top}\mathbf{b}$, whose solutions are precisely the least squares solutions of the original, generally inconsistent, system $A\mathbf{x} = \mathbf{b}$; when $A^{\top}A$ is nonsingular, the unique solution is $\mathbf{x} = (A^{\top}A)^{-1}A^{\top}\mathbf{b}$, and this expression coincides with $A^{+}\mathbf{b}$, the pseudoinverse solution obtained in the previous section, whenever $A$ has full column rank.
\end{example}

\subsection{Multivariable Extrema and Second-Derivative Tests}

\begin{definition}
A function $f$ of two variables has a local maximum at $(a,b)$ if $f(x,y) \leq f(a,b)$ for every $(x,y)$ in some neighborhood $U(a,b)$; the number $f(a,b)$ is then called a local maximum value. An analogous definition, with the inequality reversed, defines a local minimum and the corresponding local minimum value.
\end{definition}

As in the single-variable theory, the vanishing of an appropriate derivative provides a necessary condition for the presence of a local extremum, and this condition is furnished by the following two-variable generalization of Fermat's Theorem.

\begin{theorem}[Fermat's Theorem for Functions of Two Variables]
If $f$ has a local maximum or minimum at $(a,b)$, and the first-order partial derivatives of $f$ exist there, then $\nabla f(a,b) = (0,0)$, that is, $f_x(a,b) = 0$ and $f_y(a,b) = 0$.

\end{theorem}

The idea underlying the proof is that, if $f$ has a local maximum at $(a,b)$, then the restriction of $f$ along the line $y=b$, namely the single-variable function $g(x) = f(x,b)$, itself has a local maximum at $x=a$; the single-variable Fermat's Theorem therefore yields $g'(a) = 0$, and since $g'(a) = f_x(a,b)$, one concludes that $f_x(a,b) = 0$.

An entirely analogous argument, applied to the restriction of $f$ along the line $x=a$, yields $f_y(a,b) = 0$.

\begin{definition}
A point $(a,b)$ is called a critical point of $f$ if $\nabla f(a,b) = (0,0)$.
\end{definition}

Fermat's Theorem, restated in this terminology, asserts that if $f$ has a local maximum or minimum at $(a,b)$, and the partial derivatives of $f$ exist at $(a,b)$, then $(a,b)$ must be a critical point of $f$. It is important to observe, however, that, exactly as in the single-variable theory, not every critical point corresponds to a local maximum or minimum; a critical point may instead be a saddle point, at which the function increases along some directions and decreases along others.

In order to determine whether a given critical point actually corresponds to a local extremum, or instead to a saddle point, one relies on the following test, formulated in terms of the second-order partial derivatives of $f$.

\begin{theorem}[Second Derivatives Test]
Suppose the second partial derivatives of $f$ are continuous on a neighborhood $U(a,b)$, and suppose that $(a,b)$ is a critical point of $f$. Let
\[
D = f_{xx}(a,b) \, f_{yy}(a,b) - \big[ f_{xy}(a,b) \big]^2.
\]

If $D < 0$, then $(a,b)$ is neither a local maximum nor a local minimum, but a saddle point. If $D > 0$ and $f_{xx}(a,b) > 0$, then $f(a,b)$ is a local minimum. If $D > 0$ and $f_{xx}(a,b) < 0$, then $f(a,b)$ is a local maximum.
\end{theorem}

\begin{remark}
If $D = 0$, the test is inconclusive: the function $f$ could have a local maximum, a local minimum, or a saddle point at $(a,b)$, and no further information can be extracted from the sign of $D$ alone. This ambiguity is illustrated by the three functions
\[
\begin{gathered}
z = x^4+y^4, \\[0.5em]
z = -x^4-y^4, \\[0.5em]
z = x^4-y^4.
\end{gathered}
\]
All of these expressions have a critical point with $D=0$ at the origin, yet exhibit, respectively, a local minimum, a local maximum, and a saddle point there.
\end{remark}

The four second-order partial derivatives of $f$ at $(a,b)$ can be conveniently organized into a matrix, called the Hessian matrix of $f$ at $(a,b)$,
\[
H_f(a,b) = \begin{pmatrix} f_{xx}(a,b) & f_{xy}(a,b) \\ f_{yx}(a,b) & f_{yy}(a,b) \end{pmatrix}.
\]
The quantity $D$ appearing in the Second Derivatives Test is precisely the determinant of this Hessian matrix. Consider
\[
f(x,y) = x^4+y^4-4xy+1.
\]
Its gradient vanishes when
\[
\begin{aligned}
f_x&=4x^3-4y=0,\\
f_y&=4y^3-4x=0.
\end{aligned}
\]
Thus $y=x^3$ and $x=y^3$, so $x=x^9$. Hence
$x\in\{-1,0,1\}$ and the critical points are
\[
(-1,-1),\qquad(0,0),\qquad(1,1).
\]
The Hessian and its determinant are
\[
H_f(x,y)=\begin{pmatrix}12x^2&-4\\-4&12y^2\end{pmatrix},
\qquad D(x,y)=144x^2y^2-16.
\]
At $(0,0)$, $D=-16<0$, so the origin is a saddle point. At
$(1,1)$ and $(-1,-1)$, $D=128>0$ and $f_{xx}=12>0$, so both are
local minima. Their common value is
\[
f(1,1)=f(-1,-1)=1+1-4+1=-1.
\]
There are no local maxima among the critical points.

Figure~\ref{fig:atlas-saddle-surface} makes the opposite curvatures at a saddle point visible.

\begin{figure}[!htbp]
\centering
\begin{tikzpicture}[>=Stealth]
\begin{axis}[width=11cm,height=7cm,axis lines=middle,ticklabel style={font=\small},label style={font=\small},legend style={font=\scriptsize,draw=black!20,fill=white},view={45}{28},xlabel={$x$},ylabel={$y$},zlabel={$z$},colormap={soft}{color=(white);color=(figblue!45)},xtick={-1,0,1},ytick={-1,0,1},ztick={-1,0,1}]
\addplot3[surf,shader=interp,opacity=0.65,domain=-1.2:1.2,y domain=-1.2:1.2,samples=19,samples y=19] {x^2-y^2};
\addplot3[figblue,very thick,domain=-1.2:1.2,samples=60,samples y=1] ({x},{0},{x^2});
\addplot3[figred,very thick,domain=-1.2:1.2,samples=60,samples y=1] ({0},{x},{-x^2});
\end{axis}
\end{tikzpicture}
\caption{The surface $z=x^2-y^2$ rises along the $x$-axis and falls along the $y$-axis. The origin is stationary but is neither a local minimum nor a local maximum. The blue and red section curves display the two opposing behaviors.}
\label{fig:atlas-saddle-surface}
\end{figure}

\section{Differentiability, Directional Derivatives, and Jacobians}

\subsection{Tangent Planes and Differentiability}

Before introducing the notion of differentiability for a function of two variables, it is useful to recall that any plane in $\mathbb{R}^3$ passing through a point $P(x_0,y_0,z_0)$ admits an equation of the form $A(x-x_0) + B(y-y_0) + C(z-z_0) = 0$, or equivalently $Ax+By+Cz=D$; the particular planes $z=z_0$, $y=y_0$, and $x=x_0$ arise as special cases of this general form.

Let $S$ denote the graph of a function $z = f(x,y)$, regarded as a surface in $\mathbb{R}^3$, and let $C_1$ and $C_2$ denote the traces of this surface in the planes $y=b$ and $x=a$ respectively, that is, the curves obtained by intersecting $S$ with these two coordinate planes. The partial derivatives $f_x(a,b)$ and $f_y(a,b)$ admit a direct geometric interpretation as the slopes of the tangent lines $T_1$ and $T_2$ to the traces $C_1$ and $C_2$ at the common point $P(a,b,c)$, where $c = f(a,b)$; this geometric interpretation is well illustrated by the graph of $f(x,y) = 4-x^2-2y^2$, together with the corresponding tangent lines at the point $(1,1)$.

\begin{remark}
A function may possess both partial derivatives at a point $(a,b)$ without being continuous at that same point. This somewhat surprising phenomenon is illustrated by the function
\[
f(x,y) =
\begin{cases}
\dfrac{xy}{x^2+y^2} & (x,y) \neq (0,0), \\[1mm]
0 & (x,y) = (0,0)
\end{cases}.
\]
Both partial derivatives of this function exist at the origin, and in fact $f_x(0,0) = f_y(0,0) = 0$; however, $f(x,y) = 1/2$ at every point on the line $y=x$ other than the origin, while $f(x,y) = 0$ at every point on the line $y=0$, so that $f$ fails to be continuous at the origin, since it does not approach a common limiting value along these two distinct lines.
\end{remark}

The plane containing the two tangent lines $T_1$ and $T_2$ discussed above is called the tangent plane to the graph of $f$ at the point $(a,b,f(a,b))$, and its equation is given explicitly by
\[
z = f(a,b) + f_x(a,b)(x-a) + f_y(a,b)(y-b).
\]

Recall that, for a function of a single variable, differentiability of $f$ at $x=a$ means precisely that the tangent line $L(x) = f(a) + f'(a)(x-a)$ provides the best possible linear approximation of $f$ near $x=a$, in the sense that
\[
\lim_{x \to a} [f(x)-L(x)]/(x-a) = 0.
\]

This notion extends naturally to functions of two variables, with the tangent line replaced by the tangent plane introduced above. Denote by $L(x,y)$ the linear function whose graph is precisely this tangent plane, that is,
\[
L(x,y) = f(a,b) + f_x(a,b)(x-a) + f_y(a,b)(y-b).
\]

\begin{definition}
We say that $f$ is differentiable at $(a,b)$ if the tangent plane provides a good approximation of the function near $(a,b)$, in the sense that
\[
\lim_{(x,y) \to (a,b)} \frac{f(x,y) - L(x,y)}{\sqrt{(x-a)^2+(y-b)^2}} = 0.
\]

\end{definition}

As already observed, a function may possess partial derivatives at a point without being continuous there; however, this pathology cannot occur once the function is known to be differentiable at that point, as the following theorem shows.

\begin{theorem}
If $f$ is differentiable at $(a,b)$, then $f$ is continuous at $(a,b)$.

\end{theorem}

While the definition of differentiability given above is precise, it is often cumbersome to verify directly from the definition. The following theorem provides a convenient sufficient condition for differentiability, expressed purely in terms of the continuity of the partial derivatives, and is the criterion most frequently applied in practice.

\begin{theorem}[Sufficient Condition for Differentiability]
If the partial derivatives $f_x$ and $f_y$ exist on a neighborhood $U(a,b)$ and are continuous at $(a,b)$, then $f$ is differentiable at $(a,b)$.

\end{theorem}

\subsection{Directional Derivatives and the Gradient}

The partial derivatives $f_x(x_0,y_0)$ and $f_y(x_0,y_0)$ represent the rates of change of $z=f(x,y)$ at the point $(x_0,y_0)$ specifically in the $x$- and $y$-directions, that is, in the directions of the unit vectors $\mathbf{i}$ and $\mathbf{j}$ respectively. It is natural to seek a corresponding notion of rate of change in an arbitrary direction, specified by an arbitrary unit vector $\mathbf{u} = a\mathbf{i} + b\mathbf{j}$, satisfying $a^2+b^2=1$.

\begin{definition}[Directional Derivative]
The directional derivative of $f$ at $(x_0,y_0)$ in the direction of the unit vector $\mathbf{u} = a\mathbf{i}+b\mathbf{j}$, denoted $D_{\mathbf{u}} f(x_0,y_0)$, is defined as \[
D_{\mathbf{u}} f(x_0,y_0) = \lim_{t \to 0} \frac{f(x_0+ta, y_0+tb) - f(x_0,y_0)}{t},
\] provided this limit exists and is finite. Geometrically, this quantity represents the slope of the tangent line $T$ at the point $P(x_0,y_0,z_0)$ to the trace $C$ obtained by intersecting the graph of $f$ with the vertical plane through $(x_0,y_0)$ in the direction $\mathbf{u}$.
\end{definition}

When $f$ is known to be differentiable at $(x_0,y_0)$, its directional derivatives in every direction can be computed very efficiently by means of the following formula, which reduces the computation of a directional derivative to a simple dot product involving the gradient vector.

\begin{theorem}[Gradient Formula]
Suppose $f$ is differentiable at $(x_0,y_0)$. Then $f$ has a directional derivative $D_{\mathbf{u}} f(x_0,y_0)$ in the direction of every unit vector $\mathbf{u} = a\mathbf{i}+b\mathbf{j}$, given by
\[
\begin{aligned}
D_{\mathbf{u}} f(x_0,y_0) &= \nabla f(x_0,y_0) \cdot \mathbf{u} \\[0.5em]
&= f_x(x_0,y_0) \, a + f_y(x_0,y_0) \, b.
\end{aligned}
\]

\end{theorem}

The proof of this formula relies on the Chain Rule for partial derivatives of composite functions, a tool that extends the single-variable Chain Rule to functions of several variables and that will be developed in subsequent material.

The Gradient Formula established above has an important consequence concerning the direction in which a function increases most rapidly at a given point.

\begin{theorem}
The maximum value of the directional derivative $D_{\mathbf{u}} f(x_0,y_0)$, taken over all unit vectors $\mathbf{u}$, equals $\| \nabla f(x_0,y_0) \|$, and this maximum value is attained precisely when $\mathbf{u}$ has the same direction as the gradient vector $\nabla f(x_0,y_0)$, that is, when
\[
\mathbf{u} = \frac{\nabla f(x_0,y_0)}{\| \nabla f(x_0,y_0) \|}.
\]

\end{theorem}

For the function $f(x,y) = x^2 y + \sqrt{y}$ at $(2,1)$, the gradient is
\[
\nabla f(x,y)=\left(2xy,\,x^2+\frac{1}{2\sqrt y}\right),
\]
and hence
\[
\nabla f(2,1)=\left(4,\frac92\right).
\]
The maximum rate of change is the gradient norm,
\[
\left\|\nabla f(2,1)\right\|
=\sqrt{4^2+\left(\frac92\right)^2}
=\frac{\sqrt{145}}{2},
\]
and it occurs in the unit direction
\[
\mathbf{u}_{\max}
=\frac{\nabla f(2,1)}{\|\nabla f(2,1)\|}
=\left(\frac8{\sqrt{145}},\frac9{\sqrt{145}}\right).
\]

A further important geometric property of the gradient vector, complementing the result just stated, concerns its relationship to the level curves of $f$: the gradient vector $\nabla f(x_0,y_0)$ is always perpendicular to the level curve of $f$ that passes through the point $(x_0,y_0)$, providing a further geometric characterization of the direction of steepest ascent described above.

Figure~\ref{fig:ch9-gradient} illustrates the relationship between level curves, the gradient, and directional derivatives. For the scalar field $f(x,y)=x^2+y^2$ the level curves are concentric circles. The gradient $\nabla f$ is perpendicular to the
level curve at the point and maximizes the directional derivative $D_{\bu} f = \nabla f \cdot \bu$ when $\bu \parallel \nabla f$.

\begin{figure}[!htbp]
\centering
\begin{tikzpicture}[scale=1.15,>=Stealth]
\draw[->] (-3.3,0) -- (3.3,0) node[right,font=\scriptsize]{$x$};
\draw[->] (0,-3.3) -- (0,3.3) node[above,font=\scriptsize]{$y$};
\foreach \r/\c in {0.71/figblue!25,1.22/figblue!40,1.73/figblue!55,2.24/figblue!70,2.74/figblue!85,3.24/figblue}{
\draw[\c,thick,domain=0:360,samples=100,smooth,variable=\t]
plot ({\r*cos(\t)},{\r*sin(\t)});
}
\node[circle,fill=black,inner sep=1.2pt] at (1.2,0.9) {};
\draw[vecB] (1.2,0.9) -- (2.16,1.62) node[above right,font=\scriptsize]{$\nabla f$};
\draw[vecC] (1.2,0.9) -- (2.0,0.05) node[below right,yshift=-2pt,font=\scriptsize]{$\bu$};
\draw[figblue,dashed] (0.05,1.5) -- (2.35,0.3);
\node[figblue,above left,fill=white,inner sep=1.5pt,font=\scriptsize]
  at (0.45,1.55) {tangent};
\end{tikzpicture}
\caption[Level curves, gradient, and directional derivative]{Concentric level curves of the scalar field, with the gradient shown in red and a comparison direction in green. The gradient is normal to the level curve at the marked point and gives the direction of greatest increase.}
\label{fig:ch9-gradient}
\end{figure}

\subsection{Vector-Valued Maps and the Jacobian}

Differential calculus, as developed for functions of a single variable and subsequently for scalar fields of several variables, associates to a differentiable function a linear approximation of its local behavior. The same idea extends to functions taking values in $\mathbb{R}^m$ rather than in $\mathbb{R}$, and the matrix encoding this linear approximation is called the Jacobian matrix.

\begin{definition}[Jacobian Matrix]
Given a function $\mathbf{f} : \Omega \subseteq \mathbb{R}^n \to \mathbb{R}^m$, with components $\mathbf{f} = (f_1, \dots, f_m)$, each admitting all first-order partial derivatives at $\mathbf{x} \in \Omega$, the Jacobian matrix of $\mathbf{f}$ at $\mathbf{x}$ is the matrix $J_{\mathbf{f}}(\mathbf{x}) \in M_{m,n}$ whose $(i,j)$ entry is the partial derivative $\partial f_i / \partial x_j$, evaluated at $\mathbf{x}$.
\end{definition}

As with scalar fields, the mere existence of the partial derivatives comprising the Jacobian matrix does not by itself guarantee that this matrix furnishes a genuine local linear approximation of $\mathbf{f}$; the appropriate strengthening of this notion is differentiability.

\begin{definition}[Differentiability of a Vector-Valued Function]
The function $\mathbf{f} : \Omega \to \mathbb{R}^m$ is differentiable at $\mathbf{x} \in \Omega$ if there exists a matrix $J \in M_{m,n}$ such that
\[
\lim_{\mathbf{h} \to \vec{0}} \frac{\|\mathbf{f}(\mathbf{x}+\mathbf{h}) - \mathbf{f}(\mathbf{x}) - J\mathbf{h}\|}{\|\mathbf{h}\|} = 0.
\]

\end{definition}

\begin{theorem}
If $\mathbf{f}$ is differentiable at $\mathbf{x}$, then the matrix $J$ in the definition above is unique and equals $J_{\mathbf{f}}(\mathbf{x})$.

\end{theorem}

\begin{proof}
Applying the defining limit with $\mathbf{h} = t\mathbf{e}_j$ for $t \to 0$ shows that the $j$-th column of $J$ equals the directional derivative of $\mathbf{f}$ in the direction $\mathbf{e}_j$, which is precisely $\partial \mathbf{f}/\partial x_j (\mathbf{x})$, the $j$-th column of $J_{\mathbf{f}}(\mathbf{x})$; since this holds for every $j = 1, \dots, n$, the two matrices coincide, and uniqueness follows since a matrix is determined by its columns.
\end{proof}

\begin{example}
Consider the transformation from polar to Cartesian coordinates,
\[
\mathbf{f}(r,\theta) = (r\cos\theta, r\sin\theta).
\]
Its Jacobian matrix is
\[
J_{\mathbf{f}}(r,\theta) = \begin{pmatrix} \cos\theta & -r\sin\theta \\ \sin\theta & r\cos\theta \end{pmatrix}.
\]
Its determinant is
\[
\begin{aligned}
\operatorname{Det}(J_{\mathbf{f}}) &= r\cos^2\theta + r\sin^2\theta \\[0.5em]
&= r.
\end{aligned}
\]
This recovers the familiar Jacobian factor of polar coordinates used in the change of variables formula for double integrals. Since this determinant is nonzero for $r \neq 0$, the Inverse Function Theorem guarantees that $\mathbf{f}$ is locally invertible near every point with $r \neq 0$. On a domain with a chosen angular branch, for example $\theta\in(-\pi,\pi)$ with the nonpositive $x$-axis removed from the Cartesian image, an inverse is
\[
\begin{gathered}
r = (x^2+y^2)^{1/2}, \\[0.5em]
\theta = \operatorname{atan2}(y,x)
\end{gathered}
\]
The function $\operatorname{atan2}$ is discontinuous across its branch
cut, so this formula does not define a single globally smooth inverse on
$\mathbb{R}^2\setminus\{0\}$.
\end{example}

The chain rule established for functions of one and several variables generalizes to compositions of vector-valued functions, taking the particularly clean form of a matrix product.

\begin{theorem}[Chain Rule, Matrix Form]
Let $\mathbf{f} : \mathbb{R}^n \to \mathbb{R}^m$ be differentiable at $\mathbf{x}$ and $\mathbf{g} : \mathbb{R}^m \to \mathbb{R}^p$ be differentiable at $\mathbf{f}(\mathbf{x})$. Then $\mathbf{g} \circ \mathbf{f}$ is differentiable at $\mathbf{x}$, and
\[
J_{\mathbf{g} \circ \mathbf{f}}(\mathbf{x}) = J_{\mathbf{g}}(\mathbf{f}(\mathbf{x})) \, J_{\mathbf{f}}(\mathbf{x}).
\]

\end{theorem}

\begin{proof}
Writing $\mathbf{y} = \mathbf{f}(\mathbf{x})$, the componentwise chain rule for scalar functions of several variables, applied to each component $g_i \circ \mathbf{f}$ in turn, gives, for each $i = 1, \dots, p$ and $j = 1, \dots, n$,
\[
\frac{\partial (g_i \circ \mathbf{f})}{\partial x_j}(\mathbf{x}) = \sum_{k=1}^{m} \frac{\partial g_i}{\partial y_k}(\mathbf{y}) \frac{\partial f_k}{\partial x_j}(\mathbf{x}).
\]
This is precisely the $(i,j)$ entry of the matrix product
\[
J_{\mathbf{g}}(\mathbf{f}(\mathbf{x})) J_{\mathbf{f}}(\mathbf{x});
\]
since this identity holds entrywise, the two matrices coincide.
\end{proof}

When $n=m$, the Jacobian matrix is square, and its determinant, called the Jacobian determinant, measures the local factor by which $\mathbf{f}$ distorts $n$-dimensional volume near a given point; its sign, moreover, indicates whether $\mathbf{f}$ preserves or reverses orientation locally. As already illustrated, the nonvanishing of this determinant at a point is, by the Inverse Function Theorem, sufficient to guarantee that $\mathbf{f}$ admits a differentiable local inverse near that point, whose Jacobian is then obtained by inverting $J_{\mathbf{f}}$: differentiating the identity
\[
\mathbf{f}^{-1}(\mathbf{f}(\mathbf{x})) = \mathbf{x}
\]
via the Chain Rule gives
\[
\begin{gathered}
J_{\mathbf{f}^{-1}}(\mathbf{f}(\mathbf{x})) J_{\mathbf{f}}(\mathbf{x}) = I, \\[0.5em]
\text{so that}\quad J_{\mathbf{f}^{-1}}(\mathbf{f}(\mathbf{x})) = J_{\mathbf{f}}(\mathbf{x})^{-1}.
\end{gathered}
\]

When $m=1$, the function $\mathbf{f}$ reduces to a scalar field $f : \Omega \to \mathbb{R}$, and its Jacobian matrix $J_f(\mathbf{x}) \in M_{1,n}$ is the single row vector $( \partial f/\partial x_1, \dots, \partial f/\partial x_n )$, that is, precisely the transpose of the gradient $\nabla f(\mathbf{x})$ introduced earlier for scalar fields. The Chain Rule in matrix form then recovers, as a special case, the familiar formula for the directional derivative: if $\boldsymbol{\gamma}(t)$ is a differentiable curve with $\boldsymbol{\gamma}(0) = \mathbf{x}$ and $\boldsymbol{\gamma}'(0) = \mathbf{u}$, then
\[
\begin{aligned}
\frac{d}{dt}\Big|_{t=0} f(\boldsymbol{\gamma}(t)) &= J_f(\mathbf{x}) \boldsymbol{\gamma}'(0) \\[0.5em]
&= \nabla f(\mathbf{x}) \cdot \mathbf{u}.
\end{aligned}
\]
This recovers the directional derivative formula $D_{\mathbf{u}}f(\mathbf{x}) = \nabla f(\mathbf{x}) \cdot \mathbf{u}$ established previously for unit vectors $\mathbf{u}$. Figure~\ref{fig:ch10-jacobian-polar} relates the polar coordinate directions to the columns of the Jacobian. The polar grid uses coordinates $(r,\theta)$ on the Cartesian plane; the highlighted
curvilinear cell has sides $dr$ and $r\,d\theta$ along the directions $\partial/\partial r$, $\partial/\partial\theta$: the columns of the Jacobian $J_f$, with $\det J_f = r$.

\begin{figure}[!htbp]
\centering
\begin{tikzpicture}[scale=1.3,>=Stealth]
\draw[->] (-2.3,0) -- (2.3,0) node[right,font=\scriptsize]{$x$};
\draw[->] (0,-2.3) -- (0,2.3) node[above,font=\scriptsize]{$y$};
\foreach \r in {0.5,1,1.5,2} \draw[gridline] (0,0) circle (\r);
\foreach \t in {0,30,...,330} \draw[gridline] (0,0) -- (\t:2.1);
\fill[figblue,opacity=0.35] (55:1) -- (55:1.9) arc (55:95:1.9) -- (95:1) arc (95:55:1);
\draw[vecB] (55:1) -- (55:1.9) node[midway,right=3pt,font=\scriptsize]{$\partial/\partial r$};
\draw[vecC] (55:1) arc[start angle=55,end angle=82,radius=1]
  node[pos=0.65,above left=2pt,font=\scriptsize]{$\partial/\partial\theta$};
\node[below left,font=\scriptsize] at (55:0.92) {$(r,\theta)$};
\end{tikzpicture}
\caption[Polar grid and Jacobian]{Polar coordinate grid on Cartesian axes, with a curvilinear cell highlighted in blue. The red and green arrows indicate radial and angular coordinate directions, respectively; the Jacobian accounts for the corresponding local area scaling.}
\label{fig:ch10-jacobian-polar}
\end{figure}

\section{Multilinear Maps and Tensors}

Vectors, linear functionals, and matrices, though introduced separately over the course of this document, are in fact instances of a single unifying algebraic structure: the tensor. Tensor calculus organizes and generalizes these objects, and the operations acting upon them, within one coherent framework, one that becomes indispensable once coordinates are allowed to change in ways more general than the orthonormal changes of basis considered earlier.

\begin{definition}[Tensor]
Given a finite-dimensional vector space $V$ with dual space $V^{*}$, a tensor of type $(r,s)$ on $V$ is a multilinear map
\[
T : \underbrace{V^{*} \times \cdots \times V^{*}}_{r} \times \underbrace{V \times \cdots \times V}_{s} \longrightarrow \mathbb{R}.
\]
That is, a map that is linear separately in each of its $r+s$ arguments.
\end{definition}

This single definition subsumes the objects encountered throughout this document as special cases distinguished only by the pair $(r,s)$. A scalar is a tensor of type $(0,0)$. A vector $\mathbf{v} \in V$ corresponds to a tensor of type $(1,0)$, acting on a covector $\boldsymbol{\varphi} \in V^{*}$ by $\mathbf{v}(\boldsymbol{\varphi}) = \boldsymbol{\varphi}(\mathbf{v})$; this identification of $V$ with the dual of $V^{*}$ is legitimate in finite dimension, where $V$ and its double dual $V^{**}$ are canonically isomorphic. A linear functional, or covector, is itself a tensor of type $(0,1)$.

A linear map $L : V \to V$, represented in earlier sections by a matrix $A$, corresponds to a tensor of type $(1,1)$, acting on a pair $(\boldsymbol{\varphi}, \mathbf{v})$ by $\boldsymbol{\varphi}(L\mathbf{v})$; one may verify directly that this assignment is linear in $\boldsymbol{\varphi}$, since $\boldsymbol{\varphi}$ itself is linear, and linear in $\mathbf{v}$, since $L$ is linear.

\subsection{Components and the Einstein Summation Convention}

Once a basis $\mathbf{e}_1, \dots, \mathbf{e}_n$ of $V$, and the corresponding dual basis $\boldsymbol{\varepsilon}^1, \dots, \boldsymbol{\varepsilon}^n$ of $V^{*}$ satisfying $\boldsymbol{\varepsilon}^i(\mathbf{e}_j) = \delta^i_j$, are fixed, every tensor is fully described by its components with respect to these bases, and it becomes conventional to place indices referring to the dual space as superscripts, called contravariant indices, and indices referring to $V$ itself as subscripts, called covariant indices.

\begin{definition}[Einstein Summation Convention]
Whenever an index appears once as a superscript and once as a subscript within a single term, summation over that index, from $1$ to $n$, is implied without an explicit summation symbol.
\end{definition}

Under this convention, a vector $\mathbf{v} = v^i \mathbf{e}_i$ is described by its contravariant components $v^i$, while a covector $\boldsymbol{\varphi} = \varphi_i \boldsymbol{\varepsilon}^i$ is described by its covariant components $\varphi_i$; the pairing $\boldsymbol{\varphi}(\mathbf{v}) = \varphi_i v^i$ is a contraction, discussed further below.

\begin{theorem}[Transformation Law for Components]
Let $\hat{\mathbf{e}}_j = P_j^{\;i}\mathbf{e}_i$ define a new basis of $V$ via an invertible change of basis matrix $P$, with inverse having entries $(P^{-1})^i_{\;j}$. Then the contravariant components of a fixed vector $\mathbf{v}$ transform as $\hat{v}^i = (P^{-1})^i_{\;j}v^j$, while the covariant components of a fixed covector $\boldsymbol{\varphi}$ transform as $\hat\varphi_i = P_i^{\;j}\varphi_j$.
\end{theorem}

\begin{proof}
The vector $\mathbf{v}$ itself does not depend on the choice of basis, so
\[
\begin{aligned}
\mathbf{v} &= v^i\mathbf{e}_i \\[0.5em]
&= \hat{v}^j\hat{\mathbf{e}}_j \\[0.5em]
&= \hat{v}^j P_j^{\;i}\mathbf{e}_i.
\end{aligned}
\]

Comparing coefficients of $\mathbf{e}_i$ on both sides, since $\mathbf{e}_1,\dots,\mathbf{e}_n$ are linearly independent, gives $v^i = P_j^{\;i}\hat{v}^j$; multiplying both sides by $(P^{-1})^i_{\;k}$ and summing over $i$ recovers $\hat{v}^k = (P^{-1})^k_{\;i}v^i$, as claimed, after relabeling the free index. For the covector, the dual basis vectors transform by the inverse-transpose relation $\hat{\boldsymbol{\varepsilon}}^i = (P^{-1})^i_{\;j}\boldsymbol{\varepsilon}^j$, since this is precisely the unique transformation making
\[
\begin{aligned}
\hat{\boldsymbol{\varepsilon}}^i(\hat{\mathbf{e}}_k) &= (P^{-1})^i_{\;j}P_k^{\;l}\boldsymbol{\varepsilon}^j(\mathbf{e}_l) \\[0.5em]
&= (P^{-1})^i_{\;j}P_k^{\;j} \\[0.5em]
&= \delta^i_k.
\end{aligned}
\]
These identities hold, as required of a dual basis. Since
\[
\begin{aligned}
\boldsymbol{\varphi} &= \varphi_i\boldsymbol{\varepsilon}^i \\[0.5em]
&= \hat\varphi_j\hat{\boldsymbol{\varepsilon}}^j \\[0.5em]
&= \hat\varphi_j (P^{-1})^j_{\;i}\boldsymbol{\varepsilon}^i.
\end{aligned}
\]
Comparing coefficients gives $\varphi_i = (P^{-1})^j_{\;i}\hat\varphi_j$; multiplying by $P_i^{\;k}$ and summing over $i$, and using
\[
P_i^{\;k}(P^{-1})^j_{\;i} = \delta^{jk}.
\]
More directly, inverting the relation $\varphi_i = (P^{-1})^j_{\;i}\hat\varphi_j$ by the matrix $P$ gives $\hat\varphi_k = P_k^{\;i}\varphi_i$, as claimed.
\end{proof}

The terminology ``covariant'' and ``contravariant'' is justified precisely by this theorem: covariant components transform directly with $P$, in the same way as the basis vectors themselves, while contravariant components transform with the inverse of $P$, compensating the change of basis so that the vector $\mathbf{v}$ itself remains unchanged; this is consistent with, and generalizes, the earlier formula $\hat{\mathbf{x}} = M^{-1}\mathbf{x}$ obtained for coordinates with respect to a new basis $\mathcal{B}$.

\subsection{The Metric Tensor}

Among tensors of type $(0,2)$, one plays a distinguished geometric role: it encodes the inner product of the underlying space, and thereby provides a canonical means of converting contravariant components into covariant ones, and conversely.

\begin{definition}[Metric Tensor]
Given a basis $\mathbf{e}_1, \dots, \mathbf{e}_n$ of an inner product space $V$, the metric tensor is the tensor of type $(0,2)$ with components $g_{ij} = \mathbf{e}_i \cdot \mathbf{e}_j$.

\end{definition}

\begin{theorem}
The matrix $[g_{ij}]$ is symmetric and invertible.
\end{theorem}

\begin{proof}
Symmetry is immediate, since
\[
\begin{aligned}
g_{ij} &= \mathbf{e}_i\cdot\mathbf{e}_j \\[0.5em]
&= \mathbf{e}_j\cdot\mathbf{e}_i \\[0.5em]
&= g_{ji}.
\end{aligned}
\]
By symmetry of the inner product. For invertibility, suppose $c^j g_{ij} = 0$ for every $i$, for some scalars $c^1,\dots,c^n$; then, setting $\mathbf{w} = c^j\mathbf{e}_j$, one has
\[
\begin{aligned}
\mathbf{w}\cdot\mathbf{e}_i &= c^jg_{ij} \\[0.5em]
&= 0.
\end{aligned}
\]
This holds for every $i$, so $\mathbf{w}$ is orthogonal to every basis vector, and hence $\mathbf{w}\cdot\mathbf{w} = 0$, forcing $\mathbf{w} = \vec{0}$ by positive-definiteness of the inner product; since $\mathbf{e}_1,\dots,\mathbf{e}_n$ are linearly independent, this forces every $c^j = 0$, showing that $[g_{ij}]$ has trivial kernel, and is therefore invertible.
\end{proof}

Its inverse, denoted $g^{ij}$, is itself the component matrix of a tensor of type $(2,0)$. When the basis is orthonormal, $g_{ij} = \delta_{ij}$, and the metric tensor reduces to the identity matrix, erasing the distinction between covariant and contravariant components that is otherwise essential in a general, non-orthonormal basis.

\begin{definition}[Raising and Lowering Indices]
Given a vector with contravariant components $v^i$, its associated covariant components are $v_i = g_{ij} v^j$; conversely, given covariant components $\varphi_i$, the associated contravariant components are $\varphi^i = g^{ij}\varphi_j$.

\end{definition}

\begin{theorem}
The two operations of raising and lowering indices are mutually inverse: $g^{ij}(g_{jk}v^k) = v^i$ for every choice of components $v^k$.

\end{theorem}

\begin{proof}
Since $g^{ij}$ is, by definition, the matrix inverse of $g_{ij}$, one has $g^{ij}g_{jk} = \delta^i_k$; hence
\[
\begin{aligned}
g^{ij}(g_{jk}v^k) &= (g^{ij}g_{jk})v^k \\[0.5em]
&= \delta^i_k v^k \\[0.5em]
&= v^i.
\end{aligned}
\]
Using the Einstein summation convention to collapse the sum over the repeated index $k$.
\end{proof}

With this notation, the inner product of two vectors $\mathbf{u} = u^i\mathbf{e}_i$ and $\mathbf{v}=v^j\mathbf{e}_j$ is expressed compactly as
\[
\begin{aligned}
\mathbf{u}\cdot\mathbf{v} &= g_{ij}u^iv^j \\[0.5em]
&= u_jv^j.
\end{aligned}
\]
This recovers the ordinary dot product $u^iv^i$ once the basis is orthonormal and $g_{ij}=\delta_{ij}$.

\subsection{Tensor Product and Contraction}

Two operations generate new tensors from existing ones: the tensor product, which combines tensors of lower type into one of higher type without any loss of information, and contraction, which reduces the type of a single tensor by pairing a contravariant index with a covariant one and summing.

\begin{definition}[Tensor Product]
Given a tensor $S$ of type $(r_1,s_1)$ and a tensor $T$ of type $(r_2,s_2)$ on $V$, their tensor product $S \otimes T$ is the tensor of type $(r_1+r_2, s_1+s_2)$ defined by evaluating $S$ on the first $r_1$ dual arguments and the first $s_1$ vector arguments, and $T$ on the remaining arguments, and multiplying the two results.
\end{definition}

In components, if $S$ and $T$ have respective components
\[
\begin{gathered}
S^{i_1\cdots i_{r_1}}_{\;\;\;\;\;\;\;j_1\cdots j_{s_1}} \\
T^{k_1\cdots k_{r_2}}_{\;\;\;\;\;\;\;l_1\cdots l_{s_2}}.
\end{gathered}
\]
Then $S \otimes T$ has components equal to the ordinary product
\[
S^{i_1\cdots i_{r_1}}_{\;\;\;\;\;\;\;j_1\cdots j_{s_1}}\, T^{k_1\cdots k_{r_2}}_{\;\;\;\;\;\;\;l_1\cdots l_{s_2}},
\]
a fact which follows directly from multilinearity of the tensor product upon expanding each argument in the chosen basis.

\begin{definition}[Contraction]
Given a tensor $T$ of type $(r,s)$ with $r,s \geq 1$, the contraction of $T$ over a chosen contravariant index and a chosen covariant index is the tensor of type $(r-1,s-1)$ obtained, in components, by setting the two chosen indices equal and summing, following the Einstein convention.
\end{definition}

\begin{theorem}
The contraction of a tensor is independent of the basis used to compute it, that is, it defines a genuine tensor of type $(r-1,s-1)$.

\end{theorem}

\begin{proof}
It suffices to verify the claim for a tensor $T$ of type $(1,1)$, with components $T^i_{\;j}$, since the general case follows by treating the remaining indices as passive labels throughout the argument. Under a change of basis with matrix $P$, the transformation law established above gives $\hat{T}^i_{\;j} = (P^{-1})^i_{\;k}P_j^{\;l}T^k_{\;l}$, since $T^i_{\;j}$ transforms contravariantly in its upper index and covariantly in its lower index. The contraction is
\[
\begin{aligned}
\hat{T}^i_{\;i} &= (P^{-1})^i_{\;k}P_i^{\;l}T^k_{\;l} \\[0.5em]
&= \delta^l_k T^k_{\;l} \\[0.5em]
&= T^k_{\;k}.
\end{aligned}
\]
Using $(P^{-1})^i_{\;k}P_i^{\;l} = \delta^l_k$, the statement that $P^{-1}$ and $P$ are mutually inverse matrices; hence the contracted quantity $T^k_{\;k}$ is the same scalar regardless of the basis used to compute it, confirming that it is a well-defined, basis-independent scalar, that is, a tensor of type $(0,0)$.

\end{proof}

The contraction of a tensor of type $(1,1)$, represented in a basis by the matrix $A = [a^i_{\;j}]$, over its single pair of indices yields the scalar $a^i_{\;i}$, which is precisely the trace of $A$, now seen to be basis-independent as a direct consequence of the theorem above, recovering from a different direction the earlier observation that similar matrices, representing the same linear map in different bases, share the same trace. More generally, the divergence, gradient, and curl introduced earlier for vector fields in $\mathbb{R}^3$ reappear within this framework as particular contractions of tensors built from the components of a field and their partial derivatives, a perspective that becomes indispensable once such operators must be expressed in curvilinear, rather than Cartesian, coordinate systems, where the metric tensor $g_{ij}$ varies from point to point and must be accounted for explicitly in any differentiation of tensor components.

\newpage
\chapter{Integral Calculus in Multidimensional Spaces}

\section{Double Integrals and Planar Regions}
\label{sec:double-integrals}

The Riemann integral, as developed for functions of a single real variable, measures the signed area between the graph of a function and the interval on which it is defined, via a limiting process applied to sums that approximate this area by finitely many rectangles. The same idea extends, with only notational complications, to functions of several variables: a function of two variables may be integrated over a planar region to yield a signed volume, a function of three variables integrated over a solid region yields a signed hypervolume, and the construction generalizes without further conceptual difficulty to functions on $\mathbb{R}^n$. This chapter develops the theory of multiple integration from its definition via partitions and Darboux sums, through Fubini's Theorem, which reduces the computation of a multiple integral to a sequence of ordinary single-variable integrations, to the change of variables formula, which generalizes $u$-substitution to several dimensions and relies essentially on the Jacobian matrix introduced in the preceding chapter.

\subsection{Construction of the Double Integral on Rectangles}

As in the one-dimensional theory, the integral is defined first over the simplest possible domain, a rectangle, and is subsequently extended to more general regions.

\begin{definition}[Partition of a Rectangle]
Consider the rectangle
\[
R = [a,b]\times[c,d] \subset \mathbb{R}^2.
\]
A partition of $R$ is a pair $P = (P_x, P_y)$, where
\[
\begin{gathered}
P_x : a = x_0 < x_1 < \cdots < x_m = b \\
P_y : c = y_0 < y_1 < \cdots < y_n = d
\end{gathered}
\]
are partitions of $[a,b]$ and $[c,d]$ respectively. The partition $P$ divides $R$ into $mn$ closed subrectangles
\[
R_{ij} = [x_{i-1},x_i]\times[y_{j-1},y_j],
\]
of area
\[
\Delta A_{ij} = \Delta x_i \, \Delta y_j = (x_i - x_{i-1})(y_j-y_{j-1}).
\]
The mesh of $P$ is
\[
\|P\| = \max_{i,j} \max(\Delta x_i, \Delta y_j).
\]

\end{definition}

\begin{definition}[Darboux Sums and the Double Integral]
Let $f : R \to \mathbb{R}$ be bounded, and let $P$ be a partition of $R$ with subrectangles $R_{ij}$. Set
\[
M_{ij} = \sup_{(x,y)\in R_{ij}} f(x,y), \qquad m_{ij} = \inf_{(x,y)\in R_{ij}} f(x,y),
\]
and define the upper and lower Darboux sums
\[
U(f,P) = \sum_{i,j} M_{ij}\,\Delta A_{ij}, \qquad L(f,P) = \sum_{i,j} m_{ij}\,\Delta A_{ij}.
\]
The function $f$ is (Darboux) integrable on $R$ if
\[
\sup_P L(f,P) = \inf_P U(f,P),
\]

In this case, the common value is called the double integral of $f$ over $R$. The following two notations are equivalent:
\[
\begin{gathered}
\iint_R f \, dA, \\
\iint_R f(x,y)\, dx\, dy.
\end{gathered}
\]

\end{definition}

Since refining a partition can only decrease upper sums and increase lower sums, one has
\[
L(f,P) \leq \sup_{P'} L(f,P') \leq \inf_{P'} U(f,P') \leq U(f,P)
\]
for every partition $P$, exactly as in the one-dimensional theory; the content of integrability is that these two bounding quantities coincide.

\begin{theorem}[Continuity Implies Integrability]
If $f$ is continuous on the closed rectangle
\[
R = [a,b]\times[c,d],
\]
then $f$ is integrable on $R$.
\end{theorem}

\begin{proof}
Since $R$ is closed and bounded, it is compact, and since $f$ is continuous on $R$, it is uniformly continuous on $R$ by the Heine--Cantor theorem: given $\varepsilon>0$, there exists $\delta>0$ such that $\|(x,y)-(x',y')\| < \delta$ implies
\[
|f(x,y)-f(x',y')| < \varepsilon / \operatorname{Area}(R).
\]
Choose any partition $P$ with mesh
\[
\|P\| < \delta/\sqrt{2},
\]
so that every subrectangle $R_{ij}$ has diameter
\[
((\Delta x_i)^2 + (\Delta y_j)^2)^{1/2} < \delta.
\]

On each $R_{ij}$, which is compact, $f$ attains its supremum $M_{ij}$ and infimum $m_{ij}$ at some points $\mathbf{p}_{ij}, \mathbf{q}_{ij} \in R_{ij}$ by the Extreme Value Theorem; since $\|\mathbf{p}_{ij}-\mathbf{q}_{ij}\| < \delta$, uniform continuity gives
\[
M_{ij}-m_{ij} = f(\mathbf{p}_{ij})-f(\mathbf{q}_{ij}) < \varepsilon/\operatorname{Area}(R).
\]
Consequently
\[
U(f,P) - L(f,P) = \sum_{i,j}(M_{ij}-m_{ij})\,\Delta A_{ij} < \frac{\varepsilon}{\operatorname{Area}(R)}\sum_{i,j}\Delta A_{ij} = \varepsilon,
\]
since
\[
\sum_{i,j}\Delta A_{ij} = \operatorname{Area}(R).
\]
For every partition, the following two inequalities hold:
\[
\begin{gathered}
\sup_{P'}L(f,P') \geq L(f,P) \\
\inf_{P'}U(f,P') \leq U(f,P)
\end{gathered}
\]
For the particular partition $P$ chosen above, they imply
\[
0 \leq \inf_{P'}U(f,P') - \sup_{P'}L(f,P') \leq U(f,P)-L(f,P) < \varepsilon
\]
Since $\varepsilon>0$ was arbitrary, it follows that
\[
\inf_{P'}U(f,P') = \sup_{P'}L(f,P'),
\]
that is, $f$ is integrable on $R$.
\end{proof}

Fubini's Theorem is the principal computational tool of multivariable integration: it asserts that, under mild hypotheses, a double integral may be evaluated as an iterated sequence of two ordinary one-dimensional integrals, performed one variable at a time.

\begin{theorem}[Fubini's Theorem, Rectangular Case]
Let $f$ be continuous on
\[
R = [a,b]\times[c,d].
\]

Then, for every fixed $x \in [a,b]$, the function $y \mapsto f(x,y)$ is continuous on $[c,d]$ and hence integrable there, and the resulting function
\[
A(x) = \int_c^d f(x,y)\, dy
\]
is continuous, hence integrable, on $[a,b]$. Moreover,
\[
\iint_R f\, dA = \int_a^b A(x)\, dx = \int_a^b \left( \int_c^d f(x,y)\, dy \right) dx.
\]
The same conclusion holds with the roles of $x$ and $y$ exchanged, so that in particular
\[
\int_a^b \left( \int_c^d f(x,y)\, dy \right) dx = \int_c^d \left( \int_a^b f(x,y)\, dx \right) dy.
\]

\end{theorem}

\begin{proof}
Fix $x \in [a,b]$; since $f$ is continuous on $R$, the single-variable function $y \mapsto f(x,y)$ is continuous on $[c,d]$, and therefore integrable there by the one-dimensional theory, so $A(x)$ is well defined. Let $P = (P_x,P_y)$ be any partition of $R$, with subrectangles
\[
R_{ij} = [x_{i-1},x_i]\times[y_{j-1},y_j],
\]
and retain the notation $M_{ij}, m_{ij}$ of the preceding section. Fix $i$, and let $x \in [x_{i-1},x_i]$. For every $j$ and every $y \in [y_{j-1},y_j]$, one has
\[
m_{ij} \leq f(x,y) \leq M_{ij},
\]
since $(x,y) \in R_{ij}$; integrating this inequality over $y \in [y_{j-1},y_j]$ gives
\[
m_{ij}\,\Delta y_j \;\leq\; \int_{y_{j-1}}^{y_j} f(x,y)\, dy \;\leq\; M_{ij}\,\Delta y_j.
\]

Summing over $j = 1, \dots, n$, and using additivity of the one-dimensional integral over the subintervals of $[c,d]$,
\[
\sum_{j} m_{ij}\,\Delta y_j \;\leq\; A(x) = \int_c^d f(x,y)\, dy \;\leq\; \sum_j M_{ij}\, \Delta y_j,
\]
valid for every $x \in [x_{i-1},x_i]$. Integrating this inequality over $x \in [x_{i-1},x_i]$ and summing over $i = 1, \dots, m$ gives
\[
\sum_{i,j} m_{ij}\, \Delta x_i \, \Delta y_j \;\leq\; \int_a^b A(x)\, dx \;\leq\; \sum_{i,j} M_{ij}\,\Delta x_i\, \Delta y_j,
\]
For every partition $P$ of $R$, this gives the bounds
\[
L(f,P) \;\leq\; \int_a^b A(x)\, dx \;\leq\; U(f,P),
\]
Since $f$ is continuous on $R$, it is integrable there by the previous theorem, so
\[
\sup_P L(f,P) = \inf_P U(f,P) = \iint_R f\, dA;
\]
as the fixed number
\[
\int_a^b A(x)\, dx
\]
is squeezed between $L(f,P)$ and $U(f,P)$ for every partition $P$, it must equal this common value, giving
\[
\int_a^b A(x)\, dx = \iint_R f\, dA.
\]

Continuity of $A$ follows from uniform continuity of $f$ on the compact rectangle $R$: given $\varepsilon>0$, choose $\delta>0$ so that $\|(x,y)-(x',y)\|<\delta$ implies
\[
|f(x,y)-f(x',y)| < \varepsilon/(d-c)
\]
for all $y$; then $|x-x'|<\delta$ implies
\[
|A(x)-A(x')| \leq \int_c^d |f(x,y)-f(x',y)|\, dy < \varepsilon,
\]
establishing continuity, hence integrability, of $A$. The identity with the roles of $x$ and $y$ exchanged follows by the identical argument applied with the two coordinates interchanged throughout.
\end{proof}

\begin{example}
Let $R = [0,1]\times[0,2]$ and $f(x,y) = xy^2$. By Fubini's Theorem,
\begin{align*}
\iint_R xy^2\,\dd A
&= \int_0^1 \left(\int_0^2 xy^2\,\dd y\right)\dd x \\
&= \int_0^1 x\left[\frac{y^3}{3}\right]_{0}^{2}\dd x \\
&= \int_0^1 \frac{8x}{3}\,\dd x
= \frac{8}{3}\left[\frac{x^2}{2}\right]_{0}^{1}
= \frac{4}{3}.
\end{align*}
Integrating in the opposite order gives the same value:
\begin{align*}
\int_0^2\left(\int_0^1 xy^2\,\dd x\right)\dd y
&= \int_0^2 \frac{y^2}{2}\,\dd y \\
&= \left[\frac{y^3}{6}\right]_{0}^{2}
= \frac{4}{3},
\end{align*}
confirming the theorem.
\end{example}

\subsection{Integrability on General Planar Regions}


Most regions of interest are not rectangles, and the integral must be extended to accommodate them.

\begin{definition}[Content Zero]
A bounded set $Z \subset \mathbb{R}^2$ has content zero if, for every $\varepsilon>0$, there exist finitely many rectangles $S_1,\dots,S_k$ with
\[
\begin{gathered}
Z \subseteq \bigcup_{l=1}^{k} S_l \\
\sum_{l=1}^{k}\operatorname{Area}(S_l) < \varepsilon.
\end{gathered}
\]

\end{definition}

The graph of a continuous function on a closed bounded interval has content zero, a fact used implicitly below without separate proof, since it follows from uniform continuity in the same manner as the theorem above: covering $[a,b]$ by finitely many subintervals of total length $\varepsilon'$ and small enough that the corresponding oscillation of the function is at most $\varepsilon/(b-a)$, the rectangles built from these subintervals and this oscillation bound cover the graph with total area at most $\varepsilon$.

\begin{theorem}[Integrability with a Content-Zero Set of Discontinuities]
Let $\tilde{f} : R \to \mathbb{R}$ be bounded on a rectangle $R$, say $|\tilde{f}| \leq M$, and continuous at every point of $R$ except possibly on a set $Z \subset R$ of content zero. Then $\tilde{f}$ is integrable on $R$.
\end{theorem}

\begin{proof}
Let $\varepsilon > 0$. Since $Z$ has content zero, cover $Z$ by finitely many open rectangles $S_1,\dots,S_k$ with
\[
\sum_l \operatorname{Area}(S_l) < \varepsilon/(4M).
\]

The set $K = R \setminus \bigcup_l S_l$ is closed and bounded, hence compact, and $\tilde{f}$ is continuous at every point of $K$, since $K$ contains no point of $Z$; $\tilde{f}$ is therefore uniformly continuous on $K$. Choose a partition $P$ of $R$ fine enough that: (i) every subrectangle of $P$ meeting $\bigcup_l S_l$ is contained in a slightly enlarged union of the $S_l$ whose total area remains below $\varepsilon/(2M)$, which is possible since finitely many rectangles have a boundary of content zero and the $S_l$ may be taken slightly larger at the outset to absorb this adjustment; and (ii) the mesh of $P$ is small enough that, on every subrectangle disjoint from $\bigcup_l S_l$ (hence contained in $K$), the oscillation of $\tilde f$ is less than $\varepsilon/(2\operatorname{Area}(R))$, which is possible by uniform continuity of $\tilde f$ on $K$. Splitting the sum defining $U(f,P)-L(f,P)$ into subrectangles of the first kind and of the second kind,
\[
U(\tilde f,P) - L(\tilde f,P) = \sum_{\text{meets } \cup S_l} (M_{ij}-m_{ij})\Delta A_{ij} \;+\; \sum_{\text{in } K} (M_{ij}-m_{ij}) \Delta A_{ij}.
\]

In the first sum, $M_{ij}-m_{ij} \leq 2M$ always, so this sum is at most $2M \cdot \varepsilon/(4M) = \varepsilon/2$; in the second sum, $M_{ij}-m_{ij} < \varepsilon/(2\operatorname{Area}(R))$, so this sum is less than
\[
\varepsilon/(2\operatorname{Area}(R)) \cdot \operatorname{Area}(R) = \varepsilon/2.
\]

Hence $U(\tilde f,P)-L(\tilde f,P) < \varepsilon$; since $\varepsilon>0$ was arbitrary, $\tilde f$ is integrable on $R$, exactly as in the proof of the theorem for continuous functions above.
\end{proof}

\begin{definition}[Integral over a General Bounded Region]
Let $D \subset \mathbb{R}^2$ be bounded, with $\partial D$ of content zero, and let $f$ be bounded on $D$ and continuous except possibly on a set of content zero. Choose any rectangle $R \supseteq D$, and extend $f$ to $\tilde f$ on $R$ by setting $\tilde f = f$ on $D$ and $\tilde f = 0$ on $R \setminus D$. The double integral of $f$ over $D$ is defined as
\[
\iint_D f\, dA := \iint_R \tilde f\, dA.
\]

\end{definition}

The discontinuities of $\tilde f$ occur only within $\partial D$, together with the discontinuities of $f$ already present in $D$, together forming a set of content zero; the preceding theorem guarantees that $\tilde f$ is integrable on $R$, so this definition is meaningful, and it is independent of the choice of enclosing rectangle $R$, since $\tilde f$ vanishes outside $D$ regardless.

\begin{definition}[Type I and Type II Regions]
A region $D \subset \mathbb{R}^2$ is of Type I if
\[
D = \{(x,y) : a \leq x \leq b, \; g_1(x) \leq y \leq g_2(x)\}
\]
for continuous functions $g_1 \leq g_2$ on $[a,b]$; it is of Type II if
\[
D = \{(x,y) : c \leq y \leq d, \; h_1(y) \leq x \leq h_2(y)\}
\]
for continuous functions $h_1 \leq h_2$ on $[c,d]$.
\end{definition}

For a Type I region, $\partial D$ consists of the graphs of $g_1$ and $g_2$, each of content zero as noted above, together with the two vertical segments $x=a$ and $x=b$, also visibly of content zero; hence $\partial D$ has content zero, being a finite union of sets of content zero, and the integral over $D$ is defined. Applying Fubini's Theorem on an enclosing rectangle $R = [a,b]\times[c,d]$ to the extension $\tilde f$, and noting that $\tilde f(x,y) = 0$ for $y \notin [g_1(x),g_2(x)]$, gives the following practical evaluation formula.

\begin{theorem}[Iterated Integral over a Type I Region]
If $D$ is a Type I region as above and $f$ is continuous on $D$, then
\[
\iint_D f\, dA = \int_a^b \left( \int_{g_1(x)}^{g_2(x)} f(x,y)\, dy \right) dx.
\]
An analogous formula, with the roles of $x$ and $y$ exchanged, holds for a Type II region.
\end{theorem}

\begin{proof}
Let $R = [a,b]\times[c,d]$ enclose $D$, with $c \leq g_1(x)$ and $g_2(x) \leq d$ for all $x \in [a,b]$. By Fubini's Theorem applied to $\tilde f$ on $R$,
\[
\iint_D f \, dA = \iint_R \tilde f \, dA = \int_a^b \left( \int_c^d \tilde f(x,y)\, dy \right) dx.
\]

For fixed $x \in [a,b]$, $\tilde f(x,y) = f(x,y)$ for $y \in [g_1(x),g_2(x)]$ and $\tilde f(x,y) = 0$ otherwise, so by additivity of the one-dimensional integral over subintervals,
\[
\int_c^d \tilde f(x,y)\, dy = \int_c^{g_1(x)} 0 \, dy + \int_{g_1(x)}^{g_2(x)} f(x,y)\, dy + \int_{g_2(x)}^d 0\, dy = \int_{g_1(x)}^{g_2(x)} f(x,y)\, dy,
\]
which, substituted into the displayed identity above, gives the claimed formula.
\end{proof}

\begin{theorem}[Basic Properties of the Double Integral]
Let $f,g$ be integrable on a region $D$, and let $\alpha,\beta \in \mathbb{R}$. Then:
\begin{enumerate}
\item[(i)] \textit{Linearity}: the following identity holds:

\[
\iint_D (\alpha f + \beta g)\, dA = \alpha \iint_D f\, dA + \beta \iint_D g\, dA.
\]

\item[(ii)] \textit{Monotonicity}: if $f \leq g$ on $D$, then

\[
\iint_D f\, dA \leq \iint_D g\, dA.
\]

\item[(iii)] \textit{Additivity over regions}: if $D = D_1 \cup D_2$ with $D_1, D_2$ overlapping at most on a set of content zero, then

\[
\iint_D f\, dA = \iint_{D_1} f\, dA + \iint_{D_2} f\, dA.
\]

\end{enumerate}
\end{theorem}

\begin{proof}
(i) and (ii) follow directly from the corresponding properties of upper and lower Darboux sums: since
\[
M_{ij}[\alpha f + \beta g] = \alpha M_{ij}[f] + \beta M_{ij}[g]
\]
whenever $\alpha,\beta \geq 0$ (with the analogous relation for $m_{ij}$ when either scalar is negative, upon exchanging the roles of $M_{ij}$ and $m_{ij}$), the Darboux sums of $\alpha f + \beta g$ are the corresponding linear combinations of the Darboux sums of $f$ and $g$ for every partition, and passing to the common limiting value defining the integral yields (i); (ii) follows since $f \leq g$ on $D$ forces $L(f,P) \leq L(g,P)$ for every partition $P$ of an enclosing rectangle (applied to the zero-extensions), hence
\[
\iint_D f\,dA = \sup_P L(f,P) \leq \sup_P L(g,P) = \iint_D g\, dA.
\]

For (iii), extend $f$ by zero to a common enclosing rectangle $R$; the zero-extension of $f$ from $D$ agrees with the sum of the zero-extensions from $D_1$ and from $D_2$ except on $D_1 \cap D_2$, a set of content zero which therefore does not affect any Darboux sum in the limit, by the same estimate used in the proof of the integrability theorem above; hence
\[
\iint_D f \, dA = \iint_{D_1} f\, dA + \iint_{D_2} f \, dA
\]
by linearity applied to the (almost everywhere) sum of the two zero-extensions.
\end{proof}

Figure~\ref{fig:ch11-typeI-region} shows the vertical slices used to evaluate an integral over a Type I region:
\[
\iint_D f\,dA.
\]

\begin{figure}[!htbp]
\centering
\begin{tikzpicture}[scale=1.15,>=Stealth]
  \draw[->] (-0.4,0) -- (4.4,0) node[right,font=\scriptsize]{$x$};
  \draw[->] (0,-0.4) -- (0,3.3) node[above,font=\scriptsize]{$y$};
  \fill[figblue,opacity=0.30]
    plot[domain=0.6:3.6,smooth,variable=\x] ({\x},{0.35+0.18*(\x-0.6)*(\x-0.6)})
    -- plot[domain=3.6:0.6,smooth,variable=\x] ({\x},{2.9-0.28*(\x-2.1)*(\x-2.1)})
    -- cycle;
  \draw[vecB,thick,domain=0.6:3.6,smooth,variable=\x]
    plot ({\x},{0.35+0.18*(\x-0.6)*(\x-0.6)}) node[right,font=\scriptsize]{$y=g_1(x)$};
  \draw[vecC,thick,domain=0.6:3.6,smooth,variable=\x]
    plot ({\x},{2.9-0.28*(\x-2.1)*(\x-2.1)}) node[right,font=\scriptsize]{$y=g_2(x)$};
  \draw[dashed,gray] (0.6,0) -- (0.6,0.35) node[pos=0,below,font=\scriptsize,black]{$a$};
  \draw[dashed,gray] (3.6,0) -- (3.6,1.62);
  \node[below,font=\scriptsize] at (3.6,0) {$b$};
  \node[font=\scriptsize] at (2.0,1.5) {$D$};
\end{tikzpicture}
\caption[Type I region]{A Type I region bounded by two graphs and two vertical lines. Integrating along each vertical slice first, and then across the horizontal interval, gives the double integral over $D$.}
\label{fig:ch11-typeI-region}
\end{figure}

Figure~\ref{fig:atlas-triangular-domain} compares vertical and horizontal slices of the same integration region.

\begin{figure}[!htbp]
\centering
\begin{tikzpicture}[>=Stealth]
\begin{scope}[scale=3.6]
\fill[figblue!15] (0,0) -- (1,0) -- (0,1) -- cycle;
\draw[->] (-0.1,0) -- (1.2,0) node[right] {$x$};
\draw[->] (0,-0.1) -- (0,1.2) node[above] {$y$};
\draw[figblue,thick] (0,1) -- (1,0);
\draw[figred,very thick] (0.35,0) -- (0.35,0.65);
\node[below] at (0.35,0) {$x$};
\node[above right] at (0.35,0.65) {$1-x$};
\node[below] at (1,0) {$1$};
\node[left] at (0,1) {$1$};
\end{scope}
\begin{scope}[shift={(6,0)},scale=3.6]
\fill[figblue!15] (0,0) -- (1,0) -- (0,1) -- cycle;
\draw[->] (-0.1,0) -- (1.2,0) node[right] {$x$};
\draw[->] (0,-0.1) -- (0,1.2) node[above] {$y$};
\draw[figblue,thick] (0,1) -- (1,0);
\draw[figgreen,very thick] (0,0.4) -- (0.6,0.4);
\node[left] at (0,0.4) {$y$};
\node[above right] at (0.6,0.4) {$1-y$};
\node[below] at (1,0) {$1$};
\node[left] at (0,1) {$1$};
\end{scope}
\end{tikzpicture}
\caption{The triangular region is bounded by the coordinate axes and $x+y=1$. Vertical slices use $0\le y\le1-x$, whereas horizontal slices use $0\le x\le1-y$. Both descriptions cover the same region and lead to equivalent iterated integrals.}
\label{fig:atlas-triangular-domain}
\end{figure}

\subsection{Change of Variables in the Plane}

The one-dimensional substitution rule
\[
\int_a^b f(x)\,dx = \int_{\alpha}^{\beta} f(\varphi(u))\,\varphi'(u)\,du,
\]
valid for a differentiable bijection $\varphi$ with $\varphi(\alpha)=a$, $\varphi(\beta)=b$, admits a direct generalization to $\mathbb{R}^2$, in which the scalar factor $\varphi'(u)$ is replaced by the absolute value of the Jacobian determinant introduced in the preceding chapter.

\begin{theorem}[Change of Variables Formula, Two Dimensions]
Let $T : U \to V$,
\[
T(u,v) = (x(u,v),y(u,v)),
\]
be a bijective $C^1$ map between open sets $U,V \subset \mathbb{R}^2$, with $\det J_T(u,v) \neq 0$ for every $(u,v) \in U$. Let $D \subset U$ be a bounded region with $D' = T(D) \subset V$, and let $f$ be continuous on $D'$. Then
\[
\iint_{D'} f(x,y)\, dx\, dy = \iint_D f\big(T(u,v)\big)\, |\det J_T(u,v)|\, du\, dv.
\]

\end{theorem}

\begin{proof}[Derivation]
Partition $D$ by a fine grid of small rectangles
\[
R_{ij} = [u_{i-1},u_i]\times[v_{j-1},v_j],
\]
of dimensions $\Delta u \times \Delta v$, and fix a corner $(u_i,v_j)$ of $R_{ij}$. By differentiability of $T$ at $(u_i,v_j)$ -- precisely the statement that $J_T(u_i,v_j)$ furnishes the best local linear approximation to $T$, established in the preceding chapter -- the image of $R_{ij}$ under $T$ satisfies
\begin{align*}
T(u_i + \Delta u, v_j) &\approx T(u_i,v_j) + \frac{\partial T}{\partial u}(u_i,v_j)\,\Delta u, \\
T(u_i, v_j+\Delta v) &\approx T(u_i,v_j) + \frac{\partial T}{\partial v}(u_i,v_j)\, \Delta v,
\end{align*}
with error terms that are $o(\Delta u)$ and $o(\Delta v)$ respectively as $\Delta u, \Delta v \to 0$. Consequently, $T(R_{ij})$ is approximated, to leading order, by the parallelogram based at $T(u_i,v_j)$ and spanned by the two vectors $\partial T/\partial u\, \Delta u$ and $\partial T/\partial v\, \Delta v$, which are precisely the two columns of $J_T(u_i,v_j)$, each scaled respectively by $\Delta u$ and $\Delta v$. The area of a parallelogram spanned by two vectors equals the absolute value of the determinant of the matrix formed by those vectors as columns; hence
\begin{align*}
\operatorname{Area}\big(T(R_{ij})\big) &\;\approx\; \left| \det \begin{pmatrix} \partial T/\partial u \, \Delta u & \partial T/\partial v \, \Delta v \end{pmatrix} \right| \\
&= |\det J_T(u_i,v_j)| \, \Delta u\, \Delta v = |\det J_T(u_i,v_j)|\, \operatorname{Area}(R_{ij}),
\end{align*}
using multilinearity of the determinant to factor $\Delta u$ and $\Delta v$ out of the respective columns. As the mesh of the partition tends to zero, the images $T(R_{ij})$ tile $D'$ with overlaps and gaps of area tending to zero relative to $\operatorname{Area}(D')$, by injectivity and continuity of $T$, so that
\begin{align*}
\iint_{D'} f\, dA &= \lim_{\text{mesh}\to 0} \sum_{i,j} f\big(T(u_i,v_j)\big)\, \operatorname{Area}\big(T(R_{ij})\big) \\
&\;\approx\; \lim_{\text{mesh}\to 0} \sum_{i,j} f\big(T(u_i,v_j)\big)\, |\det J_T(u_i,v_j)|\, \Delta u \, \Delta v,
\end{align*}
and the right-hand side is precisely a Riemann sum, over the partition $\{R_{ij}\}$ of $D$, for the integrand $f(T(u,v))\,|\det J_T(u,v)|$, which is continuous on $D$ since $f$, $T$, and $\det J_T$ are all continuous, hence integrable there. Passing to the limit converts both approximate equalities into exact ones, since the accumulated error, controlled by the uniform continuity of $J_T$ on the compact closure of $D$, vanishes in the limit; we omit the fully detailed uniform error estimates, which require only a routine, if tedious, application of uniform continuity of $DT$ on compact subsets. This establishes the formula.
\end{proof}

\begin{corollary}[Polar Coordinates]
Define the polar-coordinate map by
\[
T(r,\theta) = (r\cos\theta, r\sin\theta).
\]

As computed in the example of the preceding chapter, $\det J_T(r,\theta) = r$ (see Figure~\ref{fig:ch10-jacobian-polar}), so that for a region $D'$ described in polar coordinates by
\[
D = \{(r,\theta) : \alpha \leq \theta \leq \beta, \; r_1(\theta) \leq r \leq r_2(\theta)\},
\]

\[
\iint_{D'} f(x,y)\, dx\, dy = \int_{\alpha}^{\beta} \int_{r_1(\theta)}^{r_2(\theta)} f(r\cos\theta, r\sin\theta)\; r\, dr\, d\theta.
\]

\end{corollary}

\begin{example}[The Gaussian Integral]
The improper integral
\[
I = \int_{-\infty}^{\infty} e^{-x^2}\, dx
\]
admits no elementary antiderivative, yet its value can be found exactly by passing to two dimensions and then to polar coordinates. Since
\[
I = \int_{-\infty}^\infty e^{-y^2}\,dy
\]
as well, Fubini's Theorem (extended to the improper setting, valid here since the integrand is nonnegative) gives
\[
I^2 = \left(\int_{-\infty}^{\infty} e^{-x^2}\, dx\right)\left(\int_{-\infty}^{\infty} e^{-y^2}\, dy\right) = \iint_{\mathbb{R}^2} e^{-x^2-y^2}\, dA.
\]

Passing to polar coordinates, with $x^2+y^2 = r^2$ and the region $\mathbb{R}^2$ described by $r \in [0,\infty)$, $\theta \in [0,2\pi)$,
\[
\begin{aligned}
I^2 &= \int_0^{2\pi}\int_0^{\infty} e^{-r^2}\, r\, dr\, d\theta \\
&= 2\pi \int_0^{\infty} e^{-r^2} r\, dr \\
&= 2\pi \left[ -\frac{1}{2} e^{-r^2} \right]_0^{\infty}
= 2\pi \cdot \frac{1}{2} = \pi,
\end{aligned}
\]
where the substitution $t=r^2$, $dt = 2r\,dr$ converts
\[
\int_0^\infty e^{-r^2}r\,dr
\]
into the elementary
\[
\frac12\int_0^\infty e^{-t}\,dt = \frac12.
\]
Since $I > 0$, this gives the classical identity
\[
\int_{-\infty}^{\infty} e^{-x^2}\, dx = \sqrt{\pi}.
\]

\end{example}

\section{Triple Integrals and Spatial Coordinates}

\subsection{Triple Integrals over Boxes}

The construction of Section~\ref{sec:double-integrals} extends verbatim to three dimensions, replacing rectangles by boxes and areas by volumes.

\begin{definition}[Triple Integral over a Box]
Consider the box
\[
B = [a_1,b_1]\times[a_2,b_2]\times[a_3,b_3] \subset \mathbb{R}^3,
\]

Let $f : B \to \mathbb{R}$ be bounded. A partition of $B$ is obtained by partitioning each edge, dividing $B$ into subboxes $B_{ijk}$ of volume $\Delta V_{ijk}$. With
\[
\begin{gathered}
M_{ijk} = \sup_{B_{ijk}} f \\
m_{ijk} = \inf_{B_{ijk}} f,
\end{gathered}
\]
define
\[
\begin{gathered}
U(f,P) = \sum M_{ijk}\Delta V_{ijk} \\
L(f,P) = \sum m_{ijk} \Delta V_{ijk}
\end{gathered}
\]
exactly as before; $f$ is integrable on $B$ if
\[
\sup_P L(f,P) = \inf_P U(f,P),
\]
and this common value is the triple integral
\[
\iiint_B f\, dV.
\]

\end{definition}

The proofs that continuous functions on a box are integrable, and that integrability persists under a content-zero (now, volume-zero) set of discontinuities, are entirely analogous to the two-dimensional case, replacing the Heine--Cantor and covering arguments given above by their immediate three-variable counterparts; we do not repeat the details. The corresponding version of Fubini's Theorem likewise follows the same squeeze argument as before, applied with one variable singled out at a time.

\begin{theorem}[Fubini's Theorem, Box Case]
If $f$ is continuous on
\[
B = [a_1,b_1]\times[a_2,b_2]\times[a_3,b_3],
\]
Fubini's formula gives
\[
\iiint_B f\, dV = \int_{a_1}^{b_1}\int_{a_2}^{b_2}\int_{a_3}^{b_3} f(x,y,z)\; dz\, dy\, dx,
\]

The iterated integral may equally be evaluated in any of the six orders obtained by permuting $x,y,z$.
\end{theorem}

\begin{proof}
Group the last two variables together: writing
\[
\begin{gathered}
\mathbf{y} = (y,z) \in \mathbb{R}^2 \\
R = [a_2,b_2]\times[a_3,b_3],
\end{gathered}
\]
the two-dimensional Fubini Theorem, applied to the function $\mathbf{y} \mapsto f(x,\mathbf{y})$ for fixed $x$, together with the identical squeeze argument of the rectangular case applied with $x$ playing the role of the single ``outer'' variable and $R$ playing the role of the two-dimensional ``inner'' rectangle, yields
\[
\iiint_B f\, dV = \int_{a_1}^{b_1}\left(\iint_R f(x,y,z)\, dy\,dz\right) dx;
\]
applying the two-dimensional Fubini Theorem once more to the inner double integral, for each fixed $x$, resolves it into the iterated integral over $y$ then $z$ (or $z$ then $y$), giving the stated formula. Since the grouping of variables was arbitrary, the same argument applied with a different variable singled out as ``outer'' establishes the remaining orders of integration.
\end{proof}

\subsection{Triple Integrals over General Regions}

\begin{definition}[Type I Solid Region]
A solid region $E \subset \mathbb{R}^3$ is of Type I if
\[
E = \big\{ (x,y,z) : (x,y) \in D, \; u_1(x,y) \leq z \leq u_2(x,y) \big\}
\]
for a Type I or Type II planar region $D \subset \mathbb{R}^2$ and continuous functions $u_1 \leq u_2$ on $D$. Regions defined analogously with the roles of $x,y,z$ permuted are of Type II and Type III.
\end{definition}

As in the planar case, extending $f$ by zero outside $E$ and applying Fubini's Theorem on an enclosing box yields, for a Type I region,
\[
\iiint_E f\, dV = \iint_D \left( \int_{u_1(x,y)}^{u_2(x,y)} f(x,y,z)\, dz \right) dA,
\]
and the inner double integral over $D$ is then evaluated by the planar methods of the preceding sections.

\begin{example}[Volume of a Tetrahedron]
Let $E$ be the solid tetrahedron bounded by the coordinate planes $x=0$, $y=0$, $z=0$ and the plane $x+y+z=1$. Regarding $E$ as a Type I region over the planar triangle
\[
D = \{(x,y): x\geq0,\,y\geq0,\,x+y\leq1\},
\]
The vertical bounds over this base are
\[
\begin{aligned}
0 &\leq z \\
&\leq 1-x-y,
\end{aligned}
\]

\[
\operatorname{Vol}(E) = \iiint_E dV = \iint_D (1-x-y)\, dA = \int_0^1 \int_0^{1-x} (1-x-y)\, dy\, dx.
\]
The inner integral is
\[
\int_0^{1-x}(1-x-y)\,dy = \left[(1-x)y - \tfrac{y^2}{2}\right]_0^{1-x} = (1-x)^2 - \tfrac{(1-x)^2}{2} = \tfrac{(1-x)^2}{2},
\]
so
\[
\operatorname{Vol}(E) = \int_0^1 \frac{(1-x)^2}{2}\, dx = \left[-\frac{(1-x)^3}{6}\right]_0^1 = \frac{1}{6},
\]
matching the familiar formula
\[
\operatorname{Vol} = (1/6)|abc|
\]
for a coordinate tetrahedron with
\[
\begin{aligned}
a &= b \\
&= c \\
&= 1.
\end{aligned}
\]

\end{example}

\subsection{Cylindrical Coordinates}


The three-dimensional change of variables formula is the direct generalization of the planar one, with the parallelogram of the earlier derivation replaced by a parallelepiped, whose volume is likewise the absolute value of the determinant of the matrix formed by its three spanning vectors -- here, the three columns of the Jacobian matrix, scaled by the corresponding side lengths of the small box.

\begin{theorem}[Change of Variables Formula, Three Dimensions]
Let $T : U \to V$ be a bijective $C^1$ map between open sets $U,V \subset \mathbb{R}^3$ with $\det J_T \neq 0$ on $U$. Let $E \subset U$ with $E' = T(E) \subset V$, and let $f$ be continuous on $E'$. Then
\[
\iiint_{E'} f\, dV = \iiint_E f\big(T(u,v,w)\big)\, |\det J_T(u,v,w)| \; du\, dv\, dw.
\]

\end{theorem}

\begin{proof}[Derivation]
The argument is identical, mutatis mutandis, to the two-dimensional derivation above: partitioning $E$ into small boxes, approximating the image of each box under $T$ by the parallelepiped spanned by the three columns of $J_T$ scaled by $\Delta u, \Delta v, \Delta w$, whose volume is $|\det J_T|\,\Delta u\,\Delta v\,\Delta w$, and passing to the limit of Riemann sums; we do not repeat the details.
\end{proof}

Cylindrical coordinates $(r,\theta,z)$ describe a point via its polar coordinates $(r,\theta)$ in the $xy$-plane together with its height $z$, via the map
\[
T(r,\theta,z) = (r\cos\theta, \, r\sin\theta, \, z).
\]

\begin{proposition}[Jacobian of Cylindrical Coordinates]
$\det J_T(r,\theta,z) = r$.
\end{proposition}

\begin{proof}
The Jacobian matrix, with columns indexed by $r,\theta,z$ and rows by $x,y,z$, is
\[
J_T(r,\theta,z) = \begin{pmatrix} \cos\theta & -r\sin\theta & 0 \\ \sin\theta & r\cos\theta & 0 \\ 0 & 0 & 1 \end{pmatrix}.
\]

Expanding the determinant along the third row, which has only the single nonzero entry $1$ in the $(3,3)$ position, gives
\begin{align*}
\det J_T &= 1 \cdot \det \begin{pmatrix} \cos\theta & -r\sin\theta \\ \sin\theta & r\cos\theta \end{pmatrix} = \cos\theta(r\cos\theta) - (-r\sin\theta)(\sin\theta) \\
&= r\cos^2\theta + r\sin^2\theta = r,
\end{align*}
as claimed.
\end{proof}

Hence, since $r \geq 0$ throughout the natural domain of cylindrical coordinates, $|\det J_T| = r$, and the change of variables formula reads
\[
\iiint_{E'} f(x,y,z)\, dV = \iiint_E f(r\cos\theta, r\sin\theta, z) \; r\, dr\, d\theta\, dz.
\]

\begin{example}[Volume of a Solid Cylinder by Cylindrical Coordinates]
Let $E'$ be the solid cylinder $x^2+y^2 \leq R^2$,
\[
\begin{aligned}
0 &\leq z \\
&\leq h.
\end{aligned}
\]
In cylindrical coordinates, $E$ is the box
\[
\begin{aligned}
0 &\leq r \\
&\leq R,
\end{aligned}
\]

\[
\begin{aligned}
0 &\leq \theta \\
&\leq 2\pi,
\end{aligned}
\]

\[
\begin{aligned}
0 &\leq z \\
&\leq h,
\end{aligned}
\]
so
\[
\operatorname{Vol}(E') = \int_0^{2\pi}\int_0^R \int_0^h r\, dz\, dr\, d\theta = 2\pi \cdot h \cdot \int_0^R r\, dr = 2\pi h \cdot \frac{R^2}{2} = \pi R^2 h,
\]
recovering the elementary formula for the volume of a cylinder.
\end{example}

\subsection{Spherical Coordinates}

Spherical coordinates $(\rho,\varphi,\theta)$ describe a point by its distance $\rho \geq 0$ from the origin, the polar angle $\varphi \in [0,\pi]$ measured from the positive $z$-axis, and the azimuthal angle $\theta \in [0,2\pi)$ measured in the $xy$-plane from the positive $x$-axis, via
\[
T(\rho,\varphi,\theta) = \big( \rho\sin\varphi\cos\theta, \; \rho\sin\varphi\sin\theta, \; \rho\cos\varphi \big).
\]

\begin{proposition}[Jacobian of Spherical Coordinates]

\[
\det J_T(\rho,\varphi,\theta) = \rho^2 \sin\varphi.
\]

\end{proposition}

\begin{proof}
Differentiating each component of $T$ with respect to each of $\rho,\varphi,\theta$ gives the Jacobian matrix
\[
\setlength{\arraycolsep}{3pt}
J_T = \begin{pmatrix} \sin\varphi\cos\theta & \rho\cos\varphi\cos\theta & -\rho\sin\varphi\sin\theta \\ \sin\varphi\sin\theta & \rho\cos\varphi\sin\theta & \rho\sin\varphi\cos\theta \\ \cos\varphi & -\rho\sin\varphi & 0 \end{pmatrix}.
\]

Expand the determinant along the third row. The $(3,3)$ entry is $0$, so only the first two entries of that row contribute:
\begin{align*}
\setlength{\arraycolsep}{3pt}
\det J_T &= \cos\varphi \cdot \det\begin{pmatrix} \rho\cos\varphi\cos\theta & -\rho\sin\varphi\sin\theta \\ \rho\cos\varphi\sin\theta & \rho\sin\varphi\cos\theta \end{pmatrix} \\
&\quad -\; (-\rho\sin\varphi)\cdot \det\begin{pmatrix} \sin\varphi\cos\theta & -\rho\sin\varphi\sin\theta \\ \sin\varphi\sin\theta & \rho\sin\varphi\cos\theta \end{pmatrix}.
\end{align*}
The first minor equals
\begin{align*}
&(\rho\cos\varphi\cos\theta)(\rho\sin\varphi\cos\theta) - (-\rho\sin\varphi\sin\theta)(\rho\cos\varphi\sin\theta) \\
&\quad= \rho^2\cos\varphi\sin\varphi\cos^2\theta + \rho^2\sin\varphi\cos\varphi\sin^2\theta = \rho^2\sin\varphi\cos\varphi,
\end{align*}
using $\cos^2\theta+\sin^2\theta=1$; hence the first term is
\[
\cos\varphi \cdot \rho^2\sin\varphi\cos\varphi = \rho^2\sin\varphi\cos^2\varphi.
\]
The second minor equals
\begin{align*}
&(\sin\varphi\cos\theta)(\rho\sin\varphi\cos\theta) - (-\rho\sin\varphi\sin\theta)(\sin\varphi\sin\theta) \\
&\quad= \rho\sin^2\varphi\cos^2\theta + \rho\sin^2\varphi\sin^2\theta = \rho\sin^2\varphi,
\end{align*}
so the second term is
\[
\rho\sin\varphi \cdot \rho\sin^2\varphi = \rho^2\sin^3\varphi.
\]
Adding the two contributions,
\[
\begin{aligned}
\det J_T &= \rho^2\sin\varphi\cos^2\varphi + \rho^2\sin^3\varphi \\
&= \rho^2\sin\varphi\big(\cos^2\varphi+\sin^2\varphi\big) \\
&= \rho^2\sin\varphi,
\end{aligned}
\]
as claimed.
\end{proof}

Since $\varphi \in [0,\pi]$, $\sin\varphi \geq 0$, so $|\det J_T| = \rho^2\sin\varphi$, and the change of variables formula reads
\begin{align*}
\iiint_{E'} f(x,y,z)\, dV &= \iiint_E f\big(\rho\sin\varphi\cos\theta, \rho\sin\varphi\sin\theta, \rho\cos\varphi\big) \\
&\quad\times\; \rho^2\sin\varphi\; d\rho\, d\varphi\, d\theta.
\end{align*}

\begin{example}[Volume of a Ball by Spherical Coordinates]
Let $E'$ be the solid ball $x^2+y^2+z^2 \leq R^2$. In spherical coordinates, $E$ is the box
\[
\begin{aligned}
0 &\leq \rho \\
&\leq R,
\end{aligned}
\]

\[
\begin{aligned}
0 &\leq \varphi \\
&\leq \pi,
\end{aligned}
\]

\[
\begin{aligned}
0 &\leq \theta \\
&\leq 2\pi,
\end{aligned}
\]
so
\begin{align*}
\operatorname{Vol}(E') &= \int_0^{2\pi}\!\!\int_0^{\pi}\!\!\int_0^{R} \rho^2\sin\varphi \; d\rho\, d\varphi\, d\theta \\
&= \left(\int_0^{2\pi} d\theta\right)\left(\int_0^\pi \sin\varphi\, d\varphi\right)\left(\int_0^R \rho^2\, d\rho\right) = 2\pi \cdot 2 \cdot \frac{R^3}{3} = \frac{4}{3}\pi R^3,
\end{align*}
where the three factors separate because the integrand is a product of a function of $\rho$ alone and a function of $\varphi$ alone, and Fubini's Theorem applies on the box $[0,R]\times[0,\pi]\times[0,2\pi]$; this recovers the classical formula for the volume of a ball.
\end{example}

Figure~\ref{fig:ch11-spherical} identifies the radial distance and the two angles used in spherical coordinates.

\begin{figure}[!htbp]
\centering
\begin{tikzpicture}[scale=1.5,>=Stealth]
  \coordinate (O) at (0,0);
  \coordinate (Xax) at (-2.05,-1.35);
  \coordinate (Yax) at (2.85,0);
  \coordinate (Zax) at (0,3.15);
  \coordinate (P) at (1.8,2.05);
  \coordinate (Q) at (1.8,0.48);
  \draw[->] (O) -- (Xax) node[left,font=\scriptsize]{$x$};
  \draw[->] (O) -- (Yax) node[right,font=\scriptsize]{$y$};
  \draw[->] (O) -- (Zax) node[above,font=\scriptsize]{$z$};
  \draw[dashed,gray!75] (O) -- (Q);
  \draw[dashed,gray!75] (Q) -- (P);
  \draw[vecB,thick] (O) -- (P);
  \node[figred,fill=white,inner sep=1pt,font=\scriptsize]
    at (1.18,1.48) {$\rho$};
  \fill (P) circle (1.3pt)
    node[above right,fill=white,inner sep=1pt,font=\scriptsize]{$P=(\rho,\varphi,\theta)$};
  \draw[vecC] (90:0.92) arc[start angle=90,end angle=49,radius=0.92];
  \node[figgreen,fill=white,inner sep=1pt,font=\scriptsize] at (0.33,0.95) {$\varphi$};
  \draw[gray!80,-{Stealth[length=2mm]}] (0:0.72)
    arc[start angle=0,end angle=15,radius=0.72];
  \node[gray!80,below,fill=white,inner sep=1pt,font=\scriptsize]
    at (0.72,0.18) {$\theta$};
  \fill (Q) circle (1pt) node[below right,font=\scriptsize]{$Q$};
\end{tikzpicture}
\caption[Spherical coordinates]{The spherical coordinates $(\rho,\varphi,\theta)$ of a point $P$: $\rho$ is the distance from the origin, $\varphi$ the angle from the positive $z$-axis, and $\theta$ the azimuthal angle, measured in the $xy$-plane, of the projection $Q$ of $P$.}
\label{fig:ch11-spherical}
\end{figure}

\section{Integration in Arbitrary Dimension}

\subsection{Multiple Integrals and Change of Variables}

The theory developed above for $n=2$ and $n=3$ extends, without any new idea, to functions of $n$ real variables; only the bookkeeping becomes heavier.

\begin{definition}[Integral over a Box in $\mathbb{R}^n$]
Consider the box
\[
B = [a_1,b_1]\times\cdots\times[a_n,b_n] \subset \mathbb{R}^n,
\]

Let $f : B \to \mathbb{R}$ be bounded. A partition of $B$, obtained by partitioning each edge, divides $B$ into subboxes $B_{\mathbf{i}}$, indexed by a multi-index $\mathbf{i} = (i_1,\dots,i_n)$, of $n$-dimensional volume $\Delta V_{\mathbf{i}}$. With
\[
M_{\mathbf{i}} = \sup_{B_{\mathbf{i}}} f,
\]

\[
m_{\mathbf{i}} = \inf_{B_{\mathbf{i}}} f,
\]
define
\[
\begin{gathered}
U(f,P) = \sum_{\mathbf{i}} M_{\mathbf{i}}\, \Delta V_{\mathbf{i}} \\
L(f,P) = \sum_{\mathbf{i}} m_{\mathbf{i}}\, \Delta V_{\mathbf{i}};
\end{gathered}
\]
$f$ is integrable on $B$ if
\[
\sup_P L(f,P) = \inf_P U(f,P),
\]
and this common value is denoted
\[
\int_B f\, dV.
\]

\end{definition}

Integrability of continuous functions, and of functions continuous except on a set of $n$-dimensional content zero, follow exactly as before, by the identical uniform continuity arguments applied coordinatewise; the definition of the integral over a general bounded region $D \subset \mathbb{R}^n$ with $\partial D$ of content zero, via zero-extension to an enclosing box, is likewise unchanged.

\begin{theorem}[Fubini's Theorem in $\mathbb{R}^n$]
Let $B_1 \subset \mathbb{R}^{n-1}$ be a box and define
\[
B = B_1 \times [c,d] \subset \mathbb{R}^{n-1}\times\mathbb{R} = \mathbb{R}^n,
\]
Let $f$ be continuous on $B$. Then
\[
\int_B f \, dV = \int_{B_1} \left( \int_c^d f(\mathbf{x}',x_n)\, dx_n \right) dV',
\]
and the inner $(n-1)$-dimensional integral over $B_1$ may itself be evaluated as a further iterated integral, one variable at a time, in any order.
\end{theorem}

\begin{proof}
By induction on $n$. The base case $n=2$ is the rectangular Fubini Theorem established above. For the inductive step, assume the theorem holds in dimension $n-1$; the identical squeeze argument used in the proof of the two-dimensional case, applied with $x_n$ playing the role of the single distinguished ``outer'' variable and $B_1$ playing the role of the ``inner'' rectangle -- with sums over subrectangles of $B_1$ replacing sums over the index $i$, and Darboux sums over $B_1$-partitions replacing the one-dimensional Darboux sums of that proof -- shows that
\[
A(x_n) := \int_{B_1} f(\mathbf{x}',x_n)\, dV'
\]
is well defined and continuous, and that
\[
\int_B f\, dV = \int_c^d A(x_n)\, dx_n.
\]

By the inductive hypothesis, applied to the continuous function $\mathbf{x}' \mapsto f(\mathbf{x}',x_n)$ on $B_1 \subset \mathbb{R}^{n-1}$ for each fixed $x_n$, the integral $A(x_n)$ may be evaluated as a further iterated integral over the $n-1$ remaining variables, one at a time, in any order; combined with the outer integration over $x_n$, this yields the theorem for dimension $n$, completing the induction.
\end{proof}

\begin{theorem}[Change of Variables Formula in $\mathbb{R}^n$]
Let $T : U \to V$ be a bijective $C^1$ map between open sets $U,V \subset \mathbb{R}^n$ with $\det J_T \neq 0$ on $U$, let $D \subset U$ with $D' = T(D) \subset V$, and let $f$ be continuous on $D'$. Then
\[
\int_{D'} f\, dV = \int_D f\big(T(\mathbf{u})\big) \, |\det J_T(\mathbf{u})| \; dV_{\mathbf{u}}.
\]

\end{theorem}

\begin{proof}[Derivation]
Exactly as in the two- and three-dimensional cases, partition $D$ into small $n$-dimensional boxes and approximate the image, under $T$, of each box by the parallelotope spanned by the $n$ columns of $J_T$, each scaled by the corresponding side length; the $n$-dimensional volume of this parallelotope is the absolute value of the determinant of the matrix formed by its $n$ spanning vectors, that is, $|\det J_T|$ times the volume of the original box, by the same multilinearity argument used above. Passing to the limit of the associated Riemann sums, exactly as before, yields the formula; we omit the details, which repeat those of the two-dimensional derivation verbatim with $n$ coordinates in place of two.
\end{proof}

\subsection{Radial Integrals and Volumes of Balls}

\begin{lemma}[Spherical Shell Formula in $\mathbb{R}^n$]
Let $g : [0,\infty) \to \mathbb{R}$ be continuous, and let $S_{n-1}$ denote the $(n-1)$-dimensional measure of the unit sphere
\[
\{\mathbf{x} \in \mathbb{R}^n : \|\mathbf{x}\|=1\}.
\]
Then
\[
\int_{\mathbb{R}^n} g(\|\mathbf{x}\|)\, dV = S_{n-1} \int_0^{\infty} g(\rho)\, \rho^{n-1}\, d\rho.
\]

\end{lemma}

\begin{proof}[Justification]
This is the direct $n$-dimensional generalization of polar coordinates ($n=2$, where the Jacobian factor $\rho^{n-1}=\rho$ matches the factor $r$ derived above) and spherical coordinates ($n=3$, where $\rho^{n-1}=\rho^2$ matches the factor $\rho^2\sin\varphi$ derived above upon integrating out the angular variable $\varphi,\theta$ over the full unit sphere, which by definition contributes precisely $S_2 = 4\pi$). In general, let $\boldsymbol{\omega}$ range over the unit sphere $S^{n-1}$ and use the radial representation
\[
\begin{aligned}
\mathbf{x} &= \rho\boldsymbol{\omega}, \\
\rho &= \|\mathbf{x}\| \geq 0.
\end{aligned}
\]

The change of variables formula applied to this generalized polar map yields a Jacobian factor depending on $\rho$ alone as $\rho^{n-1}$, exactly as in the cases $n=2,3$ verified explicitly above, times a factor depending only on the angular coordinates of $\boldsymbol\omega$; integrating $g(\rho)\rho^{n-1}$ against this angular factor over all angular coordinates, for each fixed $\rho$, yields by definition the total measure $S_{n-1}$ of the unit sphere, giving the stated formula.
\end{proof}

\begin{example}[Volume of the Unit Ball in $\mathbb{R}^n$]
The volume
\[
V_n = \operatorname{Vol}(\{\mathbf{x}\in\mathbb{R}^n : \|\mathbf{x}\|\leq 1\})
\]
of the unit ball in $\mathbb{R}^n$ can be computed in closed form using the Gaussian integral together with the shell formula above. By Fubini's Theorem in $\mathbb{R}^n$, the $n$-fold Gaussian integral separates into a product of $n$ one-dimensional Gaussian integrals, each equal to $\sqrt{\pi}$ by the earlier example:
\[
\int_{\mathbb{R}^n} e^{-\|\mathbf{x}\|^2}\, dV = \left( \int_{-\infty}^{\infty} e^{-t^2}\, dt \right)^n = \pi^{n/2}.
\]
On the other hand, applying the shell formula with $g(\rho) = e^{-\rho^2}$,
\[
\int_{\mathbb{R}^n} e^{-\|\mathbf{x}\|^2}\, dV = S_{n-1} \int_0^{\infty} e^{-\rho^2}\, \rho^{n-1}\, d\rho.
\]
Substituting $t = \rho^2$, $dt = 2\rho\, d\rho$, so that
\[
\rho^{n-1}\, d\rho = \rho^{n-2}\cdot \rho\, d\rho = t^{(n-2)/2} \cdot (1/2)\, dt,
\]

\[
\int_0^{\infty} e^{-\rho^2}\rho^{n-1}\, d\rho = \frac{1}{2}\int_0^{\infty} e^{-t}\, t^{n/2-1}\, dt = \frac{1}{2}\Gamma\!\left(\frac{n}{2}\right),
\]
by definition of the Gamma function
\[
\Gamma(s) = \int_0^\infty e^{-t}t^{s-1}\,dt.
\]
Equating the two expressions for
\[
\int_{\mathbb{R}^n} e^{-\|\mathbf{x}\|^2}\, dV
\]
and solving,
\[
\pi^{n/2} = S_{n-1}\cdot \frac{1}{2}\Gamma\!\left(\frac{n}{2}\right) \qquad\Longrightarrow\qquad S_{n-1} = \frac{2\pi^{n/2}}{\Gamma(n/2)}.
\]
Finally, applying the shell formula with $g \equiv 1$ but restricted to $\rho \in [0,1]$,
\[
\begin{aligned}
V_n &= \int_0^1 S_{n-1}\, \rho^{n-1}\, d\rho
= S_{n-1} \left[\frac{\rho^n}{n}\right]_0^1 \\
&= \frac{S_{n-1}}{n}
= \frac{2\pi^{n/2}}{n\,\Gamma(n/2)}
= \frac{\pi^{n/2}}{\Gamma\!\left(\frac{n}{2}+1\right)},
\end{aligned}
\]
using the recursive property
\[
\Gamma(s+1) = s\,\Gamma(s)
\]
of the Gamma function with $s=n/2$. As a consistency check, for $n=2$ this gives
\[
\begin{aligned}
V_2 &= \pi/\Gamma(2) \\
&= \pi/1 \\
&= \pi,
\end{aligned}
\]
the area of the unit disk; for $n=3$, using
\[
\begin{aligned}
\Gamma(5/2) &= (3/2)\cdot(1/2)\cdot\sqrt{\pi} \\
&= (3/4)\sqrt{\pi},
\end{aligned}
\]
one obtains
\[
\begin{aligned}
V_3 &= \pi^{3/2}/((3/4)\sqrt{\pi}) \\
&= (4/3)\pi,
\end{aligned}
\]
in agreement with the direct spherical-coordinates computation of the preceding section.
\end{example}

\section{Integral Theorems of Vector Calculus}
\label{sec:vector-integral-theorems}

The gradient, divergence, and curl introduced in Section~\ref{sec:vector-operators} describe local changes in scalar and vector fields. Recall that $\nabla\phi$ is a vector field, $\nabla\cdot\mathbf{F}$ is a scalar field, and $\nabla\times\mathbf{F}$ is a vector field. Positive divergence indicates a local source, while negative divergence indicates a local sink. The curl describes local circulation: its direction is the axis selected by the right-hand rule, and, for a velocity field, its magnitude is twice the local angular speed of the rotational part of the flow. The following theorems make these interpretations precise by connecting local derivatives to global integrals over boundaries.

\subsection{The Divergence Theorem}

\begin{definition}[Flux through a Closed Surface]
Let $\Omega\subset\R^3$ be a bounded domain with piecewise smooth boundary $\partial\Omega$, and let $\mathbf{n}(\bx)$ denote the outward unit normal. The \emph{flux} of a continuous vector field $\mathbf{F}$ through $\partial\Omega$ is
\begin{equation}
\Phi := \iint_{\partial\Omega}\mathbf{F}\cdot\mathbf{n}\,\dd S,
\end{equation}
where $\dd S$ is the surface area element. On a regular parametrized surface $\mathbf{r}(s,t)$, this element is
\[
\dd S=|\mathbf{r}_s\times\mathbf{r}_t|\,\dd s\,\dd t;
\]
the sign of the unit normal is chosen to agree with the prescribed orientation.
\end{definition}

\begin{theorem}[Divergence Theorem, Gauss--Ostrogradsky]
\label{thm:divergence}
Let $\Omega\subset\R^3$ be a bounded domain with piecewise smooth boundary, and let $\mathbf{F}$ be continuously differentiable on a neighborhood of $\overline\Omega$. Then
\begin{equation}
\iiint_\Omega\nabla\cdot\mathbf{F}\,\dd V
=\iint_{\partial\Omega}\mathbf{F}\cdot\mathbf{n}\,\dd S.
\end{equation}

\end{theorem}

\begin{proof}[Proof for Simple Domains]
By linearity, it is enough to establish the identity for each component. For the third component, we must show that
\begin{equation}
\iiint_\Omega\frac{\partial F_3}{\partial z}\,\dd V
=\iint_{\partial\Omega}F_3n_3\,\dd S.
\end{equation}
Assume that $\Omega$ is $z$-simple, with the description
\[
\Omega=\{(x,y,z):(x,y)\in D,\ g_1(x,y)\leq z\leq g_2(x,y)\}.
\]
Integrating first in $z$ and applying the fundamental theorem of calculus gives
\begin{align}
\iiint_\Omega\frac{\partial F_3}{\partial z}\,\dd V
&=\iint_D\left[\int_{g_1(x,y)}^{g_2(x,y)}\frac{\partial F_3}{\partial z}\,\dd z\right]\dd x\,\dd y\notag\\
&=\iint_D\big[F_3(x,y,g_2(x,y))-F_3(x,y,g_1(x,y))\big]\,\dd x\,\dd y.
\end{align}

The boundary consists of the upper graph, the lower graph, and possibly vertical sides. On the upper graph, $n_3\,\dd S=\dd x\,\dd y$; on the lower graph, $n_3\,\dd S=-\dd x\,\dd y$. These signs reflect the outward orientation. On the vertical sides, the normal is horizontal, so $n_3=0$. Therefore
\begin{align}
\iint_{\partial\Omega}F_3n_3\,\dd S
&=\iint_D F_3(x,y,g_2(x,y))\,\dd x\,\dd y\notag\\
&\quad-\iint_D F_3(x,y,g_1(x,y))\,\dd x\,\dd y,
\end{align}
which is exactly the volume integral above. The other components follow by exchanging the roles of $x,y,z$ on domains simple in the corresponding directions. Summing proves the theorem for domains simple in all three directions. For domains decomposable into finitely many such pieces, the fluxes across internal faces cancel because their outward normals are opposite. The extension to general piecewise smooth domains uses localization and approximation, which we omit.
\end{proof}

\begin{remark}[A Multidimensional Fundamental Theorem of Calculus]
The one-dimensional identity
\[
\int_a^b g'(x)\,\dd x=g(b)-g(a)
\]
expresses an integral of a derivative in terms of boundary values. The signs at $a$ and $b$ are precisely the outward normals $-1$ and $+1$. The divergence theorem has the same structure in three dimensions: the volume integral of a derivative becomes a boundary integral weighted by the outward normal. This explains why its proof closely follows the fundamental theorem of calculus.
\end{remark}

\begin{example}[The Continuity Equation]
\label{ex:continuity-equation}
Let $\rho(\bx,t)$ be the density of a conserved quantity, such as mass, charge, or probability, and let $\mathbf{J}(\bx,t)$ be its associated current. Conservation in a fixed volume $\Omega$ means that the rate of change of the total amount equals minus the net outward flux:
\begin{equation}
\frac{\dd}{\dd t}\iiint_\Omega\rho\,\dd V
=-\iint_{\partial\Omega}\mathbf{J}\cdot\mathbf{n}\,\dd S.
\end{equation}

Assuming sufficient regularity to differentiate under the integral and apply Theorem~\ref{thm:divergence}, we obtain
\begin{equation}
\iiint_\Omega\left(\frac{\partial\rho}{\partial t}+\nabla\cdot\mathbf{J}\right)\dd V=0.
\end{equation}

This identity holds for every admissible volume. A continuous integrand whose integral vanishes on every such volume must vanish identically: otherwise, continuity would give a neighborhood on which it has a strict, constant sign, contradicting the integral identity. Hence
\begin{equation}
\frac{\partial\rho}{\partial t}+\nabla\cdot\mathbf{J}=0,
\end{equation}
the \emph{continuity equation}.
\end{example}

Figure~\ref{fig:atlas-outward-flux} connects positive divergence inside a region with outward boundary flux.

\begin{figure}[!htbp]
\centering
\begin{tikzpicture}[scale=1.35,>=Stealth]
\fill[figblue!8] (0,0) circle (1.5);
\draw[figblue,thick] (0,0) circle (1.5);
\foreach \angle in {0,30,...,330} {
  \draw[->,figred,thick] (\angle:1.5) -- (\angle:2.05);
  \draw[->,black!40] (\angle:0.55) -- (\angle:0.95);
}
\node at (0,0) {$\Omega$};
\node[above right,figblue] at (45:1.5) {$\partial\Omega$};
\node[right,figred] at (0:2.05) {$n$};
\end{tikzpicture}
\caption{A planar slice through the unit ball for the three-dimensional field $F(x,y,z)=(x,y,z)$. The field has divergence $3$ and agrees with the outward unit normal on the unit sphere. The resulting flux $4\pi$ equals divergence times volume. Arrows show radial directions, not a surface parametrization.}
\label{fig:atlas-outward-flux}
\end{figure}

\subsection{Green's Theorem in the Plane}
\label{sec:green-plane}

We establish the planar theorem before Stokes' theorem so that the latter can be derived without a circular argument.

\begin{definition}[Line Integrals and Circulation]
For a piecewise smooth oriented curve $\gamma$ with parametrization $\mathbf{r}:[a,b]\to\R^3$, the line integral of a continuous vector field is
\begin{equation}
\int_\gamma\mathbf{F}\cdot\dd\mathbf{r}
=\int_a^b\mathbf{F}(\mathbf{r}(t))\cdot\mathbf{r}'(t)\,\dd t.
\end{equation}
For a closed curve, this is called the \emph{circulation} and is denoted by
\[
\oint_\gamma\mathbf{F}\cdot\dd\mathbf{r}.
\]
In the plane, with $\mathbf{F}=(P,Q)$, the integrand is $P\,\dd x+Q\,\dd y$.
\end{definition}

\begin{theorem}[Green's Theorem]
\label{thm:green-plane}
Let $D\subset\R^2$ be a bounded domain whose boundary is a piecewise smooth simple closed curve, oriented counterclockwise. If $P,Q$ are continuously differentiable on a neighborhood of $\overline D$, then
\begin{equation}
\oint_{\partial D}(P\,\dd x+Q\,\dd y)
=\iint_D\left(\frac{\partial Q}{\partial x}-\frac{\partial P}{\partial y}\right)\dd x\,\dd y.
\end{equation}

\end{theorem}

\begin{proof}[Proof for Simple Domains]
It is enough to prove the identities for $P\,\dd x$ and $Q\,\dd y$ separately. For a $y$-simple domain
\[
D=\{(x,y):a\leq x\leq b,\ h_1(x)\leq y\leq h_2(x)\},
\]
the fundamental theorem of calculus yields
\begin{equation}
\iint_D\frac{\partial P}{\partial y}\,\dd x\,\dd y
=\int_a^b\big[P(x,h_2(x))-P(x,h_1(x))\big]\,\dd x.
\end{equation}

Counterclockwise orientation traverses the lower boundary from $a$ to $b$ and the upper boundary from $b$ to $a$. Any vertical segments contribute nothing because $\dd x=0$. Consequently,
\begin{align}
\oint_{\partial D}P\,\dd x
&=\int_a^b P(x,h_1(x))\,\dd x-\int_a^b P(x,h_2(x))\,\dd x\notag\\
&=-\iint_D\frac{\partial P}{\partial y}\,\dd x\,\dd y.
\end{align}
Similarly, integration first in $x$ on an $x$-simple domain gives
\[
\oint_{\partial D}Q\,\dd y=\iint_D\frac{\partial Q}{\partial x}\,\dd x\,\dd y.
\]

Adding proves the formula when the domain admits both descriptions. Subdivision and cancellation of internal boundary integrals extend the argument to unions of such pieces; the general piecewise smooth case follows by approximation.
\end{proof}

\begin{remark}[The Divergence Form of Green's Theorem]
Applying the circulation formula to the rotated field $(-Q,P)$ gives
\begin{equation}
\oint_{\partial D}(P\,\dd y-Q\,\dd x)
=\iint_D\left(\frac{\partial P}{\partial x}+\frac{\partial Q}{\partial y}\right)\dd x\,\dd y.
\end{equation}
For counterclockwise orientation,
\[
\mathbf{n}\,\dd s=(\dd y,-\dd x),
\]
so the left-hand side is the outward flux of $(P,Q)$. This is the two-dimensional divergence theorem: in the plane, circulation and divergence are related by a rotation through $90^\circ$.
\end{remark}

\subsection{Stokes' Theorem}

Let $\Sigma\subset\R^3$ be an oriented piecewise smooth surface with boundary $\partial\Sigma$. The boundary orientation is induced by the right-hand rule: when the right thumb points along the chosen normal $\mathbf{n}$, the fingers indicate the positive direction around the surface.

\begin{theorem}[Kelvin--Stokes Theorem]
\label{thm:stokes}
If $\mathbf{F}$ is continuously differentiable on a neighborhood of $\Sigma$, then
\begin{equation}
\iint_\Sigma(\nabla\times\mathbf{F})\cdot\mathbf{n}\,\dd S
=\oint_{\partial\Sigma}\mathbf{F}\cdot\dd\mathbf{r}.
\end{equation}

\end{theorem}

\begin{proof}[Proof Sketch]
First suppose that $\Sigma$ is the graph
\[
\mathbf{r}(x,y)=(x,y,g(x,y))
\]
over a planar domain $D$, oriented upward, with $g$ sufficiently smooth. Writing $\mathbf{F}=(F_1,F_2,F_3)$ and evaluating its components on the graph gives
\[
\mathbf{F}\cdot\dd\mathbf{r}
=(F_1+F_3g_x)\,\dd x+(F_2+F_3g_y)\,\dd y.
\]
Set
\[
\begin{gathered}
P=F_1\circ\mathbf{r}+(F_3\circ\mathbf{r})g_x \\
Q=F_2\circ\mathbf{r}+(F_3\circ\mathbf{r})g_y.
\end{gathered}
\]
The chain rule and cancellation of mixed derivatives show that
\[
Q_x-P_y=(\nabla\times\mathbf{F})(\mathbf{r}(x,y))\cdot(-g_x,-g_y,1).
\]
The oriented surface element is
\[
\mathbf{n}\,\dd S=(-g_x,-g_y,1)\,\dd x\,\dd y,
\]

Green's theorem gives the desired equality. A general surface is treated by subdividing it into local graph patches. Adjacent patches traverse their common boundary in opposite directions, so internal contributions cancel, leaving only $\partial\Sigma$.
\end{proof}

\begin{remark}[A Common Boundary Principle]
Stokes' theorem again has the form ``integral of a derivative over a region equals integral of the original field over its boundary.'' Here the derivative is the curl, the region is a two-dimensional surface, and its boundary is a one-dimensional curve. Green's circulation theorem is the special case $\Sigma=D$ in the $xy$-plane, $\mathbf{n}=(0,0,1)$, and $\mathbf{F}=(P,Q,0)$, because $(\nabla\times\mathbf{F})\cdot\mathbf{n}=Q_x-P_y$. Thus the two forms of Green's theorem connect the planar versions of both Stokes' and Gauss' theorems.
\end{remark}

\newpage
\chapter{Mathematical Optimization and Convex Programming}

\section{Equality and Inequality Constraints}

\subsection{Equality-Constrained Extrema and Lagrange Multipliers}

The extremum tests developed for functions of several variables locate critical points of $f$ over its entire domain, but many optimization problems require $\mathbf{x}$ to satisfy one or more equality constraints, restricting the search to a lower-dimensional subset of the domain. The method of Lagrange multipliers reduces such constrained problems to an unconstrained search for critical points of an auxiliary function.

\begin{definition}[Lagrangian]
Given $f : \mathbb{R}^n \to \mathbb{R}$ and constraints
\[
g_1(\mathbf{x}) = \cdots = g_k(\mathbf{x}) = 0,
\]
the Lagrangian associated with the problem of extremizing $f$ subject to these constraints is the function
\[
\mathcal{L}(\mathbf{x}, \boldsymbol{\lambda}) = f(\mathbf{x}) - \sum_{i=1}^{k} \lambda_i g_i(\mathbf{x}),
\]
where $\boldsymbol{\lambda} = (\lambda_1, \dots, \lambda_k) \in \mathbb{R}^k$ are called the Lagrange multipliers.
\end{definition}

The central result gives a necessary condition: every constrained local
extremum satisfying the stated regularity assumptions must occur among
the points where the gradient of $f$ lies in the span of the gradients of
the constraint functions. A solution of the Lagrange equations need not,
by itself, be an extremum.

\begin{theorem}[Lagrange Multiplier Theorem]
Let $f, g_1, \dots, g_k$ be continuously differentiable, and suppose $\mathbf{x}^* \in \mathbb{R}^n$ is a local extremum of $f$ subject to $g_i(\mathbf{x}) = 0$ for $i = 1, \dots, k$, at which the gradients
\[
\nabla g_1(\mathbf{x}^*), \dots, \nabla g_k(\mathbf{x}^*)
\]
are linearly independent. Then there exists $\boldsymbol{\lambda}^* \in \mathbb{R}^k$ such that
\[
\nabla f(\mathbf{x}^*) = \sum_{i=1}^{k} \lambda_i^* \nabla g_i(\mathbf{x}^*),
\]
equivalently
\[
\nabla_{\mathbf{x}} \mathcal{L}(\mathbf{x}^*, \boldsymbol{\lambda}^*) = \vec{0}.
\]

\end{theorem}

\begin{proof}
Define the constraint set by
\[
M = \{\mathbf{x} \in \mathbb{R}^n : g_i(\mathbf{x}) = 0, \, i=1,\dots,k\}
\]
Since the gradients
\[
\nabla g_1(\mathbf{x}^*), \dots, \nabla g_k(\mathbf{x}^*)
\]
are linearly independent, the Implicit Function Theorem guarantees that $M$ is, near $\mathbf{x}^*$, an $(n-k)$-dimensional smooth surface whose tangent space $T_{\mathbf{x}^*}M$ coincides with $\operatorname{Ker}(J),$ where $J \in M_{k,n}$ is the matrix with rows
\[
\nabla g_1(\mathbf{x}^*), \dots, \nabla g_k(\mathbf{x}^*);
\]
equivalently, $T_{\mathbf{x}^*}M$ is the orthogonal complement of
\[
\operatorname{span}(\nabla g_1(\mathbf{x}^*), \dots, \nabla g_k(\mathbf{x}^*))
\]
in $\mathbb{R}^n$. Let $\mathbf{v} \in T_{\mathbf{x}^*}M$ be an arbitrary tangent vector, and let $\boldsymbol{\gamma}(t)$ be a smooth curve lying in $M$ with $\boldsymbol{\gamma}(0)=\mathbf{x}^*$ and $\boldsymbol{\gamma}'(0) = \mathbf{v}$, whose existence is again guaranteed by the Implicit Function Theorem. Since $\mathbf{x}^*$ is a local extremum of $f$ restricted to $M$, the composite function $t \mapsto f(\boldsymbol{\gamma}(t))$ has a local extremum at $t=0$, so its derivative there vanishes; by the Chain Rule, this derivative equals
\[
\nabla f(\mathbf{x}^*) \cdot \boldsymbol{\gamma}'(0) = \nabla f(\mathbf{x}^*) \cdot \mathbf{v}.
\]

Since $\mathbf{v} \in T_{\mathbf{x}^*}M$ was arbitrary, $\nabla f(\mathbf{x}^*)$ is orthogonal to every vector in $T_{\mathbf{x}^*}M$, that is, $\nabla f(\mathbf{x}^*)$ lies in the orthogonal complement of $T_{\mathbf{x}^*}M$, which is exactly
\[
\operatorname{span}(\nabla g_1(\mathbf{x}^*), \dots, \nabla g_k(\mathbf{x}^*)).
\]
Hence there exist scalars $\lambda_1^*, \dots, \lambda_k^*$ such that
\[
\nabla f(\mathbf{x}^*) = \sum_i \lambda_i^* \nabla g_i(\mathbf{x}^*),
\]
as claimed.
\end{proof}

\begin{example}
To maximize $f(x,y) = xy$ subject to the constraint
\[
g(x,y) = x+y-1 = 0,
\]
the Lagrange condition $\nabla f = \lambda \nabla g$ reads $(y,x) = \lambda(1,1)$, forcing $x=y$, which combined with $x+y=1$ gives
\[
\begin{aligned}
x &= y \\
&= 1/2,
\end{aligned}
\]
with $\lambda = 1/2$. Since $\nabla g = (1,1) \neq \vec{0}$ everywhere, the constraint qualification of the theorem holds at every point of the constraint line, and the constrained maximum value is
\[
f(1/2,1/2) = 1/4.
\]

\end{example}

Figure~\ref{fig:ch11-lagrange} illustrates tangency between a constraint and a level curve at a constrained extremum.

\begin{figure}[!htbp]
\centering
\begin{tikzpicture}[scale=3.0,>=Stealth]
  \draw[->] (-0.25,0) -- (1.45,0) node[right,font=\scriptsize]{$x$};
  \draw[->] (0,-0.25) -- (0,1.45) node[above,font=\scriptsize]{$y$};
  \foreach \c/\col in {0.08/figblue!35,0.16/figblue!60,0.25/figblue}{
    \draw[\col,thick,domain=0.12:1.35,samples=100,smooth,variable=\x]
      plot ({\x},{\c/\x});
  }
  \draw[figred,very thick] (-0.2,1.2) -- (1.2,-0.2);
  \node[figred,font=\scriptsize,anchor=south west] at (0.83,0.17)
    {$x+y=1$};
  \fill (0.5,0.5) circle (0.7pt);
  \node[font=\scriptsize,above left] at (0.5,0.5) {$(1/2,1/2)$};
  \draw[vecA] (0.5,0.5) -- (0.78,0.78)
    node[above right,font=\scriptsize] {$\nabla f$};
  \draw[vecC] (0.5,0.5) -- (0.9,0.9)
    node[right,font=\scriptsize] {$\nabla g$};
\end{tikzpicture}
\caption[Lagrange multipliers]{Level curves $xy=c$ (blue) and the constraint
  $x+y=1$ (red) from Example~9.1.1. At the constrained maximum
  $(1/2,1/2)$, the level curve $xy=1/4$ is tangent to the constraint and
  $\nabla f=(1/2,1/2)=\tfrac12\nabla g$.}
\label{fig:ch11-lagrange}
\end{figure}

\subsection{Karush--Kuhn--Tucker Conditions}

Many optimization problems of practical interest involve, in addition to equality constraints, inequality constraints of the form $g_i(\mathbf{x}) \leq 0$, which restrict the feasible region without forcing $\mathbf{x}$ onto a lower-dimensional surface everywhere. The Karush--Kuhn--Tucker conditions extend the Lagrange multiplier framework to this more general setting.

\begin{definition}[Active Constraint]
Given a feasible point $\mathbf{x}$, that is, a point satisfying $g_i(\mathbf{x}) \leq 0$ for every $i$, a constraint $g_i$ is called active at $\mathbf{x}$ if $g_i(\mathbf{x}) = 0$, and inactive if $g_i(\mathbf{x}) < 0$.
\end{definition}

\begin{definition}[KKT Conditions]
Consider the problem of minimizing $f(\mathbf{x})$ subject to $g_i(\mathbf{x}) \leq 0$ for $i=1,\dots,k$ and $h_j(\mathbf{x}) = 0$ for $j=1,\dots,l$. A point $\mathbf{x}^*$, together with multipliers $\boldsymbol{\mu}^* \in \mathbb{R}^k$ and $\boldsymbol{\lambda}^* \in \mathbb{R}^l$, is said to satisfy the KKT conditions if four requirements hold simultaneously: \textit{Stationarity}, requiring
\[
\nabla f(\mathbf{x}^*) + \sum_i \mu_i^* \nabla g_i(\mathbf{x}^*) + \sum_j \lambda_j^* \nabla h_j(\mathbf{x}^*) = \vec{0};
\]

\textit{primal feasibility}, requiring $g_i(\mathbf{x}^*) \leq 0$ and $h_j(\mathbf{x}^*) = 0$ for every $i,j$; \textit{dual feasibility}, requiring $\mu_i^* \geq 0$ for every $i$; and \textit{complementary slackness}, requiring $\mu_i^* g_i(\mathbf{x}^*) = 0$ for every $i$.
\end{definition}

Complementary slackness expresses a natural economic and geometric intuition: a multiplier $\mu_i^*$ is permitted to be nonzero only when the corresponding constraint $g_i(\mathbf{x}^*) \leq 0$ is active; constraints that are strictly satisfied, being inactive, contribute nothing to the stationarity condition, exactly as one would expect of a constraint that plays no binding role near $\mathbf{x}^*$.

\begin{theorem}[Necessity of the KKT Conditions]
Suppose $\mathbf{x}^*$ is a local minimum of $f$ subject to $g_i(\mathbf{x}) \leq 0$, $i=1,\dots,k$, and $h_j(\mathbf{x})=0$, $j=1,\dots,l$, and suppose the gradients of the active inequality constraints together with the gradients of all the equality constraints, evaluated at $\mathbf{x}^*$, are linearly independent. Then there exist multipliers $\boldsymbol{\mu}^* \geq \vec{0}$ and $\boldsymbol{\lambda}^*$ satisfying the KKT conditions at $\mathbf{x}^*$.
\end{theorem}

\begin{proof}
Denote the active inequality constraints at $\mathbf{x}^*$ by
\[
I = \{i : g_i(\mathbf{x}^*) = 0\}
\]

We first argue that $\nabla f(\mathbf{x}^*)$ must lie in the cone generated nonnegatively by $\{\nabla g_i(\mathbf{x}^*)\}_{i \in I}$ together with the span of $\{\nabla h_j(\mathbf{x}^*)\}_{j=1}^l$. Suppose, for the sake of contradiction, that this were not the case; a standard separation argument for convex cones then produces a direction $\mathbf{d} \in \mathbb{R}^n$ satisfying $\nabla f(\mathbf{x}^*) \cdot \mathbf{d} < 0$, $\nabla g_i(\mathbf{x}^*) \cdot \mathbf{d} < 0$ for every $i \in I$, and $\nabla h_j(\mathbf{x}^*) \cdot \mathbf{d} = 0$ for every $j$. Since the gradients of the active and equality constraints are linearly independent by hypothesis, the Implicit Function Theorem allows one to construct a feasible curve $\boldsymbol{\gamma}(t)$ with $\boldsymbol{\gamma}(0) = \mathbf{x}^*$ and $\boldsymbol{\gamma}'(0) = \mathbf{d}$, remaining feasible for the equality constraints exactly and for the active inequality constraints to first order, while inactive inequality constraints remain strictly satisfied for $t$ small by continuity. Along this curve,
\[
(d/dt)|_{t=0} f(\boldsymbol{\gamma}(t)) = \nabla f(\mathbf{x}^*) \cdot \mathbf{d} < 0,
\]
so $f$ strictly decreases for small $t>0$ while remaining feasible, contradicting the local minimality of $\mathbf{x}^*$. Hence
\[
\nabla f(\mathbf{x}^*) = -\sum_{i \in I} \mu_i^* \nabla g_i(\mathbf{x}^*) - \sum_j \lambda_j^* \nabla h_j(\mathbf{x}^*)
\]
for some $\mu_i^* \geq 0$, $i \in I$, and some $\lambda_j^* \in \mathbb{R}$. Setting $\mu_i^* = 0$ for every inactive constraint $i \notin I$, the stationarity condition
\[
\nabla f(\mathbf{x}^*) + \sum_{i=1}^k \mu_i^* \nabla g_i(\mathbf{x}^*) + \sum_j \lambda_j^*\nabla h_j(\mathbf{x}^*) = \vec{0}
\]
holds by construction, dual feasibility $\mu_i^* \geq 0$ holds by construction for $i \in I$ and trivially for $i \notin I$, primal feasibility holds since $\mathbf{x}^*$ is feasible by hypothesis, and complementary slackness holds since $\mu_i^* g_i(\mathbf{x}^*) = 0$ automatically, either because $\mu_i^*=0$ for $i \notin I$ or because $g_i(\mathbf{x}^*)=0$ for $i \in I$.
\end{proof}

\begin{example}
Minimize $f(x,y) = x^2+y^2$ subject to
\[
g(x,y) = 1 - x - y \leq 0.
\]
The stationarity condition reads
\[
(2x,2y) + \mu(-1,-1) = \vec{0},
\]
giving
\[
\begin{aligned}
x &= y \\
&= \mu/2.
\end{aligned}
\]
If the constraint is active, $x+y=1$ forces
\[
\begin{gathered}
\begin{aligned}
x &= y \\
&= 1/2
\end{aligned} \\
\begin{aligned}
\mu &= 1 \\
&\geq 0,
\end{aligned}
\end{gathered}
\]
satisfying dual feasibility; complementary slackness is automatic since the constraint is active, and primal feasibility holds since
\[
g(1/2,1/2)=0 \leq 0.
\]

Since $\nabla g = (-1,-1) \neq \vec{0}$, the constraint qualification holds, and $(1/2,1/2)$ with $\mu^*=1$ satisfies the KKT conditions; as the feasible region $\{x+y\geq 1\}$ is convex and $f$ is convex, the discussion below confirms that this point is indeed the global constrained minimizer.
\end{example}

Figure~\ref{fig:atlas-active-constraint} illustrates a constrained optimum on an active inequality boundary.

\begin{figure}[!htbp]
\centering
\begin{tikzpicture}[>=Stealth]
\begin{axis}[width=11cm,height=6.5cm,axis lines=middle,ticklabel style={font=\small},label style={font=\small},xlabel={$x$},ylabel={$y$},xmin=-0.2,xmax=1.6,ymin=-0.2,ymax=1.6,axis equal image,xtick={0,0.5,1,1.5},ytick={0,0.5,1,1.5}]
\path[fill=figgreen!12] (axis cs:-0.2,1.2) -- (axis cs:-0.2,1.6) -- (axis cs:1.6,1.6) -- (axis cs:1.6,-0.2) -- (axis cs:1.2,-0.2) -- cycle;
\addplot[figgreen,thick,domain=-0.2:1.2] {1-x};
\addplot[figblue,domain=0:90,samples=100] ({sqrt(0.5)*cos(x)},{sqrt(0.5)*sin(x)});
\addplot[figblue!40,domain=0:90,samples=100] ({1.15*cos(x)},{1.15*sin(x)});
\addplot[only marks,figred,mark=*] coordinates {(0.5,0.5)};
\draw[->,figred,thick] (axis cs:0.5,0.5) -- (axis cs:1.05,1.05) node[above right] {$\nabla f$};
\draw[->,black!70,thick] (axis cs:0.5,0.5) -- (axis cs:0.1,0.1) node[below left] {$\nabla g$};
\end{axis}
\end{tikzpicture}
\caption{Minimizing $f(x,y)=x^2+y^2$ subject to $x+y\ge1$ gives the
tangency point $(1/2,1/2)$ from Example~9.1.2. The shaded half-plane is
feasible. For $g=1-x-y$, stationarity
$\nabla f+\mu\nabla g=0$ holds with $\mu=1$.}
\label{fig:atlas-active-constraint}
\end{figure}

\section{Convex Optimization and Duality}

\subsection{Definitions and Differential Characterizations of Convexity}


The KKT conditions are, in general, only necessary for a local minimum, and a point satisfying them need not be a global minimizer. Under an
additional structural hypothesis, however, this gap disappears entirely, and every stationary point becomes automatically a global optimum. This
hypothesis is \textit{convexity}.

\begin{definition}[Convex Set and Convex Function]
A set $C \subseteq \mathbb{R}^n$ is convex if $t\mathbf{x} + (1-t)\mathbf{y} \in C$ for every $\mathbf{x}, \mathbf{y} \in C$ and every $t \in [0,1]$. A function $f : C \to \mathbb{R}$, defined on a convex set $C$, is convex if
\[
f(t\mathbf{x}+(1-t)\mathbf{y}) \leq t f(\mathbf{x}) + (1-t) f(\mathbf{y})
\]
for every $\mathbf{x}, \mathbf{y} \in C$ and $t \in [0,1]$; $f$ is strictly convex if this inequality is strict whenever $\mathbf{x} \neq \mathbf{y}$ and $t \in (0,1)$.
\end{definition}

For differentiable functions, convexity admits two further characterizations, one in terms of the gradient and one in terms of the Hessian, both of which considerably simplify its verification in practice.

\begin{theorem}[First-Order Characterization of Convexity]
A differentiable function $f$ on a convex open set $C$ is convex if and only if
\[
f(\mathbf{y}) \geq f(\mathbf{x}) + \nabla f(\mathbf{x}) \cdot (\mathbf{y}-\mathbf{x})
\]
for every $\mathbf{x}, \mathbf{y} \in C$.
\end{theorem}

\begin{proof}
Suppose $f$ is convex. For $t \in (0,1]$, convexity gives
\[
\begin{aligned}
f(\mathbf{x}+t(\mathbf{y}-\mathbf{x})) &= f((1-t)\mathbf{x}+t\mathbf{y}) \\
&\leq (1-t)f(\mathbf{x}) + tf(\mathbf{y}),
\end{aligned}
\]
so that
\[
\frac{f(\mathbf{x}+t(\mathbf{y}-\mathbf{x})) - f(\mathbf{x})}{t} \leq f(\mathbf{y}) - f(\mathbf{x}).
\]

Letting $t \to 0^+$, the left-hand side converges to the directional derivative $\nabla f(\mathbf{x}) \cdot (\mathbf{y}-\mathbf{x})$, giving
\[
\nabla f(\mathbf{x})\cdot(\mathbf{y}-\mathbf{x}) \leq f(\mathbf{y})-f(\mathbf{x}),
\]
which is the claimed inequality. Conversely, suppose the inequality holds for every $\mathbf{x},\mathbf{y} \in C$. Fix $\mathbf{x}, \mathbf{y} \in C$ and $t \in [0,1]$, and let $\mathbf{z} = t\mathbf{x}+(1-t)\mathbf{y}$. Applying the hypothesis twice, once with the pair $(\mathbf{z},\mathbf{x})$ and once with the pair $(\mathbf{z},\mathbf{y})$, gives
\[
\begin{gathered}
f(\mathbf{x}) \geq f(\mathbf{z}) + \nabla f(\mathbf{z})\cdot(\mathbf{x}-\mathbf{z}) \\
f(\mathbf{y}) \geq f(\mathbf{z}) + \nabla f(\mathbf{z})\cdot(\mathbf{y}-\mathbf{z}).
\end{gathered}
\]
Multiplying the first by $t$, the second by $1-t$, and adding, the gradient terms combine to
\[
\begin{aligned}
\nabla f(\mathbf{z}) \cdot (t\mathbf{x}+(1-t)\mathbf{y} - \mathbf{z}) &= \nabla f(\mathbf{z}) \cdot \vec{0} \\
&= 0,
\end{aligned}
\]
yielding
\[
\begin{aligned}
tf(\mathbf{x}) + (1-t)f(\mathbf{y}) &\geq f(\mathbf{z}) \\
&= f(t\mathbf{x}+(1-t)\mathbf{y}),
\end{aligned}
\]
which is exactly convexity.
\end{proof}

\begin{theorem}[Second-Order Characterization of Convexity]
A twice continuously differentiable function $f$ on a convex open set $C$ is convex if and only if $\nabla^2 f(\mathbf{x})$ is positive semidefinite for every $\mathbf{x} \in C$.
\end{theorem}

\begin{proof}
Suppose $\nabla^2 f(\mathbf{x})$ is positive semidefinite for every $\mathbf{x} \in C$. Fix $\mathbf{x}, \mathbf{y} \in C$; by Taylor's theorem with Lagrange remainder, applied to the function $t \mapsto f(\mathbf{x}+t(\mathbf{y}-\mathbf{x}))$, there exists $\xi \in (0,1)$ such that, writing $\mathbf{z} = \mathbf{x}+\xi(\mathbf{y}-\mathbf{x})$,
\[
f(\mathbf{y}) = f(\mathbf{x}) + \nabla f(\mathbf{x}) \cdot (\mathbf{y}-\mathbf{x}) + \tfrac{1}{2}(\mathbf{y}-\mathbf{x})^{\top} \nabla^2 f(\mathbf{z}) (\mathbf{y}-\mathbf{x}).
\]
Since $\nabla^2 f(\mathbf{z})$ is positive semidefinite, the final term is nonnegative, giving
\[
f(\mathbf{y}) \geq f(\mathbf{x}) + \nabla f(\mathbf{x})\cdot(\mathbf{y}-\mathbf{x}),
\]
which is the first-order characterization above; hence $f$ is convex. Conversely, suppose $f$ is convex but, for contradiction, $\nabla^2 f(\mathbf{x}_0)\mathbf{v} \cdot \mathbf{v} < 0$ for some $\mathbf{x}_0 \in C$ and some $\mathbf{v} \neq \vec{0}$. By continuity of $\nabla^2 f$, this negative quantity persists on a neighborhood of $\mathbf{x}_0$, so for $t$ sufficiently small, the same Taylor expansion as above, applied to $\mathbf{y} = \mathbf{x}_0+t\mathbf{v}$, gives
\[
f(\mathbf{x}_0+t\mathbf{v}) < f(\mathbf{x}_0) + t\nabla f(\mathbf{x}_0)\cdot \mathbf{v}
\]
for small $t \neq 0$, contradicting the first-order characterization of convexity established above, which is equivalent to convexity itself.
\end{proof}

\subsection{Global Optimality and Convex KKT Conditions}

The decisive advantage of convexity is that it eliminates the possibility of misleading local optima that are not globally optimal, a phenomenon that plagues general nonlinear optimization.

\begin{theorem}[Local Minima are Global for Convex Functions]
If $f : C \to \mathbb{R}$ is convex on a convex set $C$, then every local minimum of $f$ on $C$ is also a global minimum.
\end{theorem}

\begin{proof}
Suppose $\mathbf{x}^*$ is a local minimum but not a global one, so that some $\mathbf{y} \in C$ satisfies $f(\mathbf{y}) < f(\mathbf{x}^*)$. For $t \in (0,1]$, the point
\[
\mathbf{x}_t = (1-t)\mathbf{x}^* + t\mathbf{y}
\]
lies in $C$, by convexity of $C$, and $\mathbf{x}_t \to \mathbf{x}^*$ as $t \to 0^+$. By convexity of $f$,
\[
f(\mathbf{x}_t) \leq (1-t)f(\mathbf{x}^*) + t f(\mathbf{y}) = f(\mathbf{x}^*) + t(f(\mathbf{y})-f(\mathbf{x}^*)) < f(\mathbf{x}^*)
\]
for every $t \in (0,1]$, since $f(\mathbf{y}) < f(\mathbf{x}^*)$; this contradicts the assumption that $\mathbf{x}^*$ is a local minimum, since points $\mathbf{x}_t$ arbitrarily close to $\mathbf{x}^*$ achieve strictly smaller values of $f$.
\end{proof}

\begin{theorem}[Sufficiency of the KKT Conditions under Convexity]
Suppose $f$ and $g_1, \dots, g_k$ are convex and differentiable, and $h_1, \dots, h_l$ are affine. If $\mathbf{x}^*$, $\boldsymbol{\mu}^*$, $\boldsymbol{\lambda}^*$ satisfy the KKT conditions, then $\mathbf{x}^*$ is a global minimizer of $f$ over the feasible set.
\end{theorem}

\begin{proof}
Let $\mathbf{x}$ be any feasible point. By the first-order characterization of convexity applied to $f$ at $\mathbf{x}^*$,
\[
f(\mathbf{x}) \geq f(\mathbf{x}^*) + \nabla f(\mathbf{x}^*)\cdot(\mathbf{x}-\mathbf{x}^*).
\]
By stationarity,
\[
\nabla f(\mathbf{x}^*) = -\sum_i \mu_i^* \nabla g_i(\mathbf{x}^*) - \sum_j \lambda_j^* \nabla h_j(\mathbf{x}^*),
\]
so
\[
f(\mathbf{x}) - f(\mathbf{x}^*) \geq -\sum_i \mu_i^* \nabla g_i(\mathbf{x}^*)\cdot(\mathbf{x}-\mathbf{x}^*) - \sum_j \lambda_j^* \nabla h_j(\mathbf{x}^*)\cdot(\mathbf{x}-\mathbf{x}^*).
\]
Since each $h_j$ is affine,
\[
h_j(\mathbf{x}) - h_j(\mathbf{x}^*) = \nabla h_j(\mathbf{x}^*)\cdot(\mathbf{x}-\mathbf{x}^*)
\]
exactly, and this vanishes since both $\mathbf{x}$ and $\mathbf{x}^*$ are feasible, hence satisfy $h_j = 0$; the second sum therefore vanishes entirely. By the first-order characterization applied to each convex $g_i$ at $\mathbf{x}^*$,
\[
g_i(\mathbf{x}) \geq g_i(\mathbf{x}^*) + \nabla g_i(\mathbf{x}^*)\cdot(\mathbf{x}-\mathbf{x}^*),
\]
Rearranging gives
\[
\nabla g_i(\mathbf{x}^*)\cdot(\mathbf{x}-\mathbf{x}^*) \leq g_i(\mathbf{x}) - g_i(\mathbf{x}^*).
\]
Since $\mu_i^* \geq 0$ by dual feasibility, multiplying preserves the inequality direction, so
\[
\begin{aligned}
-\mu_i^*\nabla g_i(\mathbf{x}^*)\cdot(\mathbf{x}-\mathbf{x}^*) &\geq -\mu_i^*(g_i(\mathbf{x})-g_i(\mathbf{x}^*)) \\
&= -\mu_i^* g_i(\mathbf{x}) + \mu_i^*g_i(\mathbf{x}^*).
\end{aligned}
\]

By complementary slackness, $\mu_i^*g_i(\mathbf{x}^*)=0$, and by primal feasibility of $\mathbf{x}$, $g_i(\mathbf{x}) \leq 0$, so together with $\mu_i^* \geq 0$ one has $-\mu_i^*g_i(\mathbf{x}) \geq 0$. Combining,
\[
-\mu_i^*\nabla g_i(\mathbf{x}^*)\cdot(\mathbf{x}-\mathbf{x}^*) \geq 0
\]
for every $i$. Substituting back, $f(\mathbf{x}) - f(\mathbf{x}^*) \geq 0$, that is, $f(\mathbf{x}) \geq f(\mathbf{x}^*)$ for every feasible $\mathbf{x}$, establishing global optimality of $\mathbf{x}^*$.
\end{proof}

Figure~\ref{fig:atlas-convex-chord} shows the chord inequality that characterizes convexity.

\begin{figure}[!htbp]
\centering
\begin{tikzpicture}[>=Stealth]
\begin{axis}[width=11cm,height=5.8cm,axis lines=middle,ticklabel style={font=\small},label style={font=\small},legend style={font=\scriptsize,draw=black!20,fill=white},xlabel={$x$},ylabel={$f(x)$},xmin=-1.5,xmax=2.4,ymin=-0.3,ymax=4.7,xtick={-1,0.5,2},ytick={1,2.5,4}]
\addplot[figblue,very thick,domain=-1.3:2.1,samples=120] {x^2};
\addplot[figred,thick] coordinates {(-1,1)(2,4)};
\draw[figgreen,thick,<->] (axis cs:0.5,0.25) -- (axis cs:0.5,2.5);
\addplot[only marks,figblue,mark=*] coordinates {(-1,1)(2,4)(0.5,0.25)};
\addplot[only marks,figred,mark=*] coordinates {(0.5,2.5)};
\end{axis}
\end{tikzpicture}
\caption{For the convex function $f(x)=x^2$, the line segment joining two graph points lies above the graph. At the midpoint of $x_1=-1$ and $x_2=2$, the function value is $1/4$ whereas the mean of the endpoint values is $5/2$. The vertical gap illustrates the convexity inequality.}
\label{fig:atlas-convex-chord}
\end{figure}

\subsection{Lagrangian Duality}

This class of problems, encompassing linear programs, quadratic programs with positive semidefinite objective, and many others, forms the subject of convex optimization; its tractability in both theory and practice is closely tied to an associated dual problem, obtained by relaxing the constraints into the objective via their multipliers.

\begin{definition}[Lagrangian Dual Function]
Given the optimization problem of minimizing $f(\mathbf{x})$ subject to $g_i(\mathbf{x}) \leq 0$, $i=1,\dots,k$, and $h_j(\mathbf{x})=0$, $j=1,\dots,l$, over $\mathbf{x} \in D \subseteq \mathbb{R}^n$, the Lagrangian dual function is
\[
q(\boldsymbol{\mu}, \boldsymbol{\lambda}) = \inf_{\mathbf{x} \in D} \left( f(\mathbf{x}) + \sum_i \mu_i g_i(\mathbf{x}) + \sum_j \lambda_j h_j(\mathbf{x}) \right).
\]

\end{definition}

\begin{theorem}[Weak Duality]
For every feasible $\mathbf{x}$ and every $\boldsymbol{\mu} \geq \vec{0}$, $q(\boldsymbol{\mu},\boldsymbol{\lambda}) \leq f(\mathbf{x})$.
\end{theorem}

\begin{proof}
Since $\mathbf{x}$ is feasible, $g_i(\mathbf{x}) \leq 0$ and $h_j(\mathbf{x})=0$ for every $i,j$; since $\mu_i \geq 0$, $\mu_i g_i(\mathbf{x}) \leq 0$, so
\[
f(\mathbf{x}) + \sum_i \mu_i g_i(\mathbf{x}) + \sum_j \lambda_j h_j(\mathbf{x}) \leq f(\mathbf{x}).
\]

Taking the infimum over all $\mathbf{x} \in D$ on the left-hand side, in particular over the feasible $\mathbf{x}$ under consideration, gives
\[
q(\boldsymbol{\mu},\boldsymbol{\lambda}) \leq f(\mathbf{x}) + \sum_i \mu_i g_i(\mathbf{x}) + \sum_j \lambda_j h_j(\mathbf{x}) \leq f(\mathbf{x}).
\]

\end{proof}

Weak duality shows that the dual function always provides a lower bound on the optimal value of the original, or primal, problem, regardless of any convexity assumption; when the primal problem is convex and a mild regularity condition, known as Slater's condition, holds, this gap between the two problems closes entirely.

\begin{definition}[Slater's Condition]
The convex problem above satisfies Slater's condition if there exists a feasible point $\hat{\mathbf{x}} \in D$ with $g_i(\hat{\mathbf{x}}) < 0$ strictly for every $i$, and $h_j(\hat{\mathbf{x}}) = 0$ for every $j$.
\end{definition}

\begin{theorem}[Strong Duality]
If $f, g_1, \dots, g_k$ are convex, $h_1, \dots, h_l$ are affine, and Slater's condition holds, then the optimal value of the primal problem equals
\[
\sup_{\boldsymbol{\mu}\geq \vec{0}, \boldsymbol{\lambda}} q(\boldsymbol{\mu},\boldsymbol{\lambda}),
\]
and this supremum is attained.
\end{theorem}

The proof of strong duality proceeds by a separation argument entirely analogous in spirit to the one used to establish the necessity of the KKT conditions, applied this time to a suitably constructed convex set encoding the achievable values of the objective and constraint functions; we omit the details, referring the reader to the specialized literature on convex analysis, and record only the important consequence that, under the hypotheses above, solving the dual problem, which is always a convex problem regardless of the primal, is equivalent to solving the original primal problem itself.

\section{Calculus of Variations}
\label{sec:calculus-variations}

\subsection{Functionals and the Euler--Lagrange Equation}

\begin{definition}[Functional]
A \emph{functional} assigns a real number to a function in a specified admissible class. A typical example is
\begin{equation}
J[y]:=\int_a^b L\big(x,y(x),y'(x)\big)\,\dd x,
\end{equation}
where the prescribed function $L$, called the \emph{Lagrangian}, depends on position, function value, and derivative. A \emph{variational problem} seeks a function $y$ that minimizes $J[y]$, or more generally makes it stationary, subject to conditions such as $y(a)=y_a$ and $y(b)=y_b$.
\end{definition}

The optimization problems studied earlier in this chapter involve finitely many variables, with stationarity expressed by $\nabla f(\bv^*)=\mathbf{0}$, or by vanishing derivatives in all admissible directions when constraints are present. Here the variable is an entire function $y(\cdot)$, so the problem is infinite-dimensional. We seek the analogue of the zero-gradient condition for this setting.


\begin{theorem}[Euler--Lagrange Equation]
\label{thm:euler-lagrange}
Let $L$ be $C^2$ in its arguments, and let $y^*\in C^2([a,b])$ be a stationary point of
\[
J[y]=\int_a^b L(x,y,y')\,\dd x
\]
among functions with fixed endpoint values $y(a)=y_a$ and $y(b)=y_b$. In particular, this applies to a smooth local minimizer. Then $y^*$ satisfies
\begin{equation}
\frac{\partial L}{\partial y}
-\frac{\dd}{\dd x}\left(\frac{\partial L}{\partial y'}\right)=0,
\end{equation}
with the derivatives evaluated along $(x,y^*(x),{y^*}'(x))$.
\end{theorem}

\begin{proof}
\textit{Step 1: introduce a variation.} Choose any $\eta\in C^2([a,b])$ with

\[
\begin{aligned}
\eta(a) &= \eta(b) \\
&= 0,
\end{aligned}
\]
and set
\[
y_\varepsilon(x)=y^*(x)+\varepsilon\eta(x),\qquad g(\varepsilon)=J[y_\varepsilon].
\]

Every $y_\varepsilon$ satisfies the endpoint conditions. Stationarity of $y^*$ therefore implies $g'(0)=0$. This is precisely the requirement that the directional derivative vanish in every admissible direction.

\textit{Step 2: compute the first variation.} By definition,

\begin{equation}
g(\varepsilon)=\int_a^b L\big(x,y^*(x)+\varepsilon\eta(x),{y^*}'(x)+\varepsilon\eta'(x)\big)\,\dd x.
\end{equation}

Differentiation under the integral, justified by the smoothness assumptions, and the chain rule give
\begin{equation}
g'(\varepsilon)=\int_a^b\left[
\left.\frac{\partial L}{\partial y}\right|_{y_\varepsilon}\eta(x)
+\left.\frac{\partial L}{\partial y'}\right|_{y_\varepsilon}\eta'(x)\right]\dd x.
\end{equation}
At $\varepsilon=0$, this becomes
\begin{equation}
g'(0)=\int_a^b\left[\frac{\partial L}{\partial y}\eta(x)
+\frac{\partial L}{\partial y'}\eta'(x)\right]\dd x,
\end{equation}
where both partial derivatives are now evaluated along $y^*$.

\textit{Step 3: integrate by parts.} The term containing $\eta'$ satisfies

\begin{equation}
\int_a^b\frac{\partial L}{\partial y'}\eta'\,\dd x
=\left[\frac{\partial L}{\partial y'}\eta\right]_a^b
-\int_a^b\frac{\dd}{\dd x}\left(\frac{\partial L}{\partial y'}\right)\eta\,\dd x.
\end{equation}
The boundary term vanishes because
\[
\begin{aligned}
\eta(a) &= \eta(b) \\
&= 0.
\end{aligned}
\]
Hence
\begin{equation}
g'(0)=\int_a^b\left[\frac{\partial L}{\partial y}
-\frac{\dd}{\dd x}\left(\frac{\partial L}{\partial y'}\right)\right]\eta(x)\,\dd x=0
\end{equation}
for every admissible $\eta$.

\textit{Step 4: apply the fundamental lemma of the calculus of variations.} If a continuous function $h$ satisfies

\[
\int_a^b h\eta\,\dd x=0
\]
for every such variation, then $h\equiv0$. Indeed, if $h(x_0)>0$ at an interior point, continuity gives an interval $(x_0-\delta,x_0+\delta)\subset(a,b)$ where $h>0$. Choosing
\[
\eta(x)=
\begin{cases}
\big[\delta^2-(x-x_0)^2\big]^3,& |x-x_0|<\delta,\\
0,& |x-x_0|\geq\delta,
\end{cases}
\]
gives a nonnegative $C^2$ variation with
\[
\int_a^b h\eta\,\dd x>0,
\]
a contradiction. The negative case is analogous, and continuity gives the conclusion at the endpoints as well. Apply this lemma to
\[
h=\partial L/\partial y-\dd(\partial L/\partial y')/\dd x
\]
to obtain the Euler--Lagrange equation.
\end{proof}

\subsection{Applications: Geodesics and Stationary Action}

\begin{example}[Geodesics in the Plane: The Shortest Path Is Straight]
For the arc length of a graph,
\begin{equation}
J[y]=\int_a^b\sqrt{1+y'(x)^2}\,\dd x,
\end{equation}
the Lagrangian does not depend explicitly on $x$ or $y$. Thus the Euler--Lagrange equation reduces to
\[
\frac{\dd}{\dd x}\left(\frac{y'}{\sqrt{1+y'^2}}\right)=0.
\]

Since the map $t\mapsto t/(1+t^2)^{1/2}$ is strictly increasing, $y'$ must be constant, giving $y(x)=mx+q$. To verify that this stationary path really minimizes length, set
\[
m=(y_b-y_a)/(b-a)
\]
and use the convexity of
\[
\ell(t)=(1+t^2)^{1/2}:
\]

\[
\ell(y')\geq\ell(m)+\ell'(m)(y'-m).
\]
Integrating and using the endpoint values makes the last term vanish, so
\[
J[y]\geq(b-a)(1+m^2)^{1/2},
\]
with equality for the straight segment. The Euler--Lagrange equation alone is a stationarity condition, not a general guarantee of minimality.
\end{example}

\begin{example}[The Principle of Stationary Action]
In classical mechanics, a physical trajectory $\bx(t)$ with fixed endpoint positions makes the \emph{action}
\begin{equation}
S[\bx]=\int_{t_1}^{t_2}\mathcal{L}(\bx,\dot{\bx})\,\dd t
\end{equation}
stationary, where $\mathcal{L}=T-V$ is kinetic energy minus potential energy. Applying the Euler--Lagrange equation to each component yields
\[
\frac{\partial\mathcal{L}}{\partial\bx}
-\frac{\dd}{\dd t}\left(\frac{\partial\mathcal{L}}{\partial\dot{\bx}}\right)=\mathbf{0}.
\]
For
\[
\mathcal{L}=(1/2)m|\dot{\bx}|^2-V(\bx),
\]
the two derivatives are $-\nabla V$ and $m\dot{\bx}$, respectively. Therefore
\begin{equation}
m\ddot{\bx}=-\nabla V(\bx),
\end{equation}
which is Newton's second law. The variational viewpoint provides an alternative route to the equations of motion and extends naturally to constrained systems and field theories.
\end{example}
\newpage
\chapter{Functional, Fourier, and Complex Analysis}

\section{Hilbert Spaces and Square-Integrable Functions}

To rigorously generalize the concept of Euclidean space $\mathbb{R}^N$ to infinite dimensions, we must construct a vector space of functions that retains the geometric notions of length (norm) and angle (inner product), while ensuring that the space is \textit{complete}. Completeness guarantees that limits of Cauchy sequences exist within the space, which is an absolute prerequisite for performing calculus and infinite series expansions. The standard choice for physical applications is the space of square-integrable functions, denoted $L^2(\mathbb{R})$.

However, constructing this space requires addressing several foundational subtleties that are often glossed over. We will build it step-by-step, explicitly justifying every definition.

\subsection{Measure-Theoretic Construction of \texorpdfstring{$L^2(\mathbb{R})$}{L2(R)}}

Before defining the space, we must address the integral. In elementary calculus, the Riemann integral is used. However, the Riemann integral is inadequate for our purposes because the space of Riemann-integrable functions is \textit{not complete} under the $L^2$ norm. Specifically, a sequence of continuous (and thus Riemann-integrable) functions can converge in the $L^2$ norm to a highly discontinuous function that is not Riemann-integrable. To ensure completeness, we must use the \textit{Lebesgue integral}.
\begin{itemize}
\item \textit{Lebesgue measure ($\mu$)}: a function that assigns a ``length'' or ``volume'' to subsets of $\mathbb{R}$. Unlike the intuitive length of an interval, it is defined on a $\sigma$-algebra of sets, allowing us to measure highly fragmented sets (like the rational numbers, which have measure zero).
\item \textit{Measurable functions}: a function $v: \mathbb{R} \to \mathbb{C}$ is Lebesgue measurable if the preimage of any open set in $\mathbb{C}$ is a Lebesgue measurable set in $\mathbb{R}$. This is the exact class of functions for which the Lebesgue integral

\[
\int_\mathbb{R} |v(t)|^2 \dd t
\]
is well-defined.
\end{itemize}


Let $\mathcal{L}^2(\mathbb{R})$ be the set of all Lebesgue measurable functions $v: \mathbb{R} \to \mathbb{C}$ such that the integral of their squared modulus is finite:
\begin{equation}
\mathcal{L}^2(\mathbb{R}) = \left\{ v \in \mathcal{M}(\mathbb{R}) : \int_\mathbb{R} |v(t)|^2 \dd t < \infty \right\}
\end{equation}
where $\mathcal{M}(\mathbb{R})$ denotes the set of all complex-valued Lebesgue measurable functions.

\begin{remark}[The Problem of Positive Definiteness]
If we attempt to define a norm on $\mathcal{L}^2(\mathbb{R})$ as
\[
\|v\| = \left( \int_\mathbb{R} |v(t)|^2 \dd t \right)^{1/2},
\]
we encounter a fatal flaw. The integral of a nonnegative function is zero if and only if the function is zero \textit{almost everywhere} (a.e.), meaning the set of points where it is nonzero has Lebesgue measure zero.
Therefore, $\|v\| = 0$ does \textit{not} imply $v(t) = 0$ for all $t \in \mathbb{R}$; it only implies $v(t) = 0$ except possibly on a set of measure zero. This violates the axiom of positive definiteness for a norm, which strictly requires
\[
\|v\| = 0 \implies v = \mathbf{0}.
\]

\end{remark}

To resolve this, we must introduce an equivalence relation and form a quotient space.

\begin{definition}[Equivalence Relation and Quotient Space]
We define a relation $\sim$ on $\mathcal{L}^2(\mathbb{R})$ such that $v \sim u$ if and only if $v(t) = u(t)$ for almost every $t \in \mathbb{R}$. Formally, the set $\{t \in \mathbb{R} : v(t) \neq u(t)\}$ has Lebesgue measure zero.
This relation is reflexive, symmetric, and transitive, making it a valid equivalence relation. It partitions $\mathcal{L}^2(\mathbb{R})$ into disjoint equivalence classes.
We define the Hilbert space $L^2(\mathbb{R})$ as the quotient space:
\begin{equation}
\begin{aligned}
L^2(\mathbb{R}) &= \mathcal{L}^2(\mathbb{R}) / \\
&\sim
\end{aligned}
\end{equation}

An element of $L^2(\mathbb{R})$ is not a single function, but an equivalence class $[v]$ of functions that differ only on sets of measure zero. By convention, we abuse notation and write $v$ instead of $[v]$, treating functions that are equal almost everywhere as identical elements of the space.
\end{definition}

\subsection{The Inner Product and Norm}

With the elements of the space properly defined as equivalence classes, we can define the geometric structure.

\begin{definition}[Inner Product on $L^2(\mathbb{R})$]
For any two elements $v, u \in L^2(\mathbb{R})$, we define the inner product as:
\begin{equation}
\langle v, u \rangle = \int_\mathbb{R} v(t) \overline{u(t)} \dd t
\end{equation}

\end{definition}

\begin{remark}[Well-definedness and Finiteness of the Inner Product]
Before using this formula, we must prove the integral actually converges to a finite complex number. We apply the \textit{Cauchy--Schwarz inequality for integrals}, which states that for measurable functions $f, g$:
\begin{equation}
\int_\mathbb{R} |f(t) g(t)| \dd t \leq \left( \int_\mathbb{R} |f(t)|^2 \dd t \right)^{1/2} \left( \int_\mathbb{R} |g(t)|^2 \dd t \right)^{1/2}
\end{equation}
Setting $f(t) = v(t)$ and $g(t) = \overline{u(t)}$, we get:
\begin{equation}
\int_\mathbb{R} |v(t) \overline{u(t)}| \dd t \leq \|v\| \|u\| < \infty
\end{equation}
Since the integral of the absolute value is finite, the integral
\[
\int v \bar{u}
\]
is absolutely convergent and thus well-defined. Furthermore, because $v$ and $u$ are equivalence classes, changing the representative functions on a set of measure zero does not change the value of the integral, making the inner product well-defined on the quotient space.
\end{remark}

The inner product induces a norm via
\[
\|v\| = (\langle v, v \rangle)^{1/2}.
\]
Because we quotiented out the null sets,
\[
\|v\| = 0 \implies \int |v|^2 = 0 \implies v = 0
\]
a.e. $\implies v$ is the zero vector in $L^2(\mathbb{R})$. The triangle inequality for this norm follows directly from the \textit{Minkowski inequality}.

\subsection{Completeness and the Riesz--Fischer Theorem}

A normed space is a \textit{Banach space} if it is complete. A \textit{Hilbert space} is a Banach space whose norm is induced by an inner product.

\begin{definition}[Cauchy Sequence and Completeness]
A sequence $\{v_n\}$ in $L^2(\mathbb{R})$ is a \textit{Cauchy sequence} if for every $\varepsilon > 0$, there exists an integer $N$ such that for all $m, n > N$, $\|v_m - v_n\| < \varepsilon$.
The space is \textit{complete} if every Cauchy sequence converges to a limit $v \in L^2(\mathbb{R})$, meaning
\[
\lim_{n \to \infty} \|v_n - v\| = 0.
\]

\end{definition}

\begin{theorem}[Riesz--Fischer Theorem]
The space $L^2(\mathbb{R})$ is complete with respect to the norm induced by its inner product. Consequently, $L^2(\mathbb{R})$ is a Hilbert space.
\end{theorem}

\begin{proof}[Rigorous Sketch of Logic]
The original text's reliance on the Dominated Convergence Theorem (DCT) is a common logical leap; DCT requires a \textit{pre-existing} dominating function, which a general Cauchy sequence does not possess a priori. The rigorous proof constructs the limit function explicitly using a subsequence.

\textit{Step 1: Extract a rapidly converging subsequence.}
Since $\{v_n\}$ is Cauchy, we can find a subsequence $\{v_{n_k}\}$ such that
\[
\|v_{n_{k+1}} - v_{n_k}\| < 2^{-k}
\]
for all $k \geq 1$.

\textit{Step 2: Define a sequence of partial sums of differences.}
Define the partial sums
\[
g_K(t) = \sum_{k=1}^K |v_{n_{k+1}}(t) - v_{n_k}(t)|.
\]
This is a monotonically increasing sequence of nonnegative measurable functions.

\textit{Step 3: Apply the Monotone Convergence Theorem (MCT).}
By Minkowski's inequality, the $L^2$ norm of $g_K$ is bounded:
\begin{equation}
\|g_K\| \leq \sum_{k=1}^K \|v_{n_{k+1}} - v_{n_k}\| < \sum_{k=1}^K 2^{-k} < 1
\end{equation}
Thus,
\[
\int_\mathbb{R} |g_K(t)|^2 \dd t < 1.
\]
By the Monotone Convergence Theorem, the pointwise limit
\[
g(t) = \lim_{K \to \infty} g_K(t)
\]
exists, and
\[
\int_\mathbb{R} |g(t)|^2 \dd t \leq 1.
\]
Because the integral of $|g(t)|^2$ is finite, $g(t)$ must be finite almost everywhere.

\textit{Step 4: Construct the limit function.}
Since $g(t) < \infty$ a.e., the series
\[
\sum_{k=1}^\infty |v_{n_{k+1}}(t) - v_{n_k}(t)|
\]
converges absolutely a.e. This implies that the telescoping series:
\begin{equation}
v_{n_1}(t) + \sum_{k=1}^\infty \left( v_{n_{k+1}}(t) - v_{n_k}(t) \right)
\end{equation}
converges absolutely a.e. to some measurable function $v(t)$. By construction, $v_{n_k}(t) \to v(t)$ pointwise a.e.

\textit{Step 5: Prove $v \in L^2(\mathbb{R})$ and $v_n \to v$ in norm.}
Using Fatou's Lemma on the sequence $|v_{n_k} - v_m|^2$ (for fixed $m$), one can show that $v - v_m \in L^2(\mathbb{R})$, which implies $v \in L^2(\mathbb{R})$. Finally, because $\{v_n\}$ is Cauchy and a subsequence converges to $v$ in norm, the entire sequence must converge to $v$ in the $L^2$ norm.
\end{proof}

\subsection{Separability and Orthonormal Bases}

Completeness allows us to treat $L^2(\mathbb{R})$ as the rigorous infinite-dimensional analogue of $\mathbb{R}^N$. In $\mathbb{R}^N$, any vector is expanded in a finite orthonormal basis. In infinite dimensions, we require a countable basis.

\begin{definition}[Separability]
A metric space is \textit{separable} if it contains a countable, dense subset. $L^2(\mathbb{R})$ is separable because the set of step functions with rational endpoints and rational heights is countable and dense in $L^2(\mathbb{R})$.
\end{definition}

Because $L^2(\mathbb{R})$ is a separable Hilbert space, it admits a \textit{countable orthonormal basis} (more formally, a complete orthonormal system) $\{\phi_n\}_{n \in \mathbb{N}}$. ``Orthonormal'' means $\langle \phi_n, \phi_m \rangle = \delta_{nm}$ (the Kronecker delta). ``Complete'' means the only vector orthogonal to all $\phi_n$ is the zero vector.

Just as in $\mathbb{R}^N$, any $v \in L^2(\mathbb{R})$ has an expansion with coefficients given by its inner products with the basis functions:
\begin{equation}
\begin{aligned}
v &= \sum_{n=1}^\infty c_n \phi_n, \\
c_n &= \langle v, \phi_n \rangle.
\end{aligned}
\end{equation}

\begin{remark}[Nature of Convergence]
The infinite sum
\[
\sum_{n=1}^\infty c_n \phi_n
\]
does \textit{not} necessarily converge pointwise. The equality $v = \sum c_n \phi_n$ is understood in the \textit{$L^2$-norm sense} (mean-square convergence). Let
\[
S_N = \sum_{n=1}^N c_n \phi_n
\]
be the partial sum. The expansion means:
\begin{equation}
\lim_{N \to \infty} \left\| v - S_N \right\| = 0 \implies \lim_{N \to \infty} \int_\mathbb{R} \left| v(t) - \sum_{n=1}^N c_n \phi_n(t) \right|^2 \dd t = 0
\end{equation}

This is guaranteed by the completeness of the space and Bessel's inequality becoming an equality (Parseval's identity):
\[
\|v\|^2 = \sum_{n=1}^\infty |c_n|^2.
\]

\end{remark}

Figure~\ref{fig:atlas-orthogonal-modes} illustrates orthogonal oscillatory modes in a function space.

\begin{figure}[!htbp]
\centering
\begin{tikzpicture}[>=Stealth]
\begin{axis}[
  width=12.5cm,height=6.4cm,
  axis lines=left,
  ticklabel style={font=\small},label style={font=\small},
  xlabel={$x$},ylabel={$e_n(x)$},
  xmin=-0.03,xmax=1.03,ymin=-1.6,ymax=1.65,
  xtick={0,0.5,1},ytick={-1,0,1},
  legend style={at={(0.5,1.02)},anchor=south,legend columns=3,
    font=\scriptsize,draw=black!20,fill=white,column sep=1em}
]
\draw[black!20,densely dashed] (axis cs:0,-1.5) -- (axis cs:0,1.5);
\draw[black!20,densely dashed] (axis cs:1,-1.5) -- (axis cs:1,1.5);
\addplot[figblue,very thick,domain=0:1,samples=140] {sqrt(2)*sin(deg(pi*x))};
\addlegendentry{$n=1$}
\addplot[figred,thick,dashed,domain=0:1,samples=140] {sqrt(2)*sin(deg(2*pi*x))};
\addlegendentry{$n=2$}
\addplot[figgreen,thick,dashdotted,domain=0:1,samples=140] {sqrt(2)*sin(deg(3*pi*x))};
\addlegendentry{$n=3$}
\end{axis}
\end{tikzpicture}
\caption{The first three functions $\sqrt{2}\sin(n\pi x)$ on $[0,1]$, extended by zero outside this interval, are orthonormal in $L^2(\mathbb{R})$. Their different oscillation patterns have zero pairwise inner products. Orthonormality refers to an integral, not to perpendicular curves in the picture.}
\label{fig:atlas-orthogonal-modes}
\end{figure}

\section{Integral Operators and Spectral Theory}

\subsection{Hilbert--Schmidt Integral Operators}

In linear algebra, linear transformations on $\mathbb{R}^N$ are represented by matrices. In $L^2(\mathbb{R})$, a broad and highly useful class of linear operators is represented by integral kernels.

\begin{definition}[Linear Integral Operator]
Let $W(t, t')$ be a measurable function on $\mathbb{R} \times \mathbb{R}$. We define a linear operator
\[
W: L^2(\mathbb{R}) \to L^2(\mathbb{R})
\]
acting on a function $v \in L^2(\mathbb{R})$ via:
\begin{equation}
(Wv)(t) = \int_\mathbb{R} W(t, t') v(t') \dd t'
\end{equation}

\end{definition}

For $W$ to be a valid operator on $L^2(\mathbb{R})$, it must map $L^2$ functions to $L^2$ functions, and ideally, it should be a \textit{bounded operator}. An operator is bounded if there exists a constant $C > 0$ such that $\|Wv\| \leq C\|v\|$ for all $v \in L^2(\mathbb{R})$. Boundedness is equivalent to continuity in normed spaces.

\begin{definition}[Hilbert--Schmidt Kernel and Operator]
If the kernel $W(t, t')$ is square-integrable over the product space $\mathbb{R}^2$ with respect to the product Lebesgue measure $\dd t \dd t'$, meaning:
\begin{equation}
\|W\|_{HS}^2 := \iint_{\mathbb{R}^2} |W(t, t')|^2 \dd t \dd t' < \infty
\end{equation}
then $W$ is called a \textit{Hilbert--Schmidt kernel}.
\end{definition}

\begin{theorem}[Boundedness of Hilbert--Schmidt Operators]
Every Hilbert--Schmidt kernel defines a bounded linear operator from $L^2(\mathbb{R})$ to $L^2(\mathbb{R})$, with operator norm bounded by $\|W\|_{HS}$.
\end{theorem}

\begin{proof}[Explicit Derivation of Boundedness]
We must prove that if $v \in L^2(\mathbb{R})$, then
\[
\begin{gathered}
Wv \in L^2(\mathbb{R}) \\
\|Wv\| \leq \|W\|_{HS} \|v\|.
\end{gathered}
\]

First, we evaluate the pointwise absolute value of $(Wv)(t)$ using the \textit{Cauchy--Schwarz inequality for integrals} with respect to the variable $t'$:
\begin{align}
|(Wv)(t)|^2 &= \left| \int_\mathbb{R} W(t, t') v(t') \dd t' \right|^2 \\
    &\leq \left( \int_\mathbb{R} |W(t, t')|^2 \dd t' \right) \left( \int_\mathbb{R} |v(t')|^2 \dd t' \right) \\
    &= \|v\|^2 \int_\mathbb{R} |W(t, t')|^2 \dd t'
\end{align}
Notice that we factored out $\|v\|^2$ because it does not depend on $t'$. Next, to find the $L^2$ norm of $Wv$, we integrate both sides of this inequality with respect to $t$ over $\mathbb{R}$:
\begin{align}
\|Wv\|^2 = \int_\mathbb{R} |(Wv)(t)|^2 \dd t &\leq \int_\mathbb{R} \left( \|v\|^2 \int_\mathbb{R} |W(t, t')|^2 \dd t' \right) \dd t \\
    &= \|v\|^2 \iint_{\mathbb{R}^2} |W(t, t')|^2 \dd t' \dd t \\
    &= \|v\|^2 \|W\|_{HS}^2
\end{align}
Taking the square root of both sides yields:
\begin{equation}
\|Wv\| \leq \|W\|_{HS} \|v\|
\end{equation}

This explicitly proves that $W$ is a bounded linear operator, and the constant $C$ in the definition of boundedness is precisely the Hilbert--Schmidt norm of the kernel, $\|W\|_{HS}$. Furthermore, because $|Wv(t)|^2$ is bounded by an $L^1$ function of $t$, it guarantees that $Wv \in L^2(\mathbb{R})$.
\end{proof}

Figure~\ref{fig:atlas-integral-kernel} visualizes a symmetric square-integrable kernel on a bounded domain.

\begin{figure}[!htbp]
\centering
\begin{tikzpicture}[>=Stealth]
\begin{axis}[width=8cm,height=7cm,axis lines=middle,ticklabel style={font=\small},label style={font=\small},legend style={font=\scriptsize,draw=black!20,fill=white},view={0}{90},axis lines=box,xlabel={$x$},ylabel={$y$},xtick={0,0.5,1},ytick={0,0.5,1},colorbar,colormap={kernel}{color=(white);color=(figblue)},colorbar style={ylabel={$K(x,y)$},ytick={0,0.5,1}},point meta min=0,point meta max=1]
\addplot3[surf,shader=interp,domain=0:1,y domain=0:1,samples=25,samples y=25] {min(x,y)};
\end{axis}
\end{tikzpicture}
\caption{The color map represents $K(x,y)=\min(x,y)$ on the unit square, with zero extension outside it. Symmetry across the diagonal corresponds to $K(x,y)=K(y,x)$. Each fixed-$x$ slice supplies the weights used to integrate an input function; the bounded, compactly supported kernel is square-integrable.}
\label{fig:atlas-integral-kernel}
\end{figure}

\subsection{Eigenfunctions, Self-Adjointness, and Compactness}

In finite-dimensional linear algebra, the Spectral Theorem states that any real symmetric matrix (or complex Hermitian matrix) is diagonalizable by a unitary matrix. To generalize this to infinite-dimensional Hilbert spaces, we must carefully identify which properties of finite matrices survive the transition to infinite dimensions, and which new conditions are required to prevent pathological behavior.


Let $\calH$ be a separable Hilbert space with inner product $\langle \cdot, \cdot \rangle$ and induced norm $\|\cdot\|$. Let $W: \calH \to \calH$ be a bounded linear operator. We first define the adjoint and then the two crucial properties required for the spectral theorem.

\begin{definition}[Adjoint and Self-Adjoint Operator]
The \textit{adjoint} of $W$, denoted $W^\dagger$, is the unique bounded linear operator satisfying $\langle Wv, u \rangle = \langle v, W^\dagger u \rangle$ for all $v, u \in \calH$.
An operator $W$ is \textit{self-adjoint} (or Hermitian) if $W = W^\dagger$, meaning that the following identity holds for all $v, u \in \calH$:
\begin{equation}
\langle Wv, u \rangle = \langle v, Wu \rangle.
\end{equation}

\end{definition}

\begin{remark}[Why Self-Adjointness?]
Self-adjointness is the infinite-dimensional analogue of matrix symmetry and yields two key spectral consequences:
\textit{1. Real Eigenvalues:} If $Wv = \lambda v$ for $v \neq 0$, then

\[
\lambda \langle v, v \rangle = \langle Wv, v \rangle = \langle v, Wv \rangle = \langle v, \lambda v \rangle = \bar{\lambda} \langle v, v \rangle.
\]
Since $\langle v, v \rangle > 0$, we must have $\lambda = \bar{\lambda}$, so $\lambda \in \R$.
\textit{2. Orthogonal Eigenvectors:} If $Wv_1 = \lambda_1 v_1$ and $Wv_2 = \lambda_2 v_2$ with $\lambda_1 \neq \lambda_2$, then

\[
\lambda_1 \langle v_1, v_2 \rangle = \langle Wv_1, v_2 \rangle = \langle v_1, Wv_2 \rangle = \lambda_2 \langle v_1, v_2 \rangle.
\]
Thus,
\[
(\lambda_1 - \lambda_2)\langle v_1, v_2 \rangle = 0,
\]
which forces $\langle v_1, v_2 \rangle = 0$.
\end{remark}

\begin{definition}[Compact Operator]
An operator $W: \calH \to \calH$ is \textit{compact} if it maps every bounded set in $\calH$ to a relatively compact set (a set whose closure is compact). In a metric space like a Hilbert space, compactness is equivalent to sequential compactness. Therefore, $W$ is compact if and only if for every bounded sequence $\{v_n\} \subset \calH$, the image sequence $\{Wv_n\}$ contains a convergent subsequence in $\calH$.
\end{definition}

\begin{remark}[Why Compactness?]
In finite dimensions, the closed unit ball is compact (Heine--Borel theorem), so every bounded linear operator is automatically compact. In infinite dimensions, the closed unit ball is \textit{never} compact. Compact operators are those that ``compress'' the infinite-dimensional space enough to mimic finite-dimensional behavior. Without compactness, the spectrum of an operator can be continuous, meaning the operator might have no eigenvectors in $\calH$ at all.
\end{remark}

\subsection{The Spectral Theorem for Compact Self-Adjoint Operators}

\begin{theorem}[Spectral Theorem]
Let $W$ be a compact self-adjoint operator on a separable Hilbert space $\calH$. Then:
\begin{enumerate}
\item \textit{Real eigenvalues}: the eigenvalues of $W$ are all real.
\item \textit{Orthonormal eigenvectors}: there exists an orthonormal sequence of eigenvectors $\{\phi_n\}$ with corresponding real eigenvalues $\{\lambda_n\}$.
\item \textit{Eigenvalue convergence}: if the sequence of nonzero eigenvalues is infinite, then

\[
\lim_{n \to \infty} \lambda_n = 0.
\]

\item \textit{Spectral expansion}: any vector $v \in \calH$ can be expanded as:

\begin{equation}
Wv = \sum_{n} \lambda_n \langle v, \phi_n \rangle \phi_n
\end{equation}
where the series converges in the norm of $\calH$. (If $W$ is injective, $\{\phi_n\}$ forms a complete orthonormal basis for $\calH$.)
\end{enumerate}
\end{theorem}

\begin{proof}[Explicit Logic for $\lambda_n \to 0$]
We must explicitly prove why compactness forces the eigenvalues to converge to zero.
Suppose, for contradiction, that $\lambda_n \not\to 0$. Then there exists some $\varepsilon > 0$ and a subsequence of eigenvalues $\{\lambda_{n_k}\}$ such that $|\lambda_{n_k}| \geq \varepsilon$ for all $k$.
Consider the corresponding eigenvectors $\{\phi_{n_k}\}$. Because they are orthonormal, they form a bounded sequence ($\|\phi_{n_k}\| = 1$).
Now, examine the image sequence under $W$: $W\phi_{n_k} = \lambda_{n_k} \phi_{n_k}$.
For any $j \neq k$, the distance between two terms in this image sequence is:
\begin{align}
\|W\phi_{n_j} - W\phi_{n_k}\|^2 &= \|\lambda_{n_j} \phi_{n_j} - \lambda_{n_k} \phi_{n_k}\|^2 \notag \\
    &= \langle \lambda_{n_j} \phi_{n_j} - \lambda_{n_k} \phi_{n_k}, \lambda_{n_j} \phi_{n_j} - \lambda_{n_k} \phi_{n_k} \rangle \notag \\
    &= |\lambda_{n_j}|^2 \|\phi_{n_j}\|^2 + |\lambda_{n_k}|^2 \|\phi_{n_k}\|^2 - 2\Re(\lambda_{n_j} \bar{\lambda}_{n_k} \langle \phi_{n_j}, \phi_{n_k} \rangle)
\end{align}
Because the eigenvectors are orthonormal, $\langle \phi_{n_j}, \phi_{n_k} \rangle = 0$ and $\|\phi_{n_k}\| = 1$. Thus:
\begin{equation}
\begin{aligned}
\|W\phi_{n_j} - W\phi_{n_k}\|^2 &= |\lambda_{n_j}|^2 + |\lambda_{n_k}|^2 \\
&\geq \varepsilon^2 + \varepsilon^2 \\
&= 2\varepsilon^2
\end{aligned}
\end{equation}

This means the distance between any two distinct terms in the sequence $\{W\phi_{n_k}\}$ is at least $\sqrt{2}\varepsilon$. Consequently, this sequence cannot contain any Cauchy subsequence, and therefore cannot contain any convergent subsequence. This directly contradicts the definition of $W$ being a compact operator. Hence, our assumption was false, and we must have $\lambda_n \to 0$.
\end{proof}

\begin{remark}[Connection to Finite-Dimensional Diagonalization]
In finite dimensions, diagonalization is written as $W = E \Lambda E^{-1}$, where $E$ is the matrix of eigenvectors and $\Lambda$ is the diagonal matrix of eigenvalues. Because $W$ is self-adjoint, $E$ is unitary, so $E^{-1} = E^\dagger$.
The infinite-dimensional expansion
\[
Wv = \sum_n \lambda_n \langle v, \phi_n \rangle \phi_n
\]
is the exact analogue. If we define the operator $E$ by its action
\[
E c = \sum_n c_n \phi_n,
\]
then $W = E \Lambda E^\dagger$, where $\Lambda$ acts by multiplying the $n$-th component by $\lambda_n$.
\end{remark}

\subsection{Convolution Operators and the Continuous Spectrum}

When we drop the requirement of compactness, the spectral theorem fails in its discrete form. A prime example in physics and signal processing is the convolution operator.

\begin{definition}[Convolution Operator]
Let $K \in L^1(\R)$ be a kernel function. The convolution operator $W_K: L^2(\R) \to L^2(\R)$ is defined by:
\begin{equation}
\begin{aligned}
(W_K v)(t) &= (K * v)(t) \\
&= \int_\R K(t-t') v(t') \dd t'
\end{aligned}
\end{equation}

\end{definition}

\begin{remark}[Boundedness via Young's Inequality]
We must verify that $W_K$ actually maps $L^2$ to $L^2$. By \textit{Young's convolution inequality}, if $K \in L^1(\R)$ and $v \in L^2(\R)$, then
\[
\begin{gathered}
K * v \in L^2(\R) \\
\|K * v\|_{L^2} \leq \|K\|_{L^1} \|v\|_{L^2}.
\end{gathered}
\]

Thus, the condition $K \in L^1$ is sufficient to guarantee that $W_K$ is a bounded (continuous) operator on $L^2(\R)$.
\end{remark}

\begin{proposition}[Action on Complex Exponentials]
Consider the complex exponential functions $e_\omega(t) = e^{i\omega t}$ for $\omega \in \R$. The action of $W_K$ on $e_\omega$ yields:
\begin{equation}
(W_K e_\omega)(t) = \hat{K}(\omega) e_\omega(t)
\end{equation}
Here the Fourier transform of $K$ is defined by
\[
\hat{K}(\omega) = \int_\R K(\tau) e^{-i\omega \tau} \dd \tau
\]

\end{proposition}

\begin{proof}
We evaluate the integral explicitly using the substitution $\tau = t - t'$, which implies $t' = t - \tau$ and $\dd t' = -\dd \tau$. The limits of integration change from $t' \in (-\infty, \infty)$ to $\tau \in (\infty, -\infty)$:
\begin{align}
(W_K e_\omega)(t) &= \int_{-\infty}^\infty K(t-t') e^{i\omega t'} \dd t' \notag \\
    &= \int_{\infty}^{-\infty} K(\tau) e^{i\omega (t-\tau)} (-\dd \tau) \notag \\
    &= \int_{-\infty}^\infty K(\tau) e^{i\omega t} e^{-i\omega \tau} \dd \tau \notag \\
    &= e^{i\omega t} \left( \int_{-\infty}^\infty K(\tau) e^{-i\omega \tau} \dd \tau \right) \notag \\
    &= \hat{K}(\omega) e_\omega(t)
\end{align}
This shows that $e_\omega$ is an eigenfunction of the convolution operator with eigenvalue $\hat{K}(\omega)$.
\end{proof}

\begin{remark}[Why $e_\omega$ are not true Eigenvectors in $L^2$]
For $e_\omega$ to be a valid eigenvector in the Hilbert space $L^2(\R)$, it must belong to $L^2(\R)$. However, its norm is:
\begin{equation}
\begin{aligned}
\|e_\omega\|^2 &= \int_{-\infty}^\infty |e^{i\omega t}|^2 \dd t \\
&= \int_{-\infty}^\infty 1 \dd t \\
&= \infty
\end{aligned}
\end{equation}

Thus, $e_\omega \notin L^2(\R)$. They are ``generalized eigenvectors''. Because the eigenvalues $\hat{K}(\omega)$ form a continuous range (as $\omega$ varies continuously over $\R$), the operator $W_K$ possesses a \textit{continuous spectrum}.
\end{remark}

\begin{proposition}[Convolution Operators are Never Compact]
Unless $W_K$ is the zero operator, it is never compact.
\end{proposition}

\begin{proof}
We prove this by constructing a bounded sequence with no convergent subsequence in the image.
Let $v \in L^2(\R)$ be any nonzero function. Define a sequence of translated functions $v_n(t) = v(t - n)$ for $n \in \N$.
Because translation preserves the $L^2$ norm,
\[
\|v_n\|_{L^2} = \|v\|_{L^2},
\]
so $\{v_n\}$ is a bounded sequence.
The image of this sequence under $W_K$ is:
\begin{equation}
\begin{aligned}
(W_K v_n)(t) &= \int_\R K(t-t') v(t'-n) \dd t' \\
&= \int_\R K(t - n - \tau) v(\tau) \dd \tau \\
&= (W_K v)(t - n)
\end{aligned}
\end{equation}

Thus, $W_K v_n$ is simply the function $W_K v$ shifted by $n$.
For $j \neq k$, the squared distance between two terms in the image sequence is:
\begin{equation}
\|W_K v_j - W_K v_k\|^2 = \int_\R |(W_K v)(t-j) - (W_K v)(t-k)|^2 \dd t
\end{equation}

As $j, k \to \infty$, the overlap between the shifted functions $(W_K v)(t-j)$ and $(W_K v)(t-k)$ approaches zero because $W_K v \in L^2(\R)$ (its mass cannot be infinite). Therefore, the cross-terms in the squared integral vanish, leaving:
\begin{equation}
\lim_{j,k \to \infty} \|W_K v_j - W_K v_k\|^2 = \|W_K v\|^2 + \|W_K v\|^2 = 2\|W_K v\|^2 > 0
\end{equation}

Since the distance between terms does not approach zero, the sequence $\{W_K v_n\}$ contains no Cauchy subsequence, and therefore no convergent subsequence. This violates the definition of a compact operator.
\end{proof}

\begin{remark}[Diagonalization via the Fourier Transform]
Although $W_K$ lacks a discrete eigenbasis, it is still ``diagonalized'' by the Fourier transform. The \textit{Convolution Theorem} states that
\[
\widehat{K * v}(\omega) = \hat{K}(\omega) \hat{v}(\omega).
\]

In the Fourier domain, the operator $W_K$ acts on $\hat{v}$ by simple pointwise multiplication by the function $\hat{K}(\omega)$. This is the exact continuous analogue of a diagonal matrix multiplying a vector. Rigorously, this is formalized by the theory of \textit{Direct Integrals}, where the discrete sum
\[
\sum_n \lambda_n \langle v, \phi_n \rangle \phi_n
\]
is replaced by the continuous integral
\[
\int_\R \hat{K}(\omega) \hat{v}(\omega) e^{i\omega t} \frac{\dd \omega}{2\pi}.
\]

\end{remark}

\section{Sobolev Spaces}

\subsection{Unbounded Differential Operators}

In quantum mechanics and field theory, we frequently encounter differential operators, such as the momentum operator $P = -i(\dd/\dd t)$. However, differential operators are inherently \textit{unbounded} on $L^2(\R)$.

\begin{remark}[Why Differential Operators are Unbounded]
An operator $D$ is bounded if there exists a constant $C$ such that $\|Dv\| \leq C\|v\|$ for all $v$. Consider the sequence of functions
\[
v_n(t) = (1/\sqrt{n}) \sin(nt) \chi(t),
\]
where $\chi(t)$ is a smooth bump function with compact support.
The $L^2$ norm of $v_n$ scales as
\[
\|v_n\|_{L^2} \sim (1/\sqrt{n}) \to 0.
\]
However, the derivative is
\[
v_n'(t) = \sqrt{n} \cos(nt) \chi(t) + \dots,
\]
so its $L^2$ norm scales as
\[
\|v_n'\|_{L^2} \sim \sqrt{n} \to \infty.
\]
Because the ratio
\[
\|v_n'\| / \|v_n\| \sim n \to \infty,
\]
no finite constant $C$ can bound the derivative operator. An unbounded operator cannot be defined on the entire Hilbert space; it requires a restricted, dense domain.
\end{remark}

To rigorously define the domain of differential operators and ensure they are closed and self-adjoint, we introduce Sobolev spaces. These spaces restrict the domain to functions that possess a specific degree of ``smoothness'' or ``regularity''.

\subsection{Fourier Characterization of Sobolev Regularity}

Instead of defining smoothness via pointwise derivatives (which can be problematic for non-differentiable functions), we define it via the decay of the Fourier transform at high frequencies. Rapid decay in the frequency domain corresponds to smoothness in the time domain.

\begin{definition}[Sobolev Space $H^s(\R)$]
For any real number $s \geq 0$, the Sobolev space $H^s(\R)$ is defined as the set of all functions $v \in L^2(\R)$ whose Fourier transform $\hat{v}(\omega)$ satisfies:
\begin{equation}
\|v\|_{H^s}^2 := \int_\R (1 + \omega^2)^s |\hat{v}(\omega)|^2 \dd \omega < \infty
\end{equation}
This space is a Hilbert space equipped with the inner product:
\begin{equation}
\langle v, u \rangle_{H^s} = \int_\R (1 + \omega^2)^s \hat{v}(\omega) \overline{\hat{u}(\omega)} \dd \omega
\end{equation}

\end{definition}

\begin{remark}[Equivalence to Weak Derivatives for Integer $s$]
Why does this Fourier definition control derivatives? We use the property that the Fourier transform turns differentiation into algebraic multiplication. Specifically, if we define the Fourier transform as
\[
\hat{v}(\omega) = \frac{1}{\sqrt{2\pi}} \int_\R v(t) e^{-i\omega t} \dd t,
\]
then the Fourier transform of the derivative is
\[
\widehat{v'}(\omega) = i\omega \hat{v}(\omega).
\]
By the \textit{Plancherel Theorem}, the $L^2$ norm is preserved under the Fourier transform:
\[
\|\hat{f}\|_{L^2} = \|f\|_{L^2}.
\]
Let us evaluate the $H^1$ norm explicitly:
\begin{align}
\|v\|_{H^1}^2 &= \int_\R (1 + \omega^2) |\hat{v}(\omega)|^2 \dd \omega \notag \\
    &= \int_\R |\hat{v}(\omega)|^2 \dd \omega + \int_\R \omega^2 |\hat{v}(\omega)|^2 \dd \omega \notag \\
    &= \|\hat{v}\|_{L^2}^2 + \|i\omega \hat{v}\|_{L^2}^2 \notag \\
    &= \|v\|_{L^2}^2 + \|\widehat{v'}\|_{L^2}^2 \notag \\
    &= \|v\|_{L^2}^2 + \|v'\|_{L^2}^2
\end{align}

This explicitly proves that $v \in H^1(\R)$ if and only if both $v$ and its first derivative $v'$ are in $L^2(\R)$. For a general integer $s$, expanding $(1+\omega^2)^s$ yields a polynomial in $\omega^2$ of degree $s$, which corresponds exactly to controlling the $L^2$ norms of all weak derivatives up to order $s$.
\end{remark}

\begin{definition}[Weak Derivative]
Because functions in $H^s(\R)$ might not be classically differentiable everywhere, we rely on \textit{weak derivatives}. A locally integrable function $g$ is the weak derivative of $f$ if for every smooth test function $\varphi$ with compact support:
\begin{equation}
\int_\R g(t) \varphi(t) \dd t = - \int_\R f(t) \varphi'(t) \dd t
\end{equation}

The Fourier definition of $H^s(\R)$ automatically handles weak derivatives because it only requires the integral of $|\omega|^s |\hat{v}(\omega)|^2$ to be finite, bypassing the need for pointwise classical derivatives.
\end{definition}

\begin{remark}[Fractional Differentiability for Non-Integer $s$]
For non-integer $s$ (e.g., $s = 1/2$), the weight $(1+\omega^2)^s$ does not correspond to an integer number of classical derivatives. Instead, it defines \textit{fractional differentiability}.
The factor $(1+\omega^2)^s$ penalizes high frequencies smoothly rather than abruptly. This measures the ``fractional smoothness'' of the function, interpolating between integer derivative spaces. This is crucial in the analysis of partial differential equations and neural field models, where solutions often exhibit regularity that is not an integer number of derivatives, but still possesses a well-defined, quantifiable degree of smoothness that ensures the convergence of numerical schemes and the well-posedness of the equations.
\end{remark}

\newpage
\chapter{Differential Equations and Dynamical Systems}

\section{First-Order Differential Equations and Initial Value Problems}

\subsection{Differential Equations and Families of Solutions}

An ordinary differential equation of first order is an equation that involves an unknown function of a single variable, which we shall generally denote by $y = y(t)$, together with its first derivative. The general form of such an equation is
\[
y' = F(t,y),
\]
where $F(t,y)$ is a given expression in the two variables $t$ and $y$. Several important classes of first-order equations arise as special cases of this general form, according to the specific dependence of $F$ on $t$ and $y$. A pure-time equation is one of the form $y' = f(t)$, in which the right-hand side depends only on the independent variable. A separable equation is one of the form $y' = a(t) g(y)$, in which the right-hand side factors as a product of a function of $t$ alone and a function of $y$ alone. An autonomous equation is the special case of a separable equation in which the dependence on $t$ disappears entirely, so that $y' = g(y)$. Finally, a linear equation is one of the form $y' = a(t) y + f(t)$, in which the right-hand side is affine in the unknown $y$.

\begin{definition}
A function $y = y(t)$ is called a solution of the differential equation $y' = F(t,y)$ on the interval $I$ if the equation is satisfied identically, that is, if
\[
y'(t) = F(t, y(t))
\]
for every $t \in I$.
\end{definition}

As a simple illustration of this definition, one may verify directly that $y(t) = 10 e^{2t}$ is a solution of the equation $y' = 2y$, since differentiating this expression indeed yields
\[
\begin{aligned}
y'(t) &= 20 e^{2t} \\
&= 2 y(t).
\end{aligned}
\]

When asked to solve a differential equation, one is generally expected to determine all possible solutions of the equation, rather than a single particular solution; this immediately raises the natural question of how many solutions a first-order equation may possess.


To address this question, it is instructive to consider the simplest possible class of first-order equations, namely the pure-time equations $y' = f(t)$. For an arbitrary constant $C \in \mathbb{R}$, the solutions are given by
\[
\begin{aligned}
y(t) &= \int f(t) \, dt \\
&= g(t) + C,
\end{aligned}
\]
where $g$ denotes any fixed antiderivative of $f$; the equation therefore possesses infinitely many solutions, parametrized by the arbitrary constant $C$. This phenomenon is not peculiar to pure-time equations, but is in fact a general feature of first-order ordinary differential equations: such equations typically possess infinitely many solutions, depending on a single free parameter, which may be fixed by imposing an additional condition, as will be discussed below in connection with initial value problems.

Figure~\ref{fig:atlas-ode-direction-field} relates a differential equation to its direction field and family of solutions.

\begin{figure}[!htbp]
\centering
\begin{tikzpicture}[>=Stealth]
\begin{axis}[
  width=12.5cm,height=6.5cm,
  axis lines=left,
  ticklabel style={font=\small},label style={font=\small},
  xlabel={$t$},ylabel={$y$},
  xmin=0,xmax=2.1,ymin=0,ymax=3.3,
  xtick={0,0.5,1,1.5,2},ytick={0,1,2,3},
  legend style={at={(0.5,1.02)},anchor=south,legend columns=3,
    font=\scriptsize,draw=black!20,fill=white,column sep=1em}
]
\draw[black!25] (axis cs:-0.07500,0.00000) -- (axis cs:0.07500,0.00000);
\draw[black!25] (axis cs:-0.06708,0.46646) -- (axis cs:0.06708,0.53354);
\draw[black!25] (axis cs:-0.05303,0.94697) -- (axis cs:0.05303,1.05303);
\draw[black!25] (axis cs:-0.04160,1.43760) -- (axis cs:0.04160,1.56240);
\draw[black!25] (axis cs:-0.03354,1.93292) -- (axis cs:0.03354,2.06708);
\draw[black!25] (axis cs:-0.02785,2.43036) -- (axis cs:0.02785,2.56964);
\draw[black!25] (axis cs:-0.02372,2.92885) -- (axis cs:0.02372,3.07115);
\draw[black!25] (axis cs:0.17500,0.00000) -- (axis cs:0.32500,0.00000);
\draw[black!25] (axis cs:0.18292,0.46646) -- (axis cs:0.31708,0.53354);
\draw[black!25] (axis cs:0.19697,0.94697) -- (axis cs:0.30303,1.05303);
\draw[black!25] (axis cs:0.20840,1.43760) -- (axis cs:0.29160,1.56240);
\draw[black!25] (axis cs:0.21646,1.93292) -- (axis cs:0.28354,2.06708);
\draw[black!25] (axis cs:0.22215,2.43036) -- (axis cs:0.27785,2.56964);
\draw[black!25] (axis cs:0.22628,2.92885) -- (axis cs:0.27372,3.07115);
\draw[black!25] (axis cs:0.42500,0.00000) -- (axis cs:0.57500,0.00000);
\draw[black!25] (axis cs:0.43292,0.46646) -- (axis cs:0.56708,0.53354);
\draw[black!25] (axis cs:0.44697,0.94697) -- (axis cs:0.55303,1.05303);
\draw[black!25] (axis cs:0.45840,1.43760) -- (axis cs:0.54160,1.56240);
\draw[black!25] (axis cs:0.46646,1.93292) -- (axis cs:0.53354,2.06708);
\draw[black!25] (axis cs:0.47215,2.43036) -- (axis cs:0.52785,2.56964);
\draw[black!25] (axis cs:0.47628,2.92885) -- (axis cs:0.52372,3.07115);
\draw[black!25] (axis cs:0.67500,0.00000) -- (axis cs:0.82500,0.00000);
\draw[black!25] (axis cs:0.68292,0.46646) -- (axis cs:0.81708,0.53354);
\draw[black!25] (axis cs:0.69697,0.94697) -- (axis cs:0.80303,1.05303);
\draw[black!25] (axis cs:0.70840,1.43760) -- (axis cs:0.79160,1.56240);
\draw[black!25] (axis cs:0.71646,1.93292) -- (axis cs:0.78354,2.06708);
\draw[black!25] (axis cs:0.72215,2.43036) -- (axis cs:0.77785,2.56964);
\draw[black!25] (axis cs:0.72628,2.92885) -- (axis cs:0.77372,3.07115);
\draw[black!25] (axis cs:0.92500,0.00000) -- (axis cs:1.07500,0.00000);
\draw[black!25] (axis cs:0.93292,0.46646) -- (axis cs:1.06708,0.53354);
\draw[black!25] (axis cs:0.94697,0.94697) -- (axis cs:1.05303,1.05303);
\draw[black!25] (axis cs:0.95840,1.43760) -- (axis cs:1.04160,1.56240);
\draw[black!25] (axis cs:0.96646,1.93292) -- (axis cs:1.03354,2.06708);
\draw[black!25] (axis cs:0.97215,2.43036) -- (axis cs:1.02785,2.56964);
\draw[black!25] (axis cs:0.97628,2.92885) -- (axis cs:1.02372,3.07115);
\draw[black!25] (axis cs:1.17500,0.00000) -- (axis cs:1.32500,0.00000);
\draw[black!25] (axis cs:1.18292,0.46646) -- (axis cs:1.31708,0.53354);
\draw[black!25] (axis cs:1.19697,0.94697) -- (axis cs:1.30303,1.05303);
\draw[black!25] (axis cs:1.20840,1.43760) -- (axis cs:1.29160,1.56240);
\draw[black!25] (axis cs:1.21646,1.93292) -- (axis cs:1.28354,2.06708);
\draw[black!25] (axis cs:1.22215,2.43036) -- (axis cs:1.27785,2.56964);
\draw[black!25] (axis cs:1.22628,2.92885) -- (axis cs:1.27372,3.07115);
\draw[black!25] (axis cs:1.42500,0.00000) -- (axis cs:1.57500,0.00000);
\draw[black!25] (axis cs:1.43292,0.46646) -- (axis cs:1.56708,0.53354);
\draw[black!25] (axis cs:1.44697,0.94697) -- (axis cs:1.55303,1.05303);
\draw[black!25] (axis cs:1.45840,1.43760) -- (axis cs:1.54160,1.56240);
\draw[black!25] (axis cs:1.46646,1.93292) -- (axis cs:1.53354,2.06708);
\draw[black!25] (axis cs:1.47215,2.43036) -- (axis cs:1.52785,2.56964);
\draw[black!25] (axis cs:1.47628,2.92885) -- (axis cs:1.52372,3.07115);
\draw[black!25] (axis cs:1.67500,0.00000) -- (axis cs:1.82500,0.00000);
\draw[black!25] (axis cs:1.68292,0.46646) -- (axis cs:1.81708,0.53354);
\draw[black!25] (axis cs:1.69697,0.94697) -- (axis cs:1.80303,1.05303);
\draw[black!25] (axis cs:1.70840,1.43760) -- (axis cs:1.79160,1.56240);
\draw[black!25] (axis cs:1.71646,1.93292) -- (axis cs:1.78354,2.06708);
\draw[black!25] (axis cs:1.72215,2.43036) -- (axis cs:1.77785,2.56964);
\draw[black!25] (axis cs:1.72628,2.92885) -- (axis cs:1.77372,3.07115);
\draw[black!25] (axis cs:1.92500,0.00000) -- (axis cs:2.07500,0.00000);
\draw[black!25] (axis cs:1.93292,0.46646) -- (axis cs:2.06708,0.53354);
\draw[black!25] (axis cs:1.94697,0.94697) -- (axis cs:2.05303,1.05303);
\draw[black!25] (axis cs:1.95840,1.43760) -- (axis cs:2.04160,1.56240);
\draw[black!25] (axis cs:1.96646,1.93292) -- (axis cs:2.03354,2.06708);
\draw[black!25] (axis cs:1.97215,2.43036) -- (axis cs:2.02785,2.56964);
\draw[black!25] (axis cs:1.97628,2.92885) -- (axis cs:2.02372,3.07115);
\addplot[figblue,very thick,domain=0:2,samples=100] {0.2*exp(x)};
\addlegendentry{$C=0.2$}
\addplot[figred,thick,domain=0:2,samples=100] {0.3*exp(x)};
\addlegendentry{$C=0.3$}
\addplot[figgreen,thick,domain=0:2,samples=100] {0.4*exp(x)};
\addlegendentry{$C=0.4$}
\end{axis}
\end{tikzpicture}
\caption{Short segments indicate the slope $y$ prescribed by $y^{\prime}=y$. The curves are exact solutions $y(t)=Ce^t$ for three positive initial values. Each initial value selects one curve; the direction field itself represents all such solutions, not a single trajectory.}
\label{fig:atlas-ode-direction-field}
\end{figure}

\subsection{Techniques for Solving First-Order Equations}

Among the general theory of first-order equations, explicit techniques are available for finding a formula for the general solution, also called the general integral, of several important classes: pure-time equations $y' = f(t)$, separable equations $y' = a(t) g(y)$, including the autonomous case, and linear equations
\[
y' = a(t) y + f(t).
\]

We now illustrate the method applicable to separable equations through two worked examples, before stating the general procedure.

\begin{example}
Consider the equation $y' = \varepsilon y$, where $\varepsilon$ is a given constant. One immediately identifies the constant solution $y(t) \equiv 0$. Setting $y \neq 0$ and dividing by $y$, the equation becomes
\[
\frac{y'(t)}{y(t)} = \varepsilon \implies \int \frac{y'(t)}{y(t)} \, dt = \int \varepsilon \, dt.
\]

The integral on the left-hand side is elementary, yielding $\ln |y(t)| = \varepsilon t + c$ for an arbitrary constant $c$; setting $C = e^c$, we obtain $|y(t)| = C e^{\varepsilon t}$. This furnishes a branch of positive solutions $y(t) = C e^{\varepsilon t}$, for $C > 0$, together with a branch of negative solutions $y(t) = -C e^{\varepsilon t}$, again for $C > 0$; together with the constant solution $y \equiv 0$ found initially, these exhaust all solutions of the equation.
\end{example}

\begin{example}
Consider the equation $y' = t y^3$. As before, one immediately identifies the constant solution $y(t) \equiv 0$. Setting $y \neq 0$ and separating the variables by dividing by $y^3$, the equation becomes
\[
\frac{y'(t)}{y^3(t)} = t \implies \int \frac{y'(t)}{y^3(t)} \, dt = \int t \, dt.
\]
The integral on the left-hand side is again elementary, yielding
\[
-(1/(2 y^2(t))) = (t^2/2) + c;
\]
setting $C = -2c$, we obtain two branches of solutions,
\[
y(t) = \frac{1}{\sqrt{C - t^2}}, \qquad y(t) = -\frac{1}{\sqrt{C - t^2}},
\]
each defined on the maximal interval $(-\sqrt{C}, \sqrt{C})$, for every fixed $C > 0$.
\end{example}

These two examples illustrate the general method applicable to the broader class of separable equations, defined as follows.

\begin{definition}[Separable Equation]
An equation of the form $y' = a(t) g(y)$ is called a separable equation. The special case in which $a(t)$ is constant, so that the equation takes the form $y' = g(y)$, is called an autonomous equation.
\end{definition}

Assuming that $a$ and $g$ are continuous functions, the general solution of a separable equation $y' = a(t) g(y)$ can be obtained by means of a systematic four-step procedure. The first step consists of searching for equilibrium solutions, if any exist, namely constant solutions of the form $y(t) \equiv r \in \mathbb{R}$, which correspond precisely to the roots of the equation $g(r) = 0$. The second step consists of separating the two variables $t$ and $y$: assuming that $g(y) \neq 0$, one divides the equation by $g(y)$ and integrates both sides with respect to $t$, obtaining
\[
\int \frac{y'(t)}{g(y(t))} \, dt = \int a(t) \, dt.
\]

The third step consists of evaluating both integrals explicitly: denoting by $A(t)$ a given antiderivative of $a(t)$, and by $H(u)$ an antiderivative of $1/g(u)$, one obtains the implicit relation
\[
H(y(t)) = A(t) + C,
\]
for an arbitrary constant $C \in \mathbb{R}$. The fourth and final step consists of inverting this implicit relation, whenever $H$ is invertible, to obtain the explicit formula
\[
y(t) = H^{-1}(A(t) + C);
\]
it should be noted that restrictions on the domain of $t$, possibly depending on the value of $C$, may arise at this final step, as illustrated in the two examples above.

\subsection{Applications to Population Dynamics}


One of the most important applications of the theory of first-order differential equations lies in biology, where a wide variety of processes can be modeled by such equations, providing considerable insight into the dynamics of living organisms, that is, into how individuals and populations evolve over time. The earliest and simplest such model, proposed by Malthus in 1830, describes the growth of a population under idealized conditions.

Let $N(t)$ denote the number of individuals in a population at time $t$. The Malthus model rests on three assumptions: each individual produces offspring at a constant rate $\lambda$, so that the total birth rate at time $t$ equals $\lambda N(t)$; and individuals die at a constant per-capita rate $\mu$, so that the total loss rate through death at time $t$ equals $\mu N(t)$. Since the rate of change of the population size equals the total birth rate minus the total death rate, with $\varepsilon = \lambda - \mu$, one obtains the following equivalent forms of the differential equation:
\[
\begin{aligned}
N'(t) &= \lambda N(t) - \mu N(t), \\
\frac{dN}{dt} &= \varepsilon N.
\end{aligned}
\]

The constant $\varepsilon$ is called the biological potential, or Malthus parameter, of the population. This linear autonomous equation can be solved explicitly by the separation of variables technique developed above, yielding exponential growth or decay depending on the sign of $\varepsilon$. A comparison with experimental data on the growth of a population of yeast cells, for instance, shows reasonably good agreement between the model $dN/dt = 0.5\, N$, with initial condition $N(0) = 0.2 \times 10^6$ cells per milliliter, and the observed population counts, at least during the early stages of growth.


The Malthus model, while adequate for describing the early stages of population growth, fails to capture an important biological phenomenon known as the logistic effect, namely the intraspecific competition for limited resources: when the population density is high, one typically observes a decrease in the birth rate together with an increase in the mortality rate, effects that the constant biological potential of the Malthus model cannot represent. The logistic, or Verhulst, model, proposed in 1838, addresses this limitation by allowing the biological potential to depend on the population size itself.

Specifically, instead of assuming a constant biological potential, the logistic model assumes that the biological potential decreases linearly as the population size increases, starting from a value $\varepsilon > 0$ when $N = 0$ and reaching the value $0$ when $N = k$, where the constant $k$ is called the carrying capacity of the environment. This linear dependence is expressed by the formula
\[
\varepsilon(N) = \varepsilon \left( 1 - \frac{N}{k} \right),
\]
and substituting this expression for the biological potential into the growth equation yields the logistic equation
\[
\frac{dN}{dt} = \varepsilon \left( 1 - \frac{N}{k} \right) N.
\]

As with the Malthus model, a comparison with experimental data on the growth of a yeast population shows that the logistic model
\[
dN/dt = (0.55 - 0.0026\, N)\, N,
\]
again with initial condition $N(0) = 0.2 \times 10^6$ cells per milliliter, reproduces the observed data considerably more accurately than the simple Malthus model, particularly as the population approaches its carrying capacity. As exercises illustrating the qualitative behavior of the logistic equation, the reader may solve the equation $y' = y(2-y)$ explicitly by the method of separation of variables, and may likewise solve the closely related equation $y' = y(y-2)$, verifying in the process that some solutions of this second equation remain defined for all time, while others blow up in finite time. Figure~\ref{fig:ch7-logistic} connects the logistic trajectory with its phase-line description.

\begin{figure}[!htbp]
\centering
\begin{tikzpicture}
\begin{axis}[
  width=6.2cm,height=6.2cm, axis lines=left,
  xlabel={$t$}, ylabel={$S(t)$},
  xmin=0,xmax=10, ymin=0,ymax=5.6,
  ytick={5},yticklabels={$K$}, xtick=\empty,
  ]
  \addplot[figblue,very thick,domain=0:10,samples=150]{5/(1+9*exp(-0.9*x))};
  \draw[figred,dashed] (axis cs:0,5) -- (axis cs:10,5);
\end{axis}
\begin{axis}[
  at={(6.9cm,0)},
  width=6.2cm,height=6.2cm, axis lines=middle,
  xlabel={$S$}, ylabel={$\dot S$},
  xmin=0,xmax=5.6, ymin=-1,ymax=1.6,
  xtick={5},xticklabels={$K$}, ytick=\empty,
  ]
  \addplot[figgreen,very thick,domain=0:5.6,samples=100]{0.5*x*(1-x/5)};
  \node[circle,fill=white,draw=black,inner sep=1.3pt] at (axis cs:0,0) {};
  \node[below right,fill=white,inner sep=1pt,font=\scriptsize]
    at (axis cs:0.08,-0.12) {unstable};
  \node[circle,fill=black,inner sep=1.3pt] at (axis cs:5,0) {};
  \node[below left,fill=white,inner sep=1pt,font=\scriptsize]
    at (axis cs:4.95,-0.12) {stable};
  \draw[vecA,shorten >=1pt] (axis cs:1.0,0) -- (axis cs:2.2,0);
  \draw[vecA,shorten >=1pt] (axis cs:3.0,0) -- (axis cs:4.2,0);
  \draw[vecA,shorten >=1pt] (axis cs:5.55,0) -- (axis cs:5.15,0);
\end{axis}
\end{tikzpicture}
\caption[Logistic model: S-curve and phase line]{Left: solution $S(t)$ approaching
  the carrying capacity $K$ asymptotically. Right: phase portrait $\dot S$ vs $S$, with
  blue arrows indicating the flow direction; $S=0$ is unstable, $S=K$ is stable.}
\label{fig:ch7-logistic}
\end{figure}

\subsection{The Initial Value Problem}

In many applications arising from physical or biological problems, one is typically not interested in the entire family of solutions of a differential equation, but rather in the single, particular solution that satisfies an additional condition of the form $y(t_0) = y_0$ for prescribed values $t_0$ and $y_0$. Such a condition is called an initial condition, and the problem of finding the solution of the differential equation that satisfies a given initial condition is called an initial value problem, or Cauchy problem.

\begin{definition}[Cauchy Problem]
A Cauchy problem is a system consisting of a differential equation together with an initial condition, of the form
\[
\begin{cases}
y' = F(t,y), \\
y(t_0) = y_0,
\end{cases}
\]
where $t_0$ and $y_0$ are assigned values.
\end{definition}

A central theoretical question concerns whether a given Cauchy problem admits a solution, and if so, whether that solution is unique; a problem for which both properties hold is said to be well-posed.

\begin{definition}[Well-Posedness]
The Cauchy problem with initial datum $(t_0, y_0)$ is said to be well-posed if it admits a unique solution $y(\cdot)$ defined in a neighborhood of $t_0$.
\end{definition}

Well-posedness has an important geometric consequence: it implies that exactly one solution curve of the differential equation passes through the point $(t_0, y_0)$ in the $(t,y)$-plane, and consequently that two distinct solution curves can never intersect at a common point. In particular, if the Cauchy problem is well-posed for every admissible choice of initial condition, then a constant solution $y \equiv k$ can never be crossed by any other solution curve; the horizontal line $y = k$ therefore constitutes an impassable barrier for every other solution curve of the equation, a fact of considerable qualitative importance when sketching the family of solutions of a given equation.

For separable equations, well-posedness of the associated Cauchy problem can be guaranteed under mild regularity assumptions on the functions $a$ and $g$, though only in a local sense, as the following theorem makes precise.

\begin{theorem}[Local Well-Posedness for Separable Equations]
Let $a(t)$ be continuous on an interval $I = (\alpha, \beta)$, and let $g(y)$ be differentiable with $g'(y)$ continuous on an interval $J = (\gamma, \delta)$. Then, for every choice of $t_0 \in (\alpha,\beta)$ and every choice of initial datum $y_0 \in (\gamma,\delta)$, the Cauchy problem
\[
\begin{cases}
y' = a(t) g(y), \\
y(t_0) = y_0,
\end{cases}
\]
has a unique solution $y(t)$, defined on an interval $(t_0 - \delta_0, t_0 + \delta_0)$ for some $\delta_0 > 0$.
\end{theorem}

\begin{remark}
It is important to emphasize that the size of the interval of existence $\delta_0$ guaranteed by this theorem cannot, in general, be predicted theoretically from the hypotheses alone; the theorem asserts only the local existence and uniqueness of the solution, without providing quantitative information on how far this solution can be extended.
\end{remark}

The local nature of the well-posedness theorem stated above is not a mere technical artifact: it reflects genuine phenomena that can occur even for very simple separable equations, as the following two examples illustrate.

\begin{remark}
Existence of solutions is, in general, not global in time: there exist solutions of separable equations that are not defined on the whole interval $I$ where the coefficient $a(t)$ is continuous.
\end{remark}

\begin{example}[Blow-Up in Finite Time]
Consider the Cauchy problem $y' = y^2$, $y(0) = 1$. This equation is separable, with $a(t) \equiv 1$, continuous on all of $\mathbb{R}$, and $g(y) = y^2$. The problem admits a unique solution, given explicitly by $y(t) = 1/(1-t)$. Since
\[
\lim_{t \to 1^-} y(t) = +\infty,
\]
the solution blows up at the finite time $t=1$, and its maximal interval of existence is $I_{\max} = (-\infty, 1)$, strictly smaller than the interval $\mathbb{R}$ on which the coefficients of the equation are defined.
\end{example}

\begin{remark}
A lack of uniqueness may likewise occur whenever the hypotheses of the Local Well-Posedness Theorem, particularly the differentiability of $g$, are violated.
\end{remark}

\begin{example}[Lack of Uniqueness]
Consider the Cauchy problem
\[
y' = 3\sqrt[3]{y^2},
\]
$y(0) = 0$. This equation is separable, with
\[
g(y) = 3\sqrt[3]{y^2},
\]
which fails to be differentiable at $y = 0$. The null constant function $y(t) = 0$, for every $t \in \mathbb{R}$, is one solution of this problem; however, the function $y(t) = t^3$ provides another solution, defined globally on $\mathbb{R}$, distinct from the first.
\end{example}

In fact, the Cauchy problem of the preceding example admits infinitely many solutions, a phenomenon known as Peano's brush (\emph{Pennello di Peano}); as a further illustration, the function
\[
y(t) =
\begin{cases}
0 & t \leq 3, \\
(t-3)^3 & t > 3,
\end{cases}
\]
provides yet another solution of the same Cauchy problem, distinct from both of those exhibited above.

As further exercises illustrating the qualitative theory of separable equations, the reader may consider the equation $y' = t y^3$ and determine, for each real number $a$, the solution of the initial value problem with condition $y(1) = a$; observing that the size of the maximal interval of existence depends sensitively on the specific initial datum chosen, one may plot several of these solutions on a common graph in order to visualize this dependence. A closely related exercise concerns the equation
\[
y' = (1/2)(y^2 - 1),
\]
for which the solutions corresponding to the initial condition $y(0) = a$, again depending on the real parameter $a$, may be found explicitly and compared.

Figure~\ref{fig:atlas-finite-time-blowup} shows the finite right endpoint of a maximal solution interval.

\begin{figure}[!htbp]
\centering
\begin{tikzpicture}[>=Stealth]
\begin{axis}[
  width=12.3cm,height=6.3cm,
  axis lines=left,
  ticklabel style={font=\small},label style={font=\small},
  xlabel={$t$},ylabel={$y(t)$},
  xmin=-1,xmax=1.25,ymin=0,ymax=10.5,
  xtick={-1,0,1},ytick={1,3,5,7,9}
]
\addplot[figblue,very thick,-{Stealth[length=2.4mm]},domain=-1:0.89,samples=180] {1/(1-x)};
\draw[figred,dashed,thick] (axis cs:1,0) -- (axis cs:1,9.8);
\node[figred,above left,fill=white,inner sep=1pt,font=\small]
  at (axis cs:0.98,9.6) {$T_+=1$};
\addplot[only marks,figblue,mark=*] coordinates {(0,1)};
\node[figblue,above left,font=\scriptsize] at (axis cs:0,1) {$y(0)=1$};
\end{axis}
\end{tikzpicture}
\caption{The solution $y(t)=1/(1-t)$ with initial value $y(0)=1$ satisfies $y^{\prime}=y^2$ and is defined on $(-\infty,1)$. The vertical dashed line marks the blow-up time, not part of the solution. The equation has smooth coefficients, but this trajectory cannot be continued through that time with a finite value.}
\label{fig:atlas-finite-time-blowup}
\end{figure}

\section{Linear First-Order Equations}

\subsection{Global Well-Posedness}

We now turn to the class of linear first-order equations, already introduced above, which take the general form $y' = a(t) y + f(t)$. Several familiar equations from physics and other applications fall within this class, including the equation $P' = 2P$ governing simple population or capital growth, the equation $y' = y + 3x$, and the equation $y' - (1/t) y = \ln t$.

Unlike the case of general separable equations, for which well-posedness of the Cauchy problem could only be guaranteed locally in time, linear equations enjoy a considerably stronger property: well-posedness holds globally on the entire interval where the coefficients are defined. This distinctive feature, a direct consequence of the linear structure of the equation, is recorded in the following theorem.

\begin{theorem}[Global Well-Posedness for Linear Equations]
Let $I$ be an interval, and assume that $a(t)$ and $f(t)$ are both continuous on $I$. Then, for every choice of initial time $t_0 \in I$ and every choice of initial value $y_0 \in \mathbb{R}$, the Cauchy problem
\[
\begin{cases}
y' = a(t) y + f(t), \\
y(t_0) = y_0,
\end{cases}
\]
admits a unique solution $y(t)$, defined on the whole interval $I$.
\end{theorem}


\subsection{The Integrating Factor Method}

Every first-order linear differential equation can be solved explicitly by means of a systematic technique known as the integrating factor method. As a first illustration, one may consider the equation $y' + (1/x) y = 2$, which can be solved by multiplying both sides by the factor $x$; this particular example already suggests the general strategy, which we now describe in full.

To solve the linear differential equation written in the form $y' - a(t) y = f(t)$, one multiplies both sides by the integrating factor
\[
e^{-\int a(t) \, dt}
\]
and subsequently integrates both sides with respect to $t$.

\begin{proof}
Let $A(t)$ denote any fixed antiderivative of $a(t)$. Multiplying the equation
\[
y' - a(t) y = f(t)
\]
by the integrating factor $e^{-A(t)}$, we obtain
\begin{equation*}
e^{-A(t)} \big[ y' - a(t) y \big] = e^{-A(t)} f(t). \tag{$*$}
\end{equation*}

Consider now the product function $y(t) e^{-A(t)}$; its derivative, computed by the product rule together with the Chain Rule applied to $e^{-A(t)}$, is precisely
\[
\frac{d}{dt} \big[ y(t) e^{-A(t)} \big] = e^{-A(t)} \big[ y'(t) - a(t) y(t) \big],
\]
so that equation $(*)$ may be rewritten as
\[
\frac{d}{dt} \big[ y(t) e^{-A(t)} \big] = e^{-A(t)} f(t).
\]

By the definition of the indefinite integral, this identity gives the following relation with an arbitrary constant $c \in \mathbb{R}$:
\[
y(t) e^{-A(t)} + c = \int e^{-A(t)} f(t) \, dt.
\]
Equivalently, the solution is
\[
y(t) = C \, e^{A(t)} + e^{A(t)} \int e^{-A(t)} f(t) \, dt,
\]
where $C$ denotes an arbitrary real constant, absorbing the constant $-c$ from the previous line. This establishes the desired formula for the general solution.
\end{proof}

The proof just given yields an explicit formula for the general solution of any linear first-order equation with continuous coefficients, which we now record for reference.

\begin{theorem}[General Integral of the Linear Equation]
Let $a(t)$ and $f(t)$ be continuous functions on an interval $I$, and let $A(t)$ denote an antiderivative of $a$ on $I$. For an arbitrary constant $C \in \mathbb{R}$, the solutions of $y' = a(t) y + f(t)$ on $I$ and a particular solution are given by
\[
\begin{aligned}
y(t) &= C \, e^{A(t)} + y_P(t), \\
y_P(t) &= e^{A(t)} \int e^{-A(t)} f(t) \, dt.
\end{aligned}
\]

\end{theorem}

\begin{remark}
The general solution formula obtained above can be written suggestively as
\[
y(t) = y_{\text{homo}}(t) + y_P(t).
\]
The homogeneous term is
\[
y_{\text{homo}}(t) = C e^{A(t)}.
\]

This is the general solution of the associated homogeneous equation $y' = a(t) y$, obtained by setting $f \equiv 0$, and $y_P(t)$ denotes a single particular solution of the full nonhomogeneous equation. This decomposition anticipates a structural principle that will play an even more prominent role in the theory of second-order linear equations developed in the following section.
\end{remark}

\section{Second-Order Linear Differential Equations}

\subsection{Second-Order Linear Equations and Initial Data}

We now turn to the study of second-order linear differential equations, which take the general form
\[
a(t) y'' + b(t) y' + c(t) y = f(t), \qquad t \in I,
\]
where $a$, $b$, and $c$ are continuous functions on the interval $I$, called the coefficient functions, subject to the nondegeneracy condition $a(t) \neq 0$ for every $t \in I$, and where $f$ is a continuous function on $I$, called the forcing term. Equations of this type arise naturally in a wide variety of physical contexts; a particularly simple example is furnished by the equation of motion of a falling body, $m y'' = -mg$, whose general solution is readily found by direct integration to be
\[
y(t) = -(1/2) g t^2 + v_0 t + y_0,
\]
where $v_0$ and $y_0$ denote the initial velocity and initial position of the body, respectively. As for first-order linear equations, the associated initial value problem for a second-order linear equation is well-posed globally on the interval of continuity of the coefficients, subject now to the specification of two initial conditions, since the equation involves a second derivative.

\begin{theorem}[Cauchy Theorem for Second-Order Linear Equations]
Assume that $a$, $b$, $c$, and $f$ are continuous functions on an interval $I$, with $a(t) \neq 0$ throughout $I$. Then, for any given $t_0 \in I$, $y_0 \in \mathbb{R}$, and $v_0 \in \mathbb{R}$, the Cauchy problem
\[
\begin{cases}
a(t) y'' + b(t) y' + c(t) y = f(t), \\
y(t_0) = y_0, \\
y'(t_0) = v_0,
\end{cases}
\]
has a unique solution $y : I \to \mathbb{R}$.
\end{theorem}

\subsection{The Homogeneous Equation}

We first study the homogeneous equation associated with the general second-order linear equation, namely
\begin{equation*}
a(t) y'' + b(t) y' + c(t) y = 0, \qquad t \in I, \tag{$\star$}
\end{equation*}
where, as above, $a$, $b$, $c$ are continuous on $I$ with $a(t) \neq 0$ throughout. The linear structure of this equation is reflected in the following superposition principle, which shows that the set of solutions is closed under linear combination.

\begin{theorem}[Superposition Principle for the Homogeneous Equation]
If $y_1(t)$ and $y_2(t)$ are both solutions of the linear homogeneous equation $(\star)$ on the interval $I$, and $c_1, c_2$ are arbitrary constants, then the function
\[
y(t) = c_1 y_1(t) + c_2 y_2(t)
\]
is also a solution of $(\star)$ on $I$.
\end{theorem}

This superposition principle, together with a suitable notion of linear independence, allows one to describe the entire family of solutions of the homogeneous equation in terms of just two particular solutions.

\begin{theorem}[General Integral of the Homogeneous Equation]
If $y_1(t)$ and $y_2(t)$ are linearly independent solutions of the linear homogeneous equation $(\star)$ on the interval $I$, meaning that the ratio $y_1/y_2$ is not constant on $I$, then the general solution of $(\star)$ is given by
\[
y(t) = c_1 y_1(t) + c_2 y_2(t), \qquad t \in I,
\]
where $c_1$ and $c_2$ are arbitrary constants.
\end{theorem}

This result shows that, once two particular linearly independent solutions of the homogeneous equation are known, every solution of the equation is thereby determined; equivalently, one says that the set of all solutions of the second-order linear homogeneous equation $(\star)$ forms a vector space of dimension two.

In general, discovering two particular solutions of a second-order linear equation with variable coefficients can be a genuinely difficult problem. It is, however, always possible to do so explicitly when the coefficients are constant, that is, for an equation of the form $a y'' + b y' + c y = 0$ with $a, b, c \in \mathbb{R}$ and $a \neq 0$; the question of how to find two linearly independent solutions of such a constant-coefficient homogeneous equation, according to the sign of the discriminant of the associated characteristic equation, is addressed through dedicated techniques developed elsewhere in the course.

\subsection{The Nonhomogeneous Equation}

We now turn to the general, nonhomogeneous second-order linear equation,
\begin{equation*}
a(t) y'' + b(t) y' + c(t) y = f(t), \qquad t \in I, \tag{$\star\star$}
\end{equation*}
where $a$, $b$, $c$ are continuous on $I$, with $a(t) \neq 0$ throughout, and $f$ is a given continuous external forcing term. As in the homogeneous case, the linear structure of the equation gives rise to a corresponding superposition principle, which in this setting relates solutions of the equation with different forcing terms.

\begin{theorem}[Superposition Principle for the Nonhomogeneous Equation]
If $y_1(t)$ and $y_2(t)$ are solutions of the linear equation $(\star\star)$ with external forcing terms $f_1$ and $f_2$ respectively, and $c_1, c_2$ are arbitrary constants, then the function
\[
y(t) = c_1 y_1(t) + c_2 y_2(t)
\]
is a solution of $(\star\star)$ with forcing term
\[
f(t) = c_1 f_1(t) + c_2 f_2(t).
\]

\end{theorem}

The general solution of the nonhomogeneous equation can be described in terms of the solutions of the associated homogeneous equation, together with any single particular solution of the nonhomogeneous equation itself, according to the following theorem.

\begin{theorem}[General Integral of the Nonhomogeneous Equation]
The general solution of the nonhomogeneous differential equation $(\star\star)$ can be written as $y = y_{\text{homo}} + y_P$, where $y_P$ is any particular solution of $(\star\star)$, and $y_{\text{homo}}$ denotes the general solution of the associated complementary homogeneous equation $(\star)$.
\end{theorem}

This theorem shows that, once two linearly independent solutions of the associated homogeneous equation $(\star)$ are known, together with a single particular solution of the nonhomogeneous equation $(\star\star)$, every solution of $(\star\star)$ is thereby determined, exactly as in the analogous situation encountered for first-order linear equations.

As in the homogeneous case, finding a particular solution of a nonhomogeneous second-order linear equation with variable coefficients can be genuinely difficult in general. When the coefficient functions are constant, however, that is, for an equation of the form
\[
a y'' + b y' + c y = f(t)
\]
with $a, b, c \in \mathbb{R}$ and $a \neq 0$, a systematic technique known as the Method of Undetermined Coefficients is available, which allows one to guess the form of a particular solution directly from the structure of the forcing term $f(t)$, whether polynomial, exponential, or trigonometric in nature; the detailed development of this method, together with the classification of the homogeneous solutions according to the sign of the discriminant $\Delta$ of the characteristic equation, is left to the dedicated material provided for further study.

\section{Existence, Uniqueness, and Maximal Solutions}

\subsection{Initial Data and Local Lipschitz Regularity}

We consider an autonomous system of ODEs in $\R^N$:
\begin{equation}
\dot{\bv} = f(\bv), \quad \bv(0) = \bv_0
\end{equation}
where $f: \R^N \to \R^N$ is a vector field. A \textit{solution} to this initial value problem is a differentiable function $\bv: I \to \R^N$, defined on some interval $I$ containing $0$, such that
\[
(\dd/\dd t)\bv(t) = f(\bv(t))
\]
for all $t \in I$, and $\bv(0) = \bv_0$. To guarantee that solutions exist and are unique, we must impose regularity conditions on $f$. The standard condition is local Lipschitz continuity.

\begin{definition}[Locally Lipschitz Continuous]
A function $f: \R^N \to \R^N$ is \textit{locally Lipschitz continuous} if for every point $\bx \in \R^N$, there exists a neighborhood $U$ of $\bx$ and a constant $L_U > 0$ such that for all $\bx, \by \in U$:
\begin{equation}
\|f(\bx) - f(\by)\| \leq L_U \|\bx - \by\|
\end{equation}
where $\|\cdot\|$ denotes the standard Euclidean norm.
\end{definition}

\begin{remark}[Why the Lipschitz Condition?]
Continuity of $f$ alone (via Peano's existence theorem) guarantees that a solution exists, but it does not guarantee that the solution is \textit{unique}. Without uniqueness, trajectories could branch, making deterministic prediction impossible. The Lipschitz condition bounds the rate at which $f$ can change, preventing the vector field from becoming too steep and forcing trajectories to merge or branch.
\end{remark}

\subsection{Gronwall's Integral Inequality}


To prove uniqueness over the entire interval $[0, T]$, we rely on a fundamental integral inequality.

\begin{lemma}[Gronwall's Inequality]
Let $u$ and $\beta$ be continuous functions with the following domain and codomain:
\[
u, \beta: [0, T] \to [0, \infty).
\]
Suppose there exists a constant $C \geq 0$ such that the following bound holds for every $t \in [0, T]$:
\begin{equation}
u(t) \leq C + \int_0^t \beta(s) u(s) \dd s.
\end{equation}
For every $t$ in the same interval, the resulting estimate is
\begin{equation}
u(t) \leq C \exp\left( \int_0^t \beta(s) \dd s\right).
\end{equation}

\end{lemma}

\begin{proof}
We construct an explicit upper bound using the method of integrating factors.
Define the function $R(t)$ as the right-hand side of the inequality:
\begin{equation}
R(t) = C + \int_0^t \beta(s) u(s) \dd s
\end{equation}

By the given hypothesis, we immediately have $u(t) \leq R(t)$ for all $t$.
Notice that $R(0) = C$. By the Fundamental Theorem of Calculus, since the integrand $\beta(s)u(s)$ is continuous, $R(t)$ is differentiable and its derivative is:
\begin{equation}
R'(t) = \beta(t) u(t)
\end{equation}
Since $u(t) \leq R(t)$ and $\beta(t) \geq 0$, we can substitute to obtain a differential inequality for $R(t)$:
\begin{equation}
R'(t) \leq \beta(t) R(t) \implies R'(t) - \beta(t) R(t) \leq 0
\end{equation}
To solve this differential inequality, we multiply both sides by the integrating factor
\[
e^{-\int_0^t \beta(s) \dd s},
\]
which is strictly positive. Using the product rule for differentiation, we recognize the left side as the exact derivative of a product:
\begin{equation}
\begin{aligned}
\frac{\dd}{\dd t} \left[ R(t) e^{-\int_0^t \beta(s) \dd s} \right] &= R'(t) e^{-\int_0^t \beta} + R(t) \left( -\beta(t) e^{-\int_0^t \beta} \right) \\
&= e^{-\int_0^t \beta} \Big( R'(t) - \beta(t) R(t) \Big)
\end{aligned}
\end{equation}
The signs of the two factors are
\[
\begin{gathered}
e^{-\int_0^t \beta} > 0 \\
R'(t) - \beta(t) R(t) \leq 0,
\end{gathered}
\]
Hence the derivative of the product is non-positive:
\begin{equation}
\frac{\dd}{\dd t} \left[ R(t) e^{-\int_0^t \beta(s) \dd s} \right] \leq 0
\end{equation}
This implies that the function
\[
g(t) = R(t) e^{-\int_0^t \beta(s) \dd s}
\]
is monotonically non-increasing. Therefore, for any $t \geq 0$, $g(t) \leq g(0)$. Evaluating at $t=0$:
\begin{equation}
\begin{aligned}
g(0) &= R(0) e^0 \\
&= C \cdot 1 \\
&= C
\end{aligned}
\end{equation}
Thus,
\[
R(t) e^{-\int_0^t \beta(s) \dd s} \leq C.
\]
Multiplying both sides by the positive quantity
\[
e^{\int_0^t \beta(s) \dd s}
\]
yields:
\begin{equation}
R(t) \leq C \exp\left( \int_0^t \beta(s) \dd s \right)
\end{equation}
Finally, since $u(t) \leq R(t)$, we conclude that
\[
u(t) \leq C \exp\left( \int_0^t \beta(s) \dd s \right).
\]

\end{proof}

\subsection{Application of Gronwall's Inequality to Global Uniqueness}

\begin{theorem}[Global Uniqueness]
If $\bv_1(t)$ and $\bv_2(t)$ are two solutions to $\dot{\bv} = f(\bv)$ on the interval $[0, T]$ satisfying the same initial condition
\[
\begin{aligned}
\bv_1(0) &= \bv_2(0) \\
&= \bv_0,
\end{aligned}
\]
and if $f$ is globally Lipschitz continuous with constant $L$ (i.e.,
\[
\|f(\bx) - f(\by)\| \leq L\|\bx - \by\|
\]
for all $\bx, \by \in \R^N$), then $\bv_1(t) = \bv_2(t)$ for all $t \in [0, T]$.
\end{theorem}

\begin{proof}
First, we convert the differential equation into an integral equation. Since $\bv_i(t)$ is a solution,
\[
(\dd/\dd t)\bv_i(t) = f(\bv_i(t)).
\]
Integrating both sides from $0$ to $t$ and applying the Fundamental Theorem of Calculus yields:
\begin{equation}
\bv_i(t) - \bv_i(0) = \int_0^t f(\bv_i(s)) \dd s \implies \bv_i(t) = \bv_0 + \int_0^t f(\bv_i(s)) \dd s
\end{equation}
Now, consider the difference between the two solutions. Let
\[
u(t) = \|\bv_1(t) - \bv_2(t)\|.
\]
Subtracting the integral equations:
\begin{equation}
\bv_1(t) - \bv_2(t) = \int_0^t \Big( f(\bv_1(s)) - f(\bv_2(s)) \Big) \dd s
\end{equation}
Taking the norm of both sides and applying the triangle inequality for integrals,
\[
\|\int \mathbf{g}\| \leq \int \|\mathbf{g}\|
\]
gives:
\begin{equation}
\begin{aligned}
u(t) &= \left\| \int_0^t \Big( f(\bv_1(s)) - f(\bv_2(s)) \Big) \dd s \right\| \\
&\leq \int_0^t \left\| f(\bv_1(s)) - f(\bv_2(s)) \right\| \dd s
\end{aligned}
\end{equation}
We now explicitly apply the Lipschitz condition
\[
\|f(\bx) - f(\by)\| \leq L\|\bx - \by\|
\]
to the integrand:
\begin{equation}
\begin{aligned}
u(t) &\leq \int_0^t L \|\bv_1(s) - \bv_2(s)\| \dd s \\
&= L \int_0^t u(s) \dd s
\end{aligned}
\end{equation}

This is exactly the hypothesis of Gronwall's Inequality with $C = 0$ and $\beta(s) = L$. Applying the lemma:
\begin{equation}
\begin{aligned}
u(t) &\leq 0 \cdot \exp\left( \int_0^t L \dd s \right) \\
&= 0
\end{aligned}
\end{equation}
By the definition of a norm, the difference satisfies
\[
\begin{aligned}
u(t) &= \|\bv_1(t) - \bv_2(t)\| \\
&\geq 0
\end{aligned}
\]
Consequently, we must have $u(t) = 0$ for all $t \in [0, T]$. Therefore, $\bv_1(t) = \bv_2(t)$.
\end{proof}

\subsection{Maximal Interval of Existence}

Solutions to ODEs are guaranteed to exist locally. By continuously patching together local solutions, we can extend the solution until it either exists for all time or ``blows up''.

\begin{definition}[Maximal Interval of Existence]
The \textit{maximal interval of existence} $(T_-, T_+)$ is the largest open interval on which a solution $\bv(t)$ to the initial value problem is defined. A solution is called \textit{global} if $T_+ = \infty$ and $T_- = -\infty$.
\end{definition}

\begin{theorem}[Escape Lemma (Finite-Time Blow-up)]
Let $\bv(t)$ be the maximal solution defined on $(T_-, T_+)$. If the forward maximal time is finite, i.e., $T_+ < \infty$, then the trajectory must escape to infinity:
\begin{equation}
\lim_{t \to T_+^-} \|\bv(t)\| = \infty
\end{equation}

\end{theorem}

\begin{proof}
We prove this by contradiction. Suppose $T_+ < \infty$ but the trajectory does \textit{not} escape to infinity. This means there exists a sequence of times $t_n \to T_+^-$ such that $\|\bv(t_n)\|$ is bounded. More rigorously, if the limit is not infinity, there must exist a compact set $K \subset \R^N$ such that $\bv(t) \in K$ for all $t \in [0, T_+)$.

Because $K$ is compact and $f$ is locally Lipschitz, $f$ is uniformly Lipschitz on $K$ with some constant $L_K$, and $f$ is bounded on $K$, meaning there exists $M_K > 0$ such that $\|f(\bx)\| \leq M_K$ for all $\bx \in K$.
Since $\dot{\bv}(t) = f(\bv(t))$ and $\bv(t) \in K$, we have $\|\dot{\bv}(t)\| \leq M_K$. By the Mean Value Theorem, for any $t_1, t_2 \in [0, T_+)$:
\begin{equation}
\|\bv(t_1) - \bv(t_2)\| \leq M_K |t_1 - t_2|
\end{equation}

This shows that $\bv(t)$ is uniformly continuous on $[0, T_+)$. Because $\R^N$ is a complete metric space and $\bv(t)$ is uniformly continuous and bounded, the limit
\[
\bv_+ = \lim_{t \to T_+^-} \bv(t)
\]
must exist and is finite. We can now define a new initial value problem starting at time $T_+$ with the initial condition $\bv(T_+) = \bv_+$. By the local existence theorem (Picard--Lindelöf), there exists a solution $\tilde{\bv}(t)$ defined on some interval $[T_+, T_+ + \delta)$ for $\delta > 0$.
We can concatenate the original solution $\bv(t)$ and the new solution $\tilde{\bv}(t)$ to form a continuous, differentiable solution defined on $[0, T_+ + \delta)$. This strictly extends the domain of the solution past $T_+$, which directly contradicts the assumption that $(T_-, T_+)$ was the \textit{maximal} interval of existence.
Therefore, our assumption that the trajectory remains in a compact set must be false; it must escape to infinity.
\end{proof}

\begin{remark}[Practical Consequence]
The Escape Lemma provides a powerful practical tool: if you can prove that a trajectory is confined to a bounded region (for example, by finding a Lyapunov function or an invariant set), you immediately guarantee that $T_+ = \infty$, meaning the solution exists globally for all future time.
\end{remark}

\section{Stability and Local Dynamics}

\subsection{Lyapunov Stability}

When analyzing the long-term behavior of a system, we often study the behavior near an equilibrium point.

\begin{definition}[Equilibrium and Stability]
A point $\bv_\infty \in \R^N$ is an \textit{equilibrium} (or fixed point) if $f(\bv_\infty) = \mathbf{0}$. The constant function $\bv(t) = \bv_\infty$ is then a trivial solution.
The equilibrium $\bv_\infty$ is:
\begin{itemize}
\item \textit{Lyapunov stable}: if for every $\varepsilon > 0$, there exists a $\delta > 0$ such that if $\|\bv(0) - \bv_\infty\| < \delta$, then $\|\bv(t) - \bv_\infty\| < \varepsilon$ for all $t \geq 0$.
\item \textit{Asymptotically stable}: if it is Lyapunov stable, and there exists a $\delta_0 > 0$ such that if $\|\bv(0) - \bv_\infty\| < \delta_0$, then

\[
\lim_{t \to \infty} \bv(t) = \bv_\infty.
\]

\end{itemize}
\end{definition}

\begin{remark}[Geometric Meaning]
Lyapunov stability means that trajectories starting sufficiently close to the equilibrium remain arbitrarily close to it forever; they do not necessarily converge to it, but they never escape a small tube around it. Asymptotic stability adds the requirement that they actually converge to the equilibrium.
\end{remark}

To prove stability without explicitly solving the ODE, we use Lyapunov's Direct Method, which relies on energy-like scalar functions.

\begin{definition}[Lyapunov Function and Orbital Derivative]
Let $U$ be an open neighborhood of $\bv_\infty$. A continuously differentiable ($C^1$) function $V: U \to \R$ is a \textit{Lyapunov function} if:
\begin{enumerate}
\item \textit{Positive definiteness}: $V(\bv_\infty) = 0$ and $V(\bv) > 0$ for all $\bv \in U \setminus \{\bv_\infty\}$.
\item \textit{Orbital derivative}: the orbital derivative $\dot{V}(\bv) \leq 0$ for all $\bv \in U$.
\end{enumerate}
The \textit{orbital derivative} $\dot{V}(\bv)$ is the time derivative of $V$ evaluated along the trajectories of the system. By the multivariate chain rule:
\begin{equation}
\begin{aligned}
\dot{V}(\bv) &= \frac{\dd}{\dd t} V(\bv(t)) \bigg|_{t=0} \\
&= \nabla V(\bv) \cdot \dot{\bv} \\
&= \nabla V(\bv) \cdot f(\bv)
\end{aligned}
\end{equation}

\end{definition}

\begin{theorem}[Lyapunov's Direct Method]
If a Lyapunov function $V$ exists for $\bv_\infty$, then $\bv_\infty$ is Lyapunov stable. If, furthermore, $\dot{V}(\bv) < 0$ for all $\bv \in U \setminus \{\bv_\infty\}$ (strictly negative definite), then $\bv_\infty$ is asymptotically stable.
\end{theorem}

\begin{proof}[Rigorous Proof of Stability]
Let $\varepsilon > 0$ be given. We must find a $\delta > 0$ such that trajectories starting in the ball $B_\delta(\bv_\infty)$ never leave the ball $B_\varepsilon(\bv_\infty)$.
Consider the boundary of the $\varepsilon$-ball,
\[
S_\varepsilon = \{ \bv : \|\bv - \bv_\infty\| = \varepsilon \}.
\]

This set is closed and bounded, hence compact. Since $\bv_\infty \notin S_\varepsilon$ and $V$ is continuous and positive definite, $V$ attains a strictly positive minimum on $S_\varepsilon$. Let:
\begin{equation}
m = \min_{\bv \in S_\varepsilon} V(\bv) > 0
\end{equation}
Define the sublevel set
\[
U_m = \{ \bv \in U : V(\bv) < m \}.
\]

Because $V$ is continuous, $U_m$ is an open set. Since $V(\bv_\infty) = 0 < m$, $\bv_\infty \in U_m$. Therefore, there exists a $\delta > 0$ such that the ball $B_\delta(\bv_\infty) \subset U_m$.
Now, suppose $\bv(0) \in B_\delta(\bv_\infty)$. Then $V(\bv(0)) < m$.
Along the trajectory, the time derivative of $V$ is
\[
\begin{aligned}
(\dd/\dd t) V(\bv(t)) &= \dot{V}(\bv(t)) \\
&\leq 0.
\end{aligned}
\]
This means $V(\bv(t))$ is a monotonically non-increasing function of time. Thus, for all $t \geq 0$:
\begin{equation}
V(\bv(t)) \leq V(\bv(0)) < m
\end{equation}

This implies that $\bv(t) \in U_m$ for all $t \geq 0$. By the definition of $m$, any point on the boundary $S_\varepsilon$ has $V(\bv) \geq m$. Since $V(\bv(t)) < m$, the trajectory can never reach $S_\varepsilon$. Because the trajectory is continuous and starts inside $B_\varepsilon$, it can never cross $S_\varepsilon$ to exit the ball. Hence, $\|\bv(t) - \bv_\infty\| < \varepsilon$ for all $t \geq 0$, proving Lyapunov stability.
\end{proof}

\begin{proof}[Rigorous Proof of Asymptotic Stability]
We already know the equilibrium is stable, so trajectories starting sufficiently close remain close. We must now prove they converge to $\bv_\infty$.
Assume $\dot{V}(\bv) < 0$ for $\bv \neq \bv_\infty$. For a trajectory starting near $\bv_\infty$, $V(\bv(t))$ is monotonically decreasing and bounded below by $0$. By the Monotone Convergence Theorem for real sequences, the limit must exist:
\begin{equation}
\begin{aligned}
\lim_{t \to \infty} V(\bv(t)) &= c \\
&\geq 0
\end{aligned}
\end{equation}

We must show that $c = 0$. Suppose, for contradiction, that $c > 0$.
Because the trajectory is stable, it is confined to a compact set
\[
K = \{ \bv : c \leq V(\bv) \leq V(\bv(0)) \}.
\]

Note that $\bv_\infty \notin K$ because $V(\bv_\infty) = 0 < c$.
The orbital derivative $\dot{V}(\bv)$ is a continuous function on the compact set $K$. Since $\dot{V}(\bv) < 0$ everywhere on $K$, it must attain a strict maximum on $K$. Let this maximum be $-\eta$, where $\eta > 0$. Thus, for all $\bv \in K$:
\begin{equation}
\dot{V}(\bv) \leq -\eta < 0
\end{equation}
Evaluating this along the trajectory $\bv(t) \in K$:
\begin{equation}
\frac{\dd}{\dd t} V(\bv(t)) \leq -\eta
\end{equation}
Integrating both sides from $0$ to $t$:
\begin{equation}
V(\bv(t)) - V(\bv(0)) \leq -\eta t \implies V(\bv(t)) \leq V(\bv(0)) - \eta t
\end{equation}

As $t \to \infty$, the right-hand side approaches $-\infty$. This implies $V(\bv(t)) \to -\infty$, which strictly contradicts the fact that $V(\bv) \geq 0$ for all $\bv$.
Therefore, our assumption $c > 0$ must be false, and we conclude $c = 0$.
Since $V(\bv(t)) \to 0$ and $V$ is continuous and positive definite (meaning
\[
V(\bv) = 0 \iff \bv = \bv_\infty
\]
), it follows rigorously that
\[
\lim_{t \to \infty} \bv(t) = \bv_\infty.
\]

\end{proof}

\subsection{Hartman--Grobman and Center Manifolds}

When a Lyapunov function is difficult to construct, we often analyze the local behavior near an equilibrium by linearizing the system.

\begin{definition}[Jacobian and Hyperbolic Equilibrium]
The \textit{Jacobian matrix} of $f$ at $\bv_\infty$ is the $N \times N$ matrix $J$ of partial derivatives, with entries
\[
J_{ij} = (\partial f_i/\partial v_j)(\bv_\infty).
\]

The equilibrium $\bv_\infty$ is called \textit{hyperbolic} if every eigenvalue $\lambda$ of $J$ has a nonzero real part, i.e., $\Re(\lambda) \neq 0$.
\end{definition}

\begin{theorem}[Hartman--Grobman Theorem]
If $\bv_\infty$ is a hyperbolic equilibrium, then the nonlinear flow near $\bv_\infty$ is \textit{topologically conjugate} to the linear flow $\dot{\bepsilon} = J\bepsilon$.
\end{theorem}

\begin{remark}[Unpacking Topological Conjugacy]
What does ``topologically conjugate'' mean? It means there exists a homeomorphism (a continuous, bijective map with a continuous inverse) $h: U \to V$, where $U$ is a neighborhood of $\bv_\infty$ and $V$ is a neighborhood of the origin in $\R^N$, such that $h(\bv_\infty) = \mathbf{0}$, and $h$ maps trajectories of the nonlinear system to trajectories of the linear system, preserving the direction of time. Formally, if $\Phi_t(\bx)$ is the nonlinear flow and $e^{Jt}$ is the linear flow, then:
\begin{equation}
h(\Phi_t(\bx)) = e^{Jt} h(\bx)
\end{equation}

This implies that the phase portrait of the nonlinear system near a hyperbolic equilibrium is qualitatively identical to the phase portrait of its linearization. The topological structure (e.g., whether it is a node, saddle, or spiral) is completely determined by the eigenvalues of $J$.
\end{remark}

\begin{remark}[Why Hyperbolicity is Required]
If $J$ has an eigenvalue with zero real part ($\Re(\lambda) = 0$), the equilibrium is \textit{non-hyperbolic}. In this case, the linear terms are ``neutral'' (neither exponentially growing nor decaying), and the nonlinear higher-order terms in the Taylor expansion of $f$ can dominate the local dynamics. The nonlinear system could exhibit behaviors completely absent in the linearization, such as a center (closed orbits) or a weak focus. Thus, linearization is inconclusive for non-hyperbolic equilibria.
\end{remark}

To handle non-hyperbolic equilibria, we use the Center Manifold Theorem, which allows us to reduce the dimensionality of the system.

\begin{definition}[Spectral Subspaces]
Based on the eigenvalues of the Jacobian $J$, we decompose $\R^N$ into three invariant linear subspaces:
\begin{itemize}
\item \textit{Stable subspace $E^s$}: the span of all generalized eigenvectors corresponding to eigenvalues with $\Re(\lambda) < 0$.
\item \textit{Unstable subspace $E^u$}: the span of all generalized eigenvectors corresponding to eigenvalues with $\Re(\lambda) > 0$.
\item \textit{Center subspace $E^c$}: the span of all generalized eigenvectors corresponding to eigenvalues with $\Re(\lambda) = 0$.
\end{itemize}
By linear algebra, $\R^N = E^s \oplus E^u \oplus E^c$.
\end{definition}

\begin{theorem}[Center Manifold Theorem]
There exists a locally invariant manifold $W^c$ (the center manifold) passing through $\bv_\infty$ such that:
\begin{enumerate}
\item \textit{Tangency}: $W^c$ is tangent to the center subspace $E^c$ at $\bv_\infty$.
\item \textit{Local invariance}: $W^c$ is locally invariant: if a trajectory starts in $W^c$, it remains in $W^c$ for all time (as long as it stays in the neighborhood).
\item \textit{Dimension}: the dimension of $W^c$ is equal to the dimension of $E^c$.
\end{enumerate}
Furthermore, there exist locally invariant stable and unstable manifolds $W^s$ and $W^u$ tangent to $E^s$ and $E^u$, respectively.
\end{theorem}

\begin{remark}[Why is this Useful?]
The dynamics on the stable manifold $W^s$ are trivial: trajectories exponentially decay to $\bv_\infty$. The dynamics on the unstable manifold $W^u$ are also well-understood: trajectories exponentially repel from $\bv_\infty$.
The complex, interesting dynamics (bifurcations, stability transitions, limit cycles) occur entirely on the center manifold $W^c$. Because
\[
\dim(W^c) = \dim(E^c) < N
\]

 (assuming $E^s$ or $E^u$ are non-trivial), the Center Manifold Theorem allows us to reduce the study of the full $N$-dimensional nonlinear system to a lower-dimensional system restricted to $W^c$. On this reduced space, the linear terms are zero, and the dynamics are governed purely by the nonlinear terms, which can be analyzed using techniques like Lyapunov functions or normal forms.
\end{remark}

Figure~\ref{fig:ch15-phase-portraits} compares the local phase portraits of four types of planar equilibrium.

\begin{figure}[!htbp]
\centering
\begin{tikzpicture}
\begin{axis}[width=6.1cm,height=6.1cm,axis lines=middle,xmin=-2,xmax=2,ymin=-2,ymax=2,
  xtick=\empty,ytick=\empty,]
  \foreach \sx/\sy in {1/1,-1/1,1/-1,-1/-1}{
    \foreach \c in {0.4,0.8,1.4,2.2}{
      \addplot[figblue,thick,domain=0:2.2,samples=50,forget plot]
        ({\sx*\c*exp(-0.6*x)},{\sy*\c*0.4*exp(-0.3*x)});
    }
    \addplot[figblue,thick,-{Stealth[length=2.6mm]},forget plot] coordinates {
      ({\sx*1.4*exp(-0.6*1.55)},{\sy*1.4*0.4*exp(-0.3*1.55)})
      ({\sx*1.4*exp(-0.6*1.85)},{\sy*1.4*0.4*exp(-0.3*1.85)})
    };
  }
\end{axis}
\begin{axis}[at={(6.6cm,0)},width=6.1cm,height=6.1cm,axis lines=middle,xmin=-2,xmax=2,ymin=-2,ymax=2,
  xtick=\empty,ytick=\empty,]
  \foreach \c in {0.3,0.7,1.2,2}{
    \addplot[figred,thick,domain=-1.5:1.5,samples=60,forget plot] ({x},{\c/(x+0.001)});
    \addplot[figred,thick,domain=-1.5:1.5,samples=60,forget plot] ({x},{-\c/(x+0.001)});
  }
  \draw[vecB] (-1.9,-1.9) -- (-0.35,-0.35);
  \draw[vecB] (1.9,1.9) -- (0.35,0.35);
  \draw[vecB] (0.35,-0.35) -- (1.9,-1.9);
  \draw[vecB] (-0.35,0.35) -- (-1.9,1.9);
\end{axis}
\begin{axis}[at={(0,-6.6cm)},width=6.1cm,height=6.1cm,axis lines=middle,xmin=-2,xmax=2,ymin=-2,ymax=2,
  xtick=\empty,ytick=\empty,]
  \addplot[figgreen,thick,domain=0:12,samples=200,forget plot,variable=t]
    ({1.9*exp(-0.15*t)*cos(deg(t))},{1.9*exp(-0.15*t)*sin(deg(t))});
  \addplot[figgreen,thick,-{Stealth[length=2.6mm]},forget plot] coordinates {
    ({1.9*exp(-0.15*7.4)*cos(deg(7.4))},{1.9*exp(-0.15*7.4)*sin(deg(7.4))})
    ({1.9*exp(-0.15*7.8)*cos(deg(7.8))},{1.9*exp(-0.15*7.8)*sin(deg(7.8))})
  };
\end{axis}
\begin{axis}[at={(6.6cm,-6.6cm)},width=6.1cm,height=6.1cm,axis lines=middle,xmin=-2,xmax=2,ymin=-2,ymax=2,
  xtick=\empty,ytick=\empty,]
  \foreach \r in {0.5,0.9,1.3,1.7}{
    \addplot[black!55,thick,domain=0:360,samples=80,forget plot] ({\r*cos(x)},{\r*sin(x)});
  }
\end{axis}
\end{tikzpicture}
\caption[Phase portraits: node, saddle, focus, center]{Local phase portraits around the origin. Top left: stable node with two negative real eigenvalues. Top right: saddle with one negative and one positive real eigenvalue. Bottom left: stable focus with complex eigenvalues having negative real part. Bottom right: non-hyperbolic center with purely imaginary eigenvalues. Arrows indicate the direction of evolution where shown.}
\label{fig:ch15-phase-portraits}
\end{figure}

\section{Partial Differential Equations: Solution Methods}
\label{sec:pde-methods}

The ordinary differential equations studied so far describe finitely many state variables evolving in time. Spatially distributed systems require an unknown field $u(\bx,t)$ and partial derivatives in both space and time. The integral theorems, spectral methods, and distribution theory developed in earlier chapters now provide tools for solving these equations.

\subsection{Fundamental PDEs and Separation of Variables}

Three fundamental linear second-order partial differential equations are
\begin{align}
\text{Heat (diffusion) equation:}\qquad &\frac{\partial u}{\partial t}=D\,\Delta u,\\
\text{Wave equation:}\qquad &\frac{\partial^2u}{\partial t^2}=c^2\,\Delta u,\\
\text{Laplace (steady-state) equation:}\qquad &\Delta u=0,
\end{align}
Here the Laplacian is
\[
\begin{aligned}
\Delta u &= \nabla\cdot(\nabla u) \\
&= \sum_i\partial^2u/\partial x_i^2,
\end{aligned}
\]

The constant $D>0$ is the diffusion coefficient, and $c>0$ is the propagation speed. In mathematical biology, the diffusion equation models the spatial spreading of a cellular or molecular concentration. Initial and boundary conditions are needed to select a solution appropriate to the physical problem.


The basic idea is to seek a factored solution
\[
u(\bx,t)=X(\bx)T(t).
\]

Substitution into a suitable linear PDE separates the spatial and temporal dependence, linking them by a separation constant. In one spatial dimension, this produces two ordinary differential equations; in higher dimensions, the spatial factor generally satisfies an eigenvalue PDE.

\begin{example}[The Heat Equation with Homogeneous Dirichlet Conditions]
\label{ex:heat-separation}
Consider $u_t=D\,u_{xx}$ for $0<x<L$ and $t>0$, with
\[
u(0,t)=u(L,t)=0,\qquad u(x,0)=f(x).
\]

The boundary conditions are called \emph{homogeneous Dirichlet conditions} because the function vanishes at the boundary.

\textit{Step 1: separate the variables.} Substituting

\[
u(x,t)=X(x)T(t)
\]
gives $XT'=DX''T$. Wherever the factors are nonzero,
\begin{equation}
\frac{T'(t)}{D\,T(t)}=\frac{X''(x)}{X(x)}.
\end{equation}

The left-hand side depends only on $t$ and the right-hand side only on $x$. They must therefore equal a common constant, denoted by $-\lambda$:
\begin{equation}
T'=-\lambda D\,T,\qquad X''=-\lambda X.
\end{equation}

These equations are then imposed throughout the interval, including any zeros of $X$; the quotient is only a device for identifying the separated equations.

\textit{Step 2: solve the spatial eigenvalue problem.} A nonzero temporal factor requires

\[
\begin{aligned}
X(0) &= X(L) \\
&= 0.
\end{aligned}
\]
For $\lambda>0$,
\[
X(x)=A\cos(\sqrt\lambda x)+B\sin(\sqrt\lambda x).
\]

The first boundary condition gives $A=0$, and the second requires $\sin(\sqrt{\lambda} L)=0$ for a nontrivial solution. Hence
\begin{equation}
\lambda_n=\left(\frac{n\pi}{L}\right)^2,
\qquad X_n(x)=\sin\left(\frac{n\pi x}{L}\right),\qquad n=1,2,\dots.
\end{equation}
No nontrivial solution has $\lambda\le0$: integration by parts in $-X''=\lambda X$ gives
\[
\int_0^L|X'|^2\,\dd x=\lambda\int_0^L|X|^2\,\dd x,
\]
and the zero endpoint values rule out a nonzero constant function.

\textit{Step 3: solve the temporal equation and superpose.} For each eigenvalue,

\[
T_n(t)=e^{-\lambda_nDt}
\]
up to a multiplicative constant. Linearity permits the expansion
\begin{equation}
u(x,t)=\sum_{n=1}^{\infty}b_n\sin\left(\frac{n\pi x}{L}\right)e^{-\lambda_nDt}.
\end{equation}

Matching the initial data amounts to expanding $f$ in a Fourier sine series. Orthogonality on $(0,L)$ gives
\begin{equation}
b_n=\frac{2}{L}\int_0^L f(x)\sin\left(\frac{n\pi x}{L}\right)\,\dd x.
\end{equation}

For $f\in L^2(0,L)$, the initial condition is attained in the $L^2$ sense as $t\downarrow0$, and the exponential factors smooth the series for $t>0$. With suitable smoothness and endpoint compatibility, it gives the classical solution up to the initial boundary. The functions $X_n$ are eigenfunctions of the self-adjoint Dirichlet operator $-\dd^2/\dd x^2$, just as orthogonal eigenvectors arise for symmetric matrices. Thus separation of variables connects ODE methods, Fourier expansions, and operator theory.
\end{example}

Figure~\ref{fig:atlas-heat-mode-decay} illustrates the different decay rates of separated heat-equation modes.

\begin{figure}[!htbp]
\centering
\begin{tikzpicture}[>=Stealth]
\begin{axis}[width=11cm,height=5.8cm,axis lines=middle,ticklabel style={font=\small},label style={font=\small},legend style={font=\scriptsize,draw=black!20,fill=white},xlabel={$x$},ylabel={$u(x,t)$},xmin=0,xmax=1.05,ymin=0,ymax=1.3,xtick={0,0.5,1},legend columns=3,legend pos=north east]
\addplot[figblue,very thick,domain=0:1,samples=150] {sin(deg(pi*x))+0.35*sin(deg(3*pi*x))};
\addlegendentry{$t=0$}
\addplot[figred,thick,dashed,domain=0:1,samples=150] {exp(-pi^2*0.03)*sin(deg(pi*x))+0.35*exp(-9*pi^2*0.03)*sin(deg(3*pi*x))};
\addlegendentry{$t=0.03$}
\addplot[figgreen,thick,dashdotted,domain=0:1,samples=150] {exp(-pi^2*0.1)*sin(deg(pi*x))+0.35*exp(-9*pi^2*0.1)*sin(deg(3*pi*x))};
\addlegendentry{$t=0.1$}
\end{axis}
\end{tikzpicture}
\caption{Three profiles solve $u_t=u_{xx}$ on $[0,1]$ with zero boundary values. The initial profile combines the first sine mode with $0.35$ times the third. The third mode decays nine times faster in the exponent, so small-scale structure disappears before the main profile relaxes to zero.}
\label{fig:atlas-heat-mode-decay}
\end{figure}

\subsection{Dirichlet and Neumann Boundary-Value Problems}
\label{sec:pde-boundary-problems}

\begin{definition}[Dirichlet and Neumann Problems]
Let $\Omega\subset\R^N$ be a bounded connected domain with smooth boundary. A \emph{Dirichlet problem} for the Laplace equation prescribes the boundary values:
\begin{equation}
\Delta u=0\quad\text{in }\Omega,\qquad u=g\quad\text{on }\partial\Omega.
\end{equation}
A \emph{Neumann problem} instead prescribes the outward normal derivative:
\begin{equation}
\Delta u=0\quad\text{in }\Omega,\qquad
\frac{\partial u}{\partial\mathbf{n}}=h\quad\text{on }\partial\Omega,
\end{equation}
where $\partial u/\partial\mathbf{n}:=\nabla u\cdot\mathbf{n}$. This specifies a quantity proportional to the boundary flux rather than the value of $u$ itself; for a diffusive current $\mathbf{J}=-D\nabla u$, the outward flux is $-D h$.
\end{definition}

\begin{theorem}[Uniqueness for the Dirichlet Problem]
If a classical solution of the Dirichlet problem exists with sufficient regularity for integration by parts, it is unique.
\end{theorem}

\begin{proof}[The Energy Method]
Let $u_1,u_2$ be two solutions and set $w=u_1-u_2$. Then $\Delta w=0$ in $\Omega$ and $w=0$ on $\partial\Omega$. The product rule gives
\begin{equation}
\begin{aligned}
\nabla\cdot(w\nabla w) &= |\nabla w|^2+w\Delta w \\
&= |\nabla w|^2.
\end{aligned}
\end{equation}
Integrating and applying the $N$-dimensional version of the divergence theorem yields
\begin{equation}
\int_\Omega|\nabla w|^2\,\dd\bx
=\int_{\partial\Omega}w\frac{\partial w}{\partial\mathbf{n}}\,\dd S=0.
\end{equation}

The continuous nonnegative integrand must vanish identically. Thus $\nabla w=0$, and connectedness implies that $w$ is constant. Its zero boundary value forces that constant to be zero, so $u_1=u_2$.
\end{proof}

\begin{remark}[Neumann Uniqueness and Compatibility]
For two solutions of the same Neumann problem, the difference satisfies $\partial w/\partial\mathbf{n}=0$, which again makes the energy boundary term vanish. Therefore $w$ is constant, but the constant need not be zero: the solution is unique only up to an additive constant. Existence also requires the compatibility condition
\begin{equation}
\int_{\partial\Omega}h\,\dd S=\int_\Omega\Delta u\,\dd\bx=0.
\end{equation}
Prescribing, for example, the mean of $u$ fixes the additive constant.
\end{remark}

\subsection{Green's Functions}
\label{sec:green-functions}

We now consider Poisson's equation $\Delta u=-\rho$ in $\Omega$, with homogeneous Dirichlet data. The idea parallels the use of the Dirac delta in distribution theory: if the response to a point source at $\by$ is known, the response to a distributed source $\rho$ can be reconstructed by integrating these point-source responses, weighted by $\rho(\by)$.

\begin{definition}[Dirichlet Green's Function]
A \emph{Green's function} for $-\Delta$ on $\Omega$ with homogeneous Dirichlet conditions satisfies, for each fixed $\by\in\Omega$,
\begin{equation}
-\Delta_{\bx}G(\bx,\by)=\delta(\bx-\by)\quad\text{in }\Omega,
\qquad G(\bx,\by)=0\quad\text{on }\partial\Omega.
\end{equation}

The differential equation is interpreted in the distributional sense. The multidimensional Dirac delta acts on a test function by $\delta_{\by}[\varphi]=\varphi(\by)$.
\end{definition}

\begin{theorem}[Green's Function Representation]
For a sufficiently regular domain and source for which the Dirichlet Green's function representation is valid, the solution of $\Delta u=-\rho$ in $\Omega$, $u=0$ on $\partial\Omega$, is
\begin{equation}
u(\bx)=\int_\Omega G(\bx,\by)\rho(\by)\,\dd\by.
\end{equation}

\end{theorem}

\begin{proof}[Distributional Verification]
Formally applying the Laplacian gives
\begin{align}
\Delta_{\bx}u(\bx)
&=\int_\Omega\Delta_{\bx}G(\bx,\by)\rho(\by)\,\dd\by\notag\\
&=-\int_\Omega\delta(\bx-\by)\rho(\by)\,\dd\by=-\rho(\bx).
\end{align}

Because $G$ is singular at $\bx=\by$, differentiation here is understood distributionally rather than as an unqualified exchange of classical derivatives and integration. More explicitly, for a compactly supported smooth test function $\varphi$, the defining identity for $G$ and Fubini's theorem give
\begin{align}
\int_\Omega u(\bx)(-\Delta\varphi(\bx))\,\dd\bx
&=\int_\Omega\rho(\by)\left[\int_\Omega G(\bx,\by)(-\Delta\varphi(\bx))\,\dd\bx\right]\dd\by\notag\\
&=\int_\Omega\rho(\by)\varphi(\by)\,\dd\by.
\end{align}

Under the assumed regularity, the boundary trace is zero because the Green's function has zero Dirichlet trace. When the represented solution is classical, the preceding uniqueness theorem identifies it as the only solution.
\end{proof}

\begin{example}[The Upper Half-Plane and the Method of Images]
Consider the upper half-plane
\[
\Omega=\{\bx=(x_1,x_2):x_2>0\}\subset\R^2
\]

Let $\by=(y_1,y_2)$ with $y_2>0$. Reflect the source across the boundary to obtain $\by^*=(y_1,-y_2)$. Then
\begin{equation}
G(\bx,\by)=\frac{1}{2\pi}\log\frac{|\bx-\by^*|}{|\bx-\by|}
\end{equation}
is a Dirichlet Green's function for the upper half-plane. Indeed, the two-dimensional fundamental solution satisfies
\[
-\Delta_{\bx}\left(-\frac{1}{2\pi}\log|\bx-\by|\right)=\delta(\bx-\by).
\]

The added term $\log|\bx-\by^*|/(2\pi)$ is harmonic throughout $\Omega$ because its singularity lies outside the domain. On $x_2=0$, the distances to $\by$ and $\by^*$ coincide, so the logarithmic ratio vanishes. The normalization of the fundamental solution follows from the divergence theorem on a small circle around the source, just as $1/(4\pi|\bx-\by|)$ is the fundamental solution of $-\Delta$ in three dimensions.

Adding a fictitious exterior source to enforce the boundary condition is the \emph{method of images}. It is particularly useful on domains with suitable symmetries, such as half-spaces and balls, although the image location and strength depend on the geometry. On an unbounded domain, appropriate behavior at infinity must also be specified to obtain uniqueness.
\end{example}
\newpage
\section{Harmonic Analysis and Fourier Transforms}

\subsection{Fourier Transform and Inversion}

\begin{definition}[Fourier Transform]
Let $f \in L^1(\R)$ be a complex-valued, Lebesgue-integrable function. The Fourier transform of $f$, denoted $\hat{f}$, is defined as the function $\hat{f}: \R \to \C$ given by:
\begin{equation}
\hat{f}(\omega) = \int_\R f(t) e^{-i\omega t} \dd t
\end{equation}

\end{definition}

\begin{remark}[Why $L^1(\R)$?]
The requirement that $f \in L^1(\R)$ is strictly necessary to guarantee that the integral converges absolutely for every $\omega \in \R$. Because the complex exponential is bounded by unity ($|e^{-i\omega t}| = 1$), we have:
\begin{equation}
\int_\R |f(t) e^{-i\omega t}| \dd t = \int_\R |f(t)| \dd t = \|f\|_{L^1} < \infty
\end{equation}
Absolute convergence ensures that $\hat{f}(\omega)$ is a well-defined, finite complex number for all $\omega$.
\end{remark}

\begin{lemma}[Riemann--Lebesgue Lemma]
If $f \in L^1(\R)$, then its Fourier transform $\hat{f}$ is continuous, bounded, and vanishes at infinity:
\[
\lim_{|\omega| \to \infty} \hat{f}(\omega) = 0.
\]

\end{lemma}

\begin{proof}
We prove the three properties sequentially.

\textit{1. Boundedness:}
By the triangle inequality for integrals:
\begin{equation}
\begin{aligned}
|\hat{f}(\omega)| &= \left| \int_\R f(t) e^{-i\omega t} \dd t \right| \\
&\leq \int_\R |f(t)| |e^{-i\omega t}| \dd t \\
&= \int_\R |f(t)| \dd t \\
&= \|f\|_{L^1}
\end{aligned}
\end{equation}
Thus, $\hat{f}$ is bounded by the $L^1$ norm of $f$.

\textit{2. Continuity:}
Let $\omega_n$ be a sequence converging to $\omega$. The integrand $f(t) e^{-i\omega_n t}$ converges pointwise to $f(t) e^{-i\omega t}$ for every $t$. Furthermore, the absolute value of the integrand is bounded by $|f(t)|$, which is an integrable dominating function since $f \in L^1(\R)$. By the \textit{Dominated Convergence Theorem}, we can pass the limit inside the integral:
\begin{equation}
\begin{aligned}
    \lim_{n \to \infty} \hat{f}(\omega_n)
    &= \int_\R \lim_{n \to \infty} \left( f(t) e^{-i\omega_n t} \right) \dd t \\
    &= \int_\R f(t) e^{-i\omega t} \dd t = \hat{f}(\omega)
    \end{aligned}
\end{equation}
This proves $\hat{f}$ is continuous.

\textit{3. Decay at Infinity:}
We use a phase-shift trick. Let $\omega \neq 0$. We perform the change of variables $s = t + \pi/\omega$, which implies $t = s - \pi/\omega$ and $\dd t = \dd s$. Notice that
\[
\begin{aligned}
e^{-i\omega(s - \pi/\omega)} &= e^{-i\omega s} e^{i\pi} \\
&= -e^{-i\omega s}.
\end{aligned}
\]
Substituting this into the definition:
\begin{equation}
\begin{aligned}
\hat{f}(\omega) &= \int_\R f(s - \pi/\omega) e^{-i\omega (s - \pi/\omega)} \dd s \\
&= -\int_\R f(s - \pi/\omega) e^{-i\omega s} \dd s
\end{aligned}
\end{equation}

Adding this expression to the original definition of $\hat{f}(\omega)$ (and renaming the dummy variable $s$ back to $t$):
\begin{equation}
2\hat{f}(\omega) = \int_\R \left[ f(t) - f(t - \pi/\omega) \right] e^{-i\omega t} \dd t
\end{equation}
Taking the absolute value and applying the triangle inequality:
\begin{equation}
\begin{aligned}
2|\hat{f}(\omega)| &\leq \int_\R \left| f(t) - f(t - \pi/\omega) \right| \dd t \\
&= \left\| f(\cdot) - f(\cdot - \pi/\omega) \right\|_{L^1}
\end{aligned}
\end{equation}

To show this goes to zero as $|\omega| \to \infty$, we must explicitly prove that translation is continuous in the $L^1$ norm. Let $\tau = \pi/\omega$. We want to show \[
\lim_{\tau \to 0} \|f - f_\tau\|_{L^1} = 0,
\] where $f_\tau(t) = f(t-\tau)$.
Let $\varepsilon > 0$. The space of continuous functions with compact support, $C_c(\R)$, is dense in $L^1(\R)$. Thus, there exists a function $g \in C_c(\R)$ such that $\|f - g\|_{L^1} < \varepsilon/3$. By the triangle inequality:
\begin{equation}
\|f - f_\tau\|_{L^1} \leq \|f - g\|_{L^1} + \|g - g_\tau\|_{L^1} + \|g_\tau - f_\tau\|_{L^1}
\end{equation}
Because the Lebesgue measure is translation-invariant,
\[
\|g_\tau - f_\tau\|_{L^1} = \|g - f\|_{L^1} < \varepsilon/3.
\]

Since $g$ is continuous and has compact support, it is uniformly continuous. Let $K$ be a compact set containing the support of $g$ and $g_\tau$ for small $\tau$. For sufficiently small $|\tau|$, $|g(t) - g(t-\tau)| < (\varepsilon/(3 \mu(K)))$ for all $t$. Thus:
\begin{equation}
\|g - g_\tau\|_{L^1} = \int_K |g(t) - g(t-\tau)| \dd t < \mu(K) \frac{\varepsilon}{3 \mu(K)} = \varepsilon/3
\end{equation}

Combining these, $\|f - f_\tau\|_{L^1} < \varepsilon$. Since $\pi/\omega \to 0$ as $|\omega| \to \infty$, we conclude that $\|f - f_{\pi/\omega}\|_{L^1} \to 0$, which forces $\hat{f}(\omega) \to 0$.
\end{proof}

\begin{theorem}[Fourier Inversion]
If $f \in L^1(\R)$ and its Fourier transform $\hat{f} \in L^1(\R)$, then for almost every $t \in \R$:
\begin{equation}
f(t) = \frac{1}{2\pi} \int_\R \hat{f}(\omega) e^{i\omega t} \dd \omega
\end{equation}
If $f$ is continuous, this equality holds for all $t \in \R$.
\end{theorem}

\begin{proof}
The integral \[
\int \hat{f}(\omega) e^{i\omega t} \dd \omega
\] might not converge absolutely if we merely know $\hat{f} \in L^1$ and attempt to substitute the definition of $\hat{f}$ directly, because the resulting double integral may fail Fubini's conditions. To bypass this, we introduce a regularization factor $e^{-\varepsilon \omega^2}$ (Gaussian damping) and take the limit as $\varepsilon \to 0^+$. Define the regularized integral:
\begin{equation}
I_\varepsilon(t) = \frac{1}{2\pi} \int_\R \hat{f}(\omega) e^{-\varepsilon \omega^2} e^{i\omega t} \dd \omega
\end{equation}
The two facts needed for dominated convergence are
\[
\begin{gathered}
\hat{f} \in L^1 \\
|e^{-\varepsilon \omega^2} e^{i\omega t}| \leq 1,
\end{gathered}
\]

Together, they show that the integrand is dominated by $|\hat{f}(\omega)| \in L^1$. By the Dominated Convergence Theorem, as $\varepsilon \to 0^+$:
\begin{equation}
\lim_{\varepsilon \to 0^+} I_\varepsilon(t) = \frac{1}{2\pi} \int_\R \hat{f}(\omega) e^{i\omega t} \dd \omega
\end{equation}
Now, substitute the definition of \[
\hat{f}(\omega) = \int_\R f(s) e^{-i\omega s} \dd s
\] into $I_\varepsilon(t)$:
\begin{equation}
I_\varepsilon(t) = \frac{1}{2\pi} \int_\R \left[ \int_\R f(s) e^{-i\omega s} \dd s \right] e^{-\varepsilon \omega^2} e^{i\omega t} \dd \omega
\end{equation}

We must justify swapping the order of integration. The absolute value of the integrand is $|f(s)| e^{-\varepsilon \omega^2}$. The double integral of the absolute value is:
\begin{equation}
\begin{aligned}
    \int_\R \int_\R |f(s)| e^{-\varepsilon \omega^2} \dd s \dd \omega
    &= \|f\|_{L^1} \int_\R e^{-\varepsilon \omega^2} \dd \omega \\
    &= \|f\|_{L^1} \sqrt{\frac{\pi}{\varepsilon}} < \infty
    \end{aligned}
\end{equation}

Since the double integral is absolutely convergent, the \textit{Fubini--Tonelli Theorem} permits us to swap the integrals:
\begin{equation}
\begin{aligned}
I_\varepsilon(t)
    &= \int_\R f(s) \underbrace{ \left[ \frac{1}{2\pi} \int_\R e^{-\varepsilon \omega^2} e^{i\omega(t-s)} \dd \omega \right] }_{g_\varepsilon(t-s)} \dd s \\
    &= (f * g_\varepsilon)(t).
\end{aligned}
\end{equation}
Here $*$ denotes convolution. We now explicitly evaluate the inner integral $g_\varepsilon(u)$ by completing the square in the exponent:
\begin{equation}
\begin{aligned}
    -\varepsilon \omega^2 + i\omega u
    &= -\varepsilon \left( \omega^2 - \frac{i\omega u}{\varepsilon} \right) \\
    &= -\varepsilon \left( \left(\omega - \frac{iu}{2\varepsilon}\right)^2 + \frac{u^2}{4\varepsilon^2} \right) \\
    &= -\varepsilon \left(\omega - \frac{iu}{2\varepsilon}\right)^2 - \frac{u^2}{4\varepsilon}
    \end{aligned}
\end{equation}

Substituting $z = \omega - iu/(2\varepsilon)$ and shifting the contour of integration (justified by Cauchy's Integral Theorem, since $e^{-\varepsilon z^2}$ is entire and decays rapidly in horizontal strips):
\begin{equation}
\begin{aligned}
g_\varepsilon(u)
    &= \frac{1}{2\pi} e^{-u^2 / 4\varepsilon} \int_\R e^{-\varepsilon z^2} \dd z \\
    &= \frac{1}{2\pi} e^{-u^2 / 4\varepsilon} \sqrt{\frac{\pi}{\varepsilon}}
    = \frac{1}{\sqrt{4\pi\varepsilon}} e^{-u^2 / 4\varepsilon}
    \end{aligned}
\end{equation}

The family of functions $\{g_\varepsilon\}_{\varepsilon > 0}$ forms a \textit{nascent delta sequence}. We explicitly verify the three required conditions:
\begin{enumerate}
\item \textit{Non-negativity}: $g_\varepsilon(u) > 0$ for all $u \in \R$.
\item \textit{Normalization}: the following identity holds:

\[
\int_\R g_\varepsilon(u) \dd u = 1
\]

(which follows from evaluating the Gaussian integral or noting it is the inverse Fourier transform of $e^{-\varepsilon \omega^2}$ at $t=0$).
\item \textit{Concentration}: for any $\eta > 0$, let $v = u/(4\varepsilon)^{1/2}$. Then:

\begin{equation}
\int_{|u| > \eta} g_\varepsilon(u) \dd u = \frac{2}{\sqrt{\pi}} \int_{|v| > \eta/\sqrt{4\varepsilon}} e^{-v^2} \dd v
\end{equation}
As $\varepsilon \to 0^+$, the lower limit $\eta/(4\varepsilon)^{1/2} \to \infty$, so the integral approaches 0.
\end{enumerate}
By the properties of nascent delta sequences, the convolution $(f * g_\varepsilon)(t)$ converges to $f(t)$ at every point of continuity of $f$. Equating the two limits ($\varepsilon \to 0^+$) yields the inversion formula.
\end{proof}

\subsection{\texorpdfstring{$L^2$}{L2} Theory and Convolution}

\begin{theorem}[Plancherel Theorem]
The Fourier transform extends uniquely to a unitary operator (up to a scaling factor) on $L^2(\R)$. For all $f, g \in L^2(\R)$:
\begin{equation}
\int_\R f(t) \overline{g(t)} \dd t = \frac{1}{2\pi} \int_\R \hat{f}(\omega) \overline{\hat{g}(\omega)} \dd \omega
\end{equation}

\end{theorem}

\begin{proof}
We first prove the identity for $f, g \in \mathcal{S}(\R)$ (the Schwartz space), and then extend it to $L^2(\R)$.

\textit{Step 1: Setup for Schwartz functions.}
Define the time-reversed, complex-conjugated function $\tilde{g}(t) = \overline{g(-t)}$. Since $g \in \mathcal{S}(\R)$, it is clear that $\tilde{g} \in \mathcal{S}(\R)$. Let $h = f * \tilde{g}$. Because $\mathcal{S}(\R)$ is closed under convolution, $h \in \mathcal{S}(\R)$.

\textit{Step 2: Evaluate $h(0)$ in the time domain.}

\begin{equation}
h(0) = (f * \tilde{g})(0) = \int_\R f(t) \tilde{g}(-t) \dd t = \int_\R f(t) \overline{g(t)} \dd t = \langle f, g \rangle_{L^2}
\end{equation}

\textit{Step 3: Evaluate $\hat{h}(\omega)$ in the frequency domain.}
By the Convolution Theorem (proven below),
\[
\hat{h}(\omega) = \hat{f}(\omega) \hat{\tilde{g}}(\omega).
\]
We compute $\hat{\tilde{g}}(\omega)$:
\begin{equation}
\hat{\tilde{g}}(\omega) = \int_\R \overline{g(-t)} e^{-i\omega t} \dd t
\end{equation}
Substitute $s = -t$, so $\dd t = -\dd s$:
\begin{equation}
\begin{aligned}
    \hat{\tilde{g}}(\omega)
    &= \int_\R \overline{g(s)} e^{i\omega s} \dd s \\
    &= \overline{ \left( \int_\R g(s) e^{-i\omega s} \dd s \right) }
    = \overline{\hat{g}(\omega)}
    \end{aligned}
\end{equation}
Thus,
\[
\hat{h}(\omega) = \hat{f}(\omega) \overline{\hat{g}(\omega)}.
\]

\textit{Step 4: Apply the Inversion Theorem.}
Since $h \in \mathcal{S}(\R)$, both $h$ and $\hat{h}$ are in $L^1(\R)$. Applying the Fourier Inversion Theorem at $t=0$:
\begin{equation}
\begin{aligned}
h(0) &= \frac{1}{2\pi} \int_\R \hat{h}(\omega) e^{0} \dd \omega \\
    &= \frac{1}{2\pi} \int_\R \hat{f}(\omega) \overline{\hat{g}(\omega)} \dd \omega \\
    &= \frac{1}{2\pi} \langle \hat{f}, \hat{g} \rangle_{L^2}
    \end{aligned}
\end{equation}
Equating the results from Step 2 and Step 4 yields
\[
\langle f, g \rangle_{L^2} = (1/(2\pi)) \langle \hat{f}, \hat{g} \rangle_{L^2}.
\]

\textit{Step 5: Extension to $L^2(\R)$.}
The Schwartz space $\mathcal{S}(\R)$ is dense in $L^2(\R)$. On this dense subset, the Fourier transform is a linear isometry (up to the $1/(2\pi)^{1/2}$ factor). By the \textit{Bounded Linear Transformation (BLT) Theorem}, any uniformly continuous linear map from a dense subset of a metric space to a complete metric space extends uniquely to the entire space. Since $L^2(\R)$ is complete, the Fourier transform extends uniquely to a bounded linear operator on all of $L^2(\R)$. Because the isometry property holds on the dense subset, it holds on the entire space by the continuity of the norm.
\end{proof}

\begin{corollary}[Parseval's Identity]
Setting $f=g$ in Plancherel's theorem yields the conservation of energy:
\begin{equation}
\int_\R |f(t)|^2 \dd t = \frac{1}{2\pi} \int_\R |\hat{f}(\omega)|^2 \dd \omega
\end{equation}

\end{corollary}

\begin{theorem}[Convolution Theorem]
If $f, g \in L^1(\R)$, then their convolution $f*g \in L^1(\R)$ and the Fourier transform of the convolution is the pointwise product of their Fourier transforms:
\begin{equation}
\widehat{f*g}(\omega) = \hat{f}(\omega) \hat{g}(\omega)
\end{equation}

\end{theorem}

\begin{proof}
By definition,
\[
(f*g)(t) = \int_\R f(\tau) g(t-\tau) \dd \tau.
\]
We evaluate its Fourier transform:
\begin{equation}
\widehat{f*g}(\omega) = \int_\R e^{-i\omega t} \left( \int_\R f(\tau) g(t-\tau) \dd \tau \right) \dd t
\end{equation}
To swap the order of integration, we verify absolute integrability via Fubini--Tonelli:
\begin{equation}
\int_\R \int_\R |f(\tau)| |g(t-\tau)| \dd \tau \dd t = \int_\R |f(\tau)| \left( \int_\R |g(t-\tau)| \dd t \right) \dd \tau = \|f\|_{L^1} \|g\|_{L^1} < \infty
\end{equation}
Since the double integral is absolutely convergent, we swap the integrals:
\begin{equation}
\widehat{f*g}(\omega) = \int_\R f(\tau) \left( \int_\R g(t-\tau) e^{-i\omega t} \dd t \right) \dd \tau
\end{equation}

In the inner integral, we perform the substitution $u = t-\tau$, which implies $t = u+\tau$ and $\dd t = \dd u$. The limits of integration remain $-\infty$ to $\infty$:
\begin{equation}
\begin{aligned}
\int_\R g(u) e^{-i\omega (u+\tau)} \dd u &= e^{-i\omega \tau} \int_\R g(u) e^{-i\omega u} \dd u \\
&= e^{-i\omega \tau} \hat{g}(\omega)
\end{aligned}
\end{equation}
Substituting this back into the outer integral:
\begin{equation}
\begin{aligned}
    \widehat{f*g}(\omega)
    &= \int_\R f(\tau) e^{-i\omega \tau} \hat{g}(\omega) \dd \tau \\
    &= \hat{g}(\omega) \int_\R f(\tau) e^{-i\omega \tau} \dd \tau
    = \hat{g}(\omega) \hat{f}(\omega)
    \end{aligned}
\end{equation}

\end{proof}

Figure~\ref{fig:atlas-box-convolution} interprets convolution as the length of overlap between two moving intervals.

\begin{figure}[!htbp]
\centering
\begin{tikzpicture}[>=Stealth]
\begin{axis}[width=5.8cm,height=4.8cm,axis lines=middle,ticklabel style={font=\small},label style={font=\small},legend style={font=\scriptsize,draw=black!20,fill=white},xlabel={$s$},xmin=-0.6,xmax=1.3,ymin=0,ymax=1.4,xtick={-0.25,0,0.75,1},xticklabels={$t-1$,$0$,$t$,$1$},ytick=\empty]
\addplot[figblue,very thick] coordinates {(0,1)(1,1)};
\addplot[figred,very thick] coordinates {(-0.25,0.65)(0.75,0.65)};
\addplot[figgreen,line width=4pt] coordinates {(0,0.3)(0.75,0.3)};
\draw[black!40,dashed] (axis cs:0,0) -- (axis cs:0,1.2);
\draw[black!40,dashed] (axis cs:0.75,0) -- (axis cs:0.75,1.2);
\end{axis}
\begin{axis}[width=5.8cm,height=4.8cm,axis lines=middle,ticklabel style={font=\small},label style={font=\small},legend style={font=\scriptsize,draw=black!20,fill=white},at={(6.2cm,0)},xlabel={$t$},ylabel={$(f*f)(t)$},xmin=-0.3,xmax=2.3,ymin=-0.1,ymax=1.4,xtick={0,1,2},ytick={1}]
\addplot[figblue,very thick] coordinates {(-0.3,0)(0,0)(1,1)(2,0)(2.3,0)};
\addplot[only marks,figgreen,mark=*,mark size=2.5pt] coordinates {(0.75,0.75)};
\draw[figgreen,dashed] (axis cs:0.75,0) -- (axis cs:0.75,0.75);
\end{axis}
\end{tikzpicture}
\caption{Left: at shift $t=3/4$, the intervals $[0,1]$ and $[t-1,t]$ overlap on a segment of length $3/4$. Right: as $t$ varies, that overlap length forms the triangular convolution of two unit-interval indicator functions. It is zero outside $[0,2]$ and reaches one at $t=1$.}
\label{fig:atlas-box-convolution}
\end{figure}

\subsection{The Heisenberg Uncertainty Principle}


\begin{theorem}[Heisenberg Uncertainty Principle]
For $f \in L^2(\R)$ with $\|f\|_{L^2}=1$, assuming $tf(t)$ and $\omega \hat{f}(\omega)$ are in $L^2(\R)$:
\begin{equation}
\left( \int_\R t^2 |f(t)|^2 \dd t \right) \left( \frac{1}{2\pi} \int_\R \omega^2 |\hat{f}(\omega)|^2 \dd \omega \right) \geq \frac{1}{4}
\end{equation}

\end{theorem}

\begin{proof}
Define the time and frequency second moments by
\[
\begin{gathered}
\Delta t^2 = \int_\R t^2 |f(t)|^2 \dd t \\
\Delta \omega^2 = \frac{1}{2\pi} \int_\R \omega^2 |\hat{f}(\omega)|^2 \dd \omega.
\end{gathered}
\]

For the derivation, we assume $f \in \mathcal{S}(\R)$ to ensure all boundary terms vanish and derivatives are well-behaved; the result extends to the specified domain in $L^2$ by a standard density argument.

\textit{Step 1: Relate $\Delta \omega^2$ to the time-domain derivative.}
First, we establish the derivative property of the Fourier transform. Using integration by parts on $\widehat{f'}(\omega)$:
\begin{equation}
\begin{aligned}
    \widehat{f'}(\omega)
    &= \int_\R f'(t) e^{-i\omega t} \dd t \\
    &= \left[ f(t) e^{-i\omega t} \right]_{-\infty}^\infty
    + \int_\R f(t) (i\omega) e^{-i\omega t} \dd t
    \end{aligned}
\end{equation}
Because $f \in \mathcal{S}(\R)$, $f(t) \to 0$ as $|t| \to \infty$, so the boundary term vanishes. Thus,
\[
\widehat{f'}(\omega) = i\omega \hat{f}(\omega).
\]
Applying Plancherel's theorem to the function $f'$:
\begin{equation}
\begin{aligned}
    \int_\R |f'(t)|^2 \dd t
    &= \frac{1}{2\pi} \int_\R |\widehat{f'}(\omega)|^2 \dd \omega \\
    &= \frac{1}{2\pi} \int_\R |i\omega \hat{f}(\omega)|^2 \dd \omega \\
    &= \frac{1}{2\pi} \int_\R \omega^2 |\hat{f}(\omega)|^2 \dd \omega
    = \Delta \omega^2
    \end{aligned}
\end{equation}

\textit{Step 2: Integration by parts on the time second moment.}
We evaluate the integral of $t (d/dt) |f(t)|^2$. Note that by the product rule:
\begin{equation}
\begin{aligned}
\frac{d}{dt} |f(t)|^2 &= \frac{d}{dt} (f(t) \overline{f(t)}) \\
&= f'(t) \overline{f(t)} + f(t) \overline{f'(t)} \\
&= 2 \Re\left( f'(t) \overline{f(t)} \right)
\end{aligned}
\end{equation}
Integrating by parts over $\R$:
\begin{equation}
\int_\R t \frac{d}{dt} |f(t)|^2 \dd t = \left[ t |f(t)|^2 \right]_{-\infty}^\infty - \int_\R |f(t)|^2 \dd t
\end{equation}

Since $f \in \mathcal{S}(\R)$, the boundary term vanishes. The remaining integral is $-\|f\|_{L^2}^2 = -1$. Equating the two expressions for the integral:
\begin{equation}
\begin{aligned}
-1 &= \int_\R t \left( f'(t) \overline{f(t)} + f(t) \overline{f'(t)} \right) \dd t \\
&= 2 \Re \int_\R t f'(t) \overline{f(t)} \dd t
\end{aligned}
\end{equation}

\textit{Step 3: Apply the Cauchy--Schwarz inequality.}
Taking the absolute value of both sides:
\begin{equation}
\begin{aligned}
1 &= \left| 2 \Re \int_\R t f'(t) \overline{f(t)} \dd t \right| \\
&\leq 2 \left| \int_\R t f'(t) \overline{f(t)} \dd t \right| \\
&\leq 2 \int_\R |t f(t)| |f'(t)| \dd t
\end{aligned}
\end{equation}

We apply the Cauchy--Schwarz inequality to the functions $A(t) = t f(t)$ and $B(t) = f'(t)$ in $L^2(\R)$:
\begin{equation}
\begin{aligned}
\int_\R |A(t) B(t)| \dd t &\leq \|A\|_{L^2} \|B\|_{L^2} \\
&= \left( \int_\R t^2 |f(t)|^2 \dd t \right)^{1/2} \left( \int_\R |f'(t)|^2 \dd t \right)^{1/2} \\
&= \sqrt{\Delta t^2 \Delta \omega^2}
\end{aligned}
\end{equation}
Substituting this bound back into the inequality:
\begin{equation}
1 \leq 2 \sqrt{\Delta t^2 \Delta \omega^2}
\end{equation}
Squaring both sides and dividing by 4 yields the uncertainty principle: $\Delta t^2 \Delta \omega^2 \geq 1/4$.

\textit{Step 4: Determine the equality condition.}
Equality in the Cauchy--Schwarz inequality holds if and only if $B(t) = c A(t)$ for some constant $c \in \C$, meaning $f'(t) = c t f(t)$.
Furthermore, equality in the bound $\Re(z) \leq |z|$ requires $z$ to be a nonnegative real number. Here,
\[
\begin{aligned}
z &= \int t f' \bar{f} \\
&= c \int t^2 |f|^2 \\
&= c \Delta t^2.
\end{aligned}
\]

Since $1 = -2 \Re(z)$, $z$ must be real and strictly negative, which implies $c$ is a negative real number. Let $c = -\alpha$ with $\alpha > 0$.
Solving the differential equation
\[
f'(t) = -\alpha t f(t):
\]

\begin{equation}
\frac{f'(t)}{f(t)} = -\alpha t \implies \ln f(t) = -\frac{\alpha}{2} t^2 + C \implies f(t) = A e^{-\frac{\alpha}{2} t^2}
\end{equation}
This explicitly proves that the lower bound is achieved if and only if $f$ is a Gaussian function.
\end{proof}

Figure~\ref{fig:ch14-uncertainty} illustrates the trade-off between localization in time and in frequency.

\begin{figure}[!htbp]
\centering
\begin{tikzpicture}
\begin{axis}[
  width=6.5cm,height=5.8cm, axis lines=middle,
  xlabel={$t$}, ylabel={$f(t)$},
  xmin=-3,xmax=3,ymin=0,ymax=1.3, xtick=\empty,ytick=\empty,
  ]
  \addplot[figblue,very thick,domain=-3:3,samples=150] {exp(-4*x^2)};
\end{axis}
\begin{axis}[
  at={(7.2cm,0)}, width=6.5cm,height=5.8cm, axis lines=middle,
  xlabel={$\xi$}, ylabel={$\hat f(\xi)$},
  xmin=-3,xmax=3,ymin=0,ymax=1.3, xtick=\empty,ytick=\empty,
  ]
  \addplot[figred,very thick,domain=-3:3,samples=150] {exp(-0.4*x^2)};
\end{axis}
\begin{axis}[
  at={(0,-5.2cm)}, width=6.5cm,height=5.8cm, axis lines=middle,
  xlabel={$t$}, ylabel={$f(t)$},
  xmin=-3,xmax=3,ymin=0,ymax=1.3, xtick=\empty,ytick=\empty,
  ]
  \addplot[figgreen,very thick,domain=-3:3,samples=150] {exp(-0.4*x^2)};
\end{axis}
\begin{axis}[
  at={(7.2cm,-5.2cm)}, width=6.5cm,height=5.8cm, axis lines=middle,
  xlabel={$\xi$}, ylabel={$\hat f(\xi)$},
  xmin=-3,xmax=3,ymin=0,ymax=1.3, xtick=\empty,ytick=\empty,
  ]
  \addplot[figblue,very thick,domain=-3:3,samples=150] {exp(-4*x^2)};
\end{axis}
\end{tikzpicture}
\caption[Uncertainty principle: time vs frequency]{Schematic time--frequency localization using Gaussian profiles. Left column: signals in time. Right column: frequency profiles. A narrow signal corresponds to a broad frequency profile (top row), while a broad signal corresponds to a narrow frequency profile (bottom row), illustrating the uncertainty bound discussed in the text.}
\label{fig:ch14-uncertainty}
\end{figure}

\section{Holomorphic Functions, Singularities, and Residues}
\label{sec:complex-analysis}

Complex differentiation and contour integration complement the preceding Fourier methods. They also connect the planar integral theorems to harmonic functions and provide tools for evaluating real integrals and Fourier transforms.

\subsection{Holomorphic Functions and the Cauchy--Riemann Equations}

\begin{definition}[Complex Differentiability and Holomorphy]
Let $U\subset\C$ be open. A function $f:U\to\C$ is \emph{complex differentiable} at $z_0\in U$ if the limit \[
f'(z_0):=\lim_{h\to0}\frac{f(z_0+h)-f(z_0)}{h}
\] exists independently of how the complex increment $h$ approaches zero. The function is \emph{holomorphic} on $U$ if it is complex differentiable at every point of $U$; it is holomorphic near $z_0$ if this holds in a neighborhood of $z_0$.
\end{definition}

This condition is stronger than real differentiability of a map $\R^2\to\R^2$. Writing
\[
\begin{gathered}
z=x+iy \\
f(z)=u(x,y)+iv(x,y).
\end{gathered}
\]
The difference quotient must have the same limit along the real and imaginary directions. This imposes relations between the partial derivatives of $u$ and $v$.

\begin{theorem}[Cauchy--Riemann Equations]
If $f=u+iv$ is complex differentiable at $z_0=x_0+iy_0$, then the partial derivatives of $u,v$ exist at $(x_0,y_0)$ and satisfy
\begin{equation}
\frac{\partial u}{\partial x}=\frac{\partial v}{\partial y},
\qquad
\frac{\partial u}{\partial y}=-\frac{\partial v}{\partial x}.
\end{equation}

\end{theorem}

\begin{proof}
\textit{Step 1: approach along the real direction.} Taking $h=\Delta x\in\R$ gives

\begin{align}
f'(z_0)
&=\lim_{\Delta x\to0}\left[
\frac{u(x_0+\Delta x,y_0)-u(x_0,y_0)}{\Delta x}\right.\notag\\
&\hspace{35mm}\left.+i\frac{v(x_0+\Delta x,y_0)-v(x_0,y_0)}{\Delta x}\right]\notag\\
&=u_x(x_0,y_0)+iv_x(x_0,y_0).
\end{align}

\textit{Step 2: approach along the imaginary direction.} Taking $h=i\Delta y$ with $\Delta y\in\R$ gives

\begin{align}
f'(z_0)
&=\lim_{\Delta y\to0}\left[
\frac{u(x_0,y_0+\Delta y)-u(x_0,y_0)}{i\Delta y}\right.\notag\\
&\hspace{35mm}\left.+\frac{v(x_0,y_0+\Delta y)-v(x_0,y_0)}{\Delta y}\right]\notag\\
&=-iu_y(x_0,y_0)+v_y(x_0,y_0),
\end{align}
because $1/i=-i$.

\textit{Step 3: compare the two limits.} Independence of the approach direction implies

\[
u_x+iv_x=v_y-iu_y.
\]
Equality of the real and imaginary parts gives $u_x=v_y$ and $v_x=-u_y$, as required.
\end{proof}

\begin{remark}[Sufficiency and Harmonic Functions]
If $u,v$ have continuous first partial derivatives on an open set and satisfy the Cauchy--Riemann equations there, then $f=u+iv$ is holomorphic. Moreover, the real and imaginary parts of a holomorphic function are harmonic. To see the differential mechanism, assuming $u,v\in C^2$, differentiate $u_x=v_y$ in $x$ and $u_y=-v_x$ in $y$:
\[
\begin{aligned}
\Delta u &= u_{xx}+u_{yy} \\
&= v_{yx}-v_{xy} \\
&= 0.
\end{aligned}
\]

Equality of mixed partial derivatives gives the cancellation; the argument for $\Delta v=0$ is analogous. Holomorphic functions in fact possess the required higher regularity, a result not proved here. This connects complex analysis to the Laplace equation and to the Dirichlet and Neumann problems discussed in Section~\ref{sec:pde-boundary-problems}.
\end{remark}

\subsection{Singularities, Residues, and the Residue Theorem}

\begin{definition}[Isolated Singularity and Residue]
An \emph{isolated singularity} of $f$ is a point $z_0$ around which $f$ is holomorphic on a punctured disk $0<|z-z_0|<r$, while its value at $z_0$ may be undefined or fail to give a holomorphic extension. On this punctured disk, $f$ has a \emph{Laurent expansion}
\begin{equation}
f(z)=\sum_{n=-\infty}^{\infty}c_n(z-z_0)^n.
\end{equation}
The coefficient of $(z-z_0)^{-1}$ is the \emph{residue}, denoted by
\[
\operatorname{Res}(f,z_0):=c_{-1}.
\]
The existence of the Laurent expansion is a standard result whose proof we omit.
\end{definition}

\begin{theorem}[Residue Theorem]
\label{thm:residue}
Let $\gamma$ be a positively oriented piecewise smooth simple closed curve. Suppose that $f$ is holomorphic on an open neighborhood of $\gamma$ and its interior, except for finitely many isolated singularities $z_1,\dots,z_k$ strictly inside $\gamma$. Then
\begin{equation}
\oint_\gamma f(z)\,\dd z
=2\pi i\sum_{j=1}^k\operatorname{Res}(f,z_j).
\end{equation}

\end{theorem}

\begin{proof}[Proof Sketch]
The argument rests on two facts. First, Cauchy's integral theorem states that the integral around a simple closed curve vanishes if $f$ is holomorphic on a neighborhood of the curve and its entire interior. Under $C^1$ regularity, Green's theorem (Theorem~\ref{thm:green-plane}) shows this directly: writing $f(z)\,\dd z=(u\,\dd x-v\,\dd y)+i(v\,\dd x+u\,\dd y)$, the two area integrands are $-v_x-u_y$ and $u_x-v_y$, both zero by the Cauchy--Riemann equations. Second, parametrizing a positively oriented circle by $z-z_0=\varepsilon e^{i\theta}$ gives
\begin{equation}
\oint_{|z-z_0|=\varepsilon}(z-z_0)^n\,\dd z
=\begin{cases}2\pi i,&n=-1,\\0,&n\in\Z,\ n\ne-1.\end{cases}
\end{equation}

Remove disjoint small disks around the singularities and apply Cauchy's theorem to the remaining region. The integral along $\gamma$ equals the sum of the counterclockwise integrals around the small circles. Integrating each Laurent expansion term by term selects only its coefficient $c_{-1}$, yielding the stated formula.
\end{proof}

\begin{example}[Evaluating a Real Integral by Residues]
Consider
\[
\int_{-\infty}^{\infty}\frac{\dd x}{1+x^2}.
\]

Its value is already accessible through the arctangent, but it illustrates a method that also applies when elementary antiderivatives are unavailable. The complex function $f(z)=1/(1+z^2)$ has simple poles at $z=\pm i$. For $R>1$, let $\gamma_R$ consist of the segment $[-R,R]$ followed by the upper semicircle
\[
C_R = \{Re^{i\theta}:0 \leq \theta \leq \pi\}.
\]
Only the pole at $i$ is enclosed, and
\begin{equation}
\operatorname{Res}(f,i)
=\lim_{z\to i}(z-i)f(z)
=\left.\frac{1}{z+i}\right|_{z=i}=\frac{1}{2i}.
\end{equation}
The residue theorem gives
\[
\oint_{\gamma_R}f(z)\,\dd z=\pi.
\]
On the semicircle, $|1+z^2|\geq R^2-1$, so
\begin{equation}
\left|\int_{C_R}f(z)\,\dd z\right|
\leq\frac{\pi R}{R^2-1}\longrightarrow0.
\end{equation}
Letting $R\to\infty$ therefore yields
\begin{equation}
\int_{-\infty}^{\infty}\frac{\dd x}{1+x^2}=\pi.
\end{equation}

Closing a real contour with an arc whose integral vanishes is also useful for explicit Fourier-transform calculations. For integrals containing $f(x)e^{i\omega x}$, the choice of half-plane must account for the exponential decay as well as the poles of $f$.
\end{example}

Figure~\ref{fig:atlas-residue-contour} identifies the enclosed pole and contour orientation in the residue calculation.

\begin{figure}[!htbp]
\centering
\begin{tikzpicture}[scale=1.15,>=Stealth]
\draw[->,black!40] (-3,0) -- (3.2,0) node[right] {$\mathrm{Re}$};
\draw[->,black!40] (0,-1.5) -- (0,3) node[above] {$\mathrm{Im}$};
\draw[figblue,very thick] (-2.5,0) -- (2.5,0) arc[start angle=0,end angle=180,radius=2.5];
\draw[figblue,very thick,->] (-1.8,0) -- (-0.8,0);
\draw[figblue,very thick,->] (35:2.5) arc[start angle=35,end angle=65,radius=2.5];
\node[below] at (-2.5,0) {$-R$};
\node[below] at (2.5,0) {$R$};
\node[above right,figblue] at (45:2.5) {$C_R$};
\node[figred] at (0,1) {$\times$};
\node[right,figred] at (0.1,1) {$i$};
\node[black!65] at (0,-1) {$\times$};
\node[right] at (0.1,-1) {$-i$};
\end{tikzpicture}
\caption{For $1/(1+z^2)$, the real segment from $-R$ to $R$ is closed by an upper semicircle traversed counterclockwise. Only the pole at $i$ lies inside; the pole at $-i$ lies outside. For $R>1$, the arc contribution tends to zero as $R$ grows, leaving the real-line integral.}
\label{fig:atlas-residue-contour}
\end{figure}
\newpage
\chapter{Differential Equations and Dynamical Systems}

\section{First-Order Differential Equations and Initial Value Problems}

\subsection{Differential Equations and Families of Solutions}

An ordinary differential equation of first order is an equation that involves an unknown function of a single variable, which we shall generally denote by $y = y(t)$, together with its first derivative. The general form of such an equation is
\[
y' = F(t,y),
\]
where $F(t,y)$ is a given expression in the two variables $t$ and $y$. Several important classes of first-order equations arise as special cases of this general form, according to the specific dependence of $F$ on $t$ and $y$. A pure-time equation is one of the form $y' = f(t)$, in which the right-hand side depends only on the independent variable. A separable equation is one of the form $y' = a(t) g(y)$, in which the right-hand side factors as a product of a function of $t$ alone and a function of $y$ alone. An autonomous equation is the special case of a separable equation in which the dependence on $t$ disappears entirely, so that $y' = g(y)$. Finally, a linear equation is one of the form $y' = a(t) y + f(t)$, in which the right-hand side is affine in the unknown $y$.

\begin{definition}
A function $y = y(t)$ is called a solution of the differential equation $y' = F(t,y)$ on the interval $I$ if the equation is satisfied identically, that is, if
\[
y'(t) = F(t, y(t))
\]
for every $t \in I$.
\end{definition}

As a simple illustration of this definition, one may verify directly that $y(t) = 10 e^{2t}$ is a solution of the equation $y' = 2y$, since differentiating this expression indeed yields
\[
\begin{aligned}
y'(t) &= 20 e^{2t} \\
&= 2 y(t).
\end{aligned}
\]

When asked to solve a differential equation, one is generally expected to determine all possible solutions of the equation, rather than a single particular solution; this immediately raises the natural question of how many solutions a first-order equation may possess.


To address this question, it is instructive to consider the simplest possible class of first-order equations, namely the pure-time equations $y' = f(t)$. For an arbitrary constant $C \in \mathbb{R}$, the solutions are given by
\[
\begin{aligned}
y(t) &= \int f(t) \, dt \\
&= g(t) + C,
\end{aligned}
\]
where $g$ denotes any fixed antiderivative of $f$; the equation therefore possesses infinitely many solutions, parametrized by the arbitrary constant $C$. This phenomenon is not peculiar to pure-time equations, but is in fact a general feature of first-order ordinary differential equations: such equations typically possess infinitely many solutions, depending on a single free parameter, which may be fixed by imposing an additional condition, as will be discussed below in connection with initial value problems.

Figure~\ref{fig:atlas-ode-direction-field} relates a differential equation to its direction field and family of solutions.

\begin{figure}[!htbp]
\centering
\begin{tikzpicture}[>=Stealth]
\begin{axis}[
  width=12.5cm,height=6.5cm,
  axis lines=left,
  ticklabel style={font=\small},label style={font=\small},
  xlabel={$t$},ylabel={$y$},
  xmin=0,xmax=2.1,ymin=0,ymax=3.3,
  xtick={0,0.5,1,1.5,2},ytick={0,1,2,3},
  legend style={at={(0.5,1.02)},anchor=south,legend columns=3,
    font=\scriptsize,draw=black!20,fill=white,column sep=1em}
]
\draw[black!25] (axis cs:-0.07500,0.00000) -- (axis cs:0.07500,0.00000);
\draw[black!25] (axis cs:-0.06708,0.46646) -- (axis cs:0.06708,0.53354);
\draw[black!25] (axis cs:-0.05303,0.94697) -- (axis cs:0.05303,1.05303);
\draw[black!25] (axis cs:-0.04160,1.43760) -- (axis cs:0.04160,1.56240);
\draw[black!25] (axis cs:-0.03354,1.93292) -- (axis cs:0.03354,2.06708);
\draw[black!25] (axis cs:-0.02785,2.43036) -- (axis cs:0.02785,2.56964);
\draw[black!25] (axis cs:-0.02372,2.92885) -- (axis cs:0.02372,3.07115);
\draw[black!25] (axis cs:0.17500,0.00000) -- (axis cs:0.32500,0.00000);
\draw[black!25] (axis cs:0.18292,0.46646) -- (axis cs:0.31708,0.53354);
\draw[black!25] (axis cs:0.19697,0.94697) -- (axis cs:0.30303,1.05303);
\draw[black!25] (axis cs:0.20840,1.43760) -- (axis cs:0.29160,1.56240);
\draw[black!25] (axis cs:0.21646,1.93292) -- (axis cs:0.28354,2.06708);
\draw[black!25] (axis cs:0.22215,2.43036) -- (axis cs:0.27785,2.56964);
\draw[black!25] (axis cs:0.22628,2.92885) -- (axis cs:0.27372,3.07115);
\draw[black!25] (axis cs:0.42500,0.00000) -- (axis cs:0.57500,0.00000);
\draw[black!25] (axis cs:0.43292,0.46646) -- (axis cs:0.56708,0.53354);
\draw[black!25] (axis cs:0.44697,0.94697) -- (axis cs:0.55303,1.05303);
\draw[black!25] (axis cs:0.45840,1.43760) -- (axis cs:0.54160,1.56240);
\draw[black!25] (axis cs:0.46646,1.93292) -- (axis cs:0.53354,2.06708);
\draw[black!25] (axis cs:0.47215,2.43036) -- (axis cs:0.52785,2.56964);
\draw[black!25] (axis cs:0.47628,2.92885) -- (axis cs:0.52372,3.07115);
\draw[black!25] (axis cs:0.67500,0.00000) -- (axis cs:0.82500,0.00000);
\draw[black!25] (axis cs:0.68292,0.46646) -- (axis cs:0.81708,0.53354);
\draw[black!25] (axis cs:0.69697,0.94697) -- (axis cs:0.80303,1.05303);
\draw[black!25] (axis cs:0.70840,1.43760) -- (axis cs:0.79160,1.56240);
\draw[black!25] (axis cs:0.71646,1.93292) -- (axis cs:0.78354,2.06708);
\draw[black!25] (axis cs:0.72215,2.43036) -- (axis cs:0.77785,2.56964);
\draw[black!25] (axis cs:0.72628,2.92885) -- (axis cs:0.77372,3.07115);
\draw[black!25] (axis cs:0.92500,0.00000) -- (axis cs:1.07500,0.00000);
\draw[black!25] (axis cs:0.93292,0.46646) -- (axis cs:1.06708,0.53354);
\draw[black!25] (axis cs:0.94697,0.94697) -- (axis cs:1.05303,1.05303);
\draw[black!25] (axis cs:0.95840,1.43760) -- (axis cs:1.04160,1.56240);
\draw[black!25] (axis cs:0.96646,1.93292) -- (axis cs:1.03354,2.06708);
\draw[black!25] (axis cs:0.97215,2.43036) -- (axis cs:1.02785,2.56964);
\draw[black!25] (axis cs:0.97628,2.92885) -- (axis cs:1.02372,3.07115);
\draw[black!25] (axis cs:1.17500,0.00000) -- (axis cs:1.32500,0.00000);
\draw[black!25] (axis cs:1.18292,0.46646) -- (axis cs:1.31708,0.53354);
\draw[black!25] (axis cs:1.19697,0.94697) -- (axis cs:1.30303,1.05303);
\draw[black!25] (axis cs:1.20840,1.43760) -- (axis cs:1.29160,1.56240);
\draw[black!25] (axis cs:1.21646,1.93292) -- (axis cs:1.28354,2.06708);
\draw[black!25] (axis cs:1.22215,2.43036) -- (axis cs:1.27785,2.56964);
\draw[black!25] (axis cs:1.22628,2.92885) -- (axis cs:1.27372,3.07115);
\draw[black!25] (axis cs:1.42500,0.00000) -- (axis cs:1.57500,0.00000);
\draw[black!25] (axis cs:1.43292,0.46646) -- (axis cs:1.56708,0.53354);
\draw[black!25] (axis cs:1.44697,0.94697) -- (axis cs:1.55303,1.05303);
\draw[black!25] (axis cs:1.45840,1.43760) -- (axis cs:1.54160,1.56240);
\draw[black!25] (axis cs:1.46646,1.93292) -- (axis cs:1.53354,2.06708);
\draw[black!25] (axis cs:1.47215,2.43036) -- (axis cs:1.52785,2.56964);
\draw[black!25] (axis cs:1.47628,2.92885) -- (axis cs:1.52372,3.07115);
\draw[black!25] (axis cs:1.67500,0.00000) -- (axis cs:1.82500,0.00000);
\draw[black!25] (axis cs:1.68292,0.46646) -- (axis cs:1.81708,0.53354);
\draw[black!25] (axis cs:1.69697,0.94697) -- (axis cs:1.80303,1.05303);
\draw[black!25] (axis cs:1.70840,1.43760) -- (axis cs:1.79160,1.56240);
\draw[black!25] (axis cs:1.71646,1.93292) -- (axis cs:1.78354,2.06708);
\draw[black!25] (axis cs:1.72215,2.43036) -- (axis cs:1.77785,2.56964);
\draw[black!25] (axis cs:1.72628,2.92885) -- (axis cs:1.77372,3.07115);
\draw[black!25] (axis cs:1.92500,0.00000) -- (axis cs:2.07500,0.00000);
\draw[black!25] (axis cs:1.93292,0.46646) -- (axis cs:2.06708,0.53354);
\draw[black!25] (axis cs:1.94697,0.94697) -- (axis cs:2.05303,1.05303);
\draw[black!25] (axis cs:1.95840,1.43760) -- (axis cs:2.04160,1.56240);
\draw[black!25] (axis cs:1.96646,1.93292) -- (axis cs:2.03354,2.06708);
\draw[black!25] (axis cs:1.97215,2.43036) -- (axis cs:2.02785,2.56964);
\draw[black!25] (axis cs:1.97628,2.92885) -- (axis cs:2.02372,3.07115);
\addplot[figblue,very thick,domain=0:2,samples=100] {0.2*exp(x)};
\addlegendentry{$C=0.2$}
\addplot[figred,thick,domain=0:2,samples=100] {0.3*exp(x)};
\addlegendentry{$C=0.3$}
\addplot[figgreen,thick,domain=0:2,samples=100] {0.4*exp(x)};
\addlegendentry{$C=0.4$}
\end{axis}
\end{tikzpicture}
\caption{Short segments indicate the slope $y$ prescribed by $y^{\prime}=y$. The curves are exact solutions $y(t)=Ce^t$ for three positive initial values. Each initial value selects one curve; the direction field itself represents all such solutions, not a single trajectory.}
\label{fig:atlas-ode-direction-field}
\end{figure}

\subsection{Techniques for Solving First-Order Equations}

Among the general theory of first-order equations, explicit techniques are available for finding a formula for the general solution, also called the general integral, of several important classes: pure-time equations $y' = f(t)$, separable equations $y' = a(t) g(y)$, including the autonomous case, and linear equations
\[
y' = a(t) y + f(t).
\]

We now illustrate the method applicable to separable equations through two worked examples, before stating the general procedure.

\begin{example}
Consider the equation $y' = \varepsilon y$, where $\varepsilon$ is a given constant. One immediately identifies the constant solution $y(t) \equiv 0$. Setting $y \neq 0$ and dividing by $y$, the equation becomes
\[
\frac{y'(t)}{y(t)} = \varepsilon \implies \int \frac{y'(t)}{y(t)} \, dt = \int \varepsilon \, dt.
\]

The integral on the left-hand side is elementary, yielding $\ln |y(t)| = \varepsilon t + c$ for an arbitrary constant $c$; setting $C = e^c$, we obtain $|y(t)| = C e^{\varepsilon t}$. This furnishes a branch of positive solutions $y(t) = C e^{\varepsilon t}$, for $C > 0$, together with a branch of negative solutions $y(t) = -C e^{\varepsilon t}$, again for $C > 0$; together with the constant solution $y \equiv 0$ found initially, these exhaust all solutions of the equation.
\end{example}

\begin{example}
Consider the equation $y' = t y^3$. As before, one immediately identifies the constant solution $y(t) \equiv 0$. Setting $y \neq 0$ and separating the variables by dividing by $y^3$, the equation becomes
\[
\frac{y'(t)}{y^3(t)} = t \implies \int \frac{y'(t)}{y^3(t)} \, dt = \int t \, dt.
\]
The integral on the left-hand side is again elementary, yielding
\[
-(1/(2 y^2(t))) = (t^2/2) + c;
\]
setting $C = -2c$, we obtain two branches of solutions,
\[
y(t) = \frac{1}{\sqrt{C - t^2}}, \qquad y(t) = -\frac{1}{\sqrt{C - t^2}},
\]
each defined on the maximal interval $(-\sqrt{C}, \sqrt{C})$, for every fixed $C > 0$.
\end{example}

These two examples illustrate the general method applicable to the broader class of separable equations, defined as follows.

\begin{definition}[Separable Equation]
An equation of the form $y' = a(t) g(y)$ is called a separable equation. The special case in which $a(t)$ is constant, so that the equation takes the form $y' = g(y)$, is called an autonomous equation.
\end{definition}

Assuming that $a$ and $g$ are continuous functions, the general solution of a separable equation $y' = a(t) g(y)$ can be obtained by means of a systematic four-step procedure. The first step consists of searching for equilibrium solutions, if any exist, namely constant solutions of the form $y(t) \equiv r \in \mathbb{R}$, which correspond precisely to the roots of the equation $g(r) = 0$. The second step consists of separating the two variables $t$ and $y$: assuming that $g(y) \neq 0$, one divides the equation by $g(y)$ and integrates both sides with respect to $t$, obtaining
\[
\int \frac{y'(t)}{g(y(t))} \, dt = \int a(t) \, dt.
\]

The third step consists of evaluating both integrals explicitly: denoting by $A(t)$ a given antiderivative of $a(t)$, and by $H(u)$ an antiderivative of $1/g(u)$, one obtains the implicit relation
\[
H(y(t)) = A(t) + C,
\]
for an arbitrary constant $C \in \mathbb{R}$. The fourth and final step consists of inverting this implicit relation, whenever $H$ is invertible, to obtain the explicit formula
\[
y(t) = H^{-1}(A(t) + C);
\]
it should be noted that restrictions on the domain of $t$, possibly depending on the value of $C$, may arise at this final step, as illustrated in the two examples above.

\subsection{Applications to Population Dynamics}


One of the most important applications of the theory of first-order differential equations lies in biology, where a wide variety of processes can be modeled by such equations, providing considerable insight into the dynamics of living organisms, that is, into how individuals and populations evolve over time. The earliest and simplest such model, proposed by Malthus in 1830, describes the growth of a population under idealized conditions.

Let $N(t)$ denote the number of individuals in a population at time $t$. The Malthus model rests on three assumptions: each individual produces offspring at a constant rate $\lambda$, so that the total birth rate at time $t$ equals $\lambda N(t)$; and individuals die at a constant per-capita rate $\mu$, so that the total loss rate through death at time $t$ equals $\mu N(t)$. Since the rate of change of the population size equals the total birth rate minus the total death rate, with $\varepsilon = \lambda - \mu$, one obtains the following equivalent forms of the differential equation:
\[
\begin{aligned}
N'(t) &= \lambda N(t) - \mu N(t), \\
\frac{dN}{dt} &= \varepsilon N.
\end{aligned}
\]

The constant $\varepsilon$ is called the biological potential, or Malthus parameter, of the population. This linear autonomous equation can be solved explicitly by the separation of variables technique developed above, yielding exponential growth or decay depending on the sign of $\varepsilon$. A comparison with experimental data on the growth of a population of yeast cells, for instance, shows reasonably good agreement between the model $dN/dt = 0.5\, N$, with initial condition $N(0) = 0.2 \times 10^6$ cells per milliliter, and the observed population counts, at least during the early stages of growth.


The Malthus model, while adequate for describing the early stages of population growth, fails to capture an important biological phenomenon known as the logistic effect, namely the intraspecific competition for limited resources: when the population density is high, one typically observes a decrease in the birth rate together with an increase in the mortality rate, effects that the constant biological potential of the Malthus model cannot represent. The logistic, or Verhulst, model, proposed in 1838, addresses this limitation by allowing the biological potential to depend on the population size itself.

Specifically, instead of assuming a constant biological potential, the logistic model assumes that the biological potential decreases linearly as the population size increases, starting from a value $\varepsilon > 0$ when $N = 0$ and reaching the value $0$ when $N = k$, where the constant $k$ is called the carrying capacity of the environment. This linear dependence is expressed by the formula
\[
\varepsilon(N) = \varepsilon \left( 1 - \frac{N}{k} \right),
\]
and substituting this expression for the biological potential into the growth equation yields the logistic equation
\[
\frac{dN}{dt} = \varepsilon \left( 1 - \frac{N}{k} \right) N.
\]

As with the Malthus model, a comparison with experimental data on the growth of a yeast population shows that the logistic model
\[
dN/dt = (0.55 - 0.0026\, N)\, N,
\]
again with initial condition $N(0) = 0.2 \times 10^6$ cells per milliliter, reproduces the observed data considerably more accurately than the simple Malthus model, particularly as the population approaches its carrying capacity.

We now solve two representative logistic equations explicitly. For
\[
y'=y(2-y),
\]
the equilibrium solutions are $y\equiv0$ and $y\equiv2$. Away from
them, partial fractions give
\[
\frac1{y(2-y)}=\frac1{2y}+\frac1{2(2-y)},
\]
and therefore
\[
\frac12\ln|y|-\frac12\ln|2-y|=t+C.
\]
Solving for $y$ yields
\[
y(t)=\frac{2}{1+Ce^{-2t}},
\]
on every interval where the denominator is nonzero.

For
\[
y'=y(y-2),
\]
the equilibria are again $0$ and $2$, while
\[
\frac1{y(y-2)}=-\frac1{2y}+\frac1{2(y-2)}
\]
gives
\[
\ln\left|\frac{y-2}{y}\right|=2t+C,
\qquad
y(t)=\frac{2}{1-Ce^{2t}}.
\]
If $C\leq0$, the denominator never vanishes and the solution is global.
If $C>0$, it vanishes at $t=(1/2)\ln(1/C)$, where the solution blows
up. Figure~\ref{fig:ch7-logistic} connects the logistic trajectory with its phase-line description.

\begin{figure}[!htbp]
\centering
\begin{tikzpicture}
\begin{axis}[
  width=6.2cm,height=6.2cm, axis lines=left,
  xlabel={$t$}, ylabel={$S(t)$},
  xmin=0,xmax=10, ymin=0,ymax=5.6,
  ytick={5},yticklabels={$K$}, xtick=\empty,
  ]
  \addplot[figblue,very thick,domain=0:10,samples=150]{5/(1+9*exp(-0.9*x))};
  \draw[figred,dashed] (axis cs:0,5) -- (axis cs:10,5);
\end{axis}
\begin{axis}[
  at={(6.9cm,0)},
  width=6.2cm,height=6.2cm, axis lines=middle,
  xlabel={$S$}, ylabel={$\dot S$},
  xmin=0,xmax=5.6, ymin=-1,ymax=1.6,
  xtick={5},xticklabels={$K$}, ytick=\empty,
  ]
  \addplot[figgreen,very thick,domain=0:5.6,samples=100]{0.5*x*(1-x/5)};
  \node[circle,fill=white,draw=black,inner sep=1.3pt] at (axis cs:0,0) {};
  \node[below right,fill=white,inner sep=1pt,font=\scriptsize]
    at (axis cs:0.08,-0.12) {unstable};
  \node[circle,fill=black,inner sep=1.3pt] at (axis cs:5,0) {};
  \node[below left,fill=white,inner sep=1pt,font=\scriptsize]
    at (axis cs:4.95,-0.12) {stable};
  \draw[vecA,shorten >=1pt] (axis cs:1.0,0) -- (axis cs:2.2,0);
  \draw[vecA,shorten >=1pt] (axis cs:3.0,0) -- (axis cs:4.2,0);
  \draw[vecA,shorten >=1pt] (axis cs:5.55,0) -- (axis cs:5.15,0);
\end{axis}
\end{tikzpicture}
\caption[Logistic model: S-curve and phase line]{Left: solution $S(t)$ approaching
  the carrying capacity $K$ asymptotically. Right: phase portrait $\dot S$ vs $S$, with
  blue arrows indicating the flow direction; $S=0$ is unstable, $S=K$ is stable.}
\label{fig:ch7-logistic}
\end{figure}

\subsection{The Initial Value Problem}

In many applications arising from physical or biological problems, one is typically not interested in the entire family of solutions of a differential equation, but rather in the single, particular solution that satisfies an additional condition of the form $y(t_0) = y_0$ for prescribed values $t_0$ and $y_0$. Such a condition is called an initial condition, and the problem of finding the solution of the differential equation that satisfies a given initial condition is called an initial value problem, or Cauchy problem.

\begin{definition}[Cauchy Problem]
A Cauchy problem is a system consisting of a differential equation together with an initial condition, of the form
\[
\begin{cases}
y' = F(t,y), \\
y(t_0) = y_0,
\end{cases}
\]
where $t_0$ and $y_0$ are assigned values.
\end{definition}

A central theoretical question concerns whether a given Cauchy problem admits a solution, and if so, whether that solution is unique; a problem for which both properties hold is said to be well-posed.

\begin{definition}[Well-Posedness]
The Cauchy problem with initial datum $(t_0, y_0)$ is said to be well-posed if it admits a unique solution $y(\cdot)$ defined in a neighborhood of $t_0$.
\end{definition}

Well-posedness has an important geometric consequence: it implies that exactly one solution curve of the differential equation passes through the point $(t_0, y_0)$ in the $(t,y)$-plane, and consequently that two distinct solution curves can never intersect at a common point. In particular, if the Cauchy problem is well-posed for every admissible choice of initial condition, then a constant solution $y \equiv k$ can never be crossed by any other solution curve; the horizontal line $y = k$ therefore constitutes an impassable barrier for every other solution curve of the equation, a fact of considerable qualitative importance when sketching the family of solutions of a given equation.

For separable equations, well-posedness of the associated Cauchy problem can be guaranteed under mild regularity assumptions on the functions $a$ and $g$, though only in a local sense, as the following theorem makes precise.

\begin{theorem}[Local Well-Posedness for Separable Equations]
Let $a(t)$ be continuous on an interval $I = (\alpha, \beta)$, and let $g(y)$ be differentiable with $g'(y)$ continuous on an interval $J = (\gamma, \delta)$. Then, for every choice of $t_0 \in (\alpha,\beta)$ and every choice of initial datum $y_0 \in (\gamma,\delta)$, the Cauchy problem
\[
\begin{cases}
y' = a(t) g(y), \\
y(t_0) = y_0,
\end{cases}
\]
has a unique solution $y(t)$, defined on an interval $(t_0 - \delta_0, t_0 + \delta_0)$ for some $\delta_0 > 0$.
\end{theorem}

\begin{remark}
It is important to emphasize that the size of the interval of existence $\delta_0$ guaranteed by this theorem cannot, in general, be predicted theoretically from the hypotheses alone; the theorem asserts only the local existence and uniqueness of the solution, without providing quantitative information on how far this solution can be extended.
\end{remark}

The local nature of the well-posedness theorem stated above is not a mere technical artifact: it reflects genuine phenomena that can occur even for very simple separable equations, as the following two examples illustrate.

\begin{remark}
Existence of solutions is, in general, not global in time: there exist solutions of separable equations that are not defined on the whole interval $I$ where the coefficient $a(t)$ is continuous.
\end{remark}

\begin{example}[Blow-Up in Finite Time]
Consider the Cauchy problem $y' = y^2$, $y(0) = 1$. This equation is separable, with $a(t) \equiv 1$, continuous on all of $\mathbb{R}$, and $g(y) = y^2$. The problem admits a unique solution, given explicitly by $y(t) = 1/(1-t)$. Since
\[
\lim_{t \to 1^-} y(t) = +\infty,
\]
the solution blows up at the finite time $t=1$, and its maximal interval of existence is $I_{\max} = (-\infty, 1)$, strictly smaller than the interval $\mathbb{R}$ on which the coefficients of the equation are defined.
\end{example}

\begin{remark}
A lack of uniqueness may likewise occur whenever the hypotheses of the Local Well-Posedness Theorem, particularly the differentiability of $g$, are violated.
\end{remark}

\begin{example}[Lack of Uniqueness]
Consider the Cauchy problem
\[
y' = 3\sqrt[3]{y^2},
\]
$y(0) = 0$. This equation is separable, with
\[
g(y) = 3\sqrt[3]{y^2},
\]
which fails to be differentiable at $y = 0$. The null constant function $y(t) = 0$, for every $t \in \mathbb{R}$, is one solution of this problem; however, the function $y(t) = t^3$ provides another solution, defined globally on $\mathbb{R}$, distinct from the first.
\end{example}

In fact, the Cauchy problem of the preceding example admits infinitely many solutions, a phenomenon known as Peano's brush (\emph{Pennello di Peano}); as a further illustration, the function
\[
y(t) =
\begin{cases}
0 & t \leq 3, \\
(t-3)^3 & t > 3,
\end{cases}
\]
provides yet another solution of the same Cauchy problem, distinct from both of those exhibited above.

For the initial value problem
\[
y'=ty^3,\qquad y(1)=a,
\]
the case $a=0$ gives the global solution $y\equiv0$. If $a\neq0$,
separation gives
\[
-\frac1{2y^2}=\frac{t^2}{2}+C.
\]
Using $y(1)=a$ determines the constant and yields
\[
\frac1{y^2}=\frac1{a^2}+1-t^2,
\qquad
y(t)=\frac{a}{\sqrt{1+a^2(1-t^2)}}.
\]
The sign is fixed by the initial value. The maximal interval containing
$t=1$ is
\[
I_a=\left(-\sqrt{1+\frac1{a^2}},\,
\sqrt{1+\frac1{a^2}}\right).
\]
Its endpoints move toward $\pm1$ as $|a|$ grows and toward infinity as
$a\to0$, showing explicitly how the initial datum controls the lifespan.

A closely related problem is
\[
y' = (1/2)(y^2 - 1),
\]
with $y(0)=a$. The constant solutions are $y\equiv1$ and $y\equiv-1$.
For $a\neq\pm1$,
\[
\int\frac{\dd y}{y^2-1}=\frac12\int\dd t
\]
gives
\[
\ln\left|\frac{y-1}{y+1}\right|=t+C.
\]
Since $(a-1)/(a+1)$ is the value of the ratio at $t=0$, the solution is
\[
y(t)=\frac{1+C_a e^t}{1-C_a e^t},
\qquad
C_a=\frac{a-1}{a+1}.
\]
Its maximal interval is the component containing $0$ on which
$1-C_a e^t\neq0$. In particular, if $-1<a<1$, then $C_a<0$ and the
solution is global; if $|a|>1$, a finite-time pole occurs in one time
direction.

Figure~\ref{fig:atlas-finite-time-blowup} shows the finite right endpoint of a maximal solution interval.

\begin{figure}[!htbp]
\centering
\begin{tikzpicture}[>=Stealth]
\begin{axis}[
  width=12.3cm,height=6.3cm,
  axis lines=left,
  ticklabel style={font=\small},label style={font=\small},
  xlabel={$t$},ylabel={$y(t)$},
  xmin=-1,xmax=1.25,ymin=0,ymax=10.5,
  xtick={-1,0,1},ytick={1,3,5,7,9}
]
\addplot[figblue,very thick,-{Stealth[length=2.4mm]},domain=-1:0.89,samples=180] {1/(1-x)};
\draw[figred,dashed,thick] (axis cs:1,0) -- (axis cs:1,9.8);
\node[figred,above left,fill=white,inner sep=1pt,font=\small]
  at (axis cs:0.98,9.6) {$T_+=1$};
\addplot[only marks,figblue,mark=*] coordinates {(0,1)};
\node[figblue,above left,font=\scriptsize] at (axis cs:0,1) {$y(0)=1$};
\end{axis}
\end{tikzpicture}
\caption{The solution $y(t)=1/(1-t)$ with initial value $y(0)=1$ satisfies $y^{\prime}=y^2$ and is defined on $(-\infty,1)$. The vertical dashed line marks the blow-up time, not part of the solution. The equation has smooth coefficients, but this trajectory cannot be continued through that time with a finite value.}
\label{fig:atlas-finite-time-blowup}
\end{figure}

\section{Linear First-Order Equations}

\subsection{Global Well-Posedness}

We now turn to the class of linear first-order equations, already introduced above, which take the general form $y' = a(t) y + f(t)$. Several familiar equations from physics and other applications fall within this class, including the equation $P' = 2P$ governing simple population or capital growth, the equation $y' = y + 3x$, and the equation $y' - (1/t) y = \ln t$.

Unlike the case of general separable equations, for which well-posedness of the Cauchy problem could only be guaranteed locally in time, linear equations enjoy a considerably stronger property: well-posedness holds globally on the entire interval where the coefficients are defined. This distinctive feature, a direct consequence of the linear structure of the equation, is recorded in the following theorem.

\begin{theorem}[Global Well-Posedness for Linear Equations]
Let $I$ be an interval, and assume that $a(t)$ and $f(t)$ are both continuous on $I$. Then, for every choice of initial time $t_0 \in I$ and every choice of initial value $y_0 \in \mathbb{R}$, the Cauchy problem
\[
\begin{cases}
y' = a(t) y + f(t), \\
y(t_0) = y_0,
\end{cases}
\]
admits a unique solution $y(t)$, defined on the whole interval $I$.
\end{theorem}


\subsection{The Integrating Factor Method}

Every first-order linear differential equation can be solved explicitly by means of a systematic technique known as the integrating factor method. As a first illustration, one may consider the equation $y' + (1/x) y = 2$, which can be solved by multiplying both sides by the factor $x$; this particular example already suggests the general strategy, which we now describe in full.

To solve the linear differential equation written in the form $y' - a(t) y = f(t)$, one multiplies both sides by the integrating factor
\[
e^{-\int a(t) \, dt}
\]
and subsequently integrates both sides with respect to $t$.

\begin{proof}
Let $A(t)$ denote any fixed antiderivative of $a(t)$. Multiplying the equation
\[
y' - a(t) y = f(t)
\]
by the integrating factor $e^{-A(t)}$, we obtain
\begin{equation*}
e^{-A(t)} \big[ y' - a(t) y \big] = e^{-A(t)} f(t). \tag{$*$}
\end{equation*}

Consider now the product function $y(t) e^{-A(t)}$; its derivative, computed by the product rule together with the Chain Rule applied to $e^{-A(t)}$, is precisely
\[
\frac{d}{dt} \big[ y(t) e^{-A(t)} \big] = e^{-A(t)} \big[ y'(t) - a(t) y(t) \big],
\]
so that equation $(*)$ may be rewritten as
\[
\frac{d}{dt} \big[ y(t) e^{-A(t)} \big] = e^{-A(t)} f(t).
\]

By the definition of the indefinite integral, this identity gives the following relation with an arbitrary constant $c \in \mathbb{R}$:
\[
y(t) e^{-A(t)} + c = \int e^{-A(t)} f(t) \, dt.
\]
Equivalently, the solution is
\[
y(t) = C \, e^{A(t)} + e^{A(t)} \int e^{-A(t)} f(t) \, dt,
\]
where $C$ denotes an arbitrary real constant, absorbing the constant $-c$ from the previous line. This establishes the desired formula for the general solution.
\end{proof}

The proof just given yields an explicit formula for the general solution of any linear first-order equation with continuous coefficients, which we now record for reference.

\begin{theorem}[General Integral of the Linear Equation]
Let $a(t)$ and $f(t)$ be continuous functions on an interval $I$, and let $A(t)$ denote an antiderivative of $a$ on $I$. For an arbitrary constant $C \in \mathbb{R}$, the solutions of $y' = a(t) y + f(t)$ on $I$ and a particular solution are given by
\[
\begin{aligned}
y(t) &= C \, e^{A(t)} + y_P(t), \\
y_P(t) &= e^{A(t)} \int e^{-A(t)} f(t) \, dt.
\end{aligned}
\]

\end{theorem}

\begin{remark}
The general solution formula obtained above can be written suggestively as
\[
y(t) = y_{\text{homo}}(t) + y_P(t).
\]
The homogeneous term is
\[
y_{\text{homo}}(t) = C e^{A(t)}.
\]

This is the general solution of the associated homogeneous equation $y' = a(t) y$, obtained by setting $f \equiv 0$, and $y_P(t)$ denotes a single particular solution of the full nonhomogeneous equation. This decomposition anticipates a structural principle that will play an even more prominent role in the theory of second-order linear equations developed in the following section.
\end{remark}

\section{Second-Order Linear Differential Equations}

\subsection{Second-Order Linear Equations and Initial Data}

We now turn to the study of second-order linear differential equations, which take the general form
\[
a(t) y'' + b(t) y' + c(t) y = f(t), \qquad t \in I,
\]
where $a$, $b$, and $c$ are continuous functions on the interval $I$, called the coefficient functions, subject to the nondegeneracy condition $a(t) \neq 0$ for every $t \in I$, and where $f$ is a continuous function on $I$, called the forcing term. Equations of this type arise naturally in a wide variety of physical contexts; a particularly simple example is furnished by the equation of motion of a falling body, $m y'' = -mg$, whose general solution is readily found by direct integration to be
\[
y(t) = -(1/2) g t^2 + v_0 t + y_0,
\]
where $v_0$ and $y_0$ denote the initial velocity and initial position of the body, respectively. As for first-order linear equations, the associated initial value problem for a second-order linear equation is well-posed globally on the interval of continuity of the coefficients, subject now to the specification of two initial conditions, since the equation involves a second derivative.

\begin{theorem}[Cauchy Theorem for Second-Order Linear Equations]
Assume that $a$, $b$, $c$, and $f$ are continuous functions on an interval $I$, with $a(t) \neq 0$ throughout $I$. Then, for any given $t_0 \in I$, $y_0 \in \mathbb{R}$, and $v_0 \in \mathbb{R}$, the Cauchy problem
\[
\begin{cases}
a(t) y'' + b(t) y' + c(t) y = f(t), \\
y(t_0) = y_0, \\
y'(t_0) = v_0,
\end{cases}
\]
has a unique solution $y : I \to \mathbb{R}$.
\end{theorem}

\subsection{The Homogeneous Equation}

We first study the homogeneous equation associated with the general second-order linear equation, namely
\begin{equation*}
a(t) y'' + b(t) y' + c(t) y = 0, \qquad t \in I, \tag{$\star$}
\end{equation*}
where, as above, $a$, $b$, $c$ are continuous on $I$ with $a(t) \neq 0$ throughout. The linear structure of this equation is reflected in the following superposition principle, which shows that the set of solutions is closed under linear combination.

\begin{theorem}[Superposition Principle for the Homogeneous Equation]
If $y_1(t)$ and $y_2(t)$ are both solutions of the linear homogeneous equation $(\star)$ on the interval $I$, and $c_1, c_2$ are arbitrary constants, then the function
\[
y(t) = c_1 y_1(t) + c_2 y_2(t)
\]
is also a solution of $(\star)$ on $I$.
\end{theorem}

This superposition principle, together with a suitable notion of linear independence, allows one to describe the entire family of solutions of the homogeneous equation in terms of just two particular solutions.

\begin{theorem}[General Integral of the Homogeneous Equation]
If $y_1(t)$ and $y_2(t)$ are linearly independent solutions of the linear homogeneous equation $(\star)$ on the interval $I$, meaning that the ratio $y_1/y_2$ is not constant on $I$, then the general solution of $(\star)$ is given by
\[
y(t) = c_1 y_1(t) + c_2 y_2(t), \qquad t \in I,
\]
where $c_1$ and $c_2$ are arbitrary constants.
\end{theorem}

This result shows that, once two particular linearly independent solutions of the homogeneous equation are known, every solution of the equation is thereby determined; equivalently, one says that the set of all solutions of the second-order linear homogeneous equation $(\star)$ forms a vector space of dimension two.

In general, discovering two particular solutions of a second-order linear equation with variable coefficients can be a genuinely difficult problem. It is, however, always possible to do so explicitly when the coefficients are constant, that is, for an equation of the form $a y'' + b y' + c y = 0$ with $a, b, c \in \mathbb{R}$ and $a \neq 0$; the question of how to find two linearly independent solutions of such a constant-coefficient homogeneous equation, according to the sign of the discriminant of the associated characteristic equation, is addressed through dedicated techniques developed elsewhere in the course.

\subsection{The Nonhomogeneous Equation}

We now turn to the general, nonhomogeneous second-order linear equation,
\begin{equation*}
a(t) y'' + b(t) y' + c(t) y = f(t), \qquad t \in I, \tag{$\star\star$}
\end{equation*}
where $a$, $b$, $c$ are continuous on $I$, with $a(t) \neq 0$ throughout, and $f$ is a given continuous external forcing term. As in the homogeneous case, the linear structure of the equation gives rise to a corresponding superposition principle, which in this setting relates solutions of the equation with different forcing terms.

\begin{theorem}[Superposition Principle for the Nonhomogeneous Equation]
If $y_1(t)$ and $y_2(t)$ are solutions of the linear equation $(\star\star)$ with external forcing terms $f_1$ and $f_2$ respectively, and $c_1, c_2$ are arbitrary constants, then the function
\[
y(t) = c_1 y_1(t) + c_2 y_2(t)
\]
is a solution of $(\star\star)$ with forcing term
\[
f(t) = c_1 f_1(t) + c_2 f_2(t).
\]

\end{theorem}

The general solution of the nonhomogeneous equation can be described in terms of the solutions of the associated homogeneous equation, together with any single particular solution of the nonhomogeneous equation itself, according to the following theorem.

\begin{theorem}[General Integral of the Nonhomogeneous Equation]
The general solution of the nonhomogeneous differential equation $(\star\star)$ can be written as $y = y_{\text{homo}} + y_P$, where $y_P$ is any particular solution of $(\star\star)$, and $y_{\text{homo}}$ denotes the general solution of the associated complementary homogeneous equation $(\star)$.
\end{theorem}

This theorem shows that, once two linearly independent solutions of the associated homogeneous equation $(\star)$ are known, together with a single particular solution of the nonhomogeneous equation $(\star\star)$, every solution of $(\star\star)$ is thereby determined, exactly as in the analogous situation encountered for first-order linear equations.

As in the homogeneous case, finding a particular solution of a nonhomogeneous second-order linear equation with variable coefficients can be genuinely difficult in general. When the coefficient functions are constant, however, that is, for an equation of the form
\[
a y'' + b y' + c y = f(t)
\]
with $a, b, c \in \mathbb{R}$ and $a \neq 0$, a systematic technique known as the Method of Undetermined Coefficients is available, which allows one to guess the form of a particular solution directly from the structure of the forcing term $f(t)$, whether polynomial, exponential, or trigonometric in nature; the detailed development of this method, together with the classification of the homogeneous solutions according to the sign of the discriminant $\Delta$ of the characteristic equation, is left to the dedicated material provided for further study.

\section{Existence, Uniqueness, and Maximal Solutions}

\subsection{Initial Data and Local Lipschitz Regularity}

We consider an autonomous system of ODEs in $\R^N$:
\begin{equation}
\dot{\bv} = f(\bv), \quad \bv(0) = \bv_0
\end{equation}
where $f: \R^N \to \R^N$ is a vector field. A \textit{solution} to this initial value problem is a differentiable function $\bv: I \to \R^N$, defined on some interval $I$ containing $0$, such that
\[
(\dd/\dd t)\bv(t) = f(\bv(t))
\]
for all $t \in I$, and $\bv(0) = \bv_0$. To guarantee that solutions exist and are unique, we must impose regularity conditions on $f$. The standard condition is local Lipschitz continuity.

\begin{definition}[Locally Lipschitz Continuous]
A function $f: \R^N \to \R^N$ is \textit{locally Lipschitz continuous} if for every point $\bx \in \R^N$, there exists a neighborhood $U$ of $\bx$ and a constant $L_U > 0$ such that for all $\bx, \by \in U$:
\begin{equation}
\|f(\bx) - f(\by)\| \leq L_U \|\bx - \by\|
\end{equation}
where $\|\cdot\|$ denotes the standard Euclidean norm.
\end{definition}

\begin{remark}[Why the Lipschitz Condition?]
Continuity of $f$ alone (via Peano's existence theorem) guarantees that a solution exists, but it does not guarantee that the solution is \textit{unique}. Without uniqueness, trajectories could branch, making deterministic prediction impossible. The Lipschitz condition bounds the rate at which $f$ can change, preventing the vector field from becoming too steep and forcing trajectories to merge or branch.
\end{remark}

\subsection{Gronwall's Integral Inequality}


To prove uniqueness over the entire interval $[0, T]$, we rely on a fundamental integral inequality.

\begin{lemma}[Gronwall's Inequality]
Let $u$ and $\beta$ be continuous functions with the following domain and codomain:
\[
u, \beta: [0, T] \to [0, \infty).
\]
Suppose there exists a constant $C \geq 0$ such that the following bound holds for every $t \in [0, T]$:
\begin{equation}
u(t) \leq C + \int_0^t \beta(s) u(s) \dd s.
\end{equation}
For every $t$ in the same interval, the resulting estimate is
\begin{equation}
u(t) \leq C \exp\left( \int_0^t \beta(s) \dd s\right).
\end{equation}

\end{lemma}

\begin{proof}
We construct an explicit upper bound using the method of integrating factors.
Define the function $R(t)$ as the right-hand side of the inequality:
\begin{equation}
R(t) = C + \int_0^t \beta(s) u(s) \dd s
\end{equation}

By the given hypothesis, we immediately have $u(t) \leq R(t)$ for all $t$.
Notice that $R(0) = C$. By the Fundamental Theorem of Calculus, since the integrand $\beta(s)u(s)$ is continuous, $R(t)$ is differentiable and its derivative is:
\begin{equation}
R'(t) = \beta(t) u(t)
\end{equation}
Since $u(t) \leq R(t)$ and $\beta(t) \geq 0$, we can substitute to obtain a differential inequality for $R(t)$:
\begin{equation}
R'(t) \leq \beta(t) R(t) \implies R'(t) - \beta(t) R(t) \leq 0
\end{equation}
To solve this differential inequality, we multiply both sides by the integrating factor
\[
e^{-\int_0^t \beta(s) \dd s},
\]
which is strictly positive. Using the product rule for differentiation, we recognize the left side as the exact derivative of a product:
\begin{equation}
\begin{aligned}
\frac{\dd}{\dd t} \left[ R(t) e^{-\int_0^t \beta(s) \dd s} \right] &= R'(t) e^{-\int_0^t \beta} + R(t) \left( -\beta(t) e^{-\int_0^t \beta} \right) \\
&= e^{-\int_0^t \beta} \Big( R'(t) - \beta(t) R(t) \Big)
\end{aligned}
\end{equation}
The signs of the two factors are
\[
\begin{gathered}
e^{-\int_0^t \beta} > 0 \\
R'(t) - \beta(t) R(t) \leq 0,
\end{gathered}
\]
Hence the derivative of the product is non-positive:
\begin{equation}
\frac{\dd}{\dd t} \left[ R(t) e^{-\int_0^t \beta(s) \dd s} \right] \leq 0
\end{equation}
This implies that the function
\[
g(t) = R(t) e^{-\int_0^t \beta(s) \dd s}
\]
is monotonically non-increasing. Therefore, for any $t \geq 0$, $g(t) \leq g(0)$. Evaluating at $t=0$:
\begin{equation}
\begin{aligned}
g(0) &= R(0) e^0 \\
&= C \cdot 1 \\
&= C
\end{aligned}
\end{equation}
Thus,
\[
R(t) e^{-\int_0^t \beta(s) \dd s} \leq C.
\]
Multiplying both sides by the positive quantity
\[
e^{\int_0^t \beta(s) \dd s}
\]
yields:
\begin{equation}
R(t) \leq C \exp\left( \int_0^t \beta(s) \dd s \right)
\end{equation}
Finally, since $u(t) \leq R(t)$, we conclude that
\[
u(t) \leq C \exp\left( \int_0^t \beta(s) \dd s \right).
\]

\end{proof}

\subsection{Application of Gronwall's Inequality to Global Uniqueness}

\begin{theorem}[Global Uniqueness]
If $\bv_1(t)$ and $\bv_2(t)$ are two solutions to $\dot{\bv} = f(\bv)$ on the interval $[0, T]$ satisfying the same initial condition
\[
\begin{aligned}
\bv_1(0) &= \bv_2(0) \\
&= \bv_0,
\end{aligned}
\]
and if $f$ is globally Lipschitz continuous with constant $L$ (i.e.,
\[
\|f(\bx) - f(\by)\| \leq L\|\bx - \by\|
\]
for all $\bx, \by \in \R^N$), then $\bv_1(t) = \bv_2(t)$ for all $t \in [0, T]$.
\end{theorem}

\begin{proof}
First, we convert the differential equation into an integral equation. Since $\bv_i(t)$ is a solution,
\[
(\dd/\dd t)\bv_i(t) = f(\bv_i(t)).
\]
Integrating both sides from $0$ to $t$ and applying the Fundamental Theorem of Calculus yields:
\begin{equation}
\bv_i(t) - \bv_i(0) = \int_0^t f(\bv_i(s)) \dd s \implies \bv_i(t) = \bv_0 + \int_0^t f(\bv_i(s)) \dd s
\end{equation}
Now, consider the difference between the two solutions. Let
\[
u(t) = \|\bv_1(t) - \bv_2(t)\|.
\]
Subtracting the integral equations:
\begin{equation}
\bv_1(t) - \bv_2(t) = \int_0^t \Big( f(\bv_1(s)) - f(\bv_2(s)) \Big) \dd s
\end{equation}
Taking the norm of both sides and applying the triangle inequality for integrals,
\[
\|\int \mathbf{g}\| \leq \int \|\mathbf{g}\|
\]
gives:
\begin{equation}
\begin{aligned}
u(t) &= \left\| \int_0^t \Big( f(\bv_1(s)) - f(\bv_2(s)) \Big) \dd s \right\| \\
&\leq \int_0^t \left\| f(\bv_1(s)) - f(\bv_2(s)) \right\| \dd s
\end{aligned}
\end{equation}
We now explicitly apply the Lipschitz condition
\[
\|f(\bx) - f(\by)\| \leq L\|\bx - \by\|
\]
to the integrand:
\begin{equation}
\begin{aligned}
u(t) &\leq \int_0^t L \|\bv_1(s) - \bv_2(s)\| \dd s \\
&= L \int_0^t u(s) \dd s
\end{aligned}
\end{equation}

This is exactly the hypothesis of Gronwall's Inequality with $C = 0$ and $\beta(s) = L$. Applying the lemma:
\begin{equation}
\begin{aligned}
u(t) &\leq 0 \cdot \exp\left( \int_0^t L \dd s \right) \\
&= 0
\end{aligned}
\end{equation}
By the definition of a norm, the difference satisfies
\[
\begin{aligned}
u(t) &= \|\bv_1(t) - \bv_2(t)\| \\
&\geq 0
\end{aligned}
\]
Consequently, we must have $u(t) = 0$ for all $t \in [0, T]$. Therefore, $\bv_1(t) = \bv_2(t)$.
\end{proof}

\subsection{Maximal Interval of Existence}

Solutions to ODEs are guaranteed to exist locally. By continuously patching together local solutions, we can extend the solution until it either exists for all time or ``blows up''.

\begin{definition}[Maximal Interval of Existence]
The \textit{maximal interval of existence} $(T_-, T_+)$ is the largest open interval on which a solution $\bv(t)$ to the initial value problem is defined. A solution is called \textit{global} if $T_+ = \infty$ and $T_- = -\infty$.
\end{definition}

\begin{theorem}[Escape Lemma (Finite-Time Blow-up)]
Let $\bv(t)$ be the maximal solution defined on $(T_-, T_+)$. If the forward maximal time is finite, i.e., $T_+ < \infty$, then the trajectory must escape to infinity:
\begin{equation}
\lim_{t \to T_+^-} \|\bv(t)\| = \infty
\end{equation}

\end{theorem}

\begin{proof}
Suppose, to the contrary, that the norm does not tend to infinity.
Then there are times $t_n\to T_+^-$ and a constant $R$ such that
$\|\bv(t_n)\|\leq R$. Passing to a subsequence, compactness of the closed
ball in $\mathbb{R}^N$ gives
\[
\bv(t_n)\longrightarrow\bv_*.
\]
Choose a closed ball $B$ centered at $\bv_*$ on which $f$ is bounded and
Lipschitz. The quantitative local-existence theorem supplies a time
$\delta>0$, independent of all sufficiently large $n$, such that the
initial value problem starting from $(t_n,\bv(t_n))$ has a unique solution
on $[t_n,t_n+\delta]$. By uniqueness, this local solution agrees with the
original one wherever both are defined.

For sufficiently large $n$, $T_+-t_n<\delta$. The solution starting at
$(t_n,\bv(t_n))$ therefore continues the original trajectory beyond
$T_+$, contradicting maximality. Hence for every $R>0$ the trajectory
eventually leaves the ball of radius $R$, which is exactly
\[
\|\bv(t)\|\longrightarrow\infty
\qquad (t\to T_+^-).
\]
\end{proof}

\begin{remark}[Practical Consequence]
The Escape Lemma provides a powerful practical tool: if you can prove that a trajectory is confined to a bounded region (for example, by finding a Lyapunov function or an invariant set), you immediately guarantee that $T_+ = \infty$, meaning the solution exists globally for all future time.
\end{remark}

\section{Stability and Local Dynamics}

\subsection{Lyapunov Stability}

When analyzing the long-term behavior of a system, we often study the behavior near an equilibrium point.

\begin{definition}[Equilibrium and Stability]
A point $\bv_\infty \in \R^N$ is an \textit{equilibrium} (or fixed point) if $f(\bv_\infty) = \mathbf{0}$. The constant function $\bv(t) = \bv_\infty$ is then a trivial solution.
The equilibrium $\bv_\infty$ is:
\begin{itemize}
\item \textit{Lyapunov stable}: if for every $\varepsilon > 0$, there exists a $\delta > 0$ such that if $\|\bv(0) - \bv_\infty\| < \delta$, then $\|\bv(t) - \bv_\infty\| < \varepsilon$ for all $t \geq 0$.
\item \textit{Asymptotically stable}: if it is Lyapunov stable, and there exists a $\delta_0 > 0$ such that if $\|\bv(0) - \bv_\infty\| < \delta_0$, then

\[
\lim_{t \to \infty} \bv(t) = \bv_\infty.
\]

\end{itemize}
\end{definition}

\begin{remark}[Geometric Meaning]
Lyapunov stability means that trajectories starting sufficiently close to the equilibrium remain arbitrarily close to it forever; they do not necessarily converge to it, but they never escape a small tube around it. Asymptotic stability adds the requirement that they actually converge to the equilibrium.
\end{remark}

To prove stability without explicitly solving the ODE, we use Lyapunov's Direct Method, which relies on energy-like scalar functions.

\begin{definition}[Lyapunov Function and Orbital Derivative]
Let $U$ be an open neighborhood of $\bv_\infty$. A continuously differentiable ($C^1$) function $V: U \to \R$ is a \textit{Lyapunov function} if:
\begin{enumerate}
\item \textit{Positive definiteness}: $V(\bv_\infty) = 0$ and $V(\bv) > 0$ for all $\bv \in U \setminus \{\bv_\infty\}$.
\item \textit{Orbital derivative}: the orbital derivative $\dot{V}(\bv) \leq 0$ for all $\bv \in U$.
\end{enumerate}
The \textit{orbital derivative} $\dot{V}(\bv)$ is the time derivative of $V$ evaluated along the trajectories of the system. By the multivariate chain rule:
\begin{equation}
\begin{aligned}
\dot{V}(\bv) &= \frac{\dd}{\dd t} V(\bv(t)) \bigg|_{t=0} \\
&= \nabla V(\bv) \cdot \dot{\bv} \\
&= \nabla V(\bv) \cdot f(\bv)
\end{aligned}
\end{equation}

\end{definition}

\begin{theorem}[Lyapunov's Direct Method]
If a Lyapunov function $V$ exists for $\bv_\infty$, then $\bv_\infty$ is Lyapunov stable. If, furthermore, $\dot{V}(\bv) < 0$ for all $\bv \in U \setminus \{\bv_\infty\}$ (strictly negative definite), then $\bv_\infty$ is asymptotically stable.
\end{theorem}

\begin{proof}[Rigorous Proof of Stability]
Let $\varepsilon > 0$ be given. We must find a $\delta > 0$ such that trajectories starting in the ball $B_\delta(\bv_\infty)$ never leave the ball $B_\varepsilon(\bv_\infty)$.
Consider the boundary of the $\varepsilon$-ball,
\[
S_\varepsilon = \{ \bv : \|\bv - \bv_\infty\| = \varepsilon \}.
\]

This set is closed and bounded, hence compact. Since $\bv_\infty \notin S_\varepsilon$ and $V$ is continuous and positive definite, $V$ attains a strictly positive minimum on $S_\varepsilon$. Let:
\begin{equation}
m = \min_{\bv \in S_\varepsilon} V(\bv) > 0
\end{equation}
Define the sublevel set
\[
U_m = \{ \bv \in U : V(\bv) < m \}.
\]

Because $V$ is continuous, $U_m$ is an open set. Since $V(\bv_\infty) = 0 < m$, $\bv_\infty \in U_m$. Therefore, there exists a $\delta > 0$ such that the ball $B_\delta(\bv_\infty) \subset U_m$.
Now, suppose $\bv(0) \in B_\delta(\bv_\infty)$. Then $V(\bv(0)) < m$.
Along the trajectory, the time derivative of $V$ is
\[
\begin{aligned}
(\dd/\dd t) V(\bv(t)) &= \dot{V}(\bv(t)) \\
&\leq 0.
\end{aligned}
\]
This means $V(\bv(t))$ is a monotonically non-increasing function of time. Thus, for all $t \geq 0$:
\begin{equation}
V(\bv(t)) \leq V(\bv(0)) < m
\end{equation}

This implies that $\bv(t) \in U_m$ for all $t \geq 0$. By the definition of $m$, any point on the boundary $S_\varepsilon$ has $V(\bv) \geq m$. Since $V(\bv(t)) < m$, the trajectory can never reach $S_\varepsilon$. Because the trajectory is continuous and starts inside $B_\varepsilon$, it can never cross $S_\varepsilon$ to exit the ball. Hence, $\|\bv(t) - \bv_\infty\| < \varepsilon$ for all $t \geq 0$, proving Lyapunov stability.
\end{proof}

\begin{proof}[Rigorous Proof of Asymptotic Stability]
We already know the equilibrium is stable, so trajectories starting sufficiently close remain close. We must now prove they converge to $\bv_\infty$.
Assume $\dot{V}(\bv) < 0$ for $\bv \neq \bv_\infty$. For a trajectory starting near $\bv_\infty$, $V(\bv(t))$ is monotonically decreasing and bounded below by $0$. By the Monotone Convergence Theorem for real sequences, the limit must exist:
\begin{equation}
\begin{aligned}
\lim_{t \to \infty} V(\bv(t)) &= c \\
&\geq 0
\end{aligned}
\end{equation}

We must show that $c = 0$. Suppose, for contradiction, that $c > 0$.
Because the trajectory is stable, it is confined to a compact set
\[
K = \{ \bv : c \leq V(\bv) \leq V(\bv(0)) \}.
\]

Note that $\bv_\infty \notin K$ because $V(\bv_\infty) = 0 < c$.
The orbital derivative $\dot{V}(\bv)$ is a continuous function on the compact set $K$. Since $\dot{V}(\bv) < 0$ everywhere on $K$, it must attain a strict maximum on $K$. Let this maximum be $-\eta$, where $\eta > 0$. Thus, for all $\bv \in K$:
\begin{equation}
\dot{V}(\bv) \leq -\eta < 0
\end{equation}
Evaluating this along the trajectory $\bv(t) \in K$:
\begin{equation}
\frac{\dd}{\dd t} V(\bv(t)) \leq -\eta
\end{equation}
Integrating both sides from $0$ to $t$:
\begin{equation}
V(\bv(t)) - V(\bv(0)) \leq -\eta t \implies V(\bv(t)) \leq V(\bv(0)) - \eta t
\end{equation}

As $t \to \infty$, the right-hand side approaches $-\infty$. This implies $V(\bv(t)) \to -\infty$, which strictly contradicts the fact that $V(\bv) \geq 0$ for all $\bv$.
Therefore, our assumption $c > 0$ must be false, and we conclude $c = 0$.
Since $V(\bv(t)) \to 0$ and $V$ is continuous and positive definite (meaning
\[
V(\bv) = 0 \iff \bv = \bv_\infty
\]
), it follows rigorously that
\[
\lim_{t \to \infty} \bv(t) = \bv_\infty.
\]

\end{proof}

\subsection{Hartman--Grobman and Center Manifolds}

When a Lyapunov function is difficult to construct, we often analyze the local behavior near an equilibrium by linearizing the system.

\begin{definition}[Jacobian and Hyperbolic Equilibrium]
The \textit{Jacobian matrix} of $f$ at $\bv_\infty$ is the $N \times N$ matrix $J$ of partial derivatives, with entries
\[
J_{ij} = (\partial f_i/\partial v_j)(\bv_\infty).
\]

The equilibrium $\bv_\infty$ is called \textit{hyperbolic} if every eigenvalue $\lambda$ of $J$ has a nonzero real part, i.e., $\Re(\lambda) \neq 0$.
\end{definition}

\begin{theorem}[Hartman--Grobman Theorem]
If $\bv_\infty$ is a hyperbolic equilibrium, then the nonlinear flow near $\bv_\infty$ is \textit{topologically conjugate} to the linear flow $\dot{\bepsilon} = J\bepsilon$.
\end{theorem}

\begin{remark}[Unpacking Topological Conjugacy]
What does ``topologically conjugate'' mean? It means there exists a homeomorphism (a continuous, bijective map with a continuous inverse) $h: U \to V$, where $U$ is a neighborhood of $\bv_\infty$ and $V$ is a neighborhood of the origin in $\R^N$, such that $h(\bv_\infty) = \mathbf{0}$, and $h$ maps trajectories of the nonlinear system to trajectories of the linear system, preserving the direction of time. Formally, if $\Phi_t(\bx)$ is the nonlinear flow and $e^{Jt}$ is the linear flow, then:
\begin{equation}
h(\Phi_t(\bx)) = e^{Jt} h(\bx)
\end{equation}

This implies that the nonlinear and linearized systems have the same local
orbit structure up to a homeomorphism. The eigenvalues determine the
stable and unstable dimensions and hence the local topological behavior
relevant to hyperbolicity. They also determine finer linear features such
as decay rates and whether linear trajectories spiral, but such
quantitative or differentiable information is not generally preserved by
a merely topological conjugacy.
\end{remark}

\begin{remark}[Why Hyperbolicity is Required]
If $J$ has an eigenvalue with zero real part ($\Re(\lambda) = 0$), the equilibrium is \textit{non-hyperbolic}. In this case, the linear terms are ``neutral'' (neither exponentially growing nor decaying), and the nonlinear higher-order terms in the Taylor expansion of $f$ can dominate the local dynamics. The nonlinear system could exhibit behaviors completely absent in the linearization, such as a center (closed orbits) or a weak focus. Thus, linearization is inconclusive for non-hyperbolic equilibria.
\end{remark}

To handle non-hyperbolic equilibria, we use the Center Manifold Theorem, which allows us to reduce the dimensionality of the system.

\begin{definition}[Spectral Subspaces]
Based on the eigenvalues of the Jacobian $J$, we decompose $\R^N$ into three invariant linear subspaces:
\begin{itemize}
\item \textit{Stable subspace $E^s$}: the span of all generalized eigenvectors corresponding to eigenvalues with $\Re(\lambda) < 0$.
\item \textit{Unstable subspace $E^u$}: the span of all generalized eigenvectors corresponding to eigenvalues with $\Re(\lambda) > 0$.
\item \textit{Center subspace $E^c$}: the span of all generalized eigenvectors corresponding to eigenvalues with $\Re(\lambda) = 0$.
\end{itemize}
By linear algebra, $\R^N = E^s \oplus E^u \oplus E^c$.
\end{definition}

\begin{theorem}[Center Manifold Theorem]
There exists a locally invariant manifold $W^c$ (the center manifold) passing through $\bv_\infty$ such that:
\begin{enumerate}
\item \textit{Tangency}: $W^c$ is tangent to the center subspace $E^c$ at $\bv_\infty$.
\item \textit{Local invariance}: $W^c$ is locally invariant: if a trajectory starts in $W^c$, it remains in $W^c$ for all time (as long as it stays in the neighborhood).
\item \textit{Dimension}: the dimension of $W^c$ is equal to the dimension of $E^c$.
\end{enumerate}
Furthermore, there exist locally invariant stable and unstable manifolds $W^s$ and $W^u$ tangent to $E^s$ and $E^u$, respectively.
\end{theorem}

\begin{remark}[Why is this Useful?]
The stable and unstable manifolds describe directions of local
exponential attraction and repulsion. The center manifold captures the
local dynamics associated with the neutral center directions and is the
appropriate reduced object for studying local bifurcations and stability
changes that linearization cannot decide. This is a local reduction
principle, not a claim that all interesting or global dynamics of the
full system lie on $W^c$. Because
\[
\dim(W^c) = \dim(E^c) < N
\]

 (assuming $E^s$ or $E^u$ are non-trivial), the Center Manifold Theorem allows us to reduce the study of the full $N$-dimensional nonlinear system to a lower-dimensional system restricted to $W^c$. On this reduced space, the linear terms are zero, and the dynamics are governed purely by the nonlinear terms, which can be analyzed using techniques like Lyapunov functions or normal forms.
\end{remark}

Figure~\ref{fig:ch15-phase-portraits} compares the local phase portraits of four types of planar equilibrium.

\begin{figure}[!htbp]
\centering
\begin{tikzpicture}
\begin{axis}[width=6.1cm,height=6.1cm,axis lines=middle,xmin=-2,xmax=2,ymin=-2,ymax=2,
  xtick=\empty,ytick=\empty,]
  \foreach \sx/\sy in {1/1,-1/1,1/-1,-1/-1}{
    \foreach \c in {0.4,0.8,1.4,2.2}{
      \addplot[figblue,thick,domain=0:2.2,samples=50,forget plot]
        ({\sx*\c*exp(-0.6*x)},{\sy*\c*0.4*exp(-0.3*x)});
    }
    \addplot[figblue,thick,-{Stealth[length=2.6mm]},forget plot] coordinates {
      ({\sx*1.4*exp(-0.6*1.55)},{\sy*1.4*0.4*exp(-0.3*1.55)})
      ({\sx*1.4*exp(-0.6*1.85)},{\sy*1.4*0.4*exp(-0.3*1.85)})
    };
  }
\end{axis}
\begin{axis}[at={(6.6cm,0)},width=6.1cm,height=6.1cm,axis lines=middle,xmin=-2,xmax=2,ymin=-2,ymax=2,
  xtick=\empty,ytick=\empty,]
  \foreach \c in {0.3,0.7,1.2,2}{
    \addplot[figred,thick,domain=-1.5:1.5,samples=60,forget plot] ({x},{\c/(x+0.001)});
    \addplot[figred,thick,domain=-1.5:1.5,samples=60,forget plot] ({x},{-\c/(x+0.001)});
  }
  \draw[vecB] (-1.9,-1.9) -- (-0.35,-0.35);
  \draw[vecB] (1.9,1.9) -- (0.35,0.35);
  \draw[vecB] (0.35,-0.35) -- (1.9,-1.9);
  \draw[vecB] (-0.35,0.35) -- (-1.9,1.9);
\end{axis}
\begin{axis}[at={(0,-6.6cm)},width=6.1cm,height=6.1cm,axis lines=middle,xmin=-2,xmax=2,ymin=-2,ymax=2,
  xtick=\empty,ytick=\empty,]
  \addplot[figgreen,thick,domain=0:12,samples=200,forget plot,variable=t]
    ({1.9*exp(-0.15*t)*cos(deg(t))},{1.9*exp(-0.15*t)*sin(deg(t))});
  \addplot[figgreen,thick,-{Stealth[length=2.6mm]},forget plot] coordinates {
    ({1.9*exp(-0.15*7.4)*cos(deg(7.4))},{1.9*exp(-0.15*7.4)*sin(deg(7.4))})
    ({1.9*exp(-0.15*7.8)*cos(deg(7.8))},{1.9*exp(-0.15*7.8)*sin(deg(7.8))})
  };
\end{axis}
\begin{axis}[at={(6.6cm,-6.6cm)},width=6.1cm,height=6.1cm,axis lines=middle,xmin=-2,xmax=2,ymin=-2,ymax=2,
  xtick=\empty,ytick=\empty,]
  \foreach \r in {0.5,0.9,1.3,1.7}{
    \addplot[black!55,thick,domain=0:360,samples=80,forget plot] ({\r*cos(x)},{\r*sin(x)});
  }
\end{axis}
\end{tikzpicture}
\caption[Phase portraits: node, saddle, focus, center]{Local phase portraits around the origin. Top left: stable node with two negative real eigenvalues. Top right: saddle with one negative and one positive real eigenvalue. Bottom left: stable focus with complex eigenvalues having negative real part. Bottom right: non-hyperbolic center with purely imaginary eigenvalues. Arrows indicate the direction of evolution where shown.}
\label{fig:ch15-phase-portraits}
\end{figure}

\section{Partial Differential Equations: Solution Methods}
\label{sec:pde-methods}

The ordinary differential equations studied so far describe finitely many state variables evolving in time. Spatially distributed systems require an unknown field $u(\bx,t)$ and partial derivatives in both space and time. The integral theorems, spectral methods, and distribution theory developed in earlier chapters now provide tools for solving these equations.

\subsection{Fundamental PDEs and Separation of Variables}

Three fundamental linear second-order partial differential equations are
\begin{align}
\text{Heat (diffusion) equation:}\qquad &\frac{\partial u}{\partial t}=D\,\Delta u,\\
\text{Wave equation:}\qquad &\frac{\partial^2u}{\partial t^2}=c^2\,\Delta u,\\
\text{Laplace (steady-state) equation:}\qquad &\Delta u=0,
\end{align}
Here the Laplacian is
\[
\begin{aligned}
\Delta u &= \nabla\cdot(\nabla u) \\
&= \sum_i\partial^2u/\partial x_i^2,
\end{aligned}
\]

The constant $D>0$ is the diffusion coefficient, and $c>0$ is the propagation speed. In mathematical biology, the diffusion equation models the spatial spreading of a cellular or molecular concentration. Initial and boundary conditions are needed to select a solution appropriate to the physical problem.


The basic idea is to seek a factored solution
\[
u(\bx,t)=X(\bx)T(t).
\]

Substitution into a suitable linear PDE separates the spatial and temporal dependence, linking them by a separation constant. In one spatial dimension, this produces two ordinary differential equations; in higher dimensions, the spatial factor generally satisfies an eigenvalue PDE.

\begin{example}[The Heat Equation with Homogeneous Dirichlet Conditions]
\label{ex:heat-separation}
Consider $u_t=D\,u_{xx}$ for $0<x<L$ and $t>0$, with
\[
u(0,t)=u(L,t)=0,\qquad u(x,0)=f(x).
\]

The boundary conditions are called \emph{homogeneous Dirichlet conditions} because the function vanishes at the boundary.

\textit{Step 1: separate the variables.} Substituting

\[
u(x,t)=X(x)T(t)
\]
gives $XT'=DX''T$. Wherever the factors are nonzero,
\begin{equation}
\frac{T'(t)}{D\,T(t)}=\frac{X''(x)}{X(x)}.
\end{equation}

The left-hand side depends only on $t$ and the right-hand side only on $x$. They must therefore equal a common constant, denoted by $-\lambda$:
\begin{equation}
T'=-\lambda D\,T,\qquad X''=-\lambda X.
\end{equation}

These equations are then imposed throughout the interval, including any zeros of $X$; the quotient is only a device for identifying the separated equations.

\textit{Step 2: solve the spatial eigenvalue problem.} A nonzero temporal factor requires

\[
\begin{aligned}
X(0) &= X(L) \\
&= 0.
\end{aligned}
\]
For $\lambda>0$,
\[
X(x)=A\cos(\sqrt\lambda x)+B\sin(\sqrt\lambda x).
\]

The first boundary condition gives $A=0$, and the second requires $\sin(\sqrt{\lambda} L)=0$ for a nontrivial solution. Hence
\begin{equation}
\lambda_n=\left(\frac{n\pi}{L}\right)^2,
\qquad X_n(x)=\sin\left(\frac{n\pi x}{L}\right),\qquad n=1,2,\dots.
\end{equation}
No nontrivial solution has $\lambda\le0$: integration by parts in $-X''=\lambda X$ gives
\[
\int_0^L|X'|^2\,\dd x=\lambda\int_0^L|X|^2\,\dd x,
\]
and the zero endpoint values rule out a nonzero constant function.

\textit{Step 3: solve the temporal equation and superpose.} For each eigenvalue,

\[
T_n(t)=e^{-\lambda_nDt}
\]
up to a multiplicative constant. Linearity permits the expansion
\begin{equation}
u(x,t)=\sum_{n=1}^{\infty}b_n\sin\left(\frac{n\pi x}{L}\right)e^{-\lambda_nDt}.
\end{equation}

Matching the initial data amounts to expanding $f$ in a Fourier sine series. Orthogonality on $(0,L)$ gives
\begin{equation}
b_n=\frac{2}{L}\int_0^L f(x)\sin\left(\frac{n\pi x}{L}\right)\,\dd x.
\end{equation}

For $f\in L^2(0,L)$, the initial condition is attained in the $L^2$ sense as $t\downarrow0$, and the exponential factors smooth the series for $t>0$. With suitable smoothness and endpoint compatibility, it gives the classical solution up to the initial boundary. The functions $X_n$ are eigenfunctions of the self-adjoint Dirichlet operator $-\dd^2/\dd x^2$, just as orthogonal eigenvectors arise for symmetric matrices. Thus separation of variables connects ODE methods, Fourier expansions, and operator theory.
\end{example}

Figure~\ref{fig:atlas-heat-mode-decay} illustrates the different decay rates of separated heat-equation modes.

\begin{figure}[!htbp]
\centering
\begin{tikzpicture}[>=Stealth]
\begin{axis}[width=11cm,height=5.8cm,axis lines=middle,ticklabel style={font=\small},label style={font=\small},legend style={font=\scriptsize,draw=black!20,fill=white},xlabel={$x$},ylabel={$u(x,t)$},xmin=0,xmax=1.05,ymin=0,ymax=1.3,xtick={0,0.5,1},legend columns=3,legend pos=north east]
\addplot[figblue,very thick,domain=0:1,samples=150] {sin(deg(pi*x))+0.35*sin(deg(3*pi*x))};
\addlegendentry{$t=0$}
\addplot[figred,thick,dashed,domain=0:1,samples=150] {exp(-pi^2*0.03)*sin(deg(pi*x))+0.35*exp(-9*pi^2*0.03)*sin(deg(3*pi*x))};
\addlegendentry{$t=0.03$}
\addplot[figgreen,thick,dashdotted,domain=0:1,samples=150] {exp(-pi^2*0.1)*sin(deg(pi*x))+0.35*exp(-9*pi^2*0.1)*sin(deg(3*pi*x))};
\addlegendentry{$t=0.1$}
\end{axis}
\end{tikzpicture}
\caption{Three profiles solve $u_t=u_{xx}$ on $[0,1]$ with zero boundary values. The initial profile combines the first sine mode with $0.35$ times the third. The third mode decays nine times faster in the exponent, so small-scale structure disappears before the main profile relaxes to zero.}
\label{fig:atlas-heat-mode-decay}
\end{figure}

\subsection{Dirichlet and Neumann Boundary-Value Problems}
\label{sec:pde-boundary-problems}

\begin{definition}[Dirichlet and Neumann Problems]
Let $\Omega\subset\R^N$ be a bounded connected domain with smooth boundary. A \emph{Dirichlet problem} for the Laplace equation prescribes the boundary values:
\begin{equation}
\Delta u=0\quad\text{in }\Omega,\qquad u=g\quad\text{on }\partial\Omega.
\end{equation}
A \emph{Neumann problem} instead prescribes the outward normal derivative:
\begin{equation}
\Delta u=0\quad\text{in }\Omega,\qquad
\frac{\partial u}{\partial\mathbf{n}}=h\quad\text{on }\partial\Omega,
\end{equation}
where $\partial u/\partial\mathbf{n}:=\nabla u\cdot\mathbf{n}$. This specifies a quantity proportional to the boundary flux rather than the value of $u$ itself; for a diffusive current $\mathbf{J}=-D\nabla u$, the outward flux is $-D h$.
\end{definition}

\begin{theorem}[Uniqueness for the Dirichlet Problem]
If a classical solution of the Dirichlet problem exists with sufficient regularity for integration by parts, it is unique.
\end{theorem}

\begin{proof}[The Energy Method]
Let $u_1,u_2$ be two solutions and set $w=u_1-u_2$. Then $\Delta w=0$ in $\Omega$ and $w=0$ on $\partial\Omega$. The product rule gives
\begin{equation}
\begin{aligned}
\nabla\cdot(w\nabla w) &= |\nabla w|^2+w\Delta w \\
&= |\nabla w|^2.
\end{aligned}
\end{equation}
Integrating and applying the $N$-dimensional version of the divergence theorem yields
\begin{equation}
\int_\Omega|\nabla w|^2\,\dd\bx
=\int_{\partial\Omega}w\frac{\partial w}{\partial\mathbf{n}}\,\dd S=0.
\end{equation}

The continuous nonnegative integrand must vanish identically. Thus $\nabla w=0$, and connectedness implies that $w$ is constant. Its zero boundary value forces that constant to be zero, so $u_1=u_2$.
\end{proof}

\begin{remark}[Neumann Uniqueness and Compatibility]
For two solutions of the same Neumann problem, the difference satisfies $\partial w/\partial\mathbf{n}=0$, which again makes the energy boundary term vanish. Therefore $w$ is constant, but the constant need not be zero: the solution is unique only up to an additive constant. Existence also requires the compatibility condition
\begin{equation}
\int_{\partial\Omega}h\,\dd S=\int_\Omega\Delta u\,\dd\bx=0.
\end{equation}
Prescribing, for example, the mean of $u$ fixes the additive constant.
\end{remark}

\subsection{Green's Functions}
\label{sec:green-functions}

We now consider Poisson's equation $\Delta u=-\rho$ in $\Omega$, with homogeneous Dirichlet data. The idea parallels the use of the Dirac delta in distribution theory: if the response to a point source at $\by$ is known, the response to a distributed source $\rho$ can be reconstructed by integrating these point-source responses, weighted by $\rho(\by)$.

\begin{definition}[Dirichlet Green's Function]
A \emph{Green's function} for $-\Delta$ on $\Omega$ with homogeneous Dirichlet conditions satisfies, for each fixed $\by\in\Omega$,
\begin{equation}
-\Delta_{\bx}G(\bx,\by)=\delta(\bx-\by)\quad\text{in }\Omega,
\qquad G(\bx,\by)=0\quad\text{on }\partial\Omega.
\end{equation}

The differential equation is interpreted in the distributional sense. The multidimensional Dirac delta acts on a test function by $\delta_{\by}[\varphi]=\varphi(\by)$.
\end{definition}

\begin{theorem}[Green's Function Representation]
For a sufficiently regular domain and source for which the Dirichlet Green's function representation is valid, the solution of $\Delta u=-\rho$ in $\Omega$, $u=0$ on $\partial\Omega$, is
\begin{equation}
u(\bx)=\int_\Omega G(\bx,\by)\rho(\by)\,\dd\by.
\end{equation}

\end{theorem}

\begin{proof}[Distributional Verification]
Formally applying the Laplacian gives
\begin{align}
\Delta_{\bx}u(\bx)
&=\int_\Omega\Delta_{\bx}G(\bx,\by)\rho(\by)\,\dd\by\notag\\
&=-\int_\Omega\delta(\bx-\by)\rho(\by)\,\dd\by=-\rho(\bx).
\end{align}

Because $G$ is singular at $\bx=\by$, differentiation here is understood distributionally rather than as an unqualified exchange of classical derivatives and integration. More explicitly, for a compactly supported smooth test function $\varphi$, the defining identity for $G$ and Fubini's theorem give
\begin{align}
\int_\Omega u(\bx)(-\Delta\varphi(\bx))\,\dd\bx
&=\int_\Omega\rho(\by)\left[\int_\Omega G(\bx,\by)(-\Delta\varphi(\bx))\,\dd\bx\right]\dd\by\notag\\
&=\int_\Omega\rho(\by)\varphi(\by)\,\dd\by.
\end{align}

Under the assumed regularity, the boundary trace is zero because the Green's function has zero Dirichlet trace. When the represented solution is classical, the preceding uniqueness theorem identifies it as the only solution.
\end{proof}

\begin{example}[The Upper Half-Plane and the Method of Images]
Consider the upper half-plane
\[
\Omega=\{\bx=(x_1,x_2):x_2>0\}\subset\R^2
\]

Let $\by=(y_1,y_2)$ with $y_2>0$. Reflect the source across the boundary to obtain $\by^*=(y_1,-y_2)$. Then
\begin{equation}
G(\bx,\by)=\frac{1}{2\pi}\log\frac{|\bx-\by^*|}{|\bx-\by|}
\end{equation}
is a Dirichlet Green's function for the upper half-plane. Indeed, the two-dimensional fundamental solution satisfies
\[
-\Delta_{\bx}\left(-\frac{1}{2\pi}\log|\bx-\by|\right)=\delta(\bx-\by).
\]

The added term $\log|\bx-\by^*|/(2\pi)$ is harmonic throughout $\Omega$ because its singularity lies outside the domain. On $x_2=0$, the distances to $\by$ and $\by^*$ coincide, so the logarithmic ratio vanishes. The normalization of the fundamental solution follows from the divergence theorem on a small circle around the source, just as $1/(4\pi|\bx-\by|)$ is the fundamental solution of $-\Delta$ in three dimensions.

Adding a fictitious exterior source to enforce the boundary condition is the \emph{method of images}. It is particularly useful on domains with suitable symmetries, such as half-spaces and balls, although the image location and strength depend on the geometry. On an unbounded domain, appropriate behavior at infinity must also be specified to obtain uniqueness.
\end{example}
\newpage
\chapter{Probabilistic Foundations and Limit Theorems}

\section{Probability and Moment Bounds}

\subsection{Foundations of Probability Theory}

To discuss limit theorems rigorously, we must first establish the measure-theoretic foundation of probability. This framework prevents the paradoxes that arise in naive probability and allows us to handle continuous and infinite-dimensional spaces seamlessly.

\begin{definition}[Probability Space]
A \textit{probability space} is a triple $(\Omega, \mathcal{F}, \Prob)$ consisting of:
\begin{enumerate}
\item \textit{Sample space ($\Omega$)}: a nonempty set of all possible outcomes of a random experiment. Each element $\omega$ represents one individual outcome.
\item \textit{$\sigma$-algebra of events ($\mathcal{F}$)}: a collection of subsets of $\Omega$, called events, that contains $\Omega$ and is closed under complements relative to $\Omega$ and countable unions. Explicitly, for every event $A$ and every sequence of events $A_n$, these conditions are
\[
\begin{gathered}
\Omega \in \mathcal{F}, \\
\Omega \setminus A \in \mathcal{F}, \\
\bigcup_{n=1}^{\infty} A_n \in \mathcal{F}.
\end{gathered}
\]
\item \textit{Probability measure ($\Prob$)}: a function assigning to each event a number between zero and one, with total probability one:
\[
\begin{gathered}
\Prob: \mathcal{F} \to [0,1], \\
\Prob(\Omega) = 1.
\end{gathered}
\]
It is countably additive: for every sequence of pairwise disjoint events $A_n$, the probability of their union equals the sum of their probabilities:
\[
\Prob\!\left(\bigcup_{n=1}^{\infty} A_n\right)
= \sum_{n=1}^{\infty} \Prob(A_n).
\]
\end{enumerate}
\end{definition}

\begin{remark}[Why a $\sigma$-algebra?]
We cannot simply assign probabilities to \textit{all} subsets of $\Omega$ (for example, if $\Omega = \R$, the Axiom of Choice implies the existence of non-measurable sets like the Vitali set). The $\sigma$-algebra $\mathcal{F}$ restricts the domain of $\Prob$ to a collection of ``well-behaved'' sets. Closure under countable unions is strictly necessary because probability theory frequently involves limits of events (e.g., ``the event happens eventually''), and we must guarantee that the limit of a sequence of valid events is still a valid event to which we can assign a probability.
\end{remark}

\begin{definition}[Random Variable and Expectation]
A \textit{random variable} $S$ is a measurable function
\[
S: (\Omega, \mathcal{F}) \to (\R, \mathcal{B}(\R)),
\]
where $\mathcal{B}(\R)$ is the Borel $\sigma$-algebra on $\R$. Measurability means that for any Borel set $B$, the preimage
\[
S^{-1}(B) = \{\omega \in \Omega : S(\omega) \in B\}
\]
is an element of $\mathcal{F}$.
For the purposes of this chapter, we restrict attention to integrable
random variables, meaning that
\[
\E[|S|]=\int_\Omega |S(\omega)|\,\dd\Prob(\omega)<\infty.
\]
For such a random variable, the \textit{expectation} (or expected value)
is the Lebesgue integral
\begin{equation}
\E[S] = \int_\Omega S(\omega) \dd \Prob(\omega)
\end{equation}
and is finite. More generally, extended expectations can be defined from
the positive and negative parts when they are not both infinite.
\end{definition}

\begin{remark}[Why Measurability?]
We require $S$ to be measurable so that we can calculate probabilities of events defined by $S$. For instance, to compute $\Prob(S \leq x)$, we need the set $\{\omega : S(\omega) \leq x\}$ to be in $\mathcal{F}$ so that the measure $\Prob$ can be applied to it. Without measurability, the expectation integral would be undefined.
\end{remark}

Figure~\ref{fig:atlas-conditional-events} visualizes conditioning as restricting attention to a specified event.

\begin{figure}[!htbp]
\centering
\begin{tikzpicture}[scale=1.45,>=Stealth]
\draw[black!60] (-2.5,-1.6) rectangle (2.5,1.6);
\fill[figblue!15] (0.65,0) circle (1.25);
\begin{scope}
\clip (-0.65,0) circle (1.25);
\fill[figgreen!45] (0.65,0) circle (1.25);
\end{scope}
\draw[figred,thick] (-0.65,0) circle (1.25);
\draw[figblue,thick] (0.65,0) circle (1.25);
\node at (-1.25,0.25) {$A$};
\node at (1.25,0.25) {$B$};
\node at (0,0) {$A\cap B$};
\node at (2.2,1.3) {$\Omega$};
\end{tikzpicture}
\caption{The outer rectangle represents the sample space. Conditioning on $B$ restricts attention to the blue disk; its overlap with $A$ is highlighted in green. For $\Prob(B)>0$, the conditional probability is $\Prob(A\mid B)=\Prob(A\cap B)/\Prob(B)$. Areas are schematic and need not be proportional to probabilities.}
\label{fig:atlas-conditional-events}
\end{figure}

\subsection{Markov and Chebyshev Inequalities}

These inequalities provide deterministic bounds on the tails of probability distributions using only the moments (expectation and variance) of the random variable.

\begin{proposition}[Markov's Inequality]
Let $S$ be a nonnegative random variable ($S \geq 0$ almost surely) and let $a > 0$. Then:
\begin{equation}
\Prob(S \geq a) \leq \frac{\E[S]}{a} % <--- QUI USA \Prob
\end{equation}

\end{proposition}

\begin{proof}
We introduce the \textit{indicator function} $\mathbf{1}_{S \geq a}(\omega)$, which equals $1$ if $S(\omega) \geq a$ and $0$ otherwise.
Observe the pointwise inequality for all $\omega \in \Omega$:
\begin{equation}
S(\omega) \geq a \cdot \mathbf{1}_{S \geq a}(\omega)
\end{equation}

This holds because if $S(\omega) \geq a$, the right side is $a$, and $S(\omega) \geq a$; if $S(\omega) < a$, the right side is $0$, and $S(\omega) \geq 0$ by assumption.
Since the Lebesgue integral preserves inequalities for nonnegative functions, we integrate both sides with respect to $\Prob$ over $\Omega$:
\begin{equation}
\int_\Omega S(\omega) \dd \Prob(\omega) \geq \int_\Omega a \cdot \mathbf{1}_{S \geq a}(\omega) \dd \Prob(\omega)
\end{equation}

The left side is exactly $\E[S]$. On the right side, $a$ is a constant and can be pulled out. The integral of an indicator function over $\Omega$ is precisely the probability of the set it indicates:
\begin{equation}
\begin{aligned}
\E[S] &\geq a \int_\Omega \mathbf{1}_{S \geq a}(\omega) \dd \Prob(\omega) \\
&= a \cdot \Prob(S \geq a)
\end{aligned}
\end{equation}
Dividing both sides by $a > 0$ yields the result.
\end{proof}

\begin{definition}[Variance]
The \textit{variance} of a random variable $S$ with finite expectation $\mu = \E[S]$ is defined as the expected squared deviation from the mean:
\begin{equation}
\var(S) = \E[(S - \E[S])^2]
\end{equation}

\end{definition}

\begin{theorem}[Chebyshev's Inequality]
If $S$ has finite expectation $\mu$ and finite variance $\sigma^2 = \var(S) < \infty$, then for any $a > 0$:
\begin{equation}
\Prob(|S - \mu| \geq a) \leq \frac{\sigma^2}{a^2}
\end{equation}

\end{theorem}

\begin{proof}
Define a new random variable $X = (S - \mu)^2$. Since it is a square, $X \geq 0$ almost surely.
Notice that the event $\{|S - \mu| \geq a\}$ is mathematically equivalent to the event $\{(S - \mu)^2 \geq a^2\}$, which is exactly $\{X \geq a^2\}$.
We apply Markov's Inequality to the nonnegative random variable $X$ with the threshold $a^2 > 0$:
\begin{equation}
\Prob(X \geq a^2) \leq \frac{\E[X]}{a^2}
\end{equation}
Substituting back $X = (S - \mu)^2$ and recognizing that
\[
\begin{aligned}
\E[X] &= \E[(S - \mu)^2] \\
&= \var(S) \\
&= \sigma^2,
\end{aligned}
\]
we obtain:
\begin{equation}
\Prob(|S - \mu| \geq a) \leq \frac{\sigma^2}{a^2}
\end{equation}

\end{proof}

\section{Limit Theorems for Sums of Random Variables}

\subsection{Weak Law of Large Numbers}

The Weak Law of Large Numbers (WLLN) formalizes the intuitive idea that the average of a large number of independent observations converges to the expected value.

\begin{definition}[Independent and Identically Distributed (i.i.d.)]
A sequence of random variables $S_1, S_2, \dots$ is:
\begin{itemize}
\item \textit{Identically distributed}: if they all share the same cumulative distribution function, and consequently, the same expectation $\mu$ and variance $\sigma^2$.
\item \textit{Independent}: if the joint distribution of any finite subset is the product of their marginal distributions. Crucially for expectations, this implies that for $i \neq j$, the expectation of the product factors:

\[
\E[S_i S_j] = \E[S_i]\E[S_j].
\]
Consequently, their covariance is zero:
\[
\E[(S_i - \mu)(S_j - \mu)] = 0.
\]

\end{itemize}
\end{definition}

\begin{theorem}[Weak Law of Large Numbers]
Let $S_1, S_2, \dots$ be a sequence of i.i.d. random variables with mean $\mu$ and finite variance $\sigma^2$. Let the sample mean be
\[
z_m = (1/m) \sum_{i=1}^m S_i.
\]
Then for any $\varepsilon > 0$:
\begin{equation}
\lim_{m \to \infty} \Prob(|z_m - \mu| \geq \varepsilon) = 0
\end{equation}

\end{theorem}

\begin{proof}
We first explicitly calculate the expectation and variance of $z_m$.
By the linearity of the expectation operator:
\begin{equation}
\begin{aligned}
\E[z_m] &= \E\left[ \frac{1}{m} \sum_{i=1}^m S_i \right] \\
&= \frac{1}{m} \sum_{i=1}^m \E[S_i] \\
&= \frac{1}{m} (m \mu) \\
&= \mu
\end{aligned}
\end{equation}
Next, we calculate the variance of $z_m$. By definition:
\begin{equation}
\begin{aligned}
\var(z_m) &= \E[(z_m - \E[z_m])^2] \\
&= \E\left[ \left( \frac{1}{m} \sum_{i=1}^m (S_i - \mu) \right)^2 \right]
\end{aligned}
\end{equation}
Expanding the square of the sum yields a double sum:
\begin{equation}
\var(z_m) = \frac{1}{m^2} \E\left[ \sum_{i=1}^m \sum_{j=1}^m (S_i - \mu)(S_j - \mu) \right]
\end{equation}
Using the linearity of expectation again, we bring the expectation inside the sums:
\begin{equation}
\var(z_m) = \frac{1}{m^2} \sum_{i=1}^m \sum_{j=1}^m \E[(S_i - \mu)(S_j - \mu)]
\end{equation}
We evaluate the terms in the double sum based on the indices:
\begin{itemize}
\item \textit{Diagonal terms}: if $i = j$, the term is

\[
\begin{aligned}
\E[(S_i - \mu)^2] &= \var(S_i) \\
&= \sigma^2.
\end{aligned}
\]
There are exactly $m$ such terms on the diagonal of the double sum.
\item \textit{Off-diagonal terms}: if $i \neq j$, the term is $\E[(S_i - \mu)(S_j - \mu)]$. Because $S_i$ and $S_j$ are independent, this factors into

\[
\begin{aligned}
\E[S_i - \mu]\E[S_j - \mu] &= 0 \cdot 0 \\
&= 0.
\end{aligned}
\]
All $m^2 - m$ off-diagonal terms vanish.
\end{itemize}
Summing these up, we get:
\begin{equation}
\begin{aligned}
\var(z_m) &= \frac{1}{m^2} (m \sigma^2) \\
&= \frac{\sigma^2}{m}
\end{aligned}
\end{equation}
Now, we apply Chebyshev's Inequality to the random variable $z_m$ with threshold $\varepsilon > 0$:
\begin{equation}
\begin{aligned}
\Prob(|z_m - \mu| \geq \varepsilon) &\leq \frac{\var(z_m)}{\varepsilon^2} \\
&= \frac{\sigma^2}{m \varepsilon^2}
\end{aligned}
\end{equation}
Since $\sigma^2$ and $\varepsilon^2$ are fixed finite constants, taking the limit as $m \to \infty$ gives:
\begin{equation}
\lim_{m \to \infty} \frac{\sigma^2}{m \varepsilon^2} = 0
\end{equation}
Because probabilities are nonnegative, we conclude that
\[
\lim_{m \to \infty} \Prob(|z_m - \mu| \geq \varepsilon) = 0.
\]

\end{proof}

\subsection{Characteristic Functions and the Central Limit Theorem}


While the WLLN tells us the sample mean converges to a constant, the Central Limit Theorem (CLT) describes the exact distribution of the fluctuations around that constant.

\begin{definition}[Convergence in Distribution]
A sequence of random variables $Y_m$ \textit{converges in distribution} to a random variable $Y$, denoted $Y_m \xrightarrow{d} Y$, if their cumulative distribution functions (CDFs) $F_{Y_m}(y)$ converge to $F_Y(y)$ at every point $y$ where $F_Y(y)$ is continuous.
\end{definition}

\begin{remark}[Why use Characteristic Functions?]
Proving convergence of CDFs directly is analytically intractable. Instead, we use the \textit{Characteristic Function} (CF), which is the Fourier transform of the probability measure. The CF always exists, uniquely determines the probability distribution, and converts the convolution of independent random variables into simple algebraic multiplication, making limit calculations straightforward.
\end{remark}

\begin{definition}[Characteristic Function]
The characteristic function of a random variable $S$ is defined as:
\begin{equation}
\varphi_S(t) = \E[e^{itS}] = \int_\Omega e^{itS(\omega)} \dd \Prob(\omega), \quad t \in \R
\end{equation}

\end{definition}

\begin{theorem}[Lindeberg--Lévy Central Limit Theorem]
Let $S_1, S_2, \dots$ be a sequence of i.i.d. random variables with mean $\mu$ and finite variance $\sigma^2 > 0$. Let
\[
z_m = (1/m) \sum_{i=1}^m S_i.
\]

Then the normalized sample mean converges in distribution to a standard normal random variable:
\begin{equation}
Y_m = \frac{\sqrt{m}(z_m - \mu)}{\sigma} \xrightarrow{d} \mathcal{N}(0, 1)
\end{equation}

\end{theorem}

\subsection{Proof of the Central Limit Theorem}

\begin{proof}
We will prove this by showing that the characteristic function of $Y_m$ converges pointwise to the characteristic function of the standard normal distribution $\mathcal{N}(0,1)$.

\textit{Step 1: Express $Y_m$ as a sum of centered variables.}
Let $X_i = S_i - \mu$. The variables $X_i$ are i.i.d. with $\E[X_i] = 0$ and $\var(X_i) = \sigma^2$. We can rewrite $Y_m$ as:
\begin{equation}
\begin{aligned}
Y_m &= \frac{\sqrt{m}}{\sigma} \left( \frac{1}{m} \sum_{i=1}^m X_i \right) \\
&= \frac{1}{\sigma\sqrt{m}} \sum_{i=1}^m X_i
\end{aligned}
\end{equation}

\textit{Step 2: Relate the CF of $Y_m$ to the CF of $X_1$.}
Using the properties of characteristic functions:
\begin{itemize}
\item \textit{Scaling}: the following identity holds:

\[
\begin{aligned}
\varphi_{cX}(t) &= \E[e^{it(cX)}] \\
&= \varphi_X(ct).
\end{aligned}
\]

\item \textit{Independence}: if $X$ and $Y$ are independent,

\[
\begin{aligned}
\varphi_{X+Y}(t) &= \E[e^{itX}e^{itY}] \\
&= \E[e^{itX}]\E[e^{itY}] \\
&= \varphi_X(t)\varphi_Y(t).
\end{aligned}
\]

\end{itemize}
Applying these properties to $Y_m$:
\begin{equation}
\begin{aligned}
\varphi_{Y_m}(t) &= \varphi_{\sum X_i / (\sigma\sqrt{m})}(t) \\
&= \left[ \varphi_{X_1}\left( \frac{t}{\sigma\sqrt{m}} \right) \right]^m
\end{aligned}
\end{equation}

\textit{Step 3: Taylor expand the characteristic function of $X_1$.}
We expand $\varphi_{X_1}(u)$ around $u=0$. To do this rigorously, we differentiate under the integral sign. This is justified by the Dominated Convergence Theorem because the derivatives of the integrand are bounded by $|X_1|^k$, and we assume $\E[|X_1|^2] < \infty$.
\begin{align*}
\varphi_{X_1}(u) &= \E[e^{iuX_1}] \\
    \varphi_{X_1}'(u) &= \E[iX_1 e^{iuX_1}] \implies \varphi_{X_1}'(0) = i\E[X_1] = 0 \\
    \varphi_{X_1}''(u) &= \E[(iX_1)^2 e^{iuX_1}] = -\E[X_1^2 e^{iuX_1}] \implies \varphi_{X_1}''(0) = -\E[X_1^2] = -\var(X_1) = -\sigma^2
\end{align*}
By Taylor's theorem with the Peano form of the remainder, we have:
\begin{equation}
\begin{aligned}
\varphi_{X_1}(u) &= \varphi_{X_1}(0) + u\varphi_{X_1}'(0) + \frac{u^2}{2}\varphi_{X_1}''(0) + o(u^2) \\
&= 1 - \frac{\sigma^2 u^2}{2} + o(u^2)
\end{aligned}
\end{equation}
where $o(u^2)$ denotes a function such that
\[
\lim_{u \to 0} o(u^2)/u^2 = 0.
\]

\textit{Step 4: Substitute and take the limit.}
We substitute
\[
u = (t/(\sigma\sqrt{m}))
\]
into the Taylor expansion. Note that as $m \to \infty$, $u \to 0$:
\begin{equation}
\begin{aligned}
\varphi_{X_1}\left( \frac{t}{\sigma\sqrt{m}} \right) &= 1 - \frac{\sigma^2}{2} \left( \frac{t}{\sigma\sqrt{m}} \right)^2 + o\left( \frac{t^2}{\sigma^2 m} \right) \\
&= 1 - \frac{t^2}{2m} + o\left( \frac{1}{m} \right)
\end{aligned}
\end{equation}
Now, we substitute this back into the expression for $\varphi_{Y_m}(t)$:
\begin{equation}
\varphi_{Y_m}(t) = \left[ 1 - \frac{t^2}{2m} + o\left( \frac{1}{m} \right) \right]^m
\end{equation}

To evaluate the limit as $m \to \infty$, we rewrite the expression using the exponential and natural logarithm:
\begin{equation}
\varphi_{Y_m}(t) = \exp\left( m \ln\left( 1 - \frac{t^2}{2m} + o\left( \frac{1}{m} \right) \right) \right)
\end{equation}
Using the Taylor expansion $\ln(1 + x) = x + o(x)$ for small $x$, where
\[
x = -(t^2/(2m)) + o( (1/m) ):
\]

\begin{equation}
\begin{aligned}
m \ln\left( 1 - \frac{t^2}{2m} + o\left( \frac{1}{m} \right) \right) &= m \left( -\frac{t^2}{2m} + o\left( \frac{1}{m} \right) \right) \\
&= -\frac{t^2}{2} + m \cdot o\left( \frac{1}{m} \right)
\end{aligned}
\end{equation}
The remainder term satisfies
\[
\lim_{m \to \infty} m \cdot o(1/m) = 0,
\]
Therefore, the exponent converges to $-t^2/2$. Therefore:
\begin{equation}
\lim_{m \to \infty} \varphi_{Y_m}(t) = e^{-t^2/2}
\end{equation}

\textit{Step 5: Apply Lévy's Continuity Theorem.}
\textit{Lévy's Continuity Theorem} states that if a sequence of characteristic functions $\varphi_{Y_m}(t)$ converges pointwise to a limit function $\varphi(t)$ that is continuous at $t=0$, then $Y_m$ converges in distribution to a random variable $Y$ whose characteristic function is $\varphi(t)$.
The limit function $e^{-t^2/2}$ is continuous everywhere. Furthermore, it is exactly the characteristic function of the standard normal distribution $\mathcal{N}(0,1)$.
Thus, we conclude that $Y_m \xrightarrow{d} \mathcal{N}(0, 1)$. \end{proof}

\subsection{Normal Approximation and Interpretation}

Figure~\ref{fig:ch16-clt} illustrates the approach to a normal distribution described by the Central Limit Theorem:
\[
z_m \xrightarrow{d} N(0,1).
\]

\begin{figure}[!htbp]
\centering
\begin{tikzpicture}
\begin{axis}[
  width=13cm,height=8cm, axis lines=left,
  xlabel={$z=((\sqrt{m}(\bar X_m-\mu))/\sigma)$}, ylabel={density},
  xmin=-4,xmax=4,ymin=0,ymax=0.62,
  legend pos=north east, legend style={font=\scriptsize, draw=black!30}
  ]
  \addplot[black!45,thick,const plot,domain=-4:4,samples=9]
    coordinates {(-3.5,0.05)(-2.5,0.09)(-1.5,0.22)(-0.5,0.4)(0.5,0.4)(1.5,0.22)(2.5,0.09)(3.5,0.05)};
  \addlegendentry{histogram of $z_m$, moderate $m$}
  \addplot[figred,very thick,domain=-4:4,samples=150] {1/sqrt(2*pi)*exp(-x^2/2)};
  \addlegendentry{$N(0,1)$ limit}
\end{axis}
\end{tikzpicture}
\caption[Central Limit Theorem]{Illustrative histogram of a standardized sample mean for a moderate sample size, compared with the standard normal limiting density (red). The horizontal coordinate is $z=\sqrt{m}(\bar X_m-\mu)/\sigma$.}
\label{fig:ch16-clt}
\end{figure}

\section{Stochastic Differential Equations}
\label{sec:stochastic-differential-equations}

The ordinary differential equations studied in the preceding chapter describe deterministic systems: under existence and uniqueness hypotheses, an initial condition determines a trajectory on its interval of existence. Many biophysical systems, however, experience intrinsic fluctuations, such as thermal noise, fluctuations in molecular counts in a genetic circuit, or synaptic noise in a neural network. To model these effects, we extend $\dot{\bv}=f(\bv)$ by introducing a random driving term. The probability concepts and Gaussian limits developed in this chapter provide the background for this extension.

\subsection{Brownian Motion}

\begin{definition}[Wiener Process]
\label{def:wiener-process}
A \emph{standard Wiener process}, or \emph{standard Brownian motion}, is a real-valued stochastic process $W(t)$, $t\ge0$, with the following properties:
\begin{enumerate}
\item \textit{Initial value}: $W(0)=0$ almost surely.
\item \textit{Independent increments}: it has independent increments: for any $0\leq t_0<t_1<\cdots<t_m$, the random variables $W(t_j)-W(t_{j-1})$, $j=1,\dots,m$, are jointly independent.
\item \textit{Gaussian increments}: for $0\leq s<t$, the increment $W(t)-W(s)$ has distribution $\mathcal{N}(0,t-s)$.
\item \textit{Continuity}: its sample paths $t\mapsto W(t)$ are almost surely continuous.
\end{enumerate}
\end{definition}

\begin{remark}[Continuity without Differentiability]
Brownian paths are almost surely nowhere differentiable, a result we do not prove here. The scale of an increment gives some intuition: over an interval of length $\Delta t$, its standard deviation is $(\Delta t)^{1/2}$, so a typical difference quotient has magnitude of order
\[
\frac{\sqrt{\Delta t}}{\Delta t}=\frac{1}{\sqrt{\Delta t}},
\]
which grows without bound as $\Delta t\to0$. This scaling is a heuristic, not a proof of nowhere differentiability. It explains why $\dot W$ cannot be treated as an ordinary function and why an equation written informally as $\dot{\bv}=f(\bv)+\dot{\bW}$ is not a classical ODE. A rigorous treatment requires stochastic integration.
\end{remark}

Figure~\ref{fig:atlas-brownian-paths} compares independent discretized realizations of Brownian motion.

\begin{figure}[!htbp]
\centering
\begin{tikzpicture}[>=Stealth]
\begin{axis}[
  width=12.5cm,height=6.2cm,
  axis lines=left,
  ticklabel style={font=\small},
  label style={font=\small},
  xlabel={$t$},ylabel={$W_t$},
  xmin=0,xmax=1.03,ymin=-1.4,ymax=2.2,
  xtick={0,0.25,0.5,0.75,1},
  grid=major,grid style={black!10}
]
\addplot[figblue,thin] coordinates {(0.000000,0.000000) (0.003906,0.071894) (0.007812,0.099667) (0.011719,0.161176) (0.015625,0.276561) (0.019531,0.303321) (0.023438,0.274530) (0.027344,0.308431) (0.031250,0.357354) (0.035156,0.294141) (0.039062,0.318164) (0.042969,0.348022) (0.046875,0.360721) (0.050781,0.378277) (0.054688,0.370466) (0.058594,0.371051) (0.062500,0.385740) (0.066406,0.327751) (0.070312,0.381763) (0.074219,0.372351) (0.078125,0.331046) (0.082031,0.286042) (0.085938,0.360965) (0.089844,0.381335) (0.093750,0.303530) (0.097656,0.275398) (0.101562,0.209094) (0.105469,0.158379) (0.109375,0.180587) (0.113281,0.226890) (0.117188,0.279170) (0.121094,0.268315) (0.125000,0.255850) (0.128906,0.201120) (0.132812,0.233901) (0.136719,0.157671) (0.140625,0.161345) (0.144531,0.182987) (0.148438,0.094929) (0.152344,0.174125) (0.156250,0.086354) (0.160156,0.102647) (0.164062,0.134677) (0.167969,0.138649) (0.171875,0.130281) (0.175781,0.164821) (0.179688,0.086271) (0.183594,0.114667) (0.187500,0.111874) (0.191406,0.138953) (0.195312,0.263080) (0.199219,0.174867) (0.203125,0.063584) (0.207031,0.158739) (0.210938,0.105439) (0.214844,0.201024) (0.218750,0.122955) (0.222656,0.242608) (0.226562,0.425723) (0.230469,0.408087) (0.234375,0.487504) (0.238281,0.499184) (0.242188,0.464558) (0.246094,0.438329) (0.250000,0.380200) (0.253906,0.392881) (0.257812,0.340354) (0.261719,0.409444) (0.265625,0.420920) (0.269531,0.352419) (0.273438,0.292158) (0.277344,0.358440) (0.281250,0.412231) (0.285156,0.439973) (0.289062,0.475860) (0.292969,0.483186) (0.296875,0.583534) (0.300781,0.560372) (0.304688,0.573706) (0.308594,0.669594) (0.312500,0.560663) (0.316406,0.596844) (0.320312,0.564603) (0.324219,0.561020) (0.328125,0.548369) (0.332031,0.475480) (0.335938,0.498073) (0.339844,0.475058) (0.343750,0.397360) (0.347656,0.301083) (0.351562,0.328202) (0.355469,0.415408) (0.359375,0.391643) (0.363281,0.320465) (0.367188,0.429462) (0.371094,0.427974) (0.375000,0.430546) (0.378906,0.425323) (0.382812,0.415855) (0.386719,0.460336) (0.390625,0.312285) (0.394531,0.224310) (0.398438,0.206260) (0.402344,0.267705) (0.406250,0.274579) (0.410156,0.080395) (0.414062,0.118034) (0.417969,-0.060501) (0.421875,-0.165487) (0.425781,-0.184497) (0.429688,-0.250202) (0.433594,-0.270557) (0.437500,-0.355921) (0.441406,-0.344766) (0.445312,-0.324919) (0.449219,-0.394724) (0.453125,-0.380942) (0.457031,-0.383763) (0.460938,-0.469585) (0.464844,-0.544962) (0.468750,-0.531145) (0.472656,-0.482526) (0.476562,-0.457292) (0.480469,-0.369933) (0.484375,-0.267920) (0.488281,-0.277632) (0.492188,-0.265407) (0.496094,-0.215602) (0.500000,-0.237320) (0.503906,-0.175597) (0.507812,-0.204590) (0.511719,-0.197919) (0.515625,-0.181113) (0.519531,-0.170213) (0.523438,-0.179235) (0.527344,-0.161139) (0.531250,-0.126996) (0.535156,-0.021619) (0.539062,-0.088090) (0.542969,-0.140985) (0.546875,-0.182087) (0.550781,-0.176136) (0.554688,-0.188205) (0.558594,-0.081948) (0.562500,-0.227149) (0.566406,-0.233622) (0.570312,-0.164457) (0.574219,-0.093145) (0.578125,-0.221449) (0.582031,-0.229593) (0.585938,-0.293614) (0.589844,-0.381834) (0.593750,-0.425104) (0.597656,-0.473408) (0.601562,-0.495901) (0.605469,-0.378299) (0.609375,-0.369934) (0.613281,-0.389069) (0.617188,-0.429705) (0.621094,-0.421121) (0.625000,-0.449874) (0.628906,-0.522024) (0.632812,-0.535298) (0.636719,-0.635955) (0.640625,-0.572943) (0.644531,-0.687222) (0.648438,-0.587140) (0.652344,-0.451273) (0.656250,-0.443199) (0.660156,-0.377786) (0.664062,-0.307570) (0.667969,-0.347673) (0.671875,-0.348133) (0.675781,-0.411539) (0.679688,-0.490655) (0.683594,-0.511667) (0.687500,-0.476269) (0.691406,-0.493906) (0.695312,-0.523734) (0.699219,-0.571607) (0.703125,-0.628278) (0.707031,-0.516322) (0.710938,-0.513697) (0.714844,-0.521439) (0.718750,-0.585341) (0.722656,-0.568698) (0.726562,-0.556772) (0.730469,-0.572136) (0.734375,-0.494649) (0.738281,-0.466995) (0.742188,-0.517893) (0.746094,-0.593066) (0.750000,-0.598301) (0.753906,-0.552629) (0.757812,-0.579400) (0.761719,-0.509875) (0.765625,-0.553263) (0.769531,-0.612072) (0.773438,-0.579023) (0.777344,-0.469131) (0.781250,-0.525759) (0.785156,-0.489410) (0.789062,-0.438775) (0.792969,-0.467512) (0.796875,-0.479576) (0.800781,-0.553124) (0.804688,-0.629102) (0.808594,-0.543822) (0.812500,-0.531688) (0.816406,-0.504896) (0.820312,-0.538428) (0.824219,-0.446751) (0.828125,-0.431713) (0.832031,-0.554976) (0.835938,-0.482365) (0.839844,-0.503309) (0.843750,-0.530958) (0.847656,-0.628338) (0.851562,-0.544194) (0.855469,-0.557405) (0.859375,-0.615450) (0.863281,-0.507438) (0.867188,-0.503442) (0.871094,-0.575805) (0.875000,-0.616645) (0.878906,-0.635226) (0.882812,-0.694043) (0.886719,-0.680323) (0.890625,-0.749864) (0.894531,-0.707605) (0.898438,-0.778608) (0.902344,-0.700066) (0.906250,-0.804586) (0.910156,-0.893449) (0.914062,-0.811714) (0.917969,-0.735168) (0.921875,-0.809518) (0.925781,-0.833310) (0.929688,-0.856793) (0.933594,-0.864099) (0.937500,-0.908484) (0.941406,-0.889948) (0.945312,-0.905708) (0.949219,-0.916522) (0.953125,-0.974218) (0.957031,-1.024776) (0.960938,-1.200692) (0.964844,-1.166844) (0.968750,-1.107111) (0.972656,-1.082881) (0.976562,-1.115680) (0.980469,-1.058859) (0.984375,-1.035515) (0.988281,-1.119239) (0.992188,-1.118584) (0.996094,-1.126853) (1.000000,-1.262117)};
\addplot[figred,thin] coordinates {(0.000000,0.000000) (0.003906,-0.029506) (0.007812,-0.065051) (0.011719,-0.028742) (0.015625,-0.096435) (0.019531,-0.150042) (0.023438,-0.154057) (0.027344,-0.168918) (0.031250,-0.175222) (0.035156,-0.111769) (0.039062,-0.149804) (0.042969,-0.256598) (0.046875,-0.215331) (0.050781,-0.184882) (0.054688,-0.136478) (0.058594,-0.206721) (0.062500,-0.323506) (0.066406,-0.295969) (0.070312,-0.286195) (0.074219,-0.287495) (0.078125,-0.240818) (0.082031,-0.220457) (0.085938,-0.252248) (0.089844,-0.299812) (0.093750,-0.394050) (0.097656,-0.343627) (0.101562,-0.357341) (0.105469,-0.445727) (0.109375,-0.444332) (0.113281,-0.461245) (0.117188,-0.417181) (0.121094,-0.500581) (0.125000,-0.463202) (0.128906,-0.439790) (0.132812,-0.430240) (0.136719,-0.400264) (0.140625,-0.431789) (0.144531,-0.388948) (0.148438,-0.337485) (0.152344,-0.372051) (0.156250,-0.461285) (0.160156,-0.440345) (0.164062,-0.441948) (0.167969,-0.385368) (0.171875,-0.409401) (0.175781,-0.456165) (0.179688,-0.458066) (0.183594,-0.504559) (0.187500,-0.511522) (0.191406,-0.516361) (0.195312,-0.541767) (0.199219,-0.561876) (0.203125,-0.598022) (0.207031,-0.619152) (0.210938,-0.626959) (0.214844,-0.671837) (0.218750,-0.730047) (0.222656,-0.724819) (0.226562,-0.642867) (0.230469,-0.582404) (0.234375,-0.552744) (0.238281,-0.603631) (0.242188,-0.672372) (0.246094,-0.637026) (0.250000,-0.578936) (0.253906,-0.464766) (0.257812,-0.381258) (0.261719,-0.292269) (0.265625,-0.324989) (0.269531,-0.314480) (0.273438,-0.264114) (0.277344,-0.264684) (0.281250,-0.248511) (0.285156,-0.195882) (0.289062,-0.061264) (0.292969,-0.043557) (0.296875,0.034533) (0.300781,0.063224) (0.304688,0.048614) (0.308594,-0.004877) (0.312500,0.037370) (0.316406,0.119303) (0.320312,0.105191) (0.324219,0.061354) (0.328125,0.134040) (0.332031,0.204421) (0.335938,0.215187) (0.339844,0.178180) (0.343750,0.142985) (0.347656,0.083378) (0.351562,0.060888) (0.355469,0.164955) (0.359375,0.175233) (0.363281,0.143365) (0.367188,0.253995) (0.371094,0.250465) (0.375000,0.214465) (0.378906,0.201689) (0.382812,0.186901) (0.386719,0.224599) (0.390625,0.167184) (0.394531,0.107884) (0.398438,0.027083) (0.402344,0.138455) (0.406250,0.201816) (0.410156,0.175921) (0.414062,0.225826) (0.417969,0.086679) (0.421875,0.026409) (0.425781,-0.010201) (0.429688,0.067482) (0.433594,0.054060) (0.437500,0.114557) (0.441406,0.115057) (0.445312,0.200676) (0.449219,0.154276) (0.453125,0.141920) (0.457031,0.135491) (0.460938,0.233644) (0.464844,0.149857) (0.468750,0.111019) (0.472656,0.079022) (0.476562,0.136081) (0.480469,0.138743) (0.484375,0.258287) (0.488281,0.235928) (0.492188,0.275948) (0.496094,0.211159) (0.500000,0.214171) (0.503906,0.183368) (0.507812,0.210278) (0.511719,0.147170) (0.515625,0.182961) (0.519531,0.171410) (0.523438,0.112894) (0.527344,0.099114) (0.531250,0.214013) (0.535156,0.335778) (0.539062,0.305042) (0.542969,0.347105) (0.546875,0.334276) (0.550781,0.428011) (0.554688,0.299347) (0.558594,0.327492) (0.562500,0.341076) (0.566406,0.355604) (0.570312,0.408944) (0.574219,0.460406) (0.578125,0.549397) (0.582031,0.583684) (0.585938,0.655660) (0.589844,0.656731) (0.593750,0.696730) (0.597656,0.650395) (0.601562,0.732357) (0.605469,0.769263) (0.609375,0.792091) (0.613281,0.787495) (0.617188,0.773844) (0.621094,0.784915) (0.625000,0.887730) (0.628906,0.881640) (0.632812,0.867941) (0.636719,0.877520) (0.640625,0.968315) (0.644531,0.894307) (0.648438,0.785656) (0.652344,0.900927) (0.656250,0.921730) (0.660156,1.029922) (0.664062,1.089172) (0.667969,1.062291) (0.671875,1.151485) (0.675781,1.067422) (0.679688,1.137168) (0.683594,1.140709) (0.687500,1.175807) (0.691406,1.223792) (0.695312,1.171859) (0.699219,1.272209) (0.703125,1.208656) (0.707031,1.209317) (0.710938,1.298264) (0.714844,1.344554) (0.718750,1.249169) (0.722656,1.303576) (0.726562,1.314609) (0.730469,1.427168) (0.734375,1.499224) (0.738281,1.584170) (0.742188,1.588172) (0.746094,1.744125) (0.750000,1.755908) (0.753906,1.652448) (0.757812,1.640733) (0.761719,1.723463) (0.765625,1.781676) (0.769531,1.771427) (0.773438,1.682514) (0.777344,1.575211) (0.781250,1.517077) (0.785156,1.559594) (0.789062,1.452903) (0.792969,1.511255) (0.796875,1.529966) (0.800781,1.509829) (0.804688,1.380237) (0.808594,1.473233) (0.812500,1.586517) (0.816406,1.617732) (0.820312,1.682910) (0.824219,1.745530) (0.828125,1.724981) (0.832031,1.681721) (0.835938,1.727572) (0.839844,1.729809) (0.843750,1.684408) (0.847656,1.813063) (0.851562,1.883828) (0.855469,1.899424) (0.859375,1.907413) (0.863281,1.908964) (0.867188,1.894269) (0.871094,1.894773) (0.875000,1.891761) (0.878906,1.846694) (0.882812,1.835272) (0.886719,1.862396) (0.890625,1.847036) (0.894531,1.864544) (0.898438,1.938020) (0.902344,2.006924) (0.906250,1.944791) (0.910156,1.973791) (0.914062,1.919343) (0.917969,1.903685) (0.921875,1.882549) (0.925781,1.800078) (0.929688,1.880092) (0.933594,1.781018) (0.937500,1.761642) (0.941406,1.770016) (0.945312,1.776495) (0.949219,1.791739) (0.953125,1.819722) (0.957031,1.777823) (0.960938,1.835601) (0.964844,1.903809) (0.968750,1.960148) (0.972656,1.932963) (0.976562,1.972191) (0.980469,2.030314) (0.984375,1.886979) (0.988281,1.835040) (0.992188,1.912010) (0.996094,1.848262) (1.000000,1.880328)};
\addplot[figgreen,thin] coordinates {(0.000000,0.000000) (0.003906,-0.025111) (0.007812,-0.005208) (0.011719,0.057699) (0.015625,0.103643) (0.019531,0.175262) (0.023438,0.300538) (0.027344,0.364271) (0.031250,0.297284) (0.035156,0.293449) (0.039062,0.276558) (0.042969,0.297778) (0.046875,0.352644) (0.050781,0.336403) (0.054688,0.341060) (0.058594,0.390238) (0.062500,0.399815) (0.066406,0.477564) (0.070312,0.505813) (0.074219,0.594553) (0.078125,0.651313) (0.082031,0.720662) (0.085938,0.718457) (0.089844,0.778104) (0.093750,0.712839) (0.097656,0.785709) (0.101562,0.774404) (0.105469,0.772335) (0.109375,0.736525) (0.113281,0.756319) (0.117188,0.810040) (0.121094,0.774600) (0.125000,0.772764) (0.128906,0.834047) (0.132812,0.813085) (0.136719,0.699305) (0.140625,0.766272) (0.144531,0.770126) (0.148438,0.771713) (0.152344,0.898191) (0.156250,0.879438) (0.160156,0.888041) (0.164062,0.973387) (0.167969,1.106829) (0.171875,1.238212) (0.175781,1.237260) (0.179688,1.187050) (0.183594,1.240801) (0.187500,1.131164) (0.191406,1.132831) (0.195312,1.152186) (0.199219,1.152715) (0.203125,1.123587) (0.207031,1.175339) (0.210938,1.197038) (0.214844,1.120756) (0.218750,1.164631) (0.222656,1.132019) (0.226562,1.170531) (0.230469,1.247165) (0.234375,1.301283) (0.238281,1.277580) (0.242188,1.397103) (0.246094,1.405889) (0.250000,1.368024) (0.253906,1.288213) (0.257812,1.336854) (0.261719,1.276878) (0.265625,1.348805) (0.269531,1.389297) (0.273438,1.331032) (0.277344,1.335090) (0.281250,1.400331) (0.285156,1.282302) (0.289062,1.294682) (0.292969,1.122304) (0.296875,1.109572) (0.300781,1.082782) (0.304688,1.004933) (0.308594,0.931148) (0.312500,0.881930) (0.316406,0.901384) (0.320312,0.930065) (0.324219,0.970049) (0.328125,1.015457) (0.332031,1.052842) (0.335938,1.016752) (0.339844,1.014984) (0.343750,0.988403) (0.347656,1.044343) (0.351562,0.925725) (0.355469,0.878376) (0.359375,0.810663) (0.363281,0.758180) (0.367188,0.863064) (0.371094,0.797596) (0.375000,0.932530) (0.378906,0.903269) (0.382812,0.869623) (0.386719,0.853793) (0.390625,0.861174) (0.394531,0.844708) (0.398438,0.838100) (0.402344,0.881782) (0.406250,1.028804) (0.410156,0.960426) (0.414062,0.926870) (0.417969,0.902357) (0.421875,0.966700) (0.425781,0.947902) (0.429688,0.978309) (0.433594,1.044760) (0.437500,1.044206) (0.441406,1.006303) (0.445312,1.018223) (0.449219,1.023594) (0.453125,1.003573) (0.457031,0.978695) (0.460938,0.923673) (0.464844,0.843002) (0.468750,0.818019) (0.472656,0.832027) (0.476562,0.812912) (0.480469,0.749909) (0.484375,0.755257) (0.488281,0.773499) (0.492188,0.821470) (0.496094,0.822808) (0.500000,0.804076) (0.503906,0.812094) (0.507812,0.876290) (0.511719,0.736135) (0.515625,0.825891) (0.519531,0.728027) (0.523438,0.644779) (0.527344,0.716917) (0.531250,0.686794) (0.535156,0.736577) (0.539062,0.569293) (0.542969,0.488366) (0.546875,0.465007) (0.550781,0.482931) (0.554688,0.447180) (0.558594,0.470264) (0.562500,0.486013) (0.566406,0.484168) (0.570312,0.517742) (0.574219,0.412646) (0.578125,0.392537) (0.582031,0.421263) (0.585938,0.424241) (0.589844,0.474981) (0.593750,0.500363) (0.597656,0.461650) (0.601562,0.452201) (0.605469,0.482847) (0.609375,0.472603) (0.613281,0.520173) (0.617188,0.577951) (0.621094,0.568870) (0.625000,0.569530) (0.628906,0.656062) (0.632812,0.636640) (0.636719,0.672636) (0.640625,0.626296) (0.644531,0.590166) (0.648438,0.640416) (0.652344,0.636585) (0.656250,0.601023) (0.660156,0.497326) (0.664062,0.456873) (0.667969,0.351289) (0.671875,0.321702) (0.675781,0.345032) (0.679688,0.293851) (0.683594,0.398429) (0.687500,0.408629) (0.691406,0.374026) (0.695312,0.417987) (0.699219,0.427224) (0.703125,0.517577) (0.707031,0.581326) (0.710938,0.665200) (0.714844,0.642112) (0.718750,0.542446) (0.722656,0.504294) (0.726562,0.513517) (0.730469,0.572497) (0.734375,0.589831) (0.738281,0.553430) (0.742188,0.570038) (0.746094,0.610921) (0.750000,0.732581) (0.753906,0.773528) (0.757812,0.837833) (0.761719,0.895476) (0.765625,0.852360) (0.769531,0.870582) (0.773438,0.892371) (0.777344,0.767355) (0.781250,0.723896) (0.785156,0.738528) (0.789062,0.827874) (0.792969,0.786093) (0.796875,0.739173) (0.800781,0.746967) (0.804688,0.863414) (0.808594,0.847187) (0.812500,0.827386) (0.816406,0.762280) (0.820312,0.815711) (0.824219,0.797078) (0.828125,0.861175) (0.832031,0.745486) (0.835938,0.736451) (0.839844,0.889338) (0.843750,0.934403) (0.847656,0.953715) (0.851562,0.956925) (0.855469,0.852220) (0.859375,0.793411) (0.863281,0.775036) (0.867188,0.796473) (0.871094,0.790292) (0.875000,0.808401) (0.878906,0.822130) (0.882812,0.789090) (0.886719,0.668311) (0.890625,0.682052) (0.894531,0.677788) (0.898438,0.667126) (0.902344,0.701464) (0.906250,0.699755) (0.910156,0.627214) (0.914062,0.559267) (0.917969,0.536071) (0.921875,0.591080) (0.925781,0.648641) (0.929688,0.601301) (0.933594,0.613979) (0.937500,0.549664) (0.941406,0.529575) (0.945312,0.500962) (0.949219,0.462865) (0.953125,0.486680) (0.957031,0.639620) (0.960938,0.688515) (0.964844,0.570401) (0.968750,0.670939) (0.972656,0.624456) (0.976562,0.649979) (0.980469,0.641179) (0.984375,0.665896) (0.988281,0.726512) (0.992188,0.772373) (0.996094,0.763951) (1.000000,0.743884)};
\end{axis}
\end{tikzpicture}
\caption{Three reproducibly simulated paths start at zero and use independent Gaussian increments of variance $\Delta t=1/256$. Straight segments interpolate the sampled values for display. These finite-resolution drawings do not imply differentiability of Brownian paths and are not experimental observations.}
\label{fig:atlas-brownian-paths}
\end{figure}

\subsection{The Langevin Equation}

\begin{definition}[It\^o Stochastic Differential Equation]
\label{def:langevin-equation}
An It\^o \emph{Langevin equation} has the form
\begin{equation}
\dd\bv(t)=f(\bv(t))\,\dd t+\sigma(\bv(t))\,\dd\bW(t),
\end{equation}
where $\bv(t)\in\R^d$, $f:\R^d\to\R^d$ is the deterministic \emph{drift}, $\bW$ is an $m$-dimensional Wiener process with independent components, and $\sigma:\R^d\to\R^{d\times m}$ is the noise-amplitude matrix. The diffusion covariance is $\sigma\sigma^{\top}$. The equation is interpreted in integral form:
\begin{equation}
\bv(t)=\bv(0)+\int_0^t f(\bv(s))\,\dd s
+\int_0^t\sigma(\bv(s))\,\dd\bW(s).
\end{equation}

The last term is an \emph{It\^o stochastic integral}. For a suitable adapted integrand, it is constructed as a limit of non-anticipating sums; for regular integrands, these take the left-endpoint form
\[
\sum_j\sigma(\bv(t_j))\big(\bW(t_{j+1})-\bW(t_j)\big).
\]

Unlike ordinary Riemann integration, the evaluation convention matters. A different convention, such as Stratonovich integration, generally changes the drift when $\sigma$ depends on the state.
\end{definition}

For example, global Lipschitz and linear-growth assumptions on $f$ and $\sigma$, together with a square-integrable initial state independent of future Brownian increments, ensure a unique global strong solution. We do not develop the existence theory here; these assumptions indicate the role played by regularity conditions, analogous to those required for deterministic ODEs.

\begin{example}[The Ornstein--Uhlenbeck Process]
\label{ex:ornstein-uhlenbeck}
In one dimension, linear relaxation toward zero with rate $\gamma>0$ and constant noise amplitude $\sigma>0$ gives
\begin{equation}
\dd v=-\gamma v\,\dd t+\sigma\,\dd W.
\end{equation}

This is the \emph{Ornstein--Uhlenbeck process}, the classical Langevin model for a particle velocity subject to viscous friction and random molecular impacts. Its drift is the same linear relaxation encountered in deterministic first-order equations, while its noise term represents the fluctuating forcing. It also describes the subthreshold dynamics of a leaky integrate-and-fire neuron driven by Gaussian noise after shifting the deterministic equilibrium to zero; a full spiking model additionally requires a threshold and reset rule.
\end{example}

\section{Diffusion Equations and Stationary Probability Flows}

\subsection{The Fokker--Planck Equation and Infinitesimal Moments}


Instead of following a single random trajectory, we may study the time-dependent probability density $p(v,t)$ of the state across all noise realizations. Under suitable assumptions, this density satisfies a deterministic PDE.

\begin{theorem}[Fokker--Planck Equation for an It\^o Diffusion]
\label{thm:fokker-planck}
Let $v(t)$ solve the one-dimensional It\^o equation
\[
\dd v=f(v)\,\dd t+\sigma(v)\,\dd W.
\]

Assume that the coefficients and the probability density $p(v,t)$ are sufficiently regular for the following derivatives to exist classically. Then
\begin{equation}
\frac{\partial p}{\partial t}
=-\frac{\partial}{\partial v}\big[f(v)p\big]
+\frac12\frac{\partial^2}{\partial v^2}\big[\sigma(v)^2p\big].
\end{equation}

With less regularity, the equation is interpreted weakly; an arbitrary stochastic differential equation need not have a smooth density.
\end{theorem}

\begin{proof}[Heuristic Derivation via the Kramers--Moyal Expansion]
\leavevmode\par\nobreak
\textit{Step 1: express evolution through transition moments.} For a Markov process, the conditional distribution of a future state depends on the past only through the present state. Under conditions that justify the moment expansion, the forward equation can formally be written as

\begin{equation}
\frac{\partial p}{\partial t}
=\sum_{n=1}^{\infty}\frac{(-1)^n}{n!}
\frac{\partial^n}{\partial v^n}\big[D_n(v)p\big].
\end{equation}
The coefficients in this expansion are defined by
\begin{equation}
D_n(v):=\lim_{\Delta t\downarrow0}\frac{1}{\Delta t}
\E\big[(v(t+\Delta t)-v(t))^n\mid v(t)=v\big].
\end{equation}

These are conditional raw moments of the increment per unit time, not moments centered about its conditional mean. The Kramers--Moyal series is used here as a formal tool, rather than asserted as a convergent representation for every Markov process.

\textit{Step 2: compute the infinitesimal moments.} Over a short interval, freezing the coefficients at the current state gives the leading approximation

\begin{align}
v(t+\Delta t)-v(t)
&=\int_t^{t+\Delta t}f(v(s))\,\dd s
+\int_t^{t+\Delta t}\sigma(v(s))\,\dd W(s)\notag\\
&\approx f(v)\Delta t+\sigma(v)\Delta W,
\end{align}
where $\Delta W\sim\mathcal{N}(0,\Delta t)$ independently of the current state. All expectations below are conditional on $v(t)=v$. Since $\E[\Delta W]=0$,
\begin{equation}
\begin{aligned}
D_1(v) &= \lim_{\Delta t\downarrow0}\frac{f(v)\Delta t}{\Delta t} \\
&= f(v).
\end{aligned}
\end{equation}
For the second moment, the frozen-coefficient approximation gives
\begin{align}
\E\big[(v(t+\Delta t)-v(t))^2\mid v(t)=v\big]
&\approx f(v)^2\Delta t^2
+2f(v)\sigma(v)\Delta t\,\E[\Delta W]\notag\\
&\quad+\sigma(v)^2\E[(\Delta W)^2]\notag\\
&=f(v)^2\Delta t^2+\sigma(v)^2\Delta t,
\end{align}
because
\[
\begin{aligned}
\E[(\Delta W)^2] &= \var(\Delta W) \\
&= \Delta t.
\end{aligned}
\]
Dividing by $\Delta t$ and taking the limit yields
\begin{equation}
D_2(v)=\sigma(v)^2.
\end{equation}

For $n\ge3$, the higher increment moments are $o(\Delta t)$ under suitable moment bounds, so $D_n(v)=0$. The Gaussian scaling underlying this fact is
\[
\E[|\Delta W|^n]=O((\Delta t)^{n/2});
\]
for example,
\[
\begin{gathered}
\E[(\Delta W)^3]=0 \\
\E[(\Delta W)^4]=3\Delta t^2.
\end{gathered}
\]

The neglected terms must be controlled at the level of conditional moments to turn this heuristic into a proof.

\textit{Step 3: retain the first two terms.} Substitution into the formal expansion gives

\begin{align}
\frac{\partial p}{\partial t}
&=-\frac{\partial}{\partial v}(D_1p)
+\frac12\frac{\partial^2}{\partial v^2}(D_2p)\notag\\
&=-\frac{\partial}{\partial v}\big[f(v)p\big]
+\frac12\frac{\partial^2}{\partial v^2}\big[\sigma(v)^2p\big],
\end{align}
which is the Fokker--Planck equation. A rigorous weak derivation uses It\^o's formula on smooth compactly supported test functions in place of the formal series.
\end{proof}

\subsection{Probability Currents and Stationary Densities}

\begin{remark}[A Continuity Equation with Diffusive Current]
The Fokker--Planck equation is a continuity equation of the type derived in Example~\ref{ex:continuity-equation}:
\begin{equation}
\frac{\partial p}{\partial t}+\frac{\partial J}{\partial v}=0,
\qquad J(v,t)=f(v)p(v,t)-\frac12\frac{\partial}{\partial v}\big[\sigma(v)^2p(v,t)\big].
\end{equation}

The first term in the probability current is deterministic transport by the drift; the second is a diffusive flux. For constant $\sigma$ and zero drift, this reduces to the heat equation with $D=\sigma^2/2$, studied in Section~\ref{sec:pde-methods}. Thus random trajectories produce the same diffusion equation for a probability density that describes heat conduction or molecular spreading in other contexts. Conservation of total probability requires suitable boundary conditions, such as zero flux at finite boundaries or vanishing current at infinity.
\end{remark}

\begin{example}[Noise in a Genetic Circuit or a Neuron]
For the Ornstein--Uhlenbeck process of Example~\ref{ex:ornstein-uhlenbeck}, the density satisfies
\begin{equation}
\frac{\partial p}{\partial t}
=\gamma\frac{\partial}{\partial v}(vp)
+\frac{\sigma^2}{2}\frac{\partial^2p}{\partial v^2}.
\end{equation}
On the whole real line, the normalized stationary density is
\begin{equation}
p_\infty(v)=\sqrt{\frac{\gamma}{\pi\sigma^2}}
\exp\left(-\frac{\gamma v^2}{\sigma^2}\right),
\qquad \var(v)=\frac{\sigma^2}{2\gamma}.
\end{equation}
Indeed,
\[
p_\infty'=-(2\gamma v/\sigma^2)p_\infty,
\]
so the stationary current
\[
J_\infty=-\gamma v p_\infty-(\sigma^2/2)p_\infty'
\]
vanishes identically. The restoring drift narrows the distribution, while noise broadens it. This balance models fluctuations of a subthreshold membrane potential around its equilibrium, or small fluctuations of gene expression around a stable mean. In the latter interpretation, $v$ is a deviation from the mean, not an unrestricted physical molecular count; threshold/reset dynamics or positivity constraints require additional modeling.
\end{example}
\clearpage
\chapter*{References and Recommended Reading}
\addcontentsline{toc}{chapter}{References and Recommended Reading}
\markboth{References and Recommended Reading}{References and Recommended Reading}

The theoretical framework and operational tools developed in this compendium build upon the foundations of real and complex analysis and extend them through linear and functional analysis, optimization, differential equations, Fourier methods, probability, stochastic processes, and mathematical methods for machine learning. The following works provide the principal references for the material presented here and complementary routes for further study.

\begin{enumerate}[label={\textup{[\arabic*]}},ref=\arabic*,leftmargin=*,itemsep=0.5em]
\item \label{ref:bramanti}
Bramanti, M., Pagani, C.~D., and Salsa, S. (2008--2009). \emph{Analisi matematica 1} (2008) and \emph{Analisi matematica 2} (2009). Zanichelli.

\item \label{ref:rudin-principles}
Rudin, W. (1976). \emph{Principles of Mathematical Analysis} (3rd ed.). McGraw--Hill.

\item \label{ref:apostol}
Apostol, T.~M. (1974). \emph{Mathematical Analysis} (2nd ed.). Addison--Wesley.

\item \label{ref:stein-complex}
Stein, E.~M., and Shakarchi, R. (2003). \emph{Complex Analysis} (Princeton Lectures in Analysis, Vol.~2). Princeton University Press.

\item \label{ref:strang}
Strang, G. (2016). \emph{Introduction to Linear Algebra} (5th ed.). Wellesley--Cambridge Press.

\item \label{ref:axler}
Axler, S. (2015). \emph{Linear Algebra Done Right} (3rd ed.). Springer.

\item \label{ref:boyd-convex}
Boyd, S., and Vandenberghe, L. (2004). \emph{Convex Optimization}. Cambridge University Press. \href{https://doi.org/10.1017/CBO9780511804441}{doi:10.1017/CBO9780511804441}.

\item \label{ref:nocedal}
Nocedal, J., and Wright, S.~J. (2006). \emph{Numerical Optimization} (2nd ed.). Springer.

\item \label{ref:brezis}
Brezis, H. (2011). \emph{Functional Analysis, Sobolev Spaces and Partial Differential Equations}. Springer. \href{https://doi.org/10.1007/978-0-387-70914-7}{doi:10.1007/978-0-387-70914-7}.

\item \label{ref:evans}
Evans, L.~C. (2010). \emph{Partial Differential Equations} (2nd ed.). American Mathematical Society. \href{https://doi.org/10.1090/gsm/019}{doi:10.1090/gsm/019}.

\item \label{ref:stein-fourier}
Stein, E.~M., and Shakarchi, R. (2003). \emph{Fourier Analysis: An Introduction}. Princeton University Press.

\item \label{ref:kolmogorov}
Kolmogorov, A.~N., and Fomin, S.~V. (1975). \emph{Introductory Real Analysis}. Dover Publications.

\item \label{ref:hirsch-smale}
Hirsch, M.~W., Smale, S., and Devaney, R.~L. (2012). \emph{Differential Equations, Dynamical Systems, and an Introduction to Chaos} (3rd ed.). Academic Press.

\item \label{ref:strogatz}
Strogatz, S.~H. (2015). \emph{Nonlinear Dynamics and Chaos} (2nd ed.). Westview Press.

\item \label{ref:billingsley}
Billingsley, P. (1995). \emph{Probability and Measure} (3rd ed.). John Wiley \& Sons.

\item \label{ref:shiryaev}
Shiryaev, A.~N. (1996). \emph{Probability} (2nd ed.). Springer. \href{https://doi.org/10.1007/978-1-4757-2539-1}{doi:10.1007/978-1-4757-2539-1}.

\item \label{ref:oksendal}
Øksendal, B. (2003). \emph{Stochastic Differential Equations: An Introduction with Applications} (6th ed.). Springer. \href{https://doi.org/10.1007/978-3-642-14394-6}{doi:10.1007/978-3-642-14394-6}.

\item \label{ref:gardiner}
Gardiner, C. (2009). \emph{Stochastic Methods: A Handbook for the Natural and Social Sciences} (4th ed.). Springer.

\item \label{ref:bishop}
Bishop, C.~M. (2006). \emph{Pattern Recognition and Machine Learning}. Springer.

\item \label{ref:goodfellow}
Goodfellow, I., Bengio, Y., and Courville, A. (2016). \emph{Deep Learning}. MIT Press.

\item \label{ref:murphy}
Murphy, K.~P. (2022). \emph{Probabilistic Machine Learning: An Introduction}. MIT Press.
\end{enumerate}
\end{document}
